\documentclass[11pt,reqno]{amsart}
\usepackage{iftex}
\ifPDFTeX
  \usepackage[T1]{fontenc}
  \usepackage[utf8]{inputenc}
  \usepackage{lmodern}
\else
  \usepackage[no-math]{fontspec}
\fi
\usepackage{amsmath}
\usepackage{amssymb}
\usepackage{amsfonts}
\usepackage{mathrsfs}
\usepackage{amsthm}
\usepackage{mathtools}
\usepackage{geometry}
\usepackage{enumitem}
\usepackage{longtable}
\usepackage{booktabs}
\usepackage{array}
\usepackage{microtype}
\usepackage{graphicx}
\graphicspath{{../figures/}{figures/}}
\usepackage{listings}
\lstset{basicstyle=\ttfamily\scriptsize,breaklines=true,columns=fullflexible,keepspaces=true,
  frame=single,framerule=0.3pt,showstringspaces=false,keywordstyle=\color{NavyBlue},
  commentstyle=\color{OliveGreen},stringstyle=\color{BrickRed},upquote=true}
\lstdefinestyle{out}{basicstyle=\ttfamily\scriptsize,language={},keywordstyle={},commentstyle={},stringstyle={}}
\usepackage{float}
\usepackage[dvipsnames]{xcolor}
\usepackage[pagebackref=true]{hyperref}
\usepackage{bookmark}
\usepackage{pdfcomment}

% Hover previews for internal references (no layout change).
\makeatletter
\let\CG@origref\ref
\let\CG@origeqref\eqref
\newcommand{\CG@gettip}[1]{\@ifundefined{CGTip@#1}{Reference: #1}{\csname CGTip@#1\endcsname}}
\DeclareRobustCommand{\CGref}[1]{\pdftooltip{\CG@origref{#1}}{\CG@gettip{#1}}}
\DeclareRobustCommand{\CGeqref}[1]{\pdftooltip{\CG@origeqref{#1}}{\CG@gettip{#1}}}
\makeatother
% Tooltip texts shown when hovering over internal references.
\makeatletter
\expandafter\gdef\csname CGTip@sec:notation\endcsname{Section : Notation (p. 5)}
\expandafter\gdef\csname CGTip@sec:introduction\endcsname{Section 1: Introduction (p. 10)}
\expandafter\gdef\csname CGTip@subsec:background\endcsname{Section 1.1: Background (p. 10)}
\expandafter\gdef\csname CGTip@eq:mertens-def\endcsname{Eq. (1.1), p. 10: M(x):=sum\_n<= xmu(n).}
\expandafter\gdef\csname CGTip@subsec:midpoint-system\endcsname{Section 1.2: The midpoint RVM system (p. 10)}
\expandafter\gdef\csname CGTip@eq:intro-clock\endcsname{Eq. (1.2), p. 10: N\_RVM(T)=(T)/(2pi)log(T)/(2pi e)+(7)/(8),}
\expandafter\gdef\csname CGTip@eq:intro-lambert\endcsname{Eq. (1.3), p. 10: tau\_k=(2pi(k-11/8))/(W\_0((k-11/8)/e)) .}
\expandafter\gdef\csname CGTip@eq:intro-zeta-rvm\endcsname{Eq. (1.4), p. 10: P\_RVM\textasciicircum{}crit(s)=prod\_k>=1(1+(s(s-1))/((1)/(4)+tau\_k\textasciicircum{}2)), zeta\_RVM(s)=(P\_RVM\textasciicircum{}crit(s))/(G(s)), G(s)=s(s-1)Gamma(s/2)pi\textasciicircum{}-s/2.}
\expandafter\gdef\csname CGTip@eq:intro-pi\endcsname{Eq. (1.5), p. 11: pi\_RVM(x)=sum\_1<= m<=log\_2x(mu(m))/(m) J\_RVM\textasciicircum{}*(x\textasciicircum{}1/m), J\_RVM\textasciicircum{}*(x)=sum\_1<= m<=log\_2x(1)/(m) pi\_RVM(x\textasciicircum{}1/m).}
\expandafter\gdef\csname CGTip@eq:intro-real-primes\endcsname{Eq. (1.6), p. 11: p\_j\textasciicircum{}RVM:=pi\_RVM\textasciicircum{}-1(j), pi\_RVM(p\_j\textasciicircum{}RVM)=j,}
\expandafter\gdef\csname CGTip@eq:intro-transition\endcsname{Eq. (1.7), p. 11: zeta\_RVM(s)=varpi\_RVM(1+w(s))(1+O(1/|t|)), w(s)=exp-2pi i N\_RVM(-i(s-(1)/(2))),}
\expandafter\gdef\csname CGTip@eq:intro-Mbound\endcsname{Eq. (1.8), p. 12: M\_RVM\textasciicircum{}crit(x)=L+O\_eps(x\textasciicircum{}1/2+eps)}
\expandafter\gdef\csname CGTip@eq:intro-muh\endcsname{Eq. (1.9), p. 12: mu\_h(n)=mu(n)\textasciicircum{}2varsigma\_n, varsigma\_n in -1,+1, M\_h(n):=sum\_m<= nmu\_h(m),}
\expandafter\gdef\csname CGTip@eq:intro-quantizer\endcsname{Eq. (1.10), p. 12: varsigma\_n= +1, M\_h(n-1)<M(n),; -1, M\_h(n-1)>=M(n),}
\expandafter\gdef\csname CGTip@subsec:core\endcsname{Section 1.3: The core of the method (p. 12)}
\expandafter\gdef\csname CGTip@eq:intro-family\endcsname{Eq. (1.11), p. 12: gamma\_k(lambda)=(1-lambda)gamma\_k+lambdatau\_k, zeta\_lambda(s)=(1)/(G(s))prod\_k>=1(1+(s(s-1))/((1)/(4)+gamma\_k(lambda)\textasciicircum{}2)),}
\expandafter\gdef\csname CGTip@eq:intro-hamburger\endcsname{Eq. (1.12), p. 12: zeta\_Z1=zeta <=> (D): zeta\_Z1 is an absolutely convergent Dirichlet series in Re s>1 .}
\expandafter\gdef\csname CGTip@eq:intro-varpi-Z1\endcsname{Eq. (1.13), p. 13: varpi\_Z1=lim\_sigma->+infzeta\_Z1(sigma)=varpi\_RVMlim\_K->infprod\_k<= K((1)/(4)+tau\_k\textasciicircum{}2)/((1)/(4)+gamma\_k\textasciicircum{}2) >=1,}
\expandafter\gdef\csname CGTip@eq:intro-step3\endcsname{Eq. (1.14), p. 13: zeta\_Z1=zeta => M(x)=O\_eps(x\textasciicircum{}1/2+eps) (eps>0),}
\expandafter\gdef\csname CGTip@eq:intro-prime-chain\endcsname{Eq. (1.15), p. 13: tau\_k->zeta\_RVM-> J\_RVM\_RVM->(dpi\_RVM, p\_j\textasciicircum{}RVM, c\_n\textasciicircum{}RVM),}
\expandafter\gdef\csname CGTip@eq:intro-M-chain\endcsname{Eq. (1.16), p. 13: F\_RVM=zeta\_RVM\textasciicircum{}-1-> M\_RVM\textasciicircum{}crit->(M,mu\_h,M\_h),}
\expandafter\gdef\csname CGTip@eq:intro-family-chain\endcsname{Eq. (1.17), p. 13: zeta\_RVM=zeta\_1->zeta\_0=zeta\_Z1=?zeta => M(x)=O\_eps(x\textasciicircum{}1/2+eps).}
\expandafter\gdef\csname CGTip@subsec:status\endcsname{Section 1.4: Originality, scope and status (p. 13)}
\expandafter\gdef\csname CGTip@subsec:organization\endcsname{Section 1.5: Organization (p. 16)}
\expandafter\gdef\csname CGTip@sec:clock\endcsname{Section 2: The midpoint Riemann--von Mangoldt clock (p. 17)}
\expandafter\gdef\csname CGTip@eq:zero-count\endcsname{Eq. (2.1), p. 17: N\_zeta(T):=no.rho: 0<Imrho<= T=no.k:gamma\_k<= T.}
\expandafter\gdef\csname CGTip@eq:RvM-full\endcsname{Eq. (2.2), p. 17: N\_zeta(T)=(T)/(2pi)log(T)/(2pi e)+(7)/(8)+S(T)+O((1)/(T)), S(T):=(1)/(pi)((1)/(2)+iT)=O(log T),}
\expandafter\gdef\csname CGTip@eq:clock\endcsname{Eq. (2.3), p. 17: N\_RVM(T):=(T)/(2pi)log(T)/(2pi e)+(7)/(8) .}
\expandafter\gdef\csname CGTip@eq:midpoint-condition\endcsname{Eq. (2.4), p. 17: N\_RVM(tau\_k)=k-(1)/(2), k>=1,}
\expandafter\gdef\csname CGTip@prop:lambert\endcsname{Proposition 2.1 (Lambert inversion) (p. 17): For every k>=1, tau\_k=(2pi(k-11/8))/(W\_0((k-11/8)/e)), The quotient is positive also for k=1.}
\expandafter\gdef\csname CGTip@eq:lambert\endcsname{Eq. (2.5), p. 17: tau\_k=(2pi(k-11/8))/(W\_0((k-11/8)/e)),}
\expandafter\gdef\csname CGTip@eq:lambert-proof-y\endcsname{Eq. (2.6), p. 17: y(log y-1)=k-(11)/(8).}
\expandafter\gdef\csname CGTip@eq:clock-derivative\endcsname{Eq. (2.7), p. 18: N\_RVM'(T)=(1)/(2pi)log(T)/(2pi),}
\expandafter\gdef\csname CGTip@eq:clock-growth\endcsname{Eq. (2.8), p. 18: tau\_k (k)/(log k) (k>=2).}
\expandafter\gdef\csname CGTip@eq:RvM-formula\endcsname{Eq. (2.9), p. 18: N\_zeta(T)=N\_RVM(T)+S(T)+O(T\textasciicircum{}-1), S(T)=O(log T) (T>=2).}
\expandafter\gdef\csname CGTip@eq:symmetric-count\endcsname{Eq. (2.10), p. 18: N\_zeta\textasciicircum{}*(T):=(1)/(2)(N\_zeta(T+0)+N\_zeta(T-0))}
\expandafter\gdef\csname CGTip@eq:eta-def\endcsname{Eq. (2.11), p. 18: eta\_k:=tau\_k-gamma\_k .}
\expandafter\gdef\csname CGTip@lem:bounded-displacement\endcsname{Lemma 2.2 (Bounded displacement) (p. 18): There is a constant C\_eta such that |eta\_k|<= C\_eta (k>=1).}
\expandafter\gdef\csname CGTip@eq:bounded-displacement\endcsname{Eq. (2.12), p. 18: |eta\_k|<= C\_eta (k>=1).}
\expandafter\gdef\csname CGTip@eq:block-count\endcsname{Eq. (2.13), p. 18: |N\_zeta\textasciicircum{}*(gamma\_k)-(k-(1)/(2))|<=(m(gamma\_k))/(2) .}
\expandafter\gdef\csname CGTip@eq:clock-discrepancy\endcsname{Eq. (2.14), p. 18: |N\_RVM(gamma\_k)-N\_RVM(tau\_k)|=O(loggamma\_k).}
\expandafter\gdef\csname CGTip@sec:zeta-rvm\endcsname{Section 3: The pure critical RVM zeta function (p. 18)}
\expandafter\gdef\csname CGTip@eq:q-def\endcsname{Eq. (3.1), p. 18: q(s):=s(s-1),}
\expandafter\gdef\csname CGTip@eq:Pcrit\endcsname{Eq. (3.2), p. 18: P\_RVM\textasciicircum{}crit(s):=prod\_k>=1(1+(q(s))/((1)/(4)+tau\_k\textasciicircum{}2)),}
\expandafter\gdef\csname CGTip@eq:Gdef\endcsname{Eq. (3.3), p. 18: G(s):=s(s-1)Gamma(s/2)pi\textasciicircum{}-s/2,}
\expandafter\gdef\csname CGTip@eq:zetaRVM\endcsname{Eq. (3.4), p. 18: xi\_RVM(s):=(1)/(2)P\_RVM\textasciicircum{}crit(s), zeta\_RVM(s):=(P\_RVM\textasciicircum{}crit(s))/(G(s)).}
\expandafter\gdef\csname CGTip@thm:canonical\endcsname{Theorem 3.1 (Canonical characterization of zeta\_RVM) (p. 18): The product CGeqrefeq:Pcrit converges absolutely and locally uniformly in s in C. It defines an entire function with P\_RVM\textasciicircum{}crit(1-s)=P\_RVM\textasciicircum{}crit(s), P\_RVM\textasciicircum{}crit(conj s)=conj P\_RVM\textasciicircum{}crit(s), P\_RVM\textasciicircum{}crit(0)=P\_RVM\textasciicircum{}crit(1)=1 . Its zeros are exactly the simple zeros rho\_k\textasciicircum{}RVM=(1)/(2)+itau\_k, conj rho\_k\textasciicircum{}RV...}
\expandafter\gdef\csname CGTip@eq:P-symmetry\endcsname{Eq. (3.5), p. 18: P\_RVM\textasciicircum{}crit(1-s)=P\_RVM\textasciicircum{}crit(s), P\_RVM\textasciicircum{}crit(conj s)=conj P\_RVM\textasciicircum{}crit(s), P\_RVM\textasciicircum{}crit(0)=P\_RVM\textasciicircum{}crit(1)=1 .}
\expandafter\gdef\csname CGTip@eq:synthetic-divisor\endcsname{Eq. (3.6), p. 18: rho\_k\textasciicircum{}RVM=(1)/(2)+itau\_k, conj rho\_k\textasciicircum{}RVM=(1)/(2)-itau\_k, k>=1 .}
\expandafter\gdef\csname CGTip@eq:canonical-normal\endcsname{Eq. (3.7), p. 19: sup\_s in K|(q(s))/((1)/(4)+tau\_k\textasciicircum{}2)|<<\_Ktau\_k\textasciicircum{}-2,}
\expandafter\gdef\csname CGTip@eq:factor-zero\endcsname{Eq. (3.8), p. 19: (s-(1)/(2))\textasciicircum{}2+tau\_k\textasciicircum{}2=0 .}
\expandafter\gdef\csname CGTip@sec:exp-sums\endcsname{Section 4: RVM exponential sums, centroids, and the height transform (p. 19)}
\expandafter\gdef\csname CGTip@eq:S-RVM\endcsname{Eq. (4.1), p. 19: S\_RVM(N;u):=sum\_k<= Ne\textasciicircum{}iutau\_k, C\_RVM\textasciicircum{}(N;u):=(S\_RVM(N;u))/(N),}
\expandafter\gdef\csname CGTip@eq:path-centroid\endcsname{Eq. (4.2), p. 19: C\_RVM\textasciicircum{}(N;u):=(1)/(N)sum\_n<= NS\_RVM(n;u) =sum\_k<= N(N-k+1)/(N)e\textasciicircum{}iutau\_k.}
\expandafter\gdef\csname CGTip@thm:exp-sum\endcsname{Theorem 4.1 (Uniform logarithmic exponential-sum bound) (p. 19): Let 0<u\_1<= u\_2<inf. Then sup\_u in [u\_1,u\_2]|S\_RVM(N;u)|<<\_u\_1,u\_21+log N .}
\expandafter\gdef\csname CGTip@eq:exp-sum-bound\endcsname{Eq. (4.3), p. 19: sup\_u in [u\_1,u\_2]|S\_RVM(N;u)|<<\_u\_1,u\_21+log N .}
\expandafter\gdef\csname CGTip@eq:inverse-clock\endcsname{Eq. (4.4), p. 19: T'(y)=(2pi)/(log(T(y)/2pi)), T''(y)=-(4pi\textasciicircum{}2)/(T(y)log\textasciicircum{}3(T(y)/2pi)) .}
\expandafter\gdef\csname CGTip@eq:rvm-gap-bounds\endcsname{Eq. (4.5), p. 19: (2pi)/(\_k+1)<=tau\_k+1-tau\_k<=(2pi)/(\_k),}
\expandafter\gdef\csname CGTip@eq:rvm-second-difference\endcsname{Eq. (4.6), p. 19: |tau\_k+2-2tau\_k+1+tau\_k|<=(4pi\textasciicircum{}2)/(tau\_k\_k\textasciicircum{}3).}
\expandafter\gdef\csname CGTip@eq:tau-ell-k\endcsname{Eq. (4.7), p. 19: tau\_k\_k k (k>=2).}
\expandafter\gdef\csname CGTip@eq:Ek-identity\endcsname{Eq. (4.8), p. 19: E\_k=b\_k\textasciicircum{} sbp(E\_k+1-E\_k).}
\expandafter\gdef\csname CGTip@eq:summation-by-parts\endcsname{Eq. (4.9), p. 20: sum\_k=k\_0\textasciicircum{}NE\_k=b\_N\textasciicircum{} sbpE\_N+1-b\_k\_0\textasciicircum{} sbpE\_k\_0+sum\_k=k\_0+1\textasciicircum{}N(b\_k-1\textasciicircum{} sbp-b\_k\textasciicircum{} sbp)E\_k .}
\expandafter\gdef\csname CGTip@eq:kappa-bounds\endcsname{Eq. (4.10), p. 20: |b\_k\textasciicircum{} sbp|<=(pi)/(2Delta\_k)<=(\_k+1)/(4u\_1), |b\_k-1\textasciicircum{} sbp-b\_k\textasciicircum{} sbp|<=|Delta\_k-Delta\_k-1| |b\_k-1\textasciicircum{} sbp||b\_k\textasciicircum{} sbp| <=(pi\textasciicircum{}2u\_2)/(4u\_1\textasciicircum{}2) (\_k+1\textasciicircum{}2)/(tau\_k-1\_k-1\textasciicircum{}3)<<\_u\_1,u\_2(1)/(tau\_k\_k)<<(1)/(k),}
\expandafter\gdef\csname CGTip@cor:point-centroid\endcsname{Corollary 4.2 (Point centroid) (p. 20): Uniformly for u in [u\_1,u\_2](0,inf), C\_RVM\textasciicircum{}(N;u)=O\_u\_1,u\_2 ((1+log N)/(N))->0 . Theorem CGrefthm:exp-sum alone does not imply convergence of the path centroid CGeqrefeq:path-centroid.}
\expandafter\gdef\csname CGTip@eq:point-centroid\endcsname{Eq. (4.11), p. 20: C\_RVM\textasciicircum{}(N;u)=O\_u\_1,u\_2 ((1+log N)/(N))->0 .}
\expandafter\gdef\csname CGTip@eq:n-RVM\endcsname{Eq. (4.12), p. 20: n\_RVM(H):=no.k:tau\_k<= H.}
\expandafter\gdef\csname CGTip@eq:MH-RVM\endcsname{Eq. (4.13), p. 20: M\_H\textasciicircum{}RVM(x):=(1)/(n\_RVM(H))sum\_tau\_k<= Hx\textasciicircum{}itau\_k,}
\expandafter\gdef\csname CGTip@eq:MH-genuine\endcsname{Eq. (4.14), p. 20: M\_H(x):=(1)/(N\_zeta(H))sum\_gamma\_k<= Hx\textasciicircum{}igamma\_k.}
\expandafter\gdef\csname CGTip@cor:MH-decay\endcsname{Corollary 4.3 (Synthetic height-transform decay) (p. 20): For every fixed x>1, with N=n\_RVM(H), M\_H\textasciicircum{}RVM(x)=O\_x ((1+log N)/(N))->0 (H->inf).}
\expandafter\gdef\csname CGTip@eq:MH-decay\endcsname{Eq. (4.15), p. 20: M\_H\textasciicircum{}RVM(x)=O\_x ((1+log N)/(N))->0 (H->inf).}
\expandafter\gdef\csname CGTip@rem:MH-not-mertens\endcsname{Remark 4.4 (p. 20): Neither CGeqrefeq:MH-RVM nor CGeqrefeq:MH-genuine is the Mertens function; both are normalized phase averages.}
\expandafter\gdef\csname CGTip@sec:explicit-formula\endcsname{Section 5: The RVM explicit formula, J\_RVMJ\_RVM, and pi\_RVMpi\_RVM (p. 20)}
\expandafter\gdef\csname CGTip@eq:Ei-def\endcsname{Eq. (5.1), p. 20: Ei(w):=int\_-inf\textasciicircum{}a(e\textasciicircum{}v+ib)/(v+ib) dv, Ei(conj w):=conj Ei(w) .}
\expandafter\gdef\csname CGTip@eq:J-RVM\endcsname{Eq. (5.2), p. 20: J\_RVM(x):=li(x)-2Resum\_k>=1Ei (((1)/(2)+itau\_k)log x)-log2+int\_x\textasciicircum{}inf(dy)/(y(y\textasciicircum{}2-1)log y),}
\expandafter\gdef\csname CGTip@lem:Ei\endcsname{Lemma 5.1 (Asymptotics of Ei off the real axis) (p. 20): For w=a+ib, b>0, Ei(w)=(e\textasciicircum{}w)/(w)+varrho(w), |varrho(w)|<=(e\textasciicircum{}a)/(b\textasciicircum{}2).}
\expandafter\gdef\csname CGTip@eq:Ei-remainder\endcsname{Eq. (5.3), p. 20: Ei(w)=(e\textasciicircum{}w)/(w)+varrho(w), |varrho(w)|<=(e\textasciicircum{}a)/(b\textasciicircum{}2).}
\expandafter\gdef\csname CGTip@thm:J-exists\endcsname{Theorem 5.2 (Compact-uniform existence of the RVM zero sum) (p. 21): The series sum\_kReEi(((1)/(2)+itau\_k)u) converges uniformly for u in every compact interval [u\_1,u\_2](0,inf). Hence J\_RVM is a continuous function on (1,inf).}
\expandafter\gdef\csname CGTip@eq:Ei-asymptotic\endcsname{Eq. (5.4), p. 21: ReEi(((1)/(2)+itau\_k)u)=(e\textasciicircum{}u/2)/(u) (sin(utau\_k))/(tau\_k)+O\_u\_1,u\_2(tau\_k\textasciicircum{}-2),}
\expandafter\gdef\csname CGTip@eq:abel-weights\endcsname{Eq. (5.5), p. 21: (1)/(tau\_K)-(1)/(tau\_K+1)<=(tau\_K+1-tau\_K)/(tau\_K\textasciicircum{}2)<=(2pi)/(\_Ktau\_K\textasciicircum{}2)<<(log K)/(K\textasciicircum{}2)}
\expandafter\gdef\csname CGTip@eq:J-star\endcsname{Eq. (5.6), p. 21: J\_RVM\textasciicircum{}*(x):= J\_RVM(x), x>=2,; 0, 0<x<2,}
\expandafter\gdef\csname CGTip@eq:pi-RVM\endcsname{Eq. (5.7), p. 21: pi\_RVM(x):=sum\_1<= m<=log\_2x(mu(m))/(m) J\_RVM\textasciicircum{}*(x\textasciicircum{}1/m).}
\expandafter\gdef\csname CGTip@prop:finite-inversion\endcsname{Proposition 5.3 (Exact finite prime-power inversion) (p. 21): For every x>=2, J\_RVM\textasciicircum{}*(x)=sum\_1<= m<=log\_2x(1)/(m) pi\_RVM(x\textasciicircum{}1/m).}
\expandafter\gdef\csname CGTip@eq:J-pi-inverse\endcsname{Eq. (5.8), p. 21: J\_RVM\textasciicircum{}*(x)=sum\_1<= m<=log\_2x(1)/(m) pi\_RVM(x\textasciicircum{}1/m).}
\expandafter\gdef\csname CGTip@eq:mobius-coefficient\endcsname{Eq. (5.9), p. 21: (1)/(r)sum\_d| rmu(d),}
\expandafter\gdef\csname CGTip@subsec:square-root\endcsname{Section 5.1: The square-root estimate (p. 21)}
\expandafter\gdef\csname CGTip@eq:A-RVM\endcsname{Eq. (5.10), p. 21: A\_RVM(u):=sum\_k>=1(e\textasciicircum{}iutau\_k)/(tau\_k),}
\expandafter\gdef\csname CGTip@lem:A-RVM\endcsname{Lemma 5.4 (The weighted clock sum) (p. 21): For every u>=log2 the series CGeqrefeq:A-RVM converges, and |A\_RVM(u)|<=(log\textasciicircum{}2(2+2u))/(4pi)+(log(2+2u))/(2pi)+C\_A, with an absolute constant C\_A.}
\expandafter\gdef\csname CGTip@eq:A-RVM-bound\endcsname{Eq. (5.11), p. 21: |A\_RVM(u)|<=(log\textasciicircum{}2(2+2u))/(4pi)+(log(2+2u))/(2pi)+C\_A,}
\expandafter\gdef\csname CGTip@eq:A-head\endcsname{Eq. (5.12), p. 22: sum\_tau\_k<= Y\_0(1)/(tau\_k)=(n\_RVM(Y\_0))/(Y\_0)+int\_tau\_1\textasciicircum{}Y\_0(n\_RVM(t))/(t\textasciicircum{}2) dt<=((Y\_0))/(2pi)+int\_tau\_1\textasciicircum{}Y\_0((t))/(2pi t) dt<=((Y\_0))/(2pi)+((Y\_0)\textasciicircum{}2)/(4pi),}
\expandafter\gdef\csname CGTip@eq:phi-second\endcsname{Eq. (5.13), p. 22: phi''(y)=-(2pi u)/(T (T)\textasciicircum{}3), T=T(y-(1)/(2)).}
\expandafter\gdef\csname CGTip@eq:block-vdc\endcsname{Eq. (5.14), p. 22: |sum\_a<k<= be\textasciicircum{}iutau\_k|<= A\_ vdC(256((Y\_j(Y\_j+1))/(2pi)+1)\_j\textasciicircum{}1/2+\_j\textasciicircum{}-1/2)=:P\_j,}
\expandafter\gdef\csname CGTip@eq:block-weighted\endcsname{Eq. (5.15), p. 22: (P\_j)/(Y\_j)<= A\_ vdC((16)/(2pi)((pi u)/(Y\_j(Y\_j+1)))\textasciicircum{}1/2+(16)/(Y\_j)+(((Y\_j+1)\textasciicircum{}3)/(pi uY\_j))\textasciicircum{}1/2).}
\expandafter\gdef\csname CGTip@thm:Lambert-estimate\endcsname{Theorem 5.5 (Square-root estimate for J\_RVM and pi\_RVM) (p. 22): For x>=8, J\_RVM(x)-li(x)=O (( x(loglog x)\textasciicircum{}2)/(log x)), pi\_RVM(x)-li(x)=O (( x(loglog x)\textasciicircum{}2)/(log x)), with absolute implied constants. No hypothesis on zeta is used.}
\expandafter\gdef\csname CGTip@eq:J-baseline\endcsname{Eq. (5.16), p. 22: J\_RVM(x)-li(x)=O (( x(loglog x)\textasciicircum{}2)/(log x)),}
\expandafter\gdef\csname CGTip@eq:pi-baseline\endcsname{Eq. (5.17), p. 22: pi\_RVM(x)-li(x)=O (( x(loglog x)\textasciicircum{}2)/(log x)),}
\expandafter\gdef\csname CGTip@eq:Ei-split\endcsname{Eq. (5.18), p. 22: Ei(rho\_k\textasciicircum{}RVMu)=(e\textasciicircum{}rho\_k\textasciicircum{}RVMu)/(rho\_k\textasciicircum{}RVMu)+varrho\_k, |varrho\_k|<=(e\textasciicircum{}u/2)/(tau\_k\textasciicircum{}2u\textasciicircum{}2).}
\expandafter\gdef\csname CGTip@eq:Ei-sum-A\endcsname{Eq. (5.19), p. 23: sum\_k>=1Ei(rho\_k\textasciicircum{}RVMu)=(e\textasciicircum{}u/2)/(u)(-i A\_RVM(u)+sum\_kr\_ke\textasciicircum{}iutau\_k)+sum\_kvarrho\_k =(e\textasciicircum{}u/2)/(u)(-i A\_RVM(u)+O(1)).}
\expandafter\gdef\csname CGTip@eq:pi-minus-li\endcsname{Eq. (5.20), p. 23: pi\_RVM(x)-li(x)=(J\_RVM(x)-li(x))+sum\_2<= m<=log\_2x(mu(m))/(m) J\_RVM\textasciicircum{}*(x\textasciicircum{}1/m).}
\expandafter\gdef\csname CGTip@cor:three-limits\endcsname{Corollary 5.6 (The three limits) (p. 23): As x->inf, (li(x))/(pi\_RVM(x))->1, (pi(x))/(li(x))->1, (pi(x))/(pi\_RVM(x))->1 . More precisely, for some absolute c\_0>0, pi(x)-pi\_RVM(x)=O (x e\textasciicircum{}-c\_0sqrt(log x)).}
\expandafter\gdef\csname CGTip@eq:three-limits\endcsname{Eq. (5.21), p. 23: (li(x))/(pi\_RVM(x))->1, (pi(x))/(li(x))->1, (pi(x))/(pi\_RVM(x))->1 .}
\expandafter\gdef\csname CGTip@eq:pi-minus-piRVM\endcsname{Eq. (5.22), p. 23: pi(x)-pi\_RVM(x)=O (x e\textasciicircum{}-c\_0sqrt(log x)).}
\expandafter\gdef\csname CGTip@sec:integer-weights\endcsname{Section 6: Integer weights and the continuous-to-atomic problem (p. 23)}
\expandafter\gdef\csname CGTip@eq:cn\endcsname{Eq. (6.1), p. 23: c\_n\textasciicircum{}RVM:=pi\_RVM(n)-pi\_RVM(n-1), n>=2,}
\expandafter\gdef\csname CGTip@eq:weighted-zeta\endcsname{Eq. (6.2), p. 23: zeta\_Z,RVM(s):=prod\_n>=2(1-n\textasciicircum{}-s)\textasciicircum{}-c\_n\textasciicircum{}RVM.}
\expandafter\gdef\csname CGTip@eq:JZ\endcsname{Eq. (6.3), p. 23: J\_Z,RVM(x):=sum\_1<= m<=log\_2x(1)/(m) pi\_RVM(floor( x\textasciicircum{}1/m)),}
\expandafter\gdef\csname CGTip@eq:J-star-Z\endcsname{Eq. (6.4), p. 23: J\_RVM\textasciicircum{}*,Z(x):=J\_RVM\textasciicircum{}*(floor( x)).}
\expandafter\gdef\csname CGTip@thm:unit-cell\endcsname{Theorem 6.1 (Lipschitz control of J\_RVM) (p. 23): There is an absolute constant C\_J such that |J\_RVM(x+theta)-J\_RVM(x)|<= C\_J theta (x>=2, 0<theta<=1).}
\expandafter\gdef\csname CGTip@eq:unit-cell\endcsname{Eq. (6.5), p. 23: |J\_RVM(x+theta)-J\_RVM(x)|<= C\_J theta (x>=2, 0<theta<=1).}
\expandafter\gdef\csname CGTip@eq:Dk-def\endcsname{Eq. (6.6), p. 24: D\_k:=Ei(rho\_k\textasciicircum{}RVMu\_2)-Ei(rho\_k\textasciicircum{}RVMu\_1)=int\_u\_1\textasciicircum{}u\_2(e\textasciicircum{}v/2e\textasciicircum{}ivtau\_k)/(v) dv .}
\expandafter\gdef\csname CGTip@eq:damped-sum\endcsname{Eq. (6.7), p. 24: sum\_ke\textasciicircum{}-alphatau\_kD\_k=int\_u\_1\textasciicircum{}u\_2(e\textasciicircum{}v/2)/(v) S\_alpha(v) dv, S\_alpha(v):=sum\_k>=1g(tau\_k), g(tau):=e\textasciicircum{}(iv-alpha)tau.}
\expandafter\gdef\csname CGTip@eq:S-eps-split\endcsname{Eq. (6.8), p. 24: S\_alpha(v)=int\_tau\_0\textasciicircum{}inf g dN\_RVM+g(tau\_0)N\_RVM(tau\_0)+int\_tau\_0\textasciicircum{}infpsi(N\_RVM(tau)+(1)/(2))g'(tau) dtau .}
\expandafter\gdef\csname CGTip@eq:K-harmonics\endcsname{Eq. (6.9), p. 24: int\_tau\_0\textasciicircum{}infpsi(N\_RVM+(1)/(2))g' dtau=-sum\_m>=1((-1)\textasciicircum{}m)/(2pi i m)(K\_m\textasciicircum{}+-K\_m\textasciicircum{}-), K\_m\textasciicircum{}+-:=int\_tau\_0\textasciicircum{}inf(iv-alpha)e\textasciicircum{}-alphataue\textasciicircum{}iPsi\_m\textasciicircum{}+-(tau) dtau,}
\expandafter\gdef\csname CGTip@eq:K-minus-bound\endcsname{Eq. (6.10), p. 24: |K\_m\textasciicircum{}-|<=2(v+1)((8)/(mlog2)+8sqrt((2T\_m)/(m))) =2(v+1)((8)/(mlog2)+16 (e\textasciicircum{}v/(2m))/( m)).}
\expandafter\gdef\csname CGTip@eq:S-eps-bound\endcsname{Eq. (6.11), p. 24: |S\_alpha(v)|<= C (v+1) e\textasciicircum{}v/2 (v>=log2, 0<alpha<=1),}
\expandafter\gdef\csname CGTip@eq:Dk-sum-bound\endcsname{Eq. (6.12), p. 25: |sum\_ke\textasciicircum{}-alphatau\_kD\_k|<= C'(u\_2-u\_1) x<= C'theta,}
\expandafter\gdef\csname CGTip@cor:cn-bounded\endcsname{Corollary 6.2 (Bounded cell masses) (p. 25): pi\_RVM(y)-pi\_RVM(floor( y))=O(1), c\_n\textasciicircum{}RVM=O(1).}
\expandafter\gdef\csname CGTip@eq:unit-cell-consequence\endcsname{Eq. (6.13), p. 25: pi\_RVM(y)-pi\_RVM(floor( y))=O(1), c\_n\textasciicircum{}RVM=O(1).}
\expandafter\gdef\csname CGTip@cor:ZRVM-Dirichlet\endcsname{Corollary 6.3 (The integer-weighted zeta is a Dirichlet series) (p. 25): The product CGeqrefeq:weighted-zeta converges absolutely for Re s>1 and zeta\_Z,RVM(s)=sum\_N>=1(a\_N\textasciicircum{}Z)/(N\textasciicircum{}s), a\_1\textasciicircum{}Z=1, sum\_N>=1(|a\_N\textasciicircum{}Z|)/(N\textasciicircum{}sigma)<inf (sigma>1). The inverse is also an absolutely convergent Dirichlet series on the whole half-plane: (1)/(zeta\_Z,RVM(s))=sum\_N>=1(conj a\_N\textasciicircum{}Z)/(N\textasciicircum{}s), c...}
\expandafter\gdef\csname CGTip@eq:ZRVM-Dirichlet\endcsname{Eq. (6.14), p. 25: zeta\_Z,RVM(s)=sum\_N>=1(a\_N\textasciicircum{}Z)/(N\textasciicircum{}s), a\_1\textasciicircum{}Z=1, sum\_N>=1(|a\_N\textasciicircum{}Z|)/(N\textasciicircum{}sigma)<inf (sigma>1).}
\expandafter\gdef\csname CGTip@eq:ZRVM-inverse-Dirichlet\endcsname{Eq. (6.15), p. 25: (1)/(zeta\_Z,RVM(s))=sum\_N>=1(conj a\_N\textasciicircum{}Z)/(N\textasciicircum{}s), conj a\_1\textasciicircum{}Z=1, sum\_N>=1(|conj a\_N\textasciicircum{}Z|)/(N\textasciicircum{}sigma)<inf (sigma>1).}
\expandafter\gdef\csname CGTip@eq:log-ZRVM\endcsname{Eq. (6.16), p. 25: logzeta\_Z,RVM(s)=sum\_n>=2c\_n\textasciicircum{}RVMsum\_m>=1(n\textasciicircum{}-ms)/(m)=sum\_N>=2(b\_N)/(N\textasciicircum{}s), b\_N:=sum\_n\textasciicircum{}m=N(c\_n\textasciicircum{}RVM)/(m),}
\expandafter\gdef\csname CGTip@thm:MS\endcsname{Theorem 6.4 (Mellin--Stieltjes identification) (p. 25): There is a function J\_RVM\textasciicircum{} M on (1,inf) with J\_RVM\textasciicircum{} M(x)x\textasciicircum{}-c in L\textasciicircum{}2((1,inf),dx/x) for every c>1 such that logzeta\_RVM(s)=sint\_1\textasciicircum{}inf J\_RVM\textasciicircum{} M(x) x\textasciicircum{}-s-1 dx (Re s>1), and J\_RVM\textasciicircum{} M-J\_RVM is constant on (1,inf). In particular J\_RVM\textasciicircum{} M=J\_RVM\textasciicircum{}*+C\_RVM on [2,inf) for a constant C\_RVM.}
\expandafter\gdef\csname CGTip@eq:MS-RVM\endcsname{Eq. (6.17), p. 25: logzeta\_RVM(s)=sint\_1\textasciicircum{}inf J\_RVM\textasciicircum{} M(x) x\textasciicircum{}-s-1 dx (Re s>1),}
\expandafter\gdef\csname CGTip@cor:weighted-factorization\endcsname{Corollary 6.5 (Weighted factorization) (p. 26): J\_Z,RVM(x)-J\_RVM\textasciicircum{}*(x)=O(1+loglog x), and zeta\_Z,RVM(s)=zeta\_RVM(s) e\textasciicircum{}h\_RVM(s), where h\_RVM(s):=sint\_1\textasciicircum{}inf(J\_Z,RVM(x)-J\_RVM\textasciicircum{} M(x))x\textasciicircum{}-s-1 dx is holomorphic for Re s>0. The two functions have the same zeros in Re s>0, and zeta\_Z,RVM inherits through e\textasciicircum{}h\_RVM the pole of zeta\_RVM at s=1; compare the m...}
\expandafter\gdef\csname CGTip@eq:JZ-diff\endcsname{Eq. (6.18), p. 26: J\_Z,RVM(x)-J\_RVM\textasciicircum{}*(x)=O(1+loglog x),}
\expandafter\gdef\csname CGTip@eq:weighted-factorization\endcsname{Eq. (6.19), p. 26: zeta\_Z,RVM(s)=zeta\_RVM(s) e\textasciicircum{}h\_RVM(s),}
\expandafter\gdef\csname CGTip@eq:hRVM\endcsname{Eq. (6.20), p. 26: h\_RVM(s):=sint\_1\textasciicircum{}inf(J\_Z,RVM(x)-J\_RVM\textasciicircum{} M(x))x\textasciicircum{}-s-1 dx}
\expandafter\gdef\csname CGTip@lem:cell-cancel\endcsname{Lemma 6.6 (Exact first-harmonic cell cancellation) (p. 26): For every integer n, int\_n-1\textasciicircum{}n2cos(2pi x) dx=0 .}
\expandafter\gdef\csname CGTip@eq:cell-cancel\endcsname{Eq. (6.21), p. 26: int\_n-1\textasciicircum{}n2cos(2pi x) dx=0 .}
\expandafter\gdef\csname CGTip@hyp:cell-remainder\endcsname{Hypothesis 6.7 (Uniform cell remainder away from the exceptional powers of 2; proved in Corollary CGrefcor:clock-cells) (p. 26): For n>=3 with n2\textasciicircum{}j:mu(j)!=0, c\_n\textasciicircum{}RVM-int\_n-1\textasciicircum{}n((1)/(log x)+2cos(2pi x))dx=o ((1)/(log n)).}
\expandafter\gdef\csname CGTip@eq:cell-remainder\endcsname{Eq. (6.22), p. 26: c\_n\textasciicircum{}RVM-int\_n-1\textasciicircum{}n((1)/(log x)+2cos(2pi x))dx=o ((1)/(log n)).}
\expandafter\gdef\csname CGTip@prop:eventual-positive\endcsname{Proposition 6.8 (Eventual positivity) (p. 26): Hypothesis CGrefhyp:cell-remainder implies c\_n\textasciicircum{}RVM (1)/(log n), hence 0<c\_n\textasciicircum{}RVM<1 for all sufficiently large n in that set.}
\expandafter\gdef\csname CGTip@eq:eventual-positive\endcsname{Eq. (6.23), p. 26: c\_n\textasciicircum{}RVM (1)/(log n),}
\expandafter\gdef\csname CGTip@sec:mertens-transform\endcsname{Section 7: The critical RVM Mertens transform (p. 26)}
\expandafter\gdef\csname CGTip@eq:F-RVM\endcsname{Eq. (7.1), p. 26: F\_RVM(s)=hat M\_RVM\textasciicircum{}crit(s):=(1)/(zeta\_RVM(s))=(G(s))/(P\_RVM\textasciicircum{}crit(s)) .}
\expandafter\gdef\csname CGTip@eq:Mcrit-formal\endcsname{Eq. (7.2), p. 26: M\_RVM\textasciicircum{}crit(x)=(1)/(2pi i)int\_(c)F\_RVM(s)(x\textasciicircum{}s)/(s) ds ;}
\expandafter\gdef\csname CGTip@prop:Mcrit-poles\endcsname{Proposition 7.1 (Synthetic pole abscissa) (p. 27): F\_RVM=G/P\_RVM\textasciicircum{}crit is meromorphic on C and holomorphic in Re s>(1)/(2). The poles of this meromorphic function are the simple poles (1)/(2)+- itau\_k (k>=1) and s=-2,-4,-6,... and s=1 is a simple zero of F\_RVM.}
\expandafter\gdef\csname CGTip@eq:Mcrit-poles\endcsname{Eq. (7.3), p. 27: (1)/(2)+- itau\_k (k>=1) and s=-2,-4,-6,...}
\expandafter\gdef\csname CGTip@prop:Gaussian\endcsname{Proposition 7.2 (Gaussian square-root bound under vertical growth) (p. 27): Fix c>1 and 0<eps<c-(1)/(2). Assume, uniformly for (1)/(2)+eps<=sigma<= c, F\_RVM(sigma+it)=exp(o(t\textasciicircum{}2)) (|t|->inf). Define M\_RVM,G\textasciicircum{}crit(x):=(1)/(2pi i)int\_(c)e\textasciicircum{}s\textasciicircum{}2F\_RVM(s)(x\textasciicircum{}s)/(s) ds . Then M\_RVM,G\textasciicircum{}crit(x)=O\_eps(x\textasciicircum{}1/2+eps).}
\expandafter\gdef\csname CGTip@eq:subgaussian\endcsname{Eq. (7.4), p. 27: F\_RVM(sigma+it)=exp(o(t\textasciicircum{}2)) (|t|->inf).}
\expandafter\gdef\csname CGTip@eq:MG\endcsname{Eq. (7.5), p. 27: M\_RVM,G\textasciicircum{}crit(x):=(1)/(2pi i)int\_(c)e\textasciicircum{}s\textasciicircum{}2F\_RVM(s)(x\textasciicircum{}s)/(s) ds .}
\expandafter\gdef\csname CGTip@eq:MG-bound\endcsname{Eq. (7.6), p. 27: M\_RVM,G\textasciicircum{}crit(x)=O\_eps(x\textasciicircum{}1/2+eps).}
\expandafter\gdef\csname CGTip@sec:transition\endcsname{Section 8: The transition theorem and the unsmoothed synthetic Mertens bound (p. 27)}
\expandafter\gdef\csname CGTip@eq:chi\endcsname{Eq. (8.1), p. 27: chi(s):=2\textasciicircum{}spi\textasciicircum{}s-1sin ((pi s)/(2))Gamma(1-s)=pi\textasciicircum{}s-1/2(Gamma((1-s)/2))/(Gamma(s/2)),}
\expandafter\gdef\csname CGTip@eq:transition-zw\endcsname{Eq. (8.2), p. 27: s=(1)/(2)+z, z=sigma'+it, tau\_*:=-iz=t-isigma', w(s):=exp-2pi i N\_RVM(tau\_*).}
\expandafter\gdef\csname CGTip@eq:chi-size\endcsname{Eq. (8.3), p. 27: |chi(sigma+it)|=((|t|)/(2pi))\textasciicircum{}1/2-sigma(1+O(|t|\textasciicircum{}-1)).}
\expandafter\gdef\csname CGTip@thm:transition\endcsname{Theorem 8.1 (Transition evaluation of the midpoint product) (p. 27): There is a constant varpi\_RVM>0 with the following properties. For every B>0 there are t\_0(B) and C\_B such that, for 0<sigma'<= B and t>= t\_0(B), zeta\_RVM(s)=varpi\_RVM(1+w(s))(1+E\_0(s)), |E\_0(s)|<=(C\_B)/(t) . For every 0<sigma'\_1<sigma'\_2<(1)/(2) there are t\_0 and C such that, for sigma'\_1<=sigma...}
\expandafter\gdef\csname CGTip@eq:transition-w\endcsname{Eq. (8.4), p. 27: zeta\_RVM(s)=varpi\_RVM(1+w(s))(1+E\_0(s)), |E\_0(s)|<=(C\_B)/(t) .}
\expandafter\gdef\csname CGTip@eq:transition-chi\endcsname{Eq. (8.5), p. 27: zeta\_RVM(s)=varpi\_RVM(1+chi(s))(1+E(s)), |E(s)|<=(C)/(t) .}
\expandafter\gdef\csname CGTip@subsec:transition-proof\endcsname{Section 8.1: Proof of Theorem thm:transition (p. 28)}
\expandafter\gdef\csname CGTip@eq:transition-product\endcsname{Eq. (8.6), p. 28: (P\_RVM\textasciicircum{}crit(s))/(P\_RVM\textasciicircum{}crit((1)/(2)))=prod\_k>=1(1+(z\textasciicircum{}2)/(tau\_k\textasciicircum{}2))=expLambda(z), Lambda(z):=sum\_k>=1f(tau\_k), f(tau):=Log(1+(z\textasciicircum{}2)/(tau\textasciicircum{}2)).}
\expandafter\gdef\csname CGTip@eq:fprime\endcsname{Eq. (8.7), p. 28: f'(tau)=-(2z\textasciicircum{}2)/(tau(tau\textasciicircum{}2+z\textasciicircum{}2))=-(2)/(tau)+(1)/(tau-iz)+(1)/(tau+iz).}
\expandafter\gdef\csname CGTip@eq:sawtooth-count\endcsname{Eq. (8.8), p. 28: n\_RVM(tau)=floor(N\_RVM(tau)+(1)/(2))=N\_RVM(tau)-psi(N\_RVM(tau)+(1)/(2)), n\_RVM(tau\_0)=0 .}
\expandafter\gdef\csname CGTip@eq:transition-split\endcsname{Eq. (8.9), p. 28: Lambda(z)=Lambda\_ sm(z)+Lambda\_ saw(z), Lambda\_ sm(z):=int\_tau\_0\textasciicircum{}inf f dN\_RVM+f(tau\_0)N\_RVM(tau\_0),; Lambda\_ saw(z):=int\_tau\_0\textasciicircum{}infpsi(N\_RVM(tau)+(1)/(2))f'(tau) dtau .}
\expandafter\gdef\csname CGTip@lem:TR-smooth\endcsname{Lemma 8.2 (Smooth part) (p. 28): With log G(s):=Log(z\textasciicircum{}2-(1)/(4))+LogGamma((s)/(2))-(s)/(2)logpi, Lambda\_ sm(z)-log G(s)=nu\_3+O(t\textasciicircum{}-1), nu\_3 in R independent of z .}
\expandafter\gdef\csname CGTip@eq:smooth-difference\endcsname{Eq. (8.10), p. 28: Lambda\_ sm(z)-log G(s)=nu\_3+O(t\textasciicircum{}-1), nu\_3 in R independent of z .}
\expandafter\gdef\csname CGTip@eq:smooth-J-ibp\endcsname{Eq. (8.11), p. 28: Lambda\_ sm(z)=-int\_tau\_0\textasciicircum{}infN\_RVM f' dtau=int\_tau\_0\textasciicircum{}infN\_RVM(tau)((2)/(tau)-(2tau)/(tau\textasciicircum{}2+z\textasciicircum{}2))dtau .}
\expandafter\gdef\csname CGTip@eq:J2\endcsname{Eq. (8.12), p. 28: (7)/(8)Log(1+(z\textasciicircum{}2)/(tau\_0\textasciicircum{}2))=(7)/(4)Log z-(7)/(4)logtau\_0+O(t\textasciicircum{}-2).}
\expandafter\gdef\csname CGTip@eq:I-def\endcsname{Eq. (8.13), p. 28: I(z):=(1)/()\_0\textasciicircum{}inflog(tau)/(2pi e) (z\textasciicircum{}2)/(tau\textasciicircum{}2+z\textasciicircum{}2) dtau .}
\expandafter\gdef\csname CGTip@eq:J-asymp\endcsname{Eq. (8.14), p. 28: Lambda\_ sm(z)=(z)/(2)Log(z)/(2pi e)+(7)/(4)Log z+nu\_1+O(t\textasciicircum{}-2), nu\_1=-(7)/(4)logtau\_0-(1)/()\_0\textasciicircum{}tau\_0log(tau)/(2pi e) dtau .}
\expandafter\gdef\csname CGTip@eq:stirling-loggamma\endcsname{Eq. (8.15), p. 29: LogGamma((s)/(2))=(z)/(2)Log(z)/(2)-(1)/(4)Log z+(1)/(4)log2-(z)/(2)+(1)/(2)log2pi+O(t\textasciicircum{}-1),}
\expandafter\gdef\csname CGTip@eq:G-asymp\endcsname{Eq. (8.16), p. 29: log G(s)=(z)/(2)Log(z)/(2pi e)+(7)/(4)Log z+nu\_2+O(t\textasciicircum{}-1), nu\_2=(1)/(4)log2+(1)/(2)log2pi-(1)/(4)logpi .}
\expandafter\gdef\csname CGTip@lem:TR-clock\endcsname{Lemma 8.3 (The clock as a holomorphic phase) (p. 29): N\_RVM(tau)=(tau)/(2pi)(Log(tau)/(2pi)-1)+(7)/(8) is holomorphic on Retau>0, N\_RVM(conj tau)=conj N\_RVM(tau), and for Retau>=tau\_0, 2pi ImN\_RVM(tau)>=(log2) Imtau (Imtau>=0), 2pi ImN\_RVM(tau)<=(log2) Imtau (Imtau<=0).}
\expandafter\gdef\csname CGTip@eq:clock-halfplane-bound\endcsname{Eq. (8.17), p. 29: 2pi ImN\_RVM(tau)>=(log2) Imtau (Imtau>=0), 2pi ImN\_RVM(tau)<=(log2) Imtau (Imtau<=0).}
\expandafter\gdef\csname CGTip@eq:clock-im-identity\endcsname{Eq. (8.18), p. 29: 2\_RVM(tau)=Ylog(|tau|)/(2pi e)+X =Y[log(X)/(2pi e)+(1)/(2)log(1+v\textasciicircum{}2)+( v)/(v)]>= Y[log2+varphi(v)],}
\expandafter\gdef\csname CGTip@lem:TR-harmonic\endcsname{Lemma 8.4 (Abel-regularized harmonic sum) (p. 29): Lambda\_ saw(z)=Log(1+w(s))+nu\_4+O(t\textasciicircum{}-2), nu\_4 in R independent of z .}
\expandafter\gdef\csname CGTip@eq:Psi-asymp\endcsname{Eq. (8.19), p. 29: Lambda\_ saw(z)=Log(1+w(s))+nu\_4+O(t\textasciicircum{}-2), nu\_4 in R independent of z .}
\expandafter\gdef\csname CGTip@eq:psi-r\endcsname{Eq. (8.20), p. 29: psi\_r(y):=-sum\_m>=1r\textasciicircum{}m(sin2pi my)/(pi m)=(1)/()(1-re\textasciicircum{}2pi iy).}
\expandafter\gdef\csname CGTip@eq:Psi-r-limit\endcsname{Eq. (8.21), p. 29: Lambda\_ saw(z)=lim\_r1Lambda\_r(z), Lambda\_r(z):=int\_tau\_0\textasciicircum{}infpsi\_r(N\_RVM+(1)/(2))f' dtau .}
\expandafter\gdef\csname CGTip@eq:Psi-r-series\endcsname{Eq. (8.22), p. 29: Lambda\_r(z)=-sum\_m>=1((-r)\textasciicircum{}m)/(2pi i m)(I\_m\textasciicircum{}+-I\_m\textasciicircum{}-), I\_m\textasciicircum{}+-:=int\_tau\_0\textasciicircum{}inf e\textasciicircum{}+-2pi imN\_RVM(tau)f'(tau) dtau .}
\expandafter\gdef\csname CGTip@eq:ray-integrals\endcsname{Eq. (8.23), p. 29: R\_m\textasciicircum{}+:=iint\_0\textasciicircum{}inf e\textasciicircum{}2pi imN\_RVM(tau\_0+iy)f'(tau\_0+iy) dy, R\_m\textasciicircum{}-:=-iint\_0\textasciicircum{}inf e\textasciicircum{}-2pi imN\_RVM(tau\_0-iy)f'(tau\_0-iy) dy .}
\expandafter\gdef\csname CGTip@eq:lower-contour\endcsname{Eq. (8.24), p. 30: R\_m\textasciicircum{}--I\_m\textasciicircum{}-=2pi i e\textasciicircum{}-2pi imN\_RVM(tau\_*)=2pi i w\textasciicircum{}m .}
\expandafter\gdef\csname CGTip@eq:Psi-r-closed\endcsname{Eq. (8.25), p. 30: Lambda\_r(z)=-sum\_m>=1((-rw)\textasciicircum{}m)/(m)+Phi\_r(z)=Log(1+rw)+Phi\_r(z), Phi\_r(z):=-sum\_m>=1((-r)\textasciicircum{}m)/(2pi i m)(R\_m\textasciicircum{}+-R\_m\textasciicircum{}-).}
\expandafter\gdef\csname CGTip@eq:ray-bound\endcsname{Eq. (8.26), p. 30: |R\textasciicircum{}+-,(z)\_m|<=(4)/(t\textasciicircum{}2)int\_0\textasciicircum{}inf2\textasciicircum{}-my(tau\_0+y) dy=(4)/(t\textasciicircum{}2)((tau\_0)/(mlog2)+(1)/((mlog2)\textasciicircum{}2)).}
\expandafter\gdef\csname CGTip@eq:assembly\endcsname{Eq. (8.27), p. 30: zeta\_RVM(s) =(P\_RVM\textasciicircum{}crit(s))/(G(s))=P\_RVM\textasciicircum{}crit((1)/(2))exp(Lambda\_ sm(z)-log G(s)+Lambda\_ saw(z)); =P\_RVM\textasciicircum{}crit((1)/(2)) e\textasciicircum{}nu\_3+nu\_4(1+w)(1+O(t\textasciicircum{}-1)),}
\expandafter\gdef\csname CGTip@eq:varpi-def\endcsname{Eq. (8.28), p. 30: varpi\_RVM:=P\_RVM\textasciicircum{}crit((1)/(2)) e\textasciicircum{}nu\_3+nu\_4>0,}
\expandafter\gdef\csname CGTip@eq:theta-def\endcsname{Eq. (8.29), p. 30: vartheta(v):=(1)/(2i)(LogGamma((1)/(4)+(iv)/(2))-LogGamma((1)/(4)-(iv)/(2)))-(v)/(2)logpi,}
\expandafter\gdef\csname CGTip@eq:clock-theta\endcsname{Eq. (8.30), p. 30: N\_RVM(v)=(vartheta(v))/(pi)+1+O(|v|\textasciicircum{}-1).}
\expandafter\gdef\csname CGTip@eq:w-chi\endcsname{Eq. (8.31), p. 30: w=e\textasciicircum{}-2ivartheta(tau\_*)-2pi i+O(1/t)=chi(s)(1+O(t\textasciicircum{}-1)).}
\expandafter\gdef\csname CGTip@hyp:varpi-less-one\endcsname{Hypothesis 8.5 (Strict inequality for the transition constant) (p. 30): varpi\_RVM<1.}
\expandafter\gdef\csname CGTip@lem:TR-boundary\endcsname{Lemma 8.6 (Extension to the critical line) (p. 31): The function E\_0=zeta\_RVM/(varpi\_RVM(1+w))-1 extends holomorphically to |sigma'|<(1)/(2), t>t\_0, where t\_0=t\_0((1)/(4)) is the threshold of Theorem CGrefthm:transition(a) with B=(1)/(4), and |E\_0|<= C/t for 0<=sigma'<=(1)/(4).}
\expandafter\gdef\csname CGTip@cor:derivative-zeros\endcsname{Corollary 8.7 (Derivative at the zeros) (p. 31): For k with tau\_k>t\_0, zeta\_RVM'(rho\_k\textasciicircum{}RVM)=varpi\_RVMlog(tau\_k)/(2pi)(1+O((1)/(tau\_k))).}
\expandafter\gdef\csname CGTip@eq:derivative-zeros\endcsname{Eq. (8.32), p. 31: zeta\_RVM'(rho\_k\textasciicircum{}RVM)=varpi\_RVMlog(tau\_k)/(2pi)(1+O((1)/(tau\_k))).}
\expandafter\gdef\csname CGTip@rem:TR-audit\endcsname{Remark 8.8 (Structure of the proof) (p. 31): The delicate points are handled as follows: the smooth part is computed with principal branches by the identity theorem (Lemma CGreflem:TR-smooth); the clock estimate is quantified (Lemma CGreflem:TR-clock); the sawtooth series is replaced by its Abel means, termwise integration is justified for ...}
\expandafter\gdef\csname CGTip@cor:vertical-limit\endcsname{Corollary 8.9 (Vertical constant limit) (p. 32): For every fixed sigma>(1)/(2), zeta\_RVM(sigma+it)->varpi\_RVM, F\_RVM(sigma+it)-> L:=(1)/(varpi\_RVM) (|t|->inf). More precisely, with sigma'=sigma-(1)/(2), |w(sigma+it)|<= (|t|/2pi)\textasciicircum{}-sigma'(1+O\_B(|t|\textasciicircum{}-1)) (0<sigma'<= B), and, with alpha\_c:=min(c-(1)/(2),1) for c>(1)/(2), F\_RVM(c+it)-L=O\_c(|t|\textasciicircum{}-alph...}
\expandafter\gdef\csname CGTip@eq:vertical-limit\endcsname{Eq. (8.33), p. 32: zeta\_RVM(sigma+it)->varpi\_RVM, F\_RVM(sigma+it)-> L:=(1)/(varpi\_RVM) (|t|->inf).}
\expandafter\gdef\csname CGTip@eq:w-decay\endcsname{Eq. (8.34), p. 32: |w(sigma+it)|<= (|t|/2pi)\textasciicircum{}-sigma'(1+O\_B(|t|\textasciicircum{}-1)) (0<sigma'<= B),}
\expandafter\gdef\csname CGTip@eq:FminusL-all\endcsname{Eq. (8.35), p. 32: F\_RVM(c+it)-L=O\_c(|t|\textasciicircum{}-alpha\_c).}
\expandafter\gdef\csname CGTip@subsec:halfplane\endcsname{Section 8.2: Bounds on the whole half-plane (p. 32)}
\expandafter\gdef\csname CGTip@eq:saw-primitive\endcsname{Eq. (8.36), p. 32: s\_1(tau):=int\_tau\_0\textasciicircum{}taupsi(N\_RVM(v)+(1)/(2)) dv (tau>=tau\_0).}
\expandafter\gdef\csname CGTip@lem:saw-primitive\endcsname{Lemma 8.10 (Bounded primitive of the midpoint sawtooth) (p. 32): |s\_1(tau)|<=(pi)/(4log2) for all tau>=tau\_0.}
\expandafter\gdef\csname CGTip@thm:halfplane-RVM\endcsname{Theorem 8.11 (Two-sided bound and limit at +inf) (p. 32): There is an absolute constant R\_0 such that, for s=(1)/(2)+z, z=sigma'+it, sigma'>(1)/(2) and |z|>= R\_0, zeta\_RVM(s)=varpi\_RVM e\textasciicircum{}theta\_RVM(s), |theta\_RVM(s)|<=(pi\textasciicircum{}2)/(2log2)*(1)/(sigma')+(C)/(|z|). Consequently, for every sigma\_0>1 there is C(sigma\_0) with C(sigma\_0)\textasciicircum{}-1<=|zeta\_RVM(s)|<= C(sigma\_0...}
\expandafter\gdef\csname CGTip@eq:halfplane-RVM\endcsname{Eq. (8.37), p. 32: zeta\_RVM(s)=varpi\_RVM e\textasciicircum{}theta\_RVM(s), |theta\_RVM(s)|<=(pi\textasciicircum{}2)/(2log2)*(1)/(sigma')+(C)/(|z|).}
\expandafter\gdef\csname CGTip@eq:halfplane-two-sided\endcsname{Eq. (8.38), p. 32: C(sigma\_0)\textasciicircum{}-1<=|zeta\_RVM(s)|<= C(sigma\_0) (Re s>=sigma\_0),}
\expandafter\gdef\csname CGTip@eq:nu4-closed\endcsname{Eq. (8.39), p. 32: nu\_4=-2int\_tau\_0\textasciicircum{}inf(s\_1(tau))/(tau\textasciicircum{}2) dtau .}
\expandafter\gdef\csname CGTip@eq:smooth-exact\endcsname{Eq. (8.40), p. 32: Lambda\_ sm(z)=(z)/(2)Log(z)/(2pi e)+(7)/(8)Log(1+(z\textasciicircum{}2)/(tau\_0\textasciicircum{}2))-(1)/()\_0\textasciicircum{}tau\_0log(tau)/(2pi e) (z\textasciicircum{}2)/(tau\textasciicircum{}2+z\textasciicircum{}2) dtau .}
\expandafter\gdef\csname CGTip@eq:saw-ibp\endcsname{Eq. (8.41), p. 32: Lambda\_ saw(z)=-int\_tau\_0\textasciicircum{}infs\_1f'' dtau, f''(tau)=(2)/(tau\textasciicircum{}2)-(1)/((tau-iz)\textasciicircum{}2)-(1)/((tau+iz)\textasciicircum{}2).}
\expandafter\gdef\csname CGTip@eq:saw-bound\endcsname{Eq. (8.42), p. 33: |Lambda\_ saw(z)+2int\_tau\_0\textasciicircum{}inf(s\_1(tau))/(tau\textasciicircum{}2) dtau|<=(pi)/(4log2)*(2pi)/(sigma') .}
\expandafter\gdef\csname CGTip@thm:no-Dirichlet-RVM\endcsname{Theorem 8.12 (No absolutely convergent Dirichlet expansion) (p. 33): Neither zeta\_RVM nor F\_RVM=1/zeta\_RVM is represented by a Dirichlet series sum\_na\_nn\textasciicircum{}-s converging absolutely in some half-plane Re s>sigma\_0.}
\expandafter\gdef\csname CGTip@eq:mean-square\endcsname{Eq. (8.43), p. 33: (1)/(2T)int\_-T\textasciicircum{}T|F\_RVM(sigma+it)-L|\textasciicircum{}2dt =sum\_m,n a\_mconj a\_n(mn)\textasciicircum{}-sigma (1)/(2T)int\_-T\textasciicircum{}T((n)/(m))\textasciicircum{}itdt ,}
\expandafter\gdef\csname CGTip@eq:mean-square-limit\endcsname{Eq. (8.44), p. 33: lim\_T->inf(1)/(2T)int\_-T\textasciicircum{}T|F\_RVM(sigma+it)-L|\textasciicircum{}2dt=sum\_n>=1| a\_n|\textasciicircum{}2n\textasciicircum{}-2sigma.}
\expandafter\gdef\csname CGTip@thm:unsmoothed-RVM\endcsname{Theorem 8.13 (Unsmoothed synthetic Mertens bound) (p. 33): For every x>1 and every c>(1)/(2) the symmetric Perron integral M\_RVM\textasciicircum{}crit(x):=lim\_U->inf(1)/(2pi i)int\_c-iU\textasciicircum{}c+iUF\_RVM(s)(x\textasciicircum{}s)/(s) ds exists, is independent of c>(1)/(2), and equals the absolutely convergent expression M\_RVM\textasciicircum{}crit(x)=L+(1)/(2pi i)int\_(c)(F\_RVM(s)-L)(x\textasciicircum{}s)/(s) ds . For every eps in ...}
\expandafter\gdef\csname CGTip@eq:Mcrit-symmetric\endcsname{Eq. (8.45), p. 33: M\_RVM\textasciicircum{}crit(x):=lim\_U->inf(1)/(2pi i)int\_c-iU\textasciicircum{}c+iUF\_RVM(s)(x\textasciicircum{}s)/(s) ds}
\expandafter\gdef\csname CGTip@eq:Mcrit-renorm\endcsname{Eq. (8.46), p. 33: M\_RVM\textasciicircum{}crit(x)=L+(1)/(2pi i)int\_(c)(F\_RVM(s)-L)(x\textasciicircum{}s)/(s) ds .}
\expandafter\gdef\csname CGTip@eq:Mcrit-unsmoothed-bound\endcsname{Eq. (8.47), p. 33: M\_RVM\textasciicircum{}crit(x)=L+O\_eps(x\textasciicircum{}1/2+eps).}
\expandafter\gdef\csname CGTip@eq:bromwich\endcsname{Eq. (8.48), p. 33: lim\_U->inf(1)/(2pi i)int\_c-iU\textasciicircum{}c+iU(x\textasciicircum{}s)/(s) ds=1 .}
\expandafter\gdef\csname CGTip@eq:Mcrit-final-estimate\endcsname{Eq. (8.49), p. 34: |M\_RVM\textasciicircum{}crit(x)-L|<=(x\textasciicircum{}1/2+eps)/(2pi)int\_-inf\textasciicircum{}inf(|F\_RVM((1)/(2)+eps+it)-L|)/(|(1)/(2)+eps+it|) dt,}
\expandafter\gdef\csname CGTip@rem:unsmoothed-scope\endcsname{Remark 8.14 (Scope) (p. 34): Theorem CGrefthm:unsmoothed-RVM concerns 1/zeta\_RVM. The transfer to M goes through the endpoint of the lambda-family and is Theorem CGrefthm:step3.}
\expandafter\gdef\csname CGTip@sec:comparison\endcsname{Section 9: The exact Hadamard product and the comparison with zetazeta (p. 34)}
\expandafter\gdef\csname CGTip@eq:Pzeta\endcsname{Eq. (9.1), p. 34: P\_zeta(s):=G(s)zeta(s)=2xi(s), xi(s)=(1)/(2)s(s-1)pi\textasciicircum{}-s/2Gamma(s/2)zeta(s).}
\expandafter\gdef\csname CGTip@thm:Hadamard\endcsname{Theorem 9.1 (Exact Hadamard product) (p. 34): Let rho run over the nontrivial zeros of zeta with Imrho>0, with multiplicity, ordered by Imrho. The product P\_zeta(s)=prod\_Imrho>0(1+(q(s))/(rho(1-rho))) converges absolutely and locally uniformly in s in C, and the identity holds with no exponential or constant factor.}
\expandafter\gdef\csname CGTip@eq:Hadamard-exact\endcsname{Eq. (9.2), p. 34: P\_zeta(s)=prod\_Imrho>0(1+(q(s))/(rho(1-rho)))}
\expandafter\gdef\csname CGTip@eq:hadamard-general\endcsname{Eq. (9.3), p. 34: 2xi(s)=e\textasciicircum{}Bsprod\_rho(1-(s)/(rho))e\textasciicircum{}s/rho,}
\expandafter\gdef\csname CGTip@eq:pair-factor\endcsname{Eq. (9.4), p. 34: (1-(s)/(rho))(1-(s)/(1-rho))e\textasciicircum{}s/rho+s/(1-rho) =(1+(q(s))/(rho(1-rho)))e\textasciicircum{}s/(rho(1-rho)).}
\expandafter\gdef\csname CGTip@eq:hadamard-exp\endcsname{Eq. (9.5), p. 34: 2xi(s)=e\textasciicircum{}B's Pi\_zeta(s), Pi\_zeta(s):=prod\_Imrho>0(1+(q(s))/(rho(1-rho))), B'=B+sum\_Imrho>0(1)/(rho(1-rho)) .}
\expandafter\gdef\csname CGTip@eq:Q-RVM\endcsname{Eq. (9.6), p. 34: Q\_RVM(s):=(zeta(s))/(zeta\_RVM(s))=(P\_zeta(s))/(P\_RVM\textasciicircum{}crit(s)).}
\expandafter\gdef\csname CGTip@prop:M-transform-relation\endcsname{Proposition 9.2 (Exact transform relation) (p. 35): For Re s>1, and by meromorphic continuation, F\_RVM(s)=Q\_RVM(s) (1)/(zeta(s)) . If the inverse Mellin transform Q\_RVM\textasciicircum{}v of Q\_RVM and the inverse defining M\_RVM\textasciicircum{}crit exist in a common distribution space with Mellin convolution, then M\_RVM\textasciicircum{}crit=Q\_RVM\textasciicircum{}v*Mel M .}
\expandafter\gdef\csname CGTip@eq:M-transform-relation\endcsname{Eq. (9.7), p. 35: F\_RVM(s)=Q\_RVM(s) (1)/(zeta(s)) .}
\expandafter\gdef\csname CGTip@eq:M-convolution\endcsname{Eq. (9.8), p. 35: M\_RVM\textasciicircum{}crit=Q\_RVM\textasciicircum{}v*Mel M .}
\expandafter\gdef\csname CGTip@cor:Q-control\endcsname{Corollary 9.3 (All growth of zeta to the right of (1)/(2) lies in Q\_RVM) (p. 35): For every c'>(1)/(2) and 0<eps<c'-(1)/(2) there are t\_0=t\_0(eps,c') and 0<c\_-(eps,c')<= c\_+(eps,c') such that c\_-<=|zeta\_RVM(sigma+it)|<= c\_+ ((1)/(2)+eps<=sigma<= c', |t|>= t\_0), and therefore, in the same region, c\_-|Q\_RVM(s)|<=|zeta(s)|<= c\_+|Q\_RVM(s)|, (1)/(c\_+|Q\_RVM(s)|)<=|(1)/(zeta(s))|<=(1...}
\expandafter\gdef\csname CGTip@eq:zetaRVM-two-sided\endcsname{Eq. (9.9), p. 35: c\_-<=|zeta\_RVM(sigma+it)|<= c\_+ ((1)/(2)+eps<=sigma<= c', |t|>= t\_0),}
\expandafter\gdef\csname CGTip@eq:zeta-asymp-Q\endcsname{Eq. (9.10), p. 35: c\_-|Q\_RVM(s)|<=|zeta(s)|<= c\_+|Q\_RVM(s)|, (1)/(c\_+|Q\_RVM(s)|)<=|(1)/(zeta(s))|<=(1)/(c\_-|Q\_RVM(s)|).}
\expandafter\gdef\csname CGTip@eq:prime-defect\endcsname{Eq. (9.11), p. 35: d\_n\textasciicircum{}RVM:=1\_n prime-c\_n\textasciicircum{}RVM.}
\expandafter\gdef\csname CGTip@eq:prime-defect-cumulative\endcsname{Eq. (9.12), p. 35: sum\_2<= n<= Nd\_n\textasciicircum{}RVM=pi(N)-pi\_RVM(N).}
\expandafter\gdef\csname CGTip@sec:lambda-family\endcsname{Section 10: The lambdalambda-family from zeta\_RVMzeta\_RVM to zeta\_Z1zeta\_Z1 (p. 35)}
\expandafter\gdef\csname CGTip@eq:gamma-lambda\endcsname{Eq. (10.1), p. 35: gamma\_k(lambda):=(1-lambda)gamma\_k+lambdatau\_k=gamma\_k+lambdaeta\_k,}
\expandafter\gdef\csname CGTip@eq:P-lambda\endcsname{Eq. (10.2), p. 35: P\_lambda(s):=prod\_k>=1(1+(q(s))/((1)/(4)+gamma\_k(lambda)\textasciicircum{}2)),}
\expandafter\gdef\csname CGTip@eq:zeta-lambda\endcsname{Eq. (10.3), p. 35: zeta\_lambda(s):=(P\_lambda(s))/(G(s)), F\_lambda(s):=(1)/(zeta\_lambda(s)),}
\expandafter\gdef\csname CGTip@eq:family-endpoints\endcsname{Eq. (10.4), p. 35: P\_Z1:=P\_0, zeta\_Z1:=zeta\_0, zeta\_1=zeta\_RVM.}
\expandafter\gdef\csname CGTip@thm:path\endcsname{Theorem 10.1 (The lambda-family) (p. 36): The product CGeqrefeq:P-lambda converges absolutely and locally uniformly jointly on [0,1]xC. For every lambda, P\_lambda is entire of order exactly one, P\_lambda(1-s)=P\_lambda(s), P\_lambda(conj s)=conj P\_lambda(s), P\_lambda(0)=P\_lambda(1)=1, and its zeros, with multiplicity, are exactly (1)/(2)+-...}
\expandafter\gdef\csname CGTip@eq:family-divisor\endcsname{Eq. (10.5), p. 36: (1)/(2)+- igamma\_k(lambda), k>=1 .}
\expandafter\gdef\csname CGTip@eq:lambda-normal\endcsname{Eq. (10.6), p. 36: sum\_k>=1sup\_s in K, 0<=lambda<=1|(q(s))/((1)/(4)+gamma\_k(lambda)\textasciicircum{}2)|<inf,}
\expandafter\gdef\csname CGTip@eq:order-one\endcsname{Eq. (10.7), p. 36: logmax\_|s|<= R\_ fam|P\_lambda(s)|<=sum\_k>=1log(1+(R\_ fam\textasciicircum{}2+R\_ fam)/((1)/(4)+gamma\_k(lambda)\textasciicircum{}2))=O(R\_ famlog R\_ fam),}
\expandafter\gdef\csname CGTip@eq:ordinate-gaps\endcsname{Eq. (10.8), p. 36: gamma\_k+1(lambda)-gamma\_k(lambda)=(1-lambda)(gamma\_k+1-gamma\_k)+lambda(tau\_k+1-tau\_k)>=lambda(tau\_k+1-tau\_k)>0,}
\expandafter\gdef\csname CGTip@rem:only-gamma-tau\endcsname{Remark 10.2 (All differences lie in gamma\_k-tau\_k) (p. 36): For lambda,lambda' in [0,1] and s off the divisors, (zeta\_lambda(s))/(zeta\_lambda'(s))=prod\_k>=1(((1)/(4)+(gamma\_k+lambda'eta\_k)\textasciicircum{}2)((1)/(4)+(gamma\_k+lambdaeta\_k)\textasciicircum{}2+q(s)))/(((1)/(4)+(gamma\_k+lambdaeta\_k)\textasciicircum{}2)((1)/(4)+(gamma\_k+lambda'eta\_k)\textasciicircum{}2+q(s))) . The gamma factor cancels; the members differ only...}
\expandafter\gdef\csname CGTip@eq:ratio-members\endcsname{Eq. (10.9), p. 36: (zeta\_lambda(s))/(zeta\_lambda'(s))=prod\_k>=1(((1)/(4)+(gamma\_k+lambda'eta\_k)\textasciicircum{}2)((1)/(4)+(gamma\_k+lambdaeta\_k)\textasciicircum{}2+q(s)))/(((1)/(4)+(gamma\_k+lambdaeta\_k)\textasciicircum{}2)((1)/(4)+(gamma\_k+lambda'eta\_k)\textasciicircum{}2+q(s))) .}
\expandafter\gdef\csname CGTip@thm:generator\endcsname{Theorem 10.3 (Infinitesimal generator) (p. 36): Put A\_k,lambda:=(1)/(4)+gamma\_k(lambda)\textasciicircum{}2. Off the divisor of P\_lambda, V\_lambda(s):=d\_ P\_lambda(s)=-2q(s)sum\_k>=1(gamma\_k(lambda) eta\_k)/(A\_k,lambda(A\_k,lambda+q(s))), the series converging locally uniformly on compact sets separated from the divisor. Since G does not depend on lambda, d\_\_lambda...}
\expandafter\gdef\csname CGTip@eq:generator\endcsname{Eq. (10.10), p. 36: V\_lambda(s):=d\_ P\_lambda(s)=-2q(s)sum\_k>=1(gamma\_k(lambda) eta\_k)/(A\_k,lambda(A\_k,lambda+q(s))),}
\expandafter\gdef\csname CGTip@eq:generator-zeta\endcsname{Eq. (10.11), p. 36: d\_\_lambda=V\_lambda, d\_ F\_lambda=-V\_lambda .}
\expandafter\gdef\csname CGTip@eq:mode-derivative\endcsname{Eq. (10.12), p. 36: d\_(1+(q)/(A\_k,lambda))=-(2q gamma\_k(lambda)eta\_k)/(A\_k,lambda(A\_k,lambda+q)) .}
\expandafter\gdef\csname CGTip@prop:integrated-multiplier\endcsname{Proposition 10.4 (Integrated multiplier) (p. 37): On every component of C\_lambda' in [0,1](1)/(2)+- igamma\_k(lambda'):k>=1, zeta\_lambda(s)=zeta\_Z1(s)exp(int\_0\textasciicircum{}\_lambda'(s) dlambda'), F\_lambda(s)=F\_Z1(s)exp(-int\_0\textasciicircum{}\_lambda'(s) dlambda'), and the total multiplier is Pi\_ tr(s):=exp(-int\_0\textasciicircum{}1V\_lambda'(s) dlambda')=(P\_Z1(s))/(P\_RVM\textasciicircum{}crit(s)), zeta\_Z1=zet...}
\expandafter\gdef\csname CGTip@eq:integrated-zeta\endcsname{Eq. (10.13), p. 37: zeta\_lambda(s)=zeta\_Z1(s)exp(int\_0\textasciicircum{}\_lambda'(s) dlambda'),}
\expandafter\gdef\csname CGTip@eq:integrated-inverse\endcsname{Eq. (10.14), p. 37: F\_lambda(s)=F\_Z1(s)exp(-int\_0\textasciicircum{}\_lambda'(s) dlambda'),}
\expandafter\gdef\csname CGTip@eq:transport-multiplier\endcsname{Eq. (10.15), p. 37: Pi\_ tr(s):=exp(-int\_0\textasciicircum{}1V\_lambda'(s) dlambda')=(P\_Z1(s))/(P\_RVM\textasciicircum{}crit(s)), zeta\_Z1=zeta\_RVM Pi\_ tr .}
\expandafter\gdef\csname CGTip@eq:mode-integral\endcsname{Eq. (10.16), p. 37: int\_0\textasciicircum{}1d\_lambda'log(1+(q)/(A\_k,lambda'))dlambda'==log(1+q/((1)/(4)+tau\_k\textasciicircum{}2))/(1+q/((1)/(4)+gamma\_k\textasciicircum{}2))+-od2pi i.}
\expandafter\gdef\csname CGTip@cor:endpoint-expansion\endcsname{Corollary 10.5 (Expansion at the endpoint) (p. 37): Let K be compact, disjoint from the zeros of P\_Z1 and from s=1 (equivalently one may state the expansion for the product ratio P\_lambda/P\_Z1 without this exclusion). Uniformly for s in K, as lambda0, zeta\_lambda(s)=zeta\_Z1(s)(1+\_0(s)+O\_K(lambda\textasciicircum{}2)), F\_lambda(s)=F\_Z1(s)(1-\_0(s)+O\_K(lambda\textasciicircum{}2)).}
\expandafter\gdef\csname CGTip@eq:endpoint-zeta-expansion\endcsname{Eq. (10.17), p. 37: zeta\_lambda(s)=zeta\_Z1(s)(1+\_0(s)+O\_K(lambda\textasciicircum{}2)),}
\expandafter\gdef\csname CGTip@eq:endpoint-inverse-expansion\endcsname{Eq. (10.18), p. 37: F\_lambda(s)=F\_Z1(s)(1-\_0(s)+O\_K(lambda\textasciicircum{}2)).}
\expandafter\gdef\csname CGTip@prop:residues\endcsname{Proposition 10.6 (Residues along the family) (p. 37): Suppose gamma\_k is a simple ordinate (it occurs once in the list). Put rho\_k,lambda=(1)/(2)+igamma\_k(lambda) and, for 0<lambda<=1, r\_k,lambda:=Res\_s=rho\_k,lambda(1)/(s zeta\_lambda(s))=(G(rho\_k,lambda))/(rho\_k,lambdaP\_lambda'(rho\_k,lambda)). Then r\_k,lambda->(G((1)/(2)+igamma\_k))/(((1)/(2)+igamma\_...}
\expandafter\gdef\csname CGTip@eq:residue-transport\endcsname{Eq. (10.19), p. 37: r\_k,lambda:=Res\_s=rho\_k,lambda(1)/(s zeta\_lambda(s))=(G(rho\_k,lambda))/(rho\_k,lambdaP\_lambda'(rho\_k,lambda)).}
\expandafter\gdef\csname CGTip@eq:residue-limit\endcsname{Eq. (10.20), p. 37: r\_k,lambda->(G((1)/(2)+igamma\_k))/(((1)/(2)+igamma\_k)P\_Z1'((1)/(2)+igamma\_k)) (lambda0).}
\expandafter\gdef\csname CGTip@rem:multiplicity\endcsname{Remark 10.7 (Multiplicity) (p. 37): For 0<lambda<=1 every zero of zeta\_lambda is simple. At lambda=0 a repeated ordinate of zeta appears in zeta\_Z1 with the same multiplicity, and the residues of Proposition CGrefprop:residues are replaced by the full principal part. Multiplicity plays no role in Step 3: Theorem CGrefthm:step3 uses...}
\expandafter\gdef\csname CGTip@sec:Z1-step3\endcsname{Section 11: The endpoint zeta\_Z1zeta\_Z1, Hamburger's theorem, and Step 3 (p. 38)}
\expandafter\gdef\csname CGTip@subsec:Z1-class\endcsname{Section 11.1: The analytic class of the endpoint (p. 38)}
\expandafter\gdef\csname CGTip@prop:Z1-class\endcsname{Proposition 11.1 (Analytic class of zeta\_Z1) (p. 38): The endpoint P\_Z1(s)=prod\_k>=1(1+(q(s))/((1)/(4)+gamma\_k\textasciicircum{}2)), zeta\_Z1(s)=(P\_Z1(s))/(G(s)), has the following properties. P\_Z1 is entire of order exactly one, P\_Z1(1-s)=P\_Z1(s), P\_Z1(0)=P\_Z1(1)=1; the nontrivial zeros of zeta\_Z1 are exactly (1)/(2)+- igamma\_k, k>=1, with the multiplicities of the ...}
\expandafter\gdef\csname CGTip@eq:Z1-product\endcsname{Eq. (11.1), p. 38: P\_Z1(s)=prod\_k>=1(1+(q(s))/((1)/(4)+gamma\_k\textasciicircum{}2)), zeta\_Z1(s)=(P\_Z1(s))/(G(s)),}
\expandafter\gdef\csname CGTip@eq:Z1-divisor\endcsname{Eq. (11.2), p. 38: (1)/(2)+- igamma\_k, k>=1,}
\expandafter\gdef\csname CGTip@eq:Z1-FE\endcsname{Eq. (11.3), p. 38: pi\textasciicircum{}-s/2Gamma(s/2)zeta\_Z1(s)=pi\textasciicircum{}-(1-s)/2Gamma((1-s)/2)zeta\_Z1(1-s);}
\expandafter\gdef\csname CGTip@eq:G-near-one\endcsname{Eq. (11.4), p. 38: G(s)=(s-1)+O((s-1)\textasciicircum{}2) (s->1).}
\expandafter\gdef\csname CGTip@eq:Z1-FE-proof\endcsname{Eq. (11.5), p. 38: pi\textasciicircum{}-s/2Gamma(s/2)zeta\_Z1(s)=(P\_Z1(s))/(q(s)),}
\expandafter\gdef\csname CGTip@eq:Z1-entire-removal\endcsname{Eq. (11.6), p. 38: (s-1)zeta\_Z1(s)=P\_Z1(s) (pi\textasciicircum{}s/2)/(s Gamma(s/2)),}
\expandafter\gdef\csname CGTip@eq:Z1-defect-ratio\endcsname{Eq. (11.7), p. 38: R\_Z1(s):=(zeta\_Z1(s))/(zeta(s))=(P\_Z1(s))/(P\_zeta(s)),}
\expandafter\gdef\csname CGTip@subsec:Hamburger\endcsname{Section 11.2: Hamburger's theorem at the endpoint (p. 38)}
\expandafter\gdef\csname CGTip@thm:Hamburger-Z1\endcsname{Theorem 11.2 (Hamburger closure) (p. 38): Consider the condition (D) zeta\_Z1(s)=sum\_n>=1(a\_n)/(n\textasciicircum{}s) absolutely for Re s>1 . Then (D) <=> R\_Z1==1 <=> zeta\_Z1=zeta <=> RH.}
\expandafter\gdef\csname CGTip@eq:Z1-D\endcsname{Eq. (11.8), p. 38: (D) zeta\_Z1(s)=sum\_n>=1(a\_n)/(n\textasciicircum{}s) absolutely for Re s>1 .}
\expandafter\gdef\csname CGTip@eq:Hamburger-equivalence\endcsname{Eq. (11.9), p. 38: (D) <=> R\_Z1==1 <=> zeta\_Z1=zeta <=> RH.}
\expandafter\gdef\csname CGTip@eq:Hamburger-conclusion\endcsname{Eq. (11.10), p. 39: zeta\_Z1=C\_H zeta}
\expandafter\gdef\csname CGTip@eq:RH-product\endcsname{Eq. (11.11), p. 39: P\_zeta(s)=prod\_k>=1(1+(q(s))/((1)/(4)+gamma\_k\textasciicircum{}2))=P\_Z1(s)}
\expandafter\gdef\csname CGTip@subsec:D-multiplier\endcsname{Section 11.3: Condition (D) and the integer-weighted zeta (p. 39)}
\expandafter\gdef\csname CGTip@eq:K-multiplier\endcsname{Eq. (11.12), p. 39: K(s):=(zeta\_Z1(s))/(zeta\_Z,RVM(s)) .}
\expandafter\gdef\csname CGTip@thm:D-multiplier\endcsname{Theorem 11.3 (Reduction of (D) to the multiplier) (p. 39): If K has an expansion K(s)=sum\_n>=1k\_nn\textasciicircum{}-s that converges absolutely in Re s>1, then CGeqrefeq:Z1-D holds. If CGeqrefeq:Z1-D holds, then K(s)=sum\_n>=1k\_nn\textasciicircum{}-s absolutely on the whole half-plane Re s>1. Thus (D) holds if and only if K is an absolutely convergent Dirichlet series in Re s>1. Hence, b...}
\expandafter\gdef\csname CGTip@eq:K-explicit\endcsname{Eq. (11.13), p. 39: K(s)=Pi\_ tr(s) e\textasciicircum{}-h\_RVM(s)=(P\_Z1(s))/(P\_RVM\textasciicircum{}crit(s)) e\textasciicircum{}-h\_RVM(s) ,}
\expandafter\gdef\csname CGTip@rem:D-atomic\endcsname{Remark 11.4 (Continuous-to-atomic meaning of (D)) (p. 39): The prime-power count J\_Z1\textasciicircum{} M of zeta\_Z1 (Theorem CGrefthm:mellin-counts) is an L\textasciicircum{}2\_ loc count built from the zeros, and (D) makes its distributional derivative the atomic measure sum\_p,mm\textasciicircum{}-1delta\_p\textasciicircum{}m. The final step of the method is therefore the passage of zeta\_Z1 into the discrete Dirichlet ca...}
\expandafter\gdef\csname CGTip@subsec:step3\endcsname{Section 11.4: Step 3: the Mertens bound from zeta\_Z1=zetazeta\_Z1 = zeta (p. 39)}
\expandafter\gdef\csname CGTip@lem:zeta-growth\endcsname{Lemma 11.5 (Vertical growth of zeta) (p. 39): For (1)/(4)<=sigma<=4 and |t|>=2, |zeta(sigma+it)|<=15|t| .}
\expandafter\gdef\csname CGTip@eq:zeta-growth\endcsname{Eq. (11.14), p. 39: |zeta(sigma+it)|<=15|t| .}
\expandafter\gdef\csname CGTip@eq:zeta-integral-rep\endcsname{Eq. (11.15), p. 40: zeta(s)=(s)/(s-1)-sint\_1\textasciicircum{}infyy\textasciicircum{}-s-1 dy .}
\expandafter\gdef\csname CGTip@lem:log-zeta\endcsname{Lemma 11.6 (Logarithm of zeta off the critical line) (p. 40): Assume zeta(s)!=0 for Re s>(1)/(2). Let 0<delta<=(1)/(4) and omega(delta):=(log(3-2delta))/(log(3-delta))<1 . Then, with the branch of logzeta equal to sum\_n>=2Lambda(n)n\textasciicircum{}-s/log n on Re s>1, |logzeta(sigma+it)|<=((48)/()|t|)\textasciicircum{}omega(delta) (sigma>=(1)/(2)+delta, |t|>=4). Consequently, for every eps...}
\expandafter\gdef\csname CGTip@eq:three-circles-exponent\endcsname{Eq. (11.16), p. 40: omega(delta):=(log(3-2delta))/(log(3-delta))<1 .}
\expandafter\gdef\csname CGTip@eq:log-zeta-bound\endcsname{Eq. (11.17), p. 40: |logzeta(sigma+it)|<=((48)/()|t|)\textasciicircum{}omega(delta) (sigma>=(1)/(2)+delta, |t|>=4).}
\expandafter\gdef\csname CGTip@eq:inv-zeta-bound\endcsname{Eq. (11.18), p. 40: |(1)/(zeta(sigma+it))|<=|t|\textasciicircum{}eps (sigma>=(1)/(2)+delta, |t|>= t\_2).}
\expandafter\gdef\csname CGTip@eq:borel-caratheodory\endcsname{Eq. (11.19), p. 40: max\_|s-s\_0|<= r\_3|logzeta(s)|<=(2r\_3)/(R-r\_3)max\_|s-s\_0|=RRelogzeta(s)+(R+r\_3)/(R-r\_3)|logzeta(s\_0)|,}
\expandafter\gdef\csname CGTip@eq:BC-bound\endcsname{Eq. (11.20), p. 40: max\_|s-s\_0|<= r\_3|logzeta(s)|<=(12)/(delta)(log(30|t|)+logzeta(2))<=(48)/()|t| .}
\expandafter\gdef\csname CGTip@eq:three-circles\endcsname{Eq. (11.21), p. 40: M\_2<= M\_1\textasciicircum{}1-omegaM\_3\textasciicircum{}omega, omega=(log(r\_2/r\_1))/(log(r\_3/r\_1))<=(log(3-2delta))/(log(3-delta))=omega(delta).}
\expandafter\gdef\csname CGTip@lem:perron-kernel\endcsname{Lemma 11.7 (Perron kernel) (p. 40): For y>0, y!=1, c>0, T>0 [ref], |(1)/(2pi i)int\_c-iT\textasciicircum{}c+iT(y\textasciicircum{}s)/(s) ds-1\_y>1|<= y\textasciicircum{}cmin(1,(1)/(T|log y|)).}
\expandafter\gdef\csname CGTip@eq:perron-kernel\endcsname{Eq. (11.22), p. 40: |(1)/(2pi i)int\_c-iT\textasciicircum{}c+iT(y\textasciicircum{}s)/(s) ds-1\_y>1|<= y\textasciicircum{}cmin(1,(1)/(T|log y|)).}
\expandafter\gdef\csname CGTip@thm:step3\endcsname{Theorem 11.8 (Step 3) (p. 40): If zeta\_Z1=zeta, then for every eps>0 M(x)=O\_eps(x\textasciicircum{}1/2+eps) (x>=1).}
\expandafter\gdef\csname CGTip@eq:step3-bound\endcsname{Eq. (11.23), p. 40: M(x)=O\_eps(x\textasciicircum{}1/2+eps) (x>=1).}
\expandafter\gdef\csname CGTip@eq:perron-truncated\endcsname{Eq. (11.24), p. 41: M(x)=(1)/(2pi i)int\_c-iT\textasciicircum{}c+iT(x\textasciicircum{}s)/(s zeta(s)) ds+O(sum\_n>=1((x)/(n))\textasciicircum{}cmin(1,(1)/(T|log(x/n)|))).}
\expandafter\gdef\csname CGTip@eq:perron-error\endcsname{Eq. (11.25), p. 41: M(x)=(1)/(2pi i)int\_c-iT\textasciicircum{}c+iT(x\textasciicircum{}s)/(s zeta(s)) ds+O(log x).}
\expandafter\gdef\csname CGTip@eq:contour-shift\endcsname{Eq. (11.26), p. 41: int\_c-iT\textasciicircum{}c+iT=int\_(1)/(2)+delta-iT\textasciicircum{}(1)/(2)+delta+iT+int\_(1)/(2)+delta+iT\textasciicircum{}c+iT-int\_(1)/(2)+delta-iT\textasciicircum{}c-iT.}
\expandafter\gdef\csname CGTip@eq:horizontal-bound\endcsname{Eq. (11.27), p. 41: (1)/(2pi)int\_(1)/(2)+delta\textasciicircum{}c(x\textasciicircum{}sigma T\textasciicircum{}eps)/(T) dsigma<= x\textasciicircum{}cT\textasciicircum{}eps-1=e x\textasciicircum{}eps.}
\expandafter\gdef\csname CGTip@eq:left-bound\endcsname{Eq. (11.28), p. 41: (1)/(2pi)int\_-T\textasciicircum{}T(x\textasciicircum{}1/2+delta)/(|(1)/(2)+delta+it| |zeta((1)/(2)+delta+it)|) dt<= x\textasciicircum{}1/2+delta(C'(delta,eps)+int\_t\_2\textasciicircum{}Tt\textasciicircum{}eps-1dt)<<\_delta,epsx\textasciicircum{}1/2+delta+eps.}
\expandafter\gdef\csname CGTip@thm:transfer\endcsname{Theorem 11.9 (Transfer of the synthetic Mertens bound to M) (p. 41): Assume zeta\_Z1=zeta and let 0<delta<=(1)/(4), eps>0. Then: on Re s>(1)/(2), (1)/(zeta(s))=F\_RVM(s) Pi\_ tr(s)\textasciicircum{}-1, Pi\_ tr=(P\_Z1)/(P\_RVM\textasciicircum{}crit)=(zeta)/(zeta\_RVM), where Pi\_ tr is holomorphic and zero-free in Re s>(1)/(2) and Pi\_ tr(1)=1; there is t\_3=t\_3(delta,eps) such that, for (1)/(2)+delta<=sigma...}
\expandafter\gdef\csname CGTip@eq:transfer-factorization\endcsname{Eq. (11.29), p. 41: (1)/(zeta(s))=F\_RVM(s) Pi\_ tr(s)\textasciicircum{}-1, Pi\_ tr=(P\_Z1)/(P\_RVM\textasciicircum{}crit)=(zeta)/(zeta\_RVM),}
\expandafter\gdef\csname CGTip@eq:transfer-multiplier-bound\endcsname{Eq. (11.30), p. 41: |F\_RVM(sigma+it)|<=4L, |Pi\_ tr(sigma+it)\textasciicircum{}+-1|<=|t|\textasciicircum{}eps;}
\expandafter\gdef\csname CGTip@eq:transfer-conclusion\endcsname{Eq. (11.31), p. 41: M(x)=O\_eps(x\textasciicircum{}1/2+eps).}
\expandafter\gdef\csname CGTip@eq:log-zetaRVM-bounded\endcsname{Eq. (11.32), p. 42: logzeta\_RVM=logvarpi\_RVM+Log(1+w)+Log(1+E\_0)=O(1).}
\expandafter\gdef\csname CGTip@eq:log-Pi-tr\endcsname{Eq. (11.33), p. 42: |logPi\_ tr(sigma+it)|=|logzeta-logzeta\_RVM|<=((48)/()|t|)\textasciicircum{}omega(delta)+O(1)|t|}
\expandafter\gdef\csname CGTip@rem:transfer-mellin\endcsname{Remark 11.10 (The transfer in Mellin form) (p. 42): Formally, CGeqrefeq:transfer-factorization reads M=(Pi\_ tr\textasciicircum{}-1)\textasciicircum{}v*Mel M\_RVM\textasciicircum{}crit: the Mertens function is the synthetic Mertens function convolved with the inverse Mellin transform of the inverse transport multiplier. The synthetic factor F\_RVM contributes the exponent (1)/(2) and is bounded (Theo...}
\expandafter\gdef\csname CGTip@cor:step3-D\endcsname{Corollary 11.11 (Step 3 through (D)) (p. 42): Condition CGeqrefeq:Z1-D implies M(x)=O\_eps(x\textasciicircum{}1/2+eps) for every eps>0.}
\expandafter\gdef\csname CGTip@rem:Step3-status\endcsname{Remark 11.12 (Status of Step 3) (p. 42): The implication zeta\_Z1=zeta M(x)=O\_eps(x\textasciicircum{}1/2+eps) is the classical consequence RH M(x)=O\_eps(x\textasciicircum{}1/2+eps) (see [ref]) rewritten and reproved in the endpoint notation. It requires nothing beyond the identity: no separate statement about the zeros of zeta enters, because zeta=zeta\_Z1 already carries...}
\expandafter\gdef\csname CGTip@subsec:reciprocal\endcsname{Section 11.5: Reciprocal endpoint and finite-height Perron integrals (p. 42)}
\expandafter\gdef\csname CGTip@eq:FZ1\endcsname{Eq. (11.34), p. 42: F\_Z1(s):=(1)/(zeta\_Z1(s))=F\_RVM(s) Pi\_ tr(s)\textasciicircum{}-1,}
\expandafter\gdef\csname CGTip@eq:truncated-Perron\endcsname{Eq. (11.35), p. 42: M\_lambda,U(x):=(1)/(2pi i)int\_c-iU\textasciicircum{}c+iUF\_lambda(s)(x\textasciicircum{}s)/(s) ds, M\_Z1,U:=M\_0,U.}
\expandafter\gdef\csname CGTip@prop:finite-Perron\endcsname{Proposition 11.13 (Finite-height Perron integrals along the family) (p. 42): For fixed c>1, U<inf and x>0, M\_lambda,U(x)-> M\_Z1,U(x) (lambda0).}
\expandafter\gdef\csname CGTip@eq:finite-Perron-limit\endcsname{Eq. (11.36), p. 42: M\_lambda,U(x)-> M\_Z1,U(x) (lambda0).}
\expandafter\gdef\csname CGTip@thm:coefficient-criterion\endcsname{Theorem 11.14 (Reciprocal-coefficient criterion) (p. 42): Assume that in Re s>1 F\_Z1(s)=sum\_n>=1(mu\_Z1(n))/(n\textasciicircum{}s) absolutely, and let M\_Z1(x)=sum\_n<= xmu\_Z1(n), so that mu\_Z1(n)=M\_Z1(n)-M\_Z1(n-1). Then the following are equivalent: (i) M\_Z1(n)=M(n) for every n>=1; (ii) mu\_Z1=mu; (iii) F\_Z1=1/zeta in Re s>1; (iv) zeta\_Z1=zeta; (v) RH.}
\expandafter\gdef\csname CGTip@eq:Z1-reciprocal-Dirichlet\endcsname{Eq. (11.37), p. 42: F\_Z1(s)=sum\_n>=1(mu\_Z1(n))/(n\textasciicircum{}s)}
\expandafter\gdef\csname CGTip@eq:Z1-increments\endcsname{Eq. (11.38), p. 43: mu\_Z1(n)=M\_Z1(n)-M\_Z1(n-1).}
\expandafter\gdef\csname CGTip@rem:negative-convolution\endcsname{Remark 11.15 (A route that does not close Step 3) (p. 43): The following approach does not work. One may try to expand both factors of F\_Z1=F\_RVM Pi\_ tr\textasciicircum{}-1 in CGeqrefeq:FZ1 as absolutely convergent Dirichlet series and identify their Dirichlet convolution with mu. By Theorem CGrefthm:no-Dirichlet-RVM the expansion of F\_RVM does not exist in any half-plan...}
\expandafter\gdef\csname CGTip@sec:euler-products\endcsname{Section 12: The endpoint at infinity and the Euler products (p. 43)}
\expandafter\gdef\csname CGTip@subsec:value-infinity\endcsname{Section 12.1: The value of the endpoint at infinity (p. 43)}
\expandafter\gdef\csname CGTip@eq:frak-e\endcsname{Eq. (12.1), p. 43: e(tau):=N\_zeta(tau)-N\_RVM(tau), e\_1(tau):=int\_tau\_0\textasciicircum{}taue(v) dv (tau>=tau\_0=4pi).}
\expandafter\gdef\csname CGTip@eq:frak-e1-bound\endcsname{Eq. (12.2), p. 43: |e\_1(tau)|<= C\_elogtau (tau>=tau\_0).}
\expandafter\gdef\csname CGTip@thm:Z1-clock\endcsname{Theorem 12.1 (Clock decomposition of the endpoint) (p. 43): Let s=(1)/(2)+z, z=sigma'+it, sigma'>0. There is an absolute constant R\_ end>0 such that logzeta\_Z1(s)=logvarpi\_Z1-int\_tau\_0\textasciicircum{}infe\_1(tau)((1)/((tau-iz)\textasciicircum{}2)+(1)/((tau+iz)\textasciicircum{}2))dtau+O((1)/(|z|)) (|z|>= R\_ end), where varpi\_Z1:=P\_Z1((1)/(2))exp(nu\_3+2int\_tau\_0\textasciicircum{}inf(e\_1(tau))/(tau\textasciicircum{}2) dtau)>0 . Consequentl...}
\expandafter\gdef\csname CGTip@eq:Z1-clock\endcsname{Eq. (12.3), p. 43: logzeta\_Z1(s)=logvarpi\_Z1-int\_tau\_0\textasciicircum{}infe\_1(tau)((1)/((tau-iz)\textasciicircum{}2)+(1)/((tau+iz)\textasciicircum{}2))dtau+O((1)/(|z|)) (|z|>= R\_ end),}
\expandafter\gdef\csname CGTip@eq:varpi-Z1\endcsname{Eq. (12.4), p. 43: varpi\_Z1:=P\_Z1((1)/(2))exp(nu\_3+2int\_tau\_0\textasciicircum{}inf(e\_1(tau))/(tau\textasciicircum{}2) dtau)>0 .}
\expandafter\gdef\csname CGTip@eq:Z1-growth\endcsname{Eq. (12.5), p. 43: |logzeta\_Z1(s)-logvarpi\_Z1|<=(2(pi+2)C\_e)/(sigma')(1+log(2|t|+2sigma'+tau\_0))+(C)/(|z|);}
\expandafter\gdef\csname CGTip@eq:varpi-Z1-product\endcsname{Eq. (12.6), p. 43: varpi\_Z1=varpi\_RVMprod\_k>=1((1)/(4)+tau\_k\textasciicircum{}2)/((1)/(4)+gamma\_k\textasciicircum{}2):=varpi\_RVMlim\_K->infprod\_k<= K((1)/(4)+tau\_k\textasciicircum{}2)/((1)/(4)+gamma\_k\textasciicircum{}2).}
\expandafter\gdef\csname CGTip@eq:Z1-split\endcsname{Eq. (12.7), p. 44: sum\_kf(gamma\_k)=Lambda\_ sm(z)+int\_tau\_0\textasciicircum{}infe\_1f'' dtau,}
\expandafter\gdef\csname CGTip@eq:log-kernel\endcsname{Eq. (12.8), p. 44: int\_tau\_0\textasciicircum{}inf(logtau)/((tau-a)\textasciicircum{}2+sigma'\textasciicircum{}2) dtau<=(pi)/(sigma')log X+4 (1+log X)/(X)<=(pi+2)/(sigma')(1+log X).}
\expandafter\gdef\csname CGTip@eq:varpi-ratio\endcsname{Eq. (12.9), p. 44: (varpi\_Z1)/(varpi\_RVM)=(P\_Z1((1)/(2)))/(P\_RVM\textasciicircum{}crit((1)/(2)))exp(2int\_tau\_0\textasciicircum{}inf(e\_1(tau)+s\_1(tau))/(tau\textasciicircum{}2) dtau).}
\expandafter\gdef\csname CGTip@eq:count-difference-integral\endcsname{Eq. (12.10), p. 44: int\_tau\_0\textasciicircum{}inf(e\_1+s\_1)/(tau\textasciicircum{}2) dtau=lim\_X->infint\_tau\_0\textasciicircum{}X(N\_zeta(tau)-n\_RVM(tau))/(tau) dtau=lim\_K->infsum\_k<= Klog(tau\_k)/(gamma\_k).}
\expandafter\gdef\csname CGTip@thm:varpi-Z1\endcsname{Theorem 12.2 (The value at infinity decides the endpoint identity) (p. 44): varpi\_Z1>=1, and varpi\_Z1=1 if and only if zeta\_Z1=zeta.}
\expandafter\gdef\csname CGTip@eq:pair-ratio\endcsname{Eq. (12.11), p. 44: (W\_rho\textasciicircum{}2)/((W\_rho-d\_rho\textasciicircum{}2)\textasciicircum{}2+4d\_rho\textasciicircum{}2gamma\textasciicircum{}2)* c\_gamma, c\_gamma=(|rho(1-rho)|\textasciicircum{}2)/(((1)/(4)+gamma\textasciicircum{}2)\textasciicircum{}2)=1+(2d\_rho\textasciicircum{}2(gamma\textasciicircum{}2-(1)/(4))+d\_rho\textasciicircum{}4)/(((1)/(4)+gamma\textasciicircum{}2)\textasciicircum{}2)>1 .}
\expandafter\gdef\csname CGTip@cor:product-identity\endcsname{Corollary 12.3 (The endpoint identity as an identity of the displacements) (p. 45): The following are equivalent: (D); zeta\_Z1=zeta; varpi\_Z1=1; lim\_K->infprod\_k<= K((1)/(4)+gamma\_k\textasciicircum{}2)/((1)/(4)+tau\_k\textasciicircum{}2)=varpi\_RVM (ordered product). Unconditionally lim\_sigma->+infPi\_ tr(sigma)=varpi\_Z1/varpi\_RVM=varpi\_Z1L>= L, and the equality lim\_sigma->+infPi\_ tr(sigma)=L=1/varpi\_RVM is a fourt...}
\expandafter\gdef\csname CGTip@eq:product-identity\endcsname{Eq. (12.12), p. 45: lim\_K->infprod\_k<= K((1)/(4)+gamma\_k\textasciicircum{}2)/((1)/(4)+tau\_k\textasciicircum{}2)=varpi\_RVM}
\expandafter\gdef\csname CGTip@prop:Z1-bounded\endcsname{Proposition 12.4 (Boundedness of the endpoint) (p. 45): Let sigma\_0>1. Then |zeta\_Z1| is bounded above and below on Re s>=sigma\_0 if and only if sup\_sigma'>=sigma\_0-(1)/(2), t>=0|Reint\_tau\_0\textasciicircum{}inf(e\_1(tau)+s\_1(tau))/((tau-t+isigma')\textasciicircum{}2) dtau|<inf .}
\expandafter\gdef\csname CGTip@eq:conjugate-poisson\endcsname{Eq. (12.13), p. 45: sup\_sigma'>=sigma\_0-(1)/(2), t>=0|Reint\_tau\_0\textasciicircum{}inf(e\_1(tau)+s\_1(tau))/((tau-t+isigma')\textasciicircum{}2) dtau|<inf .}
\expandafter\gdef\csname CGTip@subsec:euler-products\endcsname{Section 12.2: The Euler products of zeta\_RVMzeta\_RVM and zeta\_Z1zeta\_Z1 (p. 45)}
\expandafter\gdef\csname CGTip@thm:mellin-counts\endcsname{Theorem 12.5 (Prime-power counts as Mellin transforms) (p. 45): Let X in RVM,Z1. There is a unique function J\_X\textasciicircum{} M on (1,inf) with J\_X\textasciicircum{} M(x)x\textasciicircum{}-c in L\textasciicircum{}2((1,inf),dx/x) for every c>1 and logzeta\_X(s)=sint\_1\textasciicircum{}inf J\_X\textasciicircum{} M(x) x\textasciicircum{}-s-1 dx (Re s>1), the branch of logzeta\_X being real on (1,inf). We use the notation int x\textasciicircum{}-s dJ\_X\textasciicircum{} M:=sint\_1\textasciicircum{}inf J\_X\textasciicircum{} M(x)x\textasciicircum{}-s-1dx distribut...}
\expandafter\gdef\csname CGTip@eq:mellin-count\endcsname{Eq. (12.14), p. 45: logzeta\_X(s)=sint\_1\textasciicircum{}inf J\_X\textasciicircum{} M(x) x\textasciicircum{}-s-1 dx (Re s>1),}
\expandafter\gdef\csname CGTip@eq:mellin-count-laplace\endcsname{Eq. (12.15), p. 45: (logzeta\_X(s))/(s)=int\_0\textasciicircum{}inf j(u) e\textasciicircum{}-su du (Re s>c), e\textasciicircum{}-cuj in L\textasciicircum{}2(0,inf).}
\expandafter\gdef\csname CGTip@thm:distributional-EF\endcsname{Theorem 12.6 (Distributional explicit formula) (p. 46): Let X in RVM,Z1, with ordinates kappa\_k=tau\_k for X=RVM and kappa\_k=gamma\_k for X=Z1. In D'(1,inf), (log x) (dJ\_X\textasciicircum{} M)/(dx)=1-(2)/( x) Resum\_k>=1x\textasciicircum{}ikappa\_k-(1)/(x(x\textasciicircum{}2-1)), the series converging in D'(1,inf). The explicit-formula count J\_X(x):=li(x)-2Resum\_k>=1Ei(((1)/(2)+ikappa\_k)log x)-log2+int\_x...}
\expandafter\gdef\csname CGTip@eq:distributional-EF\endcsname{Eq. (12.16), p. 46: (log x) (dJ\_X\textasciicircum{} M)/(dx)=1-(2)/( x) Resum\_k>=1x\textasciicircum{}ikappa\_k-(1)/(x(x\textasciicircum{}2-1)),}
\expandafter\gdef\csname CGTip@eq:J-X\endcsname{Eq. (12.17), p. 46: J\_X(x):=li(x)-2Resum\_k>=1Ei(((1)/(2)+ikappa\_k)log x)-log2+int\_x\textasciicircum{}inf(dy)/(y(y\textasciicircum{}2-1)log y)}
\expandafter\gdef\csname CGTip@eq:log-derivative-X\endcsname{Eq. (12.18), p. 46: -(zeta\_X')/(zeta\_X)(s)=(1)/(s-1)-sum\_k>=1((1)/(s-rho\_k)+(1)/(s-conj rho\_k))+sum\_n>=1((1)/(2n)-(1)/(s+2n))-(gamma\_E)/(2)-(1)/(2)logpi .}
\expandafter\gdef\csname CGTip@thm:cell-form\endcsname{Theorem 12.7 (Cell form of the endpoint identity) (p. 46): The following are equivalent. (D), that is, zeta\_Z1=zeta. J\_Z1\textasciicircum{} M is constant on every interval (n,n+1), n>=1, and its jumps b\_n:=v\_n-v\_n-1, where v\_n is its value on (n,n+1) and v\_0:=0, satisfy sum\_n<= X| b\_n|=O\_eps(X\textasciicircum{}1+eps) for every eps>0. For every n>=1, in D'(n,n+1), 2 Resum\_k>=1(x\textasciicircum{}igamma\_k-...}
\expandafter\gdef\csname CGTip@eq:cell-identity\endcsname{Eq. (12.19), p. 46: 2 Resum\_k>=1(x\textasciicircum{}igamma\_k-x\textasciicircum{}itau\_k)= x log x J\_RVM'(x),}
\expandafter\gdef\csname CGTip@eq:cell-difference\endcsname{Eq. (12.20), p. 47: (log x)((dJ\_Z1\textasciicircum{} M)/(dx)-J\_RVM'(x))=-(2)/( x) Resum\_k>=1(x\textasciicircum{}igamma\_k-x\textasciicircum{}itau\_k),}
\expandafter\gdef\csname CGTip@rem:cell-form\endcsname{Remark 12.8 (The final step in cell form) (p. 47): Theorem CGrefthm:cell-form is the exact form of the continuous-to-atomic passage of Remark CGrefrem:D-atomic. Both sides of CGeqrefeq:cell-identity are built from the midpoint construction: the right side from the clock alone (J\_RVM' is bounded on [2,inf) by Theorem CGrefthm:unit-cell, equal to 1...}
\expandafter\gdef\csname CGTip@sec:atomized\endcsname{Section 13: The atomized family and the programme for the final step (p. 47)}
\expandafter\gdef\csname CGTip@subsec:family-infinity\endcsname{Section 13.1: The family at infinity (p. 47)}
\expandafter\gdef\csname CGTip@prop:family-clock\endcsname{Proposition 13.1 (Clock decomposition along the family) (p. 47): There is an absolute constant C with |e\_1\textasciicircum{}(lambda)(tau)|<= Clogtau for tau>=tau\_0 and all lambda in [0,1]. Consequently Theorem CGrefthm:Z1-clock holds for every zeta\_lambda with e\_1 replaced by e\_1\textasciicircum{}(lambda), uniformly in lambda; in particular varpi\_lambda:=lim\_sigma->+infzeta\_lambda(sigma) exist...}
\expandafter\gdef\csname CGTip@eq:varpi-lambda\endcsname{Eq. (13.1), p. 47: varpi\_lambda=varpi\_RVMlim\_K->infprod\_k<= K((1)/(4)+tau\_k\textasciicircum{}2)/((1)/(4)+gamma\_k(lambda)\textasciicircum{}2), varpi\_1=varpi\_RVM, varpi\_0=varpi\_Z1.}
\expandafter\gdef\csname CGTip@eq:family-count-primitive\endcsname{Eq. (13.2), p. 48: int\_tau\_0\textasciicircum{}X(N\_lambda-n\_RVM)=sum\_k((X-gamma\_k(lambda))\textasciicircum{}+-(X-tau\_k)\textasciicircum{}+)=(1-lambda)sum\_gamma\_k(lambda),tau\_k<= Xeta\_k+O(log X),}
\expandafter\gdef\csname CGTip@eq:Upsilon\endcsname{Eq. (13.3), p. 48: (lambda):=sum\_k>=1(2gamma\_k(lambda) eta\_k)/((1)/(4)+gamma\_k(lambda)\textasciicircum{}2) (summed in the order of k).}
\expandafter\gdef\csname CGTip@prop:family-flow\endcsname{Proposition 13.2 (The value at infinity along the family) (p. 48): The series CGeqrefeq:Upsilon converges uniformly on [0,1], lambda\_lambda is C\textasciicircum{}1 on [0,1], and (d)/(dlambda)logvarpi\_lambda=-(lambda)=lim\_sigma->+infV\_lambda(sigma), '(lambda)=sum\_k>=1(2eta\_k\textasciicircum{}2((1)/(4)-gamma\_k(lambda)\textasciicircum{}2))/(((1)/(4)+gamma\_k(lambda)\textasciicircum{}2)\textasciicircum{}2)<0 , with V\_lambda the generator of Theorem C...}
\expandafter\gdef\csname CGTip@eq:varpi-derivative\endcsname{Eq. (13.4), p. 48: (d)/(dlambda)logvarpi\_lambda=-(lambda)=lim\_sigma->+infV\_lambda(sigma), '(lambda)=sum\_k>=1(2eta\_k\textasciicircum{}2((1)/(4)-gamma\_k(lambda)\textasciicircum{}2))/(((1)/(4)+gamma\_k(lambda)\textasciicircum{}2)\textasciicircum{}2)<0 ,}
\expandafter\gdef\csname CGTip@eq:flow-integral\endcsname{Eq. (13.5), p. 48: int\_0\textasciicircum{}1(lambda) dlambda=sum\_k>=1log((1)/(4)+tau\_k\textasciicircum{}2)/((1)/(4)+gamma\_k\textasciicircum{}2)=log L+logvarpi\_Z1,}
\expandafter\gdef\csname CGTip@prop:varpi-verified\endcsname{Proposition 13.3 (The value at infinity and the verified zeros) (p. 48): If every zero beta+igamma with 0<gamma<= T\_ ver has beta=(1)/(2), then 0<=logvarpi\_Z1<=0.26sum\_gamma>T\_ ver(1)/(gamma\textasciicircum{}2). With T\_ ver=3*10\textasciicircum{}12 [ref], the right side is (0.26)/(2pi T\_ ver)(log(T\_ ver)/(2pi)+1)<4*10\textasciicircum{}-13.}
\expandafter\gdef\csname CGTip@eq:varpi-verified\endcsname{Eq. (13.6), p. 48: 0<=logvarpi\_Z1<=0.26sum\_gamma>T\_ ver(1)/(gamma\textasciicircum{}2).}
\expandafter\gdef\csname CGTip@subsec:smooth-cells\endcsname{Section 13.2: Atomization by smooth cells (p. 49)}
\expandafter\gdef\csname CGTip@def:atomized\endcsname{Definition 13.4 (Atomized family) (p. 49): The cell masses, the atomized zeta function and its coefficients are c\_n(lambda):= dpi\_lambda,varphi\_n, zeta\_Z,lambda(s):=prod\_n>=2(1-n\textasciicircum{}-s)\textasciicircum{}- c\_n(lambda)=sum\_N>=1( a\_N(lambda))/(N\textasciicircum{}s), a\_N(lambda)=sum\_prod\_nn\textasciicircum{}r\_n=N prod\_n: r\_n>=1 c\_n(lambda)+r\_n-1r\_n, a\_1(lambda)=1 .}
\expandafter\gdef\csname CGTip@eq:atomized-family\endcsname{Eq. (13.7), p. 49: c\_n(lambda):= dpi\_lambda,varphi\_n, zeta\_Z,lambda(s):=prod\_n>=2(1-n\textasciicircum{}-s)\textasciicircum{}- c\_n(lambda)=sum\_N>=1( a\_N(lambda))/(N\textasciicircum{}s),}
\expandafter\gdef\csname CGTip@eq:atomized-coefficients\endcsname{Eq. (13.8), p. 49: a\_N(lambda)=sum\_prod\_nn\textasciicircum{}r\_n=N prod\_n: r\_n>=1 c\_n(lambda)+r\_n-1r\_n, a\_1(lambda)=1 .}
\expandafter\gdef\csname CGTip@eq:J-lambda-derivative\endcsname{Eq. (13.9), p. 49: (d J\_lambda)/(dlambda)(x)=2sum\_k>=1eta\_k Im(x\textasciicircum{}rho\_k,lambda)/(rho\_k,lambda), (d)/(dlambda)((log x)(dJ\_lambda)/(dx))=(2log x)/( x)sum\_k>=1eta\_ksin(gamma\_k(lambda)log x).}
\expandafter\gdef\csname CGTip@subsec:joining-factor\endcsname{Section 13.3: The joining factor cell by cell (p. 49)}
\expandafter\gdef\csname CGTip@eq:atomization-kernel\endcsname{Eq. (13.10), p. 49: a\_s(x):=sum\_n>=2varphi\_n(x) n\textasciicircum{}-s-x\textasciicircum{}-s, h\_n(lambda;s):= dpi\_lambda, varphi\_n(x)(n\textasciicircum{}-s-x\textasciicircum{}-s) .}
\expandafter\gdef\csname CGTip@eq:first-cell\endcsname{Eq. (13.11), p. 49: u\_lambda(s):=C\_lambda 2\textasciicircum{}-s+sint\_1\textasciicircum{}2J\textasciicircum{} M\_lambda(x) x\textasciicircum{}-s-1 dx, C\_lambda:=J\textasciicircum{} M\_lambda-J\_lambda ,}
\expandafter\gdef\csname CGTip@prop:first-cell\endcsname{Proposition 13.5 (The first cell) (p. 49): Let lambda in [0,1]. u\_lambda is entire, and for Re s>1 logzeta\_lambda(s)=u\_lambda(s)+ dJ\textasciicircum{}*\_lambda,x\textasciicircum{}-s, dJ\textasciicircum{}*\_lambda,x\textasciicircum{}-s:=sint\_2\textasciicircum{}inf J\_lambda(x) x\textasciicircum{}-s-1 dx . u\_lambda(sigma+it)=o(|t|) as |t|->inf, uniformly for sigma in compact sets. u\_lambda is constant if and only if J\_lambda=0 on (1,2); in tha...}
\expandafter\gdef\csname CGTip@eq:first-cell-split\endcsname{Eq. (13.12), p. 49: logzeta\_lambda(s)=u\_lambda(s)+ dJ\textasciicircum{}*\_lambda,x\textasciicircum{}-s, dJ\textasciicircum{}*\_lambda,x\textasciicircum{}-s:=sint\_2\textasciicircum{}inf J\_lambda(x) x\textasciicircum{}-s-1 dx .}
\expandafter\gdef\csname CGTip@lem:L2-partial-sums\endcsname{Lemma 13.6 (Mean-square convergence of the ordinate series) (p. 50): Let kappa be a list with |kappa\_k-tau\_k|<= C\_eta (in particular kappa=(gamma\_k(lambda)) for any lambda in [0,1]), rho\_k:=(1)/(2)+ikappa\_k, and let I(1,inf) be compact. The partial sums S\_K(x):=sum\_k<= K(x\textasciicircum{}rho\_k/rho\_k+x\textasciicircum{}conj rho\_k/conj rho\_k) and E\_K(x):=sum\_k<= K(Ei(rho\_klog x)+Ei(conj rho\_klog x...}
\expandafter\gdef\csname CGTip@eq:L2-partial-sums\endcsname{Eq. (13.13), p. 50: |S\_K'-S\_K|\_L\textasciicircum{}2(I)\textasciicircum{}2+|E\_K'-E\_K|\_L\textasciicircum{}2(I)\textasciicircum{}2 <= C\_Isum\_K<k<= K'(log(kappa\_k+2))/(kappa\_k\textasciicircum{}2) (K<K'),}
\expandafter\gdef\csname CGTip@eq:L2-partial-sums-proof\endcsname{Eq. (13.14), p. 50: intomega(u)|sum\_K<k<= K'(e\textasciicircum{}ikappa\_ku)/(rho\_k)|\textasciicircum{}2du <<sum\_K<k,k'<= K'((1+|kappa\_k-kappa\_k'|)\textasciicircum{}-2)/(kappa\_kkappa\_k'); <<sum\_K<k<= K'(1)/(kappa\_k\textasciicircum{}2)sum\_j>=0(no.k':j<=|kappa\_k'-kappa\_k|<j+1)/((1+j)\textasciicircum{}2) <<sum\_K<k<= K'(log(kappa\_k+2))/(kappa\_k\textasciicircum{}2),}
\expandafter\gdef\csname CGTip@thm:joining-factor\endcsname{Theorem 13.7 (The joining factor) (p. 51): Let h\_lambda:=log zeta\_Z,lambda-logzeta\_lambda on the half-plane where zeta\_Z,lambda converges absolutely (Theorem CGrefthm:atomized-continuation). Then zeta\_Z,lambda=zeta\_lambda e\textasciicircum{} h\_lambda, h\_lambda(s)=-u\_lambda(s)+sum\_m>=1(1)/(m)sum\_n>=2h\_n(lambda;ms), and, explicitly, h\_lambda(s)=-u\_lambda(s)...}
\expandafter\gdef\csname CGTip@eq:joining-factor\endcsname{Eq. (13.15), p. 51: zeta\_Z,lambda=zeta\_lambda e\textasciicircum{} h\_lambda, h\_lambda(s)=-u\_lambda(s)+sum\_m>=1(1)/(m)sum\_n>=2h\_n(lambda;ms),}
\expandafter\gdef\csname CGTip@eq:joining-explicit\endcsname{Eq. (13.16), p. 51: h\_lambda(s)=-u\_lambda(s)+sum\_m>=1(1)/(m)sum\_j>=1(mu(j))/(j)int\_2\textasciicircum{}infa\_ms(y\textasciicircum{}j) (1)/(log y)(1-(2)/( y)Resum\_k>=1y\textasciicircum{}igamma\_k(lambda)-(1)/(y(y\textasciicircum{}2-1)))dy,}
\expandafter\gdef\csname CGTip@eq:joining-derivative\endcsname{Eq. (13.17), p. 51: d\_lambda h\_lambda(s)=-d\_\_lambda(s)+sum\_m>=1(1)/(m)sum\_j>=1(mu(j))/(j)int\_2\textasciicircum{}infa\_ms(y\textasciicircum{}j) (2)/( y)sum\_k>=1eta\_ksin(gamma\_k(lambda)log y) dy .}
\expandafter\gdef\csname CGTip@prop:cell-explicit\endcsname{Proposition 13.8 (Explicit cell defects) (p. 51): For every n>=2, lambda in [0,1] and Re s>1, h\_n(lambda;s)=sum\_j>=1(mu(j))/(j)int\_2\textasciicircum{}infvarphi\_n(y\textasciicircum{}j)(n\textasciicircum{}-s-y\textasciicircum{}-js) (1)/(log y)(1-(2)/( y)Resum\_k>=1y\textasciicircum{}igamma\_k(lambda)-(1)/(y(y\textasciicircum{}2-1)))dy, d\_\_n(lambda;s)=sum\_j>=1(mu(j))/(j)int\_2\textasciicircum{}infvarphi\_n(y\textasciicircum{}j)(n\textasciicircum{}-s-y\textasciicircum{}-js) (2)/( y)sum\_k>=1eta\_ksin(gamma\_k(lambda)log y)...}
\expandafter\gdef\csname CGTip@eq:cell-explicit\endcsname{Eq. (13.18), p. 51: h\_n(lambda;s)=sum\_j>=1(mu(j))/(j)int\_2\textasciicircum{}infvarphi\_n(y\textasciicircum{}j)(n\textasciicircum{}-s-y\textasciicircum{}-js) (1)/(log y)(1-(2)/( y)Resum\_k>=1y\textasciicircum{}igamma\_k(lambda)-(1)/(y(y\textasciicircum{}2-1)))dy,}
\expandafter\gdef\csname CGTip@eq:cell-explicit-derivative\endcsname{Eq. (13.19), p. 51: d\_\_n(lambda;s)=sum\_j>=1(mu(j))/(j)int\_2\textasciicircum{}infvarphi\_n(y\textasciicircum{}j)(n\textasciicircum{}-s-y\textasciicircum{}-js) (2)/( y)sum\_k>=1eta\_ksin(gamma\_k(lambda)log y) dy ,}
\expandafter\gdef\csname CGTip@prop:cell-limit\endcsname{Proposition 13.9 (The limit lambda->0 and the meaning of a vanishing cell) (p. 52): For every n>=2 and Re s>1, h\_n(lambda;s)->h\_n(0;s) and c\_n(lambda)-> c\_n(0) as lambda->0\textasciicircum{}+, and h\_n(0;s)=h\_n(1;s)-int\_0\textasciicircum{}1d\_\_n(lambda;s) dlambda . Moreover h\_n(0;s)=0 for all Re s>1 if and only if varphi\_n dpi\_Z1= c\_n(0) delta\_n.}
\expandafter\gdef\csname CGTip@eq:cell-flow\endcsname{Eq. (13.20), p. 52: h\_n(0;s)=h\_n(1;s)-int\_0\textasciicircum{}1d\_\_n(lambda;s) dlambda .}
\expandafter\gdef\csname CGTip@subsec:mean-square\endcsname{Section 13.4: Mean square of the cell masses (p. 52)}
\expandafter\gdef\csname CGTip@eq:VW\endcsname{Eq. (13.21), p. 52: V\_n(gamma):=sum\_j>=1(mu(j))/(j)int\_2\textasciicircum{}infvarphi\_n(y\textasciicircum{}j) y\textasciicircum{}-1/2e\textasciicircum{}i y dy, W\_n(gamma):=sum\_j>=1(mu(j))/(j)int\_2\textasciicircum{}infvarphi\_n(y\textasciicircum{}j) (y\textasciicircum{}-1/2e\textasciicircum{}i y)/(log y) dy ,}
\expandafter\gdef\csname CGTip@eq:cell-mass-W\endcsname{Eq. (13.22), p. 52: c\_n(lambda)=m\_n-2 Resum\_k>=1W\_n(gamma\_k(lambda)), d\_lambda c\_n(lambda)=2sum\_k>=1eta\_k Im V\_n(gamma\_k(lambda)),}
\expandafter\gdef\csname CGTip@eq:cell-mass-power-two\endcsname{Eq. (13.23), p. 52: -(mu(j))/(j)int\_2\textasciicircum{}3 J\_lambda(y) (d)/(dy)varphi\_n(y\textasciicircum{}j) dy,}
\expandafter\gdef\csname CGTip@lem:mean-value-ordinates\endcsname{Lemma 13.10 (Mean value over the moving ordinates) (p. 52): There is C such that for all lambda in [0,1], T>=2, u\_0 in R and complex a\_k with |a\_k|<=1, int\_u\_0-1\textasciicircum{}u\_0+1|sum\_gamma\_k(lambda)<= Ta\_k e\textasciicircum{}igamma\_k(lambda)u|\textasciicircum{}2du<= C Tlog\textasciicircum{}2T .}
\expandafter\gdef\csname CGTip@eq:mean-value-ordinates\endcsname{Eq. (13.24), p. 52: int\_u\_0-1\textasciicircum{}u\_0+1|sum\_gamma\_k(lambda)<= Ta\_k e\textasciicircum{}igamma\_k(lambda)u|\textasciicircum{}2du<= C Tlog\textasciicircum{}2T .}
\expandafter\gdef\csname CGTip@eq:fejer-pairs\endcsname{Eq. (13.25), p. 53: <<no.(k,k'):gamma\_k(lambda)<= T, |gamma\_k(lambda)-gamma\_k'(lambda)|<=1.}
\expandafter\gdef\csname CGTip@thm:mean-square-cells\endcsname{Theorem 13.11 (Mean square of the cell masses along the family) (p. 53): For every eps>0 there is C\_eps such that for all X>=2 and all lambda in [0,1] sum\_X<n<=2X| c\_n(lambda)|\textasciicircum{}2<= C\_eps X\textasciicircum{}1+eps, sum\_X<n<=2X|d\_lambda c\_n(lambda)|\textasciicircum{}2<= C\_eps X\textasciicircum{}1+eps. In particular the same bound already includes lambda=0: sum\_X<n<=2X| c\_n(0)|\textasciicircum{}2<= C\_eps X\textasciicircum{}1+eps.}
\expandafter\gdef\csname CGTip@eq:mean-square-cells\endcsname{Eq. (13.26), p. 53: sum\_X<n<=2X| c\_n(lambda)|\textasciicircum{}2<= C\_eps X\textasciicircum{}1+eps, sum\_X<n<=2X|d\_lambda c\_n(lambda)|\textasciicircum{}2<= C\_eps X\textasciicircum{}1+eps.}
\expandafter\gdef\csname CGTip@eq:A-nj\endcsname{Eq. (13.27), p. 53: A\_n,j=sum\_kint\_2\textasciicircum{}infvarphi\_n(y\textasciicircum{}j) (y\textasciicircum{}-1/2+igamma\_k(lambda))/(log y) dy , 1<= j<= J:=floor((log(2X+1))/(log2)) .}
\expandafter\gdef\csname CGTip@eq:cell-CS\endcsname{Eq. (13.28), p. 53: |intvarphi\_n(y\textasciicircum{}j)S\_n,j,lambda(y) dy|\textasciicircum{}2<=|varphi|\_inf\textasciicircum{}2 |I\_n,j|int\_I\_n,j|S\_n,j,lambda|\textasciicircum{}2dy .}
\expandafter\gdef\csname CGTip@eq:cell-block\endcsname{Eq. (13.29), p. 53: sum\_X<n<=2X|intvarphi\_n(y\textasciicircum{}j)S dy|\textasciicircum{}2<< X\textasciicircum{}1/j-1((j)/(log X))\textasciicircum{}2int\_u\_0-1\textasciicircum{}u\_0+1|sum\_gamma\_k(lambda)<= Te\textasciicircum{}igamma\_k(lambda)u|\textasciicircum{}2du<< j\textasciicircum{}2X\textasciicircum{}1/j+2eps.}
\expandafter\gdef\csname CGTip@cor:abscissa-one\endcsname{Corollary 13.12 (Abscissa 1 for every member) (p. 54): For every lambda in [0,1] and sigma>1, sum\_n>=2| c\_n(lambda)| n\textasciicircum{}-sigma<inf ; hence zeta\_Z,lambda and 1/ zeta\_Z,lambda are absolutely convergent Dirichlet series in Re s>1, for every lambda, including the endpoint lambda=0. Moreover, uniformly in lambda, no.n<= X:| c\_n(lambda)|>n\textasciicircum{}delta<<\_eps X\textasciicircum{}1-2de...}
\expandafter\gdef\csname CGTip@eq:abscissa-one\endcsname{Eq. (13.30), p. 54: sum\_n>=2| c\_n(lambda)| n\textasciicircum{}-sigma<inf ;}
\expandafter\gdef\csname CGTip@subsec:atomized-continuation\endcsname{Section 13.5: Analytic continuation of the atomized family (p. 54)}
\expandafter\gdef\csname CGTip@thm:atomized-continuation\endcsname{Theorem 13.13 (Continuation of the atomized family) (p. 54): Let lambda in [0,1]. For every c>(1)/(2), (pi\_lambda-L) x\textasciicircum{}-c in L\textasciicircum{}2((2,inf),dx/x). | c\_n(lambda)|<= C(1+I\_n\textasciicircum{}1/2) with I\_n as below, and sum\_n| c\_n(lambda)|n\textasciicircum{}-sigma<inf for sigma>(3)/(2); by Corollary CGrefcor:abscissa-one this holds for sigma>1, and zeta\_Z,lambda is an absolutely convergent Diric...}
\expandafter\gdef\csname CGTip@eq:cell-defect-bound\endcsname{Eq. (13.31), p. 54: |h\_n(lambda;s)|<= C|s| n\textasciicircum{}-sigma-1(1+I\_n\textasciicircum{}1/2), sum\_nI\_n\textasciicircum{}1/2n\textasciicircum{}-sigma-1<inf (sigma>c).}
\expandafter\gdef\csname CGTip@thm:closing\endcsname{Theorem 13.14 (Closing criterion; via Hamburger's theorem [ref]) (p. 55): The following are equivalent. (D), that is, zeta\_Z1=zeta. h\_0==0. J\_Z1=0 on (1,2), and for every n>=2 and every Re s>1, int\_0\textasciicircum{}1d\_\_n(lambda;s) dlambda=h\_n(1;s). In this case c\_n(0)=1\_ prime(n), a\_N(0)=1 for all N, and varpi\_Z1=1.}
\expandafter\gdef\csname CGTip@eq:closing-cells\endcsname{Eq. (13.32), p. 55: int\_0\textasciicircum{}1d\_\_n(lambda;s) dlambda=h\_n(1;s).}
\expandafter\gdef\csname CGTip@rem:A6-scale\endcsname{Remark 13.15 (The pointwise size of the cell masses) (p. 55): Theorem CGrefthm:closing needs no pointwise bound on c\_n(0): its proof uses the cell masses only through CGeqrefeq:abscissa-one. The bound c\_n(0)=O(n\textasciicircum{}eps) for every n is not proved directly; Corollary CGrefcor:abscissa-one restricts the exceptions to a set of n<= X of size O(X\textasciicircum{}1-2delta+eps) for t...}
\expandafter\gdef\csname CGTip@subsec:atomized-endpoints\endcsname{Section 13.6: The two endpoints of the atomized family (p. 55)}
\expandafter\gdef\csname CGTip@prop:atomized-clock\endcsname{Proposition 13.16 (Atomization at the clock endpoint) (p. 55): At lambda=1: c\_n(1)=O(1), and zeta\_Z,1 is an absolutely convergent Dirichlet series in Re s>1; h\_1(sigma)->-logvarpi\_RVM=log L as sigma->+inf; this limit is !=0 exactly when varpi\_RVM!=1, in particular under Hypothesis CGrefhyp:varpi-less-one; h\_1 is not constant: for every sigma>1, zeta\_RVM(sigm...}
\expandafter\gdef\csname CGTip@eq:duality\endcsname{Eq. (13.33), p. 55: zeta\_lambda\_functional equation zeros (1)/(2)+- igamma\_k(lambda)* e\textasciicircum{} h\_lambda = zeta\_Z,lambda\_Dirichlet series Euler product, h\_1 not constant, (D) h\_0==0 .}
\expandafter\gdef\csname CGTip@eq:atomized-mobius\endcsname{Eq. (13.34), p. 55: (1)/( zeta\_Z,lambda(s))=prod\_n>=2(1-n\textasciicircum{}-s)\textasciicircum{} c\_n(lambda)=sum\_N>=1( mu\_lambda(N))/(N\textasciicircum{}s), mu\_lambda(N)=sum\_prod\_nn\textasciicircum{}r\_n=N prod\_n: r\_n>=1(-1)\textasciicircum{}r\_n c\_n(lambda)r\_n,}
\expandafter\gdef\csname CGTip@eq:atomized-mertens\endcsname{Eq. (13.35), p. 56: M\_lambda(x):=sum\_N<= x mu\_lambda(N), (1)/( zeta\_Z,lambda)=e\textasciicircum{}- h\_lambdaF\_lambda ,}
\expandafter\gdef\csname CGTip@prop:Mellin-Mertens\endcsname{Proposition 13.17 (Mellin form of the endpoint Mertens function) (p. 56): There are c\_1>=2 and a unique measurable B\_Z1 on (0,inf), vanishing on (0,1), with x\textasciicircum{}-cB\_Z1 in L\textasciicircum{}2((1,inf),dx/x) and F\_Z1(s)=sint\_1\textasciicircum{}inf B\_Z1(x) x\textasciicircum{}-s-1 dx=exp(-int\_[1,inf)x\textasciicircum{}-s dJ\textasciicircum{} M\_Z1(x)) (Re s>c\_1), for every c>c\_1. Moreover (D) holds if and only if B\_Z1=M almost everywhere; and if F\_Z1 has an a...}
\expandafter\gdef\csname CGTip@eq:Mellin-Mertens\endcsname{Eq. (13.36), p. 56: F\_Z1(s)=sint\_1\textasciicircum{}inf B\_Z1(x) x\textasciicircum{}-s-1 dx=exp(-int\_[1,inf)x\textasciicircum{}-s dJ\textasciicircum{} M\_Z1(x)) (Re s>c\_1),}
\expandafter\gdef\csname CGTip@subsec:hamburger-table\endcsname{Section 13.7: Hamburger's hypotheses at the two endpoints (p. 56)}
\expandafter\gdef\csname CGTip@eq:complement\endcsname{Eq. (13.37), p. 58: zeta\_lambda\_keeps the functional equation and the zeros on Re s=(1)/(2) zeta\_Z,lambda\_keeps the Dirichlet series and the Euler product, zeta\_Z,lambda=zeta\_lambda e\textasciicircum{} h\_lambda.}
\expandafter\gdef\csname CGTip@eq:closure-common\endcsname{Eq. (13.38), p. 58: h\_0==0 J\_Z1=0 on (1,2), h\_n(0;s)=0 for all n>=2, Re s>1, and varpi\_Z1=1 .}
\expandafter\gdef\csname CGTip@subsec:programme\endcsname{Section 13.8: Programme (p. 58)}
\expandafter\gdef\csname CGTip@rem:input-A7\endcsname{Remark 13.18 (Input for A7) (p. 59): Steps A1--A6 use only two properties of an ordinate list: bounded displacement from the clock and a primitive of the count difference of size O(logtau). The list (tau\_k) has both properties, and for it the corresponding function zeta\_RVM is not a Dirichlet series (Theorem CGrefthm:no-Dirichlet-RV...}
\expandafter\gdef\csname CGTip@sec:criteria\endcsname{Section 14: Criteria for the final step and the clock end of the family (p. 59)}
\expandafter\gdef\csname CGTip@eq:cell-moments\endcsname{Eq. (14.1), p. 60: varphi(xi):=int\_Rvarphi(t)e\textasciicircum{}-2pi ixi t dt, M\_r:=int\_Rt\textasciicircum{}rvarphi(t) dt, C\_r:=int\_Rt\textasciicircum{}rvarphi(t)cos(2pi t) dt .}
\expandafter\gdef\csname CGTip@eq:partition-fourier\endcsname{Eq. (14.2), p. 60: varphi(q)=int\_0\textasciicircum{}1sum\_m in Zvarphi(t-m) e\textasciicircum{}-2pi iqt dt=delta\_q0 (q in Z),}
\expandafter\gdef\csname CGTip@subsec:open-path\endcsname{Section 14.1: The open path (p. 60)}
\expandafter\gdef\csname CGTip@prop:nonconstant\endcsname{Proposition 14.1 (No member of the open path is a Dirichlet series) (p. 60): Let 0<lambda<=1. Then zeta\_lambda is not a constant multiple of a Dirichlet series sum\_na\_nn\textasciicircum{}-s converging absolutely in Re s>1. In particular the joining factor h\_lambda is not constant.}
\expandafter\gdef\csname CGTip@subsec:convexity\endcsname{Section 14.2: Convexity along the family (p. 60)}
\expandafter\gdef\csname CGTip@prop:convexity\endcsname{Proposition 14.2 (Convexity on the real axis) (p. 60): For every real sigma>1 the function lambda\_lambda(sigma) is strictly convex on [0,1]: d\_lambda\textasciicircum{}2logzeta\_lambda(sigma)=sum\_k>=1eta\_k\textasciicircum{}2(C(sigma-(1)/(2),gamma\_k(lambda))-C((1)/(2),gamma\_k(lambda)))>0, C(a,gamma):=(2(a\textasciicircum{}2-gamma\textasciicircum{}2))/((a\textasciicircum{}2+gamma\textasciicircum{}2)\textasciicircum{}2). Consequently, with the generator V\_lambda of Theore...}
\expandafter\gdef\csname CGTip@eq:convexity-second\endcsname{Eq. (14.3), p. 60: d\_lambda\textasciicircum{}2logzeta\_lambda(sigma)=sum\_k>=1eta\_k\textasciicircum{}2(C(sigma-(1)/(2),gamma\_k(lambda))-C((1)/(2),gamma\_k(lambda)))>0, C(a,gamma):=(2(a\textasciicircum{}2-gamma\textasciicircum{}2))/((a\textasciicircum{}2+gamma\textasciicircum{}2)\textasciicircum{}2).}
\expandafter\gdef\csname CGTip@eq:convexity-sandwich\endcsname{Eq. (14.4), p. 60: logzeta\_RVM(sigma)-V\_1(sigma)<logzeta\_Z1(sigma)<logzeta\_RVM(sigma)-V\_0(sigma) (sigma>1).}
\expandafter\gdef\csname CGTip@subsec:pair\endcsname{Section 14.3: Hamburger's theorem for the pair (p. 61)}
\expandafter\gdef\csname CGTip@thm:pair-criterion\endcsname{Theorem 14.3 (Hamburger's theorem for the pair (zeta\_lambda, zeta\_Z,lambda)) (p. 61): Let lambda in [0,1]. Suppose that the joining factor e\textasciicircum{} h\_lambda extends to a meromorphic function on C of finite order with finitely many poles, and that e\textasciicircum{} h\_lambda(s)=e\textasciicircum{} h\_lambda(1-s) . Then zeta\_Z,lambda=zeta. For lambda=0, moreover, zeta\_Z1=zeta and h\_0==0.}
\expandafter\gdef\csname CGTip@eq:pair-symmetry\endcsname{Eq. (14.5), p. 61: e\textasciicircum{} h\_lambda(s)=e\textasciicircum{} h\_lambda(1-s) .}
\expandafter\gdef\csname CGTip@thm:double-hamburger\endcsname{Theorem 14.4 (Real-variable and vertical forms of Hamburger's theorem for the pair) (p. 61): For lambda in [0,1], y>0 and sigma>1 put T\_lambda(y):=sum\_N>=1 a\_N(lambda) e\textasciicircum{}-pi N\textasciicircum{}2y-y\textasciicircum{}-1/2sum\_N>=1 a\_N(lambda) e\textasciicircum{}-pi N\textasciicircum{}2/y-(1)/(2)(y\textasciicircum{}-1/2-1), k\_n(lambda):=sum\_d| n a\_d(lambda) mu(n/d), V\_sigma(lambda):=sum\_n>=2|k\_n(lambda)|\textasciicircum{}2 n\textasciicircum{}-2sigma. zeta\_Z,lambda/zeta=sum\_n>=1k\_n(lambda)n\textasciicircum{}-s converges absol...}
\expandafter\gdef\csname CGTip@eq:theta-defect\endcsname{Eq. (14.6), p. 61: T\_lambda(y):=sum\_N>=1 a\_N(lambda) e\textasciicircum{}-pi N\textasciicircum{}2y-y\textasciicircum{}-1/2sum\_N>=1 a\_N(lambda) e\textasciicircum{}-pi N\textasciicircum{}2/y-(1)/(2)(y\textasciicircum{}-1/2-1),}
\expandafter\gdef\csname CGTip@eq:vertical-defect\endcsname{Eq. (14.7), p. 61: k\_n(lambda):=sum\_d| n a\_d(lambda) mu(n/d), V\_sigma(lambda):=sum\_n>=2|k\_n(lambda)|\textasciicircum{}2 n\textasciicircum{}-2sigma.}
\expandafter\gdef\csname CGTip@eq:vertical-mean\endcsname{Eq. (14.8), p. 61: lim\_T->inf(1)/(2T)int\_-T\textasciicircum{}T|( zeta\_Z,lambda(sigma+it))/(zeta(sigma+it))|\textasciicircum{}2dt=1+ V\_sigma(lambda) >=1 .}
\expandafter\gdef\csname CGTip@eq:double-hamburger-limit\endcsname{Eq. (14.9), p. 61: (D) T\_0==0 V\_sigma(0)=0\_lambda->0\textasciicircum{}+T\_lambda(y)=0 (y>0)\_lambda->0\textasciicircum{}+ V\_sigma(lambda)=0 .}
\expandafter\gdef\csname CGTip@eq:quotient-three\endcsname{Eq. (14.10), p. 62: ( zeta\_Z,lambda)/(zeta)=e\textasciicircum{} h\_lambda (zeta\_lambda)/(zeta)}
\expandafter\gdef\csname CGTip@prop:clock-asymmetry\endcsname{Proposition 14.5 (At the clock end the symmetry fails) (p. 62): The joining factor e\textasciicircum{} h\_1 has no meromorphic continuation to C that satisfies CGeqrefeq:pair-symmetry.}
\expandafter\gdef\csname CGTip@subsec:positivity\endcsname{Section 14.4: Positivity and the value at infinity (p. 62)}
\expandafter\gdef\csname CGTip@lem:summation-positive\endcsname{Lemma 14.6 (Summation formula for a representing measure) (p. 62): Let lambda in [0,1] and let A be a positive measure on [1,inf) with zeta\_lambda(s)=int\_[1,inf)u\textasciicircum{}-s d A(u) (Re s>1). Then for every even Schwartz function g g(0)+2int\_[1,inf) g d A= g(0)+2int\_[1,inf) g d A .}
\expandafter\gdef\csname CGTip@eq:representing-measure\endcsname{Eq. (14.11), p. 62: zeta\_lambda(s)=int\_[1,inf)u\textasciicircum{}-s d A(u) (Re s>1).}
\expandafter\gdef\csname CGTip@eq:summation-positive\endcsname{Eq. (14.12), p. 62: g(0)+2int\_[1,inf) g d A= g(0)+2int\_[1,inf) g d A .}
\expandafter\gdef\csname CGTip@eq:tate-local\endcsname{Eq. (14.13), p. 62: pi\textasciicircum{}-s/2Gamma((s)/(2)) g\textasciicircum{}(1-s)=pi\textasciicircum{}-(1-s)/2Gamma((1-s)/(2)) g\textasciicircum{}(s)}
\expandafter\gdef\csname CGTip@eq:summation-shift\endcsname{Eq. (14.14), p. 63: int g d A=(1)/(2) g(0)-(1)/(2) g(0)+(1)/(2pi i)int\_(1-c)zeta\_lambda(s) g\textasciicircum{}(s) ds .}
\expandafter\gdef\csname CGTip@thm:positivity\endcsname{Theorem 14.7 (Positivity criterion) (p. 63): Let lambda in [0,1]. If zeta\_lambda has a representation CGeqrefeq:representing-measure with a positive measure A, then varpi\_lambda= A(1)<=1 . In particular, if zeta\_Z1 has such a representation, then varpi\_Z1=1 and zeta\_Z1=zeta.}
\expandafter\gdef\csname CGTip@eq:positivity-bound\endcsname{Eq. (14.15), p. 63: varpi\_lambda= A(1)<=1 .}
\expandafter\gdef\csname CGTip@eq:extremal-function\endcsname{Eq. (14.16), p. 63: g(u):=((sinpi u)/(pi u))\textasciicircum{}4(1-u\textasciicircum{}2), g(xi)= B(xi)+(1)/(4pi\textasciicircum{}2) B''(xi) .}
\expandafter\gdef\csname CGTip@eq:lp-inequality\endcsname{Eq. (14.17), p. 63: 1 >= g\_epsilon(0)+2int g\_epsilon d A= g\_epsilon(0)+2int g\_epsilon d A >= g\_epsilon(0)+2varpi\_lambda g\_epsilon(1) .}
\expandafter\gdef\csname CGTip@eq:lp-value\endcsname{Eq. (14.18), p. 63: varpi\_lambda<=(1- g(0))/(2 g(1))=((1)/(3)+(1)/(2pi\textasciicircum{}2))/((1)/(3)+(1)/(2pi\textasciicircum{}2))=1 .}
\expandafter\gdef\csname CGTip@cor:positivity\endcsname{Corollary 14.8 (Positive form of the final step) (p. 63): The following are equivalent: up(i) (D); up(ii) zeta\_Z1 is the Mellin--Stieltjes transform CGeqrefeq:representing-measure of a positive measure on [1,inf); up(iii) J\textasciicircum{} M\_Z1, extended by 0 to (0,1], is non-decreasing on (0,inf).}
\expandafter\gdef\csname CGTip@lem:summation-all\endcsname{Lemma 14.9 (Mellin form and summation formula for every member) (p. 64): There is c\_1>=2 such that for every lambda in [0,1] there is a unique measurable function A\_lambda on (0,inf), vanishing on (0,1), with x\textasciicircum{}-c A\_lambda in L\textasciicircum{}2((1,inf),dx/x) for every c>c\_1 and zeta\_lambda(s)=sint\_1\textasciicircum{}inf A\_lambda(x) x\textasciicircum{}-s-1 dx (Re s>c\_1). If zeta\_lambda has a representation CGeqrefeq:...}
\expandafter\gdef\csname CGTip@eq:mellin-A\endcsname{Eq. (14.19), p. 64: zeta\_lambda(s)=sint\_1\textasciicircum{}inf A\_lambda(x) x\textasciicircum{}-s-1 dx (Re s>c\_1).}
\expandafter\gdef\csname CGTip@eq:summation-all\endcsname{Eq. (14.20), p. 64: g(0)+2 A\_lambda, g= g(0)+2 A\_lambda, g .}
\expandafter\gdef\csname CGTip@eq:parseval-A\endcsname{Eq. (14.21), p. 64: A\_lambda, g=(1)/(2pi i)int\_(c)zeta\_lambda(s) g\textasciicircum{}(s) ds .}
\expandafter\gdef\csname CGTip@prop:one-number\endcsname{Proposition 14.10 (The value at infinity as one number) (p. 64): Let g\_epsilon be as in the proof of Theorem CGrefthm:positivity, w\_epsilon:= g\_epsilon- g\_epsilon, which is non-negative on |u|>=1, and A'\_lambda(x):= A\_lambda(x)-varpi\_lambda for x>=1. For every lambda in [0,1] and epsilon>0, varpi\_lambda w\_epsilon(1)=(1)/(2)(1- g\_epsilon(0))- A'\_lambda, w\_epsil...}
\expandafter\gdef\csname CGTip@eq:one-number-eps\endcsname{Eq. (14.22), p. 64: varpi\_lambda w\_epsilon(1)=(1)/(2)(1- g\_epsilon(0))- A'\_lambda, w\_epsilon, A'\_lambda, w\_epsilon:=-int\_1\textasciicircum{}inf A'\_lambda(x) w\_epsilon'(x) dx,}
\expandafter\gdef\csname CGTip@eq:one-number\endcsname{Eq. (14.23), p. 64: varpi\_lambda-1=-( W(lambda))/( g(1)), g(1)=(1)/(6)+(1)/(4pi\textasciicircum{}2).}
\expandafter\gdef\csname CGTip@thm:real-axis\endcsname{Theorem 14.11 (The endpoint dominates zeta on the real axis) (p. 64): For every real sigma>1, zeta\_Z1(sigma) >= zeta(sigma), and equality for one sigma>1 holds if and only if zeta\_Z1=zeta. In particular prod\_k>=1(1+(2)/((1)/(4)+gamma\_k\textasciicircum{}2)) >= (pi)/(3), with equality if and only if (D).}
\expandafter\gdef\csname CGTip@eq:real-axis\endcsname{Eq. (14.24), p. 64: zeta\_Z1(sigma) >= zeta(sigma),}
\expandafter\gdef\csname CGTip@eq:sigma-two\endcsname{Eq. (14.25), p. 64: prod\_k>=1(1+(2)/((1)/(4)+gamma\_k\textasciicircum{}2)) >= (pi)/(3), with equality if and only if (D).}
\expandafter\gdef\csname CGTip@eq:pair-monotone\endcsname{Eq. (14.26), p. 65: (2)/(W)-(2(W-d\_rho\textasciicircum{}2))/((W-d\_rho\textasciicircum{}2)\textasciicircum{}2+4d\_rho\textasciicircum{}2gamma\textasciicircum{}2)=(2d\_rho\textasciicircum{}2 (4gamma\textasciicircum{}2+d\_rho\textasciicircum{}2-W))/(W((W-d\_rho\textasciicircum{}2)\textasciicircum{}2+4d\_rho\textasciicircum{}2gamma\textasciicircum{}2)),}
\expandafter\gdef\csname CGTip@prop:open-path-varpi\endcsname{Proposition 14.12 (The final step on the open path) (p. 65): Assume Hypothesis CGrefhyp:varpi-less-one, varpi\_RVM<1. Then the following are equivalent: up(i) (D); up(ii) varpi\_lambda<1 for every lambda in (0,1]; up(iii) varpi\_lambda\_m<=1 for some sequence lambda\_m->0\textasciicircum{}+; up(iv) W(lambda)>0 for every lambda in (0,1]. Without the hypothesis, (ii)(iii)(i) and ...}
\expandafter\gdef\csname CGTip@rem:hypothesis-X\endcsname{Remark 14.13 (The hypothesis added to the functional equation) (p. 65): The functional equation, the order and the zeros on the critical line do not characterize zeta: every zeta\_lambda has them. A converse theorem for zeta\_Z1 needs one further hypothesis on the function A\_Z1 of Lemma CGreflem:summation-all, and each of the following suffices and is equivalent to (D)...}
\expandafter\gdef\csname CGTip@subsec:clock-density\endcsname{Section 14.5: The density of the clock (p. 65)}
\expandafter\gdef\csname CGTip@eq:U-def\endcsname{Eq. (14.27), p. 65: U(M,t\_0):=int\_t\_0\textasciicircum{}infe\textasciicircum{}-iM(t) dt:=lim\_T->infint\_t\_0\textasciicircum{}Te\textasciicircum{}-iM(t) dt .}
\expandafter\gdef\csname CGTip@lem:stationary-U\endcsname{Lemma 14.14 (Stationary phase for) (p. 65): The limit CGeqrefeq:U-def exists, and there is an absolute constant C\_U with |U(M,t\_0)-sqrt((2pi)/(M)) e\textasciicircum{}i(M-pi/4)|<=(C\_U)/(M) (M>=1, 0<t\_0<=(1)/(2)).}
\expandafter\gdef\csname CGTip@eq:U-asymptotic\endcsname{Eq. (14.28), p. 65: |U(M,t\_0)-sqrt((2pi)/(M)) e\textasciicircum{}i(M-pi/4)|<=(C\_U)/(M) (M>=1, 0<t\_0<=(1)/(2)).}
\expandafter\gdef\csname CGTip@eq:U-substitution\endcsname{Eq. (14.29), p. 66: p( w):=(dt)/(d w)=( w)/(log t), p(0)=1, q( w):=( p( w)-1)/( w)=(1)/(log t)-(1)/( w).}
\expandafter\gdef\csname CGTip@eq:U-split\endcsname{Eq. (14.30), p. 66: e\textasciicircum{}-iM U(M,t\_0)=int\_ w\_0\textasciicircum{}inf e\textasciicircum{}-iM w\textasciicircum{}2/2 d w+int\_ w\_0\textasciicircum{}inf e\textasciicircum{}-iM w\textasciicircum{}2/2 w q( w) d w .}
\expandafter\gdef\csname CGTip@eq:clock-harmonic\endcsname{Eq. (14.31), p. 66: R\_m(y):=-2(-1)\textasciicircum{}m m\textasciicircum{}-3/2 y\textasciicircum{}(1)/(2m)-(1)/(2)cos(2pi m y\textasciicircum{}1/m-((7m+1)pi)/(4)),}
\expandafter\gdef\csname CGTip@thm:clock-density\endcsname{Theorem 14.15 (The density of the clock) (p. 66): J\_RVM is continuously differentiable on (2,inf), and for y>=16 J\_RVM'(y)=(1)/(log y)+2cos(2pi y)-(y\textasciicircum{}-1/4)/(2)cos(4pi y+(pi)/(4))+O(y\textasciicircum{}-1/3). On [2,inf) the measure dpi\_RVM is the sum of the atoms (mu(j))/(j)J\_RVM(2) delta\_2\textasciicircum{}j, j>=1, and of a density that is continuous off the points 2\textasciicircum{}j and equals...}
\expandafter\gdef\csname CGTip@eq:clock-density\endcsname{Eq. (14.32), p. 66: J\_RVM'(y)=(1)/(log y)+2cos(2pi y)-(y\textasciicircum{}-1/4)/(2)cos(4pi y+(pi)/(4))+O(y\textasciicircum{}-1/3).}
\expandafter\gdef\csname CGTip@eq:H-harmonic\endcsname{Eq. (14.33), p. 66: H\_m\textasciicircum{}+-(v):=int\_tau\_0\textasciicircum{}inf e\textasciicircum{}iPsi\_m\textasciicircum{}+-(tau) dtau ,}
\expandafter\gdef\csname CGTip@eq:S-zero\endcsname{Eq. (14.34), p. 66: sum\_kD\_k=int\_u\_1\textasciicircum{}u\_2(e\textasciicircum{}v/2)/(v) S\_0(v) dv, S\_0(v):=S\_ sm(v)+sum\_m>=1((-1)\textasciicircum{}m v)/(2pi m)(H\_m\textasciicircum{}-(v)-H\_m\textasciicircum{}+(v)),}
\expandafter\gdef\csname CGTip@eq:J-prime-exact\endcsname{Eq. (14.35), p. 66: J\_RVM'(y)=(1)/(log y)(1-(1)/(y(y\textasciicircum{}2-1))-(2)/( y) Re S\_0(log y)).}
\expandafter\gdef\csname CGTip@eq:H-by-U\endcsname{Eq. (14.36), p. 67: H\_m\textasciicircum{}-(v)= t\_m e\textasciicircum{}-7pi im/4 U(m t\_m,tau\_0/ t\_m), tau\_0/ t\_m=2e\textasciicircum{}-v/m<=(1)/(2) (y>=4\textasciicircum{}m).}
\expandafter\gdef\csname CGTip@eq:H-asymptotic\endcsname{Eq. (14.37), p. 67: ((-1)\textasciicircum{}mv)/(2pi m) H\_m\textasciicircum{}-(v)=(-1)\textasciicircum{}m (v)/(m\textasciicircum{}3/2) y\textasciicircum{}(1)/(2m)exp(i(2pi m y\textasciicircum{}1/m-((7m+1)pi)/(4)))+O((v)/(m\textasciicircum{}2)).}
\expandafter\gdef\csname CGTip@cor:clock-cells\endcsname{Corollary 14.16 (Cell masses and cell defects at the clock end) (p. 67): For n>=17, with 1\_2\textasciicircum{}j(n)=1 if n=2\textasciicircum{}j and 0 otherwise, c\_n\textasciicircum{}RVM=int\_n-1\textasciicircum{}n((1)/(log x)+2cos(2pi x))dx+1\_2\textasciicircum{}j(n) (mu(j))/(j) J\_RVM(2)+O(n\textasciicircum{}-1/4), c\_n(1)=(1)/(log n)+1\_2\textasciicircum{}j(n) (mu(j))/(j) J\_RVM(2)+O(n\textasciicircum{}-1/4), and for 0<Re s<=sigma\_1 |h\_n(1;s)|<= C\_sigma\_1(|s| n\textasciicircum{}-sigma-4/3+|s|(1+|s|) n\textasciicircum{}-sigma-2). In particu...}
\expandafter\gdef\csname CGTip@eq:clock-cell-sharp\endcsname{Eq. (14.38), p. 67: c\_n\textasciicircum{}RVM=int\_n-1\textasciicircum{}n((1)/(log x)+2cos(2pi x))dx+1\_2\textasciicircum{}j(n) (mu(j))/(j) J\_RVM(2)+O(n\textasciicircum{}-1/4),}
\expandafter\gdef\csname CGTip@eq:clock-cell-smooth\endcsname{Eq. (14.39), p. 67: c\_n(1)=(1)/(log n)+1\_2\textasciicircum{}j(n) (mu(j))/(j) J\_RVM(2)+O(n\textasciicircum{}-1/4),}
\expandafter\gdef\csname CGTip@eq:clock-defect\endcsname{Eq. (14.40), p. 67: |h\_n(1;s)|<= C\_sigma\_1(|s| n\textasciicircum{}-sigma-4/3+|s|(1+|s|) n\textasciicircum{}-sigma-2).}
\expandafter\gdef\csname CGTip@subsec:height-kernels\endcsname{Section 14.6: The cell masses in the heights of the ordinates (p. 68)}
\expandafter\gdef\csname CGTip@prop:height-kernels\endcsname{Proposition 14.17 (Height kernel of the cell masses) (p. 68): For every eps>0, lambda in [0,1] and n>=3 that is not a power of 2, c\_n(lambda)=m\_n-(2)/( n log n)sum\_k>=1varphi((gamma\_k(lambda))/(2pi n))cos(gamma\_k(lambda)log n)+E\_n(lambda), where sum\_X<n<=2X|E\_n(lambda)|\textasciicircum{}2<= C\_eps X\textasciicircum{}1/2+eps uniformly in lambda (sum over the n that are not powers of 2). The k...}
\expandafter\gdef\csname CGTip@eq:height-kernel\endcsname{Eq. (14.41), p. 68: c\_n(lambda)=m\_n-(2)/( n log n)sum\_k>=1varphi((gamma\_k(lambda))/(2pi n))cos(gamma\_k(lambda)log n)+E\_n(lambda),}
\expandafter\gdef\csname CGTip@eq:Xi-rotation\endcsname{Eq. (14.42), p. 68: Xi\_n(Omega):=sum\_k>=1varphi((tau\_k)/(2pi n))n\textasciicircum{}itau\_ke\textasciicircum{}-iOmegaeta\_k (0<=Omega<=log n).}
\expandafter\gdef\csname CGTip@prop:flow-count\endcsname{Proposition 14.18 (The flow of a cell mass and the count difference) (p. 68): For every n>=3 that is not a power of 2, c\_n(0)- c\_n(1)=-2int\_0\textasciicircum{}1sum\_k>=1eta\_k Im V\_n(gamma\_k(lambda)) dlambda=-2 Imint\_tau\_0\textasciicircum{}inf V\_n(t)(N\_zeta(t)-n\_RVM(t)) dt .}
\expandafter\gdef\csname CGTip@eq:flow-count\endcsname{Eq. (14.43), p. 68: c\_n(0)- c\_n(1)=-2int\_0\textasciicircum{}1sum\_k>=1eta\_k Im V\_n(gamma\_k(lambda)) dlambda=-2 Imint\_tau\_0\textasciicircum{}inf V\_n(t)(N\_zeta(t)-n\_RVM(t)) dt .}
\expandafter\gdef\csname CGTip@prop:defect-kernel\endcsname{Proposition 14.19 (Height kernel of the cell defects) (p. 68): For every eps>0, lambda in [0,1], Re s>0 and n>=3 that is not a power of 2, h\_n(lambda;s)=-(s n\textasciicircum{}-s-3/2)/( n)sum\_k>=1varphi'((gamma\_k(lambda))/(2pi n))sin(gamma\_k(lambda)log n)+O\_s,eps(n\textasciicircum{}-sigma-3/4+eps).}
\expandafter\gdef\csname CGTip@eq:defect-kernel\endcsname{Eq. (14.44), p. 68: h\_n(lambda;s)=-(s n\textasciicircum{}-s-3/2)/( n)sum\_k>=1varphi'((gamma\_k(lambda))/(2pi n))sin(gamma\_k(lambda)log n)+O\_s,eps(n\textasciicircum{}-sigma-3/4+eps).}
\expandafter\gdef\csname CGTip@rem:scale-test\endcsname{Remark 14.20 (The scale of the displacements) (p. 69): By Proposition CGrefprop:nonconstant, a property of the sequence (eta\_k) that yields (D) is false for every sequence (c eta\_k) with 0<= c<1. Properties that are invariant under eta\_k c eta\_k with c>0 (the signs of the eta\_k, the non-vanishing of the eta\_k, their boundedness) hold along the whole ...}
\expandafter\gdef\csname CGTip@eq:eta-exact\endcsname{Eq. (14.45), p. 69: eta\_k N\_RVM'(xi\_k)=N\_zeta\textasciicircum{}*(gamma\_k)-N\_RVM(gamma\_k) for some xi\_k between gamma\_k and tau\_k .}
\expandafter\gdef\csname CGTip@subsec:gaussian-defect\endcsname{Section 14.7: The Gaussian defect of the atomization (p. 69)}
\expandafter\gdef\csname CGTip@lem:vertical-growth\endcsname{Lemma 14.21 (Vertical growth of the joining factor) (p. 69): For every lambda in [0,1] and (1)/(2)<sigma\_0<sigma\_1, h\_lambda(sigma+it)=o(|t|) as |t|->inf, uniformly for sigma\_0<=sigma<=sigma\_1.}
\expandafter\gdef\csname CGTip@eq:cell-RL\endcsname{Eq. (14.46), p. 69: h\_n(lambda;ms)= intvarphi\_n(x)(n\textasciicircum{}-ms-x\textasciicircum{}-ms)dL(x)-int E(x)varphi\_n'(x)(n\textasciicircum{}-ms-x\textasciicircum{}-ms)dx; -msint E(x)varphi\_n(x) x\textasciicircum{}-ms-1dx .}
\expandafter\gdef\csname CGTip@eq:gaussian-defect\endcsname{Eq. (14.47), p. 69: Theta\_lambda(y):=(1)/(4pi i)int\_(c)zeta\_lambda(s) Gamma((s)/(2)) pi\textasciicircum{}-s/2 y\textasciicircum{}-s/2 ds, G\_lambda(y):=sum\_N>=1 a\_N(lambda) e\textasciicircum{}-pi N\textasciicircum{}2y-Theta\_lambda(y).}
\expandafter\gdef\csname CGTip@prop:gaussian-defect\endcsname{Proposition 14.22 (Gaussian defect) (p. 69): Let lambda in [0,1]. Theta\_lambda(y)=y\textasciicircum{}-1/2 Theta\_lambda(1/y)+(1)/(2)(y\textasciicircum{}-1/2-1) for y>0, and Theta\_lambda(y)=O\_A(y\textasciicircum{}-A) as y->inf for every A. The coefficients a\_N(lambda) satisfy the theta relation of zeta, sum\_N>=1 a\_N(lambda) e\textasciicircum{}-pi N\textasciicircum{}2y=y\textasciicircum{}-1/2sum\_N>=1 a\_N(lambda) e\textasciicircum{}-pi N\textasciicircum{}2/y+(1)/(2)(y\textasciicircum{}-1/2-1) (...}
\expandafter\gdef\csname CGTip@eq:theta-atomized\endcsname{Eq. (14.48), p. 69: sum\_N>=1 a\_N(lambda) e\textasciicircum{}-pi N\textasciicircum{}2y=y\textasciicircum{}-1/2sum\_N>=1 a\_N(lambda) e\textasciicircum{}-pi N\textasciicircum{}2/y+(1)/(2)(y\textasciicircum{}-1/2-1) (y>0),}
\expandafter\gdef\csname CGTip@eq:defect-selfdual\endcsname{Eq. (14.49), p. 69: G\_lambda(y)=y\textasciicircum{}-1/2 G\_lambda(1/y) (y>0).}
\expandafter\gdef\csname CGTip@eq:defect-expansion\endcsname{Eq. (14.50), p. 70: G\_lambda(y)=(1)/(2)(e\textasciicircum{} h\_lambda(1)-1) y\textasciicircum{}-1/2+O\_delta(y\textasciicircum{}-1/4-delta),}
\expandafter\gdef\csname CGTip@cor:residue-condition\endcsname{Corollary 14.23 (The density of the atomized integers) (p. 70): sum\_N<= x a\_N(lambda) has mean density e\textasciicircum{} h\_lambda(1) in the sense of CGeqrefeq:defect-expansion, and (D) implies h\_0(1)=0, that is, u\_0(1)=sum\_m>=1(1)/(m)sum\_n>=2h\_n(0;m).}
\expandafter\gdef\csname CGTip@eq:residue-condition\endcsname{Eq. (14.51), p. 70: h\_0(1)=0, that is, u\_0(1)=sum\_m>=1(1)/(m)sum\_n>=2h\_n(0;m).}
\expandafter\gdef\csname CGTip@eq:defect-expansion-clock\endcsname{Eq. (14.52), p. 70: G\_1(y)=(1)/(2)(e\textasciicircum{} h\_1(1)-1) y\textasciicircum{}-1/2+O\_delta(y\textasciicircum{}-delta) (y->0\textasciicircum{}+),}
\expandafter\gdef\csname CGTip@subsec:entire-lambda\endcsname{Section 14.8: Entire dependence on lambdalambda and the exceptional set (p. 70)}
\expandafter\gdef\csname CGTip@prop:entire-lambda\endcsname{Proposition 14.24 (The cell masses as entire functions of lambda) (p. 70): Let n>=3 not be a power of 2. The real-axis function lambda c\_n(lambda) has the entire extension [ C\_n(lambda):=m\_n-sum\_k(W\_n(gamma\_k+lambdaeta\_k)+conj W\_n(gamma\_k+conj lambdaeta\_k)), ] which is of exponential type at most C\_(n+1) and agrees with c\_n(lambda) for real lambda. For every r>=0, eps>0...}
\expandafter\gdef\csname CGTip@eq:lambda-derivatives\endcsname{Eq. (14.53), p. 70: d\_lambda\textasciicircum{}r c\_n(lambda)=-2 Resum\_k>=1eta\_k\textasciicircum{} r W\_n\textasciicircum{}(r)(gamma\_k(lambda)), |d\_lambda\textasciicircum{}r c\_n(lambda)|<= C\_r,eps(C\_(n+1))\textasciicircum{}r n\textasciicircum{}1/2+eps.}
\expandafter\gdef\csname CGTip@prop:exceptional-set\endcsname{Proposition 14.25 (The exceptional set) (p. 71): The set E:=lambda in [0,1]: zeta\_Z,lambda=zeta is finite, 1 E, and (D) holds if and only if 0 in E. For every lambda E the hypotheses of Theorem CGrefthm:pair-criterion do not hold together, T\_lambda==0 and V\_sigma(lambda)>0.}
\expandafter\gdef\csname CGTip@eq:exceptional-set\endcsname{Eq. (14.54), p. 71: E:=lambda in [0,1]: zeta\_Z,lambda=zeta}
\expandafter\gdef\csname CGTip@subsec:cell-mass-form\endcsname{Section 14.9: The final step in the cell masses (p. 71)}
\expandafter\gdef\csname CGTip@cor:cell-mass-form\endcsname{Corollary 14.26 (The final step is one number per cell) (p. 71): The following are equivalent: (D); c\_n(0)=1\_ prime(n) for every n>=2; there is a list kappa in O (Section CGrefsec:lists) with c\_n(kappa)=1\_ prime(n) for every n>=2.}
\expandafter\gdef\csname CGTip@eq:cell-mass-euler\endcsname{Eq. (14.55), p. 71: zeta\_Z,kappa(s)=prod\_n>=2(1-n\textasciicircum{}-s)\textasciicircum{}- c\_n(kappa)=prod\_p(1-p\textasciicircum{}-s)\textasciicircum{}-1=zeta(s) (Re s>1),}
\expandafter\gdef\csname CGTip@prop:witness\endcsname{Proposition 14.27 (Every point of the exceptional set is a witness) (p. 72): For the set E of CGeqrefeq:exceptional-set, E!=(D), so that either E= or 0 in E . If (D) holds, then E0=lambda in (0,1]: D\_n(lambda)=0 for every n>=2, D\_n(lambda):= c\_n(lambda)- c\_n(0), and for n>=3 not a power of 2, the real-axis function D\_n(lambda)=-2Resum\_k>=1(W\_n(gamma\_k+lambdaeta\_k)-W\_n(gam...}
\expandafter\gdef\csname CGTip@eq:witness\endcsname{Eq. (14.56), p. 72: E!=(D), so that either E= or 0 in E .}
\expandafter\gdef\csname CGTip@eq:witness-zeros\endcsname{Eq. (14.57), p. 72: E0=lambda in (0,1]: D\_n(lambda)=0 for every n>=2, D\_n(lambda):= c\_n(lambda)- c\_n(0),}
\expandafter\gdef\csname CGTip@rem:dual-function\endcsname{Remark 14.28 (The dual function of the atomized family) (p. 72): For Re s<(1)/(2) put B\textasciicircum{}v\_lambda(s):=zeta\_lambda(s) e\textasciicircum{} h\_lambda(1-s). By the functional equation of zeta\_lambda, pi\textasciicircum{}-s/2Gamma((s)/(2))B\textasciicircum{}v\_lambda(s)=pi\textasciicircum{}-(1-s)/2Gamma((1-s)/(2)) zeta\_Z,lambda(1-s) (Re s<(1)/(2)), so B\textasciicircum{}v\_lambda is the function attached to zeta\_Z,lambda by the functional equation. To ...}
\expandafter\gdef\csname CGTip@eq:dual-function\endcsname{Eq. (14.58), p. 72: pi\textasciicircum{}-s/2Gamma((s)/(2))B\textasciicircum{}v\_lambda(s)=pi\textasciicircum{}-(1-s)/2Gamma((1-s)/(2)) zeta\_Z,lambda(1-s) (Re s<(1)/(2)),}
\expandafter\gdef\csname CGTip@prop:per-cell\endcsname{Proposition 14.29 (The identity that remains, cell by cell) (p. 72): Let n>=17 not be a power of 2. Then c\_n(0)=1\_ prime(n) if and only if 1\_ prime(n)- c\_n(1)=-2 Imint\_tau\_0\textasciicircum{}inf V\_n(t)(N\_zeta(t)-n\_RVM(t)) dt, c\_n(1)=(1)/(log n)+O(n\textasciicircum{}-1/4).}
\expandafter\gdef\csname CGTip@eq:per-cell\endcsname{Eq. (14.59), p. 72: 1\_ prime(n)- c\_n(1)=-2 Imint\_tau\_0\textasciicircum{}inf V\_n(t)(N\_zeta(t)-n\_RVM(t)) dt, c\_n(1)=(1)/(log n)+O(n\textasciicircum{}-1/4).}
\expandafter\gdef\csname CGTip@subsec:symmetric-factors\endcsname{Section 14.10: Symmetric joining factors (p. 72)}
\expandafter\gdef\csname CGTip@rem:pair-no-slack\endcsname{Remark 14.30 (The symmetry criterion has no slack at lambda=0) (p. 72): At lambda=0 the hypothesis of Theorem CGrefthm:pair-criterion is equivalent to its conclusion. If (D) holds, then h\_0==0 (Theorem CGrefthm:closing), and e\textasciicircum{} h\_0=1 is symmetric and of finite order; hence e\textasciicircum{} h\_0 symmetric of finite order h\_0==0(D), and in the same way T\_0==0 and V\_sigma(0)=0 are equ...}
\expandafter\gdef\csname CGTip@eq:pair-no-slack\endcsname{Eq. (14.60), p. 72: e\textasciicircum{} h\_0 symmetric of finite order h\_0==0(D),}
\expandafter\gdef\csname CGTip@prop:dichotomy\endcsname{Proposition 14.31 (Symmetric joining factors in general) (p. 73): Let f be meromorphic in C with (s-1)f(s) entire of order at most one, and let f satisfy the functional equation CGeqrefeq:Z1-FE. Let g=f e\textasciicircum{}h, where e\textasciicircum{}h is entire, without zeros, of finite order, and e\textasciicircum{}h(s)=e\textasciicircum{}h(1-s). If g==0 and g(s)=sum\_n>=1a\_nn\textasciicircum{}-s converges absolutely in Re s>1, then g=C\_Hzeta w...}
\expandafter\gdef\csname CGTip@rem:dichotomy-counterexample\endcsname{Remark 14.32 (Order one cannot be dropped) (p. 73): For real nu\_5!=0 put f(s)=zeta(s)exp(-nu\_5(s-(1)/(2))\textasciicircum{}2) and e\textasciicircum{}h(s)=exp(nu\_5(s-(1)/(2))\textasciicircum{}2). Then f is meromorphic of order two with the functional equation, the pole and the zeros of zeta; e\textasciicircum{}h is entire, without zeros, of finite order and symmetric; g=f e\textasciicircum{}h=zeta is an absolutely convergent Dirich...}
\expandafter\gdef\csname CGTip@subsec:halfline\endcsname{Section 14.11: The half-line form of Hamburger's theorem for a pair (p. 73)}
\expandafter\gdef\csname CGTip@thm:ladder\endcsname{Theorem 14.33 (The ladder) (p. 73): Let d[1,inf)->R be measurable with int\_1\textasciicircum{}inf| d(x)| x\textasciicircum{}- e\_1-1 dx<inf for some e\_1>=0, and put h(sigma):=\_1\textasciicircum{}inf d(x) x\textasciicircum{}-sigma-1 dx (sigma> e\_1). For every integer N>=2, h(sigma)=O(sigma N\textasciicircum{}-sigma) (sigma->+inf) d(x)=0 for almost every x in [1,N).}
\expandafter\gdef\csname CGTip@eq:ladder-h\endcsname{Eq. (14.61), p. 73: h(sigma):=\_1\textasciicircum{}inf d(x) x\textasciicircum{}-sigma-1 dx (sigma> e\_1).}
\expandafter\gdef\csname CGTip@eq:ladder\endcsname{Eq. (14.62), p. 73: h(sigma)=O(sigma N\textasciicircum{}-sigma) (sigma->+inf) d(x)=0 for almost every x in [1,N).}
\expandafter\gdef\csname CGTip@eq:ladder-Y\endcsname{Eq. (14.63), p. 73: Y(s):=int\_0\textasciicircum{}log Nd(u) e\textasciicircum{}s(log N-u) du .}
\expandafter\gdef\csname CGTip@eq:ladder-Y-real\endcsname{Eq. (14.64), p. 73: Y(sigma)=(N\textasciicircum{}sigma h(sigma))/(sigma)-int\_log N\textasciicircum{}inf d(u) e\textasciicircum{}-sigma(u-log N) du ,}
\expandafter\gdef\csname CGTip@thm:halfline-hamburger\endcsname{Theorem 14.34 (Hamburger's theorem for a pair, half-line form) (p. 74): Let f satisfy (H0) and the functional equation CGeqrefeq:Z1-FE, and suppose that for Re s>1 log f(s)=sint\_1\textasciicircum{}inf J\_f(x) x\textasciicircum{}-s-1 dx, int\_1\textasciicircum{}inf|J\_f(x)| x\textasciicircum{}-sigma-1 dx<inf (sigma>1), with J\_f real and measurable. Let g(s)=sum\_n>=1a\_nn\textasciicircum{}-s converge absolutely in Re s>1, with a\_1=1, and let J\_g(x):=sum\_2<...}
\expandafter\gdef\csname CGTip@eq:halfline-f\endcsname{Eq. (14.65), p. 74: log f(s)=sint\_1\textasciicircum{}inf J\_f(x) x\textasciicircum{}-s-1 dx, int\_1\textasciicircum{}inf|J\_f(x)| x\textasciicircum{}-sigma-1 dx<inf (sigma>1),}
\expandafter\gdef\csname CGTip@cor:ladder-pair\endcsname{Corollary 14.35 (The ladder for the pair (zeta\_lambda, zeta\_Z,lambda)) (p. 74): Let lambda in [0,1] and N>=2. Then h\_lambda(sigma)=O(sigma N\textasciicircum{}-sigma) as sigma->+inf if and only if the Mellin count J\textasciicircum{} M\_lambda of zeta\_lambda coincides almost everywhere on [1,N) with the step function sum\_n\textasciicircum{}m<= x c\_n(lambda)/m; in particular J\textasciicircum{} M\_lambda then vanishes on (1,2) and is constant be...}
\expandafter\gdef\csname CGTip@eq:ladder-D\endcsname{Eq. (14.66), p. 74: (D) h\_0(sigma)=O\_N(sigma N\textasciicircum{}-sigma) for every N>=2 ,}
\expandafter\gdef\csname CGTip@subsec:bridge\endcsname{Section 14.12: The bridge theorem: the phase of the joining factor on the critical line (p. 74)}
\expandafter\gdef\csname CGTip@thm:bridge\endcsname{Theorem 14.36 (Bridge theorem, phase form) (p. 74): Let f be meromorphic in C, real on the real axis, with (H0), the functional equation CGeqrefeq:Z1-FE, and lim\_sigma->+inff(sigma)=varpi\_f>0. Let g(s)=sum\_n>=1a\_nn\textasciicircum{}-s converge absolutely in Re s>1, with real a\_n and a\_1=1. Suppose that g=f e\textasciicircum{}h in Re s>(1)/(2) with h holomorphic there and real on (...}
\expandafter\gdef\csname CGTip@eq:bridge-B2\endcsname{Eq. (14.67), p. 74: |h(s)|<= C (1+|s|)\textasciicircum{} e\_1(1+(sigma-(1)/(2))\textasciicircum{}- e\_2) (Re s>(1)/(2)).}
\expandafter\gdef\csname CGTip@eq:bridge-mean\endcsname{Eq. (14.68), p. 75: log|h(s)|<=(1)/(2pi)int\_0\textasciicircum{}2pilog|h(s+e\textasciicircum{}iphi)| dphi<=log C+ e\_1log(2+|s|)+log2+( e\_2)/(2pi)int\_0\textasciicircum{}2pilog\textasciicircum{}+(1)/(|sigma-(1)/(2)+cosphi|) dphi .}
\expandafter\gdef\csname CGTip@cor:bridge-family\endcsname{Corollary 14.37 (The bridge theorem for the family) (p. 75): For 0<lambda<=1 the joining factor h\_lambda does not satisfy both up(B1) and up(B2). (D) holds if and only if h\_0 satisfies up(B1) and up(B2).}
\expandafter\gdef\csname CGTip@eq:phase-defect\endcsname{Eq. (14.69), p. 75: v\_lambda(t):=lim\_sigma->(1)/(2)\textasciicircum{}+Im h\_lambda(sigma+it)=lim\_sigma->(1)/(2)\textasciicircum{}+(Imlog zeta\_Z,lambda(sigma+it)-Imlogzeta\_lambda(sigma+it)).}
\expandafter\gdef\csname CGTip@prop:continuous-phase\endcsname{Proposition 14.38 (The phase of the continuous side) (p. 75): Let lambda in [0,1] and t>0 not an ordinate gamma\_k(lambda). With the branch of logzeta\_lambda that is real on (1,inf), continued along any path in the open quadrant Re s>(1)/(2), Im s>0 (the continuation is well defined there, since the zeros of zeta\_lambda lie on its boundary line Re s=(1)/(2) ...}
\expandafter\gdef\csname CGTip@eq:continuous-phase\endcsname{Eq. (14.70), p. 75: lim\_sigma->(1)/(2)\textasciicircum{}+Imlogzeta\_lambda(sigma+it)=-vartheta(t)+pi(N\_lambda(t)-1), N\_lambda(t)=no.k:gamma\_k(lambda)<= t,}
\expandafter\gdef\csname CGTip@eq:generator-phase\endcsname{Eq. (14.71), p. 75: lim\_sigma->(1)/(2)\textasciicircum{}+ImV\_lambda(sigma+it)=-\_k>=1eta\_k delta(t-gamma\_k(lambda)).}
\expandafter\gdef\csname CGTip@prop:one-ordinate\endcsname{Proposition 14.39 (The cell masses are sums of a one-ordinate function) (p. 76): For n>=2 there is a real function o\_n(kappa,kappa') of two real variables, depending only on n and the cell varphi\_n, such that for all lists kappa,kappa' in O c\_n(kappa)- c\_n(kappa')=sum\_k>=1 o\_n(kappa\_k,kappa'\_k), o\_n(kappa,kappa')=-2 Re(W\_n(kappa)-W\_n(kappa')) if n is not a power of 2, and, wi...}
\expandafter\gdef\csname CGTip@eq:one-ordinate\endcsname{Eq. (14.72), p. 76: c\_n(kappa)- c\_n(kappa')=sum\_k>=1 o\_n(kappa\_k,kappa'\_k), o\_n(kappa,kappa')=-2 Re(W\_n(kappa)-W\_n(kappa')) if n is not a power of 2,}
\expandafter\gdef\csname CGTip@eq:one-ordinate-log\endcsname{Eq. (14.73), p. 76: log zeta\_Z,kappa(s)-log zeta\_Z,kappa'(s)=sum\_n>=2(sum\_k>=1 o\_n(kappa\_k,kappa'\_k)) L\_n(s),}
\expandafter\gdef\csname CGTip@eq:one-ordinate-logzeta\endcsname{Eq. (14.74), p. 76: logzeta\_kappa(s)-logzeta\_kappa'(s)=sum\_k>=1log(1+q(s)/((1)/(4)+kappa\_k\textasciicircum{}2))/(1+q(s)/((1)/(4)+kappa'\_k\textasciicircum{}2)) .}
\expandafter\gdef\csname CGTip@eq:one-ordinate-joining\endcsname{Eq. (14.75), p. 76: j(kappa,kappa';s):=sum\_n>=2 o\_n(kappa,kappa') L\_n(s)-log(1+q(s)/((1)/(4)+kappa\textasciicircum{}2))/(1+q(s)/((1)/(4)+kappa'\textasciicircum{}2)), h\_kappa- h\_kappa'=sum\_k>=1j(kappa\_k,kappa'\_k;s) formally.}
\expandafter\gdef\csname CGTip@heur:phase-rotation\endcsname{Heuristic 14.40 (What happens to the phase defect along the path) (p. 76): The following is a computation by stationary phase without remainder terms; it is not proved. For n much larger than kappa the cell resolves the oscillation, Re W\_n(kappa) n\textasciicircum{}-1/2cos( n)/log n, and for n much smaller than kappa/2pi the kernel W\_n(kappa) is negligible because varphi is smooth (Prop...}
\expandafter\gdef\csname CGTip@eq:phase-rotation\endcsname{Eq. (14.76), p. 77: Imj(kappa,* ;(1)/(2)+it)amplitudexcos(2\_RVM(kappa)+phase(t)), |j(kappa,* ;sigma+it)|\textasciicircum{}-1/2-sigma,}
\expandafter\gdef\csname CGTip@subsec:transport\endcsname{Section 14.13: Transport from the clock on the rungs (p. 77)}
\expandafter\gdef\csname CGTip@eq:transport-defs\endcsname{Eq. (14.77), p. 77: c\_ f(t):=int\_0\textasciicircum{}inf f(u)cos(tu) du, k(u):=(1)/(2)(e\textasciicircum{}u/2-(e\textasciicircum{}-u/2)/(e\textasciicircum{}2u-1)), k\_ f:=int\_0\textasciicircum{}inf f(u) k(u) du .}
\expandafter\gdef\csname CGTip@prop:rung-functional\endcsname{Proposition 14.41 (The rung functional) (p. 77): For every list kappa in O (Section CGrefsec:lists), R\_ f(kappa):=(1)/(2)int\_1\textasciicircum{}inf f(log x) x\textasciicircum{}-1/2log x dJ\textasciicircum{} M\_kappa(x)= k\_ f-sum\_k>=1 c\_ f(kappa\_k), the series converging absolutely. Along the family, with R\_ f(lambda):= R\_ f((gamma\_k(lambda))), R\_ f(lambda)= R\_ f(1)- S\_ f(lambda), S\_ f(lambda):=s...}
\expandafter\gdef\csname CGTip@eq:rung-functional\endcsname{Eq. (14.78), p. 77: R\_ f(kappa):=(1)/(2)int\_1\textasciicircum{}inf f(log x) x\textasciicircum{}-1/2log x dJ\textasciicircum{} M\_kappa(x)= k\_ f-sum\_k>=1 c\_ f(kappa\_k),}
\expandafter\gdef\csname CGTip@eq:transport\endcsname{Eq. (14.79), p. 77: R\_ f(lambda)= R\_ f(1)- S\_ f(lambda), S\_ f(lambda):=sum\_k>=1( c\_ f(gamma\_k(lambda))- c\_ f(tau\_k)), S\_ f(1)=0 .}
\expandafter\gdef\csname CGTip@eq:transport-D\endcsname{Eq. (14.80), p. 77: (D) => S\_ f(0)= R\_ f(1),}
\expandafter\gdef\csname CGTip@thm:transport-band\endcsname{Theorem 14.42 (The transport is band-limited in lambda) (p. 77): Let f in C\_c\textasciicircum{}inf(0,log N). The function lambda S\_ f(lambda) extends to an entire function with | S\_ f(lambda)|<= C\_ f(1+|lambda|log(2+|lambda|)) e\textasciicircum{}C\_ N |Imlambda| (lambda in C), and for every j>=1 d\_lambda\textasciicircum{}j S\_ f(1)= m\_j( f):=sum\_k>=1eta\_k\textasciicircum{} j c\_ f\textasciicircum{}(j)(tau\_k), S\_ f(0)=sum\_j>=1((-1)\textasciicircum{}j)/(j!) m\_j( f)...}
\expandafter\gdef\csname CGTip@eq:transport-band\endcsname{Eq. (14.81), p. 77: | S\_ f(lambda)|<= C\_ f(1+|lambda|log(2+|lambda|)) e\textasciicircum{}C\_ N |Imlambda| (lambda in C),}
\expandafter\gdef\csname CGTip@eq:transport-moments\endcsname{Eq. (14.82), p. 77: d\_lambda\textasciicircum{}j S\_ f(1)= m\_j( f):=sum\_k>=1eta\_k\textasciicircum{} j c\_ f\textasciicircum{}(j)(tau\_k), S\_ f(0)=sum\_j>=1((-1)\textasciicircum{}j)/(j!) m\_j( f).}
\expandafter\gdef\csname CGTip@prop:clock-defect-poisson\endcsname{Proposition 14.43 (The defect of the clock on a rung) (p. 78): Let t\_0 in (2pi,tau\_0) be the point with N\_RVM(t\_0)=0, and let f in C\_c\textasciicircum{}inf(0,inf). Then sum\_k>=1 c\_ f(tau\_k)= I\_0+2sum\_m>=1(-1)\textasciicircum{}m I\_m, I\_m:=int\_t\_0\textasciicircum{}inf c\_ f(t)cos(2pi mN\_RVM(t))N\_RVM'(t) dt=O\_ f(m\textasciicircum{}-2), I\_0=-int\_0\textasciicircum{}inf( f(u))/(2pi u)(sin(ut\_0)log(t\_0)/(2pi)+(pi)/(2)-int\_0\textasciicircum{}ut\_0(sin y)/(y) dy)du, an...}
\expandafter\gdef\csname CGTip@eq:clock-defect-poisson\endcsname{Eq. (14.83), p. 78: sum\_k>=1 c\_ f(tau\_k)= I\_0+2sum\_m>=1(-1)\textasciicircum{}m I\_m, I\_m:=int\_t\_0\textasciicircum{}inf c\_ f(t)cos(2pi mN\_RVM(t))N\_RVM'(t) dt=O\_ f(m\textasciicircum{}-2),}
\expandafter\gdef\csname CGTip@eq:clock-defect-I0\endcsname{Eq. (14.84), p. 78: I\_0=-int\_0\textasciicircum{}inf( f(u))/(2pi u)(sin(ut\_0)log(t\_0)/(2pi)+(pi)/(2)-int\_0\textasciicircum{}ut\_0(sin y)/(y) dy)du,}
\expandafter\gdef\csname CGTip@subsec:multiscale-jacobian\endcsname{Section 14.14: The multiscale Jacobian of the midpoint clock (p. 78)}
\expandafter\gdef\csname CGTip@eq:clock-jacobian\endcsname{Eq. (14.85), p. 78: (J\_Tv)(u):=-usum\_tau\_k<= Tv\_ksin(tau\_k u).}
\expandafter\gdef\csname CGTip@thm:clock-jacobian\endcsname{Theorem 14.44 (Multiscale Jacobian of the midpoint clock) (p. 78): The following statements hold. For every T>tau\_2, Delta\_T>= (2pi)/(log(T/2pi)). Fix a>0 and \_ win>0, and put [ L\_T:=(1+\_ win)log(T/2pi), I\_T=[a,a+L\_T]. ] There is c\_a,\_ win>0, independent of T, such that for all finitely supported v=(v\_k) on k:tau\_k<= T, |J\_Tv|\_L\textasciicircum{}2(I\_T)\textasciicircum{}2 >= c\_a,\_ winL\_Tsum\_tau\_k...}
\expandafter\gdef\csname CGTip@eq:clock-gap-lower\endcsname{Eq. (14.86), p. 79: Delta\_T>= (2pi)/(log(T/2pi)).}
\expandafter\gdef\csname CGTip@eq:multiscale-frame\endcsname{Eq. (14.87), p. 79: |J\_Tv|\_L\textasciicircum{}2(I\_T)\textasciicircum{}2 >= c\_a,\_ winL\_Tsum\_tau\_k<= T|v\_k|\textasciicircum{}2.}
\expandafter\gdef\csname CGTip@eq:fixed-window-pair\endcsname{Eq. (14.88), p. 79: |Jv\textasciicircum{}(k)|\_L\textasciicircum{}2(0,U) <= (U\textasciicircum{}5/2)/(sqrt(10))(tau\_k+1-tau\_k)->0.}
\expandafter\gdef\csname CGTip@eq:smooth-tail-clock\endcsname{Eq. (14.89), p. 79: sum\_tau\_k>T| c\_ f(kappa\_k)- c\_ f(tau\_k)| <<\_A, f,C\_eta T\textasciicircum{}1-Alog T.}
\expandafter\gdef\csname CGTip@eq:multiscale-residual\endcsname{Eq. (14.90), p. 80: varepsilon\_r(T,U):= \_h,b| D\_RVM+usum\_tau\_k<= Th\_ksin(tau\_k u) -sum\_log n<= Ub\_ndelta(u-log n) |\_H\textasciicircum{}-r(0,U).}
\expandafter\gdef\csname CGTip@sec:lists\endcsname{Section 15: The space of ordinate lists and the existence form of the final step (p. 80)}
\expandafter\gdef\csname CGTip@subsec:lists\endcsname{Section 15.1: Lists, their zeta functions and their atomizations (p. 80)}
\expandafter\gdef\csname CGTip@eq:list-space\endcsname{Eq. (15.1), p. 80: tau\_0\textasciicircum{}+<=kappa\_1<=kappa\_2<=*s, |kappa\_k-tau\_k|<= C\_eta, |int\_tau\_0\textasciicircum{}X(N\_kappa(tau)-N\_RVM(tau))dtau|<= Clog X (X>=tau\_0),}
\expandafter\gdef\csname CGTip@eq:list-objects\endcsname{Eq. (15.2), p. 80: zeta\_kappa(s):=(1)/(G(s))prod\_k>=1(1+(q(s))/((1)/(4)+kappa\_k\textasciicircum{}2)), c\_n(kappa), a\_N(kappa), k\_N(kappa), zeta\_Z,kappa=zeta\_kappa e\textasciicircum{} h\_kappa, V\_sigma(kappa),}
\expandafter\gdef\csname CGTip@thm:transfer-O\endcsname{Theorem 15.1 (Transfer to O) (p. 81): Let kappa in O. The following statements, proved above for the lists (gamma\_k(lambda))\_k>=1, hold for kappa with the same proofs, and their constants depend only on tau\_0\textasciicircum{}+, C\_eta and C in CGeqrefeq:list-space and on the parameters of the statement (sigma, eps, n, N), never on the list: C=C(tau\_0...}
\expandafter\gdef\csname CGTip@thm:rigidity\endcsname{Theorem 15.2 (Rigidity of the coincidence) (p. 81): The set C:=kappa in O: h\_kappa is constant of the lists whose continuous and atomized zeta functions coincide up to a constant factor has at most one element. If kappa in C, then zeta\_kappa= zeta\_Z,kappa=zeta, kappa is the list of the ordinates of zeta, and (D) holds. Conversely, (D) implies C=(g...}
\expandafter\gdef\csname CGTip@eq:coincidence-set\endcsname{Eq. (15.3), p. 81: C:=kappa in O: h\_kappa is constant}
\expandafter\gdef\csname CGTip@prop:lists-compact\endcsname{Proposition 15.3 (Compactness and continuity) (p. 82): The space O and the functions CGeqrefeq:list-objects have the following properties. O is compact in the product topology of prod\_k[tau\_k-C\_eta,tau\_k+C\_eta]. For every n>=2, N>=1 and sigma>1, the functions kappa c\_n(kappa), a\_N(kappa), k\_N(kappa) and V\_sigma(kappa) are continuous on O. For every r...}
\expandafter\gdef\csname CGTip@eq:lists-compact-dominant\endcsname{Eq. (15.4), p. 82: | F\_kappa(t)|<= H(t):=(C'\_sigma (1+log(2+|t|)))/(1+|t|), H in L\textasciicircum{}2(R),}
\expandafter\gdef\csname CGTip@cor:existence-form\endcsname{Corollary 15.4 (Existence form of the final step) (p. 82): For every sigma>1 the minimum min\_O V\_sigma is attained, and the following are equivalent: up(i) (D); up(ii) C!=; up(iii) there is kappa in O (Section CGrefsec:lists) with zeta\_Z,kappa=zeta; up(iv) min\_O V\_sigma=0.}
\expandafter\gdef\csname CGTip@rem:existence-input\endcsname{Remark 15.5 (What an existence proof uses) (p. 83): Compactness, continuity and convexity do not decide whether min\_O V\_sigma is 0: they do not involve the form of the functional equation. The existence statement uses the functional equation with the factor Gamma((s)/(2))pi\textasciicircum{}-s/2, for which a function with both an absolutely convergent Dirichlet se...}
\expandafter\gdef\csname CGTip@eq:first-variation\endcsname{Eq. (15.5), p. 83: (d c\_n(kappa))/(dkappa\_k)=2 Im V\_n(kappa\_k),}
\expandafter\gdef\csname CGTip@subsec:free-masses\endcsname{Section 15.2: The rung system with free masses (p. 83)}
\expandafter\gdef\csname CGTip@eq:rung-system\endcsname{Eq. (15.6), p. 83: sum\_k>=1cos(kappa\_ku)+sum\_n>=2 a\_n delta(u-log n)= k(u) in D'(0,inf).}
\expandafter\gdef\csname CGTip@thm:free-masses\endcsname{Theorem 15.6 (Existence form with free masses) (p. 83): (D) holds if and only if CGeqrefeq:rung-system has a solution (kappa,( a\_n)) with kappa in O (Section CGrefsec:lists). The solution is then unique: kappa=(gamma\_k), and a\_n=(log p)/(2 n) if n=p\textasciicircum{}m is a prime power, a\_n=0 otherwise.}
\expandafter\gdef\csname CGTip@eq:free-masses\endcsname{Eq. (15.7), p. 83: a\_n=(log p)/(2 n) if n=p\textasciicircum{}m is a prime power, a\_n=0 otherwise.}
\expandafter\gdef\csname CGTip@eq:free-masses-dJ\endcsname{Eq. (15.8), p. 83: (log x) dJ\textasciicircum{} M\_kappa(x)=sum\_n>=22 n a\_n delta\_n in D'(1,inf).}
\expandafter\gdef\csname CGTip@eq:free-masses-factor\endcsname{Eq. (15.9), p. 83: zeta\_kappa(s)=exp(nu\_6+nu\_7 2\textasciicircum{}-s) zeta\_Z,kappa(s) (Re s>1)}
\expandafter\gdef\csname CGTip@thm:nondegenerate\endcsname{Theorem 15.7 (Nondegeneracy of the solution) (p. 84): Let (kappa,( a\_n)) be a solution of CGeqrefeq:rung-system with kappa in O (Section CGrefsec:lists). Let (\_k) be a bounded real sequence and ( a\_n)\_n>=2 real numbers with a\_n=O\_eps(n\textasciicircum{}-1/2+eps) for every eps>0, such that sum\_k>=1\_k usin(kappa\_ku)=sum\_n>=2 a\_n delta(u-log n) in D'(0,inf), the deriva...}
\expandafter\gdef\csname CGTip@eq:rung-linearized\endcsname{Eq. (15.10), p. 84: sum\_k>=1\_k usin(kappa\_ku)=sum\_n>=2 a\_n delta(u-log n) in D'(0,inf),}
\expandafter\gdef\csname CGTip@eq:nondeg-H1\endcsname{Eq. (15.11), p. 84: H\_1(z):=sum\_k>=1\_kint\_0\textasciicircum{}inf(1- x\_0(u)) usin(kappa\_ku) e\textasciicircum{}-zu du=(1)/(2) J,(1- x\_0) e\textasciicircum{}-z|u| .}
\expandafter\gdef\csname CGTip@eq:nondeg-split\endcsname{Eq. (15.12), p. 84: H\_1(z)=sum\_k>=1\_k (2kappa\_kz)/((z\textasciicircum{}2+kappa\_k\textasciicircum{}2)\textasciicircum{}2)- y(z), y(z):=sum\_k>=1\_kint\_0\textasciicircum{}log2 x\_0(u) usin(kappa\_ku) e\textasciicircum{}-zu du .}
\expandafter\gdef\csname CGTip@eq:nondeg-H\endcsname{Eq. (15.13), p. 85: H(z)=nu\_8+sum\_k>=1\_k((kappa\_k)/(z\textasciicircum{}2+kappa\_k\textasciicircum{}2)-(1)/(kappa\_k)),}
\expandafter\gdef\csname CGTip@eq:nondeg-order\endcsname{Eq. (15.14), p. 85: max\_|s-(1)/(2)|<= r\_j|(s-1) Z(s)|<= C\_ Hlog\textasciicircum{}2r\_j*max\_|s-(1)/(2)|=r\_j|(s-1)zeta(s)|<=exp(2r\_j\textasciicircum{}1+eps)}
\expandafter\gdef\csname CGTip@rem:nondegenerate-scope\endcsname{Remark 15.8 (Scope of Theorem CGrefthm:nondegenerate) (p. 86): The theorem is conditional on the existence of a solution; it shows that a solution is isolated to first order, the Jacobian of CGeqrefeq:rung-system being injective there on the directions considered, except for the splitting of multiple ordinates. This is a necessary endpoint nondegeneracy cond...}
\expandafter\gdef\csname CGTip@subsec:gradient\endcsname{Section 15.3: The gradient of the vertical defect (p. 86)}
\expandafter\gdef\csname CGTip@eq:K-kappa\endcsname{Eq. (15.15), p. 86: K\_kappa:=( zeta\_Z,kappa)/(zeta)=sum\_n>=1k\_n(kappa) n\textasciicircum{}-s, V\_sigma(kappa)= K\_kappa, K\_kappa\_sigma-1, r\_m(kappa):=Re K\_kappa L\_m, K\_kappa\_sigma ,}
\expandafter\gdef\csname CGTip@prop:critical-points\endcsname{Proposition 15.9 (Critical points of the vertical defect) (p. 86): Let sigma>1 and kappa in O (Section CGrefsec:lists). Then | r\_m(kappa)|<=2m\textasciicircum{}-sigma(1+ V\_sigma(kappa)), | u\_m(t)|<= Cm\textasciicircum{}-1/2 off the powers of 2 and | u\_m(t)|<= C on them, and (d V\_sigma)/(dkappa\_k)(kappa)=2 Q\_kappa(kappa\_k), Q\_kappa(t):=sum\_m>=2 r\_m(kappa) u\_m(t), the series converging absolutely ...}
\expandafter\gdef\csname CGTip@eq:gradient\endcsname{Eq. (15.16), p. 86: (d V\_sigma)/(dkappa\_k)(kappa)=2 Q\_kappa(kappa\_k), Q\_kappa(t):=sum\_m>=2 r\_m(kappa) u\_m(t),}
\expandafter\gdef\csname CGTip@sec:models\endcsname{Section 16: Acceptance criteria and models in which the final step is a theorem (p. 87)}
\expandafter\gdef\csname CGTip@subsec:acceptance\endcsname{Section 16.1: The acceptance criterion: window and frequencies (p. 87)}
\expandafter\gdef\csname CGTip@eq:model-product\endcsname{Eq. (16.1), p. 87: C\_kappa( z):= zprod\_k>=1(1-( z\textasciicircum{}2)/(kappa\_k\textasciicircum{}2)),}
\expandafter\gdef\csname CGTip@eq:model-logder\endcsname{Eq. (16.2), p. 87: (C\_kappa')/(C\_kappa)( z)=(1)/( z)+sum\_k>=1(2 z)/( z\textasciicircum{}2-kappa\_k\textasciicircum{}2)=-iint\_0\textasciicircum{}inf e\textasciicircum{}i zu(1+2sum\_k>=1cos(kappa\_ku))du ,}
\expandafter\gdef\csname CGTip@thm:acceptance\endcsname{Theorem 16.1 (Acceptance criterion) (p. 87): For a list kappa as above the following are equivalent: the spectral density is atomic on the frequencies: there are real w\_l with sum\_l in F|w\_l|e\textasciicircum{}-nu\_10l<inf for some admissible nu\_10 and 1+2sum\_k>=1cos(kappa\_ku)=sum\_l in Fw\_l delta(u-l) in D'(0,inf); the continuous function has an Euler form: ...}
\expandafter\gdef\csname CGTip@eq:acceptance-density\endcsname{Eq. (16.3), p. 87: 1+2sum\_k>=1cos(kappa\_ku)=sum\_l in Fw\_l delta(u-l) in D'(0,inf);}
\expandafter\gdef\csname CGTip@eq:acceptance-euler\endcsname{Eq. (16.4), p. 87: e\textasciicircum{}ik z C\_kappa( z)=c\_0exp(sum\_l in Fb\_l e\textasciicircum{}il z), b\_l=-(w\_l)/(l);}
\expandafter\gdef\csname CGTip@eq:acceptance-window\endcsname{Eq. (16.5), p. 88: e\textasciicircum{}ik z C\_kappa( z)=sum\_l in F\textasciicircum{}+[0,2k]c\_l e\textasciicircum{}il z, c\_0!=0 .}
\expandafter\gdef\csname CGTip@subsec:graph-model\endcsname{Section 16.2: A non-periodic accepted list: one unitary matrix (p. 88)}
\expandafter\gdef\csname CGTip@eq:graph-matrices\endcsname{Eq. (16.6), p. 88: U:= v sqrt(1-v\textasciicircum{}2) sqrt(1-v\textasciicircum{}2) -v}
\expandafter\gdef\csname CGTip@eq:graph-det\endcsname{Eq. (16.7), p. 88: D( z):=(I- U X( z))=1-v e\textasciicircum{}2ip z+v e\textasciicircum{}2iq z-e\textasciicircum{}2i(p+q) z=-2i e\textasciicircum{}i(p+q) z S( z),}
\expandafter\gdef\csname CGTip@eq:graph-sine\endcsname{Eq. (16.8), p. 88: S( z):=sin((p+q) z)+vsin((p-q) z).}
\expandafter\gdef\csname CGTip@prop:graph-model\endcsname{Proposition 16.2 (The two-edge model) (p. 88): All zeros of S are real. Each interval between consecutive points (m+(1)/(2))pi/k, m in Z, contains a zero, and the number of zeros in [0,T], with multiplicity, is kT/pi+O(1); hence the positive zeros form a list kappa with |kappa\_k-kpi/k| bounded. The set of zeros is not periodic. S(- z)=-S( z),...}
\expandafter\gdef\csname CGTip@eq:graph-euler\endcsname{Eq. (16.9), p. 89: logD( z)=-sum\_r>=1(1)/(r)tr( U X( z))\textasciicircum{}r=sum\_l in Fb\_l e\textasciicircum{}il z, F=2(mp+m'q):m,m'>=00,}
\expandafter\gdef\csname CGTip@eq:graph-eigenphase\endcsname{Eq. (16.10), p. 89: theta\_j'(t)=2 Dv\_j,v\_j=2(p|v\_j,1|\textasciicircum{}2+q|v\_j,2|\textasciicircum{}2) >= 2min(p,q)>0 .}
\expandafter\gdef\csname CGTip@subsec:finite-euler-model\endcsname{Section 16.3: A finite Euler product: the constructions of this paper with the final step proved (p. 90)}
\expandafter\gdef\csname CGTip@eq:finite-euler\endcsname{Eq. (16.11), p. 90: F(s):=(1-2\textasciicircum{}-s)(1-3\textasciicircum{}-s)=1-2\textasciicircum{}-s-3\textasciicircum{}-s+6\textasciicircum{}-s=sum\_n|6mu(n) n\textasciicircum{}-s.}
\expandafter\gdef\csname CGTip@eq:finite-euler-zeros\endcsname{Eq. (16.12), p. 90: y\_k\_k>=1=(2pi m)/(log2)\_m>=1(2pi m)/(log3)\_m>=1 in increasing order,}
\expandafter\gdef\csname CGTip@eq:finite-euler-hadamard\endcsname{Eq. (16.13), p. 90: F(s)=6\textasciicircum{}-s/2 (log2log3) s\textasciicircum{}2prod\_k>=1(1+(s\textasciicircum{}2)/(y\_k\textasciicircum{}2)).}
\expandafter\gdef\csname CGTip@eq:finite-euler-list\endcsname{Eq. (16.14), p. 90: F\_kappa(s):=6\textasciicircum{}-s/2 (log2log3) s\textasciicircum{}2prod\_k>=1(1+(s\textasciicircum{}2)/(kappa\_k\textasciicircum{}2)), F\_kappa(-s)=6\textasciicircum{}s F\_kappa(s),}
\expandafter\gdef\csname CGTip@eq:finite-euler-clock\endcsname{Eq. (16.15), p. 90: F\_1(s)=(log2log3)/(2)*(s\textasciicircum{}2 (1+6\textasciicircum{}-s))/(1+s\textasciicircum{}2log\textasciicircum{}26/pi\textasciicircum{}2).}
\expandafter\gdef\csname CGTip@thm:finite-euler\endcsname{Theorem 16.3 (The final step in the model) (p. 90): up(Poisson, prime by prime.) In D'(R), 2+2sum\_k>=1cos(y\_ku)=sum\_p in 2,3 sum\_m in Zcos(2pi mu)/(log p)=sum\_p in 2,3log psum\_m in Zdelta(u-mlog p). The spectral density of the list (y\_k) is atomic on the points mlog p: the continuous function built from the list is the Dirichlet series with Euler ...}
\expandafter\gdef\csname CGTip@eq:finite-euler-poisson\endcsname{Eq. (16.16), p. 90: 2+2sum\_k>=1cos(y\_ku)=sum\_p in 2,3 sum\_m in Zcos(2pi mu)/(log p)=sum\_p in 2,3log psum\_m in Zdelta(u-mlog p).}
\expandafter\gdef\csname CGTip@eq:finite-euler-coefficients\endcsname{Eq. (16.17), p. 90: (d)/(dsigma) A\_lambda,T(sigma)=(-i)/(2T)[f(sigma+it)e\textasciicircum{}ilambda t]\_t=-T\textasciicircum{}t=T-lambda A\_lambda,T(sigma), so |(d)/(dsigma)(e\textasciicircum{}lambdasigma A\_lambda,T(sigma))|<=(M)/(T)e\textasciicircum{}lambdasigma,}
\expandafter\gdef\csname CGTip@eq:finite-euler-spectrum\endcsname{Eq. (16.18), p. 91: e\textasciicircum{}lambdasigma\_1 A\_lambda(sigma\_1)=e\textasciicircum{}lambdasigma\_2 A\_lambda(sigma\_2) (lambda in R).}
\expandafter\gdef\csname CGTip@eq:finite-euler-superposition\endcsname{Eq. (16.19), p. 91: y\_k=(2pi)/(log2) N (2pi)/(log3) N ,}
\expandafter\gdef\csname CGTip@subsec:three-abscissas\endcsname{Section 16.4: The diagonal operator over all primes and the three abscissas (p. 91)}
\expandafter\gdef\csname CGTip@eq:prime-det\endcsname{Eq. (16.20), p. 91: (I- X\_P(s))=prod\_p(1-p\textasciicircum{}-s)=(1)/(zeta(s)) (Re s>1).}
\expandafter\gdef\csname CGTip@eq:prime-det2\endcsname{Eq. (16.21), p. 91: D\_2(s):=\_2(I- X\_P(s))=prod\_p(1-p\textasciicircum{}-s)e\textasciicircum{}p\textasciicircum{}-s=exp(-sum\_m>=2(1)/(m) P(ms)) (Re s>(1)/(2)),}
\expandafter\gdef\csname CGTip@prop:no-finite-rank\endcsname{Proposition 16.4 (No common object of finite rank in the window of zeta) (p. 91): Let X>=2. There is no function f==0 that satisfies (H0) and the functional equation CGeqrefeq:Z1-FE and is an absolutely convergent Dirichlet series in Re s>1 supported on the integers all of whose prime factors are at most X.}
\expandafter\gdef\csname CGTip@eq:three-abscissas\endcsname{Eq. (16.22), p. 92: contraction: sigma>0; convergence of the determinant: sigma>1; axis of symmetry: sigma=(1)/(2) .}
\expandafter\gdef\csname CGTip@sec:routeB-finite\endcsname{Section 17: Route B for finite sections: compactness, prescribed masses, certificates, and a Kadec--Avdonin criterion with blocks of variable length (p. 92)}
\expandafter\gdef\csname CGTip@subsec:rb-intro\endcsname{Section 17.1: Introduction to the section (p. 92)}
\expandafter\gdef\csname CGTip@subsec:rb-setting\endcsname{Section 17.2: Setting (p. 93)}
\expandafter\gdef\csname CGTip@eq:rb-NRVM\endcsname{Eq. (17.1), p. 93: N\_RVM(t):=(t)/(2pi)log(t)/(2pi e)+(7)/(8), N\_RVM'(t)=(1)/(2pi)log(t)/(2pi) (t>0).}
\expandafter\gdef\csname CGTip@eq:rb-Ntau-floor\endcsname{Eq. (17.2), p. 93: n\_RVM(t)=floor(N\_RVM(t)+(1)/(2)), |n\_RVM(t)-N\_RVM(t)|<=(1)/(2) (t>=2pi).}
\expandafter\gdef\csname CGTip@def:rb-lists\endcsname{Definition 17.1 (Admissible lists) (p. 93): For C>0, a list =(kappa\_k)\_k>=1 is C-admissible if it is real and nondecreasing, kappa\_1>=2pi, and |kappa\_k-tau\_k|<= C for all k. Put N\_(t):=no.k:kappa\_k<= t (with multiplicity) and P\_(s):=prod\_k>=1(1+(q(s))/((1)/(4)+kappa\_k\textasciicircum{}2)), zeta\_:=(P\_)/(G). The clock list =(tau\_k) gives zeta\_=zeta\_RVM.}
\expandafter\gdef\csname CGTip@eq:rb-Pk\endcsname{Eq. (17.3), p. 93: P\_(s):=prod\_k>=1(1+(q(s))/((1)/(4)+kappa\_k\textasciicircum{}2)), zeta\_:=(P\_)/(G).}
\expandafter\gdef\csname CGTip@eq:rb-rung\endcsname{Eq. (17.4), p. 93: sum\_k>=1cos(kappa\_ku)+sum\_n>=2 a\_n delta(u-log n)= k(u) in D'(0,inf),}
\expandafter\gdef\csname CGTip@eq:rb-residual\endcsname{Eq. (17.5), p. 93: (, a),varphi:=sum\_k>=1 c\_varphi(kappa\_k)+sum\_n>=2 a\_nvarphi(log n)- k,varphi, c\_varphi(t):=int\_0\textasciicircum{}infvarphi(u)cos(tu) du .}
\expandafter\gdef\csname CGTip@subsec:rb-routeB\endcsname{Section 17.3: Route B: compactness, selection and the canonical box (p. 93)}
\expandafter\gdef\csname CGTip@eq:rb-rho\endcsname{Eq. (17.6), p. 93: s\_m(, a):=max\_varphi in T\_m|(, a),varphi| .}
\expandafter\gdef\csname CGTip@def:rb-B\endcsname{Definition 17.2 (Statement B) (p. 93): Statement B(C, M) says that there are finite sections (\textasciicircum{}(m), a\textasciicircum{}(m)) in K(C, M) with s\_m(\textasciicircum{}(m), a\textasciicircum{}(m))->0.}
\expandafter\gdef\csname CGTip@lem:rb-canonical-box\endcsname{Lemma 17.3 (Canonical box of the exact solution) (p. 93): Let M\_*:=max\_n>=2Lambda(n)/(2 n). Then M\_*=log7/(27)=0.367742...<1/e. If condition (D) holds, the pair (, a\textasciicircum{}zeta), with a\textasciicircum{}zeta\_n:=Lambda(n)/(2 n), solves CGeqrefeq:rb-rung and lies in K(C\_eta, M\_*), where C\_eta is the constant of CGreflem:bounded-displacement.}
\expandafter\gdef\csname CGTip@cor:rb-D-implies-B\endcsname{Corollary 17.4 ((D) implies B) (p. 94): If (D) holds, then B(C, M) holds for every C>= C\_eta and M>= M\_*.}
\expandafter\gdef\csname CGTip@prop:rb-compactness\endcsname{Proposition 17.5 (Compactness and passage to the limit) (p. 94): Let (\textasciicircum{}(m), a\textasciicircum{}(m)) in K(C, M) with s\_m(\textasciicircum{}(m), a\textasciicircum{}(m))->0. Then a subsequence converges coordinatewise to some (, a) in K(C, M), and every such limit satisfies R(, a)=0 in D'(0,inf).}
\expandafter\gdef\csname CGTip@hyp:rb-U\endcsname{Hypothesis 17.6 (Uniqueness in the box) (p. 94): U(C, M): the only (, a) in K(C, M) with R(, a)=0 is (, a\textasciicircum{}zeta).}
\expandafter\gdef\csname CGTip@rem:rb-U-status\endcsname{Remark 17.7 (What is proved and what is open in U) (p. 94): CGrefthm:free-masses proves uniqueness for lists in O. Besides the tube and the ordering, O (CGrefeq:list-space) carries a third condition, |int(N\_-N\_RVM)|<= Clog X, and a coordinatewise limit of sections does not inherit it. Hypothesis CGrefhyp:rb-U is therefore the statement that CGrefthm:free-...}
\expandafter\gdef\csname CGTip@cor:rb-selection\endcsname{Corollary 17.8 (Selection of the zeta list) (p. 94): Assume B(C, M) and U(C, M), and let (\textasciicircum{}(m), a\textasciicircum{}(m)) be any sequence in K(C, M) with s\_m->0. Then the whole sequence converges coordinatewise to (, a\textasciicircum{}zeta). In particular R(, a\textasciicircum{}zeta)=0, that is, (D) holds.}
\expandafter\gdef\csname CGTip@rem:rb-B-status\endcsname{Remark 17.9 (Logical status of route B) (p. 95): Corollaries CGrefcor:rb-D-implies-B and CGrefcor:rb-selection give (D) B(C\_eta, M\_*) and [B(C, M) and U(C, M)] (D). Route B therefore consists of two independent pieces. The first is existence, B, which is not weaker than (D). The second is uniqueness without the third condition of O, Hypothesis ...}
\expandafter\gdef\csname CGTip@subsec:rb-prescribed\endcsname{Section 17.3.1: Route B with prescribed von Mangoldt masses (p. 95)}
\expandafter\gdef\csname CGTip@eq:rb-laplace-def\endcsname{Eq. (17.7), p. 95: LT(s):=, j(u) e\textasciicircum{}-zu,}
\expandafter\gdef\csname CGTip@lem:rb-archimedean\endcsname{Lemma 17.10 (The Archimedean distribution) (p. 95): Let g(u):=(e\textasciicircum{}-u/2)/(e\textasciicircum{}2u-1)-(1)/(2u) for u>0; g is bounded, g(0\textasciicircum{}+)=-(3)/(4), and g(u)=-(1)/(2u)+O(e\textasciicircum{}-5u/2) as u->inf, in particular bounded. Put c\_K:=-(1)/(4)log(2pi e\textasciicircum{}gamma\_E), gamma\_E Euler's constant, and K:=(1)/(2)e\textasciicircum{}u/21\_+-(1)/(2) g 1\_+-(1)/(4) fp (1\_+)/(u)+c\_K delta\_0 . Then K is a distribut...}
\expandafter\gdef\csname CGTip@eq:rb-K-dist\endcsname{Eq. (17.8), p. 95: K:=(1)/(2)e\textasciicircum{}u/21\_+-(1)/(2) g 1\_+-(1)/(4) fp (1\_+)/(u)+c\_K delta\_0 .}
\expandafter\gdef\csname CGTip@eq:rb-LK\endcsname{Eq. (17.9), p. 95: LK(s)=(1)/(2) (G')/(G)(s) (Re s>1).}
\expandafter\gdef\csname CGTip@eq:rb-aux-12\endcsname{Eq. (17.10), p. 95: (1)/(2)(Gamma')/(Gamma)((s)/(2))=int\_0\textasciicircum{}inf((e\textasciicircum{}-2u)/(2u)-(e\textasciicircum{}-su)/(1-e\textasciicircum{}-2u))du =int\_0\textasciicircum{}inf((e\textasciicircum{}-2u)/(2u)-e\textasciicircum{}-zu(e\textasciicircum{}-u/2)/(e\textasciicircum{}2u-1))du-(1)/(s),}
\expandafter\gdef\csname CGTip@eq:rb-aux-13\endcsname{Eq. (17.11), p. 96: (1)/(2)(G')/(G)(s)=(1)/(2s)+(1)/(2(s-1))-(1)/(4)logpi+(1)/(4)(Gamma')/(Gamma)((s)/(2)) =L((1)/(2)e\textasciicircum{}u/21\_+)(s)-(1)/(2)L(g1\_+)(s)+(1)/(4)log(z)/(2)-(1)/(4)logpi .}
\expandafter\gdef\csname CGTip@lem:rb-laplace\endcsname{Lemma 17.11 (Laplace transform of the residual) (p. 96): Let be C-admissible and put H\_(u):=sum\_kkappa\_k\textasciicircum{}-2cos(kappa\_ku), which converges absolutely and uniformly on R since H\_(0)=sum\_kkappa\_k\textasciicircum{}-2<inf, and S\_:=-d\_u\textasciicircum{}2(H\_1\_+)+H\_(0) delta\_0', A\_zeta:=sum\_n>=2 a\textasciicircum{}zeta\_n delta\_log n, E\_:=S\_+A\_zeta-K . Then E\_ is a distribution of local order at most 2 (its re...}
\expandafter\gdef\csname CGTip@eq:rb-S-A-dist\endcsname{Eq. (17.12), p. 96: S\_:=-d\_u\textasciicircum{}2(H\_1\_+)+H\_(0) delta\_0', A\_zeta:=sum\_n>=2 a\textasciicircum{}zeta\_n delta\_log n, E\_:=S\_+A\_zeta-K .}
\expandafter\gdef\csname CGTip@eq:rb-LR\endcsname{Eq. (17.13), p. 96: LE\_(s)=hat R(s)=(1)/(2)((P\_')/(P\_)-(xi')/(xi))(s), xi:=(1)/(2)Gzeta ,}
\expandafter\gdef\csname CGTip@eq:rb-aux-14\endcsname{Eq. (17.14), p. 96: LS\_(s)=-z\textasciicircum{}2sum\_k(z)/(kappa\_k\textasciicircum{}2(z\textasciicircum{}2+kappa\_k\textasciicircum{}2))+zsum\_k(1)/(kappa\_k\textasciicircum{}2) =sum\_k(z)/(z\textasciicircum{}2+kappa\_k\textasciicircum{}2)=(1)/(2)(P\_')/(P\_)(s),}
\expandafter\gdef\csname CGTip@rem:rb-boundary-term\endcsname{Remark 17.12 (The boundary term) (p. 96): Without the correction H\_(0)delta\_0' in CGeqrefeq:rb-S-A-dist, the transform of -d\_u\textasciicircum{}2(H\_1\_+) is (1)/(2)P'\_/P\_-(s-(1)/(2))H\_(0): the second derivative of the cut-off function produces H\_(0)delta\_0'=H\_(0)delta\_0'. Either normalisation gives the same restriction to (0,inf) and differs by a polynomi...}
\expandafter\gdef\csname CGTip@prop:rb-prescribed-uniqueness\endcsname{Proposition 17.13 (Uniqueness with prescribed masses) (p. 97): Let be C-admissible for some C, and suppose R(, a\textasciicircum{}zeta)=0 in D'(0,inf). Then all nontrivial zeros of zeta lie on Re s=(1)/(2), and =.}
\expandafter\gdef\csname CGTip@def:rb-Bprime\endcsname{Definition 17.14 (Statement B') (p. 97): Statement B'(C) says that there are C-admissible lists \textasciicircum{}(m), each equal to the clock beyond some finite index, with s\_m(\textasciicircum{}(m), a\textasciicircum{}zeta)->0.}
\expandafter\gdef\csname CGTip@cor:rb-Bprime\endcsname{Corollary 17.15 (Route B') (p. 97): If B'(C) holds for some C, then every sequence as in Definition CGrefdef:rb-Bprime converges coordinatewise to , and all nontrivial zeros of zeta lie on Re s=(1)/(2). If all nontrivial zeros of zeta lie on Re s=(1)/(2), then B'(C) holds for every C>= C\_eta. Hence the Riemann hypothesis is equival...}
\expandafter\gdef\csname CGTip@rem:rb-prescribed\endcsname{Remark 17.16 (What prescribing the masses changes, and what it does not) (p. 97): Hypothesis CGrefhyp:rb-U is no longer needed: uniqueness holds for every list at bounded distance from the clock, not only in O, and it follows from Hadamard's theorem and the functional equation instead of CGrefthm:free-masses. Route B' has a single open piece, the existence statement B', which ...}
\expandafter\gdef\csname CGTip@subsec:rb-staged\endcsname{Section 17.3.2: Staged construction: a certificate theorem and a graded extension (p. 98)}
\expandafter\gdef\csname CGTip@def:rb-staged\endcsname{Definition 17.17 (Staged sections) (p. 98): Fix u\textasciicircum{}*\_1<u\textasciicircum{}*\_2<*s->inf, finite test sets T\_1\_2s supported in (0,u\textasciicircum{}*\_m) with dense union as in Section CGrefsubsec:rb-routeB, t\_m:=|T\_m|, and k\textasciicircum{}*\_m->inf. Stage m starts from a list \textasciicircum{}(m,0) that equals the clock beyond some finite index, moves the coordinates k<= k\textasciicircum{}*\_m by x in R\textasciicircum{}k\textasciicircum{}*\_m, and has mass...}
\expandafter\gdef\csname CGTip@eq:rb-Fm\endcsname{Eq. (17.15), p. 98: F\_m: R\textasciicircum{}k\textasciicircum{}*\_m-> R\textasciicircum{}t\_m, F\_m(x)\_varphi:=(\textasciicircum{}(m,0)+x, a\textasciicircum{}zeta),varphi (varphi in T\_m),}
\expandafter\gdef\csname CGTip@eq:rb-jacobian\endcsname{Eq. (17.16), p. 98: ( J\_m)\_varphi,k=d\_x\_k F\_m(0)\_varphi=-int\_0\textasciicircum{}infvarphi(u) usin(kappa\textasciicircum{}(m,0)\_ku) du (varphi in T\_m, k<= k\textasciicircum{}*\_m).}
\expandafter\gdef\csname CGTip@eq:rb-eta-H\endcsname{Eq. (17.17), p. 98: h\_m:=| F\_m(0)|\_inf, H\_m:=k\textasciicircum{}*\_mmax\_varphi in T\_mint\_0\textasciicircum{}inf|varphi(u)| u\textasciicircum{}2 du .}
\expandafter\gdef\csname CGTip@lem:rb-matching\endcsname{Lemma 17.18 (Monotone matching) (p. 98): Let (a\_k)\_k>=1 be real with |a\_k-tau\_k|<= C for all k, and let (a\_(k)) be its nondecreasing rearrangement. Then |a\_(k)-tau\_k|<= C for all k. If moreover a\_k=tau\_k for k>k\textasciicircum{}*, then a\_(k)=tau\_k for all k>k\textasciicircum{}*+j, where j:=no.k>k\textasciicircum{}*:tau\_k<=tau\_k\textasciicircum{}*+C<inf.}
\expandafter\gdef\csname CGTip@lem:rb-lipschitz\endcsname{Lemma 17.19 (Lipschitz bound for the Jacobian) (p. 98): For all x,x' in R\textasciicircum{}k\textasciicircum{}*\_m, |D F\_m(x)-D F\_m(x')|\_inf->inf<= H\_m|x-x'|\_inf, and | F\_m(x)- F\_m(0)- J\_mx|\_inf<=(1)/(2) H\_m|x|\_inf\textasciicircum{}2 .}
\expandafter\gdef\csname CGTip@eq:rb-taylor\endcsname{Eq. (17.18), p. 98: | F\_m(x)- F\_m(0)- J\_mx|\_inf<=(1)/(2) H\_m|x|\_inf\textasciicircum{}2 .}
\expandafter\gdef\csname CGTip@thm:rb-certificate\endcsname{Theorem 17.20 (Certificate for one stage) (p. 98): Suppose J\_m has a right inverse Q\_m: R\textasciicircum{}t\_m-> R\textasciicircum{}k\textasciicircum{}*\_m, J\_m Q\_m=I\_t\_m, with | Q\_m|\_inf->inf<= L\_m, and l\_m:=2 L\_m\textasciicircum{}2 H\_m h\_m<1 . Then the map T\_m(y):=y- F\_m( Q\_my) is a contraction of the ball |y|\_inf<=2 h\_m R\textasciicircum{}t\_m with constant l\_m, and x\_*:= Q\_my\_*, y\_* its unique fixed point in that ball, satisfie...}
\expandafter\gdef\csname CGTip@eq:rb-Am\endcsname{Eq. (A_m), p. 98: l\_m:=2 L\_m\textasciicircum{}2 H\_m h\_m<1 .}
\expandafter\gdef\csname CGTip@eq:rb-cert-conclusion\endcsname{Eq. (17.19), p. 98: F\_m(x\_*)=0, |x\_*|\_inf<=2 L\_m h\_m .}
\expandafter\gdef\csname CGTip@rem:rb-certificate-errors\endcsname{Remark 17.21 (Evaluation errors and order of the ordinates) (p. 99): Theorem CGrefthm:rb-certificate is stated for exact values of h\_m, L\_m and H\_m. If the residual is only known up to an error, | F\_m(0)- F\_m(0)|\_inf\_m, the theorem applies with h\_m replaced by | F\_m(0)|\_inf+epsilon\_m; if the computed matrix Q\_m satisfies only |I- J\_m Q\_m|\_inf->inf<=delta\_m<1, then...}
\expandafter\gdef\csname CGTip@thm:rb-staged\endcsname{Theorem 17.22 (Staged certificate implies route B') (p. 99): Let C<tau\_1-2pi. Start at \textasciicircum{}(1,0):= and, for m>=1, let \textasciicircum{}(m) be the nondecreasing rearrangement of \textasciicircum{}(m,0)+x\_*\textasciicircum{}(m) given by Theorem CGrefthm:rb-certificate, and \textasciicircum{}(m+1,0):=\textasciicircum{}(m). Suppose up(A\_m) holds for every m and sum\_j>=12 L\_j h\_j=sup\_m>=1 sum\_j<= m2 L\_j h\_j<= C . Then every \textasciicircum{}(m) is C-admissible a...}
\expandafter\gdef\csname CGTip@eq:rb-B\endcsname{Eq. (B), p. 99: sum\_j>=12 L\_j h\_j=sup\_m>=1 sum\_j<= m2 L\_j h\_j<= C .}
\expandafter\gdef\csname CGTip@rem:rb-certificate-status\endcsname{Remark 17.23 (What the certificate theorem does and does not do) (p. 99): Theorem CGrefthm:rb-staged replaces the extension lemma of version 3 of the companion note (there Hypothesis E) by two computable quantities per stage, l\_m=2 L\_m\textasciicircum{}2 H\_m h\_m and 2 L\_m h\_m. It requires no contraction factor q, no separate control of the perturbation of old tests, and no prior ``entr...}
\expandafter\gdef\csname CGTip@prop:rb-graded\endcsname{Proposition 17.24 (Graded extension) (p. 100): Let s\_m>=0 satisfy s\_m+1<= q\_m s\_m+ h\_m with 0<= q\_m<=1 and sum\_m h\_m<inf, and assume prod\_j=L\textasciicircum{}mq\_j->0 (m->inf) for every L>=1 . Then s\_m->0. Condition CGeqrefeq:rb-graded-tail holds if q\_m<=q<1 for all large m, and more generally if sum\_m(1-q\_m)=inf. If the lists \textasciicircum{}(m) producing s\_m= s\_m(\textasciicircum{}(m), a\textasciicircum{}...}
\expandafter\gdef\csname CGTip@eq:rb-graded-tail\endcsname{Eq. (17.20), p. 100: prod\_j=L\textasciicircum{}mq\_j->0 (m->inf) for every L>=1 .}
\expandafter\gdef\csname CGTip@rem:rb-graded-status\endcsname{Remark 17.25 (Tube budget, and the two geometries) (p. 100): In Proposition CGrefprop:rb-graded the tube |kappa\textasciicircum{}(m)\_k-tau\_k|<= C, uniform in m, is a separate hypothesis: summability of the coordinate increments sum\_m|kappa\textasciicircum{}(m+1)\_k-kappa\textasciicircum{}(m)\_k| gives convergence of each coordinate but not that the limit stays in a prescribed tube, so a construction must car...}
\expandafter\gdef\csname CGTip@heur:rb-frame\endcsname{Heuristic 17.26 (Frame bound from exponential sums of the clock) (p. 100): In the opposite regime k\textasciicircum{}*\_m t\_m one may take Q\_m:= J\_m\textasciicircum{} T( J\_m J\_m\textasciicircum{} T)\textasciicircum{}-1. If |(1)/(K\_m) J\_m J\_m\textasciicircum{} T-D|\_inf->inf<d\_- for a diagonal D>= d\_-I, and |( J\_m)\_varphi,k|<=a\_J, then L\_m<= t\_mA/(k\textasciicircum{}*\_m(d\_--epsilon))=O(t\_m/k\textasciicircum{}*\_m) and L\_m\textasciicircum{}2 H\_m=O(t\_m\textasciicircum{}2B/k\textasciicircum{}*\_m) with B\_0:=max\_|varphi|u\textasciicircum{}2, so (A\_m) becomes k\textasciicircum{}*...}
\expandafter\gdef\csname CGTip@subsec:rb-gonek\endcsname{Section 17.3.3: Explicit lists in the tube: Gonek's symmetrised Euler products (p. 101)}
\expandafter\gdef\csname CGTip@eq:rb-alphaX\endcsname{Eq. (17.21), p. 101: p\_X(t):= P\_X((1)/(2)+it)=-sum\_n<= X\textasciicircum{}2(Lambda\_X(n)sin(tlog n))/( nlog n),}
\expandafter\gdef\csname CGTip@eq:rb-aux-16\endcsname{Eq. (17.22), p. 101: m\_X:=sum\_n<= X\textasciicircum{}2(Lambda\_X(n))/( nlog n), n\_X:=sum\_n<= X\textasciicircum{}2(Lambda\_X(n))/( n),}
\expandafter\gdef\csname CGTip@eq:rb-gonek-phase\endcsname{Eq. (17.23), p. 101: \_X(t):=vartheta(t)+p\_X(t) in (pi)/(2)+pi Z .}
\expandafter\gdef\csname CGTip@lem:rb-gonek-tube\endcsname{Lemma 17.27 (Gonek lists lie in a tube around the clock) (p. 101): Fix X>=2 and put t\_X:=max(2pi e\textasciicircum{}2n\_X+2, t\_1), where t\_1 is such that |vartheta'(t)-(1)/(2)log(t)/(2pi)|<=(1)/(2) for t>= t\_1. Then \_X is strictly increasing on [t\_X,inf) with \_X'>=(1)/(2), the solutions of CGeqrefeq:rb-gonek-phase in [t\_X,inf) form a strictly increasing sequence t\_1\textasciicircum{}(X)<t\_2\textasciicircum{}(X)<*...}
\expandafter\gdef\csname CGTip@eq:rb-gonek-list\endcsname{Eq. (17.24), p. 101: kappa\textasciicircum{}(X)\_k:= tau\_k, k<= j\_X, t\textasciicircum{}(X)\_k-j\_X, k>j\_X,}
\expandafter\gdef\csname CGTip@eq:rb-gonek-tube\endcsname{Eq. (17.25), p. 101: |kappa\textasciicircum{}(X)\_k-tau\_k|<=(2m\_X+o(1))/(log(tau\_k/2pi)) (k->inf).}
\expandafter\gdef\csname CGTip@eq:rb-aux-15\endcsname{Eq. (17.26), p. 101: N\_RVM(kappa\textasciicircum{}(X)\_k)=k-(1)/(2)-(p\_X(kappa\textasciicircum{}(X)\_k))/(pi)+r\_vartheta(kappa\textasciicircum{}(X)\_k) (k>j\_X),}
\expandafter\gdef\csname CGTip@eq:rb-aux-17\endcsname{Eq. (17.27), p. 102: |N\_RVM(kappa\textasciicircum{}(X)\_k)-N\_RVM(tau\_k)|<=(m\_X)/(pi)+(c\_r)/(kappa\textasciicircum{}(X)\_k)}
\expandafter\gdef\csname CGTip@rem:rb-gonek-status\endcsname{Remark 17.28 (What the Gonek lists give and what remains) (p. 102): The lists \textasciicircum{}(X) use the primes up to X\textasciicircum{}2 and the clock, and nothing about the zeros of zeta. Their prime side is built in: formally, sum\_kdelta(t-kappa\textasciicircum{}(X)\_k)=(\_X'(t))/(pi)[1+2sum\_j>=1(-1)\textasciicircum{}jcos(2j\_X(t))] where \_X'>0, and the j=0 term carries p\_X'/pi=-(1)/()\_n<= X\textasciicircum{}2Lambda\_X(n)n\textasciicircum{}-1/2cos(tlog n), who...}
\expandafter\gdef\csname CGTip@subsec:rb-apriori\endcsname{Section 17.4: An a priori bound for quasi-solutions (p. 102)}
\expandafter\gdef\csname CGTip@def:rb-laplace-residual\endcsname{Definition 17.29 (Laplace residual) (p. 102): For Re s>1 put hat R(s):=(1)/(2) (P\_')/(P\_)(s)-(1)/(2) (G')/(G)(s)+sum\_n=2\textasciicircum{}n\textasciicircum{}* a\_n n\textasciicircum{}(1)/(2)-s. For T>=0 put E(T):=int\_-T\textasciicircum{}T|hat R(2+iv)| dv, upsilon:=sup\_T>=2( E(T))/(log T) in [0,inf]. Since int\_0\textasciicircum{}inf e\textasciicircum{}-(s-(1)/(2))ucos(kappa u) du=(s-(1)/(2))/((s-(1)/(2))\textasciicircum{}2+kappa\textasciicircum{}2)=(1)/(2)(d)/(ds)log(1+(q(s))/...}
\expandafter\gdef\csname CGTip@eq:rb-Rhat\endcsname{Eq. (17.28), p. 102: hat R(s):=(1)/(2) (P\_')/(P\_)(s)-(1)/(2) (G')/(G)(s)+sum\_n=2\textasciicircum{}n\textasciicircum{}* a\_n n\textasciicircum{}(1)/(2)-s.}
\expandafter\gdef\csname CGTip@eq:rb-E\endcsname{Eq. (17.29), p. 102: E(T):=int\_-T\textasciicircum{}T|hat R(2+iv)| dv, upsilon:=sup\_T>=2( E(T))/(log T) in [0,inf].}
\expandafter\gdef\csname CGTip@lem:rb-upsilon-infinite\endcsname{Lemma 17.30 (When the Laplace residual grows logarithmically) (p. 103): For a finite section, upsilon<inf if and only if a\_n=0 for every n. More precisely, with D(v):=sum\_n=2\textasciicircum{}n\textasciicircum{}* a\_nn\textasciicircum{}-3/2-iv, E(T)=int\_-T\textasciicircum{}T|D(v)| dv+O\_C,k\textasciicircum{}*(log T), liminf\_T->inf(1)/(2T)int\_-T\textasciicircum{}T|D(v)| dv>=(sum\_n<= n\textasciicircum{}*| a\_n|\textasciicircum{}2n\textasciicircum{}-3)/(sum\_n<= n\textasciicircum{}*| a\_n|n\textasciicircum{}-3/2)>0 ( a!=0).}
\expandafter\gdef\csname CGTip@eq:rb-upsilon-split\endcsname{Eq. (17.30), p. 103: E(T)=int\_-T\textasciicircum{}T|D(v)| dv+O\_C,k\textasciicircum{}*(log T), liminf\_T->inf(1)/(2T)int\_-T\textasciicircum{}T|D(v)| dv>=(sum\_n<= n\textasciicircum{}*| a\_n|\textasciicircum{}2n\textasciicircum{}-3)/(sum\_n<= n\textasciicircum{}*| a\_n|n\textasciicircum{}-3/2)>0 ( a!=0).}
\expandafter\gdef\csname CGTip@rem:rb-count-scope\endcsname{Remark 17.31 (Scope of the counting bound) (p. 103): Lemma CGreflem:rb-upsilon-infinite restricts Theorem CGrefthm:rb-count and Corollary CGrefcor:rb-displacement, as stated for finite sections, to the sections without masses, for which |N\_-n\_RVM|<= k\textasciicircum{}*+1 is immediate. The reason is that on a vertical line the Laplace transform of the prime side is...}
\expandafter\gdef\csname CGTip@eq:rb-beta\endcsname{Eq. (17.31), p. 103: B(C):=sum\_k>=1 sup\_x in [max(tau\_k-C,2pi), tau\_k+C] sup\_s in J|f(x;s)-f(tau\_k;s)| .}
\expandafter\gdef\csname CGTip@lem:rb-PL\endcsname{Lemma 17.32 (Phragm'en--Lindel"of, two forms) (p. 103): up(a) Let S:=l\_1<=Re s<=l\_2, Im s>= t\_1 have width a:=l\_2-l\_1. Let u be upper semicontinuous on S and subharmonic in its interior, with u<= M\_0 on d S and u(s)<= A\_0e\textasciicircum{}cIm s for some c<pi/a. Then u<= M\_0 on S.; up(b) Let g be holomorphic in Re s>1, continuous on Re s>=2, with |g|<=1 on Re s=2 and ...}
\expandafter\gdef\csname CGTip@lem:rb-factor\endcsname{Lemma 17.33 (The Dirichlet factor) (p. 104): Put D\_ a(s):=sum\_n=2\textasciicircum{}n\textasciicircum{}*(2 a\_n)/(log n)n\textasciicircum{}(1)/(2)-s and Y:=zeta\_e\textasciicircum{}-D\_ a. Then Y is holomorphic and zero-free in Re s>1, and |D\_ a|<= Mc\_ on Re s>=2. For the branch of log Y that is real on (1,inf), (log Y)'=2 hat R, |log Y(2)|<= B\_0:=|logzeta\_RVM(2)|+B(C)+ Mc\_, and, for all real t, |log Y(2+it)|<=...}
\expandafter\gdef\csname CGTip@eq:rb-logY-deriv\endcsname{Eq. (17.32), p. 104: (log Y)'=2 hat R, |log Y(2)|<= B\_0:=|logzeta\_RVM(2)|+B(C)+ Mc\_,}
\expandafter\gdef\csname CGTip@eq:rb-line\endcsname{Eq. (17.33), p. 104: |log Y(2+it)|<= B\_0+ E(|t|)<= B\_0+(2+|t|).}
\expandafter\gdef\csname CGTip@lem:rb-halfplane\endcsname{Lemma 17.34 (Bound on the half-plane) (p. 104): With W:=log|zeta\_|=log|P\_|-log|G| and B\_1:= B\_0+(22), B\_2:= B\_1+ Mc\_, for Re s>=2, |log|Y(s)||<= B\_1+|s-1|, | W(s)|<= B\_2+|s-1| .}
\expandafter\gdef\csname CGTip@eq:rb-halfplane\endcsname{Eq. (17.34), p. 104: |log|Y(s)||<= B\_1+|s-1|, | W(s)|<= B\_2+|s-1| .}
\expandafter\gdef\csname CGTip@lem:rb-strip\endcsname{Lemma 17.35 (Bound on the half-strip) (p. 104): Let S:=-1<=Re s<=2, Im s>=1 and B\_3:=max( B\_2+(3)/(4)log2, c\_tau+B(C)). Then W(s)<= B\_3+((3)/(2)+upsilon)log|s+2| (s in S).}
\expandafter\gdef\csname CGTip@eq:rb-strip\endcsname{Eq. (17.35), p. 104: W(s)<= B\_3+((3)/(2)+upsilon)log|s+2| (s in S).}
\expandafter\gdef\csname CGTip@eq:rb-G-ratio\endcsname{Eq. (17.36), p. 104: |(G(2+it))/(G(-1+it))|=(|1+it|)/(2) pi\textasciicircum{}-3/2((t)/(2)(pi t)/(2))\textasciicircum{}1/2<=(1+t)\textasciicircum{}3/2 (t>=1).}
\expandafter\gdef\csname CGTip@lem:rb-backlund\endcsname{Lemma 17.36 (Backlund--Jensen on a horizontal segment) (p. 105): There are absolute constants T\_0>=3 and c\_5 such that, for T>= T\_0 with Tkappa\_k, the variation of P\_ along the segment from 2+iT to (1)/(2)+iT is at most pi(n\_T+1), where n\_T<=(((5)/(2)+upsilon)log(T+8)+ E(T)+ B\_0+ B\_3+ Mc\_+2c\_5)/(log(4/3)) .}
\expandafter\gdef\csname CGTip@eq:rb-nu\endcsname{Eq. (17.37), p. 105: n\_T<=(((5)/(2)+upsilon)log(T+8)+ E(T)+ B\_0+ B\_3+ Mc\_+2c\_5)/(log(4/3)) .}
\expandafter\gdef\csname CGTip@lem:rb-argument\endcsname{Lemma 17.37 (Argument principle) (p. 105): For T>0, Tkappa\_k, let G\_T be the path 2->2+iT->(1)/(2)+iT. Then pi N\_(T)=Delta\_ G\_T P\_ =pi+vartheta(T)+Delta\_ G\_T\_, where vartheta(T)=(pi\textasciicircum{}-iT/2Gamma((1)/(4)+(iT)/(2))) is the Riemann--Siegel theta function, defined by continuous variation from vartheta(0)=0.}
\expandafter\gdef\csname CGTip@eq:rb-arg-principle\endcsname{Eq. (17.38), p. 105: pi N\_(T)=Delta\_ G\_T P\_ =pi+vartheta(T)+Delta\_ G\_T\_,}
\expandafter\gdef\csname CGTip@thm:rb-count\endcsname{Theorem 17.38 (Counting bound for quasi-solutions) (p. 105): There are absolute constants T\_0 and c\_5 such that, for every finite section in K(C, M) and every T>= T\_0, |N\_(T)-n\_RVM(T)|<=(c\_1+c\_2upsilon)log(T+8)+c\_3 B+c\_4, with c\_1=(5)/(2log(4/3))=8.690..., c\_2=(2)/(log(4/3))+(1)/(pi)=7.270..., c\_3=(1)/(log(4/3))+(1)/(pi)=3.794..., c\_4:=(3)/(2)+(c\_6)/(pi)+(...}
\expandafter\gdef\csname CGTip@eq:rb-count\endcsname{Eq. (17.39), p. 105: |N\_(T)-n\_RVM(T)|<=(c\_1+c\_2upsilon)log(T+8)+c\_3 B+c\_4,}
\expandafter\gdef\csname CGTip@eq:rb-a-constants\endcsname{Eq. (17.40), p. 106: c\_1=(5)/(2log(4/3))=8.690..., c\_2=(2)/(log(4/3))+(1)/(pi)=7.270..., c\_3=(1)/(log(4/3))+(1)/(pi)=3.794...,}
\expandafter\gdef\csname CGTip@eq:rb-aux-1\endcsname{Eq. (17.41), p. 106: |N\_(T)-n\_RVM(T)|<= n\_T+1+( E(T)+ Mc\_+c\_6)/(pi)+(1)/(T)+(1)/(2) .}
\expandafter\gdef\csname CGTip@cor:rb-displacement\endcsname{Corollary 17.39 (Displacement bound) (p. 106): Assume T\_0>=2pi e. For every k with tau\_k-C>= T\_0, |kappa\_k-tau\_k|<=(2pi[(c\_1+c\_2upsilon)log(tau\_k+C+8)+c\_3 B+c\_4])/(log((tau\_k-C)/2pi)) . In particular limsup\_k->inf|kappa\_k-tau\_k|<=2pi(c\_1+c\_2upsilon) uniformly over sections with given C, M,upsilon. For every p>0 there is k\_0(C, M,upsilon,p) su...}
\expandafter\gdef\csname CGTip@eq:rb-disp\endcsname{Eq. (17.42), p. 106: |kappa\_k-tau\_k|<=(2pi[(c\_1+c\_2upsilon)log(tau\_k+C+8)+c\_3 B+c\_4])/(log((tau\_k-C)/2pi)) .}
\expandafter\gdef\csname CGTip@rem:rb-tube\endcsname{Remark 17.40 (What Corollary cor:rb-displacement does not give) (p. 106): The proof uses the imposed tube |kappa\_k-tau\_k|<= C twice: through B(C) in B, and to replace D(tau\_k+q) by D(tau\_k+C). The conclusion is that, inside any finite tube, the coordinates of sufficiently high index lie in the smaller tube of radius 2pi(c\_1+c\_2upsilon)+p. So the constraint is inactive ...}
\expandafter\gdef\csname CGTip@heur:rb-window\endcsname{Heuristic 17.41 (What a finite section controls) (p. 106): The quantity in Theorem CGrefthm:rb-count is E(T), the L\textasciicircum{}1 norm on Re s=2 of the Laplace transform of the residual. It is not the Euclidean residual that a numerical section minimises over finitely many tests. Let the tests be j\_0(u)varphi(u) with a fixed window j\_0 in C\_c\textasciicircum{}inf(0,inf), j\_0==1 on [...}
\expandafter\gdef\csname CGTip@heur:rb-informative\endcsname{Heuristic 17.42 (Informative part of a section) (p. 107): Where the atoms are denser than the resolution of the tests, e\textasciicircum{}-u<<pi/w\textasciicircum{}, the masses alone can approximate a smooth target by a Riemann sum a\_n g(log n)/n without moving any ordinate. That part of the window then carries no information on the ordinates. Counting unknowns against equations in the ...}
\expandafter\gdef\csname CGTip@subsec:rb-riesz\endcsname{Section 17.5: A Kadec--Avdonin criterion with blocks of variable length (p. 107)}
\expandafter\gdef\csname CGTip@subsec:rb-generating\endcsname{Section 17.5.1: Discrepancy and generating function (p. 107)}
\expandafter\gdef\csname CGTip@def:rb-discrepancy\endcsname{Definition 17.43 (Discrepancy, Poisson and conjugate operators) (p. 107): Let X=(x\_k)\_k in Z be a real nondecreasing sequence with |x\_k-(k-(1)/(2))|<= c\_Xlog(2+|k|) (k in Z). Its discrepancy with respect to the half-integers is \_X(x):=sum\_k in Z(1\_[x\_k,inf)(x)-1\_[k-(1)/(2),inf)(x)). The k-th term is +-1 on the interval between x\_k and k-(1)/(2) and 0 elsewhere. By CGeq...}
\expandafter\gdef\csname CGTip@eq:rb-H\endcsname{Eq. (17.43), p. 107: |x\_k-(k-(1)/(2))|<= c\_Xlog(2+|k|) (k in Z).}
\expandafter\gdef\csname CGTip@eq:rb-psi\endcsname{Eq. (17.44), p. 107: \_X(x):=sum\_k in Z(1\_[x\_k,inf)(x)-1\_[k-(1)/(2),inf)(x)).}
\expandafter\gdef\csname CGTip@eq:rb-PQ\endcsname{Eq. (17.45), p. 107: P\_yg(x):=(1)/()(y g(s) ds)/((x-s)\textasciicircum{}2+y\textasciicircum{}2), Y\_yg(x):=(1)/() g(s)[(x-s)/((x-s)\textasciicircum{}2+y\textasciicircum{}2)+(s)/(1+s\textasciicircum{}2)]ds .}
\expandafter\gdef\csname CGTip@lem:rb-generating\endcsname{Lemma 17.44 (Generating function) (p. 107): Let X=(x\_k)\_k in Z satisfy CGeqrefeq:rb-H, understood as a sequence indexed by Z in which a point of multiplicity m occurs m times. Let F\_0(z):=cospi z, and fix z\_0 in R(X( Z+(1)/(2))). The series G\_X(z):=sum\_k in Z((1)/(z-x\_k)-(1)/(z-k+(1)/(2))) converges absolutely and locally uniformly off X( ...}
\expandafter\gdef\csname CGTip@eq:rb-aux-2\endcsname{Eq. (17.46), p. 107: G\_X(z):=sum\_k in Z((1)/(z-x\_k)-(1)/(z-k+(1)/(2)))}
\expandafter\gdef\csname CGTip@eq:rb-aux-3\endcsname{Eq. (17.47), p. 107: F\_X(z):=F\_0(z)expint\_z\_0\textasciicircum{}zG\_X(z) dz}
\expandafter\gdef\csname CGTip@eq:rb-F-formula\endcsname{Eq. (17.48), p. 107: log|F\_X(x+iy)|=log|F\_0(x+iy)|+pi Y\_y\_X(x)+c\_X'' ;}
\expandafter\gdef\csname CGTip@eq:rb-Q-growth\endcsname{Eq. (17.49), p. 108: | Y\_y\_X(x)|<= c\_X'''(1+log(1+y\textasciicircum{}-2))log\textasciicircum{}2(4+|z|) (z=x+iy, y>0).}
\expandafter\gdef\csname CGTip@subsec:rb-criterion\endcsname{Section 17.5.2: The criterion (p. 108)}
\expandafter\gdef\csname CGTip@eq:rb-aux-4\endcsname{Eq. (17.50), p. 108: [w]\_A\_2:=sup\_I((1)/(|I|)int\_Iw)((1)/(|I|)int\_Iw\textasciicircum{}-1)<inf ,}
\expandafter\gdef\csname CGTip@thm:rb-criterion\endcsname{Theorem 17.45 (Block (1)/(4)-criterion) (p. 108): Let X satisfy CGeqrefeq:rb-H and suppose \_X=A+B with measurable A,B such that: |A|\_inf<= d\_<(1)/(4); sup\_x| Y\_1B(x)-c\_B|<= U for some constant c\_B. Then w(x):=|F\_X(x+i)|\textasciicircum{}2 is in (A\_2) and [w]\_A\_2<=(pi\textasciicircum{}2\textasciicircum{}2(pi) e\textasciicircum{}4pi U)/(cos\textasciicircum{}2(2\_)) .}
\expandafter\gdef\csname CGTip@eq:rb-A2const\endcsname{Eq. (17.51), p. 108: [w]\_A\_2<=(pi\textasciicircum{}2\textasciicircum{}2(pi) e\textasciicircum{}4pi U)/(cos\textasciicircum{}2(2\_)) .}
\expandafter\gdef\csname CGTip@eq:rb-aux-5\endcsname{Eq. (17.52), p. 108: |W(o)|=e\textasciicircum{}2pi Y\_1+y\_1A(x\_1), (o)=-2pi P\_1+y\_1A(x\_1) in [-2\_,2\_].}
\expandafter\gdef\csname CGTip@eq:rb-aux-6\endcsname{Eq. (17.53), p. 109: (1)/(|I|)int\_I|W|<=(1)/(cosb\_) (1)/(|I|)int\_IReW<=(pi)/(cosb\_)ReW(o\_I)<=(pi)/(cosb\_)|W(o\_I)| .}
\expandafter\gdef\csname CGTip@cor:rb-riesz\endcsname{Corollary 17.46 (Riesz bases) (p. 109): Under the hypotheses of Theorem CGrefthm:rb-criterion: if X is uniformly discrete (in particular without repeated points), e\textasciicircum{}ix\_kv\_k in Z is a Riesz basis of L\textasciicircum{}2(-pi,pi); if X is a finite union of uniformly discrete sets, the generalised divided differences of e\textasciicircum{}ix\_kv in the sense of [ref] form a...}
\expandafter\gdef\csname CGTip@subsec:rb-variable\endcsname{Section 17.5.3: Kadec, Avdonin, and variable block length (p. 109)}
\expandafter\gdef\csname CGTip@prop:rb-variable-blocks\endcsname{Proposition 17.47 (Bounded conjugate for block-supported oscillation) (p. 109): Let J\_m\_m in Z be a partition of R into consecutive intervals of lengths n\_m>=1. Let B=sum\_mB\_m with suppB\_m J\_m, int B\_m=0 and |B\_m|<= D'\_m, sup\_mD'\_m<inf. Put U\_ loc:=sup\_m sup\_x in J\_m-1 J\_m J\_m+1|int\_J\_m((x-s) B\_m(s))/(pi((x-s)\textasciicircum{}2+1)) ds|, O:=sup\_m\_0sum\_|m-m\_0|>=2(n\_m\textasciicircum{}2)/(dist(J\_m\_0,J\_m)\textasciicircum{}2). I...}
\expandafter\gdef\csname CGTip@eq:rb-L0\endcsname{Eq. (17.54), p. 109: U\_ loc:=sup\_m sup\_x in J\_m-1 J\_m J\_m+1|int\_J\_m((x-s) B\_m(s))/(pi((x-s)\textasciicircum{}2+1)) ds|, O:=sup\_m\_0sum\_|m-m\_0|>=2(n\_m\textasciicircum{}2)/(dist(J\_m\_0,J\_m)\textasciicircum{}2).}
\expandafter\gdef\csname CGTip@cor:rb-kadec-avdonin\endcsname{Corollary 17.48 (Kadec and Avdonin, with the constant (1)/(4)) (p. 109): up(a) [ref] If |d\_k|<= d\_<(1)/(4) for all k, then e\textasciicircum{}ix\_kv is a Riesz basis of L\textasciicircum{}2(-pi,pi).; up(b) [ref] If X is uniformly discrete, sup\_k|d\_k|<= D, and for a fixed n the means of d\_k over consecutive blocks of n indices are at most d\_<(1)/(4) in modulus, then e\textasciicircum{}ix\_kv is a Riesz basis of L\textasciicircum{}2(-pi,pi).}
\expandafter\gdef\csname CGTip@thm:rb-variable-N\endcsname{Theorem 17.49 (Avdonin's theorem with variable block length) (p. 110): Let X be uniformly discrete with sup\_k|d\_k|<= D. Let B\_m be a partition of Z into consecutive blocks of n\_m indices, and J\_m:=cup\_k in B\_mI\_k. Let conj d\_m be the mean of d\_k over B\_m, D\_m:=max\_k in B\_m|d\_k-conj d\_m|, and n\_m\textasciicircum{}*:=n\_m-1+n\_m+n\_m+1. Assume sup\_m|conj d\_m|<= d\_<(1)/(4); O<inf, with O ...}
\expandafter\gdef\csname CGTip@eq:rb-U-variable\endcsname{Eq. (17.55), p. 110: U=16(1+D)\textasciicircum{}3+(6 V)/(pi)+( V O)/(22) .}
\expandafter\gdef\csname CGTip@rem:rb-price\endcsname{Remark 17.50 (The price of variable length) (p. 110): For fixed n, conditions (b) and (c) hold automatically, and Theorem CGrefthm:rb-variable-N is Avdonin's theorem. If n\_m->inf, the averaging condition (a) is still only |conj d\_m|<= d\_<(1)/(4). However, (c) then requires the oscillation inside the block to decay like 1/log n\_m. In Theorem CGrefthm...}
\expandafter\gdef\csname CGTip@subsec:rb-application\endcsname{Section 17.6: Application to lists near the RvM clock (p. 110)}
\expandafter\gdef\csname CGTip@subsec:rb-unfolding\endcsname{Section 17.6.1: Unfolding (p. 110)}
\expandafter\gdef\csname CGTip@def:rb-unfolding\endcsname{Definition 17.51 (Unfolded list) (p. 110): Let t\_0:=9.6769... be the zero of N\_RVM in (2pi,inf). For a C-admissible list with kappa\_1>t\_0 (true for and ), put x\_k:=N\_RVM(kappa\_k)>0 for k>=1 and x\_1-k:=-x\_k. The resulting symmetric sequence X\_ has discrepancy \_:=\_X\_, and F\_:=F\_X\_ with z\_0=0. For the clock, X\_= Z+(1)/(2) and F\_=F\_0.}
\expandafter\gdef\csname CGTip@lem:rb-psi-bound\endcsname{Lemma 17.52 (Unfolded discrepancy) (p. 111): For a C-admissible list with kappa\_1>t\_0, |x\_k-(k-(1)/(2))|<=(C)/(2pi)log(tau\_k+C)/(2pi) for k>=1, so CGeqrefeq:rb-H holds. Moreover, for x=N\_RVM(t) with x>0, \_(x)=N\_(t)-n\_RVM(t), |N\_(t)-n\_RVM(t)|<=(C)/()(t+C)/(2pi)+2, and \_(-x)=-\_(x) for almost every x.}
\expandafter\gdef\csname CGTip@eq:rb-psi-t\endcsname{Eq. (17.56), p. 111: \_(x)=N\_(t)-n\_RVM(t), |N\_(t)-n\_RVM(t)|<=(C)/()(t+C)/(2pi)+2,}
\expandafter\gdef\csname CGTip@prop:rb-O-blocks\endcsname{Proposition 17.53 (The condition of O gives the averaging hypothesis with blocks log\textasciicircum{}2T) (p. 111): Let be C-admissible with kappa\_1>t\_0, and assume the third condition of O (CGrefeq:list-space) in the form |int\_2pi\textasciicircum{}R(N\_(t)-N\_RVM(t))dt|<= C\_2log R (R>=3), with a constant C\_2>0; this is no restriction, since CGeqrefeq:rb-O3 with one constant holds with every larger one, and it is needed because ...}
\expandafter\gdef\csname CGTip@eq:rb-O3\endcsname{Eq. (17.57), p. 111: |int\_2pi\textasciicircum{}R(N\_(t)-N\_RVM(t))dt|<= C\_2log R (R>=3),}
\expandafter\gdef\csname CGTip@eq:rb-block-mean\endcsname{Eq. (17.58), p. 111: |A\_|\_inf<=(3)/(16), |J\_m|=(8C\_2)/()\textasciicircum{}2T\_m (1+o(1)) (m->inf).}
\expandafter\gdef\csname CGTip@eq:rb-aux-7\endcsname{Eq. (17.59), p. 111: int\_J\_ dx =int\_T\textasciicircum{}T+H(N\_-n\_RVM)(t) N\_RVM'(t) dt; =N\_RVM'(T) int\_T\textasciicircum{}T+H (N\_-n\_RVM)+int\_T\textasciicircum{}T+H (N\_-n\_RVM)(N\_RVM'(t)-N\_RVM'(T)) dt .}
\expandafter\gdef\csname CGTip@eq:rb-aux-8\endcsname{Eq. (17.60), p. 111: |(1)/(|J|)int\_J\_|<=(2C\_2log(T+H))/(H)+(5)/(4N\_RVM'(T)H)+(H)/(2pi TN\_RVM'(T))((C)/()(T+H+C)/(2pi)+2).}
\expandafter\gdef\csname CGTip@rem:rb-not-equivalent\endcsname{Remark 17.54 (Correction of an earlier formulation) (p. 112): Proposition CGrefprop:rb-O-blocks is an implication. The averaging condition on blocks of length log\textasciicircum{}2T does not imply CGeqrefeq:rb-O3: summing |int\_J\_m|<=(3)/(16)|J\_m| over the blocks up to height R gives only O(Rlog R), not O(log R). The two conditions do not coincide.}
\expandafter\gdef\csname CGTip@cor:rb-one-condition\endcsname{Corollary 17.55 (Reduction to one condition) (p. 112): For every as in Proposition CGrefprop:rb-O-blocks: if Y\_1B\_ is bounded up to a constant, then |F\_(x+i)|\textasciicircum{}2 in (A\_2), with constant CGeqrefeq:rb-A2const for d\_=(3)/(16). Corollary CGrefcor:rb-riesz(a) then holds for X\_ if it is uniformly discrete, and Corollary CGrefcor:rb-riesz(b) if it is a finit...}
\expandafter\gdef\csname CGTip@subsec:rb-weight\endcsname{Section 17.6.2: The weight is a quotient of synthetic zetas (p. 112)}
\expandafter\gdef\csname CGTip@prop:rb-weight\endcsname{Proposition 17.56 (Weight identity) (p. 112): Let be C-admissible and satisfy CGeqrefeq:rb-O3. For 0<<=1 and t>=2pi, log|(zeta\_)/(zeta\_RVM)((1)/(2)++it)|=int\_2pi\textasciicircum{}inf((t-s) (N\_-n\_RVM)(s))/((t-s)\textasciicircum{}2+\textasciicircum{}2) ds+c\_+m\_,(t), where the integral converges as an improper integral and c\_ is a constant. The remainder satisfies |m\_,(t)|<= c(C,C\_2)log(2+t)/(1...}
\expandafter\gdef\csname CGTip@eq:rb-weight-identity\endcsname{Eq. (17.61), p. 112: log|(zeta\_)/(zeta\_RVM)((1)/(2)++it)|=int\_2pi\textasciicircum{}inf((t-s) (N\_-n\_RVM)(s))/((t-s)\textasciicircum{}2+\textasciicircum{}2) ds+c\_+m\_,(t),}
\expandafter\gdef\csname CGTip@eq:rb-aux-9\endcsname{Eq. (17.62), p. 112: log|(P\_)/(P\_)(s)|=sum\_k[l\_(t-kappa\_k)-l\_(t-tau\_k)]+sum\_k[l\_(t+kappa\_k)-l\_(t+tau\_k)]-sum\_klog((1)/(4)+kappa\_k\textasciicircum{}2)/((1)/(4)+tau\_k\textasciicircum{}2).}
\expandafter\gdef\csname CGTip@rem:rb-local-scale\endcsname{Remark 17.57 (Local scale) (p. 112): With =1 and \_T:=/N\_RVM'(T)=2pi/log(T/2pi) (the notation of the proof of Theorem CGrefthm:rb-obstruction), the substitution x=N\_RVM(t) maps the kernel ((t-s) ds)/((t-s)\textasciicircum{}2+\_T\textasciicircum{}2) to ((x-x') dx')/((x-x')\textasciicircum{}2+1) to first order near t=T. Comparing CGeqrefeq:rb-weight-identity with CGeqrefeq:rb-F-formula ...}
\expandafter\gdef\csname CGTip@subsec:rb-obstruction\endcsname{Section 17.6.3: Obstruction for the zeta list (p. 112)}
\expandafter\gdef\csname CGTip@thm:rb-obstruction\endcsname{Theorem 17.58 (No (A\_2) for the unfolded zeta ordinates) (p. 112): Let X\_ be the unfolded list of the ordinates gamma\_k>0 of the nontrivial zeros of zeta, counted with multiplicity, and let F be any entire function of exponential type whose zero set is X\_, with multiplicities. Then |F(x+i)|\textasciicircum{}2(A\_2) for every real !=0. The proof does not assume the Riemann hypothe...}
\expandafter\gdef\csname CGTip@eq:rb-aux-10\endcsname{Eq. (17.63), p. 113: (1)/(|I|)int\_I|f-f\_I|\textasciicircum{}2<=4[w\_0]\_A\_2 for every interval I .}
\expandafter\gdef\csname CGTip@eq:rb-BMO-L2\endcsname{Eq. (17.64), p. 113: (1)/(|I|)int\_I| P\_-( P\_)\_I|\textasciicircum{}2<=c\_ JN for every interval I .}
\expandafter\gdef\csname CGTip@eq:rb-aux-11\endcsname{Eq. (17.65), p. 114: int\_T\textasciicircum{}2T|S-c|\textasciicircum{}2>=int\_T\textasciicircum{}2TS\textasciicircum{}2-(1)/(T)(int\_T\textasciicircum{}2TS)\textasciicircum{}2=(1)/(2pi\textasciicircum{}2)Tloglog T (1+o(1))}
\expandafter\gdef\csname CGTip@lem:rb-smoothing\endcsname{Lemma 17.59 (Smoothing S at scale 1/log T) (p. 114): Let >0 be fixed, \_T:=2pi /log(T/2pi), and S:=0 on (-inf,t\_0). Then | P\_\_T* S-S|\_L\textasciicircum{}2[T,2T]=O\_( T), unconditionally.}
\expandafter\gdef\csname CGTip@eq:rb-fujii\endcsname{Eq. (17.66), p. 114: int\_T\textasciicircum{}2T(S(t+u)-S(t))\textasciicircum{}2dt<= c\_FTlog(2+ulog T),}
\expandafter\gdef\csname CGTip@rem:rb-fujii\endcsname{Remark 17.60 (A second proof) (p. 114): Lemma CGreflem:rb-smoothing also follows from Selberg's mean-square approximation of S(t) by S\_x\_0(t):=-(1)/()\_p<= xp\textasciicircum{}-1/2sin(tlog p), x\_0=T\textasciicircum{}delta [ref], and the Montgomery--Vaughan mean value theorem [ref]. Indeed P\_\_T* S\_x\_0-S\_x\_0=-(1)/()\_p<= xp\textasciicircum{}-1/2(e\textasciicircum{}-\_Tlog p-1)sin(tlog p) has mean square << ...}
\expandafter\gdef\csname CGTip@cor:rb-obstruction\endcsname{Corollary 17.61 (Consequences for the zeta list) (p. 114): up(a) e\textasciicircum{}ix\_kxi\_x\_k in X\_ is not a Riesz basis of L\textasciicircum{}2(-pi,pi).; up(b) If X\_ is a finite union of uniformly discrete sets, its generalised divided differences are not a Riesz basis of L\textasciicircum{}2(-pi,pi).; up(c) For every decomposition \_=A+B with |A|\_inf<(1)/(4), Y\_1B is unbounded. In particular this holds...}
\expandafter\gdef\csname CGTip@heur:rb-graduated\endcsname{Heuristic 17.62 (Graduated version) (p. 115): The proof of Theorem CGrefthm:rb-obstruction uses intervals I\_T of length Tlog T. On a window of L unfolded points at height T, the mean square oscillation of P\_ is governed by the variance of increments of S over t-intervals of length L/N\_RVM'(T). By Fujii's asymptotics this is expected to be of...}
\expandafter\gdef\csname CGTip@rem:rb-routeB-stability\endcsname{Remark 17.63 (Bearing on route B) (p. 115): Theorem CGrefthm:rb-obstruction rules out a uniform Riesz-basis (Fredholm) structure for the unfolded zeta list on the whole line, and hence a linear stability argument with constants independent of the height. It says nothing about statement B or B'. The compactness step of route B (Proposition ...}
\expandafter\gdef\csname CGTip@subsec:rb-numerics\endcsname{Section 17.7: Numerical certificates for finite sections (p. 115)}
\expandafter\gdef\csname CGTip@subsec:rb-num-setup\endcsname{Section 17.7.1: Set-up (p. 115)}
\expandafter\gdef\csname CGTip@subsec:rb-num-init\endcsname{Section 17.7.2: Initialisation: the clock is outside the basin, the Gonek lists are inside (p. 115)}
\expandafter\gdef\csname CGTip@fig:rb-gonek\endcsname{Figure 1 (p. 116): Distance to the zeta ordinates, in units of the local mean spacing 2pi/log(gamma\_k/2pi), of the midpoint clock and of the Gonek lists \textasciicircum{}(X) for X=3,8,20 (k<=160 for X=3,20; k<=400 for X=8). The dashed line is the observed radius of the Newton basin. The ordinates gamma\_k are used here only to meas...}
\expandafter\gdef\csname CGTip@subsec:rb-num-kantorovich\endcsname{Section 17.7.3: Certificates at the Gonek list: the direct bounds fail (p. 116)}
\expandafter\gdef\csname CGTip@subsec:rb-num-certified\endcsname{Section 17.7.4: Floating-point candidate certificates for exact square sections (p. 117)}
\expandafter\gdef\csname CGTip@eq:rb-num-banach\endcsname{Eq. (17.67), p. 117: l(r):=| Q J'|\_inf->inf r+(1)/(2) || Q| A\textasciicircum{}(3)(r)|\_inf->inf r\textasciicircum{}2<1, h<=(1- l(r)) r,}
\expandafter\gdef\csname CGTip@eq:rb-num-scaling\endcsname{Eq. (17.68), p. 117: L\_m T\textasciicircum{}1.8, | Q J'|\_inf->inf T\textasciicircum{}2.1, margin T\textasciicircum{}-2.1 (T<=487).}
\expandafter\gdef\csname CGTip@subsec:rb-num-entry\endcsname{Section 17.7.5: The entry step (p. 118)}
\expandafter\gdef\csname CGTip@fig:rb-scaling\endcsname{Figure 2 (p. 118): Constants of the Banach certificate CGeqrefeq:rb-num-banach for the square sections of Section CGrefsubsec:rb-num-certified (solid lines, with fitted slopes in log T), and the band-only entry test of Section CGrefsubsec:rb-num-entry (dashed): the reciprocal of the band margin and the reciprocal o...}
\expandafter\gdef\csname CGTip@subsec:rb-num-reading\endcsname{Section 17.7.6: Reading of the computations (p. 118)}
\expandafter\gdef\csname CGTip@subsec:rb-status\endcsname{Section 17.8: Status of the statements (p. 119)}
\expandafter\gdef\csname CGTip@sec:numerics\endcsname{Section 18: Numerical illustration with 2 001 0522001052 zeros (p. 119)}
\expandafter\gdef\csname CGTip@subsec:num-clock\endcsname{Section 18.1: The clock and the zeros (p. 120)}
\expandafter\gdef\csname CGTip@exp:clock\endcsname{Experiment 18.1 (Clock, midpoints and displacements) (p. 120): Object: N\_RVM, N\_zeta, tau\_k, eta\_k=tau\_k-gamma\_k. Range: k<= N. Expected: CGeqrefeq:midpoint-condition to rounding, and |eta\_k|<= C\_eta (Lemma CGreflem:bounded-displacement). Observed: max\_k|N\_RVM(tau\_k)-(k-(1)/(2))|=6.985*10\textasciicircum{}-10; max\_k|eta\_k|=0.9937 (at k=34), mean |eta\_k|=0.1512, root mean squ...}
\expandafter\gdef\csname CGTip@fig:clock-zoom\endcsname{Figure 3 (p. 120): The midpoint clock N\_RVM (red), the zero staircase N\_zeta (black), the midpoint ordinates tau\_k (red points) and the points (gamma\_k,k-(1)/(2)) (squares). Left: the first nine zeros. Right: the last twelve zeros of the file.}
\expandafter\gdef\csname CGTip@fig:remainder\endcsname{Figure 4 (p. 121): The remainder N\_zeta(T)-N\_RVM(T)=S(T)+O(1/T) of CGeqrefeq:RvM-full for T<=200.}
\expandafter\gdef\csname CGTip@subsec:num-family\endcsname{Section 18.2: The lambdalambda-family and the endpoint (p. 120)}
\expandafter\gdef\csname CGTip@fig:family\endcsname{Figure 5 (p. 121): Top: eta\_k for every 50th k<= N and its distribution. Bottom: moving ordinates in clock units, k<=30, and the functions Z\_lambda(t) on the critical line.}
\expandafter\gdef\csname CGTip@exp:Z1\endcsname{Experiment 18.2 (The endpoint zeta\_Z1 against zeta) (p. 121): Object: zeta\_Z1=P\_Z1/G computed from the N ordinates gamma\_k, and zeta computed by mpmath. Range: sigma in 2,1.25,0.75, 1<= t<=2000 (499 points per line), and the critical line 0.5<= t<=60. Expected: if zeta\_Z1=zeta, the ratio is 1 up to the truncation of the product. Observed: max|zeta\_Z1/zeta-1...}
\expandafter\gdef\csname CGTip@fig:Z1\endcsname{Figure 6 (p. 122): Left: |zeta\_Z1/zeta-1| on three vertical lines. Right: Z\_Z1 and Hardy's Z on the critical line.}
\expandafter\gdef\csname CGTip@exp:product-identity\endcsname{Experiment 18.3 (The value at infinity: the product identity) (p. 121): Object: the ordered partial products P\_K=prod\_k<= K((1)/(4)+gamma\_k\textasciicircum{}2)/((1)/(4)+tau\_k\textasciicircum{}2) of CGeqrefeq:product-identity. Range: K<= N. Expected: by Corollary CGrefcor:product-identity, zeta\_Z1=zeta is the statement P\_K->varpi\_RVM. Observed: P\_K=0.96256, 0.96248, 0.96229, 0.961951, 0.961930, 0.9619...}
\expandafter\gdef\csname CGTip@fig:product-identity\endcsname{Figure 7 (p. 122): The partial products P\_K of CGeqrefeq:product-identity against varpi\_RVM (dashed). Left: K<=10\textasciicircum{}4. Right: (P\_K-varpi\_RVM)*10\textasciicircum{}6 for 10\textasciicircum{}5<= K<= N. The value of the endpoint at +inf is varpi\_Z1=varpi\_RVM/lim\_KP\_K (Theorem CGrefthm:Z1-clock).}
\expandafter\gdef\csname CGTip@subsec:num-transition\endcsname{Section 18.3: The transition theorem and zeta\_RVMzeta\_RVM against zetazeta (p. 122)}
\expandafter\gdef\csname CGTip@exp:transition\endcsname{Experiment 18.4 (Transition theorem) (p. 122): Object: zeta\_RVM, w(s), varpi\_RVM, L. Range: sigma' in 0.25,0.5,1.5, 10<= t<=10\textasciicircum{}4. Expected: Theorem CGrefthm:transition and Corollary CGrefcor:vertical-limit: t |zeta\_RVM/(varpi\_RVM(1+w))-1| bounded, |w|=(t/2pi)\textasciicircum{}-sigma'(1+O(1/t)). Observed: Richardson extrapolation on Re s=2 (which presupposes a...}
\expandafter\gdef\csname CGTip@fig:transition\endcsname{Figure 8 (p. 123): Top: the transition error multiplied by t, and the rescaled vertical limit t\textasciicircum{}min(sigma',1)|1/zeta\_RVM-L|. Bottom: |zeta| and |zeta\_RVM| on sigma=2 and sigma=0.75.}
\expandafter\gdef\csname CGTip@subsec:num-expsums\endcsname{Section 18.4: Exponential sums (p. 123)}
\expandafter\gdef\csname CGTip@exp:expsums\endcsname{Experiment 18.5 (Exponential sums over the two spectra) (p. 123): Object: S(N;log x)=sum\_k<= Nx\textasciicircum{}io\_k for o\_k=tau\_k and o\_k=gamma\_k. Range: N<=2 001 052, x in 2,3,4,6,10. Expected: |S\_RVM|<<1+log N (Theorem CGrefthm:exp-sum); for gamma\_k, Landau's term -(T)/(2pi)(Lambda(x))/( x). Observed: max\_N|S\_RVM(N;log x)|=3.57,3.98,3.95,4.53,9.70; at N=2 001 052, |sum x\textasciicircum{}ig...}
\expandafter\gdef\csname CGTip@fig:expsums\endcsname{Figure 9 (p. 123): The midpoint spectrum carries no prime frequencies (left); the spectrum of zeta resonates at prime powers (right).}
\expandafter\gdef\csname CGTip@subsec:num-primes\endcsname{Section 18.5: Prime counts: pipi and pi\_RVMpi\_RVM (p. 123)}
\expandafter\gdef\csname CGTip@exp:primes\endcsname{Experiment 18.6 (Prime counts and RVM primes) (p. 123): Object: J\_RVM, pi\_RVM, c\_n\textasciicircum{}RVM, p\_j\textasciicircum{}RVM against J, pi, the primes. Range: 2<= x<=100 with all N ordinates of each list (3000 points); 2<= x<=1000 with the first K=17 061 ordinates and a smooth cut (40 000 points). Expected: the explicit formula with the gamma\_k reproduces pi(x); pi\_RVM is continu...}
\expandafter\gdef\csname CGTip@fig:prime-counts\endcsname{Figure 10 (p. 124): pi(x) (black staircase), the explicit formula with the zeros of zeta (blue), pi\_RVM(x) (red) and li(x)-li(2) (green). The distance between the red and green curves is the quantity bounded in Theorem CGrefthm:Lambert-estimate; the ratios are in Figure CGreffig:pi-ratios.}
\expandafter\gdef\csname CGTip@fig:rvm-primes\endcsname{Figure 11 (p. 124): Left: cell masses c\_n\textasciicircum{}RVM, the prime indicator and 1/log n. Right: real RVM prime levels against the primes up to 1000.}
\expandafter\gdef\csname CGTip@exp:pi-ratios\endcsname{Experiment 18.7 (The three ratios and the square-root scale) (p. 124): Object: li/pi\_RVM, pi/li, pi/pi\_RVM, and (pi\_RVM-li)/B, (pi-li)/B with B(x)= x(loglog x)\textasciicircum{}2/log x. Range: the grid of Experiment CGrefexp:primes on [8,1000] (K=17 061 ordinates). Expected: the three ratios tend to 1 (Corollary CGrefcor:three-limits) and (pi\_RVM-li)/B stays bounded (Theorem CGrefth...}
\expandafter\gdef\csname CGTip@fig:pi-ratios\endcsname{Figure 12 (p. 125): Left: the three ratios of Corollary CGrefcor:three-limits. Right: pi\_RVM-li and pi-li in units of the scale B(x) of Theorem CGrefthm:Lambert-estimate. The ratio pi/pi\_RVM (black) is the closest of the three to 1; the ratios with li approach 1 at the rate (1)/(2)li( x)/li(x) of the bias J-pi.}
\expandafter\gdef\csname CGTip@subsec:num-mertens\endcsname{Section 18.6: The Mertens functions and the ternary division (p. 125)}
\expandafter\gdef\csname CGTip@eq:ternary-update\endcsname{Eq. (18.1), p. 125: M\_h(n)=M\_h(n-1)+mu(n)\textasciicircum{}2 varsigma\_n, |M\_h(n)-M(n)| small,}
\expandafter\gdef\csname CGTip@tab:ternary\endcsname{Table 8 (p. 126): The ternary division for n<=15 (the first 40 values are in Appendix CGrefapp:out-s05-mertens). The zero columns of mu and mu\_h coincide; the signs of mu\_h are chosen by comparing M\_h(n-1) with M(n). For n<=12 the two walks agree; at n=13 the target -1.85 lies above M\_h(12)=-2, so mu\_h(13)=+1 whil...}
\expandafter\gdef\csname CGTip@exp:mertens\endcsname{Experiment 18.8 (Mertens functions and the ternary walk) (p. 125): Object: M, M\_RVM\textasciicircum{}crit, M\_h; the explicit formulas. Range: n<=10\textasciicircum{}6; explicit formulas on [2,2000] with all N ordinates of each list. Expected: the explicit formula with the zeros of zeta reproduces M; M\_RVM\textasciicircum{}crit(x)=O(x\textasciicircum{}1/2+eps) (Theorem CGrefthm:unsmoothed-RVM); M\_h tracks M\_RVM\textasciicircum{}crit with the dens...}
\expandafter\gdef\csname CGTip@fig:mertens\endcsname{Figure 13 (p. 126): Left: M (black), M\_RVM\textasciicircum{}crit (red) and M\_h (green) at every integer n<=2000. Right: range of M and mean of M\_RVM\textasciicircum{}crit in blocks of 500 integers up to 10\textasciicircum{}6, with +- n.}
\expandafter\gdef\csname CGTip@fig:ternary\endcsname{Figure 14 (p. 127): Left: mu(n) and mu\_h(n) for n<=60: same zeros, different signs. Right: M, the target M\_RVM\textasciicircum{}crit and M\_h for n<=300.}
\expandafter\gdef\csname CGTip@fig:explicit\endcsname{Figure 15 (p. 127): Left: M(x), its explicit formula from the zeros of zeta and M\_RVM\textasciicircum{}crit(x) from the tau\_k on [2,300]. Right: normalized derivatives at the zeros: erratic for zeta, constant varpi\_RVM for zeta\_RVM.}
\expandafter\gdef\csname CGTip@rem:anticorrelation\endcsname{Remark 18.9 (The opposite phase of M and M\_RVM\textasciicircum{}crit between 10\textasciicircum{}3 and 2*10\textasciicircum{}3) (p. 127): In the left panel of Figure CGreffig:mertens the curves of M and M\_RVM\textasciicircum{}crit are in opposite phase between n10\textasciicircum{}3 and n2*10\textasciicircum{}3. Heuristically, the visible anticorrelation is largely explained by the first mode, 2 x Re(x\textasciicircum{}igamma\_1/(rho\_1zeta'(rho\_1))) and 2 x Re(x\textasciicircum{}itau\_1/(rho\_1\textasciicircum{}RVMzeta\_RVM'(rho\_1\textasciicircum{}RVM)...}
\expandafter\gdef\csname CGTip@subsec:num-rungs\endcsname{Section 18.7: The rung equations: transport, the defect of the clock, and the solution from the clock (p. 127)}
\expandafter\gdef\csname CGTip@exp:transport\endcsname{Experiment 18.10 (Transport on the first rung) (p. 127): Object: R\_ f(1), S\_ f(lambda) and the moments m\_j( f) of CGeqrefeq:rung-functional--CGeqrefeq:transport-moments. Range: the first 20 000 ordinates of each list; max\_k>=1000| c\_ f(tau\_k)|=2.5*10\textasciicircum{}-10 and max\_k>=3000| c\_ f(tau\_k)|=2.1*10\textasciicircum{}-13. Expected: if zeta\_Z1=zeta, then S\_ f(0)= R\_ f(1) (Proposi...}
\expandafter\gdef\csname CGTip@tab:transport\endcsname{Table 9 (p. 128): Experiment CGrefexp:transport. Left: the moments m\_j( f)=sum\_keta\_k\textasciicircum{} j c\textasciicircum{}(j)\_ f(tau\_k) and the partial sums of the series CGeqrefeq:transport-moments for S\_ f(0), compared with the defect R\_ f(1)=-0.042463282161 of the clock. Right: the rung functional along the family.}
\expandafter\gdef\csname CGTip@exp:clock-defect\endcsname{Experiment 18.11 (The defect of the clock by Poisson summation) (p. 128): Object: the two parts of R\_ f(1) in Proposition CGrefprop:clock-defect-poisson. Expected: R\_ f(1)=( k\_ f- I\_0)-2sum\_m>=1(-1)\textasciicircum{}m I\_m. Observed: t\_0=9.676906787; I\_0=0.063533498702 from CGeqrefeq:clock-defect-I0; the mismatch of the smooth densities is k\_ f- I\_0=+0.070110032214, and the oscillating ...}
\expandafter\gdef\csname CGTip@exp:clock-to-zeros\endcsname{Experiment 18.12 (From the clock to the ordinates and to the masses of the primes) (p. 128): Object: a finite section of the system CGeqrefeq:rung-system: the unknowns are the ordinates kappa\_k below a height and the masses a\_n, 2<= n<=11; the equations are the pairings of CGeqrefeq:rung-system with omega\_ test(u)cos(omega u) and omega\_ test(u)sin(omega u), where omega\_ test is a fixed s...}
\expandafter\gdef\csname CGTip@tab:clock-to-zeros\endcsname{Table 10 (p. 129): Experiment CGrefexp:clock-to-zeros. Solution of a finite section of the rung system CGeqrefeq:rung-system, started at the clock. Left: the ordinates. Right: the masses at the points log n, which are not assumed, against log p/(2 n) at prime powers n=p\textasciicircum{}m.}
\expandafter\gdef\csname CGTip@exp:multiscale-jacobian\endcsname{Experiment 18.13 (Multiscale Jacobian of the clock) (p. 129): Object: the minimum midpoint-clock gap, its bound CGeqrefeq:clock-gap-lower, the Ingham critical window 2pi/Delta\_T, and the adjacent-pair bound CGeqrefeq:fixed-window-pair. Input: none of the genuine zeta ordinates; the program tts12\_multiscale\_jacobian.py generates the tau\_k directly from the L...}
\expandafter\gdef\csname CGTip@fig:multiscale-gap\endcsname{Figure 16 (p. 130): Minimum gap of the midpoint clock up to height T and the explicit lower bound 2pi/log(T/2pi). Their ratio is already 1+2*10\textasciicircum{}-8 at K=2 001 052.}
\expandafter\gdef\csname CGTip@fig:fixed-window-degeneracy\endcsname{Figure 17 (p. 130): Upper bound in CGeqrefeq:fixed-window-pair for adjacent clock frequencies on four fixed u-windows. The decay is slow, of order 1/logtau\_k, but it is enough to exclude a height-uniform lower frame bound on any fixed window.}
\expandafter\gdef\csname CGTip@sec:conclusions\endcsname{Section 19: Conclusions and open problems (p. 131)}
\expandafter\gdef\csname CGTip@rem:open-problems\endcsname{Remark 19.1 (Open problems) (p. 135): Prove the identity zeta\_Z1=zeta, equivalently condition (D), equivalently an absolutely convergent Dirichlet expansion of K=zeta\_Z1/zeta\_Z,RVM in Re s>1, equivalently the product identity CGeqrefeq:product-identity, equivalently the cell identity CGeqrefeq:cell-identity with the jump bound of The...}
\expandafter\gdef\csname CGTip@rem:not-proved-list\endcsname{Remark 19.2 (Statements derived in part or conjectural, not proved here) (p. 135): The following statements arose in the work on the final step; none of them is used in this paper. Rates in lambda. Off the powers of 2 the formulas of Proposition CGrefprop:entire-lambda suggest h\_lambda(s)- h\_0(s)=O\_s(lambda), |h\_n(lambda;s)-h\_n(0;s)|<<\_s,epslambda n\textasciicircum{}-sigma-(1)/(2)+eps; the cell...}
\expandafter\gdef\csname CGTip@eq:rates-lambda\endcsname{Eq. (19.1), p. 135: h\_lambda(s)- h\_0(s)=O\_s(lambda), |h\_n(lambda;s)-h\_n(0;s)|<<\_s,epslambda n\textasciicircum{}-sigma-(1)/(2)+eps;}
\expandafter\gdef\csname CGTip@app:code\endcsname{Section A: Python programs (p. 136)}
\expandafter\gdef\csname CGTip@app:code-rvm-common\endcsname{Section A.1: ttrvm\_common.py (p. 136)}
\expandafter\gdef\csname CGTip@app:code-s01-clock-and-family\endcsname{Section A.2: tts01\_clock\_and\_family.py (p. 139)}
\expandafter\gdef\csname CGTip@app:code-s02-transition\endcsname{Section A.3: tts02\_transition.py (p. 141)}
\expandafter\gdef\csname CGTip@app:code-s03-exponential-sums\endcsname{Section A.4: tts03\_exponential\_sums.py (p. 143)}
\expandafter\gdef\csname CGTip@app:code-s04-primes\endcsname{Section A.5: tts04\_primes.py (p. 144)}
\expandafter\gdef\csname CGTip@app:code-s05-mertens\endcsname{Section A.6: tts05\_mertens.py (p. 147)}
\expandafter\gdef\csname CGTip@app:code-s06-prime-ratios\endcsname{Section A.7: tts06\_prime\_ratios.py (p. 151)}
\expandafter\gdef\csname CGTip@app:code-s07-product-identity\endcsname{Section A.8: tts07\_product\_identity.py (p. 152)}
\expandafter\gdef\csname CGTip@app:code-s08-transport-moments\endcsname{Section A.9: tts08\_transport\_moments.py (p. 153)}
\expandafter\gdef\csname CGTip@app:code-s09-clock-defect-split\endcsname{Section A.10: tts09\_clock\_defect\_split.py (p. 155)}
\expandafter\gdef\csname CGTip@app:code-s10-clock-to-zeros\endcsname{Section A.11: tts10\_clock\_to\_zeros.py (p. 156)}
\expandafter\gdef\csname CGTip@app:code-s11-rvm-z1\endcsname{Section A.12: tts11\_rvm\_z1\_difference\_analysis.py (p. 157)}
\expandafter\gdef\csname CGTip@app:code-s12-multiscale\endcsname{Section A.13: tts12\_multiscale\_jacobian.py (p. 162)}
\expandafter\gdef\csname CGTip@app:code-s13-bounded\endcsname{Section A.14: tts13\_nested\_bounded\_endpoint.py (p. 165)}
\expandafter\gdef\csname CGTip@app:output\endcsname{Section B: Console output of the programs (p. 168)}
\expandafter\gdef\csname CGTip@app:out-s01-clock-and-family\endcsname{Section B.1: tts01\_clock\_and\_family (p. 169)}
\expandafter\gdef\csname CGTip@app:out-s01b-clock-zoom\endcsname{Section B.2: tts01b\_clock\_zoom (p. 169)}
\expandafter\gdef\csname CGTip@app:out-s02-transition\endcsname{Section B.3: tts02\_transition (p. 170)}
\expandafter\gdef\csname CGTip@app:out-s03-exponential-sums\endcsname{Section B.4: tts03\_exponential\_sums (p. 170)}
\expandafter\gdef\csname CGTip@app:out-s04-primes\endcsname{Section B.5: tts04\_primes (p. 170)}
\expandafter\gdef\csname CGTip@app:out-s05-mertens\endcsname{Section B.6: tts05\_mertens (p. 171)}
\expandafter\gdef\csname CGTip@app:out-s06-prime-ratios\endcsname{Section B.7: tts06\_prime\_ratios (p. 172)}
\expandafter\gdef\csname CGTip@app:out-s07-product-identity\endcsname{Section B.8: tts07\_product\_identity (p. 173)}
\expandafter\gdef\csname CGTip@app:out-s08-transport-moments\endcsname{Section B.9: tts08\_transport\_moments (p. 173)}
\expandafter\gdef\csname CGTip@app:out-s09-clock-defect-split\endcsname{Section B.10: tts09\_clock\_defect\_split (p. 174)}
\expandafter\gdef\csname CGTip@app:out-s10-clock-to-zeros\endcsname{Section B.11: tts10\_clock\_to\_zeros (p. 174)}
\expandafter\gdef\csname CGTip@app:code-routeB\endcsname{Section C: Python programs of Section CGrefsec:routeB-finite (p. 178)}
\expandafter\gdef\csname CGTip@app:code-routeB-core\endcsname{Section C.1: ttcore.py (p. 178)}
\expandafter\gdef\csname CGTip@app:code-routeB-zeros\endcsname{Section C.2: ttzeros.py (p. 178)}
\expandafter\gdef\csname CGTip@app:code-routeB-gonek-fast\endcsname{Section C.3: ttgonek\_fast.py (p. 178)}
\expandafter\gdef\csname CGTip@app:code-routeB-horizonG\endcsname{Section C.4: tthorizonG.py (p. 179)}
\expandafter\gdef\csname CGTip@app:code-routeB-cert7\endcsname{Section C.5: ttcert7.py (p. 179)}
\expandafter\gdef\csname CGTip@app:code-routeB-entry\endcsname{Section C.6: ttentry.py (p. 180)}
\expandafter\gdef\csname CGTip@app:code-routeB-laplace-check\endcsname{Section C.7: ttlaplace\_check.py (p. 181)}
\expandafter\gdef\csname CGTip@app:code-routeB-figs\endcsname{Section C.8: ttfigs.py (p. 182)}
\expandafter\gdef\csname CGTip@app:out-routeB\endcsname{Section D: Console output of the programs of Section CGrefsec:routeB-finite (p. 182)}
\expandafter\gdef\csname CGTip@app:out-routeB-results\endcsname{Section D.1: ttRESULTS.txt (p. 182)}
\expandafter\gdef\csname CGTip@app:out-routeB-cert7\endcsname{Section D.2: ttout\_cert7.txt (p. 183)}
\expandafter\gdef\csname CGTip@app:out-routeB-entry\endcsname{Section D.3: ttout\_entry.txt (p. 184)}
\expandafter\gdef\csname CGTip@app:out-routeB-certG\endcsname{Section D.4: ttout\_certG.txt (p. 184)}
\expandafter\gdef\csname CGTip@app:out-routeB-cert3\endcsname{Section D.5: ttout\_cert3.txt (p. 185)}
\expandafter\gdef\csname CGTip@app:out-routeB-gonek\endcsname{Section D.6: ttout\_gonek.txt (p. 185)}
\expandafter\gdef\csname CGTip@app:out-routeB-laplace-check\endcsname{Section D.7: ttout\_laplace\_check.txt (p. 185)}
\expandafter\gdef\csname CGTip@sec:appendix-endpoint-v18\endcsname{Section E: Endpoint-only closure, interpolation tests, and autonomous residual questions (p. 186)}
\expandafter\gdef\csname CGTip@sec:v18-addendum\endcsname{Section E: Endpoint-only closure, interpolation tests, and autonomous residual questionssec:appendix-endpoint-v18 (p. 186)}
\expandafter\gdef\csname CGTip@subsec:three-interpolations-v18\endcsname{Section E.1: Three interpolations and their status (p. 186)}
\expandafter\gdef\csname CGTip@eq:additive-family-v18\endcsname{Eq. (E.1), p. 186: G\_varepsilon(s):=(1-varepsilon)zeta\_RVM(s)+varepsilonzeta(s), 0<=varepsilon<=1.}
\expandafter\gdef\csname CGTip@prop:additive-level-v18\endcsname{Proposition E.1 (The additive family is a level-set reformulation) (p. 186): Away from the zeros of zeta\_RVM, G\_varepsilon(s)=zeta\_RVM(s)((1-varepsilon)+\_RVM(s)). Hence for 0<varepsilon<1 every zero of G\_varepsilon which is not a zero of zeta\_RVM lies on the level set Q\_RVM(s)=-(1-varepsilon)/(varepsilon). On the critical line e\textasciicircum{}ivartheta(t)G\_varepsilon((1)/(2)+it) is rea...}
\expandafter\gdef\csname CGTip@eq:additive-level-v18\endcsname{Eq. (E.2), p. 186: G\_varepsilon(s)=zeta\_RVM(s)((1-varepsilon)+\_RVM(s)).}
\expandafter\gdef\csname CGTip@eq:Q-negative-level-v18\endcsname{Eq. (E.3), p. 186: Q\_RVM(s)=-(1-varepsilon)/(varepsilon).}
\expandafter\gdef\csname CGTip@subsec:first-transition-v18\endcsname{Section E.2: The clock as the first transition term (p. 186)}
\expandafter\gdef\csname CGTip@eq:rvm-first-rs-v18\endcsname{Eq. (E.4), p. 186: zeta\_RVM(s)=varpi\_RVM(1+chi(s))(1+O(|t|\textasciicircum{}-1)).}
\expandafter\gdef\csname CGTip@eq:first-rs-hardy-v18\endcsname{Eq. (E.5), p. 186: e\textasciicircum{}ivartheta(t)(1+chi((1)/(2)+it))=2cosvartheta(t).}
\expandafter\gdef\csname CGTip@eq:finite-additive-v18\endcsname{Eq. (E.6), p. 187: P\_varepsilon,M(s):=1+\_2<= n<= Mn\textasciicircum{}-s, R\_varepsilon,M(s):= P\_varepsilon,M(s)+chi(s) P\_varepsilon,M(1-s),}
\expandafter\gdef\csname CGTip@eq:finite-additive-hardy-v18\endcsname{Eq. (E.7), p. 187: e\textasciicircum{}ivartheta(t) R\_varepsilon,M((1)/(2)+it) =2Re (e\textasciicircum{}ivartheta(t) P\_varepsilon,M((1)/(2)+it)).}
\expandafter\gdef\csname CGTip@eq:dominance-v18\endcsname{Eq. (E.8), p. 187: \_2<= n<= Mn\textasciicircum{}-1/2<1,}
\expandafter\gdef\csname CGTip@subsec:live-v18\endcsname{Section E.3: What remains genuinely live (p. 187)}
\expandafter\gdef\csname CGTip@eq:clock-injectivity-v18\endcsname{Eq. (E.9), p. 187: H(z)=nu\_8+sum\_k>=1\_k((tau\_k)/(z\textasciicircum{}2+tau\_k\textasciicircum{}2)-(1)/(tau\_k)) =sum\_n>=2( a\_n)/(log n) n\textasciicircum{}-z}
\expandafter\gdef\csname CGTip@subsec:rung-apriori-v18\endcsname{Section E.4: A priori bounded displacement for exact rung solutions (p. 187)}
\expandafter\gdef\csname CGTip@prop:rung-apriori-v18\endcsname{Proposition E.2 (A priori displacement bound for an exact rung solution) (p. 188): Fix t in [0,1]. Let kappa=(kappa\_k)\_k>=1 satisfy 0<kappa\_1<=kappa\_2<=*s and kappa\_k->inf, and let ( a\_n)\_n>=2 be real masses satisfying, in D'(0,inf), sum\_k>=1cos(kappa\_k u) -(1- t)sum\_k>=1cos(tau\_k u)- t k(u) =-sum\_n>=2 a\_n delta(u-log n). Assume that for every epsilon\_0>0 the local count and th...}
\expandafter\gdef\csname CGTip@eq:rung-apriori-system-v18\endcsname{Eq. (E.10), p. 188: sum\_k>=1cos(kappa\_k u) -(1- t)sum\_k>=1cos(tau\_k u)- t k(u) =-sum\_n>=2 a\_n delta(u-log n).}
\expandafter\gdef\csname CGTip@eq:rung-apriori-density-v18\endcsname{Eq. (E.11), p. 188: N\_kappa(T+1)-N\_kappa(T)<<\_epsilon\_0(1+T)\textasciicircum{}epsilon\_0 (T>=0),}
\expandafter\gdef\csname CGTip@eq:rung-apriori-mass-v18\endcsname{Eq. (E.12), p. 188: sum\_X<n<=2X| a\_n|\textasciicircum{}2<<\_epsilon\_0X\textasciicircum{}1+epsilon\_0 (X>=1).}
\expandafter\gdef\csname CGTip@eq:rung-apriori-b-v18\endcsname{Eq. (E.13), p. 188: b:=lim\_sigma->+inf[sum\_k>=1(sigma-(1)/(2))/((sigma-(1)/(2))\textasciicircum{}2+kappa\_k\textasciicircum{}2)-(1- t)sum\_k>=1(sigma-(1)/(2))/((sigma-(1)/(2))\textasciicircum{}2+tau\_k\textasciicircum{}2)- t (G')/(2G)(sigma)]}
\expandafter\gdef\csname CGTip@eq:rung-apriori-count-v18\endcsname{Eq. (E.14), p. 188: |N\_kappa(T)-n\_RVM(T)-(2b)/(pi)T|<= C\textasciicircum{}*log(2+T) (T>=0)}
\expandafter\gdef\csname CGTip@eq:rung-apriori-bound-v18\endcsname{Eq. (E.15), p. 188: |kappa\_k-tau\_k+(4b kappa\_k)/(log(kappa\_k/2pi))|<= C\textasciicircum{}* (k>=1).}
\expandafter\gdef\csname CGTip@eq:rung-laplace-S-v18\endcsname{Eq. (E.16), p. 188: LS\_kappa(s)=sum\_k>=1(s-(1)/(2))/((s-(1)/(2))\textasciicircum{}2+kappa\_k\textasciicircum{}2)=(1)/(2) (P\_kappa')/(P\_kappa)(s) (Re s>(1)/(2)).}
\expandafter\gdef\csname CGTip@eq:rung-E-v18\endcsname{Eq. (E.17), p. 189: E:=S\_kappa-(1- t)S\_- t K+A}
\expandafter\gdef\csname CGTip@eq:rung-laplace-identity-v18\endcsname{Eq. (E.18), p. 189: (1)/(2)(P\_kappa')/(P\_kappa)(s)-(1- t) (1)/(2)((P\textasciicircum{}crit\_RVM)')/(P\textasciicircum{}crit\_RVM)(s)- t (1)/(2)(G')/(G)(s)+sum\_n>=2 a\_nn\textasciicircum{}(1)/(2)-s=p(s-(1)/(2)) (Re s>(3)/(2)).}
\expandafter\gdef\csname CGTip@eq:rung-three-factor-v18\endcsname{Eq. (E.19), p. 189: P\_kappa(s) =C\_ t(P\_RVM\textasciicircum{}crit(s))\textasciicircum{}1- t G(s)\textasciicircum{} t Y\_ t(s) e\textasciicircum{}2b s =C\_ t G(s) zeta\_RVM(s)\textasciicircum{}1- t Y\_ t(s) e\textasciicircum{}2b s (Re s>(3)/(2)),}
\expandafter\gdef\csname CGTip@eq:rung-Y-v18\endcsname{Eq. (E.20), p. 189: Y\_ t(s)=exp (sum\_n>=2(2 a\_n)/(log n) n\textasciicircum{}1/2-s),}
\expandafter\gdef\csname CGTip@eq:rung-B-v18\endcsname{Eq. (E.21), p. 189: sup\_Re s>=2|log Y\_ t(s)|<= B\_2<inf,}
\expandafter\gdef\csname CGTip@eq:rung-line-v18\endcsname{Eq. (E.22), p. 190: |log|P\_kappa(s)|-log|G(s)||<= K\_ E:=|log C\_ t|+c\_RVM+B\_2+8|b| (2<=Re s<=4),}
\expandafter\gdef\csname CGTip@eq:rung-strip-v18\endcsname{Eq. (E.23), p. 190: log|P\_kappa(s)|<= K\_ E+ K\_ E'+log|s+2|+log|G(s)|-log|s-1| .}
\expandafter\gdef\csname CGTip@eq:rung-nT-v18\endcsname{Eq. (E.24), p. 190: n\_T<=(2 K\_ E+ K\_ E'+2log(T+8)+2c\_5+log2)/(log(4/3)),}
\expandafter\gdef\csname CGTip@eq:rung-count-v18\endcsname{Eq. (E.25), p. 191: |N\_kappa(T)-n\_RVM(T)-(2b)/(pi) T| <= n\_T+1+(2c\_RVM+2B\_2+c\_6)/(pi)+(1)/(T)+(1)/(2) (T>= T\_0, Tkappa\_k),}
\expandafter\gdef\csname CGTip@eq:rung-quantile-v18\endcsname{Eq. (E.26), p. 191: N\_RVM(kappa\_k)=k-(2b)/(pi)kappa\_k+O(C\textasciicircum{}*log(2+kappa\_k)),}
\expandafter\gdef\csname CGTip@cor:rung-apriori-b0\endcsname{Corollary E.3 (Exact solutions without mass at the origin) (p. 191): Under the hypotheses of Proposition CGrefprop:rung-apriori-v18, if the extension of the left side of CGeqrefeq:rung-apriori-system-v18 constructed in the proof has no mass at u=0, that is, b=0, then |N\_kappa(T)-n\_RVM(T)|<= C\textasciicircum{}*log(2+T) and |kappa\_k-tau\_k|<= C\textasciicircum{}* for all k, with C\textasciicircum{}* depending only o...}
\expandafter\gdef\csname CGTip@rem:rung-apriori-scope-v18\endcsname{Remark E.4 (p. 191): Proposition CGrefprop:rung-apriori-v18 is an a priori theorem, not an existence theorem. The exact-atomic interpolation has no established interior solution for 0<theta<1; the proposition only says that any exact solution satisfying its hypotheses lies in one uniform bounded-displacement class. T...}
\expandafter\gdef\csname CGTip@subsec:finite-diagnostics-v18\endcsname{Section E.5: Additional finite-section diagnostics (p. 192)}
\expandafter\gdef\csname CGTip@eq:bounded-section-v18\endcsname{Eq. (E.27), p. 192: |kappa\_k-tau\_k|<=1.2, |a\_n|<=0.5,}
\expandafter\gdef\csname CGTip@fig:nested-bounded-v18\endcsname{Figure 18 (p. 193): Nested bounded endpoint sections. Each curve shows the residual of one solved section on all permanent test blocks already introduced. The ordinate shifts and free masses are bounded during the solve. This is a finite computation and is not an existence proof for the infinite rung system.}
\expandafter\gdef\csname CGTip@fig:bounded-masses-v18\endcsname{Figure 19 (p. 193): Recovered free masses in the interior n<=15 of bounded endpoint sections. The von Mangoldt values are shown only as a post-solution diagnostic and are not supplied to the solver; section-boundary masses are omitted from the plot.}
\makeatother


\hypersetup{
  colorlinks=false,
  pdfborder={0 0 1},
  linkbordercolor={1 0 0},
  citebordercolor={0 0.3 1},
  urlbordercolor={0 0.6 0.6},
  filebordercolor={0 0.6 0.6},
  linktoc=all,
  bookmarksnumbered=true,
  bookmarksopen=true,
  bookmarksopenlevel=1,
  pdfstartview=FitH,
  pdftitle={The midpoint Riemann-von Mangoldt synthetic zeta system},
  pdfauthor={Rotkiv Ayn\"asac},
  pdfsubject={Midpoint RVM zeta function, the lambda-family to the endpoint Z1, Hamburger closure, and the Mertens bound},
  pdfkeywords={Riemann zeta function, Riemann-von Mangoldt clock, Hadamard product, Hamburger theorem, Mertens function}
}
% Back-references in the bibliography (placed after the DOI).
\renewcommand*{\backref}[1]{}
\renewcommand*{\backrefalt}[4]{%
  \ifcase #1 %
    \relax
  \or
    \ (Cited on p.~#2.)%
  \else
    \ (Cited on pp.~#2.)%
  \fi}

% In \cite[note]{key} the note (e.g. a section number) is linked too.
\makeatletter
\let\CG@origcite\cite
\renewcommand*{\cite}[2][]{\def\CG@citekey{#2}%
  \ifx\relax#1\relax\CG@origcite{#2}\else\CG@origcite[#1]{#2}\fi}
\def\@cite#1#2{{\@citestyle[\citeform{#1}\if@tempswa, \hyper@link{cite}{cite.\CG@citekey}{#2}\fi]}}
\makeatother

\geometry{margin=1in}
\numberwithin{equation}{section}
\setcounter{tocdepth}{3}
\setcounter{secnumdepth}{3}
\setlength{\emergencystretch}{1.5em}

\theoremstyle{plain}
\newtheorem{theorem}{Theorem}[section]
\newtheorem{proposition}[theorem]{Proposition}
\newtheorem{lemma}[theorem]{Lemma}
\newtheorem{corollary}[theorem]{Corollary}
\newtheorem{conjecture}[theorem]{Conjecture}
\theoremstyle{definition}
\newtheorem{definition}[theorem]{Definition}
\newtheorem{hypothesis}[theorem]{Hypothesis}
\newtheorem{heuristic}[theorem]{Heuristic}
\newtheorem{experiment}[theorem]{Experiment}
\theoremstyle{remark}
\newtheorem{remark}[theorem]{Remark}
\newtheorem{warning}[theorem]{Warning}

\newcommand{\R}{\mathbb R}
\newcommand{\C}{\mathbb C}
\newcommand{\N}{\mathbb N}
\newcommand{\Z}{\mathbb Z}
\newcommand{\eps}{\varepsilon}
\newcommand{\li}{\operatorname{li}}
\newcommand{\Log}{\operatorname{Log}}
\newcommand{\Ei}{\operatorname{Ei}}
\newcommand{\Res}{\operatorname*{Res}}
\newcommand{\Mell}{\ast_{\rm Mel}}
\newcommand{\RVM}{\mathrm{RVM}}
\newcommand{\crit}{\mathrm{crit}}
\newcommand{\pts}{\mathrm{pts}}
\newcommand{\pathc}{\mathrm{path}}
\newcommand{\NR}{\mathcal N_{\RVM}}
\newcommand{\vpi}{\varpi_{\RVM}}
\newcommand{\Zo}{Z1}
\providecommand{\NoCaseChange}[1]{\lowercase{#1}}

\begin{document}

\title[The midpoint RVM synthetic zeta system]{The midpoint Riemann--von Mangoldt synthetic zeta system:\\
canonical products, exponential sums, prime counts, Mertens transforms,\\
and the transformation to $\zeta_{Z1}$}
\author{Rotkiv Ayn\"asac\\[2pt]\mbox{}}
\date{12 October 2026 -- version 22}

% The abstract is typeset after \maketitle as ordinary text so that it may run over several pages (in version 18 its end was cut off by the title box).
\makeatletter
\renewenvironment{abstract}{\par\begingroup\small\noindent{\scshape Abstract.}\enspace\ignorespaces}{\par\endgroup\medskip}
\makeatother


\maketitle
\markboth{ROTKIV AYN\"ASAC}{THE MIDPOINT RVM SYNTHETIC ZETA SYSTEM}

\begin{abstract}
We study the zeta-like system generated by the midpoint inverse of the smooth Riemann--von Mangoldt counting law.  Its nontrivial zeros are $\tfrac12\pm i\tau_k$, where $\tau_k$ is the midpoint ordinate of the clock.  We prove normal convergence of the canonical product, the functional symmetry, the exact zero divisor, a uniform logarithmic bound for the exponential sums of the $\tau_k$, compact-uniform existence of the explicit-formula prime-power count with an explicit branch of $\Ei$, exact finite M\"obius inversion to a synthetic prime count $\pi_{\RVM}$, and the unconditional square-root estimate $\pi_{\RVM}(x)-\li(x)=O(\sqrt x(\log\log x)^2/\log x)$; with the prime number theorem this gives $\li/\pi_{\RVM}\to1$, $\pi/\li\to1$ and $\pi/\pi_{\RVM}\to1$.  The transition theorem $\zeta_{\RVM}=\varpi_{\RVM}(1+w)(1+O(1/|t|))$ is proved in full; it gives a constant vertical limit of $1/\zeta_{\RVM}$ on every line to the right of $\tfrac12$, excludes absolutely convergent Dirichlet expansions of $\zeta_{\RVM}^{\pm1}$, and proves the unsmoothed synthetic Mertens bound $M_{\RVM}^{\crit}(x)=L+O_\eps(x^{1/2+\eps})$, $L=1/\varpi_{\RVM}$.  A contour-free form of the same decomposition shows that $\zeta_{\RVM}$ is bounded above and below on every half-plane $\Re s\ge\sigma_0>1$ and tends to $\varpi_{\RVM}$ as $\Re s\to+\infty$.

The core of the method is a transformation of $\zeta_{\RVM}$ depending on a parameter $\lambda\in[0,1]$: the ordinates move along $\gamma_k(\lambda)=(1-\lambda)\gamma_k+\lambda\tau_k$, where $\gamma_k$ are the ordinates of the zeros of the Riemann zeta function, every member $\zeta_\lambda$ has all its nontrivial zeros on $\Re s=\tfrac12$ by construction, $\zeta_1=\zeta_{\RVM}$, and $\zeta_0=\zeta_{Z1}$.  It is impossible for $\zeta_{Z1}$, or for any $\zeta_\lambda$, to have a zero off the critical line other than the trivial zeros.  All differences between the moving zetas lie in the numbers $\gamma_k-\tau_k$.  The final step of the method is the identity $\zeta_{Z1}=\zeta$.  This identity is the single and complete test: while $\zeta_{Z1}=\zeta$ there is nothing outside $\Re s=\tfrac12$ to speak of, and $\zeta_{Z1}\neq\zeta$ is the only possible obstruction.

For completeness we reprove the classical implication RH $\Rightarrow M(x)=O_\eps(x^{1/2+\eps})$ in the endpoint notation: if $\zeta_{Z1}=\zeta$, then $M(x)=O_\eps(x^{1/2+\eps})$, with a self-contained proof (Borel--Carath\'eodory, three circles, truncated Perron).  By Hamburger's theorem the identity $\zeta_{Z1}=\zeta$ is equivalent to the condition (D) that $\zeta_{Z1}$ is an absolutely convergent Dirichlet series in $\Re s>1$.  We prove three further forms of (D): the value of the endpoint at infinity is $\varpi_{Z1}=\lim_{\sigma\to+\infty}\zeta_{Z1}(\sigma)=\varpi_{\RVM}\lim_{K\to\infty}\prod_{k\le K}(\frac14+\tau_k^2)/(\frac14+\gamma_k^2)\ge1$, and (D) holds exactly when $\varpi_{Z1}=1$, that is, when the ordered product $\lim_{K\to\infty}\prod_{k\le K}(\frac14+\gamma_k^2)/(\frac14+\tau_k^2)=\varpi_{\RVM}$; the prime-power counts of $\zeta_{\RVM}$ and $\zeta_{Z1}$ are the Mellin transforms of their logarithms and satisfy a distributional explicit formula, which proves the Mellin--Stieltjes identification for the integer weights; and (D) holds exactly when the prime-power count of $\zeta_{Z1}$ is constant between consecutive integers and its jumps satisfy $\sum_{n\le X}|\mathsf b_n|=O_\eps(X^{1+\eps})$, an identity on each cell between the displacements $\gamma_k-\tau_k$ and the clock.  Along the family, $\log\varpi_\lambda$ has derivative $-\Upsilon(\lambda)$, the generator at infinity, and (D) holds exactly when $\int_0^1\Upsilon\,d\lambda=\log(1/\varpi_{\RVM})$.  Finally we atomize every member of the family by smooth cells: the cell masses have mean square $X^{1+\eps}$ on every dyadic block, uniformly in $\lambda$, so every atomized function $\tilde\zeta_{\Z,\lambda}=\zeta_\lambda e^{\tilde h_\lambda}$, the endpoint $\lambda=0$ included, is an absolutely convergent Dirichlet series in $\Re s>1$; the joining factor $\tilde h_\lambda$ is written cell by cell and is holomorphic in $\Re s>\frac12$, and (D) is reduced to a single statement at $\lambda=0$, the vanishing of every cell defect, which is the remaining step of the programme.  We prove further that no member $\zeta_\lambda$ with $0<\lambda\le1$ is a Dirichlet series, that $\lambda\mapsto\log\zeta_\lambda(\sigma)$ is strictly convex for every real $\sigma>1$, and that $\zeta_{Z1}(\sigma)\ge\zeta(\sigma)$ for every real $\sigma>1$, with equality at one point exactly when $\zeta_{Z1}=\zeta$.  Hamburger's theorem is extended to the pair $(\zeta_\lambda,\tilde\zeta_{\Z,\lambda})$: $\tilde\zeta_{\Z,\lambda}=\zeta$ holds as soon as the joining factor is symmetric under $s\mapsto1-s$, equivalently when the atomized coefficients satisfy the theta relation of $\zeta$, equivalently when the vertical mean square of $\tilde\zeta_{\Z,\lambda}/\zeta$ on a line $\Re s=\sigma>1$, which is always at least $1$, equals $1$; at the clock end the symmetry fails, and $\tilde\zeta_{\Z,\lambda}=\zeta$ holds for at most finitely many $\lambda$.  (D) holds exactly when $\zeta_{Z1}$ is the Mellin--Stieltjes transform of a positive measure: the functional equation then gives $\varpi_{Z1}\le1$ by an extremal summation argument, and in general $\varpi_\lambda-1$ is one explicit number computed from the summation formula of $\zeta_\lambda$.  At the clock end the absolutely continuous part of $d\pi_{\RVM}$ has density $\frac1{\log y}+2\cos(2\pi y)-2^{-1/2}y^{-1/4}\cos(4\pi\sqrt y+\frac\pi4)+O(y^{-1/3})$, while $d\pi_{\RVM}$ also has the explicit atoms $\mu(j)J_{\RVM}(2)\delta_{2^j}/j$.  Finally the constructions are carried out for every list of ordinates within a bounded distance of the clock: the lists whose continuous and atomized zeta functions coincide form a set with at most one element, which is non-empty exactly when (D) holds, so that the final step is an existence statement in a compact space of lists.  In version 18 we retain these equivalences but remove the earlier continuation claim: the interpolation tests show that the exact-atomic linear path has no established interior solution, while the additive interpolation is only a level-set reformulation of $\zeta/\zeta_{\RVM}$.  We add a corrected clock-only injectivity problem, a rigorously stated a priori bounded-displacement proposition for any exact rung solution satisfying the stated local-density hypothesis, and documented bounded finite-section diagnostics.  None of these proves (D).

\emph{What is proved and what is not.}  Every statement above is proved in this paper without any hypothesis, with one exception made explicit where it occurs: the strict inequality $\varpi_{\RVM}<1$ for the transition constant (Hypothesis~\CGref{hyp:varpi-less-one}) is supported numerically and not proved, and the five statements that use it are conditional on it.  At the clock end we further prove a multiscale theorem for the linearized rung system: up to height $T$ the ordinate block has a uniform lower frame bound on $u$-windows of length $(1+\varkappa_{\rm win})\log(T/2\pi)$, no $T$-independent lower bound is possible on a fixed window, the infinite clock linearization has dense range on every finite window, and smooth transport tails beyond height $T$ are $O_{A,\mathsf f}(T^{1-A}\log T)$.  These statements use only the midpoint clock and no individual displacement $\eta_k=\tau_k-\gamma_k$.  They rule out a fixed-window implicit-function argument.  The later interpolation tests show that no multiscale continuation theorem is presently established; the remaining endpoint problem is not replaced by a path of exact intermediate objects.

The identity $\zeta_{Z1}=\zeta$, equivalently (D), equivalently each of the criteria above, is not proved; it is the open step of the method.  Accordingly this paper proves neither the Riemann hypothesis nor (D): it proves a system of equivalences, rigidity and uniqueness statements, and a priori bounds around that step, and it says at each place which statements are theorems, which are conditional on the single numerical hypothesis $\varpi_{\RVM}<1$, which are equivalent to (D), and which are computations.

A final section treats route B for finite sections of the rung system: compactness of near-solutions in a fixed box; with the von Mangoldt masses prescribed, the Laplace transform of the residual is identified as $\frac12(P_{\boldsymbol\kappa}'/P_{\boldsymbol\kappa}-\xi'/\xi)$ and every exact solution at bounded distance from the clock is the zeta list, so the finite-section statement B$'$ is equivalent to the Riemann hypothesis (an equivalence classical in substance); a certificate theorem for staged sections (Newton--Banach with a frozen right inverse); Gonek's symmetrised Euler products as explicit lists in a tube around the clock; an a priori counting bound for quasi-solutions; a Kadec--Avdonin criterion for $(A_2)$ with blocks of variable length, and the unconditional failure of $(A_2)$ for the unfolded zeta ordinates on every horizontal line; and floating-point candidate Banach certificates of exact square sections up to height $487$, not verified by interval arithmetic, with the candidate basin shrinking like $T^{-2}$.
\end{abstract}

\tableofcontents

\section*{Notation}
\label{sec:notation}

Every symbol below has a single meaning throughout the paper.  The last column links to the defining equation; hovering over the link shows the formula.  Symbols introduced inside a single proof (indices $j,k_0,m$, auxiliary variables $v,y$, thresholds $t_0,t_1,t_2$, radii $r_1,r_2,r_3$) are local to that proof and do not appear elsewhere; so are the following local symbols, which coincide in shape with global ones but never with their meaning: $E_k$ and $T_m$ in the proofs of Theorems~\CGref{thm:exp-sum} and~\CGref{thm:unit-cell}, $\phi$ in Lemmas~\CGref{lem:A-RVM} and~\CGref{lem:TR-clock}, $K_m^\pm$ and $r$ in Section~\CGref{sec:transition}, $J$, $S$, $C_0$, $E$ and $A$ in the proofs of Section~\CGref{sec:atomized}, $U$ as a window length in Subsection~\CGref{subsec:multiscale-jacobian}, $\mathsf X_{\mathfrak P}$ for the prime diagonal operator, $\mathfrak g_1$ in Subsection~\CGref{subsec:rb-application}, and $c_1$ in Propositions~\CGref{prop:Mellin-Mertens} and~\CGref{prop:one-number}.

\medskip
{\small
\begin{longtable}{>{\raggedright\arraybackslash}p{0.27\textwidth}>{\raggedright\arraybackslash}p{0.53\textwidth}>{\raggedright\arraybackslash}p{0.12\textwidth}}
\toprule
Symbol & Meaning & Defined\\
\midrule
\endfirsthead
\toprule
Symbol & Meaning & Defined\\
\midrule
\endhead
\bottomrule
\endfoot
\multicolumn{3}{l}{\emph{Variables and the Riemann zeta function}}\\
$s=\sigma+it$ & complex variable, $\sigma=\Re s$, $t=\Im s$ & --\\
$z=s-\tfrac12=\sigma'+it$ & centred variable, $\sigma'=\sigma-\tfrac12$ & \CGeqref{eq:transition-zw}\\
$q(s)=s(s-1)$ & symmetric quadratic, $q(1-s)=q(s)$ & \CGeqref{eq:q-def}\\
$\zeta$, $\rho$ & Riemann zeta function; $\rho$ runs over its nontrivial zeros with multiplicity & \S\CGref{sec:clock}\\
$\gamma_k$ & $k$th positive ordinate of the nontrivial zeros of $\zeta$, with multiplicity, nondecreasing & \CGeqref{eq:zero-count}\\
$N_\zeta(T)$, $N_\zeta^{*}(T)$ & zero count up to height $T$; its symmetric version & \CGeqref{eq:zero-count}, \CGeqref{eq:symmetric-count}\\
$S(T)$ & remainder in the Riemann--von Mangoldt formula & \CGeqref{eq:RvM-formula}\\
$G(s)$ & $s(s-1)\Gamma(s/2)\pi^{-s/2}$ & \CGeqref{eq:Gdef}\\
$P_\zeta=2\xi$ & completed product $G\zeta$ & \CGeqref{eq:Pzeta}\\
$\mu$, $M(x)$ & M\"obius function; Mertens function $\sum_{n\le x}\mu(n)$ & \CGeqref{eq:mertens-def}\\
$\pi(x)$, $\li$ & prime-counting function; $\li(x):=\operatorname{PV}\!\int_0^x dt/\log t$ & --\\
$\chi(s)$, $\vartheta$ & functional-equation factor; Riemann--Siegel theta function & \CGeqref{eq:chi}, \CGeqref{eq:theta-def}\\
\midrule
\multicolumn{3}{l}{\emph{The midpoint clock}}\\
$\NR(T)$ & midpoint Riemann--von Mangoldt clock & \CGeqref{eq:clock}\\
$\tau_k$ & midpoint ordinates: the unique solutions in $[2\pi,\infty)$ of $\NR(\tau_k)=k-\tfrac12$ & \CGeqref{eq:midpoint-condition}\\
$W_0$ & principal branch of Lambert's $W$-function & \CGeqref{eq:lambert}\\
$\mathcal T$ & inverse function of the clock & \CGeqref{eq:inverse-clock}\\
$\ell_k$ & $\log(\tau_k/2\pi)$ & text after \CGeqref{eq:inverse-clock}\\
$n_{\RVM}(H)$ & number of $\tau_k\le H$ & \CGeqref{eq:n-RVM}\\
$\eta_k$, $C_\eta$ & $\eta_k=\tau_k-\gamma_k$; a bound $|\eta_k|\le C_\eta$ & \CGeqref{eq:eta-def}, \CGeqref{eq:bounded-displacement}\\
\midrule
\multicolumn{3}{l}{\emph{The synthetic zeta function and its sums}}\\
$P_{\RVM}^{\crit}$, $\xi_{\RVM}$, $\zeta_{\RVM}$ & canonical product, completed and synthetic zeta & \CGeqref{eq:Pcrit}, \CGeqref{eq:zetaRVM}\\
$\rho_k^{\RVM}$ & synthetic zero $\tfrac12+i\tau_k$ & \CGeqref{eq:synthetic-divisor}\\
$S_{\RVM}(N;u)$, $C_{\RVM}^{\pts}$, $C_{\RVM}^{\pathc}$ & exponential sum; point and path centroids & \CGeqref{eq:S-RVM}, \CGeqref{eq:path-centroid}\\
$u_1,u_2$ & a frequency interval $[u_1,u_2]\subset(0,\infty)$ & \CGeqref{eq:exp-sum-bound}\\
$M_H(x)$, $M_H^{\RVM}(x)$ & finite-height phase averages (not Mertens functions) & \CGeqref{eq:MH-genuine}, \CGeqref{eq:MH-RVM}\\
$\Ei$ & exponential integral on $\Im w\ne0$ & \CGeqref{eq:Ei-def}\\
$J_{\RVM}$, $J_{\RVM}^{*}$, $\pi_{\RVM}$ & explicit-formula prime-power count; truncation; synthetic prime count & \CGeqref{eq:J-RVM}, \CGeqref{eq:J-star}, \CGeqref{eq:pi-RVM}\\
$\mathcal A_{\RVM}(u)$, $\ell(t)$ & weighted clock sum $\sum_ke^{iu\tau_k}/\tau_k$; $\log(t/2\pi)$ & \CGeqref{eq:A-RVM}\\
$\mathfrak s_1$; $\mathfrak e$, $\mathfrak e_1$ & primitive of the midpoint sawtooth; $N_\zeta-\NR$ and its primitive & \CGeqref{eq:saw-primitive}, \CGeqref{eq:frak-e}\\
$\varpi_{Z1}$ & value of $\zeta_{Z1}$ at $+\infty$ & \CGeqref{eq:varpi-Z1}\\
$J_X^{\rm M}$, $J_X$, $C_X$ & Mellin and explicit-formula prime-power counts of $\zeta_X$, $X\in\{\RVM,Z1\}$; their difference & \CGeqref{eq:mellin-count}, \CGeqref{eq:J-X}\\
$\mathsf b_n$ & jumps of $J_{Z1}^{\rm M}$ at the integers & Thm.~\CGref{thm:cell-form}\\
$\varpi_\lambda$; $\Upsilon(\lambda)$ & value of $\zeta_\lambda$ at $+\infty$; flow at infinity & \CGeqref{eq:varpi-lambda}, \CGeqref{eq:Upsilon}\\
$\varphi_n$, $\tilde c_n(\lambda)$, $\tilde a_N(\lambda)$ & smooth cells; atomized masses and coefficients & \CGeqref{eq:atomized-family}, \CGeqref{eq:atomized-coefficients}\\
$\tilde\zeta_{\Z,\lambda}$, $\tilde h_\lambda$ & atomized family; joining factor & \CGeqref{eq:atomized-family}, \CGeqref{eq:joining-factor}\\
$\mathfrak a_s$, $\mathfrak h_n(\lambda;s)$, $\mathfrak L$ & atomization kernel; cell defects; $\int_1^x(1-y^{-1})dy/\log y$ & \CGeqref{eq:atomization-kernel}\\
$V_n$, $W_n$ & ordinate kernels of the cell $n$ & \CGeqref{eq:VW}\\
$\tilde\mu_\lambda$, $\tilde M_\lambda$ & atomized M\"obius coefficients; summatory function & \CGeqref{eq:atomized-mobius}\\
$B_{Z1}$ & Mellin form of the endpoint Mertens function & \CGeqref{eq:Mellin-Mertens}\\
$c_n^{\RVM}$, $\zeta_{\Z,\RVM}$ & integer weights; weighted Euler product & \CGeqref{eq:cn}, \CGeqref{eq:weighted-zeta}\\
$J_{\Z,\RVM}$, $J_{\RVM}^{*,\Z}$, $h_{\RVM}$ & integer-step counts; Mellin correction & \CGeqref{eq:JZ}, \CGeqref{eq:J-star-Z}, \CGeqref{eq:hRVM}\\
$\theta$ & unit-cell step, $0<\theta\le1$ & \CGeqref{eq:unit-cell}\\
$p_j^{\RVM}$ & real RVM prime levels $\pi_{\RVM}^{-1}(j)$ on monotone branches & \CGeqref{eq:intro-real-primes}\\
$\mathcal M(n)$, $\varsigma_n$, $\mu_h$, $M_h$ & target samples of $M_{\RVM}^{\crit}$; signs; ternary walk and its sum (companion object; $h$ is a label) & \CGeqref{eq:intro-muh}, \CGeqref{eq:intro-quantizer}\\
$a_N^{\Z}$ & Dirichlet coefficients of $\zeta_{\Z,\RVM}$ & \CGeqref{eq:ZRVM-Dirichlet}\\
$C_J$ & Lipschitz constant of $J_{\RVM}$ & \CGeqref{eq:unit-cell}\\
$\mathfrak d_n^{\RVM}$ & prime defect $\mathbf 1_{\rm prime}(n)-c_n^{\RVM}$ & \CGeqref{eq:prime-defect}\\
\midrule
\multicolumn{3}{l}{\emph{Mertens transforms and the transition theorem}}\\
$F_{\RVM}=\widehat M_{\RVM}^{\crit}$ & $1/\zeta_{\RVM}$ & \CGeqref{eq:F-RVM}\\
$c$, $U$ & Perron abscissa; truncation height & \CGeqref{eq:Mcrit-symmetric}\\
$M_{\RVM}^{\crit}$, $M_{\RVM,G}^{\crit}$ & symmetric Perron inverse; Gaussian-smoothed inverse & \CGeqref{eq:Mcrit-symmetric}, \CGeqref{eq:MG}\\
$\tau_*$, $w(s)$ & transition point $-iz=t-i\sigma'$; transition variable $e^{-2\pi i\NR(\tau_*)}$ & \CGeqref{eq:transition-zw}\\
$B$; $\sigma'_1,\sigma'_2$ & upper bound for $\sigma'$; a range $[\sigma'_1,\sigma'_2]\subset(0,\tfrac12)$ & \CGeqref{eq:transition-w}, \CGeqref{eq:transition-chi}\\
$\mathcal E_0,\mathcal E$ & error factors of the transition theorem & \CGeqref{eq:transition-w}, \CGeqref{eq:transition-chi}\\
$\varpi_{\RVM}$, $L$ & transition constant $0.9619415807$ (Richardson value at the reported precision); $L=1/\varpi_{\RVM}$ & \CGeqref{eq:varpi-def}, \CGeqref{eq:vertical-limit}\\
$f,\Lambda,\Lambda_{\rm sm},\Lambda_{\rm saw}$ & logarithm of the product and its smooth and sawtooth parts (used in the transition and related clock estimates) & \CGeqref{eq:transition-product}, \CGeqref{eq:transition-split}\\
$\psi,\psi_r,\Phi_r,\nu_1,\dots,\nu_4$ & sawtooth, Abel means, ray remainder, real constants (used in the transition and related clock estimates) & \CGeqref{eq:sawtooth-count}, \CGeqref{eq:psi-r}\\
$\alpha_c$ & $\min(c-\tfrac12,1)$ & \CGeqref{eq:FminusL-all}\\
\midrule
\multicolumn{3}{l}{\emph{Comparison with $\zeta$}}\\
$\mathcal Q_{\RVM}$, $\mathcal Q_{\RVM}^{\vee}$ & quotient $\zeta/\zeta_{\RVM}$; its inverse Mellin transform & \CGeqref{eq:Q-RVM}, \CGeqref{eq:M-convolution}\\
\midrule
\multicolumn{3}{l}{\emph{The $\lambda$-family and the endpoint $\zeta_{Z1}$}}\\
$\lambda$ & family parameter; $\lambda=1$: $\zeta_{\RVM}$, $\lambda=0$: $\zeta_{Z1}$ & \CGeqref{eq:gamma-lambda}\\
$\gamma_k(\lambda)$ & moving ordinates $(1-\lambda)\gamma_k+\lambda\tau_k$ & \CGeqref{eq:gamma-lambda}\\
$P_\lambda$, $\zeta_\lambda$, $F_\lambda$ & family product, family zeta, $1/\zeta_\lambda$ & \CGeqref{eq:P-lambda}, \CGeqref{eq:zeta-lambda}\\
$A_{k,\lambda}$ & $\tfrac14+\gamma_k(\lambda)^2$ & \CGeqref{eq:generator}\\
$\mathcal V_\lambda$ & infinitesimal generator $\partial_\lambda\log P_\lambda$ & \CGeqref{eq:generator}\\
$\Pi_{\rm tr}$ & total multiplier $P_{Z1}/P_{\RVM}^{\crit}$ & \CGeqref{eq:transport-multiplier}\\
$\rho_{k,\lambda}$, $\mathfrak r_{k,\lambda}$ & moving zero $\tfrac12+i\gamma_k(\lambda)$; residue of $1/(s\zeta_\lambda)$ & \CGeqref{eq:residue-transport}\\
$P_{Z1}$, $\zeta_{Z1}$, $F_{Z1}$ & endpoint product, endpoint zeta, $1/\zeta_{Z1}$ & \CGeqref{eq:Z1-product}, \CGeqref{eq:FZ1}\\
$\mathscr R_{Z1}$ & ratio $\zeta_{Z1}/\zeta$ & \CGeqref{eq:Z1-defect-ratio}\\
$\mathcal K$ & multiplier $\zeta_{Z1}/\zeta_{\Z,\RVM}$ & \CGeqref{eq:K-multiplier}\\
(D), $a_n$ & Dirichlet condition for $\zeta_{Z1}$ and its coefficients & \CGeqref{eq:Z1-D}\\
$C_H$ & constant in Hamburger's theorem & \CGeqref{eq:Hamburger-conclusion}\\
$\mu_{Z1}$, $M_{Z1}$ & coefficients of $F_{Z1}$; summatory function & \CGeqref{eq:Z1-reciprocal-Dirichlet}\\
$M_{\lambda,U}$, $M_{Z1,U}$ & finite-height Perron integrals & \CGeqref{eq:truncated-Perron}\\
$\delta$, $\eps$, $\omega(\delta)$ & margins $\sigma\ge\tfrac12+\delta$, exponent $\eps>0$, three-circle exponent & \CGeqref{eq:three-circles-exponent}\\
\midrule
\multicolumn{3}{l}{\emph{Criteria in the half-plane, the rung system and the models}}\\
$\mathsf D_n(\lambda)$, $B^\vee_\lambda$ & $\tilde c_n(\lambda)-\tilde c_n(0)$; dual function of the atomized family & \CGeqref{eq:witness-zeros}, \CGeqref{eq:dual-function}\\
$\mathsf d$, $\mathsf Y$, $\mathsf e_1,\mathsf e_2$ & difference of two counts; auxiliary entire function; exponents & \CGeqref{eq:ladder-h}, \CGeqref{eq:ladder-Y}, \CGeqref{eq:bridge-B2}\\
$J_f$, $J_g$, $b_n$, $\varpi_f$ & counts and coefficients of a general pair $(f,g)$; value of $f$ at $+\infty$ & \CGeqref{eq:halfline-f}, Thm.~\CGref{thm:bridge}\\
\textup{(B1)}, \textup{(B2)} & the two hypotheses of the bridge theorem & Thm.~\CGref{thm:bridge}\\
$\mathsf v_\lambda$ & phase defect of the family on the critical line & \CGeqref{eq:phase-defect}\\
$\mathsf o_n$, $\mathsf L_n$, $\mathfrak j$ & one-ordinate difference of cell masses; $-\log(1-n^{-s})$; one-ordinate joining factor & \CGeqref{eq:one-ordinate}, \CGeqref{eq:one-ordinate-joining}\\
$\mathsf f$, $\mathsf c_{\mathsf f}$, $\mathsf k$, $\mathsf k_{\mathsf f}$ & test function in $u=\log x$; its cosine transform; the part of $G$ and the pole & \CGeqref{eq:transport-defs}\\
$\mathsf R_{\mathsf f}$, $\mathsf S_{\mathsf f}$, $\mathsf m_j(\mathsf f)$ & rung functional; transport; moments sampled on the clock & \CGeqref{eq:rung-functional}, \CGeqref{eq:transport}, \CGeqref{eq:transport-moments}\\
$\mathfrak t_0$, $\mathsf I_m$ & zero of the clock, $\NR(\mathfrak t_0)=0$; Poisson integrals & \CGeqref{eq:clock-defect-poisson}\\
$\mathsf a_n$, $\dot\kappa_k$, $\dot{\mathsf a}_n$ & free masses of the rung system; a direction of variation & \CGeqref{eq:rung-system}, \CGeqref{eq:rung-linearized}\\
$\mathsf J$, $\mathsf x_0$, $\mathsf H_1$, $\mathsf y$, $\mathsf H$, $\mathsf Z$ & objects of the proof of Theorem~\CGref{thm:nondegenerate} & \CGeqref{eq:nondeg-H1}--\CGeqref{eq:nondeg-H}\\
$\langle\cdot,\cdot\rangle_\sigma$, $\mathsf K_{\boldsymbol{\kappa}}$, $\mathsf r_m$, $\mathsf u_m$, $\mathsf Q_{\boldsymbol{\kappa}}$ & pairing of Dirichlet series; $\tilde\zeta_{\Z,\boldsymbol{\kappa}}/\zeta$; gradient data of $\mathsf V_\sigma$ & \CGeqref{eq:K-kappa}, \CGeqref{eq:gradient}\\
$\mathsf z$, $\mathfrak k$, $\mathfrak C_{\boldsymbol{\kappa}}$, $\mathsf N$ & variable, half-type and canonical product of the models; $e^{i\mathfrak k\mathsf z}\mathfrak C_{\boldsymbol{\kappa}}$ & \CGeqref{eq:model-product}\\
$\mathscr F$, $\mathscr F^{+}$, $\mathfrak l$, $\mathfrak w_{\mathfrak l}$, $\mathfrak b_{\mathfrak l}$, $\mathfrak c_{\mathfrak l}$ & frequencies of a discrete category; a frequency; masses, exponents, coefficients & \CGeqref{eq:acceptance-density}--\CGeqref{eq:acceptance-window}\\
$\mathfrak p,\mathfrak q,\mathfrak v$, $\mathsf U$, $\mathsf X$, $\mathfrak D$, $\mathfrak S$ & data, unitary matrix, diagonal factor, determinant and sine function of the two-edge model & \CGeqref{eq:graph-matrices}--\CGeqref{eq:graph-sine}\\
$\mathsf F$, $\mathsf F_{\boldsymbol{\kappa}}$, $\mathsf F_\lambda$, $\mathfrak y_k$, $\mathfrak x_k$ & finite Euler product, its list functions and family; its ordinates and its clock & \CGeqref{eq:finite-euler}--\CGeqref{eq:finite-euler-list}\\
$\mathscr L$, $\mathfrak P$ & $\ell^2$ over the primes; $\sum_pp^{-s}$ & \CGeqref{eq:prime-det}, \CGeqref{eq:prime-det2}\\
$\theta_{\RVM}(s)$, $\tau_0$ & bounded half-plane phase of the clock; fixed transition threshold $\tau_0=4\pi$ (not a clock ordinate) & \CGeqref{eq:halfplane-RVM}, \S\CGref{subsec:transition-proof}\\
$\mathfrak m_n$, $\mathfrak u_\lambda$, $N_\lambda$, $\mathfrak e^{(\lambda)}$, $C_\lambda$ & cell baseline; first-cell term; moving count and count defect; Mellin/explicit-count constant & \CGeqref{eq:cell-mass-W}, \CGeqref{eq:joining-factor}, \CGeqref{eq:family-count-primitive}\\
$\mathsf A_\lambda$, $\mathsf W(\lambda)$, $T_\lambda$, $k_n(\lambda)$, $\mathsf V_\sigma$ & representing function; one-number defect; theta defect; quotient coefficients; vertical defect & \CGeqref{eq:one-number}, \CGeqref{eq:theta-defect}, \CGeqref{eq:vertical-defect}\\
$\Theta_\lambda$, $\mathfrak G_\lambda$, $\mathcal J_T$, $\Delta_T$, $K(T)$, $L_T$ & auxiliary criteria and multiscale clock-Jacobian quantities & \S\CGref{sec:criteria}\\
$E_n$, $\Xi_n$, $\varepsilon_r$ & cell-error and defect-kernel quantities; local regularization parameter where displayed & \S\CGref{sec:criteria}\\
$\boldsymbol{\kappa}$, $N_{\boldsymbol{\kappa}}$, $\tau_0^+$ & generic ordinate list, its count, and lower list threshold & \CGeqref{eq:list-space}\\
$\mathcal G_\varepsilon$, $\mathsf P_{\varepsilon,M}$, $\mathsf R_{\varepsilon,M}$ & additive interpolation and finite completed Dirichlet polynomials & \CGeqref{eq:additive-family-v18}, \CGeqref{eq:finite-additive-v18}\\
$\nu_5,\dots,\nu_{12}$ & real constants, each local to the statement where it occurs & --\\
\midrule
\multicolumn{3}{l}{\emph{Route B for finite sections (Section~\CGref{sec:routeB-finite})}}\\
$\mathbb K(C,\mathbb M)$, $\mathbb M_*$, $\mathsf a^\zeta_n$ & box of $C$-admissible lists with masses $\le\mathbb M$; $\log7/(2\sqrt7)$; prescribed masses $\Lambda(n)/(2\sqrt n)$ & \S\CGref{subsec:rb-routeB}\\
$\mathfrak R(\boldsymbol\kappa,\mathsf a)$, $\mathsf c_\varphi$, $\mathscr T_m$, $\mathsf s_m$ & residual distribution of the rung system; cosine transform of a test; nested test sets; section residual & \CGeqref{eq:rb-residual}, \CGeqref{eq:rb-rho}\\
$\mathcal L$, $\mathscr K$, $\mathscr S_{\boldsymbol\kappa}$, $\mathscr A_\zeta$, $\mathscr E_{\boldsymbol\kappa}$, $\mathfrak H_{\boldsymbol\kappa}$, $\mathfrak g$, $\mathrm c_{\mathscr K}$ & Laplace transform; Archimedean, spectral, arithmetic and total residual distributions; $\sum_k\kappa_k^{-2}\cos(\kappa_ku)$; auxiliary function and constant & \CGeqref{eq:rb-laplace-def}--\CGeqref{eq:rb-S-A-dist}\\
$\mathbb F_m$, $\mathbb J_m$, $\mathbb Q_m$, $\mathbb L_m$, $\mathbb H_m$, $\mathsf h_m$, $\mathsf l_m$, $k^*_m$, $\mathrm t_m$, $\mathbb T_m$ & residual map, Jacobian, right inverse and its norm, Lipschitz constant, start residual and contraction index of stage $m$; numbers of unknowns and tests; contraction map & \CGeqref{eq:rb-Fm}--\CGeqref{eq:rb-Am}\\
$\mathrm q_m$ & graded contraction factors & Prop.~\CGref{prop:rb-graded}\\
$P_{\mathrm X}$, $\zeta_{\mathrm X}$, $\Lambda_{\mathrm X}$, $\mathrm p_{\mathrm X}$, $\mathrm m_{\mathrm X}$, $\mathrm n_{\mathrm X}$, $\varTheta_{\mathrm X}$, $\boldsymbol\kappa^{(\mathrm X)}$ & Gonek's truncated Euler product and symmetrised function with parameter $\mathrm X$; its prime phase and two constants; the Gonek phase; the Gonek list & \CGeqref{eq:rb-alphaX}--\CGeqref{eq:rb-gonek-list}\\
$\widehat{\mathsf R}$, $\mathbb E(T)$, $\upsilon$, $\mathfrak B(C)$, $\mathbb B,\mathbb B_0,\dots,\mathbb B_3$, $\mathrm c_1,\dots,\mathrm c_6$, $\mathrm c_\ell$ & Laplace residual of a section; its $L^1$ size on $\Re s=2$ and growth ratio; constants of the a priori bound & \CGeqref{eq:rb-Rhat}--\CGeqref{eq:rb-a-constants}\\
$\mathcal Y$, $\mathcal D_{\mathsf a}$, $\mathbb W$, $\mathbb S$, $\mathbb J$, $\mathfrak f$, $\mathfrak F_T$, $\mathfrak R_T$, $\mathbb G_T$, $n_T$, $\mathbb D(T)$ & Dirichlet factor and polynomial; $\log|\zeta_{\boldsymbol\kappa}|$; half-strip; compact set; log-factor; Backlund function; rectangle, path, zero count and bound of the counting argument & \S\CGref{subsec:rb-apriori}\\
$\mathscr X$, $\mathrm x_k$, $\varPsi_{\mathscr X}$, $\mathbb P_{\mathrm y}$, $\mathbb Y_{\mathrm y}$, $\mathfrak G_{\mathscr X}$, $\mathcal F_{\mathscr X}$, $\mathcal F_0$, $\mathfrak z=\mathrm x+i\mathrm y$ & unfolded list and its points; discrepancy; Poisson and conjugate operators; partial-fraction series; generating function; $\cos\pi\mathfrak z$; unfolded complex variable & \CGeqref{eq:rb-H}--\CGeqref{eq:rb-F-formula}\\
$\mathrm A$, $\mathrm B$, $\mathrm d_\star$, $\mathbb U$, $\mathcal W$, $\mathrm w$, $\mathrm d_k$, $\mathfrak Z_m$, $\mathfrak J_m$, $\mathfrak n_m$, $\mathrm D_m$, $\mathbb O$, $\mathbb V$, $\mathbb U_{\rm loc}$ & mean and oscillating parts of the discrepancy and their bounds; sector function; weight; displacements, index blocks, intervals, lengths, oscillations and the constants of the variable-length theorem & \S\CGref{subsec:rb-riesz}\\
$\hbar$, $\hbar_T$, $\mathbb P^{\NR}_t$, $\mathbb P^{\circ}_t$, $\mathrm c_{\rm JN}$, $\mathfrak I_T$, $\mathrm c_F$ & line height and its $t$-scale; the two kernels, the John--Nirenberg constant, the interval and Fujii's constant in the obstruction theorem & \S\CGref{subsec:rb-obstruction}\\
$\mathrm w_m$, $T_m$, $u^*_m$, $\mathbb J'$, $\mathbb A^{(3)}$, $\mathsf l(\mathrm r)$ & test width, height and window of stage $m$; second- and third-derivative matrices; Banach index of the computed certificates & \S\CGref{subsec:rb-numerics}\\
\end{longtable}
}

\medskip\noindent\textbf{The three counts.}  Three counting symbols occur and are never interchanged: $\NR(T)=\mathcal N_{\RVM}(T)$ is the \emph{continuous} midpoint clock \CGeqref{eq:clock}, a real-analytic increasing function; $n_{\RVM}(T)=\#\{k:\tau_k\le T\}=\lfloor\NR(T)+\tfrac12\rfloor$ \CGeqref{eq:n-RVM} is the \emph{integer} count of the clock ordinates; and $N_{\boldsymbol\kappa}(T)=\#\{k:\kappa_k\le T\}$ is the integer count of a general list (with $N_\zeta$ for the ordinates of $\zeta$ and $N_\lambda$ for the list $(\gamma_k(\lambda))$).  Differences such as $N_\zeta-\NR$ compare an integer staircase with the continuous clock; differences such as $N_{\boldsymbol\kappa}-n_{\RVM}$ compare two staircases.  When a mass $\Lambda(n)/(2\sqrt n)$ is called a ``von Mangoldt mass'' it is this scaled mass that is meant.  A distributional identity is stated with its space: $\mathcal D'(0,\infty)$ (the origin excluded), $\mathcal S'(\mathbb R)$, or distributions on $\mathbb R$ supported in $[0,\infty)$ (the origin included); the constants $\mathrm b$ of Appendix~\CGref{subsec:rung-apriori-v18} and the polynomials of Lemma~\CGref{lem:rb-laplace} live exactly in the difference between the first and the third.

\medskip\noindent\textbf{Collision check.}
Global symbols in the table have one meaning throughout the paper.  In version~22 the interpolation parameter of Appendix~\CGref{subsec:rung-apriori-v18} is written $\mathsf t\in[0,1]$ (it was $\theta$, which is the unit-cell step), and the constants of that appendix are $\mathbb K_{\rm E},\mathbb K_{\rm E}'$ (not $\mathbb K(C,\mathbb M)$, the box of Section~\CGref{sec:routeB-finite}), $B_2$ (not $\mathbb B_2$ or $\mathrm B_2$ of that section) and $C^*$; the letters $K$, $N$, $T$, $L$ and $C$ with different fonts or indices are different symbols, listed separately in the table.  Local proof variables that formerly collided with global notation have been renamed in version~18: in particular $c_0$ is the de la Vall\'ee Poussin constant, $\varkappa_j$ and $\mathcal P_j$ are local van-der-Corput parameters, $\omega_0$ and $\omega_{\rm test}$ are test weights, $d_\rho$ is a horizontal zero displacement, $T_{\rm ver}$ is the verified-zero height, $\mathcal H^\pm$ are transition half-planes, and $\Psi_m^\pm$ are local stationary phases; $R_0$ is reserved for the clock half-plane threshold and $R_{\rm end}$ for the endpoint clock-decomposition threshold.  The symbol $c$ is reserved for a Perron abscissa, $\eps$ for an exponent/margin, $\varepsilon$ for the additive-family parameter, $\theta$ for the unit-cell step, and $\vartheta$ for the Riemann--Siegel function.  The fixed number $\tau_0=4\pi$ is a transition threshold and is not a member of the sequence $(\tau_k)$.  Dotted variables $\dot\kappa_k,\dot{\mathsf a}_n$ denote a direction of variation.

\section{Introduction}
\label{sec:introduction}

\subsection{Background}
\label{subsec:background}

The theory of the Riemann zeta function joins four structures that are usually introduced in different languages.  On the spectral side there are the nontrivial zeros and their Riemann--von Mangoldt counting law; on the analytic side the completed function $\xi$, its functional equation and its Hadamard product; on the prime side the prime-power measure and the explicit formula; and on the inverse side the reciprocal Dirichlet series $1/\zeta$ and the Mertens function
\begin{equation}\label{eq:mertens-def}
 M(x):=\sum_{n\le x}\mu(n).
\end{equation}
Riemann's memoir made the first three structures inseparable, and the work of von Mangoldt, Hadamard and Hamburger supplied, respectively, the rigorous zero count and explicit formula, the entire-product factorization, and a converse characterization of $\zeta$ from its Dirichlet-series and functional-equation data \cite{Riemann1859,vonMangoldt1895,Hadamard1896,Hamburger1921,Titchmarsh1986,Edwards1974}.  This paper asks what remains of that architecture when the zero staircase is replaced, at the outset, by the exact midpoint inverse of its smooth Riemann--von Mangoldt clock, and how the resulting object is carried back to $\zeta$.

\subsection{The midpoint RVM system}
\label{subsec:midpoint-system}

The point of departure is the smooth counting function \CGeqref{eq:clock},
\begin{equation}\label{eq:intro-clock}
 \NR(T)=\frac{T}{2\pi}\log\frac{T}{2\pi e}+\frac78,
\end{equation}
and the real ordinates $\tau_k$ determined by the midpoint condition $\NR(\tau_k)=k-\tfrac12$.  The half-integer convention is the midpoint convention for a staircase and leads to the exact Lambert inversion \CGeqref{eq:lambert},
\begin{equation}\label{eq:intro-lambert}
 \tau_k=\frac{2\pi(k-11/8)}{W_0\bigl((k-11/8)/e\bigr)} .
\end{equation}
Lambert-$W$ approximations to individual zero ordinates and smooth inversions of the Riemann--von Mangoldt law occur in the literature, including explicit equations for the $n$th zero \cite{FrancaLeClair2014,Corless1996,DLMF2024}.  Here the midpoint inverse is treated as the exact spectrum of a deterministic system and is propagated, without the fluctuating term $S(T)$, through canonical products, explicit prime counting and Mertens inversion.

The first object is the canonical product and the synthetic zeta function
\begin{equation}\label{eq:intro-zeta-rvm}
 P_{\RVM}^{\crit}(s)=\prod_{k\ge1}\left(1+\frac{s(s-1)}{\tfrac14+\tau_k^2}\right),
 \qquad
 \zeta_{\RVM}(s)=\frac{P_{\RVM}^{\crit}(s)}{G(s)},
 \qquad G(s)=s(s-1)\Gamma(s/2)\pi^{-s/2}.
\end{equation}
Because $\tau_k\asymp k/\log k$, the factors form a normally convergent product in the symmetric variable $s(s-1)$; $P_{\RVM}^{\crit}$ is entire and invariant under $s\mapsto1-s$, and $\zeta_{\RVM}$ has a residue-one pole at $1$, the trivial zeros, and the simple nontrivial zeros $\rho_k^{\RVM}=\tfrac12\pm i\tau_k$ (Theorem~\CGref{thm:canonical}).  The real part $\tfrac12$ is part of the definition.

\emph{Generalized prime systems.}  In Beurling's theory the generalized primes are arbitrary reals $1<p_1\le p_2\le\cdots\to\infty$; their multiplicative semigroup defines generalized integers and a zeta function by a Dirichlet series, an Euler product or a Mellin--Stieltjes transform, separating prime distribution from the additive structure of the integers \cite{Beurling1937,BatemanDiamond1969,Diamond1970}.  Beurling systems with prescribed analytic behaviour, critical zeros and large oscillation show that zeta-like zero data alone do not determine the arithmetic realization \cite{DiamondMontgomeryVorhauer2006,BrouckeDebruyneVindas2020}.  The RVM system runs in the opposite direction: it starts from a prescribed zero clock, constructs an entire completed product first, obtains a continuous prime-power count from its explicit formula, and only afterwards asks for an atomic realization on the integers.

\emph{Continuous and discrete prime systems.}  The passage from a continuous prime system to a discrete one is classical in the theory of Beurling generalized primes \cite{Beurling1937,DiamondMontgomeryVorhauer2006,BrouckeVindas2024,DiamondZhang2016}: the atoms are placed at suitably chosen real points, and the discretization changes the zeta function by a factor holomorphic and zero-free in $\Re s>\frac12$ (compare Theorem~\CGref{thm:atomized-continuation}(iii)).  Measures with a continuous part and non-integer multiplicities also appear in the theory of complex dimensions \cite[Ch.~3]{LapidusvanFrankenhuysen2000}.  Grosswald and Schnitzer \cite{GrosswaldSchnitzer1978} replaced each prime $p_n$ by a real $q_n$ with $p_n\le q_n\le p_{n+1}$ and proved that the modified Euler product is $\zeta$ times a factor holomorphic and zero-free in $\Re s>0$, so that it has the same zeros as $\zeta$ there; Corollary~\CGref{cor:weighted-factorization} and Theorem~\CGref{thm:atomized-continuation}(iv) are of the same type for $\zeta_{\RVM}$.  Here the masses of $d\pi_\lambda$ are placed at the integers, with real exponents $\tilde c_n(\lambda)$; this yields an ordinary Dirichlet series for every $\lambda$, absolutely convergent in $\Re s>1$ (Corollary~\CGref{cor:abscissa-one}), and makes Hamburger's theorem \cite{Hamburger1921,Siegel1922} available at the endpoint.  Hamburger's theorem and its extensions \cite{BochnerChandrasekharan1956,KaczorowskiPerelli1999} need absolute convergence in $\Re s>1$, and Knopp \cite{Knopp1994} showed that a half-plane of absolute convergence further to the right does not suffice.  We are not aware of an earlier use of this combination.

\emph{Prime coordinate.}  From the RVM zeros we define the prime-power count $J_{\RVM}$ by \CGeqref{eq:J-RVM}, with the branch \CGeqref{eq:Ei-def} of $\Ei$.  Finite M\"obius inversion in the exponent variable, which relates Riemann's prime-power count to $\pi(x)$ \cite{Riemann1859,Titchmarsh1986,Edwards1974}, is used here as the \emph{definition} of the synthetic prime count \CGeqref{eq:pi-RVM},
\begin{equation}\label{eq:intro-pi}
 \pi_{\RVM}(x)=\sum_{1\le m\le\log_2x}\frac{\mu(m)}m\,J_{\RVM}^{*}(x^{1/m}),
 \qquad
 J_{\RVM}^{*}(x)=\sum_{1\le m\le\log_2x}\frac1m\,\pi_{\RVM}(x^{1/m}).
\end{equation}
The function $\pi_{\RVM}$ is continuous away from the points $2^j$ with $\mu(j)\ne0$, where it has the jump $\mu(j)J_{\RVM}(2)/j$; in particular it is not a staircase on the integers.  It supports two notions of RVM prime: the continuous measure $d\pi_{\RVM}$, and, on every interval where $\pi_{\RVM}$ is continuous and strictly increasing, the real level crossings
\begin{equation}\label{eq:intro-real-primes}
 p_j^{\RVM}:=\pi_{\RVM}^{-1}(j),\qquad \pi_{\RVM}(p_j^{\RVM})=j,
\end{equation}
which need not be integers.  The synthetic count has square-root accuracy without any hypothesis on $\zeta$: $J_{\RVM}(x)-\li(x)$ and $\pi_{\RVM}(x)-\li(x)$ are $O(\sqrt x(\log\log x)^2/\log x)$ (Theorem~\CGref{thm:Lambert-estimate}).  The proof reduces the zero sum to $\sum_ke^{iu\tau_k}/\tau_k$, which is $O(\log^2u)$ by the second-derivative test on dyadic blocks of the clock (Lemma~\CGref{lem:A-RVM}).  Hence $\li(x)/\pi_{\RVM}(x)\to1$, and with the prime number theorem $\pi(x)/\li(x)\to1$ and $\pi(x)/\pi_{\RVM}(x)\to1$ (Corollary~\CGref{cor:three-limits}, Figure~\CGref{fig:pi-ratios}).  Sampling on unit cells gives the masses $c_n^{\RVM}=\pi_{\RVM}(n)-\pi_{\RVM}(n-1)$ and the weighted Euler object $\zeta_{\Z,\RVM}$ of \CGeqref{eq:weighted-zeta}.  The masses are not $0$ or $1$ and need not be positive, but they are bounded: $J_{\RVM}$ is Lipschitz on $[2,\infty)$ (Theorem~\CGref{thm:unit-cell}), so $c_n^{\RVM}=O(1)$ and $\zeta_{\Z,\RVM}$ is an absolutely convergent Dirichlet series in $\Re s>1$ (Corollaries~\CGref{cor:cn-bounded} and~\CGref{cor:ZRVM-Dirichlet}).  The unit-periodic first harmonic of the clock integrates to zero on every integer cell (Lemma~\CGref{lem:cell-cancel}); under the cell hypothesis one obtains $c_n^{\RVM}\sim1/\log n$ (Proposition~\CGref{prop:eventual-positive}).

\emph{Mertens coordinate.}  The Mertens coordinate is the symmetric inverse Mellin object of $F_{\RVM}=1/\zeta_{\RVM}$.  A smooth/sawtooth Stieltjes decomposition of the clock product with Abel regularization of the sawtooth gives on every band to the right of $\tfrac12$
\begin{equation}\label{eq:intro-transition}
 \zeta_{\RVM}(s)=\varpi_{\RVM}\bigl(1+w(s)\bigr)\bigl(1+O(1/|t|)\bigr),
 \qquad w(s)=\exp\bigl\{-2\pi i\,\NR\bigl(-i(s-\tfrac12)\bigr)\bigr\},
\end{equation}
and $\zeta_{\RVM}=\varpi_{\RVM}(1+\chi)(1+O(1/|t|))$ on closed substrips of $\tfrac12<\Re s<1$ (Theorem~\CGref{thm:transition}).  Hence $F_{\RVM}(\sigma+it)\to L=1/\varpi_{\RVM}$ for every fixed $\sigma>\tfrac12$; neither $\zeta_{\RVM}$ nor $F_{\RVM}$ is an absolutely convergent Dirichlet series in any half-plane (Theorem~\CGref{thm:no-Dirichlet-RVM}); and after subtracting the vertical limit the Perron inverse converges absolutely, with
\begin{equation}\label{eq:intro-Mbound}
 M_{\RVM}^{\crit}(x)=L+O_\eps(x^{1/2+\eps})
\end{equation}
(Theorem~\CGref{thm:unsmoothed-RVM}).  The RVM Mertens function therefore exists as a continuous spectral inverse although no coefficient sequence for $F_{\RVM}$ exists.  Since the primitive of the midpoint sawtooth is bounded, the same decomposition without contours gives $C^{-1}\le|\zeta_{\RVM}|\le C$ on every half-plane $\Re s\ge\sigma_0>1$ and $\zeta_{\RVM}(s)\to\varpi_{\RVM}$ as $\Re s\to+\infty$ (Theorem~\CGref{thm:halfplane-RVM}).

\emph{Ternary discretization.}  An integer-valued representative of $M_{\RVM}^{\crit}$ is obtained by sampling the target $\mathcal M(n):=M_{\RVM}^{\crit}(n)$ and choosing signs on the squarefree support:
\begin{equation}\label{eq:intro-muh}
 \mu_h(n)=\mu(n)^2\varsigma_n,\qquad\varsigma_n\in\{-1,+1\},\qquad M_h(n):=\sum_{m\le n}\mu_h(m),
\end{equation}
\begin{equation}\label{eq:intro-quantizer}
 \varsigma_n=\begin{cases}+1,&M_h(n-1)<\mathcal M(n),\\-1,&M_h(n-1)\ge\mathcal M(n),\end{cases}\qquad(\mu(n)^2=1).
\end{equation}
This is a greedy sigma--delta feedback quantization; feedback quantization and noise shaping are standard themes of sigma--delta theory \cite{DaubechiesDeVore2003,GunturkThao2003,BenedettoPowellYilmaz2006}.  Here the target is a zero-generated Mertens inverse and the admissible times are the squarefree integers.  The walk $\mu_h$ is a comparison object: it is not claimed to be multiplicative or equal to $\mu$.  The two discretizations are distinct: the cell masses $c_n^{\RVM}$ come from $\pi_{\RVM}$, the signs $\varsigma_n$ track $M_{\RVM}^{\crit}$.

\emph{Comparison with $\zeta$.}  At the level of transforms the defect is the exact quotient $\mathcal Q_{\RVM}=\zeta/\zeta_{\RVM}$ of \CGeqref{eq:Q-RVM}, with $F_{\RVM}=\mathcal Q_{\RVM}\cdot(1/\zeta)$.  Since $|\zeta_{\RVM}|$ is bounded above and below on $\frac12+\eps\le\sigma\le c'$, $|t|\ge t_0$ (Corollary~\CGref{cor:Q-control}), on each bounded band $\frac12+\eps\le\sigma\le c'$, every bound for $\zeta$ or $1/\zeta$ is equivalent to the corresponding bound for $\mathcal Q_{\RVM}$ by Corollary~\CGref{cor:Q-control}; for $\sigma>c'$ one uses Theorem~\CGref{thm:halfplane-RVM}.  For $M$ itself the bound $M(x)=O_\eps(x^{1/2+\eps})$ is equivalent to the Riemann hypothesis \cite[\S14.25]{Titchmarsh1986}, and the disproof of the Mertens conjecture shows that the $\eps$ cannot simply be dropped \cite{OdlyzkoTeRiele1985}.

\subsection{The core of the method}
\label{subsec:core}

The transfer from $\zeta_{\RVM}$ to $\zeta$ is a transformation of $\zeta_{\RVM}$ depending on a parameter $\lambda\in[0,1]$.  The ordinates move along
\begin{equation}\label{eq:intro-family}
 \gamma_k(\lambda)=(1-\lambda)\gamma_k+\lambda\tau_k,
 \qquad
 \zeta_\lambda(s)=\frac1{G(s)}\prod_{k\ge1}\left(1+\frac{s(s-1)}{\tfrac14+\gamma_k(\lambda)^2}\right),
\end{equation}
where $\gamma_k$ are the ordinates of the nontrivial zeros of $\zeta$; thus $\zeta_1=\zeta_{\RVM}$ and $\zeta_0=\zeta_{Z1}$ (Section~\CGref{sec:lambda-family}).
\begin{enumerate}[label=\textup{(\roman*)}]
\item Every member $\zeta_\lambda$ has its nontrivial zeros exactly at $\tfrac12\pm i\gamma_k(\lambda)$.  The real part $\tfrac12$ is part of the definition of the family, like the gamma factor, and never changes along it: it is impossible for $\zeta_{Z1}$, or for any $\zeta_\lambda$, to have a zero off the critical line other than the trivial zeros.  For $0<\lambda\le1$ all zeros are simple.
\item All differences between the moving zetas lie in the numbers $\gamma_k-\tau_k$, which are bounded (Lemma~\CGref{lem:bounded-displacement}); the quotient of two members is an explicit product in $\gamma_k$ and $\lambda(\tau_k-\gamma_k)$ (Remark~\CGref{rem:only-gamma-tau}).
\item The final step is the identity $\zeta_{Z1}=\zeta$.  It is the single and complete test: while $\zeta_{Z1}=\zeta$ there is nothing outside $\Re s=\tfrac12$ to speak of, since the identity itself places every nontrivial zero of $\zeta$ at $\tfrac12\pm i\gamma_k$, and $\zeta_{Z1}\ne\zeta$ is the only possible obstruction.
\end{enumerate}
The endpoint has the functional equation, the pole, the residue and the order of $\zeta$ (Proposition~\CGref{prop:Z1-class}).  Hamburger's theorem belongs to the family of converse theorems that recover Dirichlet series from analytic continuation, functional equation and growth \cite{Hamburger1921,Perelli2016}; applied to $\zeta_{Z1}$ it gives (Theorem~\CGref{thm:Hamburger-Z1})
\begin{equation}\label{eq:intro-hamburger}
 \zeta_{Z1}=\zeta
 \quad\Longleftrightarrow\quad
 \text{(D): }\zeta_{Z1}\text{ is an absolutely convergent Dirichlet series in }\Re s>1 .
\end{equation}
With the integer-weighted zeta $\zeta_{\Z,\RVM}$, which is a Dirichlet series, (D) is equivalent to an absolutely convergent Dirichlet expansion of the explicit multiplier $\mathcal K=\zeta_{Z1}/\zeta_{\Z,\RVM}=\Pi_{\rm tr}e^{-h_{\RVM}}$ (Theorem~\CGref{thm:D-multiplier}); since $1/\zeta_{\Z,\RVM}$ is also an absolutely convergent Dirichlet series in $\Re s>1$, the equivalence holds on the whole half-plane.  Section~\CGref{sec:euler-products} gives two sharper forms of (D).  First, the clock decomposition applied to the ordinates $\gamma_k$ gives the value of the endpoint at infinity,
\begin{equation}\label{eq:intro-varpi-Z1}
 \varpi_{Z1}=\lim_{\sigma\to+\infty}\zeta_{Z1}(\sigma)=\varpi_{\RVM}\lim_{K\to\infty}\prod_{k\le K}\frac{\frac14+\tau_k^2}{\frac14+\gamma_k^2}\ \ge1,
\end{equation}
and $\zeta_{Z1}=\zeta$ holds if and only if $\varpi_{Z1}=1$ (Theorems~\CGref{thm:Z1-clock} and~\CGref{thm:varpi-Z1}, Corollary~\CGref{cor:product-identity}): the final step is the single ordered-limit identity $\lim_{K\to\infty}\prod_{k\le K}(\frac14+\gamma_k^2)/(\frac14+\tau_k^2)=\varpi_{\RVM}$ between the displacements and the transition constant.  Second, the prime-power counts $J_{\RVM}^{\rm M}$ and $J_{Z1}^{\rm M}$ are the Mellin transforms of $\log\zeta_{\RVM}$ and $\log\zeta_{Z1}$ and satisfy a distributional explicit formula (Theorems~\CGref{thm:mellin-counts} and~\CGref{thm:distributional-EF}); this proves the Mellin--Stieltjes identification of the integer weights (Theorem~\CGref{thm:MS}), and (D) holds if and only if $J_{Z1}^{\rm M}$ is constant between consecutive integers and its jumps satisfy the growth hypothesis of Theorem~\CGref{thm:cell-form}.  Third, along the family $\frac d{d\lambda}\log\varpi_\lambda=-\Upsilon(\lambda)$, the generator at infinity, and (D) is the identity $\int_0^1\Upsilon(\lambda)\,d\lambda=\log L$ (Proposition~\CGref{prop:family-flow}).  Section~\CGref{sec:atomized} atomizes every member: $\tilde\zeta_{\Z,\lambda}=\zeta_\lambda e^{\tilde h_\lambda}$ keeps the Euler product and the Dirichlet series while $\zeta_\lambda$ keeps the functional equation; the joining factor $\tilde h_\lambda$ is written cell by cell and is holomorphic in $\Re s>\frac12$ (Theorems~\CGref{thm:joining-factor} and~\CGref{thm:atomized-continuation}); the cell masses have mean square $X^{1+\eps}$ on $(X,2X]$ uniformly in $\lambda$, so every $\tilde\zeta_{\Z,\lambda}$, the endpoint included, is an absolutely convergent Dirichlet series in $\Re s>1$ (Theorem~\CGref{thm:mean-square-cells}, Corollary~\CGref{cor:abscissa-one}); and (D) is equivalent to the vanishing of every cell defect at $\lambda=0$ (Theorem~\CGref{thm:closing}).  At the clock endpoint the joining factor $\tilde h_1$ is not constant, and the endpoint Mertens function has a Mellin form $B_{Z1}$ with $B_{Z1}=M$ exactly when (D) holds (Propositions~\CGref{prop:atomized-clock} and~\CGref{prop:Mellin-Mertens}).  The remaining steps form the programme of Subsection~\CGref{subsec:programme}.  In RVM language, the final step asks whether the zero-generated analytic object $\zeta_{Z1}$ belongs to the discrete Dirichlet category; if it does, Hamburger uniqueness makes it $\zeta$, and the prime and Mertens coordinates become those of the integers.  Step~3 is Theorem~\CGref{thm:step3}:
\begin{equation}\label{eq:intro-step3}
 \zeta_{Z1}=\zeta
 \quad\Longrightarrow\quad
 M(x)=O_\eps\bigl(x^{1/2+\eps}\bigr)\qquad(\eps>0),
\end{equation}
with a self-contained proof by the Borel--Carath\'eodory lemma, Hadamard's three-circle theorem and the truncated Perron formula.  Theorem~\CGref{thm:transfer} rewrites the same endpoint consequence through the synthetic factorization.  Its estimate for $M$ is a corollary of Theorem~\CGref{thm:step3}; the synthetic Mertens bound is not an independent input to the classical implication.  The identity $\zeta_{Z1}=\zeta$ itself, equivalently (D), is not proved in this paper; it is the open step of the method (Remark~\CGref{rem:Step3-status}).

The relations among the principal objects are
\begin{equation}\label{eq:intro-prime-chain}
 \{\tau_k\}\longrightarrow\zeta_{\RVM}\longrightarrow J_{\RVM}\longleftrightarrow\pi_{\RVM}\longrightarrow\bigl(d\pi_{\RVM},\,p_j^{\RVM},\,c_n^{\RVM}\bigr),
\end{equation}
\begin{equation}\label{eq:intro-M-chain}
 F_{\RVM}=\zeta_{\RVM}^{-1}\longrightarrow M_{\RVM}^{\crit}\longrightarrow(\mathcal M,\mu_h,M_h),
\end{equation}
\begin{equation}\label{eq:intro-family-chain}
 \zeta_{\RVM}=\zeta_1\xrightarrow{\ \lambda\downarrow0\ }\zeta_0=\zeta_{Z1}\overset{?}{=}\zeta\ \Longrightarrow\ M(x)=O_\eps(x^{1/2+\eps}).
\end{equation}
The first chain is the prime coordinate, the second the Mertens coordinate, the third the $\lambda$-family; they share the midpoint spectrum and solve different realization problems.

\subsection{Originality, scope and status}
\label{subsec:status}

\emph{What this paper does not prove.}  Before the list of results: this paper does not prove the Riemann hypothesis; it does not prove (D), the identity $\zeta_{Z1}=\zeta$, or any of the statements shown below to be equivalent to it; it does not prove the existence of a solution of the rung system \CGeqref{eq:rung-system} in $\mathscr O$, nor of an interior exact-atomic path, nor the convergence of the finite sections of Section~\CGref{sec:routeB-finite} to a distributional limit; it does not prove the inequality $\varpi_{\RVM}<1$ (Hypothesis~\CGref{hyp:varpi-less-one}); and it contains no computer-assisted proof, the computations being floating-point experiments.  The only points left open by the method are these: the endpoint identity (D) in any of its equivalent forms; the existence of a solution of the infinite rung system; the injectivity of the clock-only linearization (Theorem~\CGref{thm:nondegenerate} proves injectivity at a solution, not at the clock); and the passage to the limit from finite sections.  Every ``equivalent form'' below is a reduction to (D), not a resolution of it.

The ingredients are standard or belong to neighbouring theories: the Riemann--von Mangoldt law, Lambert inversion, Hadamard products, explicit formulas, M\"obius and Perron inversion, Beurling generalized primes, sigma--delta feedback, and Hamburger's converse theorem.  The contribution is their assembly into one midpoint RVM space with four coordinated levels: (i) an exact midpoint zero clock and its canonical zeta-like product; (ii) a continuous prime-power and prime coordinate, with real prime levels and non-integral cell masses; (iii) a continuous synthetic Mertens inverse with its ternary feedback discretization; (iv) the $\lambda$-family from the midpoint spectrum to the ordinates of $\zeta$, ending in the single identity $\zeta_{Z1}=\zeta$ and the Mertens bound that follows from it.  To the author's knowledge the cited literature contains the individual mechanisms but not this chain.  Standard material is cited at the point of use; statements without citation are original.

\begingroup\small
\begin{longtable}{p{0.50\textwidth}p{0.40\textwidth}}
\toprule
Statement & Status\\
\midrule
\endfirsthead
\toprule
Statement & Status\\
\midrule
\endhead
\bottomrule
\endfoot
Midpoint clock, Lambert inversion, bounded displacement & theorem\\
Canonical product $P_{\RVM}^{\crit}$ and $\zeta_{\RVM}$ & theorem\\
$S_{\RVM}(N;u)\ll1+\log N$; centroid and $M_H^{\RVM}$ decay & theorem\\
Compact-uniform existence of $J_{\RVM}$; finite inversion & theorem\\
Real prime levels $p_j^{\RVM}$ & defined on monotone branches of $\pi_{\RVM}$\\
Square-root estimate for $J_{\RVM},\pi_{\RVM}$ & theorem\\
$\li/\pi_{\RVM}\to1$; $\pi/\pi_{\RVM}\to1$ & theorem (with the prime number theorem)\\
$J_{\RVM}$ Lipschitz, $c_n^{\RVM}=O(1)$ & theorem\\
$\zeta_{\Z,\RVM}$ absolutely convergent Dirichlet series in $\Re s>1$ & theorem\\
Mellin--Stieltjes identification; $\zeta_{\Z,\RVM}=\zeta_{\RVM}e^{h_{\RVM}}$ & theorem\\
$|\zeta|\asymp|\mathcal Q_{\RVM}|$ for $\sigma\ge\frac12+\eps$, $|t|\ge t_0$ & theorem\\
Transition theorem $\varpi_{\RVM}(1+w)$, $\varpi_{\RVM}(1+\chi)$ & theorem\\
$C^{-1}\le|\zeta_{\RVM}|\le C$ on $\Re s\ge\sigma_0>1$; $\zeta_{\RVM}\to\varpi_{\RVM}$ at $+\infty$ & theorem\\
Unsmoothed $M_{\RVM}^{\crit}=L+O_\eps(x^{1/2+\eps})$ & theorem\\
Dirichlet expansion of $\zeta_{\RVM}^{\pm1}$ & impossible (theorem)\\
Ternary walk $\mu_h$, $M_h$ & definition; comparison object\\
Exact Hadamard product of $2\xi$ & theorem\\
$\lambda$-family, generator, multiplier, simplicity for $\lambda>0$ & theorem\\
Analytic class of $\zeta_{Z1}$ & theorem\\
$\zeta_{Z1}=\zeta\Leftrightarrow$ (D) $\Leftrightarrow$ RH & theorem\\
(D) $\Leftrightarrow$ Dirichlet expansion of $\mathcal K=\zeta_{Z1}/\zeta_{\Z,\RVM}$ in $\Re s>1$ & theorem\\
$\varpi_{Z1}=\varpi_{\RVM}\lim_{K\to\infty}\prod_{k\le K}\frac{1/4+\tau_k^2}{1/4+\gamma_k^2}$ (ordered product); (D) $\Leftrightarrow\varpi_{Z1}=1$ & theorem\\
Mellin prime-power counts; distributional explicit formula & theorem\\
(D) $\Leftrightarrow$ $J_{Z1}^{\rm M}$ constant between integers with the jump-growth bound of Theorem~\CGref{thm:cell-form} & theorem\\
$\frac d{d\lambda}\log\varpi_\lambda=-\Upsilon$; (D) $\Leftrightarrow\int_0^1\Upsilon=\log L$ & theorem\\
Atomized family; joining factor holomorphic in $\Re s>\frac12$ & theorem\\
Atomized family: Dirichlet series in $\Re s>1$ for every $\lambda$ & theorem\\
(D) $\Leftrightarrow$ vanishing cell defects at $\lambda=0$ & theorem\\
$\zeta_\lambda$ not a Dirichlet series in $\Re s>1$ for $0<\lambda\le1$ & theorem\\
$\lambda\mapsto\log\zeta_\lambda(\sigma)$ strictly convex for real $\sigma>1$ & theorem\\
Hamburger's theorem for the pair $(\zeta_\lambda,\tilde\zeta_{\Z,\lambda})$ with symmetric joining factor & theorem\\
(D) $\Leftrightarrow$ $\zeta_{Z1}$ has a positive representing measure; then $\varpi_{Z1}\le1$ & theorem\\
Density of $d\pi_{\RVM}$; cell masses and defects at $\lambda=1$ & theorem\\
Height kernel of the cell masses; Gaussian defect of the atomization & theorem\\
$\zeta_{Z1}(\sigma)\ge\zeta(\sigma)$ for real $\sigma>1$; (D) $\Leftrightarrow$ equality at one $\sigma$ & theorem\\
Hamburger's theorem for the pair in theta form and in vertical mean-square form & theorem\\
Symmetry of the joining factor fails at $\lambda=1$; $\tilde\zeta_{\Z,\lambda}=\zeta$ for finitely many $\lambda$ at most & theorem\\
Summation formula for every $\zeta_\lambda$; $\varpi_\lambda-1$ as one number $\mathsf W(\lambda)$ & theorem\\
(D) $\Leftrightarrow$ $\varpi_\lambda<1$ on the open path $0<\lambda\le1$ & theorem under Hypothesis~\CGref{hyp:varpi-less-one} ($\varpi_{\RVM}<1$: numerically supported, not proved)\\
Space $\mathscr O$ of ordinate lists: rigidity of the coincidence; (D) $\Leftrightarrow$ existence of a coincidence list & theorem\\
(D) $\Leftrightarrow$ $\tilde c_n(0)=\mathbf 1_{\rm prime}(n)$ for every $n$; every point of $\mathsf E$ is a witness of (D) & theorem\\
Symmetric joining factors: order one is needed; the symmetry criterion has no slack at $\lambda=0$ & theorem\\
Ladder: $\tilde h_\lambda(\sigma)=O(\sigma N^{-\sigma})$ $\Leftrightarrow$ no mass of $dJ^{\rm M}_\lambda$ outside the integers below $N$ & theorem\\
Bridge theorem in phase form; (D) $\Leftrightarrow$ \textup{(B1)}, \textup{(B2)} for $\tilde h_0$; impossible for $\lambda>0$ & theorem\\
Phase of $\zeta_\lambda$ on the critical line; cell masses as sums of a one-ordinate function & theorem\\
Transport $\mathsf S_{\mathsf f}(\lambda)$: entire of exponential type; moment series on the clock; defect of the clock by Poisson summation & theorem\\
Multiscale clock Jacobian: lower frame bound on $U\asymp\log T$, no fixed-window bound, dense range, smooth tails & theorem; quantitative nonlinear continuation open\\
(D) $\Leftrightarrow$ solvability of the rung system with free masses; the solution is nondegenerate & theorem\\
Gradient of $\mathsf V_\sigma$ on $\mathscr O$ & theorem\\
Acceptance criterion (window and frequencies); two-edge model; the finite Euler product over $2,3$ with the final step proved & theorem\\
No object of finite rank in the window of $\zeta$; the three abscissas of the diagonal operator over the primes & theorem\\
Phase defect along the path (rotation of the one-ordinate terms) & heuristic\\
Endpoint-only interpolation audit; clock-only injectivity problem; a priori rung displacement bound & Appendix~\CGref{sec:appendix-endpoint-v18}: theorem/proposition/open problem as labelled\\
Route B: compactness of near-solutions in a box; selection under uniqueness in the box & theorem (Section~\CGref{sec:routeB-finite})\\
Laplace transform of the residual with prescribed masses; uniqueness; RH $\Leftrightarrow$ B$'(C)$ & theorem; classical in substance\\
Certificate theorem for staged sections; graded extension; Gonek lists in a tube & theorem; the certificate hypotheses (A$_m$), (B) and (G1), (G2) are open and imply RH\\
Counting bound $|N_{\boldsymbol\kappa}-n_{\RVM}|\le(\mathrm c_1+\mathrm c_2\upsilon)\log(T+8)+\mathrm c_3\mathbb B+\mathrm c_4$ for quasi-solutions & theorem; void for finite sections with masses, since then $\upsilon=\infty$ (Lemma~\CGref{lem:rb-upsilon-infinite})\\
Block Kadec--Avdonin criterion; variable-length Avdonin theorem & theorem; novelty not checked\\
No $(A_2)$ for the unfolded zeta ordinates on any horizontal line & theorem, unconditional\\
Candidate Banach certificates of exact square sections up to $T=487$; basin $\sim T^{-2}$ & computation (floating point, Gaussian tests; not verified)\\
$\zeta_{Z1}=\zeta\Rightarrow M(x)=O_\eps(x^{1/2+\eps})$ & theorem (Step~3)\\
$\zeta_{Z1}=\zeta$ & open\\
\end{longtable}
\endgroup

\emph{What is proved and what is not.}  Four groups of statements are proved in full.  \emph{(1) The synthetic system.}  The clock and its Lambert inversion, the canonical product and $\zeta_{\RVM}$, the logarithmic bound for the exponential sums, the prime coordinate $J_{\RVM}$, $\pi_{\RVM}$ with its square-root estimate and its density (Theorem~\CGref{thm:clock-density}), the transition theorem, the impossibility of a Dirichlet expansion of $\zeta_{\RVM}^{\pm1}$ and the synthetic Mertens bound.  \emph{(2) The family and its endpoint.}  Every $\zeta_\lambda$ has the functional equation, order one and all its nontrivial zeros on $\Re s=\frac12$; $\zeta_{Z1}=\zeta$ is equivalent to (D), to $\varpi_{Z1}=1$ and to the constancy of $J^{\rm M}_{Z1}$ between consecutive integers together with the jump-growth bound of Theorem~\CGref{thm:cell-form}; $\varpi_{Z1}\ge1$ and $\zeta_{Z1}(\sigma)\ge\zeta(\sigma)$ for real $\sigma>1$; $\log\zeta_\lambda(\sigma)$ is strictly convex in $\lambda$; and $\zeta_{Z1}=\zeta$ implies $M(x)=O_\eps(x^{1/2+\eps})$ (Step~3).  \emph{(3) The atomized family and the criteria.}  Every $\tilde\zeta_{\Z,\lambda}$ is an absolutely convergent Dirichlet series in $\Re s>1$, zero-free in $\Re s>\frac12$, joined to $\zeta_\lambda$ by an explicit factor; no $\zeta_\lambda$ with $\lambda>0$ is a Dirichlet series; Hamburger's theorem holds for the pair in three forms (Theorems~\CGref{thm:pair-criterion}, \CGref{thm:double-hamburger}); the positivity criterion (Theorem~\CGref{thm:positivity}); and the rigidity of the coincidence in the space of ordinate lists (Theorem~\CGref{thm:rigidity}).  \emph{(4) Criteria written where the joining factor is defined, and models.}  The reduction of (D) to the cell masses (Corollary~\CGref{cor:cell-mass-form}); the ladder and the half-line form of Hamburger's theorem for a pair (Theorems~\CGref{thm:ladder}, \CGref{thm:halfline-hamburger}); the bridge theorem in phase form (Theorem~\CGref{thm:bridge}, Corollary~\CGref{cor:bridge-family}); the transport from the clock and the exact defect of the clock (Theorem~\CGref{thm:transport-band}, Proposition~\CGref{prop:clock-defect-poisson}); the existence form with free masses and the nondegeneracy of the solution (Theorems~\CGref{thm:free-masses}, \CGref{thm:nondegenerate}); the acceptance criterion and the models in which the final step is a theorem (Theorems~\CGref{thm:acceptance}, \CGref{thm:finite-euler} for the primes $2,3$, Proposition~\CGref{prop:graph-model}).

\emph{Not proved:} the identity $\zeta_{Z1}=\zeta$.  Each of the following statements is proved to be equivalent to it, and none of them is proved: (D); $\varpi_{Z1}\le1$; $\zeta_{Z1}(\sigma)\le\zeta(\sigma)$ for one $\sigma>1$; $\varpi_\lambda<1$ for all $\lambda\in(0,1]$ (this one under Hypothesis~\CGref{hyp:varpi-less-one}); $J^{\rm M}_{Z1}$ non-decreasing; $\mathsf W(0)\ge0$ \CGeqref{eq:one-number}; $\tilde h_0$ constant; the theta relation $T_0\equiv0$ \CGeqref{eq:theta-defect}; $\mathsf V_\sigma(0)=0$ \CGeqref{eq:vertical-defect}; the existence of a list in $\mathscr O$ whose continuous and atomized zeta functions coincide (Corollary~\CGref{cor:existence-form}); $\tilde c_n(0)=\mathbf 1_{\rm prime}(n)$ for every $n\ge2$ (Corollary~\CGref{cor:cell-mass-form}); $\tilde h_0(\sigma)=O_N(\sigma N^{-\sigma})$ for every $N$ \CGeqref{eq:ladder-D}; the conditions \textup{(B1)} and \textup{(B2)} for $\tilde h_0$ (Corollary~\CGref{cor:bridge-family}); $\mathsf S_{\mathsf f}(0)=\mathsf R_{\mathsf f}(1)$ for every $\mathsf f$ whose support avoids the points $\log n$, together with the cell jump-growth condition (Proposition~\CGref{prop:rung-functional}); the solvability of the rung system \CGeqref{eq:rung-system} in $\mathscr O$ (Theorem~\CGref{thm:free-masses}).  The symmetry of the joining factor $e^{\tilde h_0}$ with finite order is a sufficient condition (Theorem~\CGref{thm:pair-criterion}), and it is also necessary (Remark~\CGref{rem:pair-no-slack}).

\subsection{Organization}
\label{subsec:organization}

Section~\CGref{sec:clock} treats the clock, its Lambert inversion and the bounded displacement; Section~\CGref{sec:zeta-rvm} the canonical product; Section~\CGref{sec:exp-sums} the exponential sums; Sections~\CGref{sec:explicit-formula} and~\CGref{sec:integer-weights} the prime coordinate and the integer weights; Sections~\CGref{sec:mertens-transform} and~\CGref{sec:transition} the Mertens transform and the transition theorem; Section~\CGref{sec:comparison} the exact Hadamard product and the comparison with $\zeta$; Section~\CGref{sec:lambda-family} the $\lambda$-family; Section~\CGref{sec:Z1-step3} the endpoint, Hamburger's theorem and Step~3; Section~\CGref{sec:euler-products} the value of the endpoint at infinity and the Euler products of $\zeta_{\RVM}$ and $\zeta_{Z1}$; Section~\CGref{sec:atomized} the atomized family, Hamburger's hypotheses at the two endpoints and the programme for the final step; Section~\CGref{sec:criteria} the open path, convexity, Hamburger's theorem for the pair, the positivity criterion, the density of the clock, the Gaussian defect, the reduction to the cell masses, the half-line form, the bridge theorem and the transport from the clock; Section~\CGref{sec:lists} the space of ordinate lists, the rigidity of the coincidence, the existence form of the final step, the rung system with free masses and the gradient of the vertical defect; Section~\CGref{sec:models} the acceptance criterion and the models in which the final step is a theorem; Section~\CGref{sec:routeB-finite} route B for finite sections, the prescribed-mass uniqueness, the staged certificates, the Gonek lists, the a priori counting bound, the block Kadec--Avdonin criterion and its failure for the zeta list, with the certificate computations; Section~\CGref{sec:numerics} the numerical illustration; Section~\CGref{sec:conclusions} the conclusions and open problems; Appendix~\CGref{sec:appendix-endpoint-v18} records the endpoint-only interpolation audit, the autonomous clock-rigidity question, the a priori rung bound, and the documented bounded finite-section diagnostics.

\section{The midpoint Riemann--von Mangoldt clock}
\label{sec:clock}

Let $\rho$ run over the nontrivial zeros of $\zeta$, counted with multiplicity.  No nontrivial zero is real, so the positive ordinates can be listed with multiplicity as $0<\gamma_1\le\gamma_2\le\cdots$.  Put
\begin{equation}\label{eq:zero-count}
 N_\zeta(T):=\#\{\rho:\ 0<\Im\rho\le T\}=\#\{k:\gamma_k\le T\}.
\end{equation}
The Riemann--von Mangoldt formula \cite[Thms.~9.3--9.4]{Titchmarsh1986}, \cite{vonMangoldt1895} states that for $T\ge2$, $T\ne\gamma_k$,
\begin{equation}\label{eq:RvM-full}
 N_\zeta(T)=\frac{T}{2\pi}\log\frac{T}{2\pi e}+\frac78+S(T)+O\Bigl(\frac1T\Bigr),
 \qquad
 S(T):=\frac1\pi\arg\zeta\bigl(\tfrac12+iT\bigr)=O(\log T),
\end{equation}
where $\arg\zeta(\frac12+iT)$ is obtained by continuous variation of the argument along the segments from $2$ to $2+iT$ and from $2+iT$ to $\frac12+iT$, starting from $\arg\zeta(2)=0$.  The $O(1/T)$ term comes from Stirling's formula for $\Gamma$ and is explicit; $S(T)$ is the fluctuating part, with a jump $+m$ at an ordinate of multiplicity $m$.  The \emph{midpoint RVM clock} is the smooth part of \CGeqref{eq:RvM-full}:
\begin{equation}\label{eq:clock}
 \NR(T):=\frac{T}{2\pi}\log\frac{T}{2\pi e}+\frac78 .
\end{equation}
It is strictly increasing on $[2\pi,\infty)$.  The midpoint ordinates $\tau_k$ are the unique solutions in $[2\pi,\infty)$ of
\begin{equation}\label{eq:midpoint-condition}
 \NR(\tau_k)=k-\frac12,
 \qquad k\ge1,
\end{equation}
that is, the abscissae on the increasing branch where the smooth count crosses the midpoints $k-\frac12$ of the steps of a unit staircase.  For $k=1$ the same level has a second positive solution in $(0,2\pi)$ on the $W_{-1}$ branch; it is not part of the midpoint clock used here.  Thus $\NR$ replaces the staircase $N_\zeta$ by its smooth part, and $\tau_k$ is the point where the $k$th step would be crossed half-way (Figure~\CGref{fig:clock-zoom}).

The Lambert $W$-function is the multivalued inverse of $v\mapsto ve^v$: $W(z)e^{W(z)}=z$.  On $[-1/e,\infty)$ the equation $ve^v=z$ has exactly one solution $v\ge-1$; this defines the principal branch $W_0:[-1/e,\infty)\to[-1,\infty)$, which is real, increasing, with $W_0(0)=0$ and $W_0(z)\sim\log z-\log\log z$ as $z\to\infty$ \cite[\S4.13]{DLMF2024}, \cite{Corless1996}.

\begin{proposition}[Lambert inversion]\label{prop:lambert}
For every $k\ge1$,
\begin{equation}\label{eq:lambert}
 \tau_k=\frac{2\pi(k-11/8)}{W_0\bigl((k-11/8)/e\bigr)},
\end{equation}
The quotient is positive also for $k=1$.
\end{proposition}

\begin{proof}
Put $y=\tau/(2\pi)$.  Equation \CGeqref{eq:midpoint-condition} reads
\begin{equation}\label{eq:lambert-proof-y}
 y(\log y-1)=k-\frac{11}{8}.
\end{equation}
With $v=\log y-1$ one has $y=e^{v+1}$ and $ve^v=(k-11/8)/e$, so $v=W_0((k-11/8)/e)$ and $y=(k-11/8)/W_0((k-11/8)/e)$.  For $k=1$, $(k-11/8)/e=-3/(8e)\in(-1/e,0)$, where $W_0$ is real and negative; numerator and denominator are both negative.
\end{proof}

The clock derivative and the resulting growth are
\begin{equation}\label{eq:clock-derivative}
 \NR'(T)=\frac1{2\pi}\log\frac{T}{2\pi},
\end{equation}
\begin{equation}\label{eq:clock-growth}
 \tau_k\asymp\frac{k}{\log k}\qquad(k\ge2).
\end{equation}
In terms of the clock, \CGeqref{eq:RvM-full} reads
\begin{equation}\label{eq:RvM-formula}
 N_\zeta(T)=\NR(T)+S(T)+O(T^{-1}),\qquad S(T)=O(\log T)\qquad(T\ge2).
\end{equation}
Let
\begin{equation}\label{eq:symmetric-count}
 N_\zeta^{*}(T):=\tfrac12\bigl(N_\zeta(T+0)+N_\zeta(T-0)\bigr)
\end{equation}
be the symmetric count, and define the displacements
\begin{equation}\label{eq:eta-def}
 \eta_k:=\tau_k-\gamma_k .
\end{equation}

\begin{lemma}[Bounded displacement]\label{lem:bounded-displacement}
There is a constant $C_\eta$ such that
\begin{equation}\label{eq:bounded-displacement}
 |\eta_k|\le C_\eta\qquad(k\ge1).
\end{equation}
\end{lemma}

\begin{proof}
Let $\gamma=\gamma_k$ and let $m=m(\gamma)$ be the number of indices $j$ with $\gamma_j=\gamma$; these indices form a block $j_0+1,\dots,j_0+m$ containing $k$.  Then $N_\zeta(\gamma-0)=j_0$ and $N_\zeta(\gamma+0)=j_0+m$, so $N_\zeta^{*}(\gamma)=j_0+m/2$ and
\begin{equation}\label{eq:block-count}
 \bigl|N_\zeta^{*}(\gamma_k)-(k-\tfrac12)\bigr|\le\frac{m(\gamma_k)}2 .
\end{equation}
By \CGeqref{eq:RvM-formula}, $m(\gamma)\le N_\zeta(\gamma+1)-N_\zeta(\gamma-1)=O(\log\gamma)$, and $N_\zeta^{*}(\gamma)=\NR(\gamma)+O(\log\gamma)$.  With \CGeqref{eq:midpoint-condition},
\begin{equation}\label{eq:clock-discrepancy}
 \bigl|\NR(\gamma_k)-\NR(\tau_k)\bigr|=O(\log\gamma_k).
\end{equation}
Explicitly, $\NR(\tau_k)=k-\frac12$ and $\NR(\gamma_k)=N_\zeta^*(\gamma_k)+O(\log\gamma_k)=k+O(\log\gamma_k)$; since $\NR(T)=\frac T{2\pi}\log\frac T{2\pi e}+O(1)$ is increasing with $\NR(T)\asymp T\log T$, both relations give $\gamma_k\asymp\tau_k\asymp k/\log k$ and hence $\log\gamma_k\asymp\log\tau_k\asymp\log k$.  By the mean-value theorem, $\NR(\gamma_k)-\NR(\tau_k)=\NR'(\xi_k)(\gamma_k-\tau_k)$ with $\xi_k$ between $\gamma_k$ and $\tau_k$, and by \CGeqref{eq:clock-derivative} $\NR'(\xi_k)=\frac1{2\pi}\log\frac{\xi_k}{2\pi}\gg\log k$ because $\xi_k\ge\min(\gamma_k,\tau_k)\gg k/\log k$; with \CGeqref{eq:clock-discrepancy} this gives $|\gamma_k-\tau_k|=O(1)$.  Finitely many initial indices are absorbed into $C_\eta$.
\end{proof}

\section{The pure critical RVM zeta function}
\label{sec:zeta-rvm}

Put
\begin{equation}\label{eq:q-def}
 q(s):=s(s-1),
\end{equation}
\begin{equation}\label{eq:Pcrit}
 P_{\RVM}^{\crit}(s):=\prod_{k\ge1}\left(1+\frac{q(s)}{\tfrac14+\tau_k^2}\right),
\end{equation}
\begin{equation}\label{eq:Gdef}
 G(s):=s(s-1)\Gamma(s/2)\pi^{-s/2},
\end{equation}
and
\begin{equation}\label{eq:zetaRVM}
 \xi_{\RVM}(s):=\frac12P_{\RVM}^{\crit}(s),
 \qquad
 \zeta_{\RVM}(s):=\frac{P_{\RVM}^{\crit}(s)}{G(s)}.
\end{equation}

\begin{theorem}[Canonical characterization of $\zeta_{\RVM}$]\label{thm:canonical}
The product \CGeqref{eq:Pcrit} converges absolutely and locally uniformly in $s\in\C$.  It defines an entire function with
\begin{equation}\label{eq:P-symmetry}
 P_{\RVM}^{\crit}(1-s)=P_{\RVM}^{\crit}(s),
 \qquad
 P_{\RVM}^{\crit}(\bar s)=\overline{P_{\RVM}^{\crit}(s)},
 \qquad
 P_{\RVM}^{\crit}(0)=P_{\RVM}^{\crit}(1)=1 .
\end{equation}
Its zeros are exactly the simple zeros
\begin{equation}\label{eq:synthetic-divisor}
 \rho_k^{\RVM}=\frac12+i\tau_k,\qquad \overline{\rho_k^{\RVM}}=\frac12-i\tau_k,
 \qquad k\ge1 .
\end{equation}
Consequently $\zeta_{\RVM}$ has a simple pole at $s=1$ with residue $1$, the trivial zeros $s=-2,-4,-6,\dots$, and nontrivial zeros exactly at \CGeqref{eq:synthetic-divisor}; moreover $\xi_{\RVM}(s)=\xi_{\RVM}(1-s)$.
\end{theorem}

\begin{proof}
By \CGeqref{eq:clock-growth}, $\sum_k\tau_k^{-2}<\infty$.  On a compact set $K$,
\begin{equation}\label{eq:canonical-normal}
 \sup_{s\in K}\left|\frac{q(s)}{\tfrac14+\tau_k^2}\right|\ll_K\tau_k^{-2},
\end{equation}
so the Weierstrass criterion gives normal convergence.  Each factor depends on $s$ only through $q(s)$ and equals $1$ at $s=0,1$, which gives \CGeqref{eq:P-symmetry}; real coefficients give the conjugation symmetry.  The $k$th factor vanishes exactly when
\begin{equation}\label{eq:factor-zero}
 \Bigl(s-\frac12\Bigr)^2+\tau_k^2=0 .
\end{equation}
The clock is strictly increasing on $[2\pi,\infty)$ and $\tau_1>2\pi$, so the $\tau_k$ are distinct; the derivative $2s-1=\pm2i\tau_k\ne0$ of the vanishing factor and the non-vanishing of the other factors make the zeros simple.  At $s=1$, $P_{\RVM}^{\crit}(1)=1$ and $G(s)=(s-1)(1+O(s-1))$, so the residue is $1$.  At $s=-2n$ ($n\ge1$) the pole of $\Gamma(s/2)$ gives a zero of $1/G$, while $P_{\RVM}^{\crit}(-2n)\ne0$ because $q(-2n)>0$.  At $s=0$, $s\Gamma(s/2)=2\Gamma(1+s/2)\to2$, so $G(0)=-2\ne0$ and $s=0$ is not a zero.
\end{proof}

\section{RVM exponential sums, centroids, and the height transform}
\label{sec:exp-sums}

For $u\in\R$ and $N\ge1$ define
\begin{equation}\label{eq:S-RVM}
 S_{\RVM}(N;u):=\sum_{k\le N}e^{iu\tau_k},
 \qquad
 C_{\RVM}^{\pts}(N;u):=\frac{S_{\RVM}(N;u)}{N},
\end{equation}
and
\begin{equation}\label{eq:path-centroid}
 C_{\RVM}^{\pathc}(N;u):=\frac1N\sum_{n\le N}S_{\RVM}(n;u)
 =\sum_{k\le N}\frac{N-k+1}{N}e^{iu\tau_k}.
\end{equation}

\begin{theorem}[Uniform logarithmic exponential-sum bound]\label{thm:exp-sum}
Let $0<u_1\le u_2<\infty$.  Then
\begin{equation}\label{eq:exp-sum-bound}
 \sup_{u\in[u_1,u_2]}|S_{\RVM}(N;u)|\ll_{u_1,u_2}1+\log N .
\end{equation}
\end{theorem}

\begin{proof}
Let $\mathcal T:[-\tfrac18,\infty)\to[2\pi,\infty)$ be the inverse of the clock, so $\tau_k=\mathcal T(k-\tfrac12)$, and put $\ell_k=\log(\tau_k/2\pi)>0$.  Differentiating $\NR(\mathcal T(y))=y$,
\begin{equation}\label{eq:inverse-clock}
 \mathcal T'(y)=\frac{2\pi}{\log(\mathcal T(y)/2\pi)},
 \qquad
 \mathcal T''(y)=-\frac{4\pi^2}{\mathcal T(y)\log^3(\mathcal T(y)/2\pi)} .
\end{equation}
Since $\mathcal T'$ decreases and $|\mathcal T''|$ decreases,
\begin{equation}\label{eq:rvm-gap-bounds}
 \frac{2\pi}{\ell_{k+1}}\le\tau_{k+1}-\tau_k\le\frac{2\pi}{\ell_k},
\end{equation}
\begin{equation}\label{eq:rvm-second-difference}
 |\tau_{k+2}-2\tau_{k+1}+\tau_k|\le\frac{4\pi^2}{\tau_k\ell_k^3}.
\end{equation}
From \CGeqref{eq:midpoint-condition}, $\tau_k\ell_k=2\pi(k-\tfrac{11}8)+\tau_k$, hence
\begin{equation}\label{eq:tau-ell-k}
 \tau_k\ell_k\asymp k\qquad(k\ge2).
\end{equation}
Fix $u\in[u_1,u_2]$ and put $\Delta_k=u(\tau_{k+1}-\tau_k)$, $E_k=e^{iu\tau_k}$.  Choose $k_0=k_0(u_2)$ with $\ell_{k_0}\ge2u_2$; then $0<\Delta_k\le\pi$ for $k\ge k_0$.  With $\mathfrak b_k^{\rm sbp}:=(e^{i\Delta_k}-1)^{-1}$ one has $E_{k+1}=e^{i\Delta_k}E_k$, hence
\begin{equation}\label{eq:Ek-identity}
 E_k=\mathfrak b_k^{\rm sbp}(E_{k+1}-E_k).
\end{equation}
Summing \CGeqref{eq:Ek-identity} over $k_0\le k\le N$ and shifting the index in the first part gives the exact identity
\begin{equation}\label{eq:summation-by-parts}
 \sum_{k=k_0}^{N}E_k=\mathfrak b_N^{\rm sbp}E_{N+1}-\mathfrak b_{k_0}^{\rm sbp}E_{k_0}+\sum_{k=k_0+1}^{N}(\mathfrak b_{k-1}^{\rm sbp}-\mathfrak b_k^{\rm sbp})E_k .
\end{equation}
For $0<x\le\pi$, $|e^{ix}-1|=2\sin(x/2)\ge2x/\pi$.  Using \CGeqref{eq:rvm-gap-bounds} and \CGeqref{eq:rvm-second-difference},
\begin{equation}\label{eq:kappa-bounds}
 |\mathfrak b_k^{\rm sbp}|\le\frac{\pi}{2\Delta_k}\le\frac{\ell_{k+1}}{4u_1},
 \quad
 |\mathfrak b_{k-1}^{\rm sbp}-\mathfrak b_k^{\rm sbp}|\le|\Delta_k-\Delta_{k-1}|\,|\mathfrak b_{k-1}^{\rm sbp}||\mathfrak b_k^{\rm sbp}|
 \le\frac{\pi^2u_2}{4u_1^2}\,\frac{\ell_{k+1}^2}{\tau_{k-1}\ell_{k-1}^3}\ll_{u_1,u_2}\frac1{\tau_k\ell_k}\ll\frac1k,
\end{equation}
the last step by \CGeqref{eq:tau-ell-k}.  Inserting \CGeqref{eq:kappa-bounds} into \CGeqref{eq:summation-by-parts} gives $|\sum_{k_0}^{N}E_k|\ll\ell_{N+1}+\ell_{k_0+1}+\sum_{k\le N}k^{-1}\ll1+\log N$.  The indices $k<k_0$ contribute at most $k_0$.  All constants depend only on $u_1,u_2$.
\end{proof}

\begin{corollary}[Point centroid]\label{cor:point-centroid}
Uniformly for $u\in[u_1,u_2]\subset(0,\infty)$,
\begin{equation}\label{eq:point-centroid}
 C_{\RVM}^{\pts}(N;u)=O_{u_1,u_2}\!\left(\frac{1+\log N}{N}\right)\longrightarrow0 .
\end{equation}
Theorem~\CGref{thm:exp-sum} alone does not imply convergence of the path centroid \CGeqref{eq:path-centroid}.
\end{corollary}

\begin{proof}
The first assertion is \CGeqref{eq:exp-sum-bound} divided by $N$.  For the second, the identity $C_{\RVM}^{\pathc}=N^{-1}\sum_{n\le N}S_{\RVM}(n;u)$ with $S_{\RVM}(n;u)=O(\log n)$ gives only $O(\log N)$.
\end{proof}

Define the midpoint count
\begin{equation}\label{eq:n-RVM}
 n_{\RVM}(H):=\#\{k:\tau_k\le H\}.
\end{equation}
For $x>0$ and $n_{\RVM}(H)\ge1$ define
\begin{equation}\label{eq:MH-RVM}
 M_H^{\RVM}(x):=\frac1{n_{\RVM}(H)}\sum_{\tau_k\le H}x^{i\tau_k},
\end{equation}
and for the ordinates of $\zeta$
\begin{equation}\label{eq:MH-genuine}
 M_H(x):=\frac1{N_\zeta(H)}\sum_{\gamma_k\le H}x^{i\gamma_k}.
\end{equation}

\begin{corollary}[Synthetic height-transform decay]\label{cor:MH-decay}
For every fixed $x>1$, with $N=n_{\RVM}(H)$,
\begin{equation}\label{eq:MH-decay}
 M_H^{\RVM}(x)=O_x\!\left(\frac{1+\log N}{N}\right)\longrightarrow0\qquad(H\to\infty).
\end{equation}
\end{corollary}

\begin{proof}
Take $u=\log x>0$ in Theorem~\CGref{thm:exp-sum}; \CGeqref{eq:MH-RVM} is the point centroid at $N=n_{\RVM}(H)$.
\end{proof}

\begin{remark}\label{rem:MH-not-mertens}
Neither \CGeqref{eq:MH-RVM} nor \CGeqref{eq:MH-genuine} is the Mertens function; both are normalized phase averages.
\end{remark}

\section{The RVM explicit formula, \texorpdfstring{$J_{\RVM}$}{J\_RVM}, and \texorpdfstring{$\pi_{\RVM}$}{pi\_RVM}}
\label{sec:explicit-formula}

For $w=a+ib$ with $b>0$ define
\begin{equation}\label{eq:Ei-def}
 \Ei(w):=\int_{-\infty}^{a}\frac{e^{v+ib}}{v+ib}\,dv,
 \qquad
 \Ei(\bar w):=\overline{\Ei(w)} .
\end{equation}
The path is the horizontal line $\Im=b$, which avoids $0$; this is the branch of $\Ei(\rho\log x)$ used in Riemann's explicit formula \cite{Edwards1974}.  For $x>1$ put
\begin{equation}\label{eq:J-RVM}
 J_{\RVM}(x):=\li(x)-2\Re\sum_{k\ge1}\Ei\!\left(\Bigl(\frac12+i\tau_k\Bigr)\log x\right)-\log2+\int_x^\infty\frac{dy}{y(y^2-1)\log y},
\end{equation}
the zero sum being taken in the order of $k$.

\begin{lemma}[Asymptotics of $\Ei$ off the real axis]\label{lem:Ei}
For $w=a+ib$, $b>0$,
\begin{equation}\label{eq:Ei-remainder}
 \Ei(w)=\frac{e^w}{w}+\varrho(w),\qquad|\varrho(w)|\le\frac{e^a}{b^2}.
\end{equation}
\end{lemma}

\begin{proof}
Integrate \CGeqref{eq:Ei-def} by parts: $\int_{-\infty}^ae^{v+ib}(v+ib)^{-1}dv=e^w/w+\int_{-\infty}^ae^{v+ib}(v+ib)^{-2}dv$, the boundary term at $-\infty$ being $0$.  The last integral is bounded by $\int_{-\infty}^ae^v b^{-2}dv=e^a/b^2$.
\end{proof}

\begin{theorem}[Compact-uniform existence of the RVM zero sum]\label{thm:J-exists}
The series $\sum_k\Re\Ei\bigl((\tfrac12+i\tau_k)u\bigr)$ converges uniformly for $u$ in every compact interval $[u_1,u_2]\subset(0,\infty)$.  Hence $J_{\RVM}$ is a continuous function on $(1,\infty)$.
\end{theorem}

\begin{proof}
Apply Lemma~\CGref{lem:Ei} with $w=(\tfrac12+i\tau_k)u$, $a=u/2$, $b=\tau_ku$.  Since $\Re\bigl(e^{iu\tau_k}/(\tfrac12+i\tau_k)\bigr)=(\tfrac12\cos u\tau_k+\tau_k\sin u\tau_k)/(\tfrac14+\tau_k^2)$,
\begin{equation}\label{eq:Ei-asymptotic}
 \Re\Ei\bigl((\tfrac12+i\tau_k)u\bigr)=\frac{e^{u/2}}u\,\frac{\sin(u\tau_k)}{\tau_k}+O_{u_1,u_2}(\tau_k^{-2}),
\end{equation}
uniformly for $u\in[u_1,u_2]$.  The error series converges absolutely and uniformly.  For the main series, Abel summation with $S_{\RVM}(K;u)=O_{u_1,u_2}(1+\log K)$ (Theorem~\CGref{thm:exp-sum}) and
\begin{equation}\label{eq:abel-weights}
 \frac1{\tau_K}-\frac1{\tau_{K+1}}\le\frac{\tau_{K+1}-\tau_K}{\tau_K^2}\le\frac{2\pi}{\ell_K\tau_K^2}\ll\frac{\log K}{K^2}
\end{equation}
gives a uniformly convergent majorant $\sum_K(\log K)^2K^{-2}$, while the boundary term $S_{\RVM}(K;u)/\tau_K\to0$ uniformly.  Hence $\sum_ke^{iu\tau_k}/\tau_k$ converges uniformly, and so does its imaginary part.
\end{proof}

Define the support-truncated version
\begin{equation}\label{eq:J-star}
 J_{\RVM}^{*}(x):=\begin{cases}J_{\RVM}(x),&x\ge2,\\0,&0<x<2,\end{cases}
\end{equation}
set $\pi_{\RVM}(x)=0$ for $0<x<2$, and for $x\ge2$
\begin{equation}\label{eq:pi-RVM}
 \pi_{\RVM}(x):=\sum_{1\le m\le\log_2x}\frac{\mu(m)}m\,J_{\RVM}^{*}(x^{1/m}).
\end{equation}

\begin{proposition}[Exact finite prime-power inversion]\label{prop:finite-inversion}
For every $x\ge2$,
\begin{equation}\label{eq:J-pi-inverse}
 J_{\RVM}^{*}(x)=\sum_{1\le m\le\log_2x}\frac1m\,\pi_{\RVM}(x^{1/m}).
\end{equation}
\end{proposition}

\begin{proof}
All sums are finite because $x^{1/m}<2$ for $m>\log_2x$.  Substituting \CGeqref{eq:pi-RVM} into the right side of \CGeqref{eq:J-pi-inverse}, the coefficient of $J_{\RVM}^{*}(x^{1/r})$ is
\begin{equation}\label{eq:mobius-coefficient}
 \frac1r\sum_{d\mid r}\mu(d),
\end{equation}
which is $1$ for $r=1$ and $0$ otherwise.
\end{proof}

\subsection{The square-root estimate}
\label{subsec:square-root}

For $u>0$ put
\begin{equation}\label{eq:A-RVM}
 \mathcal A_{\RVM}(u):=\sum_{k\ge1}\frac{e^{iu\tau_k}}{\tau_k},
\end{equation}
summed in the order of $k$, and write $\ell(t)=\log(t/2\pi)$, so that $\NR'(t)=\ell(t)/2\pi$ and $\ell(\tau_k)=\ell_k$.

\begin{lemma}[The weighted clock sum]\label{lem:A-RVM}
For every $u\ge\log2$ the series \CGeqref{eq:A-RVM} converges, and
\begin{equation}\label{eq:A-RVM-bound}
 |\mathcal A_{\RVM}(u)|\le\frac{\log^2(2+2u)}{4\pi}+\frac{\log(2+2u)}{2\pi}+C_{\mathcal A},
\end{equation}
with an absolute constant $C_{\mathcal A}$.
\end{lemma}

\begin{proof}
Put $Y_j=2^j\cdot4\pi(1+u)$ for $j\ge0$.

\emph{Head $\tau_k\le Y_0$.}  Let $n_{\RVM}(t)=\#\{k:\tau_k\le t\}$.  By \CGeqref{eq:midpoint-condition}, $\tau_k\le t$ if and only if $k\le\NR(t)+\tfrac12$, so $n_{\RVM}(t)\le\NR(t)+\tfrac12=\frac t{2\pi}(\ell(t)-1)+\frac{11}8\le\frac{t\,\ell(t)}{2\pi}$ for $t\ge\tau_1>\frac{11\pi}4$.  Stieltjes integration gives
\begin{equation}\label{eq:A-head}
 \sum_{\tau_k\le Y_0}\frac1{\tau_k}=\frac{n_{\RVM}(Y_0)}{Y_0}+\int_{\tau_1}^{Y_0}\frac{n_{\RVM}(t)}{t^2}\,dt\le\frac{\ell(Y_0)}{2\pi}+\int_{\tau_1}^{Y_0}\frac{\ell(t)}{2\pi t}\,dt\le\frac{\ell(Y_0)}{2\pi}+\frac{\ell(Y_0)^2}{4\pi},
\end{equation}
and $\ell(Y_0)=\log(2+2u)$.

\emph{Blocks.}  Let $\mathcal B_j=\{k:Y_j<\tau_k\le Y_{j+1}\}$, let $\mathcal T$ be the inverse clock of \CGeqref{eq:inverse-clock}, and put $\phi(y)=u\,\mathcal T(y-\tfrac12)/2\pi$, so that $e^{2\pi i\phi(k)}=e^{iu\tau_k}$.  By \CGeqref{eq:inverse-clock},
\begin{equation}\label{eq:phi-second}
 \phi''(y)=-\frac{2\pi u}{\mathcal T\,\ell(\mathcal T)^3},\qquad \mathcal T=\mathcal T(y-\tfrac12).
\end{equation}
On $\mathcal B_j$ put $\varkappa_j=2\pi u/(Y_{j+1}\ell(Y_{j+1})^3)$.  Since $Y_{j+1}=2Y_j$ and $\ell(Y_{j+1})\le2\ell(Y_j)$ (because $\ell(Y_j)\ge\ell(Y_0)\ge\log2$), we have $\varkappa_j\le-\phi''\le16\varkappa_j$ at the integers of the block, and $\varkappa_j\le-\phi''\le256\varkappa_j$ on the whole real interval $[a,b]$ over which a run of $\mathcal B_j$ is summed, since that interval may begin inside the previous block, across which $-\phi''$ changes by at most another factor $16$.  The number of indices in $\mathcal B_j$ is at most $\NR(Y_{j+1})-\NR(Y_j)+1\le Y_j\ell(Y_{j+1})/2\pi+1$.  The second-derivative test for exponential sums \cite[Thm.~5.9]{Titchmarsh1986} gives, for every run $a<k\le b$ of consecutive indices of $\mathcal B_j$,
\begin{equation}\label{eq:block-vdc}
 \Bigl|\sum_{a<k\le b}e^{iu\tau_k}\Bigr|\le A_{\rm vdC}\Bigl(256\Bigl(\frac{Y_j\ell(Y_{j+1})}{2\pi}+1\Bigr)\varkappa_j^{1/2}+\varkappa_j^{-1/2}\Bigr)=:\mathcal P_j,
\end{equation}
with an absolute constant $A_{\rm vdC}$; a run of length less than $1$ has at most one term, and $1\le\varkappa_j^{-1/2}$ because $\varkappa_j\le u/(4(1+u)\log^32)<1$.  The weights $1/\tau_k$ are positive and decreasing, and $1/\tau_k<1/Y_j$ on $\mathcal B_j$, so Abel's inequality bounds the contribution of any run in $\mathcal B_j$ by $\mathcal P_j/Y_j$:
\begin{equation}\label{eq:block-weighted}
 \frac{\mathcal P_j}{Y_j}\le A_{\rm vdC}\left(\frac{16}{2\pi}\Bigl(\frac{\pi u}{Y_j\ell(Y_{j+1})}\Bigr)^{1/2}+\frac{16}{Y_j}+\Bigl(\frac{\ell(Y_{j+1})^3}{\pi uY_j}\Bigr)^{1/2}\right).
\end{equation}
With $Y_j=2^j\cdot4\pi(1+u)$ and $\ell(Y_{j+1})=\log(2^{j+2}(1+u))\le j+2+\log(1+u)$, each of the three terms is at most an absolute constant times $(j+2)^{3/2}2^{-j/2}$, uniformly in $u\ge\log2$, because $\sup_{u\ge\log2}(1+\log(1+u))^3/(u(1+u))<\infty$.  Hence $\sum_{j\ge0}\mathcal P_j/Y_j\le C_{\mathcal A}$ with $C_{\mathcal A}$ absolute.

\emph{Convergence and bound.}  Any finite range $K'<k\le K''$ of indices, with $\tau_{K'+1}\in\mathcal B_J$, consists of at most two partial runs (the initial segment of $\mathcal B_J$ and the terminal segment of the block containing $\tau_{K''}$) and of complete blocks in between; each run, partial or complete, contributes at most $\mathcal P_j/Y_j$ by \CGeqref{eq:block-vdc} and Abel's inequality, so $|\sum_{K'<k\le K''}e^{iu\tau_k}/\tau_k|\le\sum_{j\ge J}\mathcal P_j/Y_j\to0$ as $K'\to\infty$ (then $J\to\infty$), uniformly in $K''$.  The series converges, and \CGeqref{eq:A-head} together with $\sum_j\mathcal P_j/Y_j\le C_{\mathcal A}$ gives \CGeqref{eq:A-RVM-bound}.
\end{proof}

\begin{theorem}[Square-root estimate for $J_{\RVM}$ and $\pi_{\RVM}$]\label{thm:Lambert-estimate}
For $x\ge8$,
\begin{equation}\label{eq:J-baseline}
 J_{\RVM}(x)-\li(x)=O\!\left(\frac{\sqrt x(\log\log x)^2}{\log x}\right),
\end{equation}
\begin{equation}\label{eq:pi-baseline}
 \pi_{\RVM}(x)-\li(x)=O\!\left(\frac{\sqrt x(\log\log x)^2}{\log x}\right),
\end{equation}
with absolute implied constants.  No hypothesis on $\zeta$ is used.
\end{theorem}

\begin{proof}
Let $u=\log x\ge\log8$ and $\rho_k^{\RVM}=\tfrac12+i\tau_k$.  Lemma~\CGref{lem:Ei} with $w=\rho_k^{\RVM}u$, $a=u/2$, $b=\tau_ku$ gives
\begin{equation}\label{eq:Ei-split}
 \Ei(\rho_k^{\RVM}u)=\frac{e^{\rho_k^{\RVM}u}}{\rho_k^{\RVM}u}+\varrho_k,\qquad|\varrho_k|\le\frac{e^{u/2}}{\tau_k^2u^2}.
\end{equation}
Moreover $\frac1{\rho_k^{\RVM}}=-\frac i{\tau_k}+r_k$ with $r_k=\frac{i/2}{\tau_k\rho_k^{\RVM}}$, so $|r_k|\le\frac1{2\tau_k^2}$.  Since $\sum_k\tau_k^{-2}<\infty$, Lemma~\CGref{lem:A-RVM} shows that the series in \CGeqref{eq:J-RVM} converges in the order of $k$ and
\begin{equation}\label{eq:Ei-sum-A}
 \sum_{k\ge1}\Ei(\rho_k^{\RVM}u)=\frac{e^{u/2}}u\Bigl(-i\,\mathcal A_{\RVM}(u)+\sum_kr_ke^{iu\tau_k}\Bigr)+\sum_k\varrho_k
 =\frac{e^{u/2}}u\bigl(-i\,\mathcal A_{\RVM}(u)+O(1)\bigr).
\end{equation}
For $u\ge2$ one has $\log(2+2u)\le\log(2u)+\log2\le3\log u$, so \CGeqref{eq:A-RVM-bound} gives $\mathcal A_{\RVM}(u)=O(\log^2u)$.  The remaining terms $-\log2+\int_x^\infty dy/(y(y^2-1)\log y)$ of \CGeqref{eq:J-RVM} are bounded for $x\ge2$.  With $e^{u/2}=\sqrt x$ and $\log u=\log\log x$ this proves \CGeqref{eq:J-baseline}.

For \CGeqref{eq:pi-baseline}, \CGeqref{eq:pi-RVM} gives
\begin{equation}\label{eq:pi-minus-li}
 \pi_{\RVM}(x)-\li(x)=\bigl(J_{\RVM}(x)-\li(x)\bigr)+\sum_{2\le m\le\log_2x}\frac{\mu(m)}m\,J_{\RVM}^{*}(x^{1/m}).
\end{equation}
By \CGeqref{eq:J-baseline}, $\li(y)\ll y/\log y$, and the continuity of $J_{\RVM}$ on $[2,8]$ (Theorem~\CGref{thm:J-exists}), $|J^*_{\RVM}(y)|\le C\,y/\log y$ for all $y\ge2$.  The term $m=2$ is therefore $O(\sqrt x/\log x)$, and the terms $3\le m\le\log_2x$ add up to at most $C\,x^{1/3}(\log_2x)^2/\log x=O(x^{1/3}\log x)$.  Both are absorbed in \CGeqref{eq:pi-baseline}.
\end{proof}

\begin{corollary}[The three limits]\label{cor:three-limits}
As $x\to\infty$,
\begin{equation}\label{eq:three-limits}
 \frac{\li(x)}{\pi_{\RVM}(x)}\to1,\qquad
 \frac{\pi(x)}{\li(x)}\to1,\qquad
 \frac{\pi(x)}{\pi_{\RVM}(x)}\to1 .
\end{equation}
More precisely, for some absolute $c_0>0$,
\begin{equation}\label{eq:pi-minus-piRVM}
 \pi(x)-\pi_{\RVM}(x)=O\!\left(x\,e^{-c_0\sqrt{\log x}}\right).
\end{equation}
\end{corollary}

\begin{proof}
Since $\li(x)\sim x/\log x$, \CGeqref{eq:pi-baseline} gives $\pi_{\RVM}(x)=\li(x)\bigl(1+O(x^{-1/2}(\log\log x)^2)\bigr)$; in particular $\pi_{\RVM}(x)>0$ for large $x$, and the first limit follows.  The second is the prime number theorem \cite{Hadamard1896,delaValleePoussin1896}, in the form $\pi(x)=\li(x)+O(x\,e^{-c_0\sqrt{\log x}})$ \cite[Ch.~18]{Davenport1980}.  Subtracting \CGeqref{eq:pi-baseline} gives \CGeqref{eq:pi-minus-piRVM}, and dividing by $\pi_{\RVM}(x)\sim x/\log x$ gives the third limit.
\end{proof}

The first limit and \CGeqref{eq:pi-baseline} concern only the synthetic system; the second and third compare it with the primes and use the prime number theorem.  Figures~\CGref{fig:prime-counts} and~\CGref{fig:pi-ratios} show them numerically.

\section{Integer weights and the continuous-to-atomic problem}
\label{sec:integer-weights}

Define
\begin{equation}\label{eq:cn}
 c_n^{\RVM}:=\pi_{\RVM}(n)-\pi_{\RVM}(n-1),\qquad n\ge2,
\end{equation}
and, initially for $\Re s>1$,
\begin{equation}\label{eq:weighted-zeta}
 \zeta_{\Z,\RVM}(s):=\prod_{n\ge2}(1-n^{-s})^{-c_n^{\RVM}}.
\end{equation}
No positivity or $0/1$ property is built into \CGeqref{eq:cn}.  Define the integer-step prime-power count and the right-continuous extension
\begin{equation}\label{eq:JZ}
 J_{\Z,\RVM}(x):=\sum_{1\le m\le\log_2x}\frac1m\,\pi_{\RVM}\bigl(\lfloor x^{1/m}\rfloor\bigr),
\end{equation}
\begin{equation}\label{eq:J-star-Z}
 J_{\RVM}^{*,\Z}(x):=J_{\RVM}^{*}(\lfloor x\rfloor).
\end{equation}

\begin{theorem}[Lipschitz control of $J_{\RVM}$]\label{thm:unit-cell}
There is an absolute constant $C_J$ such that
\begin{equation}\label{eq:unit-cell}
 |J_{\RVM}(x+\theta)-J_{\RVM}(x)|\le C_J\,\theta\qquad(x\ge2,\ 0<\theta\le1).
\end{equation}
\end{theorem}

\begin{proof}
Put $u_1=\log x$, $u_2=\log(x+\theta)$, so $0<u_2-u_1=\log(1+\theta/x)\le\theta/x$ and $u_2-u_1\le\log2$.  The terms $\li$ and $\int_x^\infty dy/(y(y^2-1)\log y)$ of \CGeqref{eq:J-RVM} have derivatives bounded by $1/\log2$ and $3/\log2$ on $[2,\infty)$, so they change by $O(\theta)$.  It remains to bound $\sum_kD_k$, where, by \CGeqref{eq:Ei-def} and $\frac{d}{dv}\Ei(\rho v)=e^{\rho v}/v$,
\begin{equation}\label{eq:Dk-def}
 D_k:=\Ei(\rho_k^{\RVM}u_2)-\Ei(\rho_k^{\RVM}u_1)=\int_{u_1}^{u_2}\frac{e^{v/2}e^{iv\tau_k}}{v}\,dv .
\end{equation}
The series $\sum_kD_k$ converges (Theorem~\CGref{thm:J-exists} and its proof).  By Abel's theorem for the increasing exponents $\tau_k$, $\sum_kD_k=\lim_{\alpha\downarrow0}\sum_ke^{-\alpha\tau_k}D_k$, and for $\alpha>0$ Fubini gives
\begin{equation}\label{eq:damped-sum}
 \sum_ke^{-\alpha\tau_k}D_k=\int_{u_1}^{u_2}\frac{e^{v/2}}v\,S_\alpha(v)\,dv,\qquad S_\alpha(v):=\sum_{k\ge1}g(\tau_k),\quad g(\tau):=e^{(iv-\alpha)\tau}.
\end{equation}
\emph{Clock decomposition.}  With $\tau_0=4\pi$ and the sawtooth count \CGeqref{eq:sawtooth-count} (proved later in Section~\CGref{sec:transition}), Stieltjes integration exactly as in the forward reference \CGeqref{eq:transition-split} ($g(\infty)=0$, $g'\in L^1$) gives
\begin{equation}\label{eq:S-eps-split}
 S_\alpha(v)=\int_{\tau_0}^\infty g\,d\NR+g(\tau_0)\NR(\tau_0)+\int_{\tau_0}^\infty\psi\bigl(\NR(\tau)+\tfrac12\bigr)g'(\tau)\,d\tau .
\end{equation}
The middle term is at most $\tfrac12$.  For the first, two integrations by parts with $\NR'(\tau)=\frac1{2\pi}\log\frac\tau{2\pi}$, $\NR''(\tau)=\frac1{2\pi\tau}$ and $v\ge\log2$ give $|\int_{\tau_0}^\infty g\,d\NR|\le C/v$, uniformly in $\alpha\in(0,1]$.

\emph{Harmonics.}  The partial sums of $\psi(y+\frac12)=-\sum_{m\ge1}(-1)^m\sin(2\pi my)/(\pi m)$ are uniformly bounded, $\sup_{M\ge1}\sup_{y\in\mathbb R}\bigl|\sum_{m\le M}\sin(2\pi my)/m\bigr|<\infty$ (for $0<y\le\frac12$ split at $M_0:=\lfloor1/y\rfloor$: the terms $m\le M_0$ contribute at most $\sum_{m\le M_0}2\pi y\le2\pi$, and Abel summation with $|\sum_{M_0<m\le M'}\sin(2\pi my)|\le1/\sin(\pi y)\le1/(2y)$ and $1/m\le1/M_0\le2y$ bounds the rest by $2$; the sum is odd and $1$-periodic in $y$), and the series converges to $\psi(y+\frac12)$ for every $y\notin\frac12+\mathbb Z$; the exceptional set $\{\tau:\NR(\tau)+\frac12\in\mathbb Z\}$ is discrete, hence Lebesgue-null.  Since $g'\in L^1(\tau_0,\infty)$ for $\alpha>0$, dominated convergence allows termwise integration:
\begin{equation}\label{eq:K-harmonics}
 \int_{\tau_0}^\infty\psi\bigl(\NR+\tfrac12\bigr)g'\,d\tau=-\sum_{m\ge1}\frac{(-1)^m}{2\pi i\,m}\bigl(K_m^+-K_m^-\bigr),\qquad
 K_m^\pm:=\int_{\tau_0}^\infty(iv-\alpha)e^{-\alpha\tau}e^{i\Psi_m^\pm(\tau)}\,d\tau,
\end{equation}
with $\Psi_m^\pm(\tau)=v\tau\pm2\pi m\NR(\tau)$.  The weight $e^{-\alpha\tau}$ is positive and decreasing and $|iv-\alpha|\le v+1$ for $0<\alpha\le1$; the second mean-value theorem (Bonnet's form), applied to the real and to the imaginary part of $e^{i\Psi_m^\pm}$ separately, bounds $|\int_a^be^{-\alpha\tau}e^{i\Psi_m^\pm}d\tau|$ by $2e^{-\alpha a}\sup_{\xi\in[a,b]}|\int_a^\xi e^{i\Psi_m^\pm}d\tau|$, so each integral is at most $2(v+1)\sup_{\xi}|\int_a^\xi e^{i\Psi_m^\pm}d\tau|$ over the relevant interval $[a,b]$.  We use the first- and second-derivative tests \cite[Lemmas~4.2 and~4.4]{Titchmarsh1986} for a real phase $\varTheta$: $|\int_a^b e^{i\varTheta}|\le4/\mu$ if $\varTheta'$ is monotone with $|\varTheta'|\ge\mu$, and $\le8/\sqrt\lambda$ if $|\varTheta''|\ge\lambda$.

For $K_m^+$: $(\Psi_m^+)'=v+m\log\frac\tau{2\pi}\ge v+m\log2$ is increasing, so $|K_m^+|\le8(v+1)/(v+m\log2)$.

For $K_m^-$: $(\Psi_m^-)'=v-m\log\frac\tau{2\pi}$ is decreasing and vanishes at $T_m:=2\pi e^{v/m}>2\pi$, and $(\Psi_m^-)''=-m/\tau$.  On $[\tau_0,T_m/2]$ (if nonempty) $(\Psi_m^-)'\ge m\log2$; on $[2T_m,\infty)$, $(\Psi_m^-)'\le-m\log2$; on $[\max(\tau_0,T_m/2),2T_m]$, $|(\Psi_m^-)''|\ge m/(2T_m)$.  Hence
\begin{equation}\label{eq:K-minus-bound}
 |K_m^-|\le2(v+1)\Bigl(\frac{8}{m\log2}+8\sqrt{\frac{2T_m}m}\Bigr)
 =2(v+1)\Bigl(\frac{8}{m\log2}+16\sqrt\pi\,\frac{e^{v/(2m)}}{\sqrt m}\Bigr).
\end{equation}
Summing with the weights $1/(2\pi m)$ and using $\sum_mm^{-3/2}e^{v/(2m)}\le e^{v/2}+\zeta(\tfrac32)e^{v/4}$,
\begin{equation}\label{eq:S-eps-bound}
 |S_\alpha(v)|\le C\,(v+1)\,e^{v/2}\qquad(v\ge\log2,\ 0<\alpha\le1),
\end{equation}
with an absolute constant $C$.

\emph{Conclusion.}  Inserting \CGeqref{eq:S-eps-bound} into \CGeqref{eq:damped-sum}, for $v\in[u_1,u_2]$ we have $e^{v/2}\le\sqrt{2x}$ and $(v+1)/v\le1+1/\log2$, so
\begin{equation}\label{eq:Dk-sum-bound}
 \Bigl|\sum_ke^{-\alpha\tau_k}D_k\Bigr|\le C'(u_2-u_1)\,x\le C'\theta,
\end{equation}
uniformly in $\alpha$.  Letting $\alpha\downarrow0$ and taking $2\Re$ proves \CGeqref{eq:unit-cell}.
\end{proof}

The $m=1$ harmonic, with stationary point $T_1=2\pi x$, produces the unit-periodic ripple of $J_{\RVM}$ (Figure~\CGref{fig:prime-counts}); the harmonics $m\ge2$ have stationary points $2\pi x^{1/m}$ and contribute $O(x^{-1/2+1/(2m)}m^{-3/2})$ to a unit increment.

\begin{corollary}[Bounded cell masses]\label{cor:cn-bounded}
\begin{equation}\label{eq:unit-cell-consequence}
 \pi_{\RVM}(y)-\pi_{\RVM}(\lfloor y\rfloor)=O(1),\qquad c_n^{\RVM}=O(1).
\end{equation}
\end{corollary}

\begin{proof}
By \CGeqref{eq:pi-RVM}, $c_n^{\RVM}=\sum_{m\le\log_2n}\frac{\mu(m)}m\bigl(J_{\RVM}^{*}(n^{1/m})-J_{\RVM}^{*}((n-1)^{1/m})\bigr)$.  The term $m=1$ is $O(1)$ by \CGeqref{eq:unit-cell}.  For $m\ge2$ and $(n-1)^{1/m}\ge2$, \CGeqref{eq:unit-cell} with $\theta=n^{1/m}-(n-1)^{1/m}\le n^{1/m-1}$ bounds the term by $C_Jn^{1/m-1}\le C_Jn^{-1/2}$, and there are at most $\log_2n$ such terms.  The jump of $J_{\RVM}^{*}$ at $2$ enters only for $n=2^m$, once, with size $|J_{\RVM}(2)|/m$.  The same argument applies to $\pi_{\RVM}(y)-\pi_{\RVM}(\lfloor y\rfloor)$.
\end{proof}

\begin{corollary}[The integer-weighted zeta is a Dirichlet series]\label{cor:ZRVM-Dirichlet}
The product \CGeqref{eq:weighted-zeta} converges absolutely for $\Re s>1$ and
\begin{equation}\label{eq:ZRVM-Dirichlet}
 \zeta_{\Z,\RVM}(s)=\sum_{N\ge1}\frac{a_N^{\Z}}{N^s},\qquad a_1^{\Z}=1,\qquad \sum_{N\ge1}\frac{|a_N^{\Z}|}{N^\sigma}<\infty\quad(\sigma>1).
\end{equation}
The inverse is also an absolutely convergent Dirichlet series on the whole half-plane:
\begin{equation}\label{eq:ZRVM-inverse-Dirichlet}
 \frac1{\zeta_{\Z,\RVM}(s)}=\sum_{N\ge1}\frac{\bar a_N^{\Z}}{N^s},\qquad \bar a_1^{\Z}=1,\qquad \sum_{N\ge1}\frac{|\bar a_N^{\Z}|}{N^\sigma}<\infty\quad(\sigma>1).
\end{equation}
\end{corollary}

\begin{proof}
By Corollary~\CGref{cor:cn-bounded},
\begin{equation}\label{eq:log-ZRVM}
 \log\zeta_{\Z,\RVM}(s)=\sum_{n\ge2}c_n^{\RVM}\sum_{m\ge1}\frac{n^{-ms}}m=\sum_{N\ge2}\frac{b_N}{N^s},\qquad b_N:=\sum_{n^m=N}\frac{c_n^{\RVM}}m,
\end{equation}
with $\sum_N|b_N|N^{-\sigma}\le C\sum_{n\ge2}\sum_m n^{-m\sigma}/m<\infty$ for $\sigma>1$.  The exponential of an absolutely convergent Dirichlet series without constant term is an absolutely convergent Dirichlet series with $a_1^{\Z}=1$, and $\sum_N|a_N^{\Z}|N^{-\sigma}\le\exp\bigl(\sum_N|b_N|N^{-\sigma}\bigr)$.  The same applies to $1/\zeta_{\Z,\RVM}=\exp\bigl(-\sum_Nb_NN^{-s}\bigr)$, which gives \CGeqref{eq:ZRVM-inverse-Dirichlet}.
\end{proof}

Numerically the cell masses lie in $[-0.38,1.04]$ for $n\le1000$ (Experiment~\CGref{exp:primes}).

\begin{theorem}[Mellin--Stieltjes identification]\label{thm:MS}
There is a function $J_{\RVM}^{\rm M}$ on $(1,\infty)$ with $J_{\RVM}^{\rm M}(x)x^{-c}\in L^2((1,\infty),dx/x)$ for every $c>1$ such that
\begin{equation}\label{eq:MS-RVM}
 \log\zeta_{\RVM}(s)=s\int_1^\infty J_{\RVM}^{\rm M}(x)\,x^{-s-1}\,dx\qquad(\Re s>1),
\end{equation}
and $J_{\RVM}^{\rm M}-J_{\RVM}$ is constant on $(1,\infty)$.  In particular $J_{\RVM}^{\rm M}=J_{\RVM}^{*}+C_{\RVM}$ on $[2,\infty)$ for a constant $C_{\RVM}$.
\end{theorem}

The proof is given in Section~\CGref{sec:euler-products} (Theorems~\CGref{thm:mellin-counts} and~\CGref{thm:distributional-EF}); it uses the half-plane bound of Theorem~\CGref{thm:halfplane-RVM}.

\begin{corollary}[Weighted factorization]\label{cor:weighted-factorization}
\begin{equation}\label{eq:JZ-diff}
 J_{\Z,\RVM}(x)-J_{\RVM}^{*}(x)=O(1+\log\log x),
\end{equation}
and
\begin{equation}\label{eq:weighted-factorization}
 \zeta_{\Z,\RVM}(s)=\zeta_{\RVM}(s)\,e^{h_{\RVM}(s)},
\end{equation}
where
\begin{equation}\label{eq:hRVM}
 h_{\RVM}(s):=s\int_1^\infty\bigl(J_{\Z,\RVM}(x)-J_{\RVM}^{\rm M}(x)\bigr)x^{-s-1}\,dx
\end{equation}
is holomorphic for $\Re s>0$.  The two functions have the same zeros in $\Re s>0$, and $\zeta_{\Z,\RVM}$ inherits through $e^{h_{\RVM}}$ the pole of $\zeta_{\RVM}$ at $s=1$; compare the modified Euler products of Grosswald and Schnitzer \cite{GrosswaldSchnitzer1978}.
\end{corollary}

\begin{proof}
By \CGeqref{eq:unit-cell-consequence} replacing $x^{1/m}$ by its integer part costs $O(1/m)$ in the $m$th summand; summing $m^{-1}$ up to $\log_2x$ gives \CGeqref{eq:JZ-diff}.  The product \CGeqref{eq:weighted-zeta} is initially defined for $\Re s>1$.  For non-integral exponents we use the principal logarithmic branch, equivalently $-\log(1-n^{-s})=\sum_{m\ge1}n^{-ms}/m$, which is unambiguous because $|n^{-s}|<1$ there.  It gives $\log\zeta_{\Z,\RVM}(s)=s\int_1^\infty J_{\Z,\RVM}(x)x^{-s-1}dx$ for $\Re s>1$, because $\sum_{n\le y}c_n^{\RVM}=\pi_{\RVM}(\lfloor y\rfloor)$.  Subtracting \CGeqref{eq:MS-RVM} gives \CGeqref{eq:weighted-factorization} on $\Re s>1$.  On $[2,\infty)$, $J_{\Z,\RVM}-J_{\RVM}^{\rm M}=(J_{\Z,\RVM}-J_{\RVM}^{*})-C_{\RVM}=O(1+\log\log x)$ by \CGeqref{eq:JZ-diff}; on $(1,2)$, $J_{\Z,\RVM}=0$ and $J_{\RVM}^{\rm M}$ is square-integrable.  Hence \CGeqref{eq:hRVM} converges locally normally for $\Re s>0$, and \CGeqref{eq:weighted-factorization} defines the holomorphic continuation of $\zeta_{\Z,\RVM}$ from $\Re s>1$ to $\Re s>0$; exponentiation introduces no zeros.
\end{proof}

\begin{lemma}[Exact first-harmonic cell cancellation]\label{lem:cell-cancel}
For every integer $n$,
\begin{equation}\label{eq:cell-cancel}
 \int_{n-1}^{n}2\cos(2\pi x)\,dx=0 .
\end{equation}
\end{lemma}

\begin{hypothesis}[Uniform cell remainder away from the exceptional powers of $2$; proved in Corollary~\CGref{cor:clock-cells}]\label{hyp:cell-remainder}
For $n\ge3$ with $n\notin\{2^j:\mu(j)\ne0\}$,
\begin{equation}\label{eq:cell-remainder}
 c_n^{\RVM}-\int_{n-1}^{n}\left(\frac1{\log x}+2\cos(2\pi x)\right)dx=o\!\left(\frac1{\log n}\right).
\end{equation}
\end{hypothesis}

Hypothesis~\CGref{hyp:cell-remainder} is proved below, in Corollary~\CGref{cor:clock-cells}, for exactly the set of $n$ in its statement; it is stated here as a hypothesis only because its proof needs the density theorem of Section~\CGref{sec:criteria}, and the proposition that follows is therefore unconditional.

\begin{proposition}[Eventual positivity]\label{prop:eventual-positive}
Hypothesis~\CGref{hyp:cell-remainder} implies
\begin{equation}\label{eq:eventual-positive}
 c_n^{\RVM}\sim\frac1{\log n},
\end{equation}
hence $0<c_n^{\RVM}<1$ for all sufficiently large $n$ in that set.
\end{proposition}

\begin{proof}
The main-density integral is $1/\log n+O(1/(n\log^2n))$; the first harmonic vanishes by Lemma~\CGref{lem:cell-cancel}.
\end{proof}

\section{The critical RVM Mertens transform}
\label{sec:mertens-transform}

Define
\begin{equation}\label{eq:F-RVM}
 F_{\RVM}(s)=\widehat M_{\RVM}^{\crit}(s):=\frac1{\zeta_{\RVM}(s)}=\frac{G(s)}{P_{\RVM}^{\crit}(s)} .
\end{equation}
Formally the associated summatory object is the inverse Mellin transform
\begin{equation}\label{eq:Mcrit-formal}
 M_{\RVM}^{\crit}(x)=\frac1{2\pi i}\int_{(c)}F_{\RVM}(s)\frac{x^s}{s}\,ds ;
\end{equation}
its existence as a symmetric limit, for every $c>\tfrac12$, is proved in Theorem~\CGref{thm:unsmoothed-RVM}.

\begin{proposition}[Synthetic pole abscissa]\label{prop:Mcrit-poles}
$F_{\RVM}=G/P_{\RVM}^{\crit}$ is meromorphic on $\C$ and holomorphic in $\Re s>\tfrac12$.  The poles of this meromorphic function are the simple poles
\begin{equation}\label{eq:Mcrit-poles}
 \frac12\pm i\tau_k\quad(k\ge1)\qquad\text{and}\qquad s=-2,-4,-6,\dots
\end{equation}
and $s=1$ is a simple zero of $F_{\RVM}$.
\end{proposition}

\begin{proof}
This is Theorem~\CGref{thm:canonical} applied to $F_{\RVM}=G/P_{\RVM}^{\crit}$: the zeros of $P_{\RVM}^{\crit}$ are \CGeqref{eq:synthetic-divisor}, and the poles of $G$ are those of $\Gamma(s/2)$ except $s=0$, where $s\Gamma(s/2)\to2$.
\end{proof}

\begin{proposition}[Gaussian square-root bound under vertical growth]\label{prop:Gaussian}
Fix $c>1$ and $0<\eps<c-\tfrac12$.  Assume, uniformly for $\tfrac12+\eps\le\sigma\le c$,
\begin{equation}\label{eq:subgaussian}
 F_{\RVM}(\sigma+it)=\exp(o(t^2))\qquad(|t|\to\infty).
\end{equation}
Define
\begin{equation}\label{eq:MG}
 M_{\RVM,G}^{\crit}(x):=\frac1{2\pi i}\int_{(c)}e^{s^2}F_{\RVM}(s)\frac{x^s}{s}\,ds .
\end{equation}
Then
\begin{equation}\label{eq:MG-bound}
 M_{\RVM,G}^{\crit}(x)=O_\eps(x^{1/2+\eps}).
\end{equation}
\end{proposition}

\begin{proof}
On $s=\sigma+it$, $|e^{s^2}|=e^{\sigma^2-t^2}$, so \CGeqref{eq:subgaussian} gives absolute convergence on vertical lines and makes the horizontal integrals tend to $0$.  The strip $\tfrac12+\eps\le\Re s\le c$ contains no pole of $F_{\RVM}$ (Proposition~\CGref{prop:Mcrit-poles}) and does not contain $s=0$, so Cauchy's theorem moves the line to $\Re s=\tfrac12+\eps$, where $|x^s|=x^{1/2+\eps}$.
\end{proof}

Theorem~\CGref{thm:unsmoothed-RVM} below removes both the assumption \CGeqref{eq:subgaussian} and the smoothing.

\section{The transition theorem and the unsmoothed synthetic Mertens bound}
\label{sec:transition}

Let
\begin{equation}\label{eq:chi}
 \chi(s):=2^s\pi^{s-1}\sin\!\Bigl(\frac{\pi s}2\Bigr)\Gamma(1-s)=\pi^{s-1/2}\frac{\Gamma((1-s)/2)}{\Gamma(s/2)},
\end{equation}
the factor of the functional equation $\zeta(s)=\chi(s)\zeta(1-s)$ \cite[\S2.1]{Titchmarsh1986}.  Write
\begin{equation}\label{eq:transition-zw}
 s=\tfrac12+z,\quad z=\sigma'+it,\qquad
 \tau_*:=-iz=t-i\sigma',\qquad
 w(s):=\exp\bigl\{-2\pi i\,\NR(\tau_*)\bigr\}.
\end{equation}
Stirling's formula gives, uniformly on compact $\sigma$-sets,
\begin{equation}\label{eq:chi-size}
 |\chi(\sigma+it)|=\Bigl(\frac{|t|}{2\pi}\Bigr)^{1/2-\sigma}\bigl(1+O(|t|^{-1})\bigr).
\end{equation}

\begin{theorem}[Transition evaluation of the midpoint product]\label{thm:transition}
There is a constant $\varpi_{\RVM}>0$ with the following properties.
\begin{enumerate}[label=\textup{(\alph*)}]
\item For every $B>0$ there are $t_0(B)$ and $C_B$ such that, for $0<\sigma'\le B$ and $t\ge t_0(B)$,
\begin{equation}\label{eq:transition-w}
 \zeta_{\RVM}(s)=\varpi_{\RVM}\bigl(1+w(s)\bigr)\bigl(1+\mathcal E_0(s)\bigr),\qquad|\mathcal E_0(s)|\le\frac{C_B}t .
\end{equation}
\item For every $0<\sigma'_1<\sigma'_2<\tfrac12$ there are $t_0$ and $C$ such that, for $\sigma'_1\le\sigma'\le\sigma'_2$ and $t\ge t_0$,
\begin{equation}\label{eq:transition-chi}
 \zeta_{\RVM}(s)=\varpi_{\RVM}\bigl(1+\chi(s)\bigr)\bigl(1+\mathcal E(s)\bigr),\qquad|\mathcal E(s)|\le\frac Ct .
\end{equation}
\end{enumerate}
For $t<0$ the same statements hold by $\zeta_{\RVM}(\bar s)=\overline{\zeta_{\RVM}(s)}$.
\end{theorem}

The phase identity $w=\chi(s)(1+O(1/t))$ alone does not give \CGeqref{eq:transition-w}; the product has to be evaluated across the transition region $\tau_k\asymp t$.  $\Log$ denotes the principal logarithm and $\Log\Gamma$ the branch of $\log\Gamma$ on $\Re v>0$ that is real on $(0,\infty)$.  Nothing about $\zeta$ is used.

\subsection{Proof of Theorem~\ref{thm:transition}}
\label{subsec:transition-proof}

Throughout, $t>0$, $0<\sigma'\le B$, $\tau_0:=4\pi$, $t^2\ge2(\tau_0^2+B^2)$, and implied constants depend only on $B$ (and on $\sigma'_1,\sigma'_2$ in part (b)).

\subsubsection*{Step 0: reduction to a sum over the clock}
Since $q(s)=z^2-\tfrac14$, each factor of \CGeqref{eq:Pcrit} equals $(\tau_k^2+z^2)/(\tau_k^2+\tfrac14)$, so
\begin{equation}\label{eq:transition-product}
 \frac{P_{\RVM}^{\crit}(s)}{P_{\RVM}^{\crit}(\tfrac12)}=\prod_{k\ge1}\Bigl(1+\frac{z^2}{\tau_k^2}\Bigr)=\exp\Lambda(z),\qquad
 \Lambda(z):=\sum_{k\ge1}f(\tau_k),\quad f(\tau):=\Log\Bigl(1+\frac{z^2}{\tau^2}\Bigr).
\end{equation}
For $\tau>0$, $\Im(1+z^2/\tau^2)=2\sigma't/\tau^2>0$, so $f$ is smooth on $(0,\infty)$, $f(\tau)=O(|z|^2\tau^{-2})$, the series converges absolutely, and
\begin{equation}\label{eq:fprime}
 f'(\tau)=-\frac{2z^2}{\tau(\tau^2+z^2)}=-\frac2\tau+\frac1{\tau-iz}+\frac1{\tau+iz}.
\end{equation}

\subsubsection*{Step 1: the midpoint sawtooth}
The clock increases on $[2\pi,\infty)$, $\NR(\tau_0)=2\log\frac2e+\frac78\in(0,\tfrac12)$, and $\tau_1>\tau_0$.  With $\psi(y):=y-\lfloor y\rfloor-\tfrac12$, the count \CGeqref{eq:n-RVM} satisfies, for $\tau\ge\tau_0$,
\begin{equation}\label{eq:sawtooth-count}
 n_{\RVM}(\tau)=\bigl\lfloor\NR(\tau)+\tfrac12\bigr\rfloor=\NR(\tau)-\psi\bigl(\NR(\tau)+\tfrac12\bigr),\qquad n_{\RVM}(\tau_0)=0 .
\end{equation}
Stieltjes integration of $f$ against \CGeqref{eq:sawtooth-count}, followed by integration by parts in the $\psi$-part ($f(\infty)=0$, $|\psi|\le\tfrac12$, $\psi(\NR(\tau_0)+\tfrac12)=\NR(\tau_0)$), gives
\begin{equation}\label{eq:transition-split}
\begin{split}
 &\Lambda(z)=\Lambda_{\rm sm}(z)+\Lambda_{\rm saw}(z),\qquad
 \Lambda_{\rm sm}(z):=\int_{\tau_0}^\infty f\,d\NR+f(\tau_0)\NR(\tau_0),\\
 &\Lambda_{\rm saw}(z):=\int_{\tau_0}^{\infty}\psi\bigl(\NR(\tau)+\tfrac12\bigr)f'(\tau)\,d\tau .
\end{split}
\end{equation}

\begin{lemma}[Smooth part]\label{lem:TR-smooth}
With $\log G(s):=\Log(z^2-\tfrac14)+\Log\Gamma(\tfrac s2)-\tfrac s2\log\pi$,
\begin{equation}\label{eq:smooth-difference}
 \Lambda_{\rm sm}(z)-\log G(s)=\nu_3+O(t^{-1}),\qquad\nu_3\in\R\ \text{independent of }z .
\end{equation}
\end{lemma}

\begin{proof}
A second integration by parts ($f\NR\to0$ at $\infty$) cancels the boundary term:
\begin{equation}\label{eq:smooth-J-ibp}
 \Lambda_{\rm sm}(z)=-\int_{\tau_0}^\infty\NR f'\,d\tau=\int_{\tau_0}^\infty\NR(\tau)\Bigl(\frac2\tau-\frac{2\tau}{\tau^2+z^2}\Bigr)d\tau .
\end{equation}
The constant $\tfrac78$ of the clock contributes $\tfrac78\Log(1+z^2/\tau_0^2)$, because $\Log\frac{\tau^2}{\tau^2+z^2}$ is an antiderivative of the bracket on $(0,\infty)$ tending to $0$ at $\infty$.  As $0<\arg z<\frac\pi2$, $\arg(1+\tau_0^2/z^2)=O(t^{-2})$ and
\begin{equation}\label{eq:J2}
 \tfrac78\Log\Bigl(1+\frac{z^2}{\tau_0^2}\Bigr)=\tfrac74\Log z-\tfrac74\log\tau_0+O(t^{-2}).
\end{equation}
The term $\frac\tau{2\pi}\log\frac\tau{2\pi e}$ contributes $\frac1\pi\int_{\tau_0}^\infty\log\frac{\tau}{2\pi e}\frac{z^2}{\tau^2+z^2}\,d\tau$.  Put
\begin{equation}\label{eq:I-def}
 \mathcal I(z):=\frac1\pi\int_0^\infty\log\frac{\tau}{2\pi e}\,\frac{z^2}{\tau^2+z^2}\,d\tau .
\end{equation}
The integrand is holomorphic in $z$ on $\Re z>0$ (its singularities are at $z=\pm i\tau$), $O(|z|^2\tau^{-2}\log\tau)$ at $\infty$ and $O(|\log\tau|)$ at $0$, locally uniformly; so $\mathcal I$ is holomorphic on $\Re z>0$.  For real $z=x>0$ the substitution $\tau=xv$ and $\int_0^\infty\frac{dv}{1+v^2}=\frac\pi2$, $\int_0^\infty\frac{\log v}{1+v^2}\,dv=0$ give $\mathcal I(x)=\frac x2\log\frac x{2\pi e}$; by the identity theorem $\mathcal I(z)=\frac z2\Log\frac z{2\pi e}$ on $\Re z>0$.  On $[0,\tau_0]$, $|\tau^2+z^2|\ge t^2/2$, so $\frac{z^2}{\tau^2+z^2}=1+O(t^{-2})$ there.  Hence
\begin{equation}\label{eq:J-asymp}
 \Lambda_{\rm sm}(z)=\frac z2\Log\frac z{2\pi e}+\frac74\Log z+\nu_1+O(t^{-2}),\qquad
 \nu_1=-\tfrac74\log\tau_0-\tfrac1\pi\int_0^{\tau_0}\log\tfrac\tau{2\pi e}\,d\tau .
\end{equation}
Stirling's formula $\Log\Gamma(v)=(v-\tfrac12)\Log v-v+\tfrac12\log2\pi+O(|v|^{-1})$ \cite[\S5.11]{DLMF2024} with $v=\frac s2=\frac{z+1/2}2$ and $\Log v=\Log z-\log2+\frac1{2z}+O(t^{-2})$ gives
\begin{equation}\label{eq:stirling-loggamma}
 \Log\Gamma(\tfrac s2)=\tfrac z2\Log\tfrac z2-\tfrac14\Log z+\tfrac14\log2-\tfrac z2+\tfrac12\log2\pi+O(t^{-1}),
\end{equation}
and with $\Log(z^2-\tfrac14)=2\Log z+O(t^{-2})$, $-\frac s2\log\pi=-\frac z2\log\pi-\frac14\log\pi$,
\begin{equation}\label{eq:G-asymp}
 \log G(s)=\frac z2\Log\frac z{2\pi e}+\frac74\Log z+\nu_2+O(t^{-1}),\qquad
 \nu_2=\tfrac14\log2+\tfrac12\log2\pi-\tfrac14\log\pi .
\end{equation}
Subtracting gives \CGeqref{eq:smooth-difference} with $\nu_3=\nu_1-\nu_2$.
\end{proof}

\begin{lemma}[The clock as a holomorphic phase]\label{lem:TR-clock}
$\NR(\tau)=\frac\tau{2\pi}\bigl(\Log\frac\tau{2\pi}-1\bigr)+\frac78$ is holomorphic on $\Re\tau>0$, $\NR(\bar\tau)=\overline{\NR(\tau)}$, and for $\Re\tau\ge\tau_0$,
\begin{equation}\label{eq:clock-halfplane-bound}
 2\pi\,\Im\NR(\tau)\ge(\log2)\,\Im\tau\quad(\Im\tau\ge0),\qquad
 2\pi\,\Im\NR(\tau)\le(\log2)\,\Im\tau\quad(\Im\tau\le0).
\end{equation}
\end{lemma}

\begin{proof}
For $\tau=X+iY$, $X\ge\tau_0$, $Y>0$, put $v=Y/X$.  Then
\begin{equation}\label{eq:clock-im-identity}
 2\pi\Im\NR(\tau)=Y\log\frac{|\tau|}{2\pi e}+X\arg\tau
 =Y\Bigl[\log\frac{X}{2\pi e}+\tfrac12\log(1+v^2)+\frac{\arctan v}v\Bigr]\ge Y\bigl[\log2+\varphi(v)\bigr],
\end{equation}
using $X\ge4\pi$, where $\varphi(v)=\frac12\log(1+v^2)+\frac{\arctan v}v-1$.  Since $\varphi(0^+)=0$ and $\varphi'(v)=(v-\arctan v)/v^2\ge0$, $\varphi\ge0$.  The case $Y<0$ follows by conjugation.
\end{proof}

In particular $|w(s)|=e^{2\pi\Im\NR(\tau_*)}\le2^{-\sigma'}<1$, and on the vertical rays $|e^{\pm2\pi im\NR(\tau_0\pm iy)}|\le2^{-my}$.

\begin{lemma}[Abel-regularized harmonic sum]\label{lem:TR-harmonic}
\begin{equation}\label{eq:Psi-asymp}
 \Lambda_{\rm saw}(z)=\Log\bigl(1+w(s)\bigr)+\nu_4+O(t^{-2}),\qquad\nu_4\in\R\ \text{independent of }z .
\end{equation}
\end{lemma}

\begin{proof}
\emph{Regularization.}  For $0<r<1$ let
\begin{equation}\label{eq:psi-r}
 \psi_r(y):=-\sum_{m\ge1}r^m\frac{\sin2\pi my}{\pi m}=\frac1\pi\Im\Log\bigl(1-re^{2\pi iy}\bigr).
\end{equation}
Since $\Re(1-re^{2\pi iy})>0$, $|\psi_r|\le\tfrac12$, and $\psi_r(y)\to\psi(y)$ as $r\uparrow1$ for $y\notin\Z$ (Abel summation of the convergent Fourier series of $\psi$).  The set $\{\tau:\NR(\tau)+\tfrac12\in\Z\}$ is discrete and $f'\in L^1(\tau_0,\infty)$, so by dominated convergence
\begin{equation}\label{eq:Psi-r-limit}
 \Lambda_{\rm saw}(z)=\lim_{r\uparrow1}\Lambda_r(z),\qquad
 \Lambda_r(z):=\int_{\tau_0}^\infty\psi_r\bigl(\NR+\tfrac12\bigr)f'\,d\tau .
\end{equation}
\emph{Termwise integration.}  For fixed $r<1$ the series for $\psi_r(\NR+\tfrac12)$, with $\sin2\pi m(y+\tfrac12)=(-1)^m\sin2\pi my$, is dominated by $\sum_mr^m/m$; by Fubini,
\begin{equation}\label{eq:Psi-r-series}
 \Lambda_r(z)=-\sum_{m\ge1}\frac{(-r)^m}{2\pi i\,m}\bigl(I_m^+-I_m^-\bigr),\qquad
 I_m^\pm:=\int_{\tau_0}^\infty e^{\pm2\pi im\NR(\tau)}f'(\tau)\,d\tau .
\end{equation}
\emph{Contours.}  Let $\mathcal H^\pm=\{\Re\tau\ge\tau_0,\ \pm\Im\tau\ge0\}$.  On $\mathcal H^\pm$, $|e^{\pm2\pi im\NR}|\le1$ by Lemma~\CGref{lem:TR-clock}, and $|f'(\tau)|\le4|z|^2|\tau|^{-3}$ for $|\tau|\ge2|z|$, so the quarter-circle arcs of radius $R$ contribute $O(|z|^2R^{-2})\to0$.  The poles of $f'$ are $0$, $iz=-t+i\sigma'\notin\mathcal H^\pm$, and $\tau_*=-iz=t-i\sigma'\in\mathcal H^-$, the last with residue $1$ (from $\frac1{\tau+iz}$).  Define the ray integrals
\begin{equation}\label{eq:ray-integrals}
 R_m^+:=i\int_0^\infty e^{2\pi im\NR(\tau_0+iy)}f'(\tau_0+iy)\,dy,\qquad
 R_m^-:=-i\int_0^\infty e^{-2\pi im\NR(\tau_0-iy)}f'(\tau_0-iy)\,dy .
\end{equation}
On $\mathcal H^+$ Cauchy's theorem gives $I_m^+=R_m^+$.  On $\mathcal H^-$ the positively oriented boundary of the quarter disc runs down the ray, along the arc, and back along $[\tau_0,\tau_0+R]$ from right to left; letting $R\to\infty$,
\begin{equation}\label{eq:lower-contour}
 R_m^--I_m^-=2\pi i\,e^{-2\pi im\NR(\tau_*)}=2\pi i\,w^m .
\end{equation}
Therefore, using $|rw|<1$,
\begin{equation}\label{eq:Psi-r-closed}
 \Lambda_r(z)=-\sum_{m\ge1}\frac{(-rw)^m}m+\Phi_r(z)=\Log(1+rw)+\Phi_r(z),\qquad
 \Phi_r(z):=-\sum_{m\ge1}\frac{(-r)^m}{2\pi i\,m}\bigl(R_m^+-R_m^-\bigr).
\end{equation}
\emph{The rays.}  Split $f'(\tau)=-\frac2\tau+g_z(\tau)$, $g_z(\tau)=\frac{2\tau}{\tau^2+z^2}$, and correspondingly $R_m^\pm=R_m^{\pm,(0)}+R_m^{\pm,(z)}$, $\Phi_r=\Phi_r^{(0)}+\Phi_r^{(z)}$.  The ray integrands carry the factor $2^{-my}$.  Since $|2/\tau|\le2/\tau_0$ on the rays, $|R^{\pm,(0)}_m|\le2/(\tau_0m\log2)$, so the series for $\Phi_r^{(0)}$ is dominated by $\sum_mC/m^2$ uniformly in $r$, and $\Phi_r^{(0)}\to\nu_4:=\Phi_1^{(0)}$, a constant depending only on $\tau_0$.  It is real: $\NR(\bar\tau)=\overline{\NR(\tau)}$ gives $R_m^{-,(0)}=\overline{R_m^{+,(0)}}$, so each term equals $-(-r)^m\Im R_m^{+,(0)}/(\pi m)$.  For $\tau=\tau_0\pm iy$, $\Re(\tau^2+z^2)=\tau_0^2+\sigma'^2-y^2-t^2\le-\frac12(y^2+t^2)$ because $t^2\ge2(\tau_0^2+B^2)$; hence $|g_z(\tau)|\le4(\tau_0+y)t^{-2}$ and
\begin{equation}\label{eq:ray-bound}
 |R^{\pm,(z)}_m|\le\frac4{t^2}\int_0^\infty2^{-my}(\tau_0+y)\,dy=\frac4{t^2}\Bigl(\frac{\tau_0}{m\log2}+\frac1{(m\log2)^2}\Bigr).
\end{equation}
So $|\Phi_r^{(z)}(z)|\le C\,t^{-2}\sum_mm^{-2}$, uniformly in $0<r<1$ and $0<\sigma'\le B$, and $\Phi_r^{(z)}$ converges as $r\uparrow1$ by dominated convergence of the series.  Since $\Log(1+rw)\to\Log(1+w)$, letting $r\uparrow1$ in \CGeqref{eq:Psi-r-closed} and \CGeqref{eq:Psi-r-limit} gives \CGeqref{eq:Psi-asymp}.  \emph{Uniformity.}  The three limits $r\uparrow1$ taken in this proof are dominated by $r$-independent majorants: $|\psi_r|\le\frac12$ against $|f'|\in L^1$ in \CGeqref{eq:Psi-r-limit}; $\sum_mr^m/m\le\sum_m1/m$ restricted to the finitely many $m$ relevant after the bounds $|R_m^{\pm,(0)}|\le2/(\tau_0m\log2)$ and \CGeqref{eq:ray-bound}, i.e.\ the majorants $C/m^2$ and $Ct^{-2}/m^2$; and $|rw|\le|w|<1$.  Hence the remainder $O(t^{-2})$ in \CGeqref{eq:Psi-asymp} is uniform in $0<\sigma'\le B$ and $t\ge t_0(B)$, and does not depend on the regularization.
\end{proof}

\subsubsection*{Step 2: assembly}
By \CGeqref{eq:transition-product}, \CGeqref{eq:transition-split} and Lemmas~\CGref{lem:TR-smooth} and~\CGref{lem:TR-harmonic},
\begin{equation}\label{eq:assembly}
\begin{split}
 \zeta_{\RVM}(s)&=\frac{P_{\RVM}^{\crit}(s)}{G(s)}=P_{\RVM}^{\crit}(\tfrac12)\exp\bigl(\Lambda_{\rm sm}(z)-\log G(s)+\Lambda_{\rm saw}(z)\bigr)\\
 &=P_{\RVM}^{\crit}(\tfrac12)\,e^{\nu_3+\nu_4}(1+w)\bigl(1+O(t^{-1})\bigr),
\end{split}
\end{equation}
uniformly for $0<\sigma'\le B$.  This is \CGeqref{eq:transition-w} with
\begin{equation}\label{eq:varpi-def}
 \varpi_{\RVM}:=P_{\RVM}^{\crit}(\tfrac12)\,e^{\nu_3+\nu_4}>0,
\end{equation}
which depends only on the clock; its numerical value $0.9619415807\ldots$ plays no role in the proof.

\subsubsection*{Step 3: from the clock phase to \texorpdfstring{$\chi$}{chi}}
Let
\begin{equation}\label{eq:theta-def}
 \vartheta(v):=\frac1{2i}\Bigl(\Log\Gamma\bigl(\tfrac14+\tfrac{iv}2\bigr)-\Log\Gamma\bigl(\tfrac14-\tfrac{iv}2\bigr)\Bigr)-\frac v2\log\pi,
\end{equation}
holomorphic on $|\Im v|<\tfrac12$.  The identity $\chi(\tfrac12+iv)=e^{-2i\vartheta(v)}$ holds for real $v$, hence on the strip, and $s=\tfrac12+i\tau_*$ with $|\Im\tau_*|=\sigma'<\tfrac12$.  Stirling's formula in the strip gives $\vartheta(v)=\frac v2\Log\frac v{2\pi}-\frac v2-\frac\pi8+O(|v|^{-1})$, i.e.
\begin{equation}\label{eq:clock-theta}
 \NR(v)=\frac{\vartheta(v)}\pi+1+O(|v|^{-1}).
\end{equation}
Therefore
\begin{equation}\label{eq:w-chi}
 w=e^{-2i\vartheta(\tau_*)-2\pi i+O(1/t)}=\chi(s)\bigl(1+O(t^{-1})\bigr).
\end{equation}
For $\sigma'\ge\sigma'_1$, \CGeqref{eq:chi-size} gives $|\chi(s)|\le\tfrac12$ when $t\ge t_0(\sigma'_1)$, so $1+w=(1+\chi)\bigl(1+(w-\chi)/(1+\chi)\bigr)=(1+\chi)(1+O(t^{-1}))$, which with \CGeqref{eq:transition-w} gives \CGeqref{eq:transition-chi}.  This completes the proof of Theorem~\CGref{thm:transition}. \qed

\begin{hypothesis}[Strict inequality for the transition constant]\label{hyp:varpi-less-one}
$\varpi_{\RVM}<1$.
\end{hypothesis}

Hypothesis~\CGref{hyp:varpi-less-one} is supported by the Richardson value $0.9619415807\ldots$ of Section~\CGref{sec:numerics}, but no rigorous enclosure of $\varpi_{\RVM}$ is given in this paper: the constants $\nu_3$ and $\nu_4$ in \CGeqref{eq:varpi-def} are defined by convergent integrals whose certified evaluation has not been carried out. By Theorem~\CGref{thm:Z1-clock} the hypothesis is equivalent to $\lim_{K\to\infty}\prod_{k\le K}\frac{\frac14+\gamma_k^2}{\frac14+\tau_k^2}<1/\varpi_{Z1}$, an ordered product whose partial product at $K=2\,001\,052$ is $0.9619415390$ (Experiment~\CGref{exp:product-identity}) and whose tail is not controlled rigorously. The inequality is used in exactly five places, each of which says so: Proposition~\CGref{prop:atomized-clock}(ii), the sign $\mathsf W(1)>0$ in Proposition~\CGref{prop:one-number}, Proposition~\CGref{prop:open-path-varpi}, the failure of the first rung at the clock end in Corollary~\CGref{cor:ladder-pair}, and the remark after Corollary~\CGref{cor:positivity}. No other statement of the paper depends on it; in particular Theorem~\CGref{thm:transition} and every unconditional statement below are proved for the number $\varpi_{\RVM}$ whatever its value.  The ten digits quoted are a Richardson extrapolation (Experiment~\CGref{exp:transition}) and carry no error bound; the value is not certified and its agreement with the ordered product of Experiment~\CGref{exp:product-identity} is a consistency check between two floating-point computations, not a proof.

\emph{What a proof of the hypothesis would require.}  By \CGeqref{eq:varpi-def}, $\varpi_{\RVM}=P^{\crit}_{\RVM}(\frac12)e^{\nu_3+\nu_4}$ is a product of three numbers defined by the clock alone: $P^{\crit}_{\RVM}(\frac12)=\prod_k\tau_k^2/(\frac14+\tau_k^2)$, whose partial product over $k\le2\,001\,052$ is $0.9940\ldots$ and whose tail is $\exp(-\frac14\sum_{k>N}\tau_k^{-2}+O(\sum_{k>N}\tau_k^{-4}))$ with $\sum_{k>N}\tau_k^{-2}$ computable to any accuracy from the Lambert formula; and $\nu_3$, $\nu_4$, the smooth and the sawtooth constants of Lemmas~\CGref{lem:TR-smooth} and~\CGref{lem:TR-harmonic}, defined by absolutely convergent integrals along the rays $\tau_0\pm iy$ whose integrands are explicit (Stirling remainders and $2^{-my}$-damped clock phases).  Numerically $e^{\nu_3+\nu_4}\approx0.9676$, so that the margin to be certified is $1-\varpi_{\RVM}\approx0.038$, about four per cent, far above the accuracy that interval arithmetic reaches for such integrals.  A proof would therefore consist of (a) a rigorous enclosure of the finite product and of its tail, (b) rigorous enclosures of $\nu_3$ and $\nu_4$ by interval quadrature with explicit remainders for the truncation of the rays and of the sum over $m$, and (c) the verification, in interval arithmetic, that the product of the three enclosures lies below $1$.  None of (a)--(c) is done in this paper, and no analytic argument giving $\varpi_{\RVM}<1$ from the structure of the clock is known to the author: the sign of $\log\varpi_{\RVM}$ is a cancellation between a negative $\log P^{\crit}_{\RVM}(\frac12)$ and the constants $\nu_3+\nu_4<0$, and nothing in Sections~\CGref{sec:zeta-rvm}--\CGref{sec:transition} forces its sign.  The statement is kept as a hypothesis for this reason.

\begin{lemma}[Extension to the critical line]\label{lem:TR-boundary}
The function $\mathcal E_0=\zeta_{\RVM}/(\varpi_{\RVM}(1+w))-1$ extends holomorphically to $\{|\sigma'|<\tfrac12,\ t>t_0\}$, where $t_0=t_0(\tfrac14)$ is the threshold of Theorem~\CGref{thm:transition}(a) with $B=\tfrac14$, and $|\mathcal E_0|\le C/t$ for $0\le\sigma'\le\tfrac14$.
\end{lemma}

\begin{proof}
By Lemma~\CGref{lem:TR-clock}, $\Im\NR(\tau_*)\ne0$ for $\sigma'\ne0$, so $1+w$ vanishes in the strip only on $\sigma'=0$, exactly at $t=\tau_k$; these zeros are simple because $\frac{d}{dt}w=-i\log(\tau_*/2\pi)\,w\ne0$.  The zeros of $\zeta_{\RVM}$ in the strip are the same simple points.  Hence $\mathcal E_0$ has only removable singularities, and the bound extends to $\sigma'=0$ by continuity.
\end{proof}

\begin{corollary}[Derivative at the zeros]\label{cor:derivative-zeros}
For $k$ with $\tau_k>t_0$,
\begin{equation}\label{eq:derivative-zeros}
 \zeta_{\RVM}'(\rho_k^{\RVM})=\varpi_{\RVM}\log\frac{\tau_k}{2\pi}\Bigl(1+O\Bigl(\frac1{\tau_k}\Bigr)\Bigr).
\end{equation}
\end{corollary}

\begin{proof}
By Lemma~\CGref{lem:TR-boundary}, $\zeta_{\RVM}=\varpi_{\RVM}(1+w)(1+\mathcal E_0)$ holds holomorphically near $\sigma'=0$.  At $s=\rho_k^{\RVM}$, $1+w=0$, so $\zeta_{\RVM}'(\rho_k^{\RVM})=\varpi_{\RVM}w'(\rho_k^{\RVM})(1+\mathcal E_0(\rho_k^{\RVM}))$.  From \CGeqref{eq:transition-zw}, $w'(s)=-2\pi i\,\NR'(\tau_*)\cdot(-i)\,w(s)=-\log(\tau_*/2\pi)\,w(s)$, and at the zero $w=-1$, $\tau_*=\tau_k$.  Finally $|\mathcal E_0|\le C/\tau_k$.
\end{proof}

\begin{remark}[Structure of the proof]\label{rem:TR-audit}
The delicate points are handled as follows: the smooth part is computed with principal branches by the identity theorem (Lemma~\CGref{lem:TR-smooth}); the clock estimate is quantified (Lemma~\CGref{lem:TR-clock}); the sawtooth series is replaced by its Abel means, termwise integration is justified for $r<1$, and $r\uparrow1$ is taken after the contour argument uniformly in $r$ and $z$ (Lemma~\CGref{lem:TR-harmonic}); the orientation of the lower contour and the sign of the residue are written out in \CGeqref{eq:lower-contour}; the $z$-dependence of the rays is isolated in $g_z$ and is $O(t^{-2})$, the rest is the explicit real constant $\nu_4$; Stirling is used uniformly for $0<\sigma'\le B$, and the only place where $\sigma'<\tfrac12$ enters is Step~3.
\end{remark}

\begin{corollary}[Vertical constant limit]\label{cor:vertical-limit}
For every fixed $\sigma>\tfrac12$,
\begin{equation}\label{eq:vertical-limit}
 \zeta_{\RVM}(\sigma+it)\longrightarrow\varpi_{\RVM},\qquad
 F_{\RVM}(\sigma+it)\longrightarrow L:=\frac1{\varpi_{\RVM}}\qquad(|t|\to\infty).
\end{equation}
More precisely, with $\sigma'=\sigma-\tfrac12$,
\begin{equation}\label{eq:w-decay}
 |w(\sigma+it)|\le (|t|/2\pi)^{-\sigma'}(1+O_B(|t|^{-1}))\qquad(0<\sigma'\le B),
\end{equation}
and, with $\alpha_c:=\min(c-\tfrac12,1)$ for $c>\tfrac12$,
\begin{equation}\label{eq:FminusL-all}
 F_{\RVM}(c+it)-L=O_c(|t|^{-\alpha_c}).
\end{equation}
\end{corollary}

\begin{proof}
Expanding the clock at $\tau_*=t-i\sigma'$ gives $\Im\NR(\tau_*)=-\frac{\sigma'}{2\pi}\log\frac t{2\pi}+O_\sigma(t^{-1})$, hence $|w|=e^{2\pi\Im\NR(\tau_*)}=O_\sigma(t^{-\sigma'})$.  Inverting \CGeqref{eq:transition-w} for large $t$ gives $F_{\RVM}=L(1+O(|w|)+O(t^{-1}))$.
\end{proof}

\subsection{Bounds on the whole half-plane}
\label{subsec:halfplane}

The constants of Theorem~\CGref{thm:transition} depend on the width $B$ of the band.  The following estimate holds on the whole half-plane and does not use contours.  Put
\begin{equation}\label{eq:saw-primitive}
 \mathfrak s_1(\tau):=\int_{\tau_0}^{\tau}\psi\bigl(\NR(v)+\tfrac12\bigr)\,dv\qquad(\tau\ge\tau_0).
\end{equation}

\begin{lemma}[Bounded primitive of the midpoint sawtooth]\label{lem:saw-primitive}
$|\mathfrak s_1(\tau)|\le\dfrac{\pi}{4\log2}$ for all $\tau\ge\tau_0$.
\end{lemma}

\begin{proof}
Substitute $y=\NR(v)$, $dv=\mathcal T'(y)\,dy$ with $\mathcal T'(y)=2\pi/\log(\mathcal T(y)/2\pi)$ from \CGeqref{eq:inverse-clock}.  The weight $\mathcal T'$ is positive and decreasing, with $\mathcal T'\le2\pi/\log2$ on $y\ge\NR(\tau_0)$ because $\log(\tau_0/2\pi)=\log2$.  The primitive of $y\mapsto\psi(y+\tfrac12)$ is periodic with values in $[-\tfrac18,0]$, so every integral of $\psi(y+\frac12)$ over an interval has modulus at most $\tfrac18$.  The second mean-value theorem gives $|\mathfrak s_1(\tau)|\le\frac{2\pi}{\log2}\cdot\frac18$.
\end{proof}

\begin{theorem}[Two-sided bound and limit at $+\infty$]\label{thm:halfplane-RVM}
There is an absolute constant $R_0$ such that, for $s=\tfrac12+z$, $z=\sigma'+it$, $\sigma'>\tfrac12$ and $|z|\ge R_0$,
\begin{equation}\label{eq:halfplane-RVM}
 \zeta_{\RVM}(s)=\varpi_{\RVM}\,e^{\theta_{\RVM}(s)},\qquad |\theta_{\RVM}(s)|\le\frac{\pi^2}{2\log2}\cdot\frac1{\sigma'}+\frac{C}{|z|}.
\end{equation}
Consequently, for every $\sigma_0>1$ there is $C(\sigma_0)$ with
\begin{equation}\label{eq:halfplane-two-sided}
 C(\sigma_0)^{-1}\le|\zeta_{\RVM}(s)|\le C(\sigma_0)\qquad(\Re s\ge\sigma_0),
\end{equation}
and $\zeta_{\RVM}(s)\to\varpi_{\RVM}$ uniformly in $t$ as $\Re s\to+\infty$.  Moreover
\begin{equation}\label{eq:nu4-closed}
 \nu_4=-2\int_{\tau_0}^\infty\frac{\mathfrak s_1(\tau)}{\tau^2}\,d\tau .
\end{equation}
\end{theorem}

\begin{proof}
\emph{Smooth part.}  For $\Re z>0$ the computation in the proof of Lemma~\CGref{lem:TR-smooth} is exact:
\begin{equation}\label{eq:smooth-exact}
 \Lambda_{\rm sm}(z)=\frac z2\Log\frac z{2\pi e}+\frac78\Log\Bigl(1+\frac{z^2}{\tau_0^2}\Bigr)-\frac1\pi\int_0^{\tau_0}\log\frac\tau{2\pi e}\,\frac{z^2}{\tau^2+z^2}\,d\tau .
\end{equation}
For $|z|\ge2\tau_0$ one has $|\tau^2+z^2|\ge\frac34|z|^2$ on $[0,\tau_0]$, and Stirling's formula is uniform in $|\arg(s/2)|<\frac\pi2$; the expansions \CGeqref{eq:J-asymp} and \CGeqref{eq:G-asymp} were derived for $t>0$, but every step (the branches of $\mathrm{Log}(z^2)$, of $\mathrm{Log}(1+z^2/\tau_0^2)$, of $\mathrm{Log}(z^2-\frac14)$ and of $f(\tau)=\mathrm{Log}(1+z^2/\tau^2)$ on $[\tau_0,\infty)$) only uses $|\arg z|<\frac\pi2$, which holds on $\Re z>0$ for either sign of $t$; they therefore hold with $O(|z|^{-1})$ in place of $O(t^{-1})$, uniformly on $\Re z>\tfrac12$, and $\Lambda_{\rm sm}(z)-\log G(s)=\nu_3+O(|z|^{-1})$.

\emph{Sawtooth part.}  Integrating \CGeqref{eq:transition-split} by parts with \CGeqref{eq:saw-primitive} ($\mathfrak s_1(\tau_0)=0$, $f'(\tau)=O(|z|^2\tau^{-3})$),
\begin{equation}\label{eq:saw-ibp}
 \Lambda_{\rm saw}(z)=-\int_{\tau_0}^\infty\mathfrak s_1f''\,d\tau,\qquad f''(\tau)=\frac2{\tau^2}-\frac1{(\tau-iz)^2}-\frac1{(\tau+iz)^2}.
\end{equation}
Since $|\tau\mp iz|^2=(\tau\pm t)^2+\sigma'^2$ and $\int_{\R}\frac{d\tau}{(\tau-a)^2+\sigma'^2}=\frac\pi{\sigma'}$, Lemma~\CGref{lem:saw-primitive} gives
\begin{equation}\label{eq:saw-bound}
 \Bigl|\Lambda_{\rm saw}(z)+2\int_{\tau_0}^\infty\frac{\mathfrak s_1(\tau)}{\tau^2}\,d\tau\Bigr|\le\frac{\pi}{4\log2}\cdot\frac{2\pi}{\sigma'} .
\end{equation}

\emph{Assembly.}  By \CGeqref{eq:assembly}, $\log\zeta_{\RVM}(s)=\log\varpi_\infty+\theta_{\RVM}(s)$ with $\varpi_\infty:=P_{\RVM}^{\crit}(\tfrac12)\exp(\nu_3-2\int_{\tau_0}^\infty\mathfrak s_1\tau^{-2}d\tau)>0$ and $\theta_{\RVM}$ bounded as in \CGeqref{eq:halfplane-RVM}.  For fixed $\sigma'$ and $t\to\infty$, Corollary~\CGref{cor:vertical-limit} gives $|\log\varpi_{\RVM}-\log\varpi_\infty|\le\pi^2/(2\sigma'\log2)$; a priori the logarithms can differ by $2\pi i k$.  For sufficiently large $\sigma'$, the bound in \CGeqref{eq:halfplane-RVM} is $<2\pi$, forcing $k=0$; letting $\sigma'\to\infty$ then gives $\varpi_\infty=\varpi_{\RVM}$, which is \CGeqref{eq:nu4-closed} by \CGeqref{eq:varpi-def}.  On the compact set $\{\Re s\ge\sigma_0,\ |z|\le R_0\}$ the function $\zeta_{\RVM}$ is continuous and has no zeros or poles, which gives \CGeqref{eq:halfplane-two-sided}.
\end{proof}

Theorem~\CGref{thm:halfplane-RVM} extends the two-sided bound \CGeqref{eq:zetaRVM-two-sided} of Corollary~\CGref{cor:Q-control} from bands to the whole half-plane $\Re s\ge\sigma_0>1$, and identifies $\varpi_{\RVM}$ also as the limit of $\zeta_{\RVM}(\sigma)$ as $\sigma\to+\infty$.

\begin{theorem}[No absolutely convergent Dirichlet expansion]\label{thm:no-Dirichlet-RVM}
Neither $\zeta_{\RVM}$ nor $F_{\RVM}=1/\zeta_{\RVM}$ is represented by a Dirichlet series $\sum_na_nn^{-s}$ converging absolutely in some half-plane $\Re s>\sigma_0$.
\end{theorem}

\begin{proof}
Suppose $F_{\RVM}(s)=\sum_{n\ge1}a_nn^{-s}$ absolutely for $\Re s>\sigma_0$.  Fix $\sigma>\max(\sigma_0,\tfrac12)$ and put $\tilde a_1=a_1-L$, $\tilde a_n=a_n$ ($n\ge2$), so $F_{\RVM}(\sigma+it)-L=\sum_n\tilde a_nn^{-\sigma-it}$ with $\sum_n|\tilde a_n|n^{-\sigma}<\infty$.  Then
\begin{equation}\label{eq:mean-square}
 \frac1{2T}\int_{-T}^{T}\bigl|F_{\RVM}(\sigma+it)-L\bigr|^2dt
 =\sum_{m,n}\tilde a_m\overline{\tilde a_n}(mn)^{-\sigma}\,\frac1{2T}\int_{-T}^{T}\Bigl(\frac nm\Bigr)^{it}dt ,
\end{equation}
the interchange being justified by absolute convergence of the double series, uniformly in $t$.  The inner average is $1$ for $m=n$ and is bounded by $\min\bigl(1,1/(T|\log(n/m)|)\bigr)\to0$ for $m\ne n$; each term is bounded by $|\tilde a_m\tilde a_n|(mn)^{-\sigma}$, which is summable, so dominated convergence gives
\begin{equation}\label{eq:mean-square-limit}
 \lim_{T\to\infty}\frac1{2T}\int_{-T}^{T}\bigl|F_{\RVM}(\sigma+it)-L\bigr|^2dt=\sum_{n\ge1}|\tilde a_n|^2n^{-2\sigma}.
\end{equation}
By \CGeqref{eq:vertical-limit} the integrand tends to $0$ and is bounded on the line (continuous, with a finite limit), so the left side of \CGeqref{eq:mean-square-limit} is $0$.  Hence every $\tilde a_n=0$, $F_{\RVM}\equiv L$ in the half-plane and, by analytic continuation, everywhere, which contradicts $F_{\RVM}(1)=0$.  The same argument applied to $\zeta_{\RVM}-\varpi_{\RVM}$ forces an absolutely convergent Dirichlet series for $\zeta_{\RVM}$ to be the constant $\varpi_{\RVM}$, which contradicts its pole at $s=1$.
\end{proof}

\begin{theorem}[Unsmoothed synthetic Mertens bound]\label{thm:unsmoothed-RVM}
For every $x>1$ and every $c>\tfrac12$ the symmetric Perron integral
\begin{equation}\label{eq:Mcrit-symmetric}
 M_{\RVM}^{\crit}(x):=\lim_{U\to\infty}\frac1{2\pi i}\int_{c-iU}^{c+iU}F_{\RVM}(s)\frac{x^s}{s}\,ds
\end{equation}
exists, is independent of $c>\tfrac12$, and equals the absolutely convergent expression
\begin{equation}\label{eq:Mcrit-renorm}
 M_{\RVM}^{\crit}(x)=L+\frac1{2\pi i}\int_{(c)}\bigl(F_{\RVM}(s)-L\bigr)\frac{x^s}{s}\,ds .
\end{equation}
For every $\eps\in(0,\tfrac12)$,
\begin{equation}\label{eq:Mcrit-unsmoothed-bound}
 M_{\RVM}^{\crit}(x)=L+O_\eps(x^{1/2+\eps}).
\end{equation}
\end{theorem}

\begin{proof}
Fix $c>\tfrac12$.  By \CGeqref{eq:FminusL-all} and continuity of $F_{\RVM}$ on $\Re s=c$ (no poles there by Proposition~\CGref{prop:Mcrit-poles}), the integral in \CGeqref{eq:Mcrit-renorm} converges absolutely.  The Perron kernel gives, for $x>1$ and $c>0$ \cite[Ch.~17]{Davenport1980},
\begin{equation}\label{eq:bromwich}
 \lim_{U\to\infty}\frac1{2\pi i}\int_{c-iU}^{c+iU}\frac{x^s}{s}\,ds=1 .
\end{equation}
Splitting $F_{\RVM}=L+(F_{\RVM}-L)$ proves existence of \CGeqref{eq:Mcrit-symmetric} and \CGeqref{eq:Mcrit-renorm}.  If $\tfrac12<c<c'$, the integrand of \CGeqref{eq:Mcrit-symmetric} is holomorphic on the closed strip $c\le\Re s\le c'$ (Proposition~\CGref{prop:Mcrit-poles}); by Theorem~\CGref{thm:transition}(a) with $B=c'-\tfrac12$, $F_{\RVM}$ is bounded on the horizontal segments $\Im s=\pm U$, $U\ge t_0(B)$, so these contribute $O(x^{c'}U^{-1})\to0$, and Cauchy's theorem gives independence of $c$.  Take $c=\tfrac12+\eps$ in \CGeqref{eq:Mcrit-renorm}:
\begin{equation}\label{eq:Mcrit-final-estimate}
 |M_{\RVM}^{\crit}(x)-L|\le\frac{x^{1/2+\eps}}{2\pi}\int_{-\infty}^{\infty}\frac{|F_{\RVM}(\tfrac12+\eps+it)-L|}{|\tfrac12+\eps+it|}\,dt,
\end{equation}
and the integral is finite by \CGeqref{eq:FminusL-all}, the integrand being $O(|t|^{-1-\alpha_c})$.
\end{proof}

\begin{remark}[Scope]\label{rem:unsmoothed-scope}
Theorem~\CGref{thm:unsmoothed-RVM} concerns $1/\zeta_{\RVM}$.  The transfer to $M$ goes through the endpoint of the $\lambda$-family and is Theorem~\CGref{thm:step3}.
\end{remark}

\section{The exact Hadamard product and the comparison with \texorpdfstring{$\zeta$}{zeta}}
\label{sec:comparison}

Put
\begin{equation}\label{eq:Pzeta}
 P_\zeta(s):=G(s)\zeta(s)=2\xi(s),\qquad \xi(s)=\tfrac12s(s-1)\pi^{-s/2}\Gamma(s/2)\zeta(s).
\end{equation}

\begin{theorem}[Exact Hadamard product]\label{thm:Hadamard}
Let $\rho$ run over the nontrivial zeros of $\zeta$ with $\Im\rho>0$, with multiplicity, ordered by $\Im\rho$.  The product
\begin{equation}\label{eq:Hadamard-exact}
 P_\zeta(s)=\prod_{\Im\rho>0}\left(1+\frac{q(s)}{\rho(1-\rho)}\right)
\end{equation}
converges absolutely and locally uniformly in $s\in\C$, and the identity holds with no exponential or constant factor.
\end{theorem}

\begin{proof}
$\xi$ is entire of order one with $\xi(0)=\tfrac12$, so Hadamard's factorization \cite{Hadamard1896}, \cite[\S2.12]{Titchmarsh1986}, \cite[Ch.~12]{Davenport1980} gives
\begin{equation}\label{eq:hadamard-general}
 2\xi(s)=e^{Bs}\prod_{\rho}\Bigl(1-\frac s\rho\Bigr)e^{s/\rho},
\end{equation}
the product running over all nontrivial zeros and converging absolutely because $\sum_\rho|\rho|^{-2}<\infty$.  The functional equation $\xi(s)=\xi(1-s)$ makes $\rho\mapsto1-\rho$ an involution of the zero multiset; no nontrivial zero is real, since $(1-2^{1-\sigma})\zeta(\sigma)=\sum_{n\ge1}(-1)^{n-1}n^{-\sigma}>0$ for $0<\sigma<1$; so this involution exchanges $\Im\rho>0$ and $\Im\rho<0$.  Grouping $\rho$ with $1-\rho$, which absolute convergence permits,
\begin{equation}\label{eq:pair-factor}
 \Bigl(1-\frac s\rho\Bigr)\Bigl(1-\frac s{1-\rho}\Bigr)e^{s/\rho+s/(1-\rho)}
 =\Bigl(1+\frac{q(s)}{\rho(1-\rho)}\Bigr)e^{s/(\rho(1-\rho))}.
\end{equation}
Since $|\rho(1-\rho)|\ge(\Im\rho)^2$ and $\sum_{\Im\rho>0}(\Im\rho)^{-2}<\infty$, the numbers $1/(\rho(1-\rho))$ are absolutely summable, so the exponentials collect into a single factor:
\begin{equation}\label{eq:hadamard-exp}
 2\xi(s)=e^{B's}\,\Pi_\zeta(s),\qquad \Pi_\zeta(s):=\prod_{\Im\rho>0}\Bigl(1+\frac{q(s)}{\rho(1-\rho)}\Bigr),\qquad B'=B+\sum_{\Im\rho>0}\frac1{\rho(1-\rho)} .
\end{equation}
$\Pi_\zeta$ depends on $s$ only through $q(s)$, hence $\Pi_\zeta(1-s)=\Pi_\zeta(s)$.  From $\xi(1-s)=\xi(s)$ we get $e^{B's}=e^{B'(1-s)}$ on the open set where $\Pi_\zeta\ne0$, so $B'=0$. No constant factor is left either: the general factorization \CGeqref{eq:hadamard-general} carries the factor $e^{A+Bs}$ with $e^{A}=2\xi(0)=1$, which is why it was written without $A$; equivalently, at $s=0$ one has $q(0)=0$, $\Pi_\zeta(0)=1$ and $2\xi(0)=1$, so \CGeqref{eq:hadamard-exp} with $B'=0$ is \CGeqref{eq:Hadamard-exact}.
\end{proof}

Define the exact quotient
\begin{equation}\label{eq:Q-RVM}
 \mathcal Q_{\RVM}(s):=\frac{\zeta(s)}{\zeta_{\RVM}(s)}=\frac{P_\zeta(s)}{P_{\RVM}^{\crit}(s)}.
\end{equation}

\begin{proposition}[Exact transform relation]\label{prop:M-transform-relation}
For $\Re s>1$, and by meromorphic continuation,
\begin{equation}\label{eq:M-transform-relation}
 F_{\RVM}(s)=\mathcal Q_{\RVM}(s)\,\frac1{\zeta(s)} .
\end{equation}
If the inverse Mellin transform $\mathcal Q_{\RVM}^{\vee}$ of $\mathcal Q_{\RVM}$ and the inverse defining $M_{\RVM}^{\crit}$ exist in a common distribution space with Mellin convolution, then
\begin{equation}\label{eq:M-convolution}
 M_{\RVM}^{\crit}=\mathcal Q_{\RVM}^{\vee}\Mell M .
\end{equation}
\end{proposition}

\CGeqref{eq:M-transform-relation} is the identity $1/\zeta_{\RVM}=(\zeta/\zeta_{\RVM})(1/\zeta)$.  Equation~\CGeqref{eq:M-convolution} is only notation for the Mellin-convolution theorem when both inverse transforms belong to a specified common convolution space; no such space is constructed here.

\begin{corollary}[All growth of $\zeta$ to the right of $\frac12$ lies in $\mathcal Q_{\RVM}$]\label{cor:Q-control}
For every $c'>\tfrac12$ and $0<\eps<c'-\tfrac12$ there are $t_0=t_0(\eps,c')$ and $0<c_-(\eps,c')\le c_+(\eps,c')$ such that
\begin{equation}\label{eq:zetaRVM-two-sided}
 c_-\le|\zeta_{\RVM}(\sigma+it)|\le c_+\qquad(\tfrac12+\eps\le\sigma\le c',\ |t|\ge t_0),
\end{equation}
and therefore, in the same region,
\begin{equation}\label{eq:zeta-asymp-Q}
 c_-|\mathcal Q_{\RVM}(s)|\le|\zeta(s)|\le c_+|\mathcal Q_{\RVM}(s)|,\qquad
 \frac1{c_+|\mathcal Q_{\RVM}(s)|}\le\Bigl|\frac1{\zeta(s)}\Bigr|\le\frac1{c_-|\mathcal Q_{\RVM}(s)|}.
\end{equation}
Consequently:
\begin{enumerate}[label=\textup{(\roman*)}]
\item the Lindel\"of hypothesis holds if and only if $\mathcal Q_{\RVM}(\sigma+it)=O_{\sigma,\eps}(|t|^{\eps})$ for every $\sigma>\tfrac12$;
\item every bound for $\zeta$ or $1/\zeta$ on $\tfrac12+\eps\le\sigma\le c'$, $|t|\ge t_0$, is equivalent, up to the constants $c_\pm$, to the same bound for $\mathcal Q_{\RVM}$ or $1/\mathcal Q_{\RVM}$.
\end{enumerate}
\end{corollary}

\begin{proof}
By Theorem~\CGref{thm:transition}(a) with $B=c'-\tfrac12$, $\zeta_{\RVM}=\varpi_{\RVM}(1+w)(1+\mathcal E_0)$ with $|\mathcal E_0|\le C_B/|t|$ and, by the uniform form \CGeqref{eq:w-decay}, $|w|\le C_B(|t|/2\pi)^{-\eps}$; for $|t|\ge t_0$ this gives $\tfrac12\le|1+w|\le2$ and $\tfrac12\le|1+\mathcal E_0|\le2$, hence \CGeqref{eq:zetaRVM-two-sided} with $c_-=\varpi_{\RVM}/4$, $c_+=4\varpi_{\RVM}$.  Then \CGeqref{eq:zeta-asymp-Q} follows from $\zeta=\mathcal Q_{\RVM}\zeta_{\RVM}$.  For (i): the Lindel\"of hypothesis is equivalent to $\zeta(\sigma+it)=O(|t|^\eps)$ for every $\sigma>\tfrac12$ \cite[Ch.~XIII]{Titchmarsh1986}, and \CGeqref{eq:zeta-asymp-Q} transfers this to $\mathcal Q_{\RVM}$.  (ii) is \CGeqref{eq:zeta-asymp-Q}.
\end{proof}

On the prime side define
\begin{equation}\label{eq:prime-defect}
 \mathfrak d_n^{\RVM}:=\mathbf 1_{\{n\ \mathrm{prime}\}}-c_n^{\RVM}.
\end{equation}
Since $\pi_{\RVM}(1)=0$, for every integer $N\ge2$
\begin{equation}\label{eq:prime-defect-cumulative}
 \sum_{2\le n\le N}\mathfrak d_n^{\RVM}=\pi(N)-\pi_{\RVM}(N).
\end{equation}

\section{The \texorpdfstring{$\lambda$}{lambda}-family from \texorpdfstring{$\zeta_{\RVM}$}{zeta\_RVM} to \texorpdfstring{$\zeta_{Z1}$}{zeta\_Z1}}
\label{sec:lambda-family}

For $\lambda\in[0,1]$ define the moving ordinates, the family product and the family zeta function
\begin{equation}\label{eq:gamma-lambda}
 \gamma_k(\lambda):=(1-\lambda)\gamma_k+\lambda\tau_k=\gamma_k+\lambda\eta_k,
\end{equation}
\begin{equation}\label{eq:P-lambda}
 P_\lambda(s):=\prod_{k\ge1}\left(1+\frac{q(s)}{\tfrac14+\gamma_k(\lambda)^2}\right),
\end{equation}
\begin{equation}\label{eq:zeta-lambda}
 \zeta_\lambda(s):=\frac{P_\lambda(s)}{G(s)},\qquad F_\lambda(s):=\frac1{\zeta_\lambda(s)},
\end{equation}
and the endpoint
\begin{equation}\label{eq:family-endpoints}
 P_{Z1}:=P_0,\qquad \zeta_{Z1}:=\zeta_0,\qquad \zeta_1=\zeta_{\RVM}.
\end{equation}

\begin{theorem}[The $\lambda$-family]\label{thm:path}
The product \CGeqref{eq:P-lambda} converges absolutely and locally uniformly jointly on $[0,1]\times\C$.  For every $\lambda$, $P_\lambda$ is entire of order exactly one, $P_\lambda(1-s)=P_\lambda(s)$, $P_\lambda(\bar s)=\overline{P_\lambda(s)}$, $P_\lambda(0)=P_\lambda(1)=1$, and its zeros, with multiplicity, are exactly
\begin{equation}\label{eq:family-divisor}
 \frac12\pm i\gamma_k(\lambda),\qquad k\ge1 .
\end{equation}
For $0<\lambda\le1$ the ordinates $\gamma_k(\lambda)$ are strictly increasing in $k$ and all zeros of $P_\lambda$ are simple; at $\lambda=0$ the multiplicities are those of the ordinate list $(\gamma_k)$.  For every $j\ge0$, $P_\lambda^{(j)}\to P_{\lambda_0}^{(j)}$ locally uniformly as $\lambda\to\lambda_0$.  In particular no member $\zeta_\lambda$ of the family, $\zeta_{Z1}=\zeta_0$ included, can have a zero off the critical line other than $-2,-4,\dots$.
\end{theorem}

\begin{proof}
By Lemma~\CGref{lem:bounded-displacement}, $|\gamma_k(\lambda)-\gamma_k|\le C_\eta$, and $\gamma_k\asymp k/\log k$; so $\gamma_k(\lambda)\ge\tfrac12\gamma_k$ for $k\ge k_0$, uniformly in $\lambda$.  On a compact $K\subset\C$,
\begin{equation}\label{eq:lambda-normal}
 \sum_{k\ge1}\sup_{s\in K,\ 0\le\lambda\le1}\left|\frac{q(s)}{\tfrac14+\gamma_k(\lambda)^2}\right|<\infty,
\end{equation}
which gives joint normal convergence and, by Cauchy's formula, convergence of all $s$-derivatives.  For $|s|\le R_{\rm fam}$, uniformly in $\lambda$,
\begin{equation}\label{eq:order-one}
 \log\max_{|s|\le R_{\rm fam}}|P_\lambda(s)|\le\sum_{k\ge1}\log\Bigl(1+\frac{R_{\rm fam}^2+R_{\rm fam}}{\tfrac14+\gamma_k(\lambda)^2}\Bigr)=O(R_{\rm fam}\log R_{\rm fam}),
\end{equation}
so the order is at most one; the zero count is $\asymp T\log T$, so the exponent of convergence is one and the order is exactly one.  Symmetries, normalization and divisor follow factor by factor as in Theorem~\CGref{thm:canonical}.  For $0<\lambda\le1$,
\begin{equation}\label{eq:ordinate-gaps}
 \gamma_{k+1}(\lambda)-\gamma_k(\lambda)=(1-\lambda)(\gamma_{k+1}-\gamma_k)+\lambda(\tau_{k+1}-\tau_k)\ge\lambda(\tau_{k+1}-\tau_k)>0,
\end{equation}
so distinct factors vanish at distinct points and each zero is simple; at $\lambda=0$ equal ordinates give equal factors.
\end{proof}

\begin{remark}[All differences lie in $\gamma_k-\tau_k$]\label{rem:only-gamma-tau}
For $\lambda,\lambda'\in[0,1]$ and $s$ off the divisors,
\begin{equation}\label{eq:ratio-members}
 \frac{\zeta_\lambda(s)}{\zeta_{\lambda'}(s)}=\prod_{k\ge1}\frac{\bigl(\tfrac14+(\gamma_k+\lambda'\eta_k)^2\bigr)\bigl(\tfrac14+(\gamma_k+\lambda\eta_k)^2+q(s)\bigr)}{\bigl(\tfrac14+(\gamma_k+\lambda\eta_k)^2\bigr)\bigl(\tfrac14+(\gamma_k+\lambda'\eta_k)^2+q(s)\bigr)} .
\end{equation}
The gamma factor cancels; the members differ only through $\eta_k=\tau_k-\gamma_k$, and all of them coincide if every $\eta_k$ vanishes.
\end{remark}

\begin{theorem}[Infinitesimal generator]\label{thm:generator}
Put $A_{k,\lambda}:=\tfrac14+\gamma_k(\lambda)^2$.  Off the divisor of $P_\lambda$,
\begin{equation}\label{eq:generator}
 \mathcal V_\lambda(s):=\partial_\lambda\log P_\lambda(s)=-2q(s)\sum_{k\ge1}\frac{\gamma_k(\lambda)\,\eta_k}{A_{k,\lambda}\bigl(A_{k,\lambda}+q(s)\bigr)},
\end{equation}
the series converging locally uniformly on compact sets separated from the divisor.  Since $G$ does not depend on $\lambda$,
\begin{equation}\label{eq:generator-zeta}
 \partial_\lambda\log\zeta_\lambda=\mathcal V_\lambda,\qquad \partial_\lambda\log F_\lambda=-\mathcal V_\lambda .
\end{equation}
\end{theorem}

\begin{proof}
Since $\partial_\lambda A_{k,\lambda}=2\gamma_k(\lambda)\eta_k$,
\begin{equation}\label{eq:mode-derivative}
 \partial_\lambda\log\Bigl(1+\frac q{A_{k,\lambda}}\Bigr)=-\frac{2q\,\gamma_k(\lambda)\eta_k}{A_{k,\lambda}(A_{k,\lambda}+q)} .
\end{equation}
On compact sets separated from the divisor the summand is $O_K(|\eta_k|\gamma_k^{-3})$, and $\sum_k\gamma_k^{-3}<\infty$, so termwise differentiation is justified.
\end{proof}

\begin{proposition}[Integrated multiplier]\label{prop:integrated-multiplier}
On every component of $\C\setminus\bigcup_{\lambda'\in[0,1]}\{\tfrac12\pm i\gamma_k(\lambda'):k\ge1\}$,
\begin{equation}\label{eq:integrated-zeta}
 \zeta_\lambda(s)=\zeta_{Z1}(s)\exp\Bigl(\int_0^\lambda\mathcal V_{\lambda'}(s)\,d\lambda'\Bigr),
\end{equation}
\begin{equation}\label{eq:integrated-inverse}
 F_\lambda(s)=F_{Z1}(s)\exp\Bigl(-\int_0^\lambda\mathcal V_{\lambda'}(s)\,d\lambda'\Bigr),
\end{equation}
and the total multiplier is
\begin{equation}\label{eq:transport-multiplier}
 \Pi_{\rm tr}(s):=\exp\Bigl(-\int_0^1\mathcal V_{\lambda'}(s)\,d\lambda'\Bigr)=\frac{P_{Z1}(s)}{P_{\RVM}^{\crit}(s)},\qquad \zeta_{Z1}=\zeta_{\RVM}\,\Pi_{\rm tr} .
\end{equation}
The identities extend meromorphically to $\C$ through the product form.
\end{proposition}

\begin{proof}
\CGeqref{eq:integrated-zeta} and \CGeqref{eq:integrated-inverse} are the fundamental theorem of calculus applied to \CGeqref{eq:generator-zeta}.  For \CGeqref{eq:transport-multiplier}, integrate each mode with $r=\gamma_k(\lambda')$, $dr=\eta_k\,d\lambda'$:
\begin{equation}\label{eq:mode-integral}
 \int_0^1\partial_{\lambda'}\log\Bigl(1+\frac q{A_{k,\lambda'}}\Bigr)d\lambda'\equiv\log\frac{1+q/(\tfrac14+\tau_k^2)}{1+q/(\tfrac14+\gamma_k^2)}\pmod{2\pi i}.
\end{equation}
Normal convergence permits summation, and exponentiation removes the branch ambiguity.
\end{proof}

\begin{corollary}[Expansion at the endpoint]\label{cor:endpoint-expansion}
Let $K\subset\C$ be compact, disjoint from the zeros of $P_{Z1}$ and from $s=1$ (equivalently one may state the expansion for the product ratio $P_\lambda/P_{Z1}$ without this exclusion).  Uniformly for $s\in K$, as $\lambda\downarrow0$,
\begin{equation}\label{eq:endpoint-zeta-expansion}
 \zeta_\lambda(s)=\zeta_{Z1}(s)\bigl(1+\lambda\mathcal V_0(s)+O_K(\lambda^2)\bigr),
\end{equation}
\begin{equation}\label{eq:endpoint-inverse-expansion}
 F_\lambda(s)=F_{Z1}(s)\bigl(1-\lambda\mathcal V_0(s)+O_K(\lambda^2)\bigr).
\end{equation}
\end{corollary}

\begin{proof}
By Theorem~\CGref{thm:path} the zeros of $P_\lambda$ converge to those of $P_{Z1}$, so $K$ avoids the moving divisor for small $\lambda$.  The series for $\mathcal V_\lambda$ and its $\lambda$-derivative converge normally on $K$ (the differentiated summands are $O_K(\gamma_k^{-3})$), so $\mathcal V_\lambda=\mathcal V_0+O_K(\lambda)$, and \CGeqref{eq:integrated-zeta} gives $\log(\zeta_\lambda/\zeta_{Z1})=\lambda\mathcal V_0+O_K(\lambda^2)$.
\end{proof}

\begin{proposition}[Residues along the family]\label{prop:residues}
Suppose $\gamma_k$ is a simple ordinate (it occurs once in the list).  Put $\rho_{k,\lambda}=\tfrac12+i\gamma_k(\lambda)$ and, for $0<\lambda\le1$,
\begin{equation}\label{eq:residue-transport}
 \mathfrak r_{k,\lambda}:=\Res_{s=\rho_{k,\lambda}}\frac1{s\,\zeta_\lambda(s)}=\frac{G(\rho_{k,\lambda})}{\rho_{k,\lambda}P_\lambda'(\rho_{k,\lambda})}.
\end{equation}
Then
\begin{equation}\label{eq:residue-limit}
 \mathfrak r_{k,\lambda}\longrightarrow\frac{G(\tfrac12+i\gamma_k)}{(\tfrac12+i\gamma_k)P_{Z1}'(\tfrac12+i\gamma_k)}\qquad(\lambda\downarrow0).
\end{equation}
\end{proposition}

\begin{proof}
Theorem~\CGref{thm:path} gives convergence of the zero locations and local uniform convergence of $P_\lambda'$; the limit denominator is nonzero because the endpoint zero is simple, and $G$ is continuous.
\end{proof}

\begin{remark}[Multiplicity]\label{rem:multiplicity}
For $0<\lambda\le1$ every zero of $\zeta_\lambda$ is simple.  At $\lambda=0$ a repeated ordinate of $\zeta$ appears in $\zeta_{Z1}$ with the same multiplicity, and the residues of Proposition~\CGref{prop:residues} are replaced by the full principal part.  Multiplicity plays no role in Step~3: Theorem~\CGref{thm:step3} uses only that $\zeta$ does not vanish in $\Re s>\tfrac12$.
\end{remark}

\section{The endpoint \texorpdfstring{$\zeta_{Z1}$}{zeta\_Z1}, Hamburger's theorem, and Step 3}
\label{sec:Z1-step3}

\subsection{The analytic class of the endpoint}
\label{subsec:Z1-class}

\begin{proposition}[Analytic class of $\zeta_{Z1}$]\label{prop:Z1-class}
The endpoint
\begin{equation}\label{eq:Z1-product}
 P_{Z1}(s)=\prod_{k\ge1}\left(1+\frac{q(s)}{\tfrac14+\gamma_k^2}\right),\qquad \zeta_{Z1}(s)=\frac{P_{Z1}(s)}{G(s)},
\end{equation}
has the following properties.
\begin{enumerate}[label=\textup{(\roman*)}]
\item $P_{Z1}$ is entire of order exactly one, $P_{Z1}(1-s)=P_{Z1}(s)$, $P_{Z1}(0)=P_{Z1}(1)=1$;
\item the nontrivial zeros of $\zeta_{Z1}$ are exactly
\begin{equation}\label{eq:Z1-divisor}
 \frac12\pm i\gamma_k,\qquad k\ge1,
\end{equation}
with the multiplicities of the ordinate list; its other zeros are $s=-2,-4,\dots$.  It is impossible for $\zeta_{Z1}$ to vanish at any point off the critical line other than $-2,-4,\dots$: every factor of $P_{Z1}$ vanishes only at $\tfrac12\pm i\gamma_k$;
\item $\zeta_{Z1}$ has a simple pole at $s=1$ with residue $1$;
\item $\zeta_{Z1}$ satisfies the functional equation of $\zeta$:
\begin{equation}\label{eq:Z1-FE}
 \pi^{-s/2}\Gamma(s/2)\zeta_{Z1}(s)=\pi^{-(1-s)/2}\Gamma((1-s)/2)\zeta_{Z1}(1-s);
\end{equation}
\item $(s-1)\zeta_{Z1}(s)$ is entire of order exactly one.
\end{enumerate}
\end{proposition}

\begin{proof}
(i) and the divisor in (ii) are the case $\lambda=0$ of Theorem~\CGref{thm:path}; the trivial zeros and the absence of other zeros come from $1/G$ as in Theorem~\CGref{thm:canonical}.  For (iii), $P_{Z1}(1)=1$ and
\begin{equation}\label{eq:G-near-one}
 G(s)=(s-1)+O\bigl((s-1)^2\bigr)\qquad(s\to1).
\end{equation}
For (iv),
\begin{equation}\label{eq:Z1-FE-proof}
 \pi^{-s/2}\Gamma(s/2)\zeta_{Z1}(s)=\frac{P_{Z1}(s)}{q(s)},
\end{equation}
and both $P_{Z1}$ and $q$ are invariant under $s\mapsto1-s$.  For (v),
\begin{equation}\label{eq:Z1-entire-removal}
 (s-1)\zeta_{Z1}(s)=P_{Z1}(s)\,\frac{\pi^{s/2}}{s\,\Gamma(s/2)},
\end{equation}
where $s\Gamma(s/2)\to2$ at $s=0$ and $1/\Gamma$ is entire of order one; the product has order at most one, and the zero count $\asymp T\log T$ gives order at least one.
\end{proof}

Define the ratio
\begin{equation}\label{eq:Z1-defect-ratio}
 \mathscr R_{Z1}(s):=\frac{\zeta_{Z1}(s)}{\zeta(s)}=\frac{P_{Z1}(s)}{P_\zeta(s)},
\end{equation}
a meromorphic function; $\mathscr R_{Z1}\equiv1$ is the identity $\zeta_{Z1}=\zeta$.

\subsection{Hamburger's theorem at the endpoint}
\label{subsec:Hamburger}

\begin{theorem}[Hamburger closure]\label{thm:Hamburger-Z1}
Consider the condition
\begin{equation}\label{eq:Z1-D}
 \text{(D)}\qquad\zeta_{Z1}(s)=\sum_{n\ge1}\frac{a_n}{n^s}\qquad\text{absolutely for }\Re s>1 .
\end{equation}
Then
\begin{equation}\label{eq:Hamburger-equivalence}
 (\mathrm D)\quad\Longleftrightarrow\quad\mathscr R_{Z1}\equiv1\quad\Longleftrightarrow\quad\zeta_{Z1}=\zeta\quad\Longleftrightarrow\quad\mathrm{RH}.
\end{equation}
\end{theorem}

\begin{proof}
Assume (D).  Then $\zeta_{Z1}(1-s)=\sum_na_nn^{-(1-s)}$ converges absolutely for $\Re s<0$, $(s-1)\zeta_{Z1}$ is entire of finite order (Proposition~\CGref{prop:Z1-class}(v)), and \CGeqref{eq:Z1-FE} holds.  These are the hypotheses of Hamburger's theorem \cite{Hamburger1921}, \cite[\S2.13]{Titchmarsh1986}, which gives
\begin{equation}\label{eq:Hamburger-conclusion}
 \zeta_{Z1}=C_H\,\zeta
\end{equation}
for a constant $C_H$.  Both functions have residue $1$ at $s=1$, so $C_H=1$.  If $\zeta_{Z1}=\zeta$, the nontrivial zeros of $\zeta$ are those of $\zeta_{Z1}$, namely \CGeqref{eq:Z1-divisor}, which is RH.  If RH holds, every $\rho$ with $\Im\rho>0$ is $\tfrac12+i\gamma_k$, so $\rho(1-\rho)=\tfrac14+\gamma_k^2$, and Theorem~\CGref{thm:Hadamard} gives
\begin{equation}\label{eq:RH-product}
 P_\zeta(s)=\prod_{k\ge1}\Bigl(1+\frac{q(s)}{\tfrac14+\gamma_k^2}\Bigr)=P_{Z1}(s)
\end{equation}
exactly, hence $\zeta=\zeta_{Z1}$.  Finally $\zeta=\sum_nn^{-s}$ converges absolutely for $\Re s>1$, so $\zeta_{Z1}=\zeta$ gives (D).
\end{proof}

\subsection{Condition (D) and the integer-weighted zeta}
\label{subsec:D-multiplier}

By Corollary~\CGref{cor:ZRVM-Dirichlet}, $\zeta_{\Z,\RVM}$ is an absolutely convergent Dirichlet series in $\Re s>1$ with first coefficient $1$, unconditionally.  Define the multiplier
\begin{equation}\label{eq:K-multiplier}
 \mathcal K(s):=\frac{\zeta_{Z1}(s)}{\zeta_{\Z,\RVM}(s)} .
\end{equation}

\begin{theorem}[Reduction of (D) to the multiplier]\label{thm:D-multiplier}
\begin{enumerate}[label=\textup{(\roman*)}]
\item If $\mathcal K$ has an expansion $\mathcal K(s)=\sum_{n\ge1}k_nn^{-s}$ that converges absolutely in $\Re s>1$, then \CGeqref{eq:Z1-D} holds.
\item If \CGeqref{eq:Z1-D} holds, then $\mathcal K(s)=\sum_{n\ge1}k_nn^{-s}$ absolutely on the whole half-plane $\Re s>1$.
\end{enumerate}
Thus (D) holds if and only if $\mathcal K$ is an absolutely convergent Dirichlet series in $\Re s>1$.
Hence, by Theorem~\CGref{thm:Hamburger-Z1}, an absolutely convergent Dirichlet expansion of $\mathcal K$ in $\Re s>1$ gives $\zeta_{Z1}=\zeta$, RH, and, by Theorem~\CGref{thm:step3}, $M(x)=O_\eps(x^{1/2+\eps})$.
\end{theorem}

\begin{proof}
(i) $\zeta_{Z1}=\zeta_{\Z,\RVM}\,\mathcal K$ is the Dirichlet convolution of two absolutely convergent series in $\Re s>1$, and $\sum_N|(a^{\Z}*k)_N|N^{-\sigma}\le\sum_N|a_N^{\Z}|N^{-\sigma}\sum_n|k_n|n^{-\sigma}<\infty$.  (ii) By \CGeqref{eq:ZRVM-inverse-Dirichlet}, $1/\zeta_{\Z,\RVM}$ is an absolutely convergent Dirichlet series in $\Re s>1$, and so is its product $\mathcal K$ with \CGeqref{eq:Z1-D}.
\end{proof}

By \CGeqref{eq:transport-multiplier} and Corollary~\CGref{cor:weighted-factorization},
\begin{equation}\label{eq:K-explicit}
 \mathcal K(s)=\Pi_{\rm tr}(s)\,e^{-h_{\RVM}(s)}=\frac{P_{Z1}(s)}{P_{\RVM}^{\crit}(s)}\,e^{-h_{\RVM}(s)} ,
\end{equation}
an explicit function of the displacements $\gamma_k-\tau_k$ and of the cell masses $c_n^{\RVM}$.  Condition (D) is thereby located in a single object built from the two constructions of this paper: the $\lambda$-family and the integer weights.

\begin{remark}[Continuous-to-atomic meaning of (D)]\label{rem:D-atomic}
The prime-power count $J_{Z1}^{\rm M}$ of $\zeta_{Z1}$ (Theorem~\CGref{thm:mellin-counts}) is an $L^2_{\rm loc}$ count built from the zeros, and (D) makes its distributional derivative the atomic measure $\sum_{p,m}m^{-1}\delta_{p^m}$.  The final step of the method is therefore the passage of $\zeta_{Z1}$ into the discrete Dirichlet category.  Theorem~\CGref{thm:cell-form} states it exactly: (D) holds if and only if $J_{Z1}^{\rm M}$ is constant between consecutive integers and its jumps obey $\sum_{n\le X}|\mathsf b_n|=O_\eps(X^{1+\eps})$, and Corollary~\CGref{cor:product-identity} reduces it to the single identity \CGeqref{eq:product-identity} between the displacements $\gamma_k-\tau_k$ and $\varpi_{\RVM}$.
\end{remark}

\subsection{Step 3: the Mertens bound from \texorpdfstring{$\zeta_{Z1}=\zeta$}{zeta\_Z1 = zeta}}
\label{subsec:step3}

\begin{lemma}[Vertical growth of $\zeta$]\label{lem:zeta-growth}
For $\tfrac14\le\sigma\le4$ and $|t|\ge2$,
\begin{equation}\label{eq:zeta-growth}
 |\zeta(\sigma+it)|\le15|t| .
\end{equation}
\end{lemma}

\begin{proof}
For $\sigma>0$ \cite[(2.1.5)]{Titchmarsh1986},
\begin{equation}\label{eq:zeta-integral-rep}
 \zeta(s)=\frac s{s-1}-s\int_1^\infty\{y\}y^{-s-1}\,dy .
\end{equation}
For $|t|\ge2$ and $\sigma\le4$ one has $|s|\le3|t|$, $|s/(s-1)|\le|s|/|t|\le3$, and the integral is at most $1/\sigma\le4$; hence $|\zeta(s)|\le3+12|t|\le15|t|$.
\end{proof}

\begin{lemma}[Logarithm of $\zeta$ off the critical line]\label{lem:log-zeta}
Assume $\zeta(s)\ne0$ for $\Re s>\tfrac12$.  Let $0<\delta\le\tfrac14$ and
\begin{equation}\label{eq:three-circles-exponent}
 \omega(\delta):=\frac{\log(3-2\delta)}{\log(3-\delta)}<1 .
\end{equation}
Then, with the branch of $\log\zeta$ equal to $\sum_{n\ge2}\Lambda(n)n^{-s}/\log n$ on $\Re s>1$,
\begin{equation}\label{eq:log-zeta-bound}
 |\log\zeta(\sigma+it)|\le\Bigl(\frac{48}\delta\log|t|\Bigr)^{\omega(\delta)}\qquad(\sigma\ge\tfrac12+\delta,\ |t|\ge4).
\end{equation}
Consequently, for every $\eps>0$ there is $t_2=t_2(\delta,\eps)$ with
\begin{equation}\label{eq:inv-zeta-bound}
 \Bigl|\frac1{\zeta(\sigma+it)}\Bigr|\le|t|^{\eps}\qquad(\sigma\ge\tfrac12+\delta,\ |t|\ge t_2).
\end{equation}
\end{lemma}

\begin{proof}
Fix $t$ with $|t|\ge4$ and put $s_0=2+it$, $R=\tfrac32-\tfrac\delta4$, $r_3=\tfrac32-\tfrac\delta2$.  The closed disc $|s-s_0|\le R$ lies in $\tfrac12+\tfrac\delta4\le\Re s\le\tfrac72$, $|\Im s-t|\le\tfrac32$; it avoids $s=1$ and, by hypothesis, every zero of $\zeta$.  Hence $\log\zeta$ has a holomorphic branch on a neighbourhood of this disc, chosen to agree with the Dirichlet-series branch on the connected overlap with $\Re s>1$.  Such branches are defined component by component on zero-free regions; no global branch on $\Re s>\tfrac12$ is asserted.  On the disc $|\Im s|\ge|t|/2\ge2$ and $|\Im s|\le2|t|$, so Lemma~\CGref{lem:zeta-growth} gives $\Re\log\zeta=\log|\zeta|\le\log(30|t|)$ there, while $|\log\zeta(s_0)|\le\log\zeta(2)$.  The Borel--Carath\'eodory lemma \cite[\S5.5]{Titchmarsh1939}, applied with radii $r_3<R$,
\begin{equation}\label{eq:borel-caratheodory}
 \max_{|s-s_0|\le r_3}|\log\zeta(s)|\le\frac{2r_3}{R-r_3}\max_{|s-s_0|=R}\Re\log\zeta(s)+\frac{R+r_3}{R-r_3}|\log\zeta(s_0)|,
\end{equation}
together with $R-r_3=\delta/4$, $2r_3\le3$, $R+r_3\le3$ and $\log30+\log\zeta(2)<4\le3\log|t|$, gives
\begin{equation}\label{eq:BC-bound}
 \max_{|s-s_0|\le r_3}|\log\zeta(s)|\le\frac{12}\delta\bigl(\log(30|t|)+\log\zeta(2)\bigr)\le\frac{48}\delta\log|t| .
\end{equation}
Now let $\tfrac12+\delta\le\sigma<\tfrac32$ and set $r_1=\tfrac12$, $r_2=2-\sigma\in(\tfrac12,\tfrac32-\delta]$.  On $|s-s_0|=r_1$ one has $\Re s\ge\tfrac32$, so $|\log\zeta(s)|\le\log\zeta(\tfrac32)<1$.  Hadamard's three-circle theorem \cite[\S5.3]{Titchmarsh1939} gives, with $M_j=\max_{|s-s_0|=r_j}|\log\zeta|$,
\begin{equation}\label{eq:three-circles}
 M_2\le M_1^{1-\omega}M_3^{\omega},\qquad \omega=\frac{\log(r_2/r_1)}{\log(r_3/r_1)}\le\frac{\log(3-2\delta)}{\log(3-\delta)}=\omega(\delta).
\end{equation}
Since $M_1<1$ and $M_3\le\frac{48}\delta\log|t|$ with $\frac{48}\delta\log|t|>1$, \CGeqref{eq:three-circles} gives \CGeqref{eq:log-zeta-bound} at $s=\sigma+it$, which lies on the circle of radius $r_2$.  For $\sigma\ge\tfrac32$, $|\log\zeta|\le\log\zeta(\tfrac32)<1$.  Finally $|1/\zeta|=\exp(-\Re\log\zeta)\le\exp|\log\zeta|$, and $(\frac{48}\delta\log|t|)^{\omega(\delta)}\le\eps\log|t|$ for $|t|\ge t_2(\delta,\eps)$ because $\omega(\delta)<1$.
\end{proof}

\begin{lemma}[Perron kernel]\label{lem:perron-kernel}
For $y>0$, $y\ne1$, $c>0$, $T>0$ \cite[Ch.~17]{Davenport1980},
\begin{equation}\label{eq:perron-kernel}
 \left|\frac1{2\pi i}\int_{c-iT}^{c+iT}\frac{y^s}{s}\,ds-\mathbf 1_{\{y>1\}}\right|\le y^c\min\Bigl(1,\frac1{T|\log y|}\Bigr).
\end{equation}
\end{lemma}

\begin{theorem}[Step 3]\label{thm:step3}
If $\zeta_{Z1}=\zeta$, then for every $\eps>0$
\begin{equation}\label{eq:step3-bound}
 M(x)=O_\eps\bigl(x^{1/2+\eps}\bigr)\qquad(x\ge1).
\end{equation}
\end{theorem}

\begin{proof}
\emph{Non-vanishing.}  By Proposition~\CGref{prop:Z1-class}(ii), $\zeta_{Z1}$ has no zeros in $\Re s>\tfrac12$; since $\zeta=\zeta_{Z1}$, neither has $\zeta$, and Lemma~\CGref{lem:log-zeta} applies.

\emph{Truncation.}  Let $x=N+\tfrac12$ with $N\ge3$ an integer, $c=1+1/\log x$ and $T=x$.  Applying \CGeqref{eq:perron-kernel} with $y=x/n$ to each term of $\sum_n\mu(n)n^{-s}=1/\zeta(s)$, which converges absolutely on $\Re s=c$, gives
\begin{equation}\label{eq:perron-truncated}
 M(x)=\frac1{2\pi i}\int_{c-iT}^{c+iT}\frac{x^s}{s\,\zeta(s)}\,ds+O\Bigl(\sum_{n\ge1}\Bigl(\frac xn\Bigr)^c\min\Bigl(1,\frac1{T|\log(x/n)|}\Bigr)\Bigr).
\end{equation}
For $n\le x/2$ or $n\ge2x$, $|\log(x/n)|\ge\log2$, and these terms contribute at most $x^c\zeta(c)/(T\log2)\le e\,x(1+\log x)/(T\log2)$.  For $x/2<n<2x$, $(x/n)^c\le4$ and $|\log(x/n)|\ge|n-x|/(2x)$; as $|n-x|\in\{j-\tfrac12:j\ge1\}$, each value taken at most twice, these terms contribute at most $8\sum_{1\le j\le2x}\min(1,4x/(T(2j-1)))\ll1+(x/T)\log x$.  Hence
\begin{equation}\label{eq:perron-error}
 M(x)=\frac1{2\pi i}\int_{c-iT}^{c+iT}\frac{x^s}{s\,\zeta(s)}\,ds+O(\log x).
\end{equation}
\emph{Contour shift.}  Fix $0<\delta\le\tfrac14$.  The integrand is holomorphic on the closed rectangle with vertices $c\pm iT$, $\tfrac12+\delta\pm iT$: $\zeta$ has no zeros there, $s=1$ is a zero of $1/\zeta$, and $s=0$ lies outside.  By Cauchy's theorem
\begin{equation}\label{eq:contour-shift}
 \int_{c-iT}^{c+iT}=\int_{\frac12+\delta-iT}^{\frac12+\delta+iT}+\int_{\frac12+\delta+iT}^{c+iT}-\int_{\frac12+\delta-iT}^{c-iT}.
\end{equation}
For $x\ge t_2(\delta,\eps)$, \CGeqref{eq:inv-zeta-bound} on $|t|=T=x$ bounds the horizontal integrals by
\begin{equation}\label{eq:horizontal-bound}
 \frac1{2\pi}\int_{\frac12+\delta}^{c}\frac{x^\sigma T^{\eps}}{T}\,d\sigma\le x^cT^{\eps-1}=e\,x^{\eps}.
\end{equation}
On the left side, $1/\zeta$ is continuous on the compact segment $\{\tfrac12+\delta+it:|t|\le t_2\}$, bounded by a constant $C(\delta,\eps)$, and \CGeqref{eq:inv-zeta-bound} holds for $|t|\ge t_2$; hence
\begin{equation}\label{eq:left-bound}
 \frac1{2\pi}\int_{-T}^{T}\frac{x^{1/2+\delta}}{|\tfrac12+\delta+it|\,|\zeta(\tfrac12+\delta+it)|}\,dt\le x^{1/2+\delta}\Bigl(C'(\delta,\eps)+\int_{t_2}^{T}t^{\eps-1}dt\Bigr)\ll_{\delta,\eps}x^{1/2+\delta+\eps}.
\end{equation}
\emph{Conclusion.}  \CGeqref{eq:perron-error}--\CGeqref{eq:left-bound} give $M(x)\ll_{\delta,\eps}x^{1/2+\delta+\eps}$ for $x=N+\tfrac12$; taking $\delta=\eps=\eps_0/2$ (with $\eps_0\le\tfrac12$) gives $O_{\eps_0}(x^{1/2+\eps_0})$.  For arbitrary $x\ge3$, $M(x)=M(\lfloor x\rfloor+\tfrac12)$ and $\lfloor x\rfloor+\tfrac12\le2x$; for $1\le x<3$, $|M(x)|\le2$.
\end{proof}

\begin{theorem}[Transfer of the synthetic Mertens bound to $M$]\label{thm:transfer}
Assume $\zeta_{Z1}=\zeta$ and let $0<\delta\le\tfrac14$, $\eps>0$.  Then:
\begin{enumerate}[label=\textup{(\roman*)}]
\item on $\Re s>\tfrac12$,
\begin{equation}\label{eq:transfer-factorization}
 \frac1{\zeta(s)}=F_{\RVM}(s)\,\Pi_{\rm tr}(s)^{-1},
 \qquad
 \Pi_{\rm tr}=\frac{P_{Z1}}{P_{\RVM}^{\crit}}=\frac{\zeta}{\zeta_{\RVM}},
\end{equation}
where $\Pi_{\rm tr}$ is holomorphic and zero-free in $\Re s>\tfrac12$ and $\Pi_{\rm tr}(1)=1$;
\item there is $t_3=t_3(\delta,\eps)$ such that, for $\tfrac12+\delta\le\sigma\le2$ and $|t|\ge t_3$,
\begin{equation}\label{eq:transfer-multiplier-bound}
 |F_{\RVM}(\sigma+it)|\le4L,\qquad |\Pi_{\rm tr}(\sigma+it)^{\pm1}|\le|t|^{\eps};
\end{equation}
\item the following classical consequence of RH is recovered at the endpoint:
\begin{equation}\label{eq:transfer-conclusion}
 M(x)=O_\eps\bigl(x^{1/2+\eps}\bigr).
\end{equation}
\end{enumerate}
\end{theorem}

\begin{proof}
(i) By \CGeqref{eq:transport-multiplier}, $\zeta_{Z1}=\zeta_{\RVM}\Pi_{\rm tr}$; with $\zeta_{Z1}=\zeta$ this is \CGeqref{eq:transfer-factorization}.  In $\Re s>\tfrac12$, $\zeta=\zeta_{Z1}$ has no zeros (Proposition~\CGref{prop:Z1-class}(ii)) and $\zeta_{\RVM}$ has none (Theorem~\CGref{thm:canonical}); both have a simple pole of residue $1$ at $s=1$, so $\Pi_{\rm tr}$ is holomorphic there with $\Pi_{\rm tr}(1)=1$, and it is zero-free in $\Re s>\tfrac12$.

(ii) By Theorem~\CGref{thm:transition}(a) with $B=\tfrac32$, for $\delta\le\sigma'\le\tfrac32$ and $|t|\ge t_0$ we have $\zeta_{\RVM}=\varpi_{\RVM}(1+w)(1+\mathcal E_0)$ with $|w|\le C|t|^{-\delta}\le\tfrac12$ and $|\mathcal E_0|\le\tfrac12$ for $|t|$ large; hence $|F_{\RVM}|\le4L$ and
\begin{equation}\label{eq:log-zetaRVM-bounded}
 \log\zeta_{\RVM}=\log\varpi_{\RVM}+\Log(1+w)+\Log(1+\mathcal E_0)=O(1).
\end{equation}
By Lemma~\CGref{lem:log-zeta}, which applies because $\zeta=\zeta_{Z1}$ has no zeros in $\Re s>\tfrac12$, $|\log\zeta(\sigma+it)|\le(\frac{48}\delta\log|t|)^{\omega(\delta)}$ with $\omega(\delta)<1$.  Therefore
\begin{equation}\label{eq:log-Pi-tr}
 |\log\Pi_{\rm tr}(\sigma+it)|=|\log\zeta-\log\zeta_{\RVM}|\le\Bigl(\frac{48}\delta\log|t|\Bigr)^{\omega(\delta)}+O(1)\le\eps\log|t|
\end{equation}
for $|t|\ge t_3(\delta,\eps)$, which gives \CGeqref{eq:transfer-multiplier-bound}.

(iii) The required estimate for $1/\zeta$ already follows from the zero-free hypothesis $\zeta=\zeta_{Z1}$ through Lemma~\CGref{lem:log-zeta} and Theorem~\CGref{thm:step3}.  The factorization \CGeqref{eq:transfer-factorization} is therefore a reformulation of the same estimate, not an independent transfer from the synthetic Mertens bound.  By (i)--(ii), $\Pi_{\rm tr}^{-1}$ is holomorphic in $\Re s>\tfrac12$ and grows at most like $|t|^{\eps}$.  Hence $|1/\zeta(\sigma+it)|\le4L|t|^{\eps}$ for $\tfrac12+\delta\le\sigma\le2$, $|t|\ge t_3$, and the truncated Perron formula, contour shift and estimates \CGeqref{eq:perron-error}--\CGeqref{eq:left-bound} of Theorem~\CGref{thm:step3}, applied with this bound, give $M(x)\ll_{\delta,\eps}x^{1/2+\delta+\eps}$.  Choosing $\delta=\eps=\eps_0/2$ gives \CGeqref{eq:transfer-conclusion}.
\end{proof}

\begin{remark}[The transfer in Mellin form]\label{rem:transfer-mellin}
Formally, \CGeqref{eq:transfer-factorization} reads $M=(\Pi_{\rm tr}^{-1})^{\vee}\Mell M_{\RVM}^{\crit}$: the Mertens function is the synthetic Mertens function convolved with the inverse Mellin transform of the inverse transport multiplier.  The synthetic factor $F_{\RVM}$ contributes the exponent $\tfrac12$ and is bounded (Theorem~\CGref{thm:transition}); the multiplier contributes at most $|t|^{\eps}$.  All the information of the $\lambda$-family enters through $\Pi_{\rm tr}=P_{Z1}/P_{\RVM}^{\crit}$, that is, through the displacements $\gamma_k-\tau_k$.
\end{remark}

\begin{corollary}[Step 3 through (D)]\label{cor:step3-D}
Condition \CGeqref{eq:Z1-D} implies $M(x)=O_\eps(x^{1/2+\eps})$ for every $\eps>0$.
\end{corollary}

\begin{proof}
Theorem~\CGref{thm:Hamburger-Z1} and Theorem~\CGref{thm:step3}.
\end{proof}

\begin{remark}[Status of Step 3]\label{rem:Step3-status}
The implication $\zeta_{Z1}=\zeta\Rightarrow M(x)=O_\eps(x^{1/2+\eps})$ is the classical consequence RH $\Rightarrow M(x)=O_\eps(x^{1/2+\eps})$ (see \cite[\S14.25]{Titchmarsh1986}) rewritten and reproved in the endpoint notation.  It requires nothing beyond the identity: no separate statement about the zeros of $\zeta$ enters, because $\zeta=\zeta_{Z1}$ already carries the divisor \CGeqref{eq:Z1-divisor}.  The identity $\zeta_{Z1}=\zeta$ itself, equivalently (D), equivalently an absolutely convergent Dirichlet expansion of the multiplier $\mathcal K$ of \CGeqref{eq:K-multiplier} in $\Re s>1$ (Theorem~\CGref{thm:D-multiplier}), is not proved in this paper.  It is the open step of the method.
\end{remark}

\subsection{Reciprocal endpoint and finite-height Perron integrals}
\label{subsec:reciprocal}

Put
\begin{equation}\label{eq:FZ1}
 F_{Z1}(s):=\frac1{\zeta_{Z1}(s)}=F_{\RVM}(s)\,\Pi_{\rm tr}(s)^{-1},
\end{equation}
the second identity by \CGeqref{eq:transport-multiplier}.  For $c>1$ and finite $U$ define
\begin{equation}\label{eq:truncated-Perron}
 M_{\lambda,U}(x):=\frac1{2\pi i}\int_{c-iU}^{c+iU}F_\lambda(s)\frac{x^s}{s}\,ds,\qquad M_{Z1,U}:=M_{0,U}.
\end{equation}

\begin{proposition}[Finite-height Perron integrals along the family]\label{prop:finite-Perron}
For fixed $c>1$, $U<\infty$ and $x>0$,
\begin{equation}\label{eq:finite-Perron-limit}
 M_{\lambda,U}(x)\longrightarrow M_{Z1,U}(x)\qquad(\lambda\downarrow0).
\end{equation}
\end{proposition}

\begin{proof}
The segment $\{c+it:|t|\le U\}$ is compact and disjoint from every divisor \CGeqref{eq:family-divisor}, and $P_\lambda\to P_{Z1}$ uniformly on it, so $F_\lambda\to F_{Z1}$ uniformly there.
\end{proof}

\begin{theorem}[Reciprocal-coefficient criterion]\label{thm:coefficient-criterion}
Assume that in $\Re s>1$
\begin{equation}\label{eq:Z1-reciprocal-Dirichlet}
 F_{Z1}(s)=\sum_{n\ge1}\frac{\mu_{Z1}(n)}{n^s}
\end{equation}
absolutely, and let $M_{Z1}(x)=\sum_{n\le x}\mu_{Z1}(n)$, so that
\begin{equation}\label{eq:Z1-increments}
 \mu_{Z1}(n)=M_{Z1}(n)-M_{Z1}(n-1).
\end{equation}
Then the following are equivalent: (i) $M_{Z1}(n)=M(n)$ for every $n\ge1$; (ii) $\mu_{Z1}=\mu$; (iii) $F_{Z1}=1/\zeta$ in $\Re s>1$; (iv) $\zeta_{Z1}=\zeta$; (v) RH.
\end{theorem}

\begin{proof}
(i)$\Leftrightarrow$(ii) by \CGeqref{eq:Z1-increments}; (ii)$\Leftrightarrow$(iii) by uniqueness of absolutely convergent Dirichlet coefficients; (iii)$\Leftrightarrow$(iv) by analytic continuation, both functions being nonzero in $\Re s>1$; (iv)$\Leftrightarrow$(v) by Theorem~\CGref{thm:Hamburger-Z1}.
\end{proof}

\begin{remark}[A route that does not close Step 3]\label{rem:negative-convolution}
The following approach does not work.  One may try to expand both factors of $F_{Z1}=F_{\RVM}\,\Pi_{\rm tr}^{-1}$ in \CGeqref{eq:FZ1} as absolutely convergent Dirichlet series and identify their Dirichlet convolution with $\mu$.  By Theorem~\CGref{thm:no-Dirichlet-RVM} the expansion of $F_{\RVM}$ does not exist in any half-plane, so there are no coefficients to convolve.  The endpoint criterion concerns $\zeta_{Z1}$ itself, as in Theorem~\CGref{thm:Hamburger-Z1}.
\end{remark}

\section{The endpoint at infinity and the Euler products}
\label{sec:euler-products}

This section applies the clock decomposition of Theorem~\CGref{thm:halfplane-RVM} to the ordinates $\gamma_k$ of the endpoint, and constructs the prime-power counts of $\zeta_{\RVM}$ and $\zeta_{Z1}$ as Mellin transforms of their logarithms.  No hypothesis is used.

\subsection{The value of the endpoint at infinity}
\label{subsec:value-infinity}

With the zero count $N_\zeta$ of \CGeqref{eq:RvM-full} put
\begin{equation}\label{eq:frak-e}
 \mathfrak e(\tau):=N_\zeta(\tau)-\NR(\tau),\qquad \mathfrak e_1(\tau):=\int_{\tau_0}^{\tau}\mathfrak e(v)\,dv\qquad(\tau\ge\tau_0=4\pi).
\end{equation}
By \CGeqref{eq:RvM-full}, $\mathfrak e(\tau)=S(\tau)+O(1/\tau)$, and Littlewood's bound $\int_0^TS(t)\,dt=O(\log T)$ \cite[Thm.~9.9]{Titchmarsh1986} gives an absolute constant $C_{\mathfrak e}$ with
\begin{equation}\label{eq:frak-e1-bound}
 |\mathfrak e_1(\tau)|\le C_{\mathfrak e}\log\tau\qquad(\tau\ge\tau_0).
\end{equation}
For the clock the corresponding primitive is $-\mathfrak s_1$, bounded by Lemma~\CGref{lem:saw-primitive}.

\begin{theorem}[Clock decomposition of the endpoint]\label{thm:Z1-clock}
Let $s=\tfrac12+z$, $z=\sigma'+it$, $\sigma'>0$.  There is an absolute constant $R_{\rm end}>0$ such that
\begin{equation}\label{eq:Z1-clock}
 \log\zeta_{Z1}(s)=\log\varpi_{Z1}-\int_{\tau_0}^\infty\mathfrak e_1(\tau)\Bigl(\frac1{(\tau-iz)^2}+\frac1{(\tau+iz)^2}\Bigr)d\tau+O\Bigl(\frac1{|z|}\Bigr)\qquad(|z|\ge R_{\rm end}),
\end{equation}
where
\begin{equation}\label{eq:varpi-Z1}
 \varpi_{Z1}:=P_{Z1}(\tfrac12)\exp\Bigl(\nu_3+2\int_{\tau_0}^\infty\frac{\mathfrak e_1(\tau)}{\tau^2}\,d\tau\Bigr)>0 .
\end{equation}
Consequently:
\begin{enumerate}[label=\textup{(\roman*)}]
\item for $\sigma'>0$ and $|z|\ge R_{\rm end}$,
\begin{equation}\label{eq:Z1-growth}
 \bigl|\log\zeta_{Z1}(s)-\log\varpi_{Z1}\bigr|\le\frac{2(\pi+2)C_{\mathfrak e}}{\sigma'}\Bigl(1+\log\bigl(2|t|+2\sigma'+\tau_0\bigr)\Bigr)+\frac C{|z|};
\end{equation}
\item $\zeta_{Z1}(\sigma)\to\varpi_{Z1}$ as $\sigma\to+\infty$;
\item the value at infinity is the ordered product
\begin{equation}\label{eq:varpi-Z1-product}
 \varpi_{Z1}=\varpi_{\RVM}\prod_{k\ge1}\frac{\tfrac14+\tau_k^2}{\tfrac14+\gamma_k^2}:=\varpi_{\RVM}\lim_{K\to\infty}\prod_{k\le K}\frac{\tfrac14+\tau_k^2}{\tfrac14+\gamma_k^2}.
\end{equation}
\end{enumerate}
\end{theorem}

\begin{proof}
As in \CGeqref{eq:transition-product}, $\log P_{Z1}(s)-\log P_{Z1}(\frac12)=\sum_kf(\gamma_k)$ with $f(\tau)=\Log(1+z^2/\tau^2)$.  Since $\gamma_1>\tau_0$, $N_\zeta(\tau_0)=0$ and $\mathfrak e(\tau_0)=-\NR(\tau_0)$.  Stieltjes integration against $N_\zeta=\NR+\mathfrak e$ and two integrations by parts ($f\mathfrak e\to0$ and $\mathfrak e_1f'\to0$ at $\infty$, $\mathfrak e_1(\tau_0)=0$) give
\begin{equation}\label{eq:Z1-split}
 \sum_kf(\gamma_k)=\Lambda_{\rm sm}(z)+\int_{\tau_0}^\infty\mathfrak e_1f''\,d\tau,
\end{equation}
with the same $\Lambda_{\rm sm}$ as in \CGeqref{eq:transition-split}.  For $\sigma'>\tfrac12$ the proof of Theorem~\CGref{thm:halfplane-RVM} gives $\Lambda_{\rm sm}-\log G=\nu_3+O(|z|^{-1})$; for $0<\sigma'\le\tfrac12$ the same estimate holds because \CGeqref{eq:smooth-exact} is an exact formula valid on all of $\Re z>0$ (the singularities of $f'$ at $\tau=\pm iz$ lie off the positive axis), $|\tau^2+z^2|\ge\frac34|z|^2$ on $[0,\tau_0]$ for $|z|\ge2\tau_0$ whatever $\sigma'>0$, Stirling's formula with remainder $O(|v|^{-1})$ is uniform in $|\arg v|\le\frac\pi2$, and $\mathrm{Log}(z^2-\frac14)=2\,\mathrm{Log}\,z+O(|z|^{-2})$ since $\Im(z^2-\frac14)=2\sigma't$ has the sign of $t$, so no cut is crossed; an $O$-estimate is not transported by analytic continuation, and none is needed.  Thus, inserting $f''$ from \CGeqref{eq:saw-ibp} gives \CGeqref{eq:Z1-clock}.

For (i), $|\tau\mp iz|^2=(\tau\pm t)^2+\sigma'^2$.  For real $a$, splitting at $X=2|a|+2\sigma'+\tau_0$ and using $(\tau-a)^2\ge\tau^2/4$ for $\tau\ge X$,
\begin{equation}\label{eq:log-kernel}
 \int_{\tau_0}^\infty\frac{\log\tau}{(\tau-a)^2+\sigma'^2}\,d\tau\le\frac\pi{\sigma'}\log X+4\,\frac{1+\log X}X\le\frac{\pi+2}{\sigma'}(1+\log X).
\end{equation}
With \CGeqref{eq:frak-e1-bound} this gives \CGeqref{eq:Z1-growth}.  For (ii) take $t=0$: the right side of \CGeqref{eq:Z1-growth} tends to $0$ as $\sigma'\to\infty$.

For (iii), Theorem~\CGref{thm:halfplane-RVM} gives $\varpi_{\RVM}=P_{\RVM}^{\crit}(\frac12)\exp(\nu_3-2\int\mathfrak s_1\tau^{-2}d\tau)$, so
\begin{equation}\label{eq:varpi-ratio}
 \frac{\varpi_{Z1}}{\varpi_{\RVM}}=\frac{P_{Z1}(\frac12)}{P_{\RVM}^{\crit}(\frac12)}\exp\Bigl(2\int_{\tau_0}^\infty\frac{\mathfrak e_1(\tau)+\mathfrak s_1(\tau)}{\tau^2}\,d\tau\Bigr).
\end{equation}
Since $\mathfrak e+\psi(\NR+\frac12)=N_\zeta-n_{\RVM}$ and $(\mathfrak e_1+\mathfrak s_1)(X)/X\to0$,
\begin{equation}\label{eq:count-difference-integral}
 \int_{\tau_0}^\infty\frac{\mathfrak e_1+\mathfrak s_1}{\tau^2}\,d\tau=\lim_{X\to\infty}\int_{\tau_0}^X\frac{N_\zeta(\tau)-n_{\RVM}(\tau)}\tau\,d\tau=\lim_{K\to\infty}\sum_{k\le K}\log\frac{\tau_k}{\gamma_k}.
\end{equation}
In the last step, $\int_{\tau_0}^X(N_\zeta-n_{\RVM})\tau^{-1}d\tau=\sum_{\gamma_k\le X}\log\frac X{\gamma_k}-\sum_{\tau_k\le X}\log\frac X{\tau_k}$; the indices counted by only one of the two sums are $O(\log X)$ in number and, by Lemma~\CGref{lem:bounded-displacement}, have ordinates in $[X-C_\eta,X]$, so they contribute $O(\log X/X)$.  Finally $P_{Z1}(\frac12)/P_{\RVM}^{\crit}(\frac12)=\prod_k\frac{\gamma_k^2}{\frac14+\gamma_k^2}\cdot\frac{\frac14+\tau_k^2}{\tau_k^2}$, an absolutely convergent product, and multiplying by $\prod_{k\le K}\tau_k^2/\gamma_k^2$ gives \CGeqref{eq:varpi-Z1-product}.
\end{proof}

The ordinates enter $\varpi_{Z1}$ only through the displacements $\gamma_k-\tau_k$, in agreement with Remark~\CGref{rem:only-gamma-tau}.  The product in \CGeqref{eq:varpi-Z1-product} is an \emph{ordered} product: its factors are $1+O(|\eta_k|/\tau_k)$, and $\sum_k|\eta_k|/\tau_k$ is not known to converge ($\sum_k\tau_k^{-1}$ diverges), so absolute convergence is not asserted and no rearrangement is permitted.  The limit is that of the partial products over $k\le K$ in increasing height, which is the order in which \CGeqref{eq:count-difference-integral} is evaluated, and this order is part of the definition of $\varpi_{Z1}$; the equality of the ordered limit with $\varpi_{Z1}$ is what \CGeqref{eq:count-difference-integral} proves, and the value is not claimed to be independent of other groupings.

\begin{theorem}[The value at infinity decides the endpoint identity]\label{thm:varpi-Z1}
$\varpi_{Z1}\ge1$, and $\varpi_{Z1}=1$ if and only if $\zeta_{Z1}=\zeta$.
\end{theorem}

\begin{proof}
Write $\zeta_{Z1}=\zeta\,\mathscr R_{Z1}$ with $\mathscr R_{Z1}=P_{Z1}/P_\zeta$ as in \CGeqref{eq:Z1-defect-ratio}, and compare the products factor by factor, with the zeros $\rho_k=\beta_k+i\gamma_k$ listed with the ordinates (Theorem~\CGref{thm:Hadamard}).  Factors with $\rho_k(1-\rho_k)=\frac14+\gamma_k^2$ cancel.  The remaining factors of $P_\zeta$ come in pairs $\{\rho,1-\bar\rho\}$ with the same ordinate $\gamma$; with $d_\rho=\Re\rho-\frac12$ and $W_\rho=(s-\frac12)^2+\gamma^2$, the corresponding factor of $\mathscr R_{Z1}$ is
\begin{equation}\label{eq:pair-ratio}
 \frac{W_\rho^2}{(W_\rho-d_\rho^2)^2+4d_\rho^2\gamma^2}\cdot c_\gamma,\qquad
 c_\gamma=\frac{|\rho(1-\rho)|^2}{(\frac14+\gamma^2)^2}=1+\frac{2d_\rho^2(\gamma^2-\frac14)+d_\rho^4}{(\frac14+\gamma^2)^2}>1 .
\end{equation}
For real $\sigma>1$, $W_\rho\ge\gamma^2\ge\gamma_1^2$, and the logarithm $\ell_\rho(\sigma)$ of the first factor satisfies $|\ell_\rho(\sigma)|\le14d_\rho^2/W_\rho\le14d_\rho^2/\gamma^2$, while $0<\log c_\gamma\le c_\gamma-1\le3d_\rho^2/\gamma^2$.  Hence $\log\mathscr R_{Z1}(\sigma)=\sum_\rho\bigl(\ell_\rho(\sigma)+\log c_\gamma\bigr)$ is a series dominated, uniformly in $\sigma>1$, by $\sum_\rho17d_\rho^2/\gamma^2\le\frac{17}4\sum_\rho\gamma^{-2}<\infty$; each term tends to $\log c_\gamma$ as $\sigma\to\infty$ because $W_\rho\to\infty$; so, by dominated convergence for series, $\log\mathscr R_{Z1}(\sigma)\to\sum_\rho\log c_\gamma$.  With $\zeta(\sigma)\to1$ this gives $\varpi_{Z1}=\lim_{\sigma\to\infty}\zeta_{Z1}(\sigma)=\prod c_\gamma\ge1$ (an absolutely convergent product of factors $>1$), with equality exactly when there are no such pairs, that is, when $\mathscr R_{Z1}\equiv1$.
\end{proof}

\begin{corollary}[The endpoint identity as an identity of the displacements]\label{cor:product-identity}
The following are equivalent: (D); $\zeta_{Z1}=\zeta$; $\varpi_{Z1}=1$;
\begin{equation}\label{eq:product-identity}
 \lim_{K\to\infty}\prod_{k\le K}\frac{\tfrac14+\gamma_k^2}{\tfrac14+\tau_k^2}=\varpi_{\RVM}
\end{equation}
(ordered product). Unconditionally $\lim_{\sigma\to+\infty}\Pi_{\rm tr}(\sigma)=\varpi_{Z1}/\varpi_{\RVM}=\varpi_{Z1}L\ge L$, and the equality $\lim_{\sigma\to+\infty}\Pi_{\rm tr}(\sigma)=L=1/\varpi_{\RVM}$ is a fourth condition equivalent to the three above.
\end{corollary}

\begin{proof}
Theorems~\CGref{thm:Hamburger-Z1}, \CGref{thm:Z1-clock} and~\CGref{thm:varpi-Z1}, and $\Pi_{\rm tr}=\zeta_{Z1}/\zeta_{\RVM}\to\varpi_{Z1}/\varpi_{\RVM}$ by Theorems~\CGref{thm:halfplane-RVM} and~\CGref{thm:Z1-clock}(ii).
\end{proof}

\begin{proposition}[Boundedness of the endpoint]\label{prop:Z1-bounded}
Let $\sigma_0>1$.  Then $|\zeta_{Z1}|$ is bounded above and below on $\Re s\ge\sigma_0$ if and only if
\begin{equation}\label{eq:conjugate-poisson}
 \sup_{\sigma'\ge\sigma_0-\frac12,\ t\ge0}\Bigl|\Re\int_{\tau_0}^\infty\frac{\mathfrak e_1(\tau)+\mathfrak s_1(\tau)}{(\tau-t+i\sigma')^2}\,d\tau\Bigr|<\infty .
\end{equation}
\end{proposition}

\begin{proof}
$\tau+iz=\tau-t+i\sigma'$.  For $t\ge0$ the term with $(\tau-iz)^{-2}$ in \CGeqref{eq:Z1-clock} is bounded by \CGeqref{eq:frak-e1-bound}, since $|\tau-iz|\ge\tau+t$; the separate integral $\int_{\tau_0}^\infty\mathfrak s_1(\tau)(\tau+iz)^{-2}d\tau$ is bounded directly by Lemma~\CGref{lem:saw-primitive} and $\int_{\tau_0}^\infty|\tau+iz|^{-2}d\tau\ll1/\sigma'$.  By \CGeqref{eq:Z1-clock}, $\log|\zeta_{Z1}|$ is bounded on $\Re s\ge\sigma_0$, $t\ge0$, exactly when \CGeqref{eq:conjugate-poisson} holds, and $\zeta_{Z1}(\bar s)=\overline{\zeta_{Z1}(s)}$ covers $t<0$.
\end{proof}

The function in \CGeqref{eq:conjugate-poisson} is, after one integration by parts, the conjugate Poisson integral at distance $\sigma'$ of the count difference $N_\zeta-n_{\RVM}$; it depends only on the displacements $\gamma_k-\tau_k$.

\subsection{The Euler products of \texorpdfstring{$\zeta_{\RVM}$}{zeta\_RVM} and \texorpdfstring{$\zeta_{Z1}$}{zeta\_Z1}}
\label{subsec:euler-products}

\begin{theorem}[Prime-power counts as Mellin transforms]\label{thm:mellin-counts}
Let $X\in\{\RVM,Z1\}$.  There is a unique function $J_X^{\rm M}$ on $(1,\infty)$ with $J_X^{\rm M}(x)x^{-c}\in L^2((1,\infty),dx/x)$ for every $c>1$ and
\begin{equation}\label{eq:mellin-count}
 \log\zeta_X(s)=s\int_1^\infty J_X^{\rm M}(x)\,x^{-s-1}\,dx\qquad(\Re s>1),
\end{equation}
the branch of $\log\zeta_X$ being real on $(1,\infty)$.  We use the notation $\int x^{-s}\,dJ_X^{\rm M}:=s\int_1^\infty J_X^{\rm M}(x)x^{-s-1}dx$ distributionally; no bounded-variation representative is asserted.  With this convention $\zeta_X(s)=\exp\bigl(\int x^{-s}\,dJ_X^{\rm M}(x)\bigr)$.
\end{theorem}

\begin{proof}
$\zeta_X$ has no zeros or poles in $\Re s>1$, so $\log\zeta_X$ is holomorphic there.  For $c>1$, Theorem~\CGref{thm:halfplane-RVM} ($X=\RVM$) and Theorem~\CGref{thm:Z1-clock} ($X=Z1$), together with continuity on the compact part, give $|\log\zeta_X(s)|\le C_c(1+\log(|t|+2))$ on $\Re s\ge c$.  Hence $\mathsf F(s):=\log\zeta_X(s)/s$ is holomorphic on $\Re s>c$ and satisfies $\sup_{\sigma>c}\int_{\R}|\mathsf F(\sigma+it)|^2dt\le B_c^2<\infty$: $\mathsf F(\,\cdot+c)$ belongs to the Hardy space $H^2$ of the right half-plane.  By the Paley--Wiener theorem for $H^2$ of a half-plane \cite[Thm.~19.2]{Rudin1987} there is $g\in L^2(0,\infty)$ with $\mathsf F(s)=\int_0^\infty g(u)e^{-(s-c)u}du$ for $\Re s>c$, the integral converging absolutely.  Put $j(u):=g(u)e^{cu}$ for $u>0$; then
\begin{equation}\label{eq:mellin-count-laplace}
 \frac{\log\zeta_X(s)}s=\int_0^\infty j(u)\,e^{-su}\,du\qquad(\Re s>c),\qquad e^{-cu}j\in L^2(0,\infty).
\end{equation}
The function $j$ does not depend on $c$: if $c'>c$, the representation obtained from $c'$ agrees with \CGeqref{eq:mellin-count-laplace} on $\Re s>c'$, and the Laplace transform is injective on functions with $e^{-c'u}j\in L^2$ (the Fourier transform of $e^{-\sigma u}j(u)\mathbf 1_{u>0}$ is $\mathsf F(\sigma+i\cdot)$, and the Fourier--Plancherel transform is injective).  Define $J_X^{\rm M}(e^u):=j(u)$ for $u>0$, i.e.\ $J_X^{\rm M}(x)=j(\log x)$ on $(1,\infty)$; the substitution $x=e^u$ turns \CGeqref{eq:mellin-count-laplace} into \CGeqref{eq:mellin-count}, and $e^{-cu}j\in L^2(0,\infty)$ is the stated condition $J_X^{\rm M}(x)x^{-c}\in L^2((1,\infty),dx/x)$ for every $c>1$.  Uniqueness is the injectivity just used.  (The earlier formulation of this step through Cauchy's theorem on test functions was not a proof, since the cut-off $\hat\psi(\Im s)$ is not holomorphic in $s$; the Hardy-space argument replaces it.)
\end{proof}

\begin{theorem}[Distributional explicit formula]\label{thm:distributional-EF}
Let $X\in\{\RVM,Z1\}$, with ordinates $\kappa_k=\tau_k$ for $X=\RVM$ and $\kappa_k=\gamma_k$ for $X=Z1$.  In $\mathcal D'(1,\infty)$,
\begin{equation}\label{eq:distributional-EF}
 (\log x)\,\frac{dJ_X^{\rm M}}{dx}=1-\frac2{\sqrt x}\,\Re\sum_{k\ge1}x^{i\kappa_k}-\frac1{x(x^2-1)},
\end{equation}
the series converging in $\mathcal D'(1,\infty)$.  The explicit-formula count
\begin{equation}\label{eq:J-X}
 J_X(x):=\li(x)-2\Re\sum_{k\ge1}\Ei\Bigl(\bigl(\tfrac12+i\kappa_k\bigr)\log x\Bigr)-\log2+\int_x^\infty\frac{dy}{y(y^2-1)\log y}
\end{equation}
converges in $\mathcal D'(1,\infty)$ ($J_{\RVM}$ is the function \CGeqref{eq:J-RVM}), and $J_X^{\rm M}-J_X$ is a constant $C_X$ on $(1,\infty)$.
\end{theorem}

\begin{proof}
With $\rho_k=\frac12+i\kappa_k$, the factorization $\zeta_X=P_X/G$, $P_X'/P_X=\sum_k\bigl(\frac1{s-\rho_k}+\frac1{s-\bar\rho_k}\bigr)$, and the digamma expansion $\frac12\frac{\Gamma'}{\Gamma}(\frac s2)+\frac1s=-\frac{\gamma_E}2+\sum_{n\ge1}\bigl(\frac1{2n}-\frac1{s+2n}\bigr)$ give, for $\Re s>1$,
\begin{equation}\label{eq:log-derivative-X}
 -\frac{\zeta_X'}{\zeta_X}(s)=\frac1{s-1}-\sum_{k\ge1}\Bigl(\frac1{s-\rho_k}+\frac1{s-\bar\rho_k}\Bigr)+\sum_{n\ge1}\Bigl(\frac1{2n}-\frac1{s+2n}\Bigr)-\frac{\gamma_E}2-\frac12\log\pi .
\end{equation}
Each term is a Mellin--Stieltjes transform $\int_{[1,\infty)}x^{-s}d\mu$: $dx\mapsto\frac1{s-1}$, $x^{\rho-1}dx\mapsto\frac1{s-\rho}$, and the sum over the trivial zeros is kept as a whole, $\sum_n\bigl(\frac1{2n}-\frac1{s+2n}\bigr)=\int_1^\infty\frac{(1-x^{-s})\,dx}{x(x^2-1)}$, a tempered distribution in $u=\log x$ whose restriction to $(1,\infty)$ is the density $-1/(x(x^2-1))$ (the two pieces $\sum\frac1{2n}$ and $x^{-2n-1}dx$ are not separately summable, nor is $1/(x(x^2-1))$ integrable at $x=1$, so the split is made only on $(1,\infty)$); $\delta_1\mapsto1$ accounts for the constant term.  After multiplication by $x^{-c}$ and the change $u=\log x$, the Fourier transforms of the partial sums over $k$ and over $n$ converge locally uniformly with polynomial bounds, so the series converge in $\mathcal S'$, and on $(1,\infty)$ the limit measure is the right side of \CGeqref{eq:distributional-EF} times $dx$.  Differentiating \CGeqref{eq:mellin-count} in $s$ shows that $-\zeta_X'/\zeta_X$ is the distributional Mellin transform of $(\log x)\,dJ_X^{\rm M}$; injectivity gives \CGeqref{eq:distributional-EF}.  The series over $k$ converges in $\mathcal D'(1,\infty)$ because $\int\varphi(x)x^{-1/2+i\kappa}dx$ is rapidly decreasing in $\kappa$ for $\varphi\in C_c^\infty(1,\infty)$ and $\kappa_k\asymp k/\log k$.  By Lemma~\CGref{lem:Ei}, $\Ei(\rho_ku)=e^{\rho_ku}/(\rho_ku)+\varrho_k$ with $\rho_k=\tfrac12+i\kappa_k$ and $\sum_k|\varrho_k|$ locally uniformly convergent, so \CGeqref{eq:J-X} converges in $\mathcal D'(1,\infty)$ for the same reason; differentiating termwise with $\frac d{dx}\Ei(\rho\log x)=x^{\rho-1}/\log x$, $(\log x)J_X'$ equals the right side of \CGeqref{eq:distributional-EF}.  A distribution on an interval with zero derivative is constant.
\end{proof}

\begin{proof}[Proof of Theorem~\CGref{thm:MS}]
Take $X=\RVM$ in Theorems~\CGref{thm:mellin-counts} and~\CGref{thm:distributional-EF}.  By Theorem~\CGref{thm:J-exists}, $J_{\RVM}$ is continuous on $(1,\infty)$, so $J_{\RVM}^{\rm M}=J_{\RVM}+C_{\RVM}$ pointwise, and $J_{\RVM}^{*}=J_{\RVM}$ on $[2,\infty)$.
\end{proof}

\begin{theorem}[Cell form of the endpoint identity]\label{thm:cell-form}
The following are equivalent.
\begin{enumerate}[label=\textup{(\roman*)}]
\item (D), that is, $\zeta_{Z1}=\zeta$.
\item $J_{Z1}^{\rm M}$ is constant on every interval $(n,n+1)$, $n\ge1$, and its jumps $\mathsf b_n:=v_n-v_{n-1}$, where $v_n$ is its value on $(n,n+1)$ and $v_0:=0$, satisfy $\sum_{n\le X}|\mathsf b_n|=O_\eps(X^{1+\eps})$ for every $\eps>0$.
\item For every $n\ge1$, in $\mathcal D'(n,n+1)$,
\begin{equation}\label{eq:cell-identity}
 2\,\Re\sum_{k\ge1}\bigl(x^{i\gamma_k}-x^{i\tau_k}\bigr)=\sqrt x\,\log x\;J_{\RVM}'(x),
\end{equation}
and the jumps of $J_{Z1}^{\rm M}$ satisfy the bound in (ii).
\end{enumerate}
In this case $J_{Z1}^{\rm M}$ is Riemann's prime-power count $J(x)=\sum_{p^m\le x}1/m$.
\end{theorem}

\begin{proof}
(ii)$\Rightarrow$(i): $|v_n|\le\sum_{m\le n}|\mathsf b_m|\ll n^{1+\eps}$, so for $\Re s>1$ summation by parts in \CGeqref{eq:mellin-count} gives $\log\zeta_{Z1}(s)=\sum_nv_n(n^{-s}-(n+1)^{-s})=\sum_{n\ge1}\mathsf b_nn^{-s}$, absolutely convergent by partial summation.  The exponential of an absolutely convergent Dirichlet series is again an absolutely convergent Dirichlet series, so $\zeta_{Z1}$ satisfies \CGeqref{eq:Z1-D}.

(i)$\Rightarrow$(ii): by Theorem~\CGref{thm:Hamburger-Z1}, $\zeta_{Z1}=\zeta$, and $\log\zeta(s)=s\int_1^\infty J(x)x^{-s-1}dx$ for $\Re s>1$; by uniqueness in Theorem~\CGref{thm:mellin-counts}, $J_{Z1}^{\rm M}=J$, which is constant between consecutive integers, with jumps $1/m$ at $p^m$ and $\sum_{n\le X}|\mathsf b_n|=J(X)\ll X$.

(ii)$\Leftrightarrow$(iii): Theorem~\CGref{thm:distributional-EF} for $X=Z1$ and $X=\RVM$, with $dJ_{\RVM}^{\rm M}/dx=J_{\RVM}'$, gives in $\mathcal D'(1,\infty)$
\begin{equation}\label{eq:cell-difference}
 (\log x)\Bigl(\frac{dJ_{Z1}^{\rm M}}{dx}-J_{\RVM}'(x)\Bigr)=-\frac2{\sqrt x}\,\Re\sum_{k\ge1}\bigl(x^{i\gamma_k}-x^{i\tau_k}\bigr),
\end{equation}
so $dJ_{Z1}^{\rm M}/dx=0$ on $(n,n+1)$ is \CGeqref{eq:cell-identity}.
\end{proof}

\begin{remark}[The final step in cell form]\label{rem:cell-form}
Theorem~\CGref{thm:cell-form} is the exact form of the continuous-to-atomic passage of Remark~\CGref{rem:D-atomic}.  Both sides of \CGeqref{eq:cell-identity} are built from the midpoint construction: the right side from the clock alone ($J_{\RVM}'$ is bounded on $[2,\infty)$ by Theorem~\CGref{thm:unit-cell}, equal to $1/\log x$ plus the unit-periodic first harmonic and the harmonics with stationary points $2\pi x^{1/m}$, Section~\CGref{sec:integer-weights}), the left side from the displacements $\gamma_k-\tau_k$.  The final step of the method is the identity \CGeqref{eq:cell-identity} on every cell, equivalently \CGeqref{eq:product-identity}.  Numerically, the explicit formula \CGeqref{eq:J-X} with the $\gamma_k$ is constant between consecutive integers up to the truncation of the list (Experiment~\CGref{exp:primes}, Figure~\CGref{fig:prime-counts}), and the partial products of \CGeqref{eq:product-identity} approach $\varpi_{\RVM}$ (Experiment~\CGref{exp:product-identity}).
\end{remark}

\section{The atomized family and the programme for the final step}
\label{sec:atomized}

The continuous family $\zeta_\lambda$ keeps, for every $\lambda$, the functional equation and the zeros on the critical line.  Its atomization $\tilde\zeta_{\Z,\lambda}$, built from the cell masses of the prime coordinate, keeps the Euler product and the Dirichlet series.  The two are joined by a factor $e^{\tilde h_\lambda}$, and both families have a limit at $\lambda=0$.  Hamburger's theorem needs the functional equation and the Dirichlet series on one and the same function; at $\lambda=0$ this is the vanishing of the joining factor.  This section constructs the atomized family for every $\lambda$, proves its analytic properties, writes the joining factor cell by cell, and lists in Subsection~\CGref{subsec:programme} the steps that remain, in the order in which they will be followed.  The passage from a continuous prime system to a discrete one is classical for Beurling generalized primes \cite{Beurling1937,DiamondMontgomeryVorhauer2006,BrouckeVindas2024,DiamondZhang2016}; the atomization used here places the masses at the integers, which produces ordinary Dirichlet series.

Throughout, $\eta_k=\tau_k-\gamma_k$, $\gamma_k(\lambda)=(1-\lambda)\gamma_k+\lambda\tau_k$, $\rho_{k,\lambda}=\frac12+i\gamma_k(\lambda)$, and $\lambda'$ denotes an integration variable in $[0,1]$.

\subsection{The family at infinity}
\label{subsec:family-infinity}

Put $N_\lambda(\tau)=\#\{k:\gamma_k(\lambda)\le\tau\}$, $\mathfrak e^{(\lambda)}=N_\lambda-\NR$ and $\mathfrak e_1^{(\lambda)}(\tau)=\int_{\tau_0}^\tau\mathfrak e^{(\lambda)}$, so that $\mathfrak e_1^{(0)}=\mathfrak e_1$ and $\mathfrak e_1^{(1)}=-\mathfrak s_1$.

\begin{proposition}[Clock decomposition along the family]\label{prop:family-clock}
There is an absolute constant $C$ with $|\mathfrak e_1^{(\lambda)}(\tau)|\le C\log\tau$ for $\tau\ge\tau_0$ and all $\lambda\in[0,1]$.  Consequently Theorem~\CGref{thm:Z1-clock} holds for every $\zeta_\lambda$ with $\mathfrak e_1$ replaced by $\mathfrak e_1^{(\lambda)}$, uniformly in $\lambda$; in particular $\varpi_\lambda:=\lim_{\sigma\to+\infty}\zeta_\lambda(\sigma)$ exists and
\begin{equation}\label{eq:varpi-lambda}
 \varpi_\lambda=\varpi_{\RVM}\lim_{K\to\infty}\prod_{k\le K}\frac{\frac14+\tau_k^2}{\frac14+\gamma_k(\lambda)^2},\qquad \varpi_1=\varpi_{\RVM},\quad\varpi_0=\varpi_{Z1}.
\end{equation}
\end{proposition}

\begin{proof}
$\gamma_k(\lambda)$ lies between $\gamma_k$ and $\tau_k$, so $\gamma_1(\lambda)>\tau_0$.  For $X\ge\tau_0$,
\begin{equation}\label{eq:family-count-primitive}
 \int_{\tau_0}^X(N_\lambda-n_{\RVM})=\sum_k\bigl((X-\gamma_k(\lambda))^+-(X-\tau_k)^+\bigr)=(1-\lambda)\sum_{\gamma_k(\lambda),\tau_k\le X}\eta_k+O(\log X),
\end{equation}
the indices counted by only one of the two sums being $O(\log X)$ in number with ordinates in $[X-C_\eta,X]$ (Lemma~\CGref{lem:bounded-displacement}).  The case $\lambda=0$ of \CGeqref{eq:family-count-primitive} identifies the last sum with $\mathfrak e_1+\mathfrak s_1+O(\log X)=O(\log X)$, by \CGeqref{eq:frak-e1-bound} and Lemma~\CGref{lem:saw-primitive}.  Since $\mathfrak e_1^{(\lambda)}=\int(N_\lambda-n_{\RVM})-\mathfrak s_1$, the bound follows, and the proof of Theorem~\CGref{thm:Z1-clock} applies verbatim.
\end{proof}

Define the \emph{flow at infinity}
\begin{equation}\label{eq:Upsilon}
 \Upsilon(\lambda):=\sum_{k\ge1}\frac{2\gamma_k(\lambda)\,\eta_k}{\frac14+\gamma_k(\lambda)^2}\qquad(\text{summed in the order of }k).
\end{equation}

\begin{proposition}[The value at infinity along the family]\label{prop:family-flow}
The series \CGeqref{eq:Upsilon} converges uniformly on $[0,1]$, $\lambda\mapsto\log\varpi_\lambda$ is $C^1$ on $[0,1]$, and
\begin{equation}\label{eq:varpi-derivative}
 \frac{d}{d\lambda}\log\varpi_\lambda=-\Upsilon(\lambda)=\lim_{\sigma\to+\infty}\mathcal V_\lambda(\sigma),\qquad
 \Upsilon'(\lambda)=\sum_{k\ge1}\frac{2\eta_k^2\bigl(\frac14-\gamma_k(\lambda)^2\bigr)}{\bigl(\frac14+\gamma_k(\lambda)^2\bigr)^2}<0 ,
\end{equation}
with $\mathcal V_\lambda$ the generator of Theorem~\CGref{thm:generator}.  Moreover
\begin{equation}\label{eq:flow-integral}
 \int_0^1\Upsilon(\lambda)\,d\lambda=\sum_{k\ge1}\log\frac{\frac14+\tau_k^2}{\frac14+\gamma_k^2}=\log L+\log\varpi_{Z1},
\end{equation}
so that (D) holds if and only if $\int_0^1\Upsilon(\lambda)\,d\lambda=\log L$, and $\Upsilon(1)<\log L+\log\varpi_{Z1}<\Upsilon(0)$.
\end{proposition}

\begin{proof}
Put $g(u)=2u/(\frac14+u^2)$.  The series $\Upsilon(1)=\sum_kg(\tau_k)\eta_k$ converges by summation by parts against $N_\zeta-n_{\RVM}$, exactly as in \CGeqref{eq:count-difference-integral}.  Since $|g(\gamma_k(\lambda))-g(\tau_k)|\le C|\eta_k|\gamma_k^{-2}$ and $\sum_k\eta_k^2\gamma_k^{-2}<\infty$, the series $\Upsilon(\lambda)-\Upsilon(1)$ converges absolutely and uniformly; the same holds for the termwise derivative, which is the series in \CGeqref{eq:varpi-derivative}. Each term of that series is $\le0$ because $\gamma_k(\lambda)>\frac12$, and the term $k=1$ is $<0$ because $\eta_1\ne0$: the first ordinate satisfies $\gamma_1<14.14$ (its rigorous location $\gamma_1=14.1347\ldots$ is classical, \cite[Ch.~XV]{Titchmarsh1986}), whereas $\NR(14.14)=0.4499\ldots<\frac12=\NR(\tau_1)$ by \CGeqref{eq:clock}, so $\tau_1>14.14>\gamma_1$ by the monotonicity of $\NR$ on $[2\pi,\infty)$, and $\eta_1=\tau_1-\gamma_1>0$. Hence $\Upsilon'<0$; the formula alone gives only $\Upsilon'\le0$, with strict inequality exactly when $\boldsymbol\gamma\ne\boldsymbol\tau$.  Differentiating \CGeqref{eq:varpi-lambda} termwise gives $\frac d{d\lambda}\log\varpi_\lambda=-\Upsilon(\lambda)$, and the limit $\sigma\to+\infty$ in the generator \CGeqref{eq:generator} is justified by Abel summation of the ordered ordinate sum: its partial sums are controlled by the count-difference integral, while the multiplier $q/(A+q)$ is bounded and monotone in $A$.  This gives the same ordered series.  Integrating from $0$ to $1$ and using $\varpi_1=\varpi_{\RVM}$, $\varpi_0=\varpi_{Z1}$ gives \CGeqref{eq:flow-integral}; the last statement follows from Theorem~\CGref{thm:varpi-Z1} and the monotonicity of $\Upsilon$.
\end{proof}

\begin{proposition}[The value at infinity and the verified zeros]\label{prop:varpi-verified}
If every zero $\beta+i\gamma$ with $0<\gamma\le T_{\rm ver}$ has $\beta=\frac12$, then
\begin{equation}\label{eq:varpi-verified}
 0\le\log\varpi_{Z1}\le0.26\sum_{\gamma>T_{\rm ver}}\frac1{\gamma^2}.
\end{equation}
With $T_{\rm ver}=3\cdot10^{12}$ \cite{PlattTrudgian2021}, the right side is $\approx\frac{0.26}{2\pi T_{\rm ver}}\bigl(\log\frac{T_{\rm ver}}{2\pi}+1\bigr)<4\cdot10^{-13}$.
\end{proposition}

\begin{proof}
By the proof of Theorem~\CGref{thm:varpi-Z1}, $\log\varpi_{Z1}=\sum\log c_\gamma$ over the pairs $\{\rho,1-\bar\rho\}$ with $\rho(1-\rho)\ne\frac14+\gamma^2$, and all of them have $\gamma>T_{\rm ver}$.  For $|d_\rho|\le\frac12$ and $\gamma\ge T_{\rm ver}$, \CGeqref{eq:pair-ratio} gives $\log c_\gamma\le(2d_\rho^2\gamma^2+d_\rho^4)\gamma^{-4}\le0.51\,\gamma^{-2}$; a pair consists of two zeros, which gives \CGeqref{eq:varpi-verified}.  The numerical value follows from $N_\zeta(T)=\frac T{2\pi}\log\frac T{2\pi e}+O(\log T)$.
\end{proof}

\subsection{Atomization by smooth cells}
\label{subsec:smooth-cells}

Fix $\varphi_n\in C_c^\infty(n-1,n+1)$, $n\ge2$, with $\sum_{n\ge2}\varphi_n=1$ on $[2,\infty)$ and $\varphi_n(m)=\delta_{nm}$ for integers $m$.  For $\lambda\in[0,1]$ let $J_\lambda$ be the explicit-formula count \CGeqref{eq:J-X} with the ordinates $\gamma_k(\lambda)$ (a distribution on $(1,\infty)$ by Theorem~\CGref{thm:distributional-EF}, which applies to every $\zeta_\lambda$ by Theorem~\CGref{thm:Z1-clock}), $J^*_\lambda=J_\lambda\mathbf 1_{[2,\infty)}$, and $\pi_\lambda=\sum_{j\ge1}\frac{\mu(j)}jJ^*_\lambda(x^{1/j})$.

\begin{definition}[Atomized family]\label{def:atomized}
The cell masses, the atomized zeta function and its coefficients are
\begin{equation}\label{eq:atomized-family}
 \tilde c_n(\lambda):=\langle d\pi_\lambda,\varphi_n\rangle,\qquad
 \tilde\zeta_{\Z,\lambda}(s):=\prod_{n\ge2}\bigl(1-n^{-s}\bigr)^{-\tilde c_n(\lambda)}=\sum_{N\ge1}\frac{\tilde a_N(\lambda)}{N^s},
\end{equation}
\begin{equation}\label{eq:atomized-coefficients}
 \tilde a_N(\lambda)=\sum_{\prod_nn^{r_n}=N}\ \prod_{n:\,r_n\ge1}\binom{\tilde c_n(\lambda)+r_n-1}{r_n},\qquad \tilde a_1(\lambda)=1 .
\end{equation}
\end{definition}

Unlike the discretizations of Beurling systems, where the atoms are placed at chosen real points \cite{DiamondMontgomeryVorhauer2006,BrouckeVindas2024,DiamondZhang2016}, the masses of $d\pi_\lambda$ are placed at the integers with real exponents; measures with continuous part and non-integer multiplicities of this kind appear also in \cite[Ch.~3]{LapidusvanFrankenhuysen2000}.  Thus $\tilde a_p(\lambda)=\tilde c_p(\lambda)$ for primes $p$, $\tilde a_4=\tilde c_4+\frac12\tilde c_2(\tilde c_2+1)$, $\tilde a_6=\tilde c_6+\tilde c_2\tilde c_3$.  The pairing with a fixed smooth partition defines $\tilde c_n(\lambda)$ for every $\lambda$ without point values of $\pi_\lambda$, and $\lambda\mapsto\tilde c_n(\lambda)$ is continuous.  If $d\pi_0=\sum_p\delta_p$, then $\tilde c_n(0)=\mathbf 1_{\rm prime}(n)$ and $\tilde a_N(0)=1$ for all $N$.  Along the family, in $\mathcal D'(1,\infty)$,
\begin{equation}\label{eq:J-lambda-derivative}
 \frac{\partial J_\lambda}{\partial\lambda}(x)=2\sum_{k\ge1}\eta_k\,\Im\frac{x^{\rho_{k,\lambda}}}{\rho_{k,\lambda}},\qquad
 \frac{\partial}{\partial\lambda}\Bigl((\log x)\frac{dJ_\lambda}{dx}\Bigr)=\frac{2\log x}{\sqrt x}\sum_{k\ge1}\eta_k\sin\bigl(\gamma_k(\lambda)\log x\bigr).
\end{equation}

\subsection{The joining factor cell by cell}
\label{subsec:joining-factor}

Define the \emph{atomization kernel} and the \emph{cell defects}
\begin{equation}\label{eq:atomization-kernel}
 \mathfrak a_s(x):=\sum_{n\ge2}\varphi_n(x)\,n^{-s}-x^{-s},\qquad
 \mathfrak h_n(\lambda;s):=\bigl\langle d\pi_\lambda,\ \varphi_n(x)\bigl(n^{-s}-x^{-s}\bigr)\bigr\rangle .
\end{equation}
$\mathfrak a_s$ vanishes at every integer $\ge2$, and $\mathfrak h_n(\lambda;s)$ measures how far the mass of the cell $n$ is from being concentrated at $n$.

The explicit-formula count $J_\lambda$ need not vanish on $(1,2)$, whereas $\pi_\lambda$ is built from $J^*_\lambda=J_\lambda\mathbf 1_{[2,\infty)}$.  The first cell $(1,2)$ enters the joining factor through
\begin{equation}\label{eq:first-cell}
 \mathfrak u_\lambda(s):=C_\lambda\,2^{-s}+s\int_1^2J^{\rm M}_\lambda(x)\,x^{-s-1}\,dx,\qquad C_\lambda:=J^{\rm M}_\lambda-J_\lambda ,
\end{equation}
where $J^{\rm M}_\lambda$ is the Mellin count of $\zeta_\lambda$ (Theorem~\CGref{thm:mellin-counts}, which applies to every $\zeta_\lambda$ by Proposition~\CGref{prop:family-clock}) and $C_\lambda$ is the constant of Theorem~\CGref{thm:distributional-EF}.

\begin{proposition}[The first cell]\label{prop:first-cell}
Let $\lambda\in[0,1]$.
\begin{enumerate}[label=\textup{(\roman*)}]
\item $\mathfrak u_\lambda$ is entire, and for $\Re s>1$
\begin{equation}\label{eq:first-cell-split}
 \log\zeta_\lambda(s)=\mathfrak u_\lambda(s)+\langle dJ^*_\lambda,x^{-s}\rangle,\qquad \langle dJ^*_\lambda,x^{-s}\rangle:=s\int_2^\infty J_\lambda(x)\,x^{-s-1}\,dx .
\end{equation}
\item $\mathfrak u_\lambda(\sigma+it)=o(|t|)$ as $|t|\to\infty$, uniformly for $\sigma$ in compact sets.
\item $\mathfrak u_\lambda$ is constant if and only if $J_\lambda=0$ on $(1,2)$; in that case $\mathfrak u_\lambda\equiv C_\lambda$.
\end{enumerate}
For Riemann's count $J(x)=\sum_{p^m\le x}1/m$, which vanishes on $(1,2)$, the corresponding term is $0$.
\end{proposition}

\begin{proof}
(i) $J^{\rm M}_\lambda(x)x^{-c}\in L^2((1,\infty),dx/x)$ for $c>1$, so $J^{\rm M}_\lambda\in L^1(1,2)$ and the integral in \CGeqref{eq:first-cell} is entire in $s$.  By \CGeqref{eq:mellin-count}, $\log\zeta_\lambda(s)=s\int_1^2J^{\rm M}_\lambda x^{-s-1}dx+s\int_2^\infty(J_\lambda+C_\lambda)x^{-s-1}dx$, and $s\int_2^\infty C_\lambda x^{-s-1}dx=C_\lambda2^{-s}$.  The second integral is the pairing of $x^{-s}$ with the distributional derivative of $J^*_\lambda$, which vanishes on $(0,2)$.
(ii) $C_\lambda2^{-s}$ is bounded on vertical strips, and $\int_1^2J^{\rm M}_\lambda(x)x^{-s-1}dx=\int_0^{\log2}J^{\rm M}_\lambda(e^u)e^{-\sigma u}e^{-itu}du\to0$ as $|t|\to\infty$ by the Riemann--Lebesgue lemma, uniformly for $\sigma$ in compact sets.
(iii) For a constant $c$ one has $c=c\,2^{-s}+s\int_1^2c\,x^{-s-1}dx$, so $\mathfrak u_\lambda(s)-c=s\int_1^\infty H(x)x^{-s-1}dx$ with $H=(J^{\rm M}_\lambda-c)\mathbf 1_{(1,2)}+(C_\lambda-c)\mathbf 1_{[2,\infty)}$.  By the injectivity used in Theorem~\CGref{thm:mellin-counts}, $\mathfrak u_\lambda\equiv c$ if and only if $H=0$ almost everywhere, that is, $C_\lambda=c$ and $J^{\rm M}_\lambda=C_\lambda$ on $(1,2)$, which is $J_\lambda=0$ on $(1,2)$.
\end{proof}

\begin{lemma}[Mean-square convergence of the ordinate series]\label{lem:L2-partial-sums}
Let $\boldsymbol{\kappa}$ be a list with $|\kappa_k-\tau_k|\le C_\eta$ (in particular $\boldsymbol{\kappa}=(\gamma_k(\lambda))$ for any $\lambda\in[0,1]$), $\rho_k:=\frac12+i\kappa_k$, and let $I\subset(1,\infty)$ be compact. The partial sums $S_K(x):=\sum_{k\le K}\bigl(x^{\rho_k}/\rho_k+x^{\bar\rho_k}/\bar\rho_k\bigr)$ and $E_K(x):=\sum_{k\le K}\bigl(\Ei(\rho_k\log x)+\Ei(\bar\rho_k\log x)\bigr)$ converge in $L^2(I)$ as $K\to\infty$, with
\begin{equation}\label{eq:L2-partial-sums}
 \|S_{K'}-S_K\|_{L^2(I)}^2+\|E_{K'}-E_K\|_{L^2(I)}^2\ \le\ C_I\sum_{K<k\le K'}\frac{\log(\kappa_k+2)}{\kappa_k^2}\qquad(K<K'),
\end{equation}
where $C_I$ depends only on $I$ and $C_\eta$. Consequently: \textup{(a)} for every Lipschitz $\psi$ with compact support in $(1,\infty)$, the pairing $\langle dJ_X^{\rm M},\psi\rangle:=-\int J_X^{\rm M}\psi'\,dx$ equals the series obtained by integrating \CGeqref{eq:distributional-EF} term by term, which converges; if moreover $\psi'$ is of bounded variation, the series converges absolutely; \textup{(b)} $\lambda\mapsto J_\lambda$ and $\lambda\mapsto\partial_\lambda J_\lambda$ are continuous from $[0,1]$ to $L^2(I)$; \textup{(c)} the constant $C_\lambda:=J^{\rm M}_\lambda-J_\lambda$ is continuous on $[0,1]$, hence $\sup_\lambda|C_\lambda|<\infty$.
\end{lemma}

\begin{proof}
In $u=\log x$, $u\in\log I=:[u_1,u_2]$, the difference $S_{K'}-S_K$ is $e^{u/2}\sum_{K<k\le K'}\bigl(e^{i\kappa_ku}/\rho_k+e^{-i\kappa_ku}/\bar\rho_k\bigr)$. Let $\omega\in C_c^\infty(\mathbb R)$ be $\ge1$ on $[u_1,u_2]$. Expanding the square and using $|\hat\omega(\xi)|\ll_N(1+|\xi|)^{-N}$,
\begin{equation}\label{eq:L2-partial-sums-proof}
\begin{split}
 \int\omega(u)\Bigl|\sum_{K<k\le K'}\frac{e^{i\kappa_ku}}{\rho_k}\Bigr|^2du
 &\ll\sum_{K<k,k'\le K'}\frac{(1+|\kappa_k-\kappa_{k'}|)^{-2}}{\kappa_k\kappa_{k'}}\\
 &\ll\sum_{K<k\le K'}\frac1{\kappa_k^2}\sum_{j\ge0}\frac{\#\{k':j\le|\kappa_{k'}-\kappa_k|<j+1\}}{(1+j)^2}
 \ \ll\sum_{K<k\le K'}\frac{\log(\kappa_k+2)}{\kappa_k^2},
\end{split}
\end{equation}
because the tube gives $\#\{k':|\kappa_{k'}-t|\le1\}\ll\log(t+2)$ uniformly (every unit interval contains $O(\log t)$ clock points after widening by $C_\eta$), $\kappa_{k'}\asymp\kappa_k$ on the pairs that matter, and the pairs with $|\kappa_k-\kappa_{k'}|>1$ are summed with the weight $(1+|\kappa_k-\kappa_{k'}|)^{-2}$ against the same local count. The same bound holds for the $\Ei$ sums by Lemma~\CGref{lem:Ei}, since $\Ei(\rho u)=e^{\rho u}/(\rho u)+\varrho_k$ with $\sum_k|\varrho_k|$ locally uniformly convergent and $u\ge u_1>0$. This is \CGeqref{eq:L2-partial-sums}, and the right side tends to $0$ as $K\to\infty$ because $\kappa_k\asymp k/\log k$.

(a) For $\psi\in C_c^\infty(1,\infty)$ the identity is the convergence in $\mathcal D'(1,\infty)$ established in Theorem~\CGref{thm:distributional-EF}. For Lipschitz $\psi$ with compact support in $(1,\infty)$, $\psi'\in L^\infty\subset L^2(I)$, and both sides are continuous in the $L^2(I)$ topology of the partial sums: the left side is the pairing of an $L^2(I)$ limit with $\psi'$, and the terms of the right side are the pairings of the summands with $\psi'$, whose partial sums converge by \CGeqref{eq:L2-partial-sums} and the Cauchy--Schwarz inequality.  Lipschitz continuity alone gives convergence, not absolute convergence: the $k$-th term is $\int\psi'(x)\,\Re(x^{\rho_k}/\rho_k)\,dx$, of size $O(1/\kappa_k)$ in general.  When $\psi'$ is of bounded variation, one integration by parts in $u=\log x$ (against $e^{i\kappa_ku}$, with $\psi'$ absorbing the boundary terms) gains a factor $1/\kappa_k$ and the terms become $O(\kappa_k^{-2}\,\mathrm{TV}(\psi'))$, so the series converges absolutely; this is the case for the test functions $\mathfrak a_{ms}(y^j)\mathbf 1_{y\ge2}$ used below, whose derivatives are piecewise smooth.

(b) With $\boldsymbol{\kappa}=(\gamma_k(\lambda))$ and $\boldsymbol{\kappa}'=(\gamma_k(\lambda'))$, $|\gamma_k(\lambda)-\gamma_k(\lambda')|=|\lambda-\lambda'||\eta_k|\le C_\eta|\lambda-\lambda'|$; the difference of the $k$-th summands is $O(|\lambda-\lambda'|\,\log x\,x^{1/2}/\kappa_k)$ termwise, and the tail beyond $K$ is controlled uniformly in $\lambda$ by \CGeqref{eq:L2-partial-sums}. Hence $J_\lambda\to J_{\lambda'}$ in $L^2(I)$ as $\lambda\to\lambda'$; the same argument with the differentiated summands $\eta_k\,\partial_\kappa(x^{\rho}/\rho)$ gives the continuity of $\partial_\lambda J_\lambda$.

(c) By Plancherel on a fixed line $\Re s=c>1$ (proof of Proposition~\CGref{prop:cell-limit}), $\lambda\mapsto J^{\rm M}_\lambda$ is continuous into the weighted $L^2$, hence into $L^2(I)$. Pairing $J^{\rm M}_\lambda-J_\lambda$ with a fixed $\phi\in C_c^\infty(2,3)$, $\int\phi=1$, gives $C_\lambda$ as a difference of two continuous functions of $\lambda$.
\end{proof}

\begin{theorem}[The joining factor]\label{thm:joining-factor}
Let $\tilde h_\lambda:=\log\tilde\zeta_{\Z,\lambda}-\log\zeta_\lambda$ on the half-plane where $\tilde\zeta_{\Z,\lambda}$ converges absolutely (Theorem~\CGref{thm:atomized-continuation}).  Then
\begin{equation}\label{eq:joining-factor}
 \tilde\zeta_{\Z,\lambda}=\zeta_\lambda\,e^{\tilde h_\lambda},\qquad
 \tilde h_\lambda(s)=-\mathfrak u_\lambda(s)+\sum_{m\ge1}\frac1m\sum_{n\ge2}\mathfrak h_n(\lambda;ms),
\end{equation}
and, explicitly,
\begin{equation}\label{eq:joining-explicit}
 \tilde h_\lambda(s)=-\mathfrak u_\lambda(s)+\sum_{m\ge1}\frac1m\sum_{j\ge1}\frac{\mu(j)}j\int_2^\infty\mathfrak a_{ms}(y^j)\,\frac1{\log y}\Bigl(1-\frac2{\sqrt y}\Re\sum_{k\ge1}y^{i\gamma_k(\lambda)}-\frac1{y(y^2-1)}\Bigr)dy,
\end{equation}
the series over $k$ converging absolutely after integration.  Its derivative along the family is
\begin{equation}\label{eq:joining-derivative}
 \partial_\lambda\tilde h_\lambda(s)=-\partial_\lambda\mathfrak u_\lambda(s)+\sum_{m\ge1}\frac1m\sum_{j\ge1}\frac{\mu(j)}j\int_2^\infty\mathfrak a_{ms}(y^j)\,\frac2{\sqrt y}\sum_{k\ge1}\eta_k\sin\bigl(\gamma_k(\lambda)\log y\bigr)\,dy .
\end{equation}
Here $\partial_\lambda\mathfrak u_\lambda(s)=\mathcal V_\lambda(s)-\partial_\lambda\langle dJ^*_\lambda,x^{-s}\rangle$ by \CGeqref{eq:first-cell-split} and \CGeqref{eq:generator-zeta}.  Moreover $\mathfrak u_\lambda(\sigma)\to\log\varpi_\lambda$ as $\sigma\to+\infty$, and $\mathfrak u_\lambda\equiv\log\varpi_\lambda$ if $J_\lambda$ vanishes on $(1,2)$.
\end{theorem}

\begin{proof}
By the Euler product, $\log\tilde\zeta_{\Z,\lambda}(s)=\sum_{m\ge1}\frac1m\sum_n\tilde c_n(\lambda)n^{-ms}$.  By Proposition~\CGref{prop:first-cell}, $\log\zeta_\lambda(s)=\mathfrak u_\lambda(s)+\langle dJ^*_\lambda,x^{-s}\rangle$; finite inversion (Proposition~\CGref{prop:finite-inversion}, whose proof applies to $J^*_\lambda$) gives $dJ^*_\lambda=\sum_m\frac1m\,d[\pi_\lambda(x^{1/m})]$, and $\langle d[\pi_\lambda(x^{1/m})],x^{-s}\rangle=\langle d\pi_\lambda,y^{-ms}\rangle$.  Inserting $1=\sum_n\varphi_n$ on $[2,\infty)$ gives $\langle d\pi_\lambda,y^{-ms}\rangle=\sum_n\bigl(\tilde c_n(\lambda)n^{-ms}-\mathfrak h_n(\lambda;ms)\bigr)$, hence \CGeqref{eq:joining-factor}.  Letting $\sigma\to+\infty$, $\log\tilde\zeta_{\Z,\lambda}\to0$ (because $\tilde a_1=1$), $\log\zeta_\lambda\to\log\varpi_\lambda$ (Proposition~\CGref{prop:family-clock}) and the double series tends to $0$ by the bounds of Theorem~\CGref{thm:atomized-continuation}; so $\mathfrak u_\lambda(\sigma)\to\log\varpi_\lambda$, and a constant $\mathfrak u_\lambda$ (Proposition~\CGref{prop:first-cell}(iii)) equals $\log\varpi_\lambda$.  Atoms of $d\pi_\lambda$ at integers contribute nothing because the test functions vanish there.  Formula \CGeqref{eq:joining-explicit} is \CGeqref{eq:joining-factor} with $d\pi_\lambda$ written by \CGeqref{eq:distributional-EF} and the substitution $x=y^j$; the integrated terms in $k$ are $O_{m,j,s}(\gamma_k(\lambda)^{-2})$ because $\mathfrak a_{ms}(2)=0$, and the termwise integration against the Lipschitz test $\mathfrak a_{ms}(y^j)\mathbf 1_{y\ge2}$ is justified by Lemma~\CGref{lem:L2-partial-sums}(a).  The remaining double summation is absolutely controlled by $\sum_{m\ge1}m^{-1}\sum_{j\ge1}j^{-1}(mj|s|)^2 2^{-jm\sigma}<\infty$, so the exchanges of the $m$- and $j$-sums with the integrated ordinate sum are justified.  Differentiating in $\lambda$ with \CGeqref{eq:J-lambda-derivative} gives \CGeqref{eq:joining-derivative}; $\mathfrak u_\lambda(s)$ is differentiable in $\lambda$ because $J^{\rm M}_\lambda$ is the Fourier--Plancherel transform of $\log\zeta_\lambda(\sigma+it)/(\sigma+it)$, whose $\lambda$-derivative $\mathcal V_\lambda(\sigma+it)/(\sigma+it)$ is bounded and continuous in $L^2(\R)$ locally uniformly in $\lambda$, because $|\mathcal V_\lambda(\sigma+it)|\ll_\sigma\log^2(|t|+2)$ uniformly in $\lambda$ for fixed $\sigma>1$ (the $O(\log t)$ ordinates with $|\gamma_k(\lambda)-t|\le1$ contribute $O(\log t)$ each through $|\eta_k|\le C_\eta$, and the others $O(\log^2t)$ in total by summation against the count), and the limit $\lambda\to\lambda'$ is dominated and dominated convergence on the ordinate sum.  Hence $\lambda\mapsto J^{\rm M}_\lambda$ is $C^1$ in the corresponding weighted $L^2$ space; the constant $C_\lambda=J^{\rm M}_\lambda-J_\lambda$ is obtained by pairing with one fixed test function and is $C^1$ as well.  In particular $\|\partial_\lambda J_\lambda\|_{L^2(2,3)}$ is locally uniformly bounded.
\end{proof}

\begin{proposition}[Explicit cell defects]\label{prop:cell-explicit}
For every $n\ge2$, $\lambda\in[0,1]$ and $\Re s>1$,
\begin{equation}\label{eq:cell-explicit}
 \mathfrak h_n(\lambda;s)=\sum_{j\ge1}\frac{\mu(j)}j\int_2^\infty\varphi_n(y^j)\bigl(n^{-s}-y^{-js}\bigr)\,\frac1{\log y}\Bigl(1-\frac2{\sqrt y}\Re\sum_{k\ge1}y^{i\gamma_k(\lambda)}-\frac1{y(y^2-1)}\Bigr)dy,
\end{equation}
\begin{equation}\label{eq:cell-explicit-derivative}
 \partial_\lambda\mathfrak h_n(\lambda;s)=\sum_{j\ge1}\frac{\mu(j)}j\int_2^\infty\varphi_n(y^j)\bigl(n^{-s}-y^{-js}\bigr)\,\frac2{\sqrt y}\sum_{k\ge1}\eta_k\sin\bigl(\gamma_k(\lambda)\log y\bigr)\,dy ,
\end{equation}
the sums over $j$ being finite and the sums over $k$ converging absolutely after integration, because the test functions are smooth with compact support.  At $\lambda=1$ the cell defect $\mathfrak h_n(1;\cdot)$ is the dispersion of the cell $n$ for $\zeta_{\RVM}$, computed from the clock alone; the absolutely continuous part of $d\pi_{\RVM}$ has bounded density there by Theorem~\CGref{thm:unit-cell}, while the atoms at $2^j$ are treated separately, so $\mathfrak h_n(1;\cdot)\equiv0$ holds only if this density vanishes on the support of $\varphi_n$.
\end{proposition}

\begin{proof}
Insert \CGeqref{eq:distributional-EF} for $d\pi_\lambda=\sum_j\frac{\mu(j)}j\,d[J^*_\lambda(x^{1/j})]$ into the definition \CGeqref{eq:atomization-kernel} and substitute $x=y^j$; the atoms of $d\pi_\lambda$ at integers, for instance the jump of $J^*_\lambda$ at $2$, contribute nothing because $n^{-s}-x^{-s}$ vanishes at $x=n$ and $\varphi_n$ vanishes at the other integers.  Differentiation in $\lambda$ uses \CGeqref{eq:J-lambda-derivative}; the factor $\log y$ cancels.  For the last statement, $\mathfrak h_n(1;\cdot)\equiv0$ forces $\varphi_n\,d\pi_{\RVM}=\tilde c_n(1)\delta_n$ by the injectivity used in Proposition~\CGref{prop:cell-limit}, and, away from the explicit atoms at $2^j$, an absolutely continuous measure equals a multiple of $\delta_n$ only if both vanish; at the powers of $2$ the boundary atom is handled by \CGeqref{eq:cell-mass-power-two}.
\end{proof}

In the form of the path, the final step requires that in every cell the flow of the displacements cancel exactly the dispersion of the clock, \CGeqref{eq:closing-cells}.

\begin{proposition}[The limit $\lambda\to0$ and the meaning of a vanishing cell]\label{prop:cell-limit}
For every $n\ge2$ and $\Re s>1$, $\mathfrak h_n(\lambda;s)\to\mathfrak h_n(0;s)$ and $\tilde c_n(\lambda)\to\tilde c_n(0)$ as $\lambda\to0^+$, and
\begin{equation}\label{eq:cell-flow}
 \mathfrak h_n(0;s)=\mathfrak h_n(1;s)-\int_0^1\partial_\lambda\mathfrak h_n(\lambda;s)\,d\lambda .
\end{equation}
Moreover $\mathfrak h_n(0;s)=0$ for all $\Re s>1$ if and only if $\varphi_n\,d\pi_{Z1}=\tilde c_n(0)\,\delta_n$.
\end{proposition}

\begin{proof}
The test function $\varphi_n(x)(n^{-s}-x^{-s})$ is smooth with compact support, and the integrated terms over $k$ are dominated uniformly in $\lambda$ because $|\gamma_k(\lambda)-\gamma_k|\le C_\eta$; dominated convergence gives the limits for cells away from the powers of $2$.  For $n=2^j$, the boundary term \CGeqref{eq:cell-mass-power-two} also converges: by Plancherel, $\log\zeta_\lambda(\sigma+it)/(\sigma+it)\to\log\zeta_{Z1}(\sigma+it)/(\sigma+it)$ in $L^2(dt)$ on a fixed line $\sigma>1$, hence $J^{\rm M}_\lambda\to J^{\rm M}_{Z1}$ in the corresponding weighted $L^2$ space, and therefore $J_\lambda\to J_{Z1}$ in $L^2(2,3)$ by Lemma~\CGref{lem:L2-partial-sums}(b),(c).  This proves all stated limits and \CGeqref{eq:cell-flow}.  Since $\mathfrak h_n(0;s)=\tilde c_n(0)n^{-s}-\langle\varphi_n\,d\pi_{Z1},x^{-s}\rangle$, its vanishing for all $s$ says that the compactly supported distribution $\varphi_n\,d\pi_{Z1}$ has the Mellin transform of $\tilde c_n(0)\delta_n$; the Mellin transform is injective on such distributions.
\end{proof}

\subsection{Mean square of the cell masses}
\label{subsec:mean-square}

From now on the cells are translates of one function: $\varphi_n(x)=\varphi(x-n)$, where $\varphi=\beta_0/\sum_{m\in\Z}\beta_0(\cdot-m)$ for a fixed bump $\beta_0\in C_c^\infty(-1,1)$ positive on $[-\frac12,\frac12]$.  Then $\sum_{n\ge2}\varphi_n=1$ on $[2,\infty)$, $\varphi_n(m)=\delta_{nm}$, and $\|\varphi_n^{(A)}\|_\infty=\|\varphi^{(A)}\|_\infty=:C_A$ for every $A\ge0$, uniformly in $n$.  For $n\ge2$ put
\begin{equation}\label{eq:VW}
 V_n(\gamma):=\sum_{j\ge1}\frac{\mu(j)}j\int_2^\infty\varphi_n(y^j)\,y^{-1/2}e^{i\gamma\log y}\,dy,\qquad
 W_n(\gamma):=\sum_{j\ge1}\frac{\mu(j)}j\int_2^\infty\varphi_n(y^j)\,\frac{y^{-1/2}e^{i\gamma\log y}}{\log y}\,dy ,
\end{equation}
so that $W_n'=iV_n$; both are rapidly decreasing in $\gamma$ when $n$ is not a power of $2$ (for $n=2^j$ the term $j$ carries the cut of $J^*_\lambda$ at $y=2$ and is treated in the proof of Theorem~\CGref{thm:mean-square-cells}).  For $n$ not a power of $2$, the density of $d\pi_\lambda$ used in \CGeqref{eq:cell-explicit} and the derivative \CGeqref{eq:J-lambda-derivative} give, with $\eta_k=\tau_k-\gamma_k$:
\begin{equation}\label{eq:cell-mass-W}
 \tilde c_n(\lambda)=\mathfrak m_n-2\,\Re\sum_{k\ge1}W_n\bigl(\gamma_k(\lambda)\bigr),\qquad
 \partial_\lambda\tilde c_n(\lambda)=2\sum_{k\ge1}\eta_k\,\Im V_n\bigl(\gamma_k(\lambda)\bigr),
\end{equation}
where $\mathfrak m_n$ collects the terms that do not involve the ordinates and is $O(1)$.  If $n=2^j$, the $j$th M\"obius-inversion term crosses the cut at $y=2$ and contributes instead the additional boundary term
\begin{equation}\label{eq:cell-mass-power-two}
 -\frac{\mu(j)}j\int_2^3 J_\lambda(y)\,\frac d{dy}\varphi_n(y^j)\,dy,
\end{equation}
whose $\lambda$-derivative is obtained by replacing $J_\lambda$ by $\partial_\lambda J_\lambda$.

\begin{lemma}[Mean value over the moving ordinates]\label{lem:mean-value-ordinates}
There is $C$ such that for all $\lambda\in[0,1]$, $T\ge2$, $u_0\in\R$ and complex $a_k$ with $|a_k|\le1$,
\begin{equation}\label{eq:mean-value-ordinates}
 \int_{u_0-1}^{u_0+1}\Bigl|\sum_{\gamma_k(\lambda)\le T}a_k\,e^{i\gamma_k(\lambda)u}\Bigr|^2du\le C\,T\log^2T .
\end{equation}
\end{lemma}

\begin{proof}
Let $\omega_0(u)=C_0\bigl(\sin(\frac{u-u_0}2)/\frac{u-u_0}2\bigr)^2$ with $C_0$ such that $\omega_0\ge1$ on $|u-u_0|\le1$; the Fourier transform of $\omega_0$ is bounded and supported in $[-1,1]$.  Expanding the square,
\begin{align}
\int \omega_0(u)\Bigl|\sum_ka_ke^{i\gamma_k(\lambda)u}\Bigr|^2du&=\sum_{k,k'}a_k\bar a_{k'}\,e^{i(\gamma_k(\lambda)-\gamma_{k'}(\lambda))u_0}\,\widehat{\omega_0}\bigl(\gamma_k(\lambda)-\gamma_{k'}(\lambda)\bigr)\nonumber\\
&\ll\#\{(k,k'):\gamma_k(\lambda)\le T,\ |\gamma_k(\lambda)-\gamma_{k'}(\lambda)|\le1\}.\label{eq:fejer-pairs}
\end{align}
Since $|\gamma_k(\lambda)-\tau_k|\le C_\eta$, every interval $[t,t+1]$ contains $O(\log(t+2))$ of the $\gamma_k(\lambda)$, uniformly in $\lambda$, because it contains that many $\tau_k$ after widening by $C_\eta$.  The number of pairs is therefore $O(T\log^2T)$.
\end{proof}

\begin{theorem}[Mean square of the cell masses along the family]\label{thm:mean-square-cells}
For every $\eps>0$ there is $C_\eps$ such that for all $X\ge2$ and all $\lambda\in[0,1]$
\begin{equation}\label{eq:mean-square-cells}
 \sum_{X<n\le2X}|\tilde c_n(\lambda)|^2\le C_\eps X^{1+\eps},\qquad
 \sum_{X<n\le2X}|\partial_\lambda\tilde c_n(\lambda)|^2\le C_\eps X^{1+\eps}.
\end{equation}
In particular the same bound already includes $\lambda=0$: $\sum_{X<n\le2X}|\tilde c_n(0)|^2\le C_\eps X^{1+\eps}$.
\end{theorem}

\begin{proof}
Fix $\eps$, $X$, $\lambda$ and put $T=X^{1+\eps}$.  For $n$ not a power of $2$, \CGeqref{eq:cell-mass-W} gives $\tilde c_n(\lambda)-\mathfrak m_n=-2\Re\sum_{j}\frac{\mu(j)}jA_{n,j}$, with
\begin{equation}\label{eq:A-nj}
 A_{n,j}=\sum_k\int_2^\infty\varphi_n(y^j)\,\frac{y^{-1/2+i\gamma_k(\lambda)}}{\log y}\,dy ,\qquad 1\le j\le J:=\Bigl\lfloor\frac{\log(2X+1)}{\log2}\Bigr\rfloor .
\end{equation}
\emph{Large ordinates.}  In $u=\log y$ the function $\varphi_n(e^{ju})e^{u/2}/u$ has support of length $O((jn)^{-1})$, size $O(n^{1/(2j)})$ and $r$-th derivative $O_r((jn)^rn^{1/(2j)})$: by Fa\`a di Bruno's formula the $r$-th derivative of $\varphi_n(e^{ju})$ is a sum of terms $\varphi_n^{(m)}(e^{ju})(je^{ju})^mj^{r-m}$, $m\le r$, and $\|\varphi_n^{(m)}\|_\infty=C_m$ does not depend on $n$; the derivatives of $e^{u/2}/u$ are $O_r(e^{u/2})$ because $u\ge\log2$.  Integrating by parts $A$ times, each term with $\gamma_k(\lambda)>T$ is $O_A\bigl((jn/\gamma_k(\lambda))^An^{1/(2j)-1}\bigr)$, and their sum over $k$ is $O(X^{-1})$ once $A>(3+\eps)/\eps$.
\emph{The cell $n=2^j$.}  Each block contains exactly one $n=2^j$; for it the term $j$ carries the cut of $J^*_\lambda$ at $y=2$ and is treated directly.  It equals $-\frac{\mu(j)}j\int_2^3J_\lambda(y)\,\frac d{dy}\varphi_n(y^j)\,dy$; the weight has $L^2$ norm $O((jn)^{1/2})$, and $\|J_\lambda\|_{L^2(2,3)}$ is bounded uniformly in $\lambda$ (Theorem~\CGref{thm:mellin-counts} with the uniform bound of Proposition~\CGref{prop:family-clock} for $J^{\rm M}_\lambda$, and Lemma~\CGref{lem:L2-partial-sums}(c) for $\sup_\lambda|C_\lambda|$).  This term is therefore $O((n/j)^{1/2})$ and contributes $O(X)$ to the sum of squares; the same holds for $\partial_\lambda\tilde c_n$, with $\partial_\lambda J_\lambda$ in place of $J_\lambda$.
\emph{Small ordinates.}  Let $\mathcal S_{n,j,\lambda}(y)=y^{-1/2}(\log y)^{-1}\sum_{\gamma_k(\lambda)\le T}y^{i\gamma_k(\lambda)}$ and let $I_{n,j}$ be the support of $y\mapsto\varphi_n(y^j)$, of length $O(X^{1/j-1})$.  By Cauchy--Schwarz in the cell,
\begin{equation}\label{eq:cell-CS}
 \Bigl|\int\varphi_n(y^j)\mathcal S_{n,j,\lambda}(y)\,dy\Bigr|^2\le\|\varphi\|_\infty^2\,|I_{n,j}|\int_{I_{n,j}}|\mathcal S_{n,j,\lambda}|^2dy .
\end{equation}
For fixed $j$ the intervals $I_{n,j}$, $X<n\le2X$, cover each point at most twice and lie in $[(X-1)^{1/j},(2X+1)^{1/j}]$, an interval of length at most $2$ in $u=\log y$.  With $dy/y=du$ and Lemma~\CGref{lem:mean-value-ordinates} ($a_k=1$),
\begin{equation}\label{eq:cell-block}
 \sum_{X<n\le2X}\Bigl|\int\varphi_n(y^j)S\,dy\Bigr|^2\ll X^{1/j-1}\Bigl(\frac{j}{\log X}\Bigr)^2\int_{u_0-1}^{u_0+1}\Bigl|\sum_{\gamma_k(\lambda)\le T}e^{i\gamma_k(\lambda)u}\Bigr|^2du\ll j^2X^{1/j+2\eps}.
\end{equation}
\emph{Assembly.}  By Minkowski's inequality over $j$, the $\ell^2$ norm over $X<n\le2X$ of $\tilde c_n(\lambda)-\mathfrak m_n$ is $\ll\sum_{j\le J}\bigl(X^{1/(2j)+\eps}+X^{-1/2}\bigr)\ll X^{1/2+\eps}$, and $\sum_{X<n\le2X}\mathfrak m_n^2\ll X$.  This is the first bound with $3\eps$ in place of $\eps$.  Away from powers of $2$, the argument for $\partial_\lambda\tilde c_n$ is the same, by \CGeqref{eq:cell-mass-W}, with coefficients $a_k=\eta_k/C_\eta$ in Lemma~\CGref{lem:mean-value-ordinates} and without the factor $(\log y)^{-1}$.
\end{proof}

\begin{corollary}[Abscissa $1$ for every member]\label{cor:abscissa-one}
For every $\lambda\in[0,1]$ and $\sigma>1$,
\begin{equation}\label{eq:abscissa-one}
 \sum_{n\ge2}|\tilde c_n(\lambda)|\,n^{-\sigma}<\infty ;
\end{equation}
hence $\tilde\zeta_{\Z,\lambda}$ and $1/\tilde\zeta_{\Z,\lambda}$ are absolutely convergent Dirichlet series in $\Re s>1$, for every $\lambda$, including the endpoint $\lambda=0$.  Moreover, uniformly in $\lambda$, $\#\{n\le X:|\tilde c_n(\lambda)|>n^{\delta}\}\ll_\eps X^{1-2\delta+\eps}$ for every $\delta>0$.
\end{corollary}

\begin{proof}
By Cauchy--Schwarz on dyadic blocks, $\sum_{X<n\le2X}|\tilde c_n(\lambda)|\le(X\cdot C_\eps X^{1+\eps})^{1/2}$, and summing over $X=2^i$ gives \CGeqref{eq:abscissa-one} for $\sigma>1+\eps/2$, with $\eps$ arbitrary.  Then $\log\tilde\zeta_{\Z,\lambda}=\sum_m\frac1m\sum_n\tilde c_n(\lambda)n^{-ms}$ converges absolutely in $\Re s>1$, and the exponentials of $\pm$ this series are absolutely convergent Dirichlet series there.  The last statement is Chebyshev's inequality on dyadic blocks.
\end{proof}

The abscissa $\frac32$ obtained below from the primitive $\pi_\lambda-\mathfrak L$ (Theorem~\CGref{thm:atomized-continuation}(ii)) is not the true one: a bound on the size of $\pi_\lambda-\mathfrak L$ does not see the cancellation of the density of $d\pi_\lambda$ in mean square, which Theorem~\CGref{thm:mean-square-cells} captures.  Absolute convergence in $\Re s>1$ is exactly what Hamburger's theorem requires, and Knopp's obstruction (Subsection~\CGref{subsec:hamburger-table}) does not arise.

\subsection{Analytic continuation of the atomized family}
\label{subsec:atomized-continuation}

Let $\mathfrak L(x)=\int_1^x(1-y^{-1})\,dy/\log y$; then $\log\frac s{s-1}=\int_1^\infty x^{-s}\,d\mathfrak L(x)$ for $\Re s>1$ (both sides have derivative $\frac1s-\frac1{s-1}$ and tend to $0$ as $s\to+\infty$), and $\mathfrak L(x)=\li(x)-\log\log x-\gamma_E$.

\begin{theorem}[Continuation of the atomized family]\label{thm:atomized-continuation}
Let $\lambda\in[0,1]$.
\begin{enumerate}[label=\textup{(\roman*)}]
\item For every $c>\frac12$, $(\pi_\lambda-\mathfrak L)\,x^{-c}\in L^2\bigl((2,\infty),dx/x\bigr)$.
\item $|\tilde c_n(\lambda)|\le C(1+I_n^{1/2})$ with $I_n$ as below, and $\sum_n|\tilde c_n(\lambda)|n^{-\sigma}<\infty$ for $\sigma>\frac32$; by Corollary~\CGref{cor:abscissa-one} this holds for $\sigma>1$, and $\tilde\zeta_{\Z,\lambda}$ is an absolutely convergent Dirichlet series in $\Re s>1$.
\item $\tilde h_\lambda$ extends holomorphically to $\Re s>\frac12$.  Hence $\tilde\zeta_{\Z,\lambda}=\zeta_\lambda e^{\tilde h_\lambda}$ is meromorphic in $\Re s>\frac12$, with a single pole at $s=1$ and no zeros.
\item For $\lambda=1$, $\tilde h_1$ extends holomorphically to $\Re s>0$, and the zeros of $\tilde\zeta_{\Z,1}$ in $\Re s>0$ are exactly $\frac12\pm i\tau_k$.
\end{enumerate}
\end{theorem}

\begin{proof}
(i) By Proposition~\CGref{prop:family-clock}, for $c>\frac12$ there is $C_c$ with $|\log\zeta_\lambda(s)-\log\frac s{s-1}|\le C_c(1+\log(|t|+2))$ on $\Re s\ge c$; the function $\log\zeta_\lambda-\log\frac s{s-1}$ is holomorphic on $\Re s>\frac12$, where $\zeta_\lambda$ has no zeros, and its pole at $1$ is removed.  It equals $s\int_1^\infty(J_\lambda^{\rm M}-\mathfrak L)x^{-s-1}dx$ for $\Re s>1$.  The argument of Theorem~\CGref{thm:mellin-counts}, applied on the lines $\Re s=c>\frac12$, gives $(J^{\rm M}_\lambda-\mathfrak L)x^{-c}\in L^2$.  The constant $C_\lambda=J^{\rm M}_\lambda-J_\lambda$, the truncation at $2$, the terms $j\ge2$ of the M\"obius inversion and $\mathfrak L(x)-\mathfrak L(x^{1/j})$ preserve the property for $c>\frac12$: with the substitution $y=x^{1/j}$, $\int_1^\infty|(J^{\rm M}_\lambda-\mathfrak L)(x^{1/j})|^2x^{-2c}\frac{dx}x=j\int_1^\infty|(J^{\rm M}_\lambda-\mathfrak L)(y)|^2y^{-2cj}\frac{dy}y<\infty$ because $cj>\frac12$, and $\mathfrak L(x^{1/j})\ll x^{1/j}/\log x$ lies in the weighted space for $2c>1$.

(ii) Put $E=\pi_\lambda-\mathfrak L$ and $I_n=\int_{n-1}^{n+1}|E|^2dx$.  Since $\int\varphi_n'=0$, $\tilde c_n(\lambda)=\int\varphi_n\,d\mathfrak L-\int E\varphi_n'\,dx$, so $|\tilde c_n(\lambda)|\le C(1+I_n^{1/2})$.  By (i), $\sum_nI_nn^{-2c-1}<\infty$, and Cauchy--Schwarz gives $\sum_nI_n^{1/2}n^{-\sigma}<\infty$ for $\sigma>c+1$.

(iii) With $\psi_n=\varphi_n(x)(n^{-s}-x^{-s})$ one has $|\psi_n'|\le C|s|n^{-\sigma-1}$ and $\psi_n$ vanishes at the ends of the cell, so
\begin{equation}\label{eq:cell-defect-bound}
 |\mathfrak h_n(\lambda;s)|\le C|s|\,n^{-\sigma-1}\bigl(1+I_n^{1/2}\bigr),\qquad \sum_nI_n^{1/2}n^{-\sigma-1}<\infty\quad(\sigma>c).
\end{equation}
Choosing $c\in(\frac12,\sigma)$, the series \CGeqref{eq:joining-factor} converges locally uniformly on $\Re s>\frac12$; the terms $m\ge2$ are evaluated at $\Re(ms)\ge2\sigma$.  Since $\mathfrak u_\lambda$ is entire (Proposition~\CGref{prop:first-cell}), $\tilde h_\lambda$ is holomorphic on $\Re s>\frac12$, and $\zeta_\lambda$ has no zeros there.

(iv) By Corollary~\CGref{cor:cn-bounded}, $\pi_{\RVM}$ oscillates by $O(1)$ on every unit cell; integrating by parts, $|\mathfrak h_n(1;s)|\le C|s|n^{-\sigma-1}$, and the series converges for $\sigma>0$.  The factor $e^{\tilde h_1}$ has no zeros, so the zeros of $\tilde\zeta_{\Z,1}$ in $\Re s>0$ are those of $\zeta_{\RVM}$.
\end{proof}

\begin{theorem}[Closing criterion; via Hamburger's theorem \cite{Hamburger1921}]\label{thm:closing}
The following are equivalent.
\begin{enumerate}[label=\textup{(\roman*)}]
\item (D), that is, $\zeta_{Z1}=\zeta$.
\item $\tilde h_0\equiv0$.
\item $J_{Z1}=0$ on $(1,2)$, and for every $n\ge2$ and every $\Re s>1$,
\begin{equation}\label{eq:closing-cells}
 \int_0^1\partial_\lambda\mathfrak h_n(\lambda;s)\,d\lambda=\mathfrak h_n(1;s).
\end{equation}
\end{enumerate}
In this case $\tilde c_n(0)=\mathbf 1_{\rm prime}(n)$, $\tilde a_N(0)=1$ for all $N$, and $\varpi_{Z1}=1$.
\end{theorem}

\begin{proof}
(ii)$\Rightarrow$(i): $\tilde h_0\equiv0$ gives $\zeta_{Z1}=\tilde\zeta_{\Z,0}=\exp\bigl(\sum_m\frac1m\sum_n\tilde c_n(0)n^{-ms}\bigr)$, an absolutely convergent Dirichlet series in $\Re s>1$ by Corollary~\CGref{cor:abscissa-one}.  (iii)$\Rightarrow$(ii): by \CGeqref{eq:cell-flow} every $\mathfrak h_n(0;\cdot)$ vanishes, and $\mathfrak u_0\equiv\log\varpi_{Z1}$ (Theorem~\CGref{thm:joining-factor}), so $\tilde h_0\equiv-\log\varpi_{Z1}$ is constant; then $\zeta_{Z1}=\varpi_{Z1}\tilde\zeta_{\Z,0}$ is an absolutely convergent Dirichlet series in $\Re s>1$ (Corollary~\CGref{cor:abscissa-one}), Theorem~\CGref{thm:Hamburger-Z1} gives $\zeta_{Z1}=\zeta$, and comparing the first coefficients, $\varpi_{Z1}=1$.  (i)$\Rightarrow$(iii): if $\zeta_{Z1}=\zeta$, Theorem~\CGref{thm:cell-form} gives $J^{\rm M}_{Z1}=J$, and $J_{Z1}=J$ by von Mangoldt's explicit formula \cite{Edwards1974}, so $J_{Z1}$ vanishes on $(1,2)$, $d\pi_{Z1}=\sum_p\delta_p$, $\tilde c_n(0)=\mathbf 1_{\rm prime}(n)$, every $\mathfrak h_n(0;\cdot)$ vanishes, and \CGeqref{eq:cell-flow} gives \CGeqref{eq:closing-cells}.
\end{proof}

\begin{remark}[The pointwise size of the cell masses]\label{rem:A6-scale}
Theorem~\CGref{thm:closing} needs no pointwise bound on $\tilde c_n(0)$: its proof uses the cell masses only through \CGeqref{eq:abscissa-one}.  The bound $\tilde c_n(0)=O(n^{\eps})$ for every $n$ is not proved directly; Corollary~\CGref{cor:abscissa-one} restricts the exceptions to a set of $n\le X$ of size $O(X^{1-2\delta+\eps})$ for the threshold $n^{\delta}$, uniformly in $\lambda$, and under the closure every cell mass is bounded, $\tilde c_n(0)=\mathbf 1_{\rm prime}(n)$.  For $n$ not a power of $2$, \CGeqref{eq:cell-mass-W} gives the trivial pointwise bound $\tilde c_n(\lambda)=O(n^{1/2+\eps})$; the mean square \CGeqref{eq:mean-square-cells} shows that the typical cell mass is $O(n^{\eps})$.
\end{remark}

\subsection{The two endpoints of the atomized family}
\label{subsec:atomized-endpoints}

\begin{proposition}[Atomization at the clock endpoint]\label{prop:atomized-clock}
At $\lambda=1$:
\begin{enumerate}[label=\textup{(\roman*)}]
\item $\tilde c_n(1)=O(1)$, and $\tilde\zeta_{\Z,1}$ is an absolutely convergent Dirichlet series in $\Re s>1$;
\item $\tilde h_1(\sigma)\to-\log\varpi_{\RVM}=\log L$ as $\sigma\to+\infty$; this limit is $\ne0$ exactly when $\varpi_{\RVM}\ne1$, in particular under Hypothesis~\CGref{hyp:varpi-less-one};
\item $\tilde h_1$ is not constant: for every $\sigma>1$, $\zeta_{\RVM}(\sigma+it)\to\varpi_{\RVM}$ as $|t|\to\infty$, whereas $\tilde\zeta_{\Z,1}(\sigma+it)$ has no limit as $|t|\to\infty$.
\end{enumerate}
\end{proposition}

\begin{proof}
(i) By Corollary~\CGref{cor:cn-bounded}, $\pi_{\RVM}$ oscillates by $O(1)$ on every unit cell, and $\tilde c_n(1)=-\int\pi_{\RVM}\varphi_n'\,dx$ with $\int\varphi_n'=0$.  (ii) is \CGeqref{eq:joining-factor} with $\sigma\to+\infty$; the limit is $\ne0$ if and only if $\varpi_{\RVM}\ne1$.  (iii) The first limit is \CGeqref{eq:vertical-limit}.  If $\tilde\zeta_{\Z,1}(\sigma+it)$ had a limit $\ell$ as $|t|\to\infty$, the mean-square argument \CGeqref{eq:mean-square}--\CGeqref{eq:mean-square-limit} applied to $\tilde\zeta_{\Z,1}-\ell$ on the line $\Re s=\sigma$ would give $\tilde a_N(1)=0$ for $N\ge2$, so $\tilde\zeta_{\Z,1}\equiv1$, which contradicts its pole at $s=1$ (Theorem~\CGref{thm:atomized-continuation}(iii)).  If $\tilde h_1$ were constant, $\zeta_{\RVM}$ would be a constant multiple of an absolutely convergent Dirichlet series, which Theorem~\CGref{thm:no-Dirichlet-RVM} excludes.
\end{proof}

The two families thus keep complementary structures and are joined by a zero-free factor:
\begin{equation}\label{eq:duality}
 \underbrace{\zeta_\lambda}_{\substack{\text{functional equation}\\ \text{zeros }\frac12\pm i\gamma_k(\lambda)}}\cdot\ e^{\tilde h_\lambda}
 \ =\ \underbrace{\tilde\zeta_{\Z,\lambda}}_{\substack{\text{Dirichlet series}\\ \text{Euler product}}},\qquad
 \tilde h_1\ \text{not constant},\qquad (D)\iff\tilde h_0\equiv0 .
\end{equation}
At $\lambda=1$ the joining factor carries the whole difference between the clock and an Euler product; Theorem~\CGref{thm:closing} states that this difference disappears at $\lambda=0$.

The reciprocal of the atomized family is again an Euler product:
\begin{equation}\label{eq:atomized-mobius}
 \frac1{\tilde\zeta_{\Z,\lambda}(s)}=\prod_{n\ge2}\bigl(1-n^{-s}\bigr)^{\tilde c_n(\lambda)}=\sum_{N\ge1}\frac{\tilde\mu_\lambda(N)}{N^s},\qquad
 \tilde\mu_\lambda(N)=\sum_{\prod_nn^{r_n}=N}\ \prod_{n:\,r_n\ge1}(-1)^{r_n}\binom{\tilde c_n(\lambda)}{r_n},
\end{equation}
absolutely convergent in $\Re s>1$ for every $\lambda$ (Corollary~\CGref{cor:abscissa-one}), because it is the exponential of the absolutely convergent series $-\sum_m\frac1m\sum_n\tilde c_n(\lambda)n^{-ms}$.  Thus $\tilde\mu_\lambda(p)=-\tilde c_p(\lambda)$, and the atomized Mertens functions
\begin{equation}\label{eq:atomized-mertens}
 \tilde M_\lambda(x):=\sum_{N\le x}\tilde\mu_\lambda(N),\qquad \frac1{\tilde\zeta_{\Z,\lambda}}=e^{-\tilde h_\lambda}F_\lambda ,
\end{equation}
exist for every $\lambda\in[0,1]$.  Under the conditions of Theorem~\CGref{thm:closing}, $\tilde c_n(0)=\mathbf 1_{\rm prime}(n)$, so $\tilde\mu_0=\mu$ and $\tilde M_0=M$.

\begin{proposition}[Mellin form of the endpoint Mertens function]\label{prop:Mellin-Mertens}
There are $c_1\ge2$ and a unique measurable $B_{Z1}$ on $(0,\infty)$, vanishing on $(0,1)$, with $x^{-c}B_{Z1}\in L^2\bigl((1,\infty),dx/x\bigr)$ and
\begin{equation}\label{eq:Mellin-Mertens}
 F_{Z1}(s)=s\int_1^\infty B_{Z1}(x)\,x^{-s-1}\,dx=\exp\Bigl(-\int_{[1,\infty)}x^{-s}\,dJ^{\rm M}_{Z1}(x)\Bigr)\qquad(\Re s>c_1),
\end{equation}
for every $c>c_1$.  Moreover $(D)$ holds if and only if $B_{Z1}=M$ almost everywhere; and if $F_{Z1}$ has an absolutely convergent Dirichlet expansion \CGeqref{eq:Z1-reciprocal-Dirichlet}, then $B_{Z1}=M_{Z1}$.
\end{proposition}

\begin{proof}
Put $A=2(\pi+2)C_{\mathfrak e}$.  By \CGeqref{eq:Z1-growth}, for $\sigma'\ge\max(4A,R_{\rm end})$, $|F_{Z1}(s)|\le C(|t|+\sigma)^{1/4}$; take $c_1=\max(2,\frac12+4A,\frac12+R_{\rm end})$.  Then $\int_{\R}|F_{Z1}(\sigma+it)|^2|\sigma+it|^{-2}dt\le C'\sigma^{-1/2}$ for $\sigma>c_1$, and the argument of Theorem~\CGref{thm:mellin-counts} gives $B_{Z1}$, its vanishing on $(0,1)$, its independence of $c$ and its uniqueness; the second equality in \CGeqref{eq:Mellin-Mertens} is Theorem~\CGref{thm:mellin-counts}.  If $\zeta_{Z1}=\zeta$, then $F_{Z1}(s)=1/\zeta(s)=s\int_1^\infty M(x)x^{-s-1}dx$ and uniqueness gives $B_{Z1}=M$.  Conversely, $B_{Z1}=M$ gives $F_{Z1}=1/\zeta$ on $\Re s>c_1$, hence $\zeta_{Z1}=\zeta$ by analytic continuation.  Under \CGeqref{eq:Z1-reciprocal-Dirichlet}, partial summation gives $F_{Z1}(s)=s\int_1^\infty M_{Z1}(x)x^{-s-1}dx$, and uniqueness gives $B_{Z1}=M_{Z1}$.
\end{proof}

\subsection{Hamburger's hypotheses at the two endpoints}
\label{subsec:hamburger-table}

Hamburger's theorem \cite{Hamburger1921}, in the form of \cite[\S2.13]{Titchmarsh1986} (see also \cite{Siegel1922,BochnerChandrasekharan1956}), has four hypotheses: (H0) $f=G/P$ with $G$ entire of finite order and $P$ a polynomial; (H1) $f(s)=\sum_na_nn^{-s}$ absolutely for $\sigma>1$, with arbitrary complex $a_n$; (H2) $f(s)\Gamma(\frac s2)\pi^{-s/2}=g(1-s)\Gamma(\frac{1-s}2)\pi^{-(1-s)/2}$; (H3) $g(1-s)=\sum_nb_nn^{s-1}$ absolutely for $\sigma<-\alpha<0$.  It contains no hypothesis on the location of zeros, and its conclusion is $f=C\zeta$.  Absolute convergence for $\sigma>1$ in (H1) cannot be weakened to some half-plane: Knopp constructed infinitely many linearly independent solutions of (H0), (H2), (H3) that are Dirichlet series absolutely convergent only in some half-plane \cite{Knopp1994,Nakamura2020}.  Within the Selberg class, the functions of degree $1$ are classified in \cite{KaczorowskiPerelli1999}.  The table records the status at the two endpoints and is updated as the programme advances.

\begin{center}
\small
\begin{tabular}{p{0.06\textwidth}p{0.40\textwidth}p{0.44\textwidth}}
\toprule
 & $\zeta_{Z1}$ (continuous limit) & $\tilde\zeta_{\Z,0}$ (atomized limit)\\
\midrule
H0 & proved: $(s-1)\zeta_{Z1}$ entire of order $1$ (Proposition~\CGref{prop:Z1-class}) & open: meromorphic in $\Re s>\frac12$ only (Theorem~\CGref{thm:atomized-continuation})\\
H1 & open: this is (D) & proved in $\Re s>1$ (Corollary~\CGref{cor:abscissa-one})\\
H2 & proved, with $g=f$ (Proposition~\CGref{prop:Z1-class}) & open\\
H3 & follows from H1 & open\\
zeros & no hypothesis & no hypothesis\\
\bottomrule
\end{tabular}
\end{center}

\noindent The two columns meet when $\tilde h_0\equiv0$: then $\zeta_{Z1}=\tilde\zeta_{\Z,0}$ has H0--H3, and Hamburger's theorem gives $\zeta_{Z1}=\zeta$ (Theorem~\CGref{thm:closing}).  It suffices that the joining factor be symmetric (Theorem~\CGref{thm:pair-criterion}), or that $\zeta_{Z1}$ have a positive representing measure (Theorem~\CGref{thm:positivity}).

The parallel between the two families, topic by topic of the paper, is the following.

\begingroup\small
\begin{longtable}{>{\raggedright\arraybackslash}p{0.20\textwidth}>{\raggedright\arraybackslash}p{0.355\textwidth}>{\raggedright\arraybackslash}p{0.355\textwidth}}
\toprule
Topic of the paper & continuous family $\zeta_\lambda$ & atomized family $\tilde\zeta_{\Z,\lambda}$\\
\midrule
\endfirsthead
\toprule
Topic of the paper & continuous family $\zeta_\lambda$ & atomized family $\tilde\zeta_{\Z,\lambda}$\\
\midrule
\endhead
\bottomrule
\endfoot
\addlinespace[2pt]
clock, $\tau_k$, bounded displacement & same & same\\
definition & Hadamard product over $\frac12\pm i\gamma_k(\lambda)$ & Euler product over $n\ge2$ with masses $\tilde c_n(\lambda)=\langle d\pi_\lambda,\varphi_n\rangle$\\
zeros & exactly $\frac12\pm i\gamma_k(\lambda)$; impossible elsewhere (Theorem~\CGref{thm:path}) & none in $\Re s>\frac12$ for every $\lambda$, since $\tilde h_\lambda$ is holomorphic there (Theorem~\CGref{thm:atomized-continuation}(iii)); for $\lambda=1$, exactly $\frac12\pm i\tau_k$ in $\Re s>0$ (Theorem~\CGref{thm:atomized-continuation}(iv))\\
functional equation & proved for every $\lambda$ & open; the factor $e^{\tilde h_\lambda}$ carries no functional equation\\
Dirichlet series & not for $\lambda=1$ in any half-plane (Theorem~\CGref{thm:no-Dirichlet-RVM}); not in $\Re s>1$ for $0<\lambda<1$ (Proposition~\CGref{prop:nonconstant}); at $\lambda=0$ this is (D) & proved in $\Re s>1$ for every $\lambda$, including $\lambda=0$ (Corollary~\CGref{cor:abscissa-one})\\
first coefficient ($\sigma\to+\infty$) & $\varpi_\lambda=\varpi_{\RVM}\prod_k\frac{1/4+\tau_k^2}{1/4+\gamma_k(\lambda)^2}$ \CGeqref{eq:varpi-lambda} & $\tilde a_1(\lambda)=1$ for every $\lambda$\\
vertical behaviour & $\lambda=1$: $\zeta_{\RVM}(\sigma+it)\to\varpi_{\RVM}$ as $|t|\to\infty$ \CGeqref{eq:vertical-limit} & almost periodic on every line of absolute convergence; $\lambda=1$: no limit as $|t|\to\infty$ (Proposition~\CGref{prop:atomized-clock})\\
half-plane bounds & $\lambda=1$: $C^{-1}\le|\zeta_{\RVM}|\le C$ on $\Re s\ge\sigma_0>1$ (Theorem~\CGref{thm:halfplane-RVM}); every $\lambda$: $|\log\zeta_\lambda-\log\frac s{s-1}|\ll_c1+\log(|t|+2)$ on $\Re s\ge c>\frac12$ & $|\tilde\zeta_{\Z,\lambda}|\le\sum_N|\tilde a_N(\lambda)|N^{-\sigma_0}$ and $|\tilde\zeta_{\Z,\lambda}^{-1}|\le\sum_N|\tilde\mu_\lambda(N)|N^{-\sigma_0}$ on $\Re s\ge\sigma_0>1$, for every $\lambda$\\
prime coordinate & counts $J_\lambda$, $\pi_\lambda$ (distributions, \CGeqref{eq:J-X}) & masses $\tilde c_n(\lambda)$ at the integers; $\tilde c_n(1)=O(1)$; mean square $X^{1+\eps}$ on $(X,2X]$ for every $\lambda$ (Theorem~\CGref{thm:mean-square-cells})\\
reciprocal & $F_\lambda=1/\zeta_\lambda$; $F_{Z1}$ is the Mellin transform of $B_{Z1}$ (Proposition~\CGref{prop:Mellin-Mertens}) & Dirichlet series with coefficients $\tilde\mu_\lambda(N)$ \CGeqref{eq:atomized-mobius}\\
Mertens & $M^{\crit}_{\RVM}=L+O_\eps(x^{1/2+\eps})$ (Theorem~\CGref{thm:unsmoothed-RVM}); $B_{Z1}=M\iff(D)$ & $\tilde M_\lambda=\sum_{N\le x}\tilde\mu_\lambda(N)$ exists for every $\lambda$ \CGeqref{eq:atomized-mertens}; the bound $x^{1/2+\eps}$ is open (it needs control of $e^{-\tilde h_\lambda}$)\\
flow in $\lambda$ & generator $\mathcal V_\lambda$ on the Hadamard product; $\Upsilon(\lambda)$ at $+\infty$ (Proposition~\CGref{prop:family-flow}) & $\log\tilde\zeta_{\Z,\lambda}=\sum_n\tilde c_n(\lambda)\bigl(-\log(1-n^{-s})\bigr)$, each mass differentiable in $\lambda$ by \CGeqref{eq:J-lambda-derivative}; joining factor \CGeqref{eq:joining-derivative}\\
limit $\lambda\to0^+$ & $\zeta_\lambda\to\zeta_{Z1}$ & $\tilde c_n(\lambda)\to\tilde c_n(0)$, $\mathfrak h_n(\lambda;\cdot)\to\mathfrak h_n(0;\cdot)$ (Proposition~\CGref{prop:cell-limit})\\
endpoint: Hamburger's hypotheses & H0 and H2 proved; H1 is (D), open & H1 proved in $\Re s>1$ (Corollary~\CGref{cor:abscissa-one}); H0 and H2 open\\
relation between the endpoints & $\tilde\zeta_{\Z,0}=\zeta_{Z1}\,e^{\tilde h_0}$ & same\\
what closes the route & $\zeta_{Z1}$ satisfies H1 & $\tilde\zeta_{\Z,0}$ satisfies H0 and H2\\
common form of the closure & $\tilde h_0\equiv0$ & $\tilde h_0\equiv0$\\
equivalent forms of the closure & $\zeta_{Z1}$ has a positive representing measure (Theorem~\CGref{thm:positivity}) & $e^{\tilde h_0}$ of finite order and symmetric under $s\mapsto1-s$ (Theorem~\CGref{thm:pair-criterion})\\
consequence of the closure & $\zeta_{Z1}=\zeta$; Step~3 gives $M(x)=O_\eps(x^{1/2+\eps})$ (Theorem~\CGref{thm:step3}) & $\tilde a_N(0)=1$, $\tilde c_n(0)=\mathbf 1_{\rm prime}(n)$, $\tilde\mu_0=\mu$, $\tilde M_0=M$; $\tilde\zeta_{\Z,0}=\zeta=\zeta_{Z1}$\\
\end{longtable}
\endgroup

\noindent\emph{How the two routes complement each other.}  Along $\lambda$ each family keeps exactly the half of Hamburger's hypotheses that the other one lacks:
\begin{equation}\label{eq:complement}
 \underbrace{\zeta_\lambda}_{\substack{\text{keeps the functional equation}\\ \text{and the zeros on }\Re s=\frac12}}\qquad\longleftrightarrow\qquad\underbrace{\tilde\zeta_{\Z,\lambda}}_{\substack{\text{keeps the Dirichlet series}\\ \text{and the Euler product}}},\qquad \tilde\zeta_{\Z,\lambda}=\zeta_\lambda\,e^{\tilde h_\lambda}.
\end{equation}
The closure is the same in both columns, the vanishing of the joining factor at the endpoint: by \CGeqref{eq:joining-factor} and Theorem~\CGref{thm:closing},
\begin{equation}\label{eq:closure-common}
 \tilde h_0\equiv0\iff J_{Z1}=0\text{ on }(1,2),\ \ \mathfrak h_n(0;s)=0\ \ \text{for all }n\ge2,\ \Re s>1,\ \text{and}\ \varpi_{Z1}=1 .
\end{equation}
No further condition is needed in the atomized column: the cell masses have mean square $X^{1+\eps}$ uniformly in $\lambda$ (Theorem~\CGref{thm:mean-square-cells}), so $\tilde\zeta_{\Z,0}$ satisfies H1 (Corollary~\CGref{cor:abscissa-one}).

\subsection{Programme}
\label{subsec:programme}

The steps are numbered in the order in which they are followed.  The status column is updated with each step.

\begingroup\small
\begin{longtable}{p{0.05\textwidth}p{0.55\textwidth}p{0.30\textwidth}}
\toprule
Step & Statement & Status\\
\midrule
\endfirsthead
\toprule
Step & Statement & Status\\
\midrule
\endhead
\bottomrule
\endfoot
A1 & Atomization by smooth cells, coefficients $\tilde a_N(\lambda)$ & proved (Definition~\CGref{def:atomized})\\
A2 & Joining factor and explicit cell defects, \CGeqref{eq:joining-factor}--\CGeqref{eq:cell-explicit-derivative} & proved (Thm.~\CGref{thm:joining-factor}, Prop.~\CGref{prop:cell-explicit})\\
A3 & $L^2$ square-root estimate for every $\lambda$; Dirichlet series in $\Re s>\frac32$ & proved (Theorem~\CGref{thm:atomized-continuation}); superseded by A6\\
A4 & $\tilde h_\lambda$ holomorphic in $\Re s>\frac12$; in $\Re s>0$ for $\lambda=1$ & proved (Theorem~\CGref{thm:atomized-continuation})\\
A5 & Limits $\lambda\to0^+$ of $\tilde c_n$, $\mathfrak h_n$, $\varpi_\lambda$; flow at infinity $\Upsilon$ & proved (Propositions~\CGref{prop:family-flow}, \CGref{prop:cell-limit})\\
A5$'$ & Endpoints: $\tilde h_1$ not constant; atomized M\"obius coefficients $\tilde\mu_\lambda$; Mellin form $B_{Z1}$, with $B_{Z1}=M_{Z1}$ only under \CGeqref{eq:Z1-reciprocal-Dirichlet} & proved (Propositions~\CGref{prop:atomized-clock}, \CGref{prop:Mellin-Mertens})\\
A6 & Mean square of the cell masses uniformly in $\lambda$; Dirichlet series in $\Re s>1$ for every $\lambda$, including $\lambda=0$ & proved (Theorem~\CGref{thm:mean-square-cells}, Corollary~\CGref{cor:abscissa-one})\\
A7 & Cell identity \CGeqref{eq:closing-cells} for every $n\ge2$, and $J_{Z1}=0$ on the first cell $(1,2)$ & open; reduced to one number for each cell in A20\\
A8 & Equivalent scalar form: $\int_0^1\Upsilon\,d\lambda=\log L$, i.e.\ \CGeqref{eq:product-identity} & open; follows from A7\\
A9 & $\zeta_{Z1}=\zeta$ by Hamburger; $M(x)=O_\eps(x^{1/2+\eps})$ by Step~3 & proved as implication from A7 (Thms.~\CGref{thm:closing}, \CGref{thm:step3})\\
A10 & Continuation of $\tilde h_\lambda$ to $0<\Re s\le\frac12$ for $\lambda<1$ & open; not needed for A9; at $\lambda=0$ it follows from \textup{(B1)}, \textup{(B2)} (A23)\\
A11 & $\zeta_\lambda$ is not a Dirichlet series and $\tilde h_\lambda$ is not constant for $0<\lambda\le1$; convexity of $\log\zeta_\lambda(\sigma)$ in $\lambda$ & proved (Propositions~\CGref{prop:nonconstant}, \CGref{prop:convexity})\\
A12 & Hamburger's theorem for the pair: A9 follows if $e^{\tilde h_0}$ is of finite order and symmetric under $s\mapsto1-s$ & proved as implication (Theorem~\CGref{thm:pair-criterion}); at $\lambda=0$ the hypothesis is equivalent to (D) (Remark~\CGref{rem:pair-no-slack})\\
A13 & Positivity: A9 follows if $\zeta_{Z1}$ has a positive representing measure; then $\varpi_{Z1}\le1$ & proved as implication (Theorem~\CGref{thm:positivity}, Corollary~\CGref{cor:positivity})\\
A14 & Clock end: density plus explicit atoms of $d\pi_{\RVM}$, cell masses $\frac1{\log n}+\mathbf1_{2^j}(n)\frac{\mu(j)}jJ_{\RVM}(2)+O(n^{-1/4})$, defects $O(|s|n^{-\sigma-4/3})$ & proved (Theorem~\CGref{thm:clock-density}, Corollary~\CGref{cor:clock-cells})\\
A15 & Height kernel of the cell masses; flow as count difference; Gaussian defect and the density condition $\tilde h_0(1)=0$ & proved (Propositions~\CGref{prop:height-kernels}, \CGref{prop:flow-count}, \CGref{prop:gaussian-defect})\\
A16 & Hamburger's theorem for the pair without continuation: A9 follows if $T_0\equiv0$, or if $\mathsf V_\sigma(0)=0$; symmetry fails at $\lambda=1$; exceptional set finite & proved (Theorem~\CGref{thm:double-hamburger}, Propositions~\CGref{prop:clock-asymmetry}, \CGref{prop:exceptional-set})\\
A17 & $\zeta_{Z1}(\sigma)\ge\zeta(\sigma)$ for $\sigma>1$; A9 follows from the reverse inequality at one $\sigma$, or from $\varpi_\lambda\le1$ along $\lambda\to0^+$, or from $\mathsf W(0)\ge0$ & proved (Theorem~\CGref{thm:real-axis}, Propositions~\CGref{prop:one-number}, \CGref{prop:open-path-varpi})\\
A18 & Space $\mathscr O$ of ordinate lists: at most one coincidence list; A9 is equivalent to its existence & proved (Theorem~\CGref{thm:rigidity}, Corollary~\CGref{cor:existence-form})\\
A19 & Existence of a list in $\mathscr O$ with $\tilde h_{\boldsymbol{\kappa}}$ constant, equivalently $\min_{\mathscr O}\mathsf V_\sigma=0$ & open; equivalent to A7\\
A20 & A9 holds if and only if $\tilde c_n(0)=\mathbf 1_{\rm prime}(n)$ for every $n$; every $\lambda\in\mathsf E$ is a witness; the identity \CGeqref{eq:per-cell} cell by cell & proved (Corollary~\CGref{cor:cell-mass-form}, Propositions~\CGref{prop:witness}, \CGref{prop:per-cell}); the identity \CGeqref{eq:per-cell} itself is open\\
A21 & Symmetric joining factors: order one is needed, and A12 has no slack at $\lambda=0$ & proved (Proposition~\CGref{prop:dichotomy}, Remarks~\CGref{rem:pair-no-slack}, \CGref{rem:dichotomy-counterexample})\\
A22 & Half-line form: A9 follows if $\tilde h_0(\sigma)=O_N(\sigma N^{-\sigma})$ for every $N$; the $N$th rung is the constancy of $J^{\rm M}_{Z1}$ on $(N-1,N)$ & proved (Theorems~\CGref{thm:ladder}, \CGref{thm:halfline-hamburger}, Corollary~\CGref{cor:ladder-pair})\\
A23 & Bridge theorem: A9 follows if $\Im\tilde h_0$ vanishes on $\Re s=\frac12$ and $\tilde h_0$ is temperate; the two conditions hold for no $\lambda>0$ & proved (Theorem~\CGref{thm:bridge}, Corollary~\CGref{cor:bridge-family}); \textup{(B1)}, \textup{(B2)} at $\lambda=0$ open\\
A24 & Transport from the clock: $\mathsf S_{\mathsf f}$ entire of type $C_\eta\log N$, moment series on the clock, defect of the clock by Poisson summation & proved (Proposition~\CGref{prop:rung-functional}, Theorem~\CGref{thm:transport-band}, Proposition~\CGref{prop:clock-defect-poisson})\\
A24b & Multiscale Jacobian at the clock: stable finite ordinate block on windows $U\asymp\log T$, fixed-window obstruction, dense range, and rapidly summable smooth tails & proved (Theorem~\CGref{thm:clock-jacobian}); quantitative multiscale continuation to a solution open\\
A25 & Rung system with free masses: A9 holds if and only if it has a solution in $\mathscr O$; the solution is unique and nondegenerate & proved (Theorems~\CGref{thm:free-masses}, \CGref{thm:nondegenerate}); existence open\\
A26 & Coordinate gradient of $\mathsf V_\sigma$ on $\mathscr O$; A9 would follow from a boundary-aware coordinate/KKT criterion forcing every minimizer to have zero defect & proved (Proposition~\CGref{prop:critical-points}); the forcing statement is open\\
A27 & Acceptance criterion and models: the final step is a theorem for the finite Euler product over the primes $2,3$ (not for more primes, see after Theorem~\CGref{thm:finite-euler}); no object of finite rank in the window of $\zeta$ & proved (Theorems~\CGref{thm:acceptance}, \CGref{thm:finite-euler}, Propositions~\CGref{prop:graph-model}, \CGref{prop:no-finite-rank})\\
\end{longtable}
\endgroup

\begin{remark}[Input for A7]\label{rem:input-A7}
Steps A1--A6 use only two properties of an ordinate list: bounded displacement from the clock and a primitive of the count difference of size $O(\log\tau)$.  The list $(\tau_k)$ has both properties, and for it the corresponding function $\zeta_{\RVM}$ is not a Dirichlet series (Theorem~\CGref{thm:no-Dirichlet-RVM}).  Step A7 must therefore use a property of the ordinates $\gamma_k$ that the clock ordinates do not have; the natural candidate is the arithmetic of $\zeta$ carried by the $\gamma_k$, for instance Landau's formula $\sum_{0<\gamma\le T}x^\rho=-\frac T{2\pi}\Lambda(x)+O(\log T)$, in which the prime powers appear as resonances of the ordinates.  The multiscale analysis below separates a different issue: the linearized clock has enough directions, but only on windows whose length grows like $\log T$; thus the obstruction is not fixed-height injectivity but quantitative nonlinear continuation across scales.
\end{remark}

\section{Criteria for the final step and the clock end of the family}
\label{sec:criteria}

This section collects the consequences of Corollary~\CGref{cor:abscissa-one} for the final step.  Along the open path $0<\lambda\le1$ no member $\zeta_\lambda$ is a Dirichlet series (Proposition~\CGref{prop:nonconstant}); the family is convex in $\lambda$ on the real axis (Proposition~\CGref{prop:convexity}); Hamburger's theorem applies to the pair $(\zeta_\lambda,\tilde\zeta_{\Z,\lambda})$ as soon as the joining factor is symmetric (Theorem~\CGref{thm:pair-criterion}); positivity of the measure representing $\zeta_{Z1}$ gives $\varpi_{Z1}\le1$ and closes the route (Theorem~\CGref{thm:positivity}).  At the clock end the density of $d\pi_{\RVM}$ is computed explicitly (Theorem~\CGref{thm:clock-density}), and the flow of the cell masses is written in the heights of the ordinates (Propositions~\CGref{prop:height-kernels} and~\CGref{prop:flow-count}).  Subsections~\CGref{subsec:cell-mass-form}--\CGref{subsec:transport} reduce the final step to the cell masses (Corollary~\CGref{cor:cell-mass-form}), state two criteria whose hypotheses are written where the joining factor is defined, the half-line form (Theorem~\CGref{thm:halfline-hamburger}) and the bridge theorem (Theorem~\CGref{thm:bridge}), and compute the transport from the clock on each cell (Theorem~\CGref{thm:transport-band}).

Throughout this section the function $\chi$ of Subsection~\CGref{subsec:mean-square} is even, so that $\varphi$ is even, and
\begin{equation}\label{eq:cell-moments}
 \hat{\varphi}(\xi):=\int_{\R}\varphi(t)e^{-2\pi i\xi t}\,dt,\qquad \mathsf M_r:=\int_{\R}t^r\varphi(t)\,dt,\qquad \mathsf C_r:=\int_{\R}t^r\varphi(t)\cos(2\pi t)\,dt .
\end{equation}
$\hat{\varphi}$ is real and even, and $\mathsf M_1=\mathsf C_1=0$.  Since $\sum_{m\in\Z}\varphi(t-m)=1$,
\begin{equation}\label{eq:partition-fourier}
 \hat{\varphi}(q)=\int_0^1\sum_{m\in\Z}\varphi(t-m)\,e^{-2\pi iqt}\,dt=\delta_{q0}\qquad(q\in\Z),
\end{equation}
so $\mathsf M_0=1$ and $\mathsf C_0=\hat{\varphi}(1)=0$.

\subsection{The open path}
\label{subsec:open-path}

\begin{proposition}[No member of the open path is a Dirichlet series]\label{prop:nonconstant}
Let $0<\lambda\le1$.  Then $\zeta_\lambda$ is not a constant multiple of a Dirichlet series $\sum_na_nn^{-s}$ converging absolutely in $\Re s>1$.  In particular the joining factor $\tilde h_\lambda$ is not constant.
\end{proposition}

\begin{proof}
Suppose $\zeta_\lambda(s)=\sum_na_nn^{-s}$ absolutely for $\Re s>1$.  By Theorem~\CGref{thm:path}, exactly as in Proposition~\CGref{prop:Z1-class}, $(s-1)\zeta_\lambda(s)$ is entire of order one and $\zeta_\lambda$ satisfies \CGeqref{eq:Z1-FE}; these are the hypotheses of Hamburger's theorem, so $\zeta_\lambda=C\zeta$ with $C\ne0$.  The nontrivial zeros of $\zeta$ with positive ordinate are then $\frac12+i\gamma_k(\lambda)$, $k\ge1$, with multiplicity.  The ordinates of $\zeta$, listed in non-decreasing order, are the $\gamma_k$, and the $\gamma_k(\lambda)$ are increasing; hence $\gamma_k(\lambda)=\gamma_k$, that is $\lambda\eta_k=0$, for every $k$.  Since $\lambda>0$, every $\eta_k$ vanishes, all members of the family coincide (Remark~\CGref{rem:only-gamma-tau}), and $\zeta_{\RVM}=\zeta_\lambda$ is an absolutely convergent Dirichlet series, against Theorem~\CGref{thm:no-Dirichlet-RVM}.  If $\tilde h_\lambda$ were a constant $c$, then $\zeta_\lambda=e^{-c}\tilde\zeta_{\Z,\lambda}$ would be such a series by Corollary~\CGref{cor:abscissa-one}.
\end{proof}

The two columns of \CGeqref{eq:complement} exist for every $\lambda$, and for every $\lambda\in(0,1]$ they are different functions.  The joining factor converges, $\tilde h_\lambda\to\tilde h_0$ as $\lambda\to0^+$, through non-constant functions; (D) is the statement that the limit is constant.  Any property of the displacement sequence that would imply (D) must distinguish the full sequence $(\eta_k)$ from every scaled sequence $(c\,\eta_k)$ with $0\le c<1$: the latter correspond to members $\lambda>0$ of the family, for which Proposition~\CGref{prop:nonconstant} excludes closure.  Thus this observation is a necessary discrimination test for a proposed criterion, not an assertion that the actual displacement sequence already satisfies such a criterion.

\subsection{Convexity along the family}
\label{subsec:convexity}

\begin{proposition}[Convexity on the real axis]\label{prop:convexity}
For every real $\sigma>1$ the function $\lambda\mapsto\log\zeta_\lambda(\sigma)$ is strictly convex on $[0,1]$:
\begin{equation}\label{eq:convexity-second}
 \partial_\lambda^2\log\zeta_\lambda(\sigma)=\sum_{k\ge1}\eta_k^2\Bigl(\mathcal C\bigl(\sigma-\tfrac12,\gamma_k(\lambda)\bigr)-\mathcal C\bigl(\tfrac12,\gamma_k(\lambda)\bigr)\Bigr)>0,\qquad \mathcal C(a,\gamma):=\frac{2(a^2-\gamma^2)}{(a^2+\gamma^2)^2}.
\end{equation}
Consequently, with the generator $\mathcal V_\lambda$ of Theorem~\CGref{thm:generator},
\begin{equation}\label{eq:convexity-sandwich}
 \log\zeta_{\RVM}(\sigma)-\mathcal V_1(\sigma)<\log\zeta_{Z1}(\sigma)<\log\zeta_{\RVM}(\sigma)-\mathcal V_0(\sigma)\qquad(\sigma>1).
\end{equation}
\end{proposition}

\begin{proof}
With $a=\sigma-\frac12>\frac12$ one has $q(\sigma)=a^2-\frac14$, so by \CGeqref{eq:P-lambda} $\log\zeta_\lambda(\sigma)=-\log G(\sigma)+\sum_k\bigl(\log(a^2+\gamma_k(\lambda)^2)-\log(\frac14+\gamma_k(\lambda)^2)\bigr)$.  Since $\partial_\gamma^2\log(a^2+\gamma^2)=\mathcal C(a,\gamma)$ and $\partial_\lambda\gamma_k(\lambda)=\eta_k$, termwise differentiation gives the series in \CGeqref{eq:convexity-second}; its terms are $O(\eta_k^2\gamma_k^{-2})$, so it converges uniformly on $[0,1]$.  Fix $\gamma>\frac12$.  Then $\mathcal C(\frac12,\gamma)<0$; $\partial_a\mathcal C(a,\gamma)=4a(3\gamma^2-a^2)(a^2+\gamma^2)^{-3}>0$ for $0<a<\sqrt3\,\gamma$; and $\mathcal C(a,\gamma)>0$ for $a>\gamma$.  Hence $\mathcal C(a,\gamma)>\mathcal C(\frac12,\gamma)$ for every $a>\frac12$.  All ordinates satisfy $\gamma_k(\lambda)>\tau_0>\frac12$, so every term of the series is $\ge0$; the sum is $>0$ exactly when some $\eta_k\ne0$, that is, when $\boldsymbol\gamma\ne\boldsymbol\tau$, and this holds with $k=1$: $\eta_1=\tau_1-\gamma_1>0$ by the comparison $\gamma_1<14.14<\tau_1$ made in the proof of Proposition~\CGref{prop:family-flow}. So the sum is positive.  A strictly convex function $F$ on $[0,1]$ satisfies $F(1)-F'(1)<F(0)<F(1)-F'(0)$; with $F(\lambda)=\log\zeta_\lambda(\sigma)$ and $F'=\mathcal V_\lambda(\sigma)$ \CGeqref{eq:generator-zeta} this is \CGeqref{eq:convexity-sandwich}.
\end{proof}

As $\sigma\to+\infty$, \CGeqref{eq:convexity-second} becomes $-\Upsilon'(\lambda)>0$ and \CGeqref{eq:convexity-sandwich} becomes the last statement of Proposition~\CGref{prop:family-flow}.  The sign in \CGeqref{eq:convexity-second} comes from the test function $x^{-\sigma}$, whose transform is not oscillatory in the ordinate; for the cell functions $\varphi_n$ the corresponding second derivative $\partial_\lambda^2\tilde c_n(\lambda)=-2\Re\sum_k\eta_k^2W_n''(\gamma_k(\lambda))$ has terms of both signs.

\subsection{Hamburger's theorem for the pair}
\label{subsec:pair}

\begin{theorem}[Hamburger's theorem for the pair $(\zeta_\lambda,\tilde\zeta_{\Z,\lambda})$]\label{thm:pair-criterion}
Let $\lambda\in[0,1]$.  Suppose that the joining factor $e^{\tilde h_\lambda}$ extends to a meromorphic function on $\C$ of finite order with finitely many poles, and that
\begin{equation}\label{eq:pair-symmetry}
 e^{\tilde h_\lambda(s)}=e^{\tilde h_\lambda(1-s)} .
\end{equation}
Then $\tilde\zeta_{\Z,\lambda}=\zeta$.  For $\lambda=0$, moreover, $\zeta_{Z1}=\zeta$ and $\tilde h_0\equiv0$.
\end{theorem}

\begin{proof}
Put $f=\tilde\zeta_{\Z,\lambda}=\zeta_\lambda e^{\tilde h_\lambda}$.  By Theorem~\CGref{thm:path}, $(s-1)\zeta_\lambda$ is entire of order one, so $f=\mathsf G/\mathsf P$ with $\mathsf G$ entire of finite order and $\mathsf P$ a polynomial: this is (H0).  By \CGeqref{eq:Z1-FE} for $\zeta_\lambda$ and \CGeqref{eq:pair-symmetry}, $f(s)\Gamma(\frac s2)\pi^{-s/2}=f(1-s)\Gamma(\frac{1-s}2)\pi^{-(1-s)/2}$: this is (H2) with $g=f$.  (H1) is Corollary~\CGref{cor:abscissa-one}, and it gives (H3) for $g=f$.  Hamburger's theorem gives $f=C\zeta$, and $C=1$ because $\tilde a_1(\lambda)=1$.  By Theorem~\CGref{thm:atomized-continuation}(iii), $f=\zeta$ has no zeros in $\Re s>\frac12$, and by the functional equation none in $\Re s<\frac12$ apart from the trivial ones: the Riemann hypothesis holds.  For $\lambda=0$, Theorem~\CGref{thm:Hamburger-Z1} (RH $\Leftrightarrow\zeta_{Z1}=\zeta$) then gives $\zeta_{Z1}=\zeta=\tilde\zeta_{\Z,0}$, and $e^{\tilde h_0}=1$.
\end{proof}

Theorem~\CGref{thm:pair-criterion} uses the functional equation of $\zeta_\lambda$ and the Dirichlet series of $\tilde\zeta_{\Z,\lambda}$ as they stand; the only hypothesis is on the factor that joins them, and it asks for symmetry instead of constancy.  At $\lambda=1$ the two hypotheses do not hold together: the conclusion $\tilde\zeta_{\Z,1}=\zeta$ would place the zeros of $\zeta$ at $\frac12\pm i\tau_k$ (Theorem~\CGref{thm:atomized-continuation}(iv)), so $\zeta=\zeta_{\RVM}$ by Theorem~\CGref{thm:Hadamard}, against Theorem~\CGref{thm:no-Dirichlet-RVM}.

\begin{theorem}[Real-variable and vertical forms of Hamburger's theorem for the pair]\label{thm:double-hamburger}
For $\lambda\in[0,1]$, $y>0$ and $\sigma>1$ put
\begin{equation}\label{eq:theta-defect}
 T_\lambda(y):=\sum_{N\ge1}\tilde a_N(\lambda)\,e^{-\pi N^2y}-y^{-1/2}\sum_{N\ge1}\tilde a_N(\lambda)\,e^{-\pi N^2/y}-\tfrac12\bigl(y^{-1/2}-1\bigr),
\end{equation}
\begin{equation}\label{eq:vertical-defect}
 k_n(\lambda):=\sum_{d\mid n}\tilde a_d(\lambda)\,\mu(n/d),\qquad \mathsf V_\sigma(\lambda):=\sum_{n\ge2}|k_n(\lambda)|^2\,n^{-2\sigma}.
\end{equation}
\begin{enumerate}[label=\textup{(\roman*)}]
\item $\tilde\zeta_{\Z,\lambda}/\zeta=\sum_{n\ge1}k_n(\lambda)n^{-s}$ converges absolutely in $\Re s>1$, $k_1(\lambda)=1$, and
\begin{equation}\label{eq:vertical-mean}
 \lim_{T\to\infty}\frac1{2T}\int_{-T}^{T}\left|\frac{\tilde\zeta_{\Z,\lambda}(\sigma+it)}{\zeta(\sigma+it)}\right|^2dt=1+\mathsf V_\sigma(\lambda)\ \ge1 .
\end{equation}
\item For each $\lambda$ the following are equivalent: \textup{(a)} $\tilde\zeta_{\Z,\lambda}=\zeta$; \textup{(b)} $T_\lambda(y)=0$ for every $y>0$; \textup{(c)} $\mathsf V_\sigma(\lambda)=0$ for some, equivalently every, $\sigma>1$.  Each of them implies (D).
\item $\lambda\mapsto T_\lambda(y)$ and $\lambda\mapsto\mathsf V_\sigma(\lambda)$ are continuous on $[0,1]$, and
\begin{equation}\label{eq:double-hamburger-limit}
 (\mathrm D)\iff T_0\equiv0\iff\mathsf V_\sigma(0)=0\iff\lim_{\lambda\to0^+}T_\lambda(y)=0\ \ (y>0)\iff\lim_{\lambda\to0^+}\mathsf V_\sigma(\lambda)=0 .
\end{equation}
At the clock end $T_1\not\equiv0$ and $\mathsf V_\sigma(1)>0$.
\end{enumerate}
\end{theorem}

\begin{proof}
(i) $1/\zeta=\sum_n\mu(n)n^{-s}$ and $\tilde\zeta_{\Z,\lambda}$ (Corollary~\CGref{cor:abscissa-one}) converge absolutely in $\Re s>1$, and so does their product, with $k_1=\tilde a_1=1$.  The mean square of an absolutely convergent Dirichlet series on a vertical line is $\sum_n|k_n|^2n^{-2\sigma}$, as in \CGeqref{eq:mean-square}--\CGeqref{eq:mean-square-limit}.
(ii) (a)$\Leftrightarrow$(c) because $\mathsf V_\sigma(\lambda)=0$ means $k_n(\lambda)=0$ for $n\ge2$.  (a)$\Rightarrow$(b) is the theta relation $\sum_{n\ge1}e^{-\pi n^2y}=y^{-1/2}\sum_{n\ge1}e^{-\pi n^2/y}+\frac12(y^{-1/2}-1)$.  (b)$\Rightarrow$(a): the relation $T_\lambda\equiv0$ is the identity reached in the first half of the proof of Hamburger's theorem in \cite[\S2.13]{Titchmarsh1986}, with $b_n=a_n=\tilde a_n(\lambda)$; the second half of that proof uses only this identity and the absolute convergence of $\sum_na_nn^{-\sigma'}$ for $\sigma'>1$, and gives $a_k=a_{k+1}$ for every $k$, hence $\tilde a_N(\lambda)=\tilde a_1(\lambda)=1$.  If (a) holds, then by Theorem~\CGref{thm:atomized-continuation}(iii) and the functional equation the nontrivial zeros of $\zeta$ lie on $\Re s=\frac12$; they are the points $\frac12\pm i\gamma_k$, and Theorem~\CGref{thm:Hamburger-Z1} gives (D).
(iii) By Theorem~\CGref{thm:mean-square-cells}, for every $\sigma'>1$ the sums $\sum_n|\tilde c_n(\lambda)|n^{-\sigma'}$ are bounded uniformly in $\lambda$; hence so are $\sum_N|\tilde a_N(\lambda)|N^{-\sigma'}\le\exp\bigl(\sum_m\frac1m\sum_n|\tilde c_n(\lambda)|n^{-m\sigma'}\bigr)$ and $\sum_n|k_n(\lambda)|n^{-\sigma'}$.  Each $\tilde a_N(\lambda)$ and each $k_n(\lambda)$ is a polynomial in finitely many cell masses, which are continuous in $\lambda$.  The tails $\sum_{N>N_0}$ of \CGeqref{eq:theta-defect} and of \CGeqref{eq:vertical-defect} (with $1<\sigma'<\sigma$) are small uniformly in $\lambda$, which gives the continuity.  Under (D), $\tilde a_N(0)=1$ (Theorem~\CGref{thm:closing}), so $T_0\equiv0$ and $\mathsf V_\sigma(0)=0$; the converse is (ii) at $\lambda=0$, and the two limits are the values at $\lambda=0$ by continuity.  Finally $\tilde\zeta_{\Z,1}\ne\zeta$ by the remark after Theorem~\CGref{thm:pair-criterion}.
\end{proof}

Theorem~\CGref{thm:double-hamburger} needs no continuation of the joining factor: \CGeqref{eq:theta-defect} is a real function of $y$ built from the atomized coefficients alone, and \CGeqref{eq:vertical-mean} lives in $\Re s>1$, where
\begin{equation}\label{eq:quotient-three}
 \frac{\tilde\zeta_{\Z,\lambda}}{\zeta}=e^{\tilde h_\lambda}\,\frac{\zeta_\lambda}{\zeta}
\end{equation}
is a quotient of three explicit functions.  Its vertical mean is $1$ and its vertical mean square is at least $1$ for every $\lambda$; (D) is the statement that at $\lambda=0$ the mean square does not exceed $1$.

\begin{proposition}[At the clock end the symmetry fails]\label{prop:clock-asymmetry}
The joining factor $e^{\tilde h_1}$ has no meromorphic continuation to $\C$ that satisfies \CGeqref{eq:pair-symmetry}.
\end{proposition}

\begin{proof}
$e^{\tilde h_1}$ is holomorphic and zero-free in $\Re s>0$ (Theorem~\CGref{thm:atomized-continuation}(iv)).  A meromorphic continuation with \CGeqref{eq:pair-symmetry} is holomorphic and zero-free also in $\Re s<1$, hence entire and zero-free, so it equals $e^{g}$ with $g$ entire, $g=\tilde h_1$ on $\Re s>0$, and $g(s)-g(1-s)$ constant.  For $\Re s\ge\frac14$, the bound $|\mathfrak h_n(1;ms)|\le C\,m|s|\,n^{-m\sigma-1}$ of Theorem~\CGref{thm:atomized-continuation}(iv) gives $\sum_m\frac1m\sum_n|\mathfrak h_n(1;ms)|\le C|s|\sum_{n\ge2}n^{-1}(n^{1/4}-1)^{-1}$, and $|\mathfrak u_1(s)|\le C(1+|s|)$ by \CGeqref{eq:first-cell}; so $|g(s)|\le C(1+|s|)$ on $\Re s\ge\frac14$, and by the symmetry also on $\Re s\le\frac34$.  An entire function of linear growth is a polynomial of degree at most $1$, and $g(s)-g(1-s)$ constant then forces $g$ to be constant.  This contradicts Proposition~\CGref{prop:atomized-clock}(iii).
\end{proof}

At $\lambda=1$ it is therefore the symmetry \CGeqref{eq:pair-symmetry} itself that fails in Theorem~\CGref{thm:pair-criterion}, whatever the order of the continuation.

\subsection{Positivity and the value at infinity}
\label{subsec:positivity}

For an even Schwartz function $\mathsf g$ on $\R$ put $\hat{\mathsf g}(\xi)=\int_{\R}\mathsf g(u)e^{-2\pi i\xi u}du$ and $\mathsf g^{\sharp}(s)=\int_0^\infty\mathsf g(u)u^{s-1}du$.

\begin{lemma}[Summation formula for a representing measure]\label{lem:summation-positive}
Let $\lambda\in[0,1]$ and let $\mathsf A$ be a positive measure on $[1,\infty)$ with
\begin{equation}\label{eq:representing-measure}
 \zeta_\lambda(s)=\int_{[1,\infty)}u^{-s}\,d\mathsf A(u)\qquad(\Re s>1).
\end{equation}
Then for every even Schwartz function $\mathsf g$
\begin{equation}\label{eq:summation-positive}
 \mathsf g(0)+2\int_{[1,\infty)}\mathsf g\,d\mathsf A=\hat{\mathsf g}(0)+2\int_{[1,\infty)}\hat{\mathsf g}\,d\mathsf A .
\end{equation}
\end{lemma}

\begin{proof}
$\mathsf g^{\sharp}$ is meromorphic with simple poles at $s=0,-2,-4,\dots$, residue $\mathsf g(0)$ at $s=0$, and decays faster than any power of $|t|$ on vertical strips; $\mathsf g^{\sharp}(1)=\frac12\hat{\mathsf g}(0)$; and
\begin{equation}\label{eq:tate-local}
 \pi^{-s/2}\Gamma(\tfrac s2)\,\hat{\mathsf g}^{\sharp}(1-s)=\pi^{-(1-s)/2}\Gamma(\tfrac{1-s}2)\,\mathsf g^{\sharp}(s)
\end{equation}
for $0<\Re s<1$, and then by meromorphic continuation: the Fourier transform of the even homogeneous distribution $|\xi|^{-s}$ is a multiple of $|u|^{s-1}$, and the multiple is fixed by $\mathsf g(u)=e^{-\pi u^2}$, for which $\hat{\mathsf g}=\mathsf g$ and $\mathsf g^{\sharp}(s)=\frac12\pi^{-s/2}\Gamma(\frac s2)$.  Fix $1<c<2$.  By Mellin inversion and Fubini, $\int\mathsf g\,d\mathsf A=\frac1{2\pi i}\int_{(c)}\zeta_\lambda(s)\mathsf g^{\sharp}(s)\,ds$, since $\int u^{-c}d\mathsf A<\infty$.  On $\Re s\ge c$, $|\zeta_\lambda(s)|\le\zeta_\lambda(c)$; with the functional equation \CGeqref{eq:Z1-FE}, valid for every $\zeta_\lambda$ by Theorem~\CGref{thm:path}, and the finite order of $(s-1)\zeta_\lambda$, the Phragm\'en--Lindel\"of principle bounds $\zeta_\lambda$ by a power of $|t|$ on $1-c\le\Re s\le c$, $|t|\ge1$.  Moving the line to $\Re s=1-c$ crosses $s=1$, where $\zeta_\lambda$ has residue $1$, and $s=0$, where $\zeta_\lambda(0)=P_\lambda(0)/G(0)=-\frac12$:
\begin{equation}\label{eq:summation-shift}
 \int\mathsf g\,d\mathsf A=\tfrac12\hat{\mathsf g}(0)-\tfrac12\mathsf g(0)+\frac1{2\pi i}\int_{(1-c)}\zeta_\lambda(s)\mathsf g^{\sharp}(s)\,ds .
\end{equation}
By \CGeqref{eq:Z1-FE} and \CGeqref{eq:tate-local}, $\zeta_\lambda(s)\mathsf g^{\sharp}(s)=\zeta_\lambda(1-s)\hat{\mathsf g}^{\sharp}(1-s)$, and the last integral is $\frac1{2\pi i}\int_{(c)}\zeta_\lambda(s)\hat{\mathsf g}^{\sharp}(s)\,ds=\int\hat{\mathsf g}\,d\mathsf A$.
\end{proof}

\begin{theorem}[Positivity criterion]\label{thm:positivity}
Let $\lambda\in[0,1]$.  If $\zeta_\lambda$ has a representation \CGeqref{eq:representing-measure} with a positive measure $\mathsf A$, then
\begin{equation}\label{eq:positivity-bound}
 \varpi_\lambda=\mathsf A(\{1\})\le1 .
\end{equation}
In particular, if $\zeta_{Z1}$ has such a representation, then $\varpi_{Z1}=1$ and $\zeta_{Z1}=\zeta$.
\end{theorem}

\begin{proof}
By monotone convergence $\zeta_\lambda(\sigma)=\int u^{-\sigma}d\mathsf A\to\mathsf A(\{1\})$ as $\sigma\to+\infty$, and the limit is $\varpi_\lambda$ (Proposition~\CGref{prop:family-clock}).  Let $\mathsf T(\xi)=(1-|\xi|)_+$ and $\mathsf B=\mathsf T*\mathsf T$, so that $\mathsf B(\xi)=\frac23-\xi^2+\frac12|\xi|^3$ for $|\xi|\le1$, $\mathsf B(\xi)=\frac16(2-|\xi|)^3$ for $1\le|\xi|\le2$, $\mathsf B=0$ elsewhere, and $\mathsf B''(\xi)=\mathsf T(\xi+1)-2\mathsf T(\xi)+\mathsf T(\xi-1)$.  Put
\begin{equation}\label{eq:extremal-function}
 \mathsf g(u):=\Bigl(\frac{\sin\pi u}{\pi u}\Bigr)^4(1-u^2),\qquad \hat{\mathsf g}(\xi)=\mathsf B(\xi)+\frac1{4\pi^2}\mathsf B''(\xi) .
\end{equation}
Then $\mathsf g(0)=1$; $\mathsf g(u)\le0$ for $|u|\ge1$; $\hat{\mathsf g}(0)=\frac23-\frac1{2\pi^2}$ and $\hat{\mathsf g}(1)=\frac16+\frac1{4\pi^2}$; and $\hat{\mathsf g}\ge0$ on $\R$: for $0\le\xi\le1$, $\hat{\mathsf g}(\xi)=\frac23-\xi^2+\frac12\xi^3+\frac{3\xi-2}{4\pi^2}$ is positive at $\xi=0$ and $\xi=1$ and its derivative $-2\xi+\frac32\xi^2+\frac3{4\pi^2}$ changes sign once, from positive to negative, and for $1\le\xi\le2$, $\hat{\mathsf g}(\xi)=\frac16(2-\xi)^3+\frac{2-\xi}{4\pi^2}\ge0$.  For $\epsilon>0$ let $\mathsf g_\epsilon(u)=\mathsf g(u)e^{-\epsilon\pi u^2}$, an even Schwartz function with $\mathsf g_\epsilon(0)=1$, $\mathsf g_\epsilon\le0$ on $|u|\ge1$, and $\hat{\mathsf g}_\epsilon=\hat{\mathsf g}*(\epsilon^{-1/2}e^{-\pi\xi^2/\epsilon})\ge0$.  By Lemma~\CGref{lem:summation-positive}, since $\mathsf A\ge0$ is carried by $[1,\infty)$ and has mass $\varpi_\lambda$ at $u=1$,
\begin{equation}\label{eq:lp-inequality}
 1\ \ge\ \mathsf g_\epsilon(0)+2\int\mathsf g_\epsilon\,d\mathsf A=\hat{\mathsf g}_\epsilon(0)+2\int\hat{\mathsf g}_\epsilon\,d\mathsf A\ \ge\ \hat{\mathsf g}_\epsilon(0)+2\varpi_\lambda\,\hat{\mathsf g}_\epsilon(1) .
\end{equation}
Letting $\epsilon\to0$, $\hat{\mathsf g}_\epsilon\to\hat{\mathsf g}$ pointwise because $\hat{\mathsf g}$ is continuous, and
\begin{equation}\label{eq:lp-value}
 \varpi_\lambda\le\frac{1-\hat{\mathsf g}(0)}{2\hat{\mathsf g}(1)}=\frac{\frac13+\frac1{2\pi^2}}{\frac13+\frac1{2\pi^2}}=1 .
\end{equation}
The last statement follows from Theorem~\CGref{thm:varpi-Z1}.
\end{proof}

\begin{corollary}[Positive form of the final step]\label{cor:positivity}
The following are equivalent: \textup{(i)} (D); \textup{(ii)} $\zeta_{Z1}$ is the Mellin--Stieltjes transform \CGeqref{eq:representing-measure} of a positive measure on $[1,\infty)$; \textup{(iii)} $J^{\rm M}_{Z1}$, extended by $0$ to $(0,1]$, is non-decreasing on $(0,\infty)$.
\end{corollary}

\begin{proof}
(i)$\Rightarrow$(iii): then $J^{\rm M}_{Z1}=J$ (Theorem~\CGref{thm:cell-form}).  (iii)$\Rightarrow$(ii): $J^{\rm M}_{Z1}(x)x^{-c}\in L^2((1,\infty),dx/x)$ for every $c>1$, so a non-decreasing $J^{\rm M}_{Z1}$ is $O_c(x^{c})$ and $dJ^{\rm M}_{Z1}$ is a positive measure on $[1,\infty)$ with $\int x^{-\sigma}dJ^{\rm M}_{Z1}<\infty$ for $\sigma>1$.  By Theorem~\CGref{thm:mellin-counts}, $\zeta_{Z1}(s)=\exp\int x^{-s}dJ^{\rm M}_{Z1}=\int u^{-s}d\mathsf A(u)$, where $\mathsf A=\sum_{r\ge0}\frac1{r!}(dJ^{\rm M}_{Z1})^{*r}$ (multiplicative convolution powers) is a positive measure on $[1,\infty)$.  (ii)$\Rightarrow$(i) is Theorem~\CGref{thm:positivity}.
\end{proof}

In Theorem~\CGref{thm:positivity} the Dirichlet series is not used: (H1) is replaced by the positivity of the representing measure, the functional equation gives $\varpi_{Z1}\le1$ through \CGeqref{eq:lp-inequality}, and the Hadamard product gives $\varpi_{Z1}\ge1$ (Theorem~\CGref{thm:varpi-Z1}).  The function \CGeqref{eq:extremal-function} is extremal for the integers: for $\mathsf A=\sum_{n\ge1}\delta_n$ both inequalities in \CGeqref{eq:lp-inequality} are equalities in the limit, since $\mathsf g(n)=0$ for $n\ge1$ and $\hat{\mathsf g}(n)=0$ for $n\ge2$.  Under Hypothesis~\CGref{hyp:varpi-less-one}, $\varpi_{\RVM}<1$ at the clock end; unconditionally, the density of $dJ_{\RVM}$ takes negative values in every cell $(n,n+1)$ with $n$ large (Theorem~\CGref{thm:clock-density}).

\begin{lemma}[Mellin form and summation formula for every member]\label{lem:summation-all}
There is $c_1\ge2$ such that for every $\lambda\in[0,1]$ there is a unique measurable function $\mathsf A_\lambda$ on $(0,\infty)$, vanishing on $(0,1)$, with $x^{-c}\mathsf A_\lambda\in L^2\bigl((1,\infty),dx/x\bigr)$ for every $c>c_1$ and
\begin{equation}\label{eq:mellin-A}
 \zeta_\lambda(s)=s\int_1^\infty\mathsf A_\lambda(x)\,x^{-s-1}\,dx\qquad(\Re s>c_1).
\end{equation}
If $\zeta_\lambda$ has a representation \CGeqref{eq:representing-measure}, then $\mathsf A_\lambda(x)=\mathsf A([1,x])$ almost everywhere.  With the pairing $\langle\mathsf A_\lambda,\mathsf g\rangle:=-\int_1^\infty\mathsf A_\lambda(x)\,\mathsf g'(x)\,dx$, every even Schwartz function $\mathsf g$ satisfies
\begin{equation}\label{eq:summation-all}
 \mathsf g(0)+2\langle\mathsf A_\lambda,\mathsf g\rangle=\hat{\mathsf g}(0)+2\langle\mathsf A_\lambda,\hat{\mathsf g}\rangle .
\end{equation}
Equivalently, the even tempered distribution $\delta_0+d\mathsf A_\lambda(|u|)$ on $\R$ is its own Fourier transform.
\end{lemma}

\begin{proof}
By Proposition~\CGref{prop:family-clock} the bound \CGeqref{eq:Z1-growth} holds for every $\zeta_\lambda$ with constants independent of $\lambda$, so, as in the proof of Proposition~\CGref{prop:Mellin-Mertens}, $|\zeta_\lambda(s)|\le C(|t|+\sigma)^{1/4}$ for $\Re s>c_1$, and the argument of Theorem~\CGref{thm:mellin-counts} applied to $\zeta_\lambda(s)/s$ gives $\mathsf A_\lambda$, its vanishing on $(0,1)$ and its uniqueness.  For a positive measure $\mathsf A$ as in \CGeqref{eq:representing-measure}, partial integration gives \CGeqref{eq:mellin-A} with $\mathsf A([1,x])$, and uniqueness applies.

Let $c>c_1$.  In the variable $u=\log x$, the function $\mathsf A_\lambda(e^u)e^{-cu}$ has Fourier transform $\zeta_\lambda(c+it)/(c+it)$, and $-\mathsf g'(e^u)e^{(c+1)u}$ has Mellin transform $\int_0^\infty(-\mathsf g'(x))x^{s}dx=s\,\mathsf g^{\sharp}(s)$ on $\Re s=c$; Parseval's formula gives
\begin{equation}\label{eq:parseval-A}
 \langle\mathsf A_\lambda,\mathsf g\rangle=\frac1{2\pi i}\int_{(c)}\zeta_\lambda(s)\,\mathsf g^{\sharp}(s)\,ds .
\end{equation}
$\zeta_\lambda$ is bounded by a power of $|t|$ on every vertical strip (finite order, the functional equation and the bound on $\Re s\ge c$), and $\mathsf g^{\sharp}$ decays faster than any power.  Move the line to $\Re s=1-c$.  The poles of $\mathsf g^{\sharp}$ at $s=-2,-4,\dots$ are cancelled by the trivial zeros of $\zeta_\lambda$; the remaining poles are $s=1$ and $s=0$, exactly as in \CGeqref{eq:summation-shift}.  On the new line \CGeqref{eq:Z1-FE} and \CGeqref{eq:tate-local} give $\zeta_\lambda(s)\mathsf g^{\sharp}(s)=\zeta_\lambda(1-s)\hat{\mathsf g}^{\sharp}(1-s)$, and \CGeqref{eq:parseval-A} for $\hat{\mathsf g}$ gives \CGeqref{eq:summation-all}.
\end{proof}

\begin{proposition}[The value at infinity as one number]\label{prop:one-number}
Let $\mathsf g_\epsilon$ be as in the proof of Theorem~\CGref{thm:positivity}, $\mathsf w_\epsilon:=\hat{\mathsf g}_\epsilon-\mathsf g_\epsilon$, which is non-negative on $|u|\ge1$, and $\mathsf A'_\lambda(x):=\mathsf A_\lambda(x)-\varpi_\lambda$ for $x\ge1$.  For every $\lambda\in[0,1]$ and $\epsilon>0$,
\begin{equation}\label{eq:one-number-eps}
 \varpi_\lambda\,\mathsf w_\epsilon(1)=\tfrac12\bigl(1-\hat{\mathsf g}_\epsilon(0)\bigr)-\langle\mathsf A'_\lambda,\mathsf w_\epsilon\rangle,\qquad \langle\mathsf A'_\lambda,\mathsf w_\epsilon\rangle:=-\int_1^\infty\mathsf A'_\lambda(x)\,\mathsf w_\epsilon'(x)\,dx,
\end{equation}
the limit $\mathsf W(\lambda):=\lim_{\epsilon\to0}\langle\mathsf A'_\lambda,\mathsf w_\epsilon\rangle$ exists, and
\begin{equation}\label{eq:one-number}
 \varpi_\lambda-1=-\frac{\mathsf W(\lambda)}{\hat{\mathsf g}(1)},\qquad \hat{\mathsf g}(1)=\frac16+\frac1{4\pi^2}.
\end{equation}
Consequently $\mathsf W(0)\le0$; $\mathsf W(1)>0$ if and only if $\varpi_{\RVM}<1$ (Hypothesis~\CGref{hyp:varpi-less-one}); and (D) holds if and only if $\mathsf W(0)\ge0$.  If $\mathsf A_\lambda$ is non-decreasing, then $\langle\mathsf A'_\lambda,\mathsf w_\epsilon\rangle=\int_{(1,\infty)}\mathsf w_\epsilon\,d\mathsf A_\lambda\ge0$, which is Theorem~\CGref{thm:positivity}.
\end{proposition}

\begin{proof}
By \CGeqref{eq:summation-all} for $\mathsf g_\epsilon$, with $\mathsf g_\epsilon(0)=1$, $2\langle\mathsf A_\lambda,\mathsf w_\epsilon\rangle=1-\hat{\mathsf g}_\epsilon(0)$.  Since $\mathsf A_\lambda=\varpi_\lambda+\mathsf A'_\lambda$ on $[1,\infty)$ and $-\int_1^\infty\mathsf w_\epsilon'=\mathsf w_\epsilon(1)$, $\langle\mathsf A_\lambda,\mathsf w_\epsilon\rangle=\varpi_\lambda\mathsf w_\epsilon(1)+\langle\mathsf A'_\lambda,\mathsf w_\epsilon\rangle$; this is \CGeqref{eq:one-number-eps}.  As $\epsilon\to0$, $\mathsf w_\epsilon(1)=\hat{\mathsf g}_\epsilon(1)\to\hat{\mathsf g}(1)$, because $\mathsf g_\epsilon(1)=0$, and $\frac12(1-\hat{\mathsf g}_\epsilon(0))\to\frac12(1-\hat{\mathsf g}(0))=\hat{\mathsf g}(1)$ by \CGeqref{eq:lp-value}; hence the limit exists and \CGeqref{eq:one-number} holds.  The sign statements follow from \CGeqref{eq:one-number}: at $\lambda=0$ from $\varpi_{Z1}\ge1$ (Theorem~\CGref{thm:varpi-Z1}), and at $\lambda=1$ from $\varpi_1=\varpi_{\RVM}$, so that $\mathsf W(1)>0$ is the same statement as $\varpi_{\RVM}<1$.
\end{proof}

\begin{theorem}[The endpoint dominates $\zeta$ on the real axis]\label{thm:real-axis}
For every real $\sigma>1$,
\begin{equation}\label{eq:real-axis}
 \zeta_{Z1}(\sigma)\ \ge\ \zeta(\sigma),
\end{equation}
and equality for one $\sigma>1$ holds if and only if $\zeta_{Z1}=\zeta$.  In particular
\begin{equation}\label{eq:sigma-two}
 \prod_{k\ge1}\Bigl(1+\frac2{\frac14+\gamma_k^2}\Bigr)\ \ge\ \frac{\pi}{3},\qquad\text{with equality if and only if }(\mathrm D).
\end{equation}
\end{theorem}

\begin{proof}
By the proof of Theorem~\CGref{thm:varpi-Z1}, $\zeta_{Z1}(\sigma)/\zeta(\sigma)=\mathscr R_{Z1}(\sigma)$ is the convergent product of the factors \CGeqref{eq:pair-ratio}, with $W=(\sigma-\frac12)^2+\gamma^2$.  At $\sigma=1$, $W=\frac14+\gamma^2$ and the factor equals $1$ by the definition of $c_\gamma$.  The logarithmic derivative in $W$ of $W_\rho^2/((W_\rho-d_\rho^2)^2+4d_\rho^2\gamma^2)$ is
\begin{equation}\label{eq:pair-monotone}
 \frac2W-\frac{2(W-d_\rho^2)}{(W-d_\rho^2)^2+4d_\rho^2\gamma^2}=\frac{2d_\rho^2\,(4\gamma^2+d_\rho^2-W)}{W\bigl((W-d_\rho^2)^2+4d_\rho^2\gamma^2\bigr)},
\end{equation}
positive for $W<4\gamma^2+d_\rho^2$ and negative afterwards, and the factor tends to $c_\gamma>1$ as $W\to\infty$.  A function that equals $1$ at $\sigma=1$, increases and then decreases to a limit greater than $1$ is greater than $1$ for every $\sigma>1$.  Hence every factor of the convergent product $\mathscr R_{Z1}(\sigma)$ is $\ge1$, and every factor attached to a pair with $d_\rho\ne0$ is $>1$; so $\mathscr R_{Z1}(\sigma)\ge1$, with equality if and only if no such factor exists, that is, if and only if the product is empty and $\mathscr R_{Z1}\equiv1$.  For $\sigma=2$: $q(2)=2$, $G(2)=2/\pi$ and $P_\zeta(2)=G(2)\zeta(2)=\pi/3$, so \CGeqref{eq:real-axis} reads $P_{Z1}(2)\ge\pi/3$.
\end{proof}

\begin{proposition}[The final step on the open path]\label{prop:open-path-varpi}
Assume Hypothesis~\CGref{hyp:varpi-less-one}, $\varpi_{\RVM}<1$. Then the following are equivalent: \textup{(i)} (D); \textup{(ii)} $\varpi_\lambda<1$ for every $\lambda\in(0,1]$; \textup{(iii)} $\varpi_{\lambda_m}\le1$ for some sequence $\lambda_m\to0^+$; \textup{(iv)} $\mathsf W(\lambda)>0$ for every $\lambda\in(0,1]$. Without the hypothesis, (ii)$\Rightarrow$(iii)$\Rightarrow$(i) and (ii)$\Leftrightarrow$(iv) still hold, while (ii) contains $\varpi_{\RVM}<1$ as its case $\lambda=1$, so the hypothesis cannot be dropped from (i)$\Rightarrow$(ii).
\end{proposition}

\begin{proof}
$f(\lambda)=\log\varpi_\lambda$ is continuous and strictly convex on $[0,1]$ (Proposition~\CGref{prop:family-flow}: $f'=-\Upsilon$, $\Upsilon'<0$), $f(1)=\log\varpi_{\RVM}<0$ by the hypothesis, and $f(0)\ge0$ (Theorem~\CGref{thm:varpi-Z1}).  If $f(0)=0$, convexity gives $f(\lambda)<(1-\lambda)f(0)+\lambda f(1)<0$ for $0<\lambda<1$.  If $f(\lambda_m)\le0$ along $\lambda_m\to0^+$, then $f(0)\le0$ by continuity, so $f(0)=0$, which is (D) by Theorem~\CGref{thm:varpi-Z1}.  (ii)$\Leftrightarrow$(iv) is \CGeqref{eq:one-number}.
\end{proof}

A statement that holds as an identity only at the endpoint cannot be reached as a limit of the same identity along the path, where it is false (Proposition~\CGref{prop:nonconstant}); an inequality can.  Proposition~\CGref{prop:open-path-varpi} is of this form: the inequality $\varpi_\lambda\le1$ on the open path passes to the limit, and the opposite inequality $\varpi_{Z1}\ge1$ holds at the endpoint.  By Theorem~\CGref{thm:real-axis} the same scheme is available at every real $\sigma>1$: (D) holds if and only if $\zeta_{Z1}(\sigma)\le\zeta(\sigma)$ for one $\sigma>1$.

\begin{remark}[The hypothesis added to the functional equation]\label{rem:hypothesis-X}
The functional equation, the order and the zeros on the critical line do not characterize $\zeta$: every $\zeta_\lambda$ has them.  A converse theorem for $\zeta_{Z1}$ needs one further hypothesis on the function $\mathsf A_{Z1}$ of Lemma~\CGref{lem:summation-all}, and each of the following suffices and is equivalent to (D): $\mathsf A_{Z1}$ is a step function with jumps $\mathsf b_n$ at the integers satisfying $\sum_{n\le X}|\mathsf b_n|=O_\eps(X^{1+\eps})$, so that $\zeta_{Z1}$ is a Dirichlet series absolutely convergent in $\Re s>1$ (Hamburger's theorem, in the form of Theorem~\CGref{thm:cell-form}; the growth condition on the jumps is part of the hypothesis, since Hamburger's theorem requires the absolute convergence); $\mathsf A_{Z1}$ is non-decreasing (Theorem~\CGref{thm:positivity}); $\mathsf W(0)\ge0$ (Proposition~\CGref{prop:one-number}), which is the weakest of the three, being equivalent to $\varpi_{Z1}\le1$.  The inequality \CGeqref{eq:lp-inequality} is the one-dimensional case of the linear programming bound of Cohn and Elkies \cite{CohnElkies2003}.  By Lemma~\CGref{lem:summation-all} the measure $\delta_0+d\mathsf A_\lambda(|u|)$ is self-dual for every $\lambda$; measures on $\R$ whose support and spectrum are both uniformly discrete are classified in \cite{LevOlevskii2015}.
\end{remark}

\subsection{The density of the clock}
\label{subsec:clock-density}

Let $\varXi(t):=t\log t-t$ for $t>0$, and for $M\ge1$, $0<t_0\le\frac12$ let
\begin{equation}\label{eq:U-def}
 \mathcal U(M,t_0):=\int_{t_0}^{\infty}e^{-iM\varXi(t)}\,dt:=\lim_{T\to\infty}\int_{t_0}^{T}e^{-iM\varXi(t)}\,dt .
\end{equation}

\begin{lemma}[Stationary phase for $\varXi$]\label{lem:stationary-U}
The limit \CGeqref{eq:U-def} exists, and there is an absolute constant $C_{\mathcal U}$ with
\begin{equation}\label{eq:U-asymptotic}
 \Bigl|\mathcal U(M,t_0)-\sqrt{\frac{2\pi}M}\,e^{i(M-\pi/4)}\Bigr|\le\frac{C_{\mathcal U}}M\qquad(M\ge1,\ 0<t_0\le\tfrac12).
\end{equation}
\end{lemma}

\begin{proof}
$\varXi'(t)=\log t$ is increasing and unbounded, so the limit exists by the first-derivative test.  $\varXi$ is convex with minimum $\varXi(1)=-1$.  Put $\mathsf w:=\operatorname{sgn}(t-1)\sqrt{2(\varXi(t)+1)}$.  Near $t=1$, $\varXi(t)+1=\frac12(t-1)^2(1+O(t-1))$ with an analytic factor, so $t\mapsto\mathsf w$ is analytic with derivative $1$ at $t=1$; elsewhere $\mathsf w\,d\mathsf w=\log t\,dt$.  Hence $t\mapsto\mathsf w$ is an increasing analytic bijection of $(0,\infty)$ onto $(-\sqrt2,\infty)$, with
\begin{equation}\label{eq:U-substitution}
 \mathsf p(\mathsf w):=\frac{dt}{d\mathsf w}=\frac{\mathsf w}{\log t},\qquad \mathsf p(0)=1,\qquad \mathsf q(\mathsf w):=\frac{\mathsf p(\mathsf w)-1}{\mathsf w}=\frac1{\log t}-\frac1{\mathsf w}.
\end{equation}
$\mathsf q$ is analytic on $(-\sqrt2,\infty)$, bounded, tends to $0$ as $\mathsf w\to\infty$, and $\mathsf q'=-\mathsf w/(t\log^3t)+\mathsf w^{-2}$.  Away from $\mathsf w=0$ one has $\int|\mathsf w|\,|t\log^3t|^{-1}d\mathsf w=\int dt/(t\log^2t)$, which converges at $t=0$ and at $t=\infty$, and $\int\mathsf w^{-2}d\mathsf w$ converges at $\infty$; so $\int_{-\sqrt2}^\infty|\mathsf q'|\,d\mathsf w<\infty$.  With $\mathsf w_0=\mathsf w(t_0)\le\mathsf w(\frac12)=-\sqrt{1-\log2}<-\frac12$ and $e^{-iM\varXi}=e^{iM}e^{-iM\mathsf w^2/2}$,
\begin{equation}\label{eq:U-split}
 e^{-iM}\,\mathcal U(M,t_0)=\int_{\mathsf w_0}^\infty e^{-iM\mathsf w^2/2}\,d\mathsf w+\int_{\mathsf w_0}^\infty e^{-iM\mathsf w^2/2}\,\mathsf w\,\mathsf q(\mathsf w)\,d\mathsf w .
\end{equation}
The first integral is $\int_{\R}e^{-iM\mathsf w^2/2}d\mathsf w=\sqrt{2\pi/M}\,e^{-i\pi/4}$ minus $\int_{-\infty}^{\mathsf w_0}$, and the latter is at most $4/(M|\mathsf w_0|)\le8/M$ by the first-derivative test \cite[Lemma~4.2]{Titchmarsh1986}.  Since $\mathsf w\,e^{-iM\mathsf w^2/2}=\frac iM\frac d{d\mathsf w}e^{-iM\mathsf w^2/2}$, integration by parts bounds the second integral by $M^{-1}\bigl(\sup|\mathsf q|+\int|\mathsf q'|\bigr)$.
\end{proof}

For $m\ge1$ and $y>1$ put
\begin{equation}\label{eq:clock-harmonic}
 \mathcal R_m(y):=-2(-1)^m\,m^{-3/2}\,y^{\frac1{2m}-\frac12}\cos\Bigl(2\pi m\,y^{1/m}-\frac{(7m+1)\pi}4\Bigr),
\end{equation}
so that $\mathcal R_1(y)=2\cos(2\pi y)$ and $\mathcal R_2(y)=-2^{-1/2}y^{-1/4}\cos(4\pi\sqrt y+\frac\pi4)$.

\begin{theorem}[The density of the clock]\label{thm:clock-density}
$J_{\RVM}$ is continuously differentiable on $(2,\infty)$, and for $y\ge16$
\begin{equation}\label{eq:clock-density}
 J_{\RVM}'(y)=\frac1{\log y}+2\cos(2\pi y)-\frac{y^{-1/4}}{\sqrt2}\cos\Bigl(4\pi\sqrt y+\frac\pi4\Bigr)+O\bigl(y^{-1/3}\bigr).
\end{equation}
On $[2,\infty)$ the measure $d\pi_{\RVM}$ is the sum of the atoms $\frac{\mu(j)}jJ_{\RVM}(2)\,\delta_{2^j}$, $j\ge1$, and of a density that is continuous off the points $2^j$ and equals $J_{\RVM}'(y)+O(y^{-1/2})$.
\end{theorem}

\begin{proof}
\emph{An exact series.}  We use the notation of the proof of Theorem~\CGref{thm:unit-cell}: $v=\log y$, $S_\alpha(v)=\sum_ke^{(iv-\alpha)\tau_k}$, the decomposition \CGeqref{eq:S-eps-split} and the harmonics $K_m^\pm$ of \CGeqref{eq:K-harmonics}, with phases $\Psi_m^\pm(\tau)=v\tau\pm2\pi m\NR(\tau)$.  Put
\begin{equation}\label{eq:H-harmonic}
 \mathcal H_m^\pm(v):=\int_{\tau_0}^\infty e^{i\Psi_m^\pm(\tau)}\,d\tau ,
\end{equation}
improper integrals that converge, locally uniformly in $v\ge\log2$, by the first-derivative test, because $(\Psi_m^\pm)'$ is monotone and unbounded.  As $\alpha\downarrow0$, $K_m^\pm\to iv\,\mathcal H_m^\pm(v)$, and the first two terms of \CGeqref{eq:S-eps-split} tend to a continuous function $S_{\rm sm}(v)$ with $|S_{\rm sm}(v)|\le C/v+\frac12$.  The bounds for $K_m^\pm$ in that proof are uniform in the Abel parameter and, for $K_m^-$, include the factor displayed in \CGeqref{eq:K-minus-bound}; only the finite range $m\le v/\log2$ can contain the stationary point.  After this finite range is separated, the remaining bounds are summable against $1/(2\pi m)$, so dominated convergence in \CGeqref{eq:damped-sum} and \CGeqref{eq:K-harmonics}
\begin{equation}\label{eq:S-zero}
 \sum_kD_k=\int_{u_1}^{u_2}\frac{e^{v/2}}v\,S_0(v)\,dv,\qquad S_0(v):=S_{\rm sm}(v)+\sum_{m\ge1}\frac{(-1)^m\,v}{2\pi m}\bigl(\mathcal H_m^-(v)-\mathcal H_m^+(v)\bigr),
\end{equation}
the series converging uniformly on compact sets.  Thus $S_0$ is continuous, $J_{\RVM}$ is continuously differentiable on $(2,\infty)$, and by \CGeqref{eq:J-RVM}
\begin{equation}\label{eq:J-prime-exact}
 J_{\RVM}'(y)=\frac1{\log y}\Bigl(1-\frac1{y(y^2-1)}-\frac2{\sqrt y}\,\Re S_0(\log y)\Bigr).
\end{equation}

\emph{The terms without a stationary point.}  $(\Psi_m^+)'\ge v+m\log2$, so $|\mathcal H_m^+|\le4/(v+m\log2)$ and $\sum_m\frac v{2\pi m}|\mathcal H_m^+|\ll1+\log v$.  For $m>v/\log2$, $(\Psi_m^-)'\le v-m\log2<0$ on $[\tau_0,\infty)$; the first-derivative test, and the second-derivative test on $[\tau_0,2\tau_0]$ with $|(\Psi_m^-)''|=m/\tau$, give $|\mathcal H_m^-|\ll\min\bigl((m\log2-v)^{-1},m^{-1/2}\bigr)$, and $\sum_{m>v/\log2}\frac v{2\pi m}|\mathcal H_m^-|\ll1+\log v$.

\emph{The terms $3\le m\le v/\log2$.}  By the argument for \CGeqref{eq:K-minus-bound}, $|\mathcal H_m^-|\le\frac8{m\log2}+16\sqrt\pi\,e^{v/(2m)}m^{-1/2}$, so these terms contribute to $S_0$ at most $C\,v\sum_{m\ge3}\bigl(m^{-2}+m^{-3/2}e^{v/(2m)}\bigr)\ll v\,e^{v/6}$.

\emph{The terms $m=1,2$.}  The stationary point of $\Psi_m^-$ is $\mathsf t_m:=2\pi e^{v/m}=2\pi y^{1/m}$.  With $\tau=\mathsf t_mt$ and \CGeqref{eq:clock}, $\Psi_m^-(\mathsf t_mt)=-m\mathsf t_m\varXi(t)-\frac{7\pi m}4$, hence
\begin{equation}\label{eq:H-by-U}
 \mathcal H_m^-(v)=\mathsf t_m\,e^{-7\pi im/4}\,\mathcal U\bigl(m\mathsf t_m,\tau_0/\mathsf t_m\bigr),\qquad \tau_0/\mathsf t_m=2e^{-v/m}\le\tfrac12\quad(y\ge4^m).
\end{equation}
By Lemma~\CGref{lem:stationary-U} with $M=m\mathsf t_m$,
\begin{equation}\label{eq:H-asymptotic}
 \frac{(-1)^mv}{2\pi m}\,\mathcal H_m^-(v)=(-1)^m\,\frac v{m^{3/2}}\,y^{\frac1{2m}}\exp\Bigl(i\Bigl(2\pi m\,y^{1/m}-\frac{(7m+1)\pi}4\Bigr)\Bigr)+O\Bigl(\frac v{m^2}\Bigr).
\end{equation}
Inserting the four estimates into \CGeqref{eq:J-prime-exact} gives $J_{\RVM}'(y)=\frac1{\log y}+\mathcal R_1(y)+\mathcal R_2(y)+O(y^{-1/3})$, which is \CGeqref{eq:clock-density}.

\emph{The prime coordinate.}  By \CGeqref{eq:pi-RVM}, on each interval $(2^r,2^{r+1})$ the function $\pi_{\RVM}$ is the finite sum of the continuously differentiable functions $\frac{\mu(j)}jJ_{\RVM}(x^{1/j})$, $j\le r$; the terms $j\ge2$ have derivative $\frac{\mu(j)}{j^2}x^{1/j-1}J_{\RVM}'(x^{1/j})$, and $|J_{\RVM}'|\le C_J$ by \CGeqref{eq:unit-cell}, so their sum is $O(x^{-1/2})$.  At $x=2^j$ the term $j$ jumps by $\frac{\mu(j)}jJ_{\RVM}(2)$.
\end{proof}

\begin{corollary}[Cell masses and cell defects at the clock end]\label{cor:clock-cells}
For $n\ge17$, with $\mathbf 1_{2^j}(n)=1$ if $n=2^j$ and $0$ otherwise,
\begin{equation}\label{eq:clock-cell-sharp}
 c_n^{\RVM}=\int_{n-1}^n\Bigl(\frac1{\log x}+2\cos(2\pi x)\Bigr)dx+\mathbf 1_{2^j}(n)\,\frac{\mu(j)}j\,J_{\RVM}(2)+O\bigl(n^{-1/4}\bigr),
\end{equation}
\begin{equation}\label{eq:clock-cell-smooth}
 \tilde c_n(1)=\frac1{\log n}+\mathbf 1_{2^j}(n)\,\frac{\mu(j)}j\,J_{\RVM}(2)+O\bigl(n^{-1/4}\bigr),
\end{equation}
and for $0<\Re s\le\sigma_1$
\begin{equation}\label{eq:clock-defect}
 |\mathfrak h_n(1;s)|\le C_{\sigma_1}\bigl(|s|\,n^{-\sigma-4/3}+|s|(1+|s|)\,n^{-\sigma-2}\bigr).
\end{equation}
In particular Hypothesis~\CGref{hyp:cell-remainder} is proved (it holds along the integers $n$ that are not of the form $2^j$ with $\mu(j)\ne0$, which is its whole range), so Proposition~\CGref{prop:eventual-positive} is unconditional, and $c_n^{\RVM}\sim1/\log n$, $0<c_n^{\RVM}<1$ for all large such $n$; at $n=2^j$, $c_n^{\RVM}=\bigl(1+\mu(j)J_{\RVM}(2)\log2\bigr)/\log n+O(n^{-1/4})$.
\end{corollary}

\begin{proof}
\CGeqref{eq:clock-cell-sharp}: $c_n^{\RVM}=\pi_{\RVM}(n)-\pi_{\RVM}(n-1)$ is the integral of the density over $(n-1,n)$ plus the atom at $n$, and $\int_{n-1}^n(|\mathcal R_2|+O(x^{-1/3}))dx=O(n^{-1/4})$.  \CGeqref{eq:clock-cell-smooth}: $\varphi_n(2^j)=\mathbf 1_{2^j}(n)$; $\int\varphi(t)\,dt/\log(n+t)=\frac1{\log n}+O(\frac1{n\log^2n})$ because $\mathsf M_0=1$; and $\int\varphi(t)\,2\cos(2\pi(n+t))\,dt=2\mathsf C_0=0$ by \CGeqref{eq:partition-fourier}.  \CGeqref{eq:clock-defect}: the test function $\varphi_n(x)(n^{-s}-x^{-s})$ vanishes at the atoms, and $n^{-s}-(n+t)^{-s}=s\,n^{-s-1}t+O_{\sigma_1}(|s|(1+|s|)n^{-\sigma-2})$ for $|t|\le1$.  The first moment of the density is $\int t\varphi(t)\pi_{\RVM}'(n+t)\,dt=\frac{\mathsf M_1}{\log n}+2\mathsf C_1+O(n^{-1/3})=O(n^{-1/3})$, because $\mathsf M_1=\mathsf C_1=0$, $\mathcal R_2(n+t)=\mathcal R_2(n)+O(n^{-3/4})$ and the density is bounded.  The last statements follow from \CGeqref{eq:clock-cell-sharp} and Lemma~\CGref{lem:cell-cancel}.
\end{proof}

By \CGeqref{eq:clock-density}, $\frac{d\pi_{\RVM}}{dy}=\bigl(\frac1{\log y}-1\bigr)+\bigl(1+2\cos(2\pi y)\bigr)+O(y^{-1/4})$: the clock carries the first harmonic of the comb $\sum_n\delta_n=1+2\sum_{q\ge1}\cos(2\pi qy)$ of all integers, with weight $1$, on a smooth background.  This harmonic comes from the ordinates $\tau_k$ near the height $2\pi y$, the stationary point of $\Phi_1^-$; by \CGeqref{eq:partition-fourier} it does not enter the cell masses \CGeqref{eq:clock-cell-smooth}, which are $1/\log n$ to leading order, and it enters the cell defects only through the moments $\mathsf C_r$, $r\ge2$.

\subsection{The cell masses in the heights of the ordinates}
\label{subsec:height-kernels}

\begin{proposition}[Height kernel of the cell masses]\label{prop:height-kernels}
For every $\eps>0$, $\lambda\in[0,1]$ and $n\ge3$ that is not a power of $2$,
\begin{equation}\label{eq:height-kernel}
 \tilde c_n(\lambda)=\mathfrak m_n-\frac2{\sqrt n\,\log n}\sum_{k\ge1}\hat{\varphi}\Bigl(\frac{\gamma_k(\lambda)}{2\pi n}\Bigr)\cos\bigl(\gamma_k(\lambda)\log n\bigr)+E_n(\lambda),
\end{equation}
where $\sum_{X<n\le2X}|E_n(\lambda)|^2\le C_\eps X^{1/2+\eps}$ uniformly in $\lambda$ (sum over the $n$ that are not powers of $2$).  The kernel $\hat{\varphi}(\gamma/2\pi n)$ vanishes at the heights $\gamma=2\pi qn$, $q=1,2,\dots$
\end{proposition}

\begin{proof}
The term $j=1$ of $W_n$ in \CGeqref{eq:VW} is $\int\varphi(t)(n+t)^{-1/2+i\gamma}\,dt/\log(n+t)$.  For $0<\gamma\le n^{1+\eps}$, $(1+t/n)^{-1/2+i\gamma}=e^{i\gamma t/n}\bigl(1+O((1+\gamma/n)/n)\bigr)$ and $\frac1{\log(n+t)}=\frac1{\log n}(1+O(\frac1n))$ on $|t|\le1$, so this term equals $\frac{n^{-1/2+i\gamma}}{\log n}\hat{\varphi}(\frac\gamma{2\pi n})+O\bigl(n^{-3/2}(1+\gamma/n)\bigr)$, because $\int\varphi(t)e^{i\gamma t/n}dt=\hat{\varphi}(\gamma/2\pi n)$ for even $\varphi$.  Summed over the $O(n^{1+\eps}\log n)$ ordinates $\gamma_k(\lambda)\le n^{1+\eps}$ the error is $O(n^{-1/2+3\eps})$.  For $\gamma>n^{1+\eps}$ both the term and $\hat{\varphi}(\gamma/2\pi n)$ are $O_A((n/\gamma)^A)$, as in the proof of Theorem~\CGref{thm:mean-square-cells}.  The terms $j\ge2$ have sum of squares $O(X^{1/2+3\eps})$ over $X<n\le2X$ by \CGeqref{eq:cell-block}.  The zeros of the kernel are \CGeqref{eq:partition-fourier}.
\end{proof}

Since $\gamma_k(\lambda)\log n=\tau_k\log n-(1-\lambda)\eta_k\log n$, the passage along the family rotates each clock phase $n^{i\tau_k}$ by the angle $-(1-\lambda)\eta_k\log n$: up to the change of $\hat{\varphi}$ over a distance $O(1/n)$, the sum in \CGeqref{eq:height-kernel} is $\Re\,\Xi_n\bigl((1-\lambda)\log n\bigr)$ with
\begin{equation}\label{eq:Xi-rotation}
 \Xi_n(\Omega):=\sum_{k\ge1}\hat{\varphi}\Bigl(\frac{\tau_k}{2\pi n}\Bigr)n^{i\tau_k}e^{-i\Omega\eta_k}\qquad(0\le\Omega\le\log n).
\end{equation}
At $\Omega=0$ this is the clock sum, whose real part is $O(n^{1/4}\log n)$: the kernel of the term $j=1$ is, up to the $O(1/n)$ change of kernel, the ordinate part of $\int\varphi(t)J_{\RVM}'(n+t)\,dt$ in \CGeqref{eq:J-prime-exact}, and Theorem~\CGref{thm:clock-density} gives $\int\varphi(t)\bigl[2\cos(2\pi(n+t))+\mathcal R_2(n+t)+O((n+t)^{-1/3})\bigr]dt=O(n^{-1/4})$ pointwise because $C_0=0$ (the bound is not obtained from $\tilde c_n(1)$, which carries the terms $j\ge2$ whose remainder is controlled only in mean square), because the resonant heights $\tau_k\approx2\pi n$ are zeros of the kernel; the endpoint of the family is $\Omega=\log n$.

\begin{proposition}[The flow of a cell mass and the count difference]\label{prop:flow-count}
For every $n\ge3$ that is not a power of $2$,
\begin{equation}\label{eq:flow-count}
 \tilde c_n(0)-\tilde c_n(1)=-2\int_0^1\sum_{k\ge1}\eta_k\,\Im V_n\bigl(\gamma_k(\lambda)\bigr)\,d\lambda=-2\,\Im\int_{\tau_0}^\infty V_n(t)\bigl(N_\zeta(t)-n_{\RVM}(t)\bigr)\,dt .
\end{equation}
\end{proposition}

\begin{proof}
The first equality is \CGeqref{eq:cell-mass-W} and the continuity of $\tilde c_n$ at $\lambda=0$.  By \CGeqref{eq:cell-mass-W}, $\tilde c_n(0)-\tilde c_n(1)=2\Re\sum_k\bigl(W_n(\tau_k)-W_n(\gamma_k)\bigr)=2\Re\int_{\tau_0}^\infty W_n\,d(n_{\RVM}-N_\zeta)$.  Both counts vanish at $\tau_0$, their difference is $O(\log t)$ by Lemma~\CGref{lem:bounded-displacement}, and $W_n$ is rapidly decreasing; integration by parts with $W_n'=iV_n$ gives $2\Re\bigl(i\int V_n(N_\zeta-n_{\RVM})\,dt\bigr)$.
\end{proof}

\begin{proposition}[Height kernel of the cell defects]\label{prop:defect-kernel}
For every $\eps>0$, $\lambda\in[0,1]$, $\Re s>0$ and $n\ge3$ that is not a power of $2$,
\begin{equation}\label{eq:defect-kernel}
 \mathfrak h_n(\lambda;s)=-\frac{s\,n^{-s-3/2}}{\pi\log n}\sum_{k\ge1}\hat{\varphi}'\Bigl(\frac{\gamma_k(\lambda)}{2\pi n}\Bigr)\sin\bigl(\gamma_k(\lambda)\log n\bigr)+O_{s,\eps}\bigl(n^{-\sigma-3/4+\eps}\bigr).
\end{equation}
\end{proposition}

\begin{proof}
For $|t|\le1$, $n^{-s}-(n+t)^{-s}=s\,n^{-s-1}t+\psi_2(t)$, where $\psi_2$ and all its derivatives in $t$ are $O_s(n^{-\sigma-2})$.  \emph{The term $j=1$.}  The part of the density that does not involve the ordinates contributes $O_s(n^{-\sigma-2})$, because $\int t\varphi(t)\,dt/\log(n+t)=O(1/(n\log^2n))$ ($\mathsf M_1=0$).  The pairing of $\varphi\psi_2$ with the ordinate sum is $O_s(n^{-\sigma-2})$ times $O(n^{1/2+\eps})$, as in Remark~\CGref{rem:A6-scale}.  For the real function $t\varphi(t)$, the computation in the proof of Proposition~\CGref{prop:height-kernels}, with $\int t\varphi(t)e^{2\pi i\xi t}dt=-\frac i{2\pi}\hat{\varphi}'(\xi)$ for even $\varphi$ and $-2\Re\bigl(-\frac i{2\pi}n^{i\gamma}\bigr)=-\frac1\pi\sin(\gamma\log n)$, gives the sum in \CGeqref{eq:defect-kernel} with an error $O_s(n^{-\sigma-3/2+3\eps})$.  \emph{The terms $j\ge2$.}  The test function $\varphi_n(y^j)(n^{-s}-y^{-js})$ is $O(|s|n^{-\sigma-1})$ on a support of length $O(n^{1/j-1})$; with $O(n^{1+\eps}\log n)$ relevant ordinates and $y^{-1/2}=O(n^{-1/(2j)})$, these terms are $O_s(n^{-\sigma-1+1/(2j)+2\eps})$, at most $O_s(n^{-\sigma-3/4+2\eps})$.
\end{proof}

The kernel $\hat{\varphi}$ of the masses vanishes at the heights $2\pi qn$ and the kernel $\hat{\varphi}'$ of the defects in general does not: the masses and the defects of a cell read complementary heights of the ordinates.  At $\lambda=1$ the resonant heights $\tau_k\approx2\pi n$ contribute to $\sum_kn^{i\tau_k}$ the real quantity $-\sqrt n\log n\,e^{2\pi in}$ (proof of Theorem~\CGref{thm:clock-density}), which the sine in \CGeqref{eq:defect-kernel} does not see: the first harmonic of the clock is centred at the integers.

\begin{remark}[The scale of the displacements]\label{rem:scale-test}
By Proposition~\CGref{prop:nonconstant}, a property of the sequence $(\eta_k)$ that yields (D) is false for every sequence $(c\,\eta_k)$ with $0\le c<1$.  Properties that are invariant under $\eta_k\mapsto c\,\eta_k$ with $c>0$ (the signs of the $\eta_k$, the non-vanishing of the $\eta_k$, their boundedness) hold along the whole open path and do not suffice; the property fixes the size of the displacements.  In the form \CGeqref{eq:Xi-rotation} it is a condition on the angles $\eta_k\log n$ modulo $2\pi$, jointly in $k$ and $n$.  The construction contains one exact relation of this kind: for a simple ordinate, \CGeqref{eq:midpoint-condition} and \CGeqref{eq:block-count} give $\NR(\tau_k)-\NR(\gamma_k)=N_\zeta^{*}(\gamma_k)-\NR(\gamma_k)$, that is,
\begin{equation}\label{eq:eta-exact}
 \eta_k\,\NR'(\xi_k)=N_\zeta^{*}(\gamma_k)-\NR(\gamma_k)\qquad\text{for some }\xi_k\text{ between }\gamma_k\text{ and }\tau_k .
\end{equation}
\end{remark}

\subsection{The Gaussian defect of the atomization}
\label{subsec:gaussian-defect}

\begin{lemma}[Vertical growth of the joining factor]\label{lem:vertical-growth}
For every $\lambda\in[0,1]$ and $\frac12<\sigma_0<\sigma_1$, $\tilde h_\lambda(\sigma+it)=o(|t|)$ as $|t|\to\infty$, uniformly for $\sigma_0\le\sigma\le\sigma_1$.
\end{lemma}

\begin{proof}
By \CGeqref{eq:joining-factor} and Proposition~\CGref{prop:first-cell}(ii) it suffices to treat the double series.  Fix $c\in(\frac12,\sigma_0)$.  By \CGeqref{eq:cell-defect-bound}, $\sum_{m\ge1}\frac1m\sum_{n>N}|\mathfrak h_n(\lambda;ms)|\le C|s|\sum_{n>N}n^{-\sigma_0-1}(1+I_n^{1/2})$ and $\sum_{m>N}\frac1m\sum_{n\ge2}|\mathfrak h_n(\lambda;ms)|\le C|s|\sum_{n\ge2}n^{-N\sigma_0-1}(1+I_n^{1/2})$; both are at most $\epsilon_N|s|$ with $\epsilon_N\to0$ as $N\to\infty$.  For the finitely many pairs $m,n\le N$, with $E=\pi_\lambda-\mathfrak L$ as in the proof of Theorem~\CGref{thm:atomized-continuation},
\begin{equation}\label{eq:cell-RL}
\begin{split}
 \mathfrak h_n(\lambda;ms)={}&\int\varphi_n(x)\bigl(n^{-ms}-x^{-ms}\bigr)d\mathfrak L(x)-\int E(x)\varphi_n'(x)\bigl(n^{-ms}-x^{-ms}\bigr)dx\\
 &-ms\int E(x)\varphi_n(x)\,x^{-ms-1}dx .
\end{split}
\end{equation}
The first two terms are bounded in $t$, and the last integral tends to $0$ as $|t|\to\infty$ by the Riemann--Lebesgue lemma, because $E\in L^1(n-1,n+1)$.  Hence $\limsup_{|t|\to\infty}|t|^{-1}\sup_\sigma|\tilde h_\lambda(\sigma+it)|\le2\epsilon_N$ for every $N$.
\end{proof}

For $y>0$ and $c>1$ put
\begin{equation}\label{eq:gaussian-defect}
 \Theta_\lambda(y):=\frac1{4\pi i}\int_{(c)}\zeta_\lambda(s)\,\Gamma(\tfrac s2)\,\pi^{-s/2}\,y^{-s/2}\,ds,\qquad
 \mathfrak G_\lambda(y):=\sum_{N\ge1}\tilde a_N(\lambda)\,e^{-\pi N^2y}-\Theta_\lambda(y).
\end{equation}
For $\zeta$ itself $\Theta(y)=\sum_{n\ge1}e^{-\pi n^2y}$, and if $\zeta_\lambda$ has a representation \CGeqref{eq:representing-measure}, then $\Theta_\lambda(y)=\int e^{-\pi u^2y}\,d\mathsf A(u)$.  $\mathfrak G_\lambda$ is the change produced by the atomization in the Gaussian count of the integers of the system.

\begin{proposition}[Gaussian defect]\label{prop:gaussian-defect}
Let $\lambda\in[0,1]$.
\begin{enumerate}[label=\textup{(\roman*)}]
\item $\Theta_\lambda(y)=y^{-1/2}\,\Theta_\lambda(1/y)+\frac12\bigl(y^{-1/2}-1\bigr)$ for $y>0$, and $\Theta_\lambda(y)=O_A(y^{-A})$ as $y\to\infty$ for every $A$.
\item The coefficients $\tilde a_N(\lambda)$ satisfy the theta relation of $\zeta$,
\begin{equation}\label{eq:theta-atomized}
 \sum_{N\ge1}\tilde a_N(\lambda)\,e^{-\pi N^2y}=y^{-1/2}\sum_{N\ge1}\tilde a_N(\lambda)\,e^{-\pi N^2/y}+\tfrac12\bigl(y^{-1/2}-1\bigr)\qquad(y>0),
\end{equation}
if and only if
\begin{equation}\label{eq:defect-selfdual}
 \mathfrak G_\lambda(y)=y^{-1/2}\,\mathfrak G_\lambda(1/y)\qquad(y>0).
\end{equation}
Under the hypotheses of Theorem~\CGref{thm:pair-criterion}, \CGeqref{eq:defect-selfdual} holds.
\item For every $\delta>0$, as $y\to0^+$,
\begin{equation}\label{eq:defect-expansion}
 \mathfrak G_\lambda(y)=\tfrac12\bigl(e^{\tilde h_\lambda(1)}-1\bigr)\,y^{-1/2}+O_\delta\bigl(y^{-1/4-\delta}\bigr),
\end{equation}
and $e^{\tilde h_\lambda(1)}=\Res_{s=1}\tilde\zeta_{\Z,\lambda}(s)$.
\item If \CGeqref{eq:defect-selfdual} holds, then $\tilde h_\lambda(1)=0$.
\end{enumerate}
\end{proposition}

\begin{proof}
(i) $\zeta_\lambda$ is bounded by a power of $|t|$ on vertical strips (proof of Lemma~\CGref{lem:summation-positive}; on $\Re s\ge c$ use Proposition~\CGref{prop:family-clock}), and $\Gamma(\frac s2)$ decays exponentially.  Moving the line to $\Re s=1-c$ crosses $s=1$, with residue $\Gamma(\frac12)\pi^{-1/2}y^{-1/2}=y^{-1/2}$, and $s=0$, with residue $2\zeta_\lambda(0)=-1$; on the new line \CGeqref{eq:Z1-FE} and $y^{-s/2}=y^{-1/2}(1/y)^{-(1-s)/2}$ give $y^{-1/2}\Theta_\lambda(1/y)$.  The decay follows by moving the line to $\Re s=2A$.
(ii) By $e^{-x}=\frac1{2\pi i}\int_{(c/2)}\Gamma(w)x^{-w}dw$ and Corollary~\CGref{cor:abscissa-one}, $\sum_N\tilde a_N(\lambda)e^{-\pi N^2y}=\frac1{4\pi i}\int_{(c)}\tilde\zeta_{\Z,\lambda}(s)\Gamma(\frac s2)\pi^{-s/2}y^{-s/2}ds$; subtracting (i) from the theta relation gives the equivalence.  Under the hypotheses of Theorem~\CGref{thm:pair-criterion}, $\tilde\zeta_{\Z,\lambda}=\zeta$.
(iii) By (ii), $2\mathfrak G_\lambda(y)=\frac1{2\pi i}\int_{(c)}\zeta_\lambda(s)\bigl(e^{\tilde h_\lambda(s)}-1\bigr)\Gamma(\frac s2)\pi^{-s/2}y^{-s/2}ds$.  By Lemma~\CGref{lem:vertical-growth}, $|e^{\tilde h_\lambda(s)}|\le e^{o(|t|)}$ on $\frac12+2\delta\le\Re s\le c$, so the line can be moved to $\Re s=\frac12+2\delta$.  The only pole crossed is $s=1$, where $\zeta_\lambda$ has residue $1$ and $\tilde h_\lambda$ is holomorphic; the residue is $(e^{\tilde h_\lambda(1)}-1)y^{-1/2}$, and the remaining integral is $O_\delta(y^{-1/4-\delta})$.
(iv) $\sum_N\tilde a_N(\lambda)e^{-\pi N^2y}=O(e^{-\pi y})$ as $y\to\infty$, so by (i) $\mathfrak G_\lambda(y)=O_A(y^{-A})$ there; then \CGeqref{eq:defect-selfdual} gives $\mathfrak G_\lambda(y)=O_A(y^{A})$ as $y\to0^+$ for every $A$, and \CGeqref{eq:defect-expansion} forces $e^{\tilde h_\lambda(1)}=1$; $\tilde h_\lambda$ is real on the real axis.
\end{proof}

\begin{corollary}[The density of the atomized integers]\label{cor:residue-condition}
$\sum_{N\le x}\tilde a_N(\lambda)$ has mean density $e^{\tilde h_\lambda(1)}$ in the sense of \CGeqref{eq:defect-expansion}, and (D) implies
\begin{equation}\label{eq:residue-condition}
 \tilde h_0(1)=0,\qquad\text{that is,}\qquad \mathfrak u_0(1)=\sum_{m\ge1}\frac1m\sum_{n\ge2}\mathfrak h_n(0;m).
\end{equation}
\end{corollary}

\begin{proof}
Under (D), $\tilde h_0\equiv0$ (Theorem~\CGref{thm:closing}); the second form is \CGeqref{eq:joining-factor} at $s=1$.
\end{proof}

The identity $\tilde h_0(\sigma)=0$ for real $\sigma$ has two distinguished members: $\sigma=+\infty$ is $\varpi_{Z1}=1$, the first coefficient (Corollary~\CGref{cor:product-identity}), and $\sigma=1$ is \CGeqref{eq:residue-condition}, the density of the integers.

By Proposition~\CGref{prop:gaussian-defect}(i), the function \CGeqref{eq:theta-defect} is $T_\lambda(y)=\mathfrak G_\lambda(y)-y^{-1/2}\mathfrak G_\lambda(1/y)$, so \CGeqref{eq:defect-selfdual} is $T_\lambda\equiv0$.  At the clock end the expansion \CGeqref{eq:defect-expansion} holds with a smaller remainder: for every $\delta>0$,
\begin{equation}\label{eq:defect-expansion-clock}
 \mathfrak G_1(y)=\tfrac12\bigl(e^{\tilde h_1(1)}-1\bigr)\,y^{-1/2}+O_\delta\bigl(y^{-\delta}\bigr)\qquad(y\to0^+),
\end{equation}
because $\tilde h_1$ is holomorphic in $\Re s>0$ with $\tilde h_1(\sigma+it)=o(|t|)$ uniformly on $\sigma_0\le\sigma\le\sigma_1$ for every $\sigma_0>0$ (the proof of Lemma~\CGref{lem:vertical-growth} with the bound $|\mathfrak h_n(1;s)|\le C|s|n^{-\sigma-1}$ of Theorem~\CGref{thm:atomized-continuation}(iv)), and the line of integration can be moved to $\Re s=2\delta$.

\subsection{Entire dependence on \texorpdfstring{$\lambda$}{lambda} and the exceptional set}
\label{subsec:entire-lambda}

\begin{proposition}[The cell masses as entire functions of $\lambda$]\label{prop:entire-lambda}
Let $n\ge3$ not be a power of $2$.
\begin{enumerate}[label=\textup{(\roman*)}]
\item The real-axis function $\lambda\mapsto\tilde c_n(\lambda)$ has the entire extension
\[
 \widetilde C_n(\lambda):=\mathfrak m_n-\sum_k\Bigl(W_n(\gamma_k+\lambda\eta_k)+\overline{W_n(\gamma_k+\bar\lambda\eta_k)}\Bigr),
\]
which is of exponential type at most $C_\eta\log(n+1)$ and agrees with $\tilde c_n(\lambda)$ for real $\lambda$.
\item For every $r\ge0$, $\eps>0$ and $\lambda\in[0,1]$,
\begin{equation}\label{eq:lambda-derivatives}
 \partial_\lambda^r\tilde c_n(\lambda)=-2\,\Re\sum_{k\ge1}\eta_k^{\,r}\,W_n^{(r)}\bigl(\gamma_k(\lambda)\bigr),\qquad \bigl|\partial_\lambda^r\tilde c_n(\lambda)\bigr|\le C_{r,\eps}\bigl(C_\eta\log(n+1)\bigr)^r\,n^{1/2+\eps}.
\end{equation}
\item $|\tilde c_n(\lambda)-\tilde c_n(\lambda')|\le C_\eps|\lambda-\lambda'|\,n^{1/2+\eps}$ on $[0,1]$.  Hence, for sets of integers $n$ that are not powers of $2$, $\tilde c_n(0)=O_\eps(n^\eps)$ if and only if $\tilde c_n(\lambda)=O_\eps(n^\eps)$ uniformly for $0\le\lambda\le n^{-1/2}$, and $\tilde c_n(\lambda)=O_\eps(n^{\eps})$ uniformly for $1-n^{-1/2}\le\lambda\le1$.
\end{enumerate}
\end{proposition}

\begin{proof}
For complex $\gamma$, $|y^{i\gamma}|\le e^{|\Im\gamma|\log y}$ with $\log y\le\log(n+1)$ on the support of $\varphi_n(y^j)$, and repeated integration by parts in \CGeqref{eq:VW} gives $|W_n^{(r)}(\gamma)|\le C_{A,r,n}\,e^{|\Im\gamma|\log(n+1)}(1+|\Re\gamma|)^{-A}$ for every $A$.  With $\gamma=\gamma_k+\lambda\eta_k$, $|\Im\gamma|\le C_\eta|\Im\lambda|$, the two reflected series defining $\widetilde C_n$ converge locally uniformly in $\lambda\in\C$, so $\widetilde C_n$ is entire; on the real axis it equals $\tilde c_n$.  Termwise differentiation on the real axis gives the first formula in \CGeqref{eq:lambda-derivatives}.  For real $\lambda$, $W_n^{(r)}$ carries the factor $(i\log y)^r$ with $\log y\le\log(n+1)$, the ordinates $\gamma_k(\lambda)>n^{1+\eps}$ contribute $O(n^{-1})$, and each of the $O(n^{1+\eps}\log n)$ remaining ordinates contributes $O\bigl((\log(n+1))^rn^{-1+1/(2j)}\bigr)$ to the term $j$, as in Remark~\CGref{rem:A6-scale}; this is the bound in \CGeqref{eq:lambda-derivatives}.  (iii) is the case $r=1$ and Corollary~\CGref{cor:clock-cells}.
\end{proof}

The bounds \CGeqref{eq:lambda-derivatives} do not by themselves control $\tilde c_n(0)$ from $\tilde c_n(1)$: the function $n^{1/2}\sin\bigl((1-\lambda)C_\eta\log n\bigr)$ is entire of the same type, satisfies the same bounds, vanishes at $\lambda=1$ and has size $n^{1/2}$ at $\lambda=0$ for infinitely many $n$.  The Taylor expansion at the clock end, $\tilde c_n(0)=\sum_{r\ge0}\frac{(-1)^r}{r!}\partial_\lambda^r\tilde c_n(1)$, has the terms $-2\Re\sum_k\eta_k^rW_n^{(r)}(\tau_k)$, clock sums weighted by the powers of the displacements.

\begin{proposition}[The exceptional set]\label{prop:exceptional-set}
The set
\begin{equation}\label{eq:exceptional-set}
 \mathsf E:=\bigl\{\lambda\in[0,1]:\ \tilde\zeta_{\Z,\lambda}=\zeta\bigr\}
\end{equation}
is finite, $1\notin\mathsf E$, and (D) holds if and only if $0\in\mathsf E$.  For every $\lambda\notin\mathsf E$ the hypotheses of Theorem~\CGref{thm:pair-criterion} do not hold together, $T_\lambda\not\equiv0$ and $\mathsf V_\sigma(\lambda)>0$.
\end{proposition}

\begin{proof}
If $\lambda\in\mathsf E$, then $\sum_n\tilde c_n(\lambda)\sum_m\frac{n^{-ms}}m=\log\zeta(s)=\sum_p\sum_m\frac{p^{-ms}}m$ for $\Re s>1$.  Comparing the coefficients of $N^{-s}$ in increasing order of $N$ determines the exponents uniquely, so $\tilde c_n(\lambda)=\mathbf 1_{\rm prime}(n)$ for every $n$.  Choose a composite $n$, not a power of $2$, so large that $\tilde c_n(1)>0$ by \CGeqref{eq:clock-cell-smooth}.  Then $\tilde c_n$ is an entire function of $\lambda$ (Proposition~\CGref{prop:entire-lambda}) that does not vanish identically; the zeros of a nonzero entire function have no finite accumulation point, so only finitely many of them lie in the compact interval $[0,1]$; they contain $\mathsf E$, and $1$ is not one of them.  $0\in\mathsf E$ implies (D) by Theorem~\CGref{thm:double-hamburger}(ii), and (D) implies $\tilde\zeta_{\Z,0}=\zeta_{Z1}=\zeta$ by Theorem~\CGref{thm:closing}.  The last statement is Theorems~\CGref{thm:pair-criterion} and~\CGref{thm:double-hamburger}(ii).
\end{proof}

Proposition~\CGref{prop:nonconstant} excludes that $\zeta_\lambda$ and $\tilde\zeta_{\Z,\lambda}$ coincide for $\lambda>0$; Proposition~\CGref{prop:exceptional-set} excludes, outside a finite set, that $\tilde\zeta_{\Z,\lambda}$ coincides with $\zeta$.  For $0<\lambda<1$ outside $\mathsf E$ it is not decided which of the two hypotheses of Theorem~\CGref{thm:pair-criterion} fails; at $\lambda=1$ it is the symmetry (Proposition~\CGref{prop:clock-asymmetry}).


\subsection{The final step in the cell masses}
\label{subsec:cell-mass-form}

\begin{corollary}[The final step is one number per cell]\label{cor:cell-mass-form}
The following are equivalent:
\begin{enumerate}[label=\textup{(\roman*)}]
\item (D);
\item $\tilde c_n(0)=\mathbf 1_{\rm prime}(n)$ for every $n\ge2$;
\item there is a list $\boldsymbol{\kappa}\in\mathscr O$ (Section~\CGref{sec:lists}) with $\tilde c_n(\boldsymbol{\kappa})=\mathbf 1_{\rm prime}(n)$ for every $n\ge2$.
\end{enumerate}
\end{corollary}

\begin{proof}
(i)$\Rightarrow$(ii) is Theorem~\CGref{thm:closing}, and (ii)$\Rightarrow$(iii) with the list $(\gamma_k)$.  If (iii) holds, then
\begin{equation}\label{eq:cell-mass-euler}
 \tilde\zeta_{\Z,\boldsymbol{\kappa}}(s)=\prod_{n\ge2}\bigl(1-n^{-s}\bigr)^{-\tilde c_n(\boldsymbol{\kappa})}=\prod_p\bigl(1-p^{-s}\bigr)^{-1}=\zeta(s)\qquad(\Re s>1),
\end{equation}
and Corollary~\CGref{cor:existence-form}, (iii)$\Rightarrow$(i), gives (D).
\end{proof}

Step A7 of the programme asks for the vanishing of the functions $\mathfrak h_n(0;\cdot)$, one function of $s$ for each cell, and for $J_{Z1}=0$ on the first cell.  Corollary~\CGref{cor:cell-mass-form} asks for one real number for each cell; the position of the mass inside the cell and the first cell then follow from Hamburger's theorem.

\begin{proposition}[Every point of the exceptional set is a witness]\label{prop:witness}
For the set $\mathsf E$ of \CGeqref{eq:exceptional-set},
\begin{equation}\label{eq:witness}
 \mathsf E\ne\varnothing\iff(\mathrm D),\qquad\text{so that either }\mathsf E=\varnothing\ \text{ or }\ 0\in\mathsf E .
\end{equation}
If (D) holds, then
\begin{equation}\label{eq:witness-zeros}
 \mathsf E\setminus\{0\}=\bigl\{\lambda\in(0,1]:\ \mathsf D_n(\lambda)=0\ \text{for every }n\ge2\bigr\},\qquad \mathsf D_n(\lambda):=\tilde c_n(\lambda)-\tilde c_n(0),
\end{equation}
and for $n\ge3$ not a power of $2$, the real-axis function $\mathsf D_n(\lambda)=-2\Re\sum_{k\ge1}\bigl(W_n(\gamma_k+\lambda\eta_k)-W_n(\gamma_k)\bigr)$ has the entire extension obtained by replacing each real part by the corresponding holomorphic plus reflected-conjugate pair; this extension is of exponential type at most $C_\eta\log(n+1)$ that vanishes at $\lambda=0$.
\end{proposition}

\begin{proof}
If $\lambda_0\in\mathsf E$, the list $(\gamma_k(\lambda_0))$ belongs to $\mathscr O$ and its atomized zeta function is $\zeta$; Corollary~\CGref{cor:existence-form}, (iii)$\Rightarrow$(i), gives (D), and then $0\in\mathsf E$ by Proposition~\CGref{prop:exceptional-set}.  Under (D), $\tilde c_n(0)=\mathbf 1_{\rm prime}(n)$, and $\lambda\in\mathsf E$ means $\tilde c_n(\lambda)=\mathbf 1_{\rm prime}(n)$ for every $n$ (proof of Proposition~\CGref{prop:exceptional-set}).  The last statement is \CGeqref{eq:cell-mass-W} and Proposition~\CGref{prop:entire-lambda}.
\end{proof}

Thus a proof that $\mathsf E\cap(0,1]=\varnothing$ does not approach the final step: if (D) fails, $\mathsf E$ is empty, and if (D) holds, $\mathsf E\cap(0,1]=\varnothing$ is the uniqueness of the witness on the path, which is the statement that the entire functions $\mathsf D_n$ have no common zero in $(0,1]$.  This uniqueness is not proved; Proposition~\CGref{prop:exceptional-set} gives that $\mathsf E$ is finite and $1\notin\mathsf E$.

\begin{remark}[The dual function of the atomized family]\label{rem:dual-function}
For $\Re s<\frac12$ put $B^\vee_\lambda(s):=\zeta_\lambda(s)\,e^{\tilde h_\lambda(1-s)}$.  By the functional equation of $\zeta_\lambda$,
\begin{equation}\label{eq:dual-function}
 \pi^{-s/2}\Gamma\bigl(\tfrac s2\bigr)B^\vee_\lambda(s)=\pi^{-(1-s)/2}\Gamma\bigl(\tfrac{1-s}2\bigr)\,\tilde\zeta_{\Z,\lambda}(1-s)\qquad(\Re s<\tfrac12),
\end{equation}
so $B^\vee_\lambda$ is the function attached to $\tilde\zeta_{\Z,\lambda}$ by the functional equation.  To ask that $B^\vee_\lambda$ continue to $\Re s>1$ as an absolutely convergent Dirichlet series is to ask for the hypotheses of Theorem~\CGref{thm:pair-criterion}; no monotonicity in $\lambda$ is known, and a parameter such as $\inf\{\lambda>0:B^\vee_\lambda\ \text{is not a Dirichlet series}\}$ is not shown to be a transition point.
\end{remark}

\begin{proposition}[The identity that remains, cell by cell]\label{prop:per-cell}
Let $n\ge17$ not be a power of $2$.  Then $\tilde c_n(0)=\mathbf 1_{\rm prime}(n)$ if and only if
\begin{equation}\label{eq:per-cell}
 \mathbf 1_{\rm prime}(n)-\tilde c_n(1)=-2\,\Im\int_{\tau_0}^\infty V_n(t)\bigl(N_\zeta(t)-n_{\RVM}(t)\bigr)\,dt,\qquad \tilde c_n(1)=\frac1{\log n}+O\bigl(n^{-1/4}\bigr).
\end{equation}
\end{proposition}

\begin{proof}
This is \CGeqref{eq:flow-count} and \CGeqref{eq:clock-cell-smooth}.
\end{proof}

The left side of \CGeqref{eq:per-cell} is explicit.  The right side is the pairing of the count difference with the kernel $V_n$, whose term $j=1$ is $n^{-1/2+it}\hat{\varphi}(t/2\pi n)$ up to $O\bigl(n^{-3/2}(1+t/n)\bigr)$ (proof of Proposition~\CGref{prop:defect-kernel}).  For $n=2^j$ the same identity holds with the cut term of Theorem~\CGref{thm:mean-square-cells} added.  No evaluation of the right side of \CGeqref{eq:per-cell} is given in this paper.

\subsection{Symmetric joining factors}
\label{subsec:symmetric-factors}

\begin{remark}[The symmetry criterion has no slack at $\lambda=0$]\label{rem:pair-no-slack}
At $\lambda=0$ the hypothesis of Theorem~\CGref{thm:pair-criterion} is equivalent to its conclusion.  If (D) holds, then $\tilde h_0\equiv0$ (Theorem~\CGref{thm:closing}), and $e^{\tilde h_0}=1$ is symmetric and of finite order; hence
\begin{equation}\label{eq:pair-no-slack}
 e^{\tilde h_0}\ \text{symmetric of finite order}\iff\tilde h_0\equiv0\iff(\mathrm D),
\end{equation}
and in the same way $T_0\equiv0$ and $\mathsf V_\sigma(0)=0$ are equivalent forms of (D), not weaker conditions.  Moreover \CGeqref{eq:pair-symmetry} presupposes a continuation of $\tilde h_0$ across the line $\Re s=\frac12$ (step A10 of the programme): $\tilde h_0$ is holomorphic in $\Re s>\frac12$, and $s\mapsto1-s$ maps this half-plane onto $\Re s<\frac12$.  Theorems~\CGref{thm:halfline-hamburger} and~\CGref{thm:bridge} below are criteria whose hypotheses are written on the half-line $\sigma\to+\infty$ and on the closed half-plane $\Re s\ge\frac12$.
\end{remark}

\begin{proposition}[Symmetric joining factors in general]\label{prop:dichotomy}
Let $f$ be meromorphic in $\C$ with $(s-1)f(s)$ entire of order at most one, and let $f$ satisfy the functional equation \CGeqref{eq:Z1-FE}.  Let $g=f\,e^{h}$, where $e^{h}$ is entire, without zeros, of finite order, and $e^{h(s)}=e^{h(1-s)}$.  If $g\not\equiv0$ and $g(s)=\sum_{n\ge1}a_nn^{-s}$ converges absolutely in $\Re s>1$, then $g=C_H\zeta$ with $C_H\ne0$ and $e^{h}$ is constant.  If moreover $\lim_{\sigma\to+\infty}g(\sigma)=\lim_{\sigma\to+\infty}f(\sigma)\ne0$, then $e^{h}\equiv1$ and $f=g=C_H\zeta$.
\end{proposition}

\begin{proof}
$g$ satisfies (H0)--(H3) of Subsection~\CGref{subsec:hamburger-table} with the second function equal to $g$, so $g=C_H\zeta$ by Hamburger's theorem.  If $C_H=0$ then $f\,e^{h}=0$, which is excluded when $g\not\equiv0$; assume $C_H\ne0$.  Then $e^{h}\,(s-1)f(s)=C_H(s-1)\zeta(s)$ identically, and we compare the Hadamard factorizations of the two entire functions of order at most one, $(s-1)f(s)=e^{A_1+B_1s}\Pi_1(s)$ and $(s-1)\zeta(s)=e^{A_2+B_2s}\Pi_2(s)$, with $\Pi_i$ the canonical products of genus at most one over their zeros.  Since $e^{h}$ has no zeros, the two functions have the same zeros with the same multiplicities, so $\Pi_1=\Pi_2$ and $e^{h}=C_He^{(A_2-A_1)+(B_2-B_1)s}$: $h$ is a polynomial of degree at most one.  The symmetry gives $h(s)-h(1-s)\in2\pi i\Z$, which forces the degree to be zero.  The last statement compares the limits at $+\infty$.
\end{proof}

\begin{remark}[Order one cannot be dropped]\label{rem:dichotomy-counterexample}
For real $\nu_5\ne0$ put $f(s)=\zeta(s)\exp\bigl(-\nu_5(s-\tfrac12)^2\bigr)$ and $e^{h(s)}=\exp\bigl(\nu_5(s-\tfrac12)^2\bigr)$.  Then $f$ is meromorphic of order two with the functional equation, the pole and the zeros of $\zeta$; $e^{h}$ is entire, without zeros, of finite order and symmetric; $g=f\,e^{h}=\zeta$ is an absolutely convergent Dirichlet series in $\Re s>1$; and $e^{h}$ is not constant.  In Theorem~\CGref{thm:pair-criterion} the conclusion $\tilde h_0\equiv0$ uses that $(s-1)\zeta_{Z1}$ has order one (Proposition~\CGref{prop:Z1-class}).  The location of the zeros of $f$ plays no role in Proposition~\CGref{prop:dichotomy}.
\end{remark}

\subsection{The half-line form of Hamburger's theorem for a pair}
\label{subsec:halfline}

The next two theorems are stated for a general pair: a function with the functional equation and a Dirichlet series, joined on the real half-line only.

\begin{theorem}[The ladder]\label{thm:ladder}
Let $\mathsf d\colon[1,\infty)\to\R$ be measurable with $\int_1^\infty|\mathsf d(x)|\,x^{-\mathsf e_1-1}\,dx<\infty$ for some $\mathsf e_1\ge0$, and put
\begin{equation}\label{eq:ladder-h}
 h(\sigma):=\sigma\int_1^\infty\mathsf d(x)\,x^{-\sigma-1}\,dx\qquad(\sigma>\mathsf e_1).
\end{equation}
For every integer $N\ge2$,
\begin{equation}\label{eq:ladder}
 h(\sigma)=O\bigl(\sigma N^{-\sigma}\bigr)\quad(\sigma\to+\infty)\iff\mathsf d(x)=0\ \text{for almost every }x\in[1,N).
\end{equation}
\end{theorem}

\begin{proof}
If $\mathsf d=0$ almost everywhere on $[1,N)$, then $|h(\sigma)|\le\sigma N^{-(\sigma-\mathsf e_1)}\int_N^\infty|\mathsf d|x^{-\mathsf e_1-1}dx$ for $\sigma>\mathsf e_1$.  Conversely, put $d(u):=\mathsf d(e^u)$, so that $d\,e^{-\mathsf e_1u}\in L^1(0,\infty)$, and
\begin{equation}\label{eq:ladder-Y}
 \mathsf Y(s):=\int_0^{\log N}d(u)\,e^{s(\log N-u)}\,du .
\end{equation}
$\mathsf Y$ is entire of exponential type at most $\log N$, bounded on $\Re s\le0$ by $\int_0^{\log N}|d|$, and $\mathsf Y(\sigma)\to0$ as $\sigma\to-\infty$.  For real $\sigma>\mathsf e_1$,
\begin{equation}\label{eq:ladder-Y-real}
 \mathsf Y(\sigma)=\frac{N^\sigma h(\sigma)}\sigma-\int_{\log N}^\infty d(u)\,e^{-\sigma(u-\log N)}\,du ,
\end{equation}
where the first term is bounded by hypothesis and the second by $N^{\mathsf e_1}\int|d|e^{-\mathsf e_1u}du$; so $\mathsf Y$ is bounded on the real axis.  The Phragm\'en--Lindel\"of theorem in the two quadrants $\{\Re s>0,\ \pm\Im s>0\}$ (angle $\pi/2$, order one, bounded on both sides) shows that $\mathsf Y$ is bounded in $\C$, hence constant, hence $0$.  By the uniqueness of the Laplace transform on $(0,\log N)$, $d=0$ almost everywhere there.
\end{proof}

No functional equation, no continuation of $h$ and no uniformity in $t$ enter Theorem~\CGref{thm:ladder}: the hypothesis lives on $\sigma\to+\infty$.

\begin{theorem}[Hamburger's theorem for a pair, half-line form]\label{thm:halfline-hamburger}
Let $f$ satisfy (H0) and the functional equation \CGeqref{eq:Z1-FE}, and suppose that for $\Re s>1$
\begin{equation}\label{eq:halfline-f}
 \log f(s)=s\int_1^\infty J_f(x)\,x^{-s-1}\,dx,\qquad\int_1^\infty|J_f(x)|\,x^{-\sigma-1}\,dx<\infty\quad(\sigma>1),
\end{equation}
with $J_f$ real and measurable.  Let $g(s)=\sum_{n\ge1}a_nn^{-s}$ converge absolutely in $\Re s>1$, with $a_1=1$, and let $J_g(x):=\sum_{2\le n\le x}b_n$, where $\log g(s)=\sum_{n\ge2}b_nn^{-s}$ for large $\sigma$.  Put $h(\sigma):=\log g(\sigma)-\log f(\sigma)$ for large real $\sigma$.  The following are equivalent:
\begin{enumerate}[label=\textup{(\roman*)}]
\item $h(\sigma)=O_N\bigl(\sigma N^{-\sigma}\bigr)$ as $\sigma\to+\infty$, for every integer $N\ge2$;
\item $f=g$;
\item $f=g=\zeta$.
\end{enumerate}
In this case $J_f$ is almost everywhere the step function $\sum_{p^m\le x}\frac1m$.
\end{theorem}

\begin{proof}
For large $\sigma$, $\log g(\sigma)=\sigma\int_1^\infty J_g(x)x^{-\sigma-1}dx$ by partial summation, so $h$ has the form \CGeqref{eq:ladder-h} with $\mathsf d=J_g-J_f$.  (i)$\Rightarrow$(ii): Theorem~\CGref{thm:ladder} gives $\mathsf d=0$ almost everywhere, so $\log f=\log g$ for large $\sigma$ and $f=g$ by analytic continuation.  (ii)$\Rightarrow$(iii): $g$ then satisfies (H0)--(H3), so $g=C_H\zeta$, and $a_1=1$ gives $C_H=1$.  (iii)$\Rightarrow$(i) because $h\equiv0$.
\end{proof}

\begin{corollary}[The ladder for the pair $(\zeta_\lambda,\tilde\zeta_{\Z,\lambda})$]\label{cor:ladder-pair}
Let $\lambda\in[0,1]$ and $N\ge2$.  Then $\tilde h_\lambda(\sigma)=O(\sigma N^{-\sigma})$ as $\sigma\to+\infty$ if and only if the Mellin count $J^{\rm M}_\lambda$ of $\zeta_\lambda$ coincides almost everywhere on $[1,N)$ with the step function $\sum_{n^m\le x}\tilde c_n(\lambda)/m$; in particular $J^{\rm M}_\lambda$ then vanishes on $(1,2)$ and is constant between consecutive integers below $N$.  Consequently
\begin{equation}\label{eq:ladder-D}
 (\mathrm D)\iff\tilde h_0(\sigma)=O_N\bigl(\sigma N^{-\sigma}\bigr)\ \text{for every }N\ge2 ,
\end{equation}
and at the clock end the first rung already fails under Hypothesis~\CGref{hyp:varpi-less-one}: $\tilde h_1(\sigma)\to\log L=-\log\varpi_{\RVM}$, which is $\ne0$ exactly when $\varpi_{\RVM}\ne1$ (Proposition~\CGref{prop:atomized-clock}(ii)).
\end{corollary}

\begin{proof}
Theorem~\CGref{thm:ladder} with $\mathsf d=J_g-J^{\rm M}_\lambda$, where $J_g$ is the count of $\log\tilde\zeta_{\Z,\lambda}=\sum_m\frac1m\sum_n\tilde c_n(\lambda)n^{-ms}$; the integrability is Theorem~\CGref{thm:mellin-counts}, whose proof applies to every $\zeta_\lambda$, and Corollary~\CGref{cor:abscissa-one}.  \CGeqref{eq:ladder-D} is Theorem~\CGref{thm:halfline-hamburger} with Corollary~\CGref{cor:abscissa-one} and Theorem~\CGref{thm:cell-form}, and the last statement is Proposition~\CGref{prop:atomized-clock}.
\end{proof}

The joining factor therefore measures the mass of $dJ^{\rm M}_\lambda$ outside the integers, and its rate of decay on the real axis is the first integer below which there is such mass.  The $N$th rung at $\lambda=0$ is the statement that $J^{\rm M}_{Z1}$ is constant on $(N-1,N)$; it is (D) restricted to one cell.

\subsection{The bridge theorem: the phase of the joining factor on the critical line}
\label{subsec:bridge}

\begin{theorem}[Bridge theorem, phase form]\label{thm:bridge}
Let $f$ be meromorphic in $\C$, real on the real axis, with (H0), the functional equation \CGeqref{eq:Z1-FE}, and $\lim_{\sigma\to+\infty}f(\sigma)=\varpi_f>0$.  Let $g(s)=\sum_{n\ge1}a_nn^{-s}$ converge absolutely in $\Re s>1$, with real $a_n$ and $a_1=1$.  Suppose that $g=f\,e^{h}$ in $\Re s>\frac12$ with $h$ holomorphic there and real on $(1,\infty)$, and that
\begin{itemize}
\item[\textup{(B1)}] $\Im h$ extends continuously to $\Re s\ge\frac12$ and vanishes on the line $\Re s=\frac12$;
\item[\textup{(B2)}] there are $\mathsf e_1,\mathsf e_2\ge0$ with
\begin{equation}\label{eq:bridge-B2}
 |h(s)|\le C\,(1+|s|)^{\mathsf e_1}\Bigl(1+\bigl(\sigma-\tfrac12\bigr)^{-\mathsf e_2}\Bigr)\qquad(\Re s>\tfrac12).
\end{equation}
\end{itemize}
Then $h\equiv-\log\varpi_f$, $g=\zeta$ and $f=\varpi_f\,\zeta$.
\end{theorem}

\begin{proof}
\emph{Reflection.}  $v:=\Im h$ is harmonic in $\Re s>\frac12$, continuous on the closed half-plane and $0$ on its boundary.  By the reflection principle \cite[Ch.~4, \S6.5]{Ahlfors1979} the function equal to $v(s)$ for $\Re s\ge\frac12$ and to $-v(1-\bar s)$ for $\Re s<\frac12$ is harmonic in $\C$.  It is the imaginary part of an entire function, which can be chosen equal to $h$ in $\Re s>\frac12$; in $\Re s<\frac12$ it differs from $\overline{h(1-\bar s)}$ by a real constant, and this constant is $0$ because the real part of an entire function is continuous across the line.  Thus $h$ is entire and $h(1-\bar s)=\overline{h(s)}$.  Since $h$ is real on $(1,\infty)$, $h(\bar s)=\overline{h(s)}$, and therefore $h(1-s)=h(s)$.

\emph{Growth.}  By \CGeqref{eq:bridge-B2} and the symmetry, $|h(s)|\le C(1+|s|)^{\mathsf e_1}\bigl(1+|\sigma-\frac12|^{-\mathsf e_2}\bigr)$ for $\Re s\ne\frac12$.  The function $\log|h|$ is subharmonic, so for every $s$
\begin{equation}\label{eq:bridge-mean}
 \log|h(s)|\le\frac1{2\pi}\int_0^{2\pi}\log\bigl|h\bigl(s+e^{i\phi}\bigr)\bigr|\,d\phi\le\log C+\mathsf e_1\log(2+|s|)+\log2+\frac{\mathsf e_2}{2\pi}\int_0^{2\pi}\log^+\frac1{|\sigma-\frac12+\cos\phi|}\,d\phi .
\end{equation}
For real $y$, $\frac1{2\pi}\int_0^{2\pi}\log|y+\cos\phi|\,d\phi\ge-\log2$, so the last integral is at most $2\pi\bigl(\log2+\log(2+|s|)\bigr)$.  Hence $|h(s)|\le C'(2+|s|)^{\mathsf e_1+\mathsf e_2}$ in $\C$, and $h$ is a polynomial.  On the real axis $h(\sigma)=\log g(\sigma)-\log f(\sigma)\to-\log\varpi_f$ as $\sigma\to+\infty$, so the polynomial is the constant $-\log\varpi_f$.

\emph{Hamburger.}  $g=f/\varpi_f$ satisfies (H0), (H1), and (H2)--(H3) with the second function equal to $g$; so $g=C_H\zeta$, and $a_1=1$ gives $C_H=1$.
\end{proof}

The symmetry of the joining factor, its continuation to $\Re s\le\frac12$ and its finite order are not hypotheses of Theorem~\CGref{thm:bridge}: they follow from \textup{(B1)} and \textup{(B2)}, which are written on the closed half-plane $\Re s\ge\frac12$ where the joining factor of the family is defined.  Theorem~\CGref{thm:pair-criterion} assumes the continuation and the symmetry instead, and nothing like \textup{(B2)}; Theorem~\CGref{thm:bridge} derives both from \textup{(B1)}--\textup{(B2)}. Neither set of hypotheses contains the other.

\begin{corollary}[The bridge theorem for the family]\label{cor:bridge-family}
\begin{enumerate}[label=\textup{(\roman*)}]
\item For $0<\lambda\le1$ the joining factor $\tilde h_\lambda$ does not satisfy both \textup{(B1)} and \textup{(B2)}.
\item (D) holds if and only if $\tilde h_0$ satisfies \textup{(B1)} and \textup{(B2)}.
\end{enumerate}
\end{corollary}

\begin{proof}
Theorem~\CGref{thm:bridge} applies with $f=\zeta_\lambda$ (Theorem~\CGref{thm:path}, Proposition~\CGref{prop:family-clock}), $g=\tilde\zeta_{\Z,\lambda}$ (Corollary~\CGref{cor:abscissa-one}) and $h=\tilde h_\lambda$ (Theorem~\CGref{thm:atomized-continuation}(iii)).  If \textup{(B1)} and \textup{(B2)} hold, then $\zeta_\lambda=\varpi_\lambda\zeta$ is an absolutely convergent Dirichlet series in $\Re s>1$, which is excluded for $0<\lambda\le1$ by Proposition~\CGref{prop:nonconstant}; for $\lambda=0$ it is (D).  Conversely (D) gives $\tilde h_0\equiv0$ (Theorem~\CGref{thm:closing}).
\end{proof}

Define the \emph{phase defect} of the family, when the limit exists,
\begin{equation}\label{eq:phase-defect}
 \mathsf v_\lambda(t):=\lim_{\sigma\to\frac12^+}\Im\tilde h_\lambda(\sigma+it)=\lim_{\sigma\to\frac12^+}\Bigl(\Im\log\tilde\zeta_{\Z,\lambda}(\sigma+it)-\Im\log\zeta_\lambda(\sigma+it)\Bigr).
\end{equation}

\begin{proposition}[The phase of the continuous side]\label{prop:continuous-phase}
Let $\lambda\in[0,1]$ and $t>0$ not an ordinate $\gamma_k(\lambda)$.  With the branch of $\log\zeta_\lambda$ that is real on $(1,\infty)$, continued along any path in the open quadrant $\{\Re s>\frac12,\ \Im s>0\}$ (the continuation is well defined there, since the zeros of $\zeta_\lambda$ lie on its boundary line $\Re s=\frac12$ and the pole $s=1$ on its boundary too; the limit below is taken at points $\frac12+it$ that are not zeros),
\begin{equation}\label{eq:continuous-phase}
 \lim_{\sigma\to\frac12^+}\Im\log\zeta_\lambda(\sigma+it)=-\vartheta(t)+\pi\bigl(N_\lambda(t)-1\bigr),\qquad N_\lambda(t)=\#\{k:\gamma_k(\lambda)\le t\},
\end{equation}
and in the sense of distributions in $t>0$
\begin{equation}\label{eq:generator-phase}
 \lim_{\sigma\to\frac12^+}\Im\mathcal V_\lambda(\sigma+it)=-\pi\sum_{k\ge1}\eta_k\,\delta\bigl(t-\gamma_k(\lambda)\bigr).
\end{equation}
Consequently $\mathsf v_\lambda(t)=\lim_{\sigma\to\frac12^+}\Im\log\tilde\zeta_{\Z,\lambda}(\sigma+it)+\vartheta(t)-\pi\bigl(N_\lambda(t)-1\bigr)$.
\end{proposition}

\begin{proof}
$\zeta_\lambda=P_\lambda/G$.  On $s=\frac12+it$, $q(s)=-(\frac14+t^2)$ and $1+q(s)/A_{k,\lambda}=(\gamma_k(\lambda)^2-t^2)/A_{k,\lambda}$ are real, so $P_\lambda(\frac12+it)$ is real and changes sign at each $\gamma_k(\lambda)$; approaching from $\sigma>\frac12$ the argument of the factor $k$ increases by $\pi$ as $t$ passes $\gamma_k(\lambda)$, since $\Im\bigl(\gamma_k(\lambda)^2+(\sigma-\frac12+it)^2\bigr)=2t(\sigma-\frac12)>0$.  The argument of $G(\frac12+it)=-(\frac14+t^2)\Gamma(\frac14+\frac{it}2)\pi^{-1/4-it/2}$ is $\pi+\vartheta(t)$ by \CGeqref{eq:theta-def}, with the branch fixed by $\zeta_\lambda(\sigma)<0$ for $\frac12<\sigma<1$.  This gives \CGeqref{eq:continuous-phase}.  For \CGeqref{eq:generator-phase} it is preferable not to interchange a boundary limit and $\partial_\lambda$.  Put $z=\sigma'+it$ and $\gamma=\gamma_k(\lambda)$.  The $k$th term of \CGeqref{eq:generator-zeta} is
\[
 2\gamma\eta_k\left(\frac1{\gamma^2+z^2}-\frac1{\frac14+\gamma^2}\right).
\]
The second summand is real, while
\[
 \frac{2\gamma}{\gamma^2+z^2}=-i\left(\frac1{z-i\gamma}-\frac1{z+i\gamma}\right).
\]
The Poisson-kernel boundary limit on $t>0$ is therefore $-\pi\eta_k\delta(t-\gamma)$; the pole at $-\gamma$ lies outside the positive half-line.  Summing against a compactly supported test function is justified by the bounded displacement and the same mean-value/summation-by-parts estimates used for the generator (equivalently by Cauchy--Schwarz with $\sum_k\eta_k^2\gamma_k^{-2}<\infty$ after one integration by parts).  This proves \CGeqref{eq:generator-phase} directly in distributions.
\end{proof}

By Proposition~\CGref{prop:continuous-phase}, $\mathsf v_\lambda\equiv0$ asks two things of the discrete function on the critical line: jumps of $\pi$ of its argument exactly at the ordinates $\gamma_k(\lambda)$, and the phase $-\vartheta(t)$ between them; equivalently, $e^{i\vartheta(t)}\tilde\zeta_{\Z,\lambda}(\frac12+it)$ is real.  At the clock end $\tilde h_1$ is holomorphic in $\Re s>0$ (Theorem~\CGref{thm:atomized-continuation}(iv)), so $\mathsf v_1(t)=\Im\tilde h_1(\frac12+it)$ is a real-analytic function; it does not vanish identically, by Corollary~\CGref{cor:bridge-family}(i) and the bound $|\tilde h_1(s)|\le C(1+|s|)$ on $\Re s\ge\frac14$ of the proof of Proposition~\CGref{prop:clock-asymmetry}.

\begin{proposition}[The cell masses are sums of a one-ordinate function]\label{prop:one-ordinate}
For $n\ge2$ there is a real function $\mathsf o_n(\kappa,\kappa')$ of two real variables, depending only on $n$ and the cell $\varphi_n$, such that for all lists $\boldsymbol{\kappa},\boldsymbol{\kappa}'\in\mathscr O$
\begin{equation}\label{eq:one-ordinate}
 \tilde c_n(\boldsymbol{\kappa})-\tilde c_n(\boldsymbol{\kappa}')=\sum_{k\ge1}\mathsf o_n(\kappa_k,\kappa'_k),\qquad \mathsf o_n(\kappa,\kappa')=-2\,\Re\bigl(W_n(\kappa)-W_n(\kappa')\bigr)\ \text{ if }n\text{ is not a power of }2,
\end{equation}
and, with $\mathsf L_n(s):=-\log(1-n^{-s})$, for $\Re s>1$
\begin{equation}\label{eq:one-ordinate-log}
 \log\tilde\zeta_{\Z,\boldsymbol{\kappa}}(s)-\log\tilde\zeta_{\Z,\boldsymbol{\kappa}'}(s)=\sum_{n\ge2}\Bigl(\sum_{k\ge1}\mathsf o_n(\kappa_k,\kappa'_k)\Bigr)\mathsf L_n(s),
\end{equation}
\begin{equation}\label{eq:one-ordinate-logzeta}
 \log\zeta_{\boldsymbol{\kappa}}(s)-\log\zeta_{\boldsymbol{\kappa}'}(s)=\sum_{k\ge1}\log\frac{1+q(s)/(\frac14+\kappa_k^2)}{1+q(s)/(\frac14+{\kappa'_k}^{2})} .
\end{equation}
\end{proposition}

\begin{proof}
$\tilde c_n(\boldsymbol{\kappa})=\langle d\pi_{\boldsymbol{\kappa}},\varphi_n\rangle$, where $d\pi_{\boldsymbol{\kappa}}$ is obtained from $J_{\boldsymbol{\kappa}}$ by truncation at $2$ and finite M\"obius inversion, both linear, and by \CGeqref{eq:J-X} the difference $J_{\boldsymbol{\kappa}}-J_{\boldsymbol{\kappa}'}=-2\Re\sum_k\bigl(\Ei((\frac12+i\kappa_k)\log x)-\Ei((\frac12+i\kappa'_k)\log x)\bigr)$ is a sum over $k$ converging in $\mathcal D'(1,\infty)$.  This gives \CGeqref{eq:one-ordinate}, with the value \CGeqref{eq:cell-mass-W} off the powers of $2$.  \CGeqref{eq:one-ordinate-log} is the Euler product with Corollary~\CGref{cor:abscissa-one}, and \CGeqref{eq:one-ordinate-logzeta} is \CGeqref{eq:list-objects}.
\end{proof}

By Proposition~\CGref{prop:one-ordinate}, the difference of two joining factors is, term by term in $n$ and in $k$, a sum over the ordinates of one universal function of one ordinate, the joining factor of a single zero factor:
\begin{equation}\label{eq:one-ordinate-joining}
 \mathfrak j(\kappa,\kappa';s):=\sum_{n\ge2}\mathsf o_n(\kappa,\kappa')\,\mathsf L_n(s)-\log\frac{1+q(s)/(\frac14+\kappa^2)}{1+q(s)/(\frac14+{\kappa'}^{2})},\qquad
 \tilde h_{\boldsymbol{\kappa}}-\tilde h_{\boldsymbol{\kappa}'}=\sum_{k\ge1}\mathfrak j(\kappa_k,\kappa'_k;s)\ \ \text{formally.}
\end{equation}
With $\boldsymbol{\kappa}'=(\tau_k)$ this writes the phase defect as $\mathsf v_{\boldsymbol{\kappa}}=\mathsf v_1+\sum_k\Im\mathfrak j(\kappa_k,\tau_k;\frac12+it)$, the defect of the clock plus a transport term.

\begin{heuristic}[What happens to the phase defect along the path]\label{heur:phase-rotation}
The following is a computation by stationary phase without remainder terms; it is not proved.  For $n$ much larger than $\kappa$ the cell resolves the oscillation, $\Re W_n(\kappa)\approx n^{-1/2}\cos(\kappa\log n)/\log n$, and for $n$ much smaller than $\kappa/2\pi$ the kernel $W_n(\kappa)$ is negligible because $\varphi$ is smooth (Proposition~\CGref{prop:height-kernels}).  The transition is at $n\approx\kappa/2\pi$, where the phase is $\kappa\log\frac\kappa{2\pi}-\kappa=2\pi\bigl(\NR(\kappa)-\frac78\bigr)$.  This suggests
\begin{equation}\label{eq:phase-rotation}
 \Im\mathfrak j\bigl(\kappa,\cdot\,;\tfrac12+it\bigr)\approx\text{amplitude}\times\cos\bigl(2\pi\NR(\kappa)+\text{phase}(t)\bigr),\qquad\bigl|\mathfrak j(\kappa,\cdot\,;\sigma+it)\bigr|\approx\kappa^{-1/2-\sigma},
\end{equation}
so that the sum over the ordinates converges absolutely exactly for $\sigma>\frac12$, in agreement with Theorem~\CGref{thm:atomized-continuation}(iii), and is at the edge of convergence on the critical line.  On the clock, $\NR(\tau_k)=k-\frac12$ and the phases are aligned, which is why $\tilde h_1$ crosses the line.  Along the path each term turns by the angle $(1-\lambda)\eta_k\ell_k$, the rotation of \CGeqref{eq:Xi-rotation}.  At $\lambda=0$, \textup{(B1)} is the statement that the rotated sum, with the phases $e^{2\pi i\NR(\gamma_k)}$, converges on the critical line to minus the defect $\mathsf v_1$ of the clock.
\end{heuristic}

\subsection{Transport from the clock on the rungs}
\label{subsec:transport}

For a real function $\mathsf f\in C_c^\infty(0,\infty)$ put
\begin{equation}\label{eq:transport-defs}
 \mathsf c_{\mathsf f}(t):=\int_0^\infty\mathsf f(u)\cos(tu)\,du,\qquad \mathsf k(u):=\frac12\Bigl(e^{u/2}-\frac{e^{-u/2}}{e^{2u}-1}\Bigr),\qquad \mathsf k_{\mathsf f}:=\int_0^\infty\mathsf f(u)\,\mathsf k(u)\,du .
\end{equation}

\begin{proposition}[The rung functional]\label{prop:rung-functional}
For every list $\boldsymbol{\kappa}\in\mathscr O$ (Section~\CGref{sec:lists}),
\begin{equation}\label{eq:rung-functional}
 \mathsf R_{\mathsf f}(\boldsymbol{\kappa}):=\frac12\int_1^\infty\mathsf f(\log x)\,x^{-1/2}\log x\,dJ^{\rm M}_{\boldsymbol{\kappa}}(x)=\mathsf k_{\mathsf f}-\sum_{k\ge1}\mathsf c_{\mathsf f}(\kappa_k),
\end{equation}
the series converging absolutely.  Along the family, with $\mathsf R_{\mathsf f}(\lambda):=\mathsf R_{\mathsf f}\bigl((\gamma_k(\lambda))\bigr)$,
\begin{equation}\label{eq:transport}
 \mathsf R_{\mathsf f}(\lambda)=\mathsf R_{\mathsf f}(1)-\mathsf S_{\mathsf f}(\lambda),\qquad \mathsf S_{\mathsf f}(\lambda):=\sum_{k\ge1}\Bigl(\mathsf c_{\mathsf f}\bigl(\gamma_k(\lambda)\bigr)-\mathsf c_{\mathsf f}(\tau_k)\Bigr),\qquad\mathsf S_{\mathsf f}(1)=0 .
\end{equation}
If the support of $\mathsf f$ contains no point $\log n$, $n\in\N$, then
\begin{equation}\label{eq:transport-D}
 (\mathrm D)\ \Longrightarrow\ \mathsf S_{\mathsf f}(0)=\mathsf R_{\mathsf f}(1),
\end{equation}
and (D) holds if and only if $\mathsf S_{\mathsf f}(0)=\mathsf R_{\mathsf f}(1)$ for every such $\mathsf f$ and the resulting integer jumps satisfy the growth hypothesis of Theorem~\CGref{thm:cell-form}.
\end{proposition}

\begin{proof}
Pair \CGeqref{eq:distributional-EF}, which holds for every list of $\mathscr O$ with the same proof, with $\frac12\mathsf f(\log x)x^{-1/2}$ and put $u=\log x$: the term $1$ gives $\frac12\int\mathsf f(u)e^{u/2}du$, the term $x^{-1}(x^2-1)^{-1}$ gives $\frac12\int\mathsf f(u)e^{-u/2}(e^{2u}-1)^{-1}du$, and the sum gives $\sum_k\mathsf c_{\mathsf f}(\kappa_k)$; $\mathsf c_{\mathsf f}$ is rapidly decreasing.  \CGeqref{eq:transport} is the difference of two cases of \CGeqref{eq:rung-functional}.  By Theorem~\CGref{thm:cell-form}, the vanishing $\mathsf R_{\mathsf f}(0)=0$ for every $\mathsf f$ whose support avoids the points $\log n$ is exactly the constancy part of the cell criterion; (D) follows when the resulting integer jumps also satisfy $\sum_{n\le X}|\mathsf b_n|=O_\eps(X^{1+\eps})$.
\end{proof}

The number $\mathsf R_{\mathsf f}(1)$ belongs to the clock; $\mathsf S_{\mathsf f}(0)$ is the transport from the clock to the ordinates of $\zeta$.  (D) is the statement that the transport cancels the defect of the clock on every rung.

\begin{theorem}[The transport is band-limited in $\lambda$]\label{thm:transport-band}
Let $\mathsf f\in C_c^\infty(0,\log N)$.  The function $\lambda\mapsto\mathsf S_{\mathsf f}(\lambda)$ extends to an entire function with
\begin{equation}\label{eq:transport-band}
 |\mathsf S_{\mathsf f}(\lambda)|\le C_{\mathsf f}\bigl(1+|\lambda|\log(2+|\lambda|)\bigr)\,e^{C_\eta\log N\,|\Im\lambda|}\qquad(\lambda\in\C),
\end{equation}
and for every $j\ge1$
\begin{equation}\label{eq:transport-moments}
 \partial_\lambda^j\mathsf S_{\mathsf f}(1)=\mathsf m_j(\mathsf f):=\sum_{k\ge1}\eta_k^{\,j}\,\mathsf c_{\mathsf f}^{(j)}(\tau_k),\qquad
 \mathsf S_{\mathsf f}(0)=\sum_{j\ge1}\frac{(-1)^j}{j!}\,\mathsf m_j(\mathsf f).
\end{equation}
Moreover $|\mathsf m_j(\mathsf f)|\le C_{\mathsf f}(1+j)^2(C_\eta\log N)^j$ (the support of $\mathsf f$ lies in some $[u_1,\log N]$ with $u_1>0$, and $C_{\mathsf f}$ depends on it).
\end{theorem}

\begin{proof}
By the Paley--Wiener theorem, $|\mathsf c_{\mathsf f}^{(j)}(t)|\le C_{A,j}\,e^{\log N|\Im t|}(1+|t|)^{-A}$ for complex $t$ and every $A$.  With $\gamma_k(\lambda)=\gamma_k+\lambda\eta_k$, $|\Im\gamma_k(\lambda)|\le C_\eta|\Im\lambda|$ and $|\gamma_k(\lambda)|\ge\gamma_k-C_\eta|\lambda|$.  The ordinates with $\gamma_k\le2C_\eta|\lambda|$ are $O(|\lambda|\log(2+|\lambda|))$ in number and each contributes at most $C\,e^{C_\eta\log N|\Im\lambda|}$; for the others $|\mathsf c_{\mathsf f}(\gamma_k(\lambda))|\le C\,e^{C_\eta\log N|\Im\lambda|}(\gamma_k/2)^{-2}$.  This gives \CGeqref{eq:transport-band}, the termwise differentiation and the convergence of the Taylor series at $\lambda=1$ evaluated at $\lambda=0$.  Finally $\mathsf c^{(j)}_{\mathsf f}(t)=\int\mathsf f(u)u^j\cos(tu+\frac{j\pi}2)du$ and two integrations by parts give $(1+t^2)|\mathsf c_{\mathsf f}^{(j)}(t)|\le C_{\mathsf f}(1+j)^2(\log N)^j$.
\end{proof}

The scaled lists $(\tau_k-c'\eta_k)$ are the members $\lambda=1-c'$ of the family, so the scale test of Remark~\CGref{rem:scale-test} is the behaviour of the entire function $\mathsf S_{\mathsf f}$ on the real axis.  In \CGeqref{eq:transport-moments} the nodes $\tau_k$ and the weights $\mathsf c^{(j)}_{\mathsf f}(\tau_k)$ belong to the clock, and the powers $\eta_k^{\,j}$ do not; on the first rung, $N=2$, the term $j$ is at most $C_{\mathsf f}(1+j)^2(C_\eta\log2)^j/j!$.

\begin{proposition}[The defect of the clock on a rung]\label{prop:clock-defect-poisson}
Let $\mathfrak t_0\in(2\pi,\tau_0)$ be the point with $\NR(\mathfrak t_0)=0$, and let $\mathsf f\in C_c^\infty(0,\infty)$.  Then
\begin{equation}\label{eq:clock-defect-poisson}
 \sum_{k\ge1}\mathsf c_{\mathsf f}(\tau_k)=\mathsf I_0+2\sum_{m\ge1}(-1)^m\mathsf I_m,\qquad \mathsf I_m:=\int_{\mathfrak t_0}^\infty\mathsf c_{\mathsf f}(t)\cos\bigl(2\pi m\NR(t)\bigr)\NR'(t)\,dt=O_{\mathsf f}\bigl(m^{-2}\bigr),
\end{equation}
\begin{equation}\label{eq:clock-defect-I0}
 \mathsf I_0=-\int_0^\infty\frac{\mathsf f(u)}{2\pi u}\Bigl(\sin(u\mathfrak t_0)\log\frac{\mathfrak t_0}{2\pi}+\frac\pi2-\int_0^{u\mathfrak t_0}\frac{\sin y}y\,dy\Bigr)du,
\end{equation}
and therefore $\mathsf R_{\mathsf f}(1)=(\mathsf k_{\mathsf f}-\mathsf I_0)-2\sum_{m\ge1}(-1)^m\mathsf I_m$.
\end{proposition}

\begin{proof}
The function $g(v):=\mathsf c_{\mathsf f}(\mathcal T(v))$, $v\ge0$, with $\mathcal T$ the inverse of the clock on $[\mathfrak t_0,\infty)$, is continuous on $[0,\infty)$, $C^\infty$ on $(0,\infty)$ with $g'(0^+)=\mathsf c_{\mathsf f}'(\mathfrak t_0)\mathcal T'(0)$ finite, and rapidly decreasing with its derivatives, and $\tau_k=\mathcal T(k-\frac12)$.  Poisson's summation formula for the even extension of $g$ (continuous, in $L^1$, with a corner at $0$ in general, since $g'(0^+)\ne0$ unless $\mathsf c_{\mathsf f}'(\mathfrak t_0)=0$) at the half-integers gives $2\sum_{k\ge1}g(k-\frac12)=\sum_{m\in\Z}(-1)^m2\int_0^\infty g(v)\cos(2\pi mv)\,dv$, the Fourier coefficients being $O(m^{-2})$ by two integrations by parts, $\int_0^\infty g(v)\cos(2\pi mv)\,dv=-\frac{g'(0^+)}{(2\pi m)^2}-\frac1{(2\pi m)^2}\int_0^\infty g''(v)\cos(2\pi mv)\,dv$, the boundary term coming from the corner (so the decay cannot be improved to $O(m^{-3})$ in general); the substitution $v=\NR(t)$ gives \CGeqref{eq:clock-defect-poisson}.  For $m=0$, $\NR'(t)=\frac1{2\pi}\log\frac t{2\pi}$; integrating by parts with the primitive $\int\mathsf f(u)\sin(tu)u^{-1}du$ of $\mathsf c_{\mathsf f}$ and exchanging the integrals gives \CGeqref{eq:clock-defect-I0}.
\end{proof}

In $\mathsf I_m$ the phases $2\pi m\NR(t)\mp tu$ have the derivative $m\log\frac t{2\pi}\mp u$, which vanishes at $t=2\pi e^{u/m}$: the defect of the clock at the point $x=e^u$ is produced at the heights $2\pi x^{1/m}$, which is the origin of the terms $2\cos(2\pi y)$ and $y^{-1/4}\cos(4\pi\sqrt y+\frac\pi4)$ of \CGeqref{eq:clock-density}.  On the rung $(n,n+1)$ the height for $m=1$ lies in $(2\pi n,2\pi(n+1))$.  On the first rung it lies in $(2\pi,4\pi)$, below $\tau_1$ and on both sides of $\mathfrak t_0$: the defect of the clock on $(1,2)$ is an effect of the start of the clock.  Experiments~\CGref{exp:transport} and~\CGref{exp:clock-defect} evaluate \CGeqref{eq:transport}--\CGeqref{eq:clock-defect-poisson} for one function $\mathsf f$.


\subsection{The multiscale Jacobian of the midpoint clock}
\label{subsec:multiscale-jacobian}

The rung system of Subsection~\CGref{subsec:free-masses} asks for a list and free masses simultaneously.  Before attempting nonlinear continuation from the clock, one must understand the linearization with respect to the ordinates.  The point is genuinely multiscale: the gaps of the midpoint clock tend to zero, so no fixed $u$-window can carry a uniform inverse, while a window of length comparable with $\log T$ is exactly large enough to resolve the ordinates up to height $T$.

For $T>\tau_2$ put
\[
 K(T):=\max\{k:\tau_k\le T\},\qquad
 \Delta_T:=\min_{1\le k<K(T)}(\tau_{k+1}-\tau_k),
\]
and define the finite ordinate block of the linearized rung map by
\begin{equation}\label{eq:clock-jacobian}
 (\mathcal J_T\mathbf v)(u):=-u\sum_{\tau_k\le T}v_k\sin(\tau_k u).
\end{equation}
No genuine zero ordinate and no individual displacement $\eta_k=\tau_k-\gamma_k$ enters this definition.

\begin{theorem}[Multiscale Jacobian of the midpoint clock]\label{thm:clock-jacobian}
The following statements hold.
\begin{enumerate}[label=\textup{(\roman*)}]
\item For every $T>\tau_2$,
\begin{equation}\label{eq:clock-gap-lower}
 \Delta_T\ge \frac{2\pi}{\log(T/2\pi)}.
\end{equation}
\item Fix $a>0$ and $\varkappa_{\rm win}>0$, and put
\[
 L_T:=(1+\varkappa_{\rm win})\log(T/2\pi),\qquad I_T=[a,a+L_T].
\]
There is $c_{a,\varkappa_{\rm win}}>0$, independent of $T$, such that for all finitely supported $\mathbf v=(v_k)$ on $\{k:\tau_k\le T\}$,
\begin{equation}\label{eq:multiscale-frame}
 \|\mathcal J_T\mathbf v\|_{L^2(I_T)}^2
 \ge c_{a,\varkappa_{\rm win}}L_T\sum_{\tau_k\le T}|v_k|^2.
\end{equation}
\item No analogue of \CGeqref{eq:multiscale-frame} with a positive $T$-independent lower constant can hold on a fixed interval $(0,U)$.  Indeed, for the unit vector $\mathbf v^{(k)}=(e_k-e_{k+1})/\sqrt2$,
\begin{equation}\label{eq:fixed-window-pair}
 \|\mathcal J\mathbf v^{(k)}\|_{L^2(0,U)}
 \le \frac{U^{5/2}}{\sqrt{10}}(\tau_{k+1}-\tau_k)\longrightarrow0.
\end{equation}
\item For every finite $U>0$,
\[
 \overline{\operatorname{span}}\{u\sin(\tau_k u):k\ge1\}=L^2(0,U).
\]
Thus the infinite linearized clock has dense range on every finite window, despite the absence of a bounded inverse on a fixed window.
\item Let $\mathsf f\in C_c^\infty(0,U)$ and $\mathsf c_{\mathsf f}(t)=\int_0^\infty\mathsf f(u)\cos(tu)\,du$.  For every list $\boldsymbol\kappa$ satisfying $|\kappa_k-\tau_k|\le C_\eta$ and every $A>1$,
\begin{equation}\label{eq:smooth-tail-clock}
 \sum_{\tau_k>T}\left|\mathsf c_{\mathsf f}(\kappa_k)-\mathsf c_{\mathsf f}(\tau_k)\right|
 \ll_{A,\mathsf f,C_\eta} T^{1-A}\log T.
\end{equation}
\end{enumerate}
\end{theorem}

\begin{proof}
The midpoint equation is $\NR(\tau_k)=k-\frac12$ with
\[
 \NR'(t)=\frac1{2\pi}\log\frac t{2\pi}.
\]
For $\tau_{k+1}\le T$, the mean-value theorem gives
\[
 1=\NR'(\xi_k)(\tau_{k+1}-\tau_k),\qquad \tau_k<\xi_k<\tau_{k+1}\le T,
\]
which proves \CGeqref{eq:clock-gap-lower}.

For (ii), apply Ingham's inequality for non-harmonic exponentials \cite{Ingham1936} to the separated frequency set $\{\pm\tau_k:\tau_k\le T\}$.  By \CGeqref{eq:clock-gap-lower},
\[
 L_T\Delta_T\ge 2\pi(1+\varkappa_{\rm win}),
\]
so after scaling the frequencies by $\Delta_T$ the interval lies a fixed positive distance beyond Ingham's critical length.  The lower Ingham constant is therefore bounded below by a positive constant depending only on $\varkappa_{\rm win}$; rescaling gives a bound of size $L_T$ for the sine sum.  Since $u\ge a$ on $I_T$, multiplication by $u$ gives \CGeqref{eq:multiscale-frame}.

For (iii), $|\sin(\tau_{k+1}u)-\sin(\tau_k u)|\le(\tau_{k+1}-\tau_k)u$, whence
\[
 \|\mathcal J\mathbf v^{(k)}\|_{L^2(0,U)}^2
 \le\frac{(\tau_{k+1}-\tau_k)^2}{2}\int_0^Uu^4\,du
 =\frac{(\tau_{k+1}-\tau_k)^2U^5}{10}.
\]
The same mean-value formula used above shows $\tau_{k+1}-\tau_k\to0$.

For (iv), suppose $F\in L^2(0,U)$ is orthogonal to every $u\sin(\tau_k u)$.  Put $H(u)=uF(u)$ and extend $H$ oddly to $(-U,U)$.  Its Fourier transform
\[
 \widehat H(z)=\int_{-U}^{U}H(u)e^{-izu}\,du
\]
is entire of exponential type at most $U$ and vanishes at every $\pm\tau_k$.  A nonzero entire function of finite exponential type has $O(r)$ zeros in $|z|\le r$ by Jensen's formula, whereas the midpoint clock has $2\NR(r)+O(1)\asymp r\log r$ points $\pm\tau_k$ in that disk.  Hence $\widehat H\equiv0$, so $H=0$ and therefore $F=0$ almost everywhere.

For (v), the Paley--Wiener decay of the cosine transform gives $|\mathsf c'_{\mathsf f}(t)|\ll_{A,\mathsf f}(1+|t|)^{-A}$ for every $A$.  The mean-value theorem and $|\kappa_k-\tau_k|\le C_\eta$ give
\[
 |\mathsf c_{\mathsf f}(\kappa_k)-\mathsf c_{\mathsf f}(\tau_k)|\ll_{A,\mathsf f,C_\eta}\tau_k^{-A}
\]
for large $k$.  Partial summation with $\NR(t)=O(t\log t)$ yields $\sum_{\tau_k>T}\tau_k^{-A}=O(T^{1-A}\log T)$.
\end{proof}

The theorem identifies the correct scale but does not prove existence of a solution of the rung system.  In particular, a fixed-window Banach-space implicit-function argument cannot work because of \CGeqref{eq:fixed-window-pair}.  A quantitative continuation must grow the resolved window together with the spectral height.

One precise residual formulation is the following.  Fix $r>\frac12$, so point masses belong to $H^{-r}(0,U)$, and put
\[
 \mathsf D_{\RVM}(u):=\mathsf k(u)-\sum_{k\ge1}\cos(\tau_k u)
\]
in $\mathcal D'(0,\infty)$.  For $T,U>0$ define
\begin{equation}\label{eq:multiscale-residual}
 \varepsilon_r(T,U):=
 \inf_{h,b}\left\|
 \mathsf D_{\RVM}+u\sum_{\tau_k\le T}h_k\sin(\tau_k u)
 -\sum_{\log n\le U}b_n\delta(u-\log n)
 \right\|_{H^{-r}(0,U)}.
\end{equation}
Theorem~\CGref{thm:clock-jacobian} supplies the finite-height separation, dense-range and tail ingredients.  What is not proved is a bound for \CGeqref{eq:multiscale-residual}, together with control of the minimizing $h$ strong enough to dominate the quadratic remainder in
\[
 \cos((\tau_k+h_k)u)=\cos(\tau_ku)-h_ku\sin(\tau_ku)+O(h_k^2u^2),
\]
uniformly along a nested sequence with $U\asymp\log T$.  Such an estimate would turn the finite-section continuation observed numerically into an existence proof for the infinite rung system.  By Theorem~\CGref{thm:free-masses}, any resulting solution would then be the zeta ordinates with the von Mangoldt masses.  The individual displacements $\eta_k$ are not data in this formulation.

\section{The space of ordinate lists and the existence form of the final step}
\label{sec:lists}

The constructions of Sections~\CGref{sec:lambda-family}--\CGref{sec:criteria} use two properties of the ordinates (Remark~\CGref{rem:input-A7}).  This section applies them to every list with these two properties.  The final step then becomes an existence statement in a compact space, and the list of $\zeta$ is not part of its formulation.

\subsection{Lists, their zeta functions and their atomizations}
\label{subsec:lists}

Let $\mathscr O$ be the set of sequences $\boldsymbol{\kappa}=(\kappa_k)_{k\ge1}$ of real numbers with
\begin{equation}\label{eq:list-space}
 \tau_0^{+}\le\kappa_1\le\kappa_2\le\cdots,\qquad|\kappa_k-\tau_k|\le C_\eta,\qquad\Bigl|\int_{\tau_0}^{X}\bigl(N_{\boldsymbol{\kappa}}(\tau)-\NR(\tau)\bigr)d\tau\Bigr|\le C\log X\quad(X\ge\tau_0),
\end{equation}
where $\tau_0^{+}$ is a fixed number with $\tau_0<\tau_0^{+}<\min(\gamma_1,\tau_1)$, $N_{\boldsymbol{\kappa}}(\tau)=\#\{k:\kappa_k\le\tau\}$, $C_\eta$ is the constant of Lemma~\CGref{lem:bounded-displacement} and $C$ the constant of Proposition~\CGref{prop:family-clock}.  The clock $(\tau_k)$, the ordinates $(\gamma_k)$ of $\zeta$ and every list $(\gamma_k(\lambda))$, $0\le\lambda\le1$, belong to $\mathscr O$.  For $\boldsymbol{\kappa}\in\mathscr O$ (Section~\CGref{sec:lists}) define, exactly as for the family,
\begin{equation}\label{eq:list-objects}
 \zeta_{\boldsymbol{\kappa}}(s):=\frac1{G(s)}\prod_{k\ge1}\Bigl(1+\frac{q(s)}{\frac14+\kappa_k^2}\Bigr),\qquad
 \tilde c_n(\boldsymbol{\kappa}),\quad \tilde a_N(\boldsymbol{\kappa}),\quad k_N(\boldsymbol{\kappa}),\qquad \tilde\zeta_{\Z,\boldsymbol{\kappa}}=\zeta_{\boldsymbol{\kappa}}\,e^{\tilde h_{\boldsymbol{\kappa}}},\qquad \mathsf V_\sigma(\boldsymbol{\kappa}),
\end{equation}
with the cell masses of Definition~\CGref{def:atomized}, the joining factor of Theorem~\CGref{thm:joining-factor} and the function \CGeqref{eq:vertical-defect}.  The results of the preceding sections that are used below for lists are transferred to $\mathscr O$ by the following theorem, which records, result by result, which part of \CGeqref{eq:list-space} each proof uses and which constants are uniform.

\begin{theorem}[Transfer to $\mathscr O$]\label{thm:transfer-O}
Let $\boldsymbol{\kappa}\in\mathscr O$. The following statements, proved above for the lists $(\gamma_k(\lambda))_{k\ge1}$, hold for $\boldsymbol{\kappa}$ with the same proofs, and their constants depend only on $\tau_0^+$, $C_\eta$ and $C$ in \CGeqref{eq:list-space} and on the parameters of the statement ($\sigma$, $\eps$, $n$, $N$), never on the list: $C=C(\tau_0^+,C_\eta,C;\sigma,\eps,\dots)$, hence they are uniform on $\mathscr O$. This uniformity is what Proposition~\CGref{prop:lists-compact} and the compactness arguments of Section~\CGref{sec:routeB-finite} use.
\begin{enumerate}[label=\textup{(\roman*)}]
\item Theorem~\CGref{thm:path}: $\zeta_{\boldsymbol{\kappa}}$ has order one, satisfies the functional equation \CGeqref{eq:Z1-FE}, and its nontrivial zeros are exactly $\frac12\pm i\kappa_k$, with multiplicity. Input: the tube (convergence of the product and $\kappa_k\asymp\tau_k$), the ordering, and $\kappa_1\ge\tau_0^+$.
\item Theorem~\CGref{thm:Z1-clock} in its list form, with $\boldsymbol{\kappa}$ in place of $\boldsymbol\gamma$: the clock decomposition of $\zeta_{\boldsymbol{\kappa}}$, the value $\varpi_{\boldsymbol{\kappa}}$ at $+\infty$ as an ordered product, and the bounds of Proposition~\CGref{prop:family-clock}. Input: the tube and the third condition of \CGeqref{eq:list-space}, which replaces $\int_{\tau_0}^X(N_\zeta-n_{\RVM})=O(\log X)$ in every summation by parts against the count difference.
\item Theorems~\CGref{thm:mellin-counts} and~\CGref{thm:distributional-EF}: the prime-power count $J^{\rm M}_{\boldsymbol{\kappa}}$ as a Mellin transform, its $L^2_{\rm loc}$ bound, and the distributional explicit formula. Input: (i), (ii) and the Perron argument, in which the ordinates enter only through $\log\zeta_{\boldsymbol{\kappa}}$ on vertical lines, controlled by (ii).
\item Proposition~\CGref{prop:one-ordinate}: the cell masses $\tilde c_n(\boldsymbol{\kappa})$ as sums of a one-ordinate function. Input: (iii) and dominated convergence, with the domination supplied by the tube.
\item Lemma~\CGref{lem:mean-value-ordinates} with $\kappa_k$ in place of $\gamma_k(\lambda)$: the tube gives $N_{\boldsymbol{\kappa}}(t+1)-N_{\boldsymbol{\kappa}}(t)\le n_{\RVM}(t+1+C_\eta)-n_{\RVM}(t-C_\eta)\ll\log(t+2)$, with an implied constant depending on $C_\eta$ only.
\item Theorem~\CGref{thm:mean-square-cells} and Corollary~\CGref{cor:abscissa-one}: $\sum_{X<n\le2X}|\tilde c_n(\boldsymbol{\kappa})|^2\le C_\eps X^{1+\eps}$, the absolute convergence of $\tilde\zeta_{\Z,\boldsymbol{\kappa}}^{\pm1}$ in $\Re s>1$, and the exceptional-set bound. Input: (v) for the cells away from the powers of $2$; for the cells $n=2^j$ the term \CGeqref{eq:cell-mass-power-two}, whose bound needs $\|J_{\boldsymbol{\kappa}}\|_{L^2(2,3)}$ bounded uniformly on $\mathscr O$, which is (iii) together with the uniform bound of (ii); and the uniformity in $n$ of the bumps $\varphi_n$. The constant $C_\eps$ depends on $\eps$, $\tau_0^+$, $C_\eta$ and $C$ only. The statement on $\partial_\lambda\tilde c_n$ has no analogue for a single list and is not transferred.
\item Theorem~\CGref{thm:atomized-continuation}(i)--(iii) for $\tilde\zeta_{\Z,\boldsymbol{\kappa}}$. Input: (vi) and (iii).
\end{enumerate}
\end{theorem}

\begin{proof}
Each item names the only places where the cited proofs use the ordinates. They are four: the tube $|\kappa_k-\tau_k|\le C_\eta$, which gives $\kappa_k\asymp\tau_k$, $\sum_k\kappa_k^{-2}<\infty$, the local count of (v), and the domination of (iv); the ordering, used in the monotone rearrangements and in the identification of the zeros of $\zeta_{\boldsymbol{\kappa}}$ with the ordinates; the threshold $\kappa_1\ge\tau_0^+>\tau_0$, used wherever a branch of $\log\bigl(1+q(s)/(\frac14+\kappa^2)\bigr)$ is fixed on $|s-\frac12|<\tau_0$; and the third condition, used in every summation by parts over the ordered ordinates in place of Proposition~\CGref{prop:family-clock}. None of the cited proofs uses that the list is a zero set of $\zeta$ or a convex combination with the clock, except in the statements about $\lambda$-derivatives, which are excluded in (vi). Reading the proofs with these four inputs in place of their sources gives the items, and the constants are those of the inputs.
\end{proof}

\begin{theorem}[Rigidity of the coincidence]\label{thm:rigidity}
The set
\begin{equation}\label{eq:coincidence-set}
 \mathscr C:=\bigl\{\boldsymbol{\kappa}\in\mathscr O:\ \tilde h_{\boldsymbol{\kappa}}\ \text{is constant}\bigr\}
\end{equation}
of the lists whose continuous and atomized zeta functions coincide up to a constant factor has at most one element.  If $\boldsymbol{\kappa}\in\mathscr C$, then $\zeta_{\boldsymbol{\kappa}}=\tilde\zeta_{\Z,\boldsymbol{\kappa}}=\zeta$, $\boldsymbol{\kappa}$ is the list of the ordinates of $\zeta$, and (D) holds.  Conversely, (D) implies $\mathscr C=\{(\gamma_k)\}$.
\end{theorem}

\begin{proof}
If $\tilde h_{\boldsymbol{\kappa}}\equiv c$, then $\zeta_{\boldsymbol{\kappa}}=e^{-c}\tilde\zeta_{\Z,\boldsymbol{\kappa}}$ is an absolutely convergent Dirichlet series in $\Re s>1$ (Corollary~\CGref{cor:abscissa-one}), $(s-1)\zeta_{\boldsymbol{\kappa}}$ is entire of order one and $\zeta_{\boldsymbol{\kappa}}$ satisfies \CGeqref{eq:Z1-FE} (Theorem~\CGref{thm:path}).  Hamburger's theorem gives $\zeta_{\boldsymbol{\kappa}}=C_H\zeta$, and $C_H=1$ because both functions have residue $1$ at $s=1$.  Hence the nontrivial zeros of $\zeta$ with positive ordinate are the points $\frac12+i\kappa_k$, with multiplicity; this determines $\boldsymbol{\kappa}$ as the list $(\gamma_k)$, and $\zeta_{Z1}=\zeta_{\boldsymbol{\kappa}}=\zeta$.  Then $e^{-c}\tilde\zeta_{\Z,\boldsymbol{\kappa}}=\zeta$ and the first coefficients give $c=0$.  Conversely, under (D), $\tilde h_0\equiv0$ (Theorem~\CGref{thm:closing}), so $(\gamma_k)\in\mathscr C$.
\end{proof}

\begin{proposition}[Compactness and continuity]\label{prop:lists-compact}
The space $\mathscr O$ and the functions \CGeqref{eq:list-objects} have the following properties.
\begin{enumerate}[label=\textup{(\roman*)}]
\item $\mathscr O$ is compact in the product topology of $\prod_k[\tau_k-C_\eta,\tau_k+C_\eta]$.
\item For every $n\ge2$, $N\ge1$ and $\sigma>1$, the functions $\boldsymbol{\kappa}\mapsto\tilde c_n(\boldsymbol{\kappa})$, $\tilde a_N(\boldsymbol{\kappa})$, $k_N(\boldsymbol{\kappa})$ and $\mathsf V_\sigma(\boldsymbol{\kappa})$ are continuous on $\mathscr O$.
\item For every real $\sigma>1$, $\boldsymbol{\kappa}\mapsto\log\zeta_{\boldsymbol{\kappa}}(\sigma)=-\log G(\sigma)+\sum_k\log\frac{(\sigma-\frac12)^2+\kappa_k^2}{\frac14+\kappa_k^2}$ is a sum of strictly convex functions of the single coordinates $\kappa_k$.
\end{enumerate}
\end{proposition}

\begin{proof}
(i) The product $\mathcal K:=\prod_k[\tau_k-C_\eta,\tau_k+C_\eta]$ is compact by Tychonoff's theorem and metrizable, being a countable product of compact intervals. Each condition in \CGeqref{eq:list-space} defines a closed subset of $\mathcal K$: $\{\kappa_1\ge\tau_0^+\}$ is closed; the ordering is $\bigcap_k\{\kappa_k\le\kappa_{k+1}\}$, an intersection of closed sets; and for fixed $X$ the function $F_X(\boldsymbol{\kappa}):=\sum_k(X-\kappa_k)^+-\int_{\tau_0}^X\NR$ equals $\int_{\tau_0}^X(N_{\boldsymbol{\kappa}}-\NR)$, depends continuously on the finitely many coordinates with $\tau_k\le X+C_\eta$, and the third condition is $\bigcap_{X\ge\tau_0}\{|F_X|\le C\log X\}$, again an intersection of closed sets. So $\mathscr O$ is closed in $\mathcal K$, hence compact.
(ii) For $n$ not a power of $2$, $\tilde c_n(\boldsymbol{\kappa})=\mathfrak m_n-2\Re\sum_kW_n(\kappa_k)$ by \CGeqref{eq:cell-mass-W}, and the series converges uniformly on $\mathscr O$ because $W_n$ is rapidly decreasing and $|\kappa_k-\tau_k|\le C_\eta$.  For $n=2^j$ the term $j$ is $-\frac{\mu(j)}j\int_2^3J_{\boldsymbol{\kappa}}(y)\frac d{dy}\varphi_n(y^j)\,dy$ (proof of Theorem~\CGref{thm:mean-square-cells}), and $\boldsymbol{\kappa}\mapsto J_{\boldsymbol{\kappa}}\in L^2(2,3)$ is continuous, by the following three steps.  (1) Fix $\sigma>1$ and put $\mathsf F_{\boldsymbol{\kappa}}(t):=\log\zeta_{\boldsymbol{\kappa}}(\sigma+it)/(\sigma+it)$.  For fixed $t$, $\log\zeta_{\boldsymbol{\kappa}}(\sigma+it)=-\log G(\sigma+it)+\sum_k\Log\bigl(1+q(\sigma+it)/(\frac14+\kappa_k^2)\bigr)$ is a series converging uniformly on $\mathscr O$ (the $k$-th term is $O(|q|/\tau_k^2)$ on the tube), hence continuous in the product topology; so $\mathsf F_{\boldsymbol{\kappa}}(t)$ is continuous in $\boldsymbol{\kappa}$ for every $t$.  (2) By Theorem~\CGref{thm:transfer-O}(ii) (the bound of Proposition~\CGref{prop:family-clock} on $\mathscr O$), $|\log\zeta_{\boldsymbol{\kappa}}(\sigma+it)-\log\frac{s}{s-1}|\le C_\sigma(1+\log(2+|t|))$ for all $\boldsymbol{\kappa}\in\mathscr O$, so that
\begin{equation}\label{eq:lists-compact-dominant}
 |\mathsf F_{\boldsymbol{\kappa}}(t)|\le\mathsf H(t):=\frac{C'_\sigma\,(1+\log(2+|t|))}{1+|t|},\qquad \mathsf H\in L^2(\R),
\end{equation}
a pointwise majorant independent of $\boldsymbol{\kappa}$; if $\boldsymbol{\kappa}^{(m)}\to\boldsymbol{\kappa}$ in $\mathscr O$, then $|\mathsf F_{\boldsymbol{\kappa}^{(m)}}-\mathsf F_{\boldsymbol{\kappa}}|^2\le4\mathsf H^2$ and dominated convergence gives $\mathsf F_{\boldsymbol{\kappa}^{(m)}}\to\mathsf F_{\boldsymbol{\kappa}}$ in $L^2(\R)$; by Plancherel and Theorem~\CGref{thm:mellin-counts}, $J^{\rm M}_{\boldsymbol{\kappa}^{(m)}}(e^u)e^{-\sigma u}\to J^{\rm M}_{\boldsymbol{\kappa}}(e^u)e^{-\sigma u}$ in $L^2(\R,du)$, hence $J^{\rm M}_{\boldsymbol{\kappa}^{(m)}}\to J^{\rm M}_{\boldsymbol{\kappa}}$ in $L^2(2,3)$.  (3) The constant $C_{\boldsymbol{\kappa}}:=J^{\rm M}_{\boldsymbol{\kappa}}-J_{\boldsymbol{\kappa}}$ equals $\langle J^{\rm M}_{\boldsymbol{\kappa}},\phi\rangle-\langle J_{\boldsymbol{\kappa}},\phi\rangle$ for a fixed $\phi\in C_c^\infty(2,3)$ with $\int\phi=1$: the first pairing is continuous by (2); the second is the sum over $k$ of $\int\bigl(\Ei(\rho_k\log x)+\Ei(\bar\rho_k\log x)\bigr)\phi(x)\,dx$ plus terms independent of the list, a series whose tails are bounded uniformly on $\mathscr O$ by \CGeqref{eq:L2-partial-sums} (whose constant depends on $C_\eta$ only) and whose terms are continuous in $\kappa_k$, hence a continuous function on $\mathscr O$.  Therefore $J_{\boldsymbol{\kappa}}=J^{\rm M}_{\boldsymbol{\kappa}}-C_{\boldsymbol{\kappa}}$ is continuous into $L^2(2,3)$, and the cells $n=2^j$ are continuous.  $\tilde a_N$ and $k_N$ are polynomials in finitely many cell masses.  The sums $\sum_n|k_n(\boldsymbol{\kappa})|n^{-\sigma'}$, $\sigma'>1$, are bounded uniformly on $\mathscr O$ (Theorem~\CGref{thm:mean-square-cells}), so the tails of $\mathsf V_\sigma$ are uniformly small, as in the proof of Theorem~\CGref{thm:double-hamburger}(iii).
(iii) is the computation in the proof of Proposition~\CGref{prop:convexity}: $\partial_\kappa^2\log\frac{a^2+\kappa^2}{1/4+\kappa^2}=\mathcal C(a,\kappa)-\mathcal C(\frac12,\kappa)>0$ for $a>\frac12$, $\kappa>\frac12$.
\end{proof}

\begin{corollary}[Existence form of the final step]\label{cor:existence-form}
For every $\sigma>1$ the minimum $\min_{\mathscr O}\mathsf V_\sigma$ is attained, and the following are equivalent: \textup{(i)} (D); \textup{(ii)} $\mathscr C\ne\varnothing$; \textup{(iii)} there is $\boldsymbol{\kappa}\in\mathscr O$ (Section~\CGref{sec:lists}) with $\tilde\zeta_{\Z,\boldsymbol{\kappa}}=\zeta$; \textup{(iv)} $\min_{\mathscr O}\mathsf V_\sigma=0$.
\end{corollary}

\begin{proof}
A continuous function on a compact space attains its minimum (Proposition~\CGref{prop:lists-compact}).  (i)$\Leftrightarrow$(ii) is Theorem~\CGref{thm:rigidity}, and (ii)$\Rightarrow$(iii) is part of it.  (iii)$\Leftrightarrow$(iv) because $\mathsf V_\sigma(\boldsymbol{\kappa})=0$ means $k_n(\boldsymbol{\kappa})=0$ for $n\ge2$.  (iii)$\Rightarrow$(i): $\tilde\zeta_{\Z,\boldsymbol{\kappa}}=\zeta$ has no zeros in $\Re s>\frac12$ (Theorem~\CGref{thm:atomized-continuation}(iii)), so, with the functional equation, the nontrivial zeros of $\zeta$ lie on $\Re s=\frac12$: the Riemann hypothesis holds, and Theorem~\CGref{thm:Hamburger-Z1} gives (D).
\end{proof}

In this form the final step does not ask to identify $\zeta_{Z1}$ with $\zeta$: it asks for the existence of one list in $\mathscr O$ whose continuous and atomized zeta functions coincide, or whose atomized zeta function is $\zeta$; Theorem~\CGref{thm:rigidity} then identifies the list.  The path $\lambda\mapsto(\gamma_k(\lambda))$ is one segment in $\mathscr O$, from the clock to the list of $\zeta$, and along it $\mathsf V_\sigma$ is the function of Theorem~\CGref{thm:double-hamburger}.

\begin{remark}[What an existence proof uses]\label{rem:existence-input}
Compactness, continuity and convexity do not decide whether $\min_{\mathscr O}\mathsf V_\sigma$ is $0$: they do not involve the form of the functional equation.  The existence statement uses the functional equation with the factor $\Gamma(\frac s2)\pi^{-s/2}$, for which a function with both an absolutely convergent Dirichlet series and the functional equation exists, namely $\zeta$.  The first variation of the defect is explicit: by \CGeqref{eq:cell-mass-W}, for $n$ not a power of $2$,
\begin{equation}\label{eq:first-variation}
 \frac{\partial\tilde c_n(\boldsymbol{\kappa})}{\partial\kappa_k}=2\,\Im V_n(\kappa_k),
\end{equation}
so the vanishing of all coordinate derivatives of $\mathsf V_\sigma$ at a list is described by the functions $V_n$ of \CGeqref{eq:VW} alone.  Since $\mathscr O$ has empty interior in the product topology and the displayed derivatives are coordinate derivatives rather than a proved Fr\'echet gradient on a fixed Banach space, boundary minimizers must be treated by one-sided/KKT conditions.
\end{remark}


\subsection{The rung system with free masses}
\label{subsec:free-masses}

With the function $\mathsf k$ of \CGeqref{eq:transport-defs}, consider, for a list $\boldsymbol{\kappa}\in\mathscr O$ (Section~\CGref{sec:lists}) and real numbers $\mathsf a_n$, $n\ge2$, on which nothing is assumed, the system
\begin{equation}\label{eq:rung-system}
 \sum_{k\ge1}\cos(\kappa_ku)+\sum_{n\ge2}\mathsf a_n\,\delta(u-\log n)=\mathsf k(u)\qquad\text{in }\mathcal D'(0,\infty).
\end{equation}
In the variable $x=e^u$ it states that the Mellin count of $\zeta_{\boldsymbol{\kappa}}$ has no mass outside the integers, the masses at the integers being unknown.

\begin{theorem}[Existence form with free masses]\label{thm:free-masses}
(D) holds if and only if \CGeqref{eq:rung-system} has a solution $(\boldsymbol{\kappa},(\mathsf a_n))$ with $\boldsymbol{\kappa}\in\mathscr O$ (Section~\CGref{sec:lists}).  The solution is then unique: $\boldsymbol{\kappa}=(\gamma_k)$, and
\begin{equation}\label{eq:free-masses}
 \mathsf a_n=\frac{\log p}{2\sqrt n}\ \text{ if }n=p^m\text{ is a prime power},\qquad \mathsf a_n=0\ \text{ otherwise.}
\end{equation}
\end{theorem}

\begin{proof}
By \CGeqref{eq:distributional-EF} for the list $\boldsymbol{\kappa}$, with $u=\log x$, \CGeqref{eq:rung-system} is the statement
\begin{equation}\label{eq:free-masses-dJ}
 (\log x)\,dJ^{\rm M}_{\boldsymbol{\kappa}}(x)=\sum_{n\ge2}2\sqrt n\,\mathsf a_n\,\delta_n\qquad\text{in }\mathcal D'(1,\infty).
\end{equation}
If (D) holds, $J^{\rm M}_{Z1}=\sum_{p^m\le x}\frac1m$ (Theorem~\CGref{thm:cell-form}), which is \CGeqref{eq:free-masses-dJ} for the list $(\gamma_k)$ with the masses \CGeqref{eq:free-masses}.  Conversely let \CGeqref{eq:free-masses-dJ} hold.  Then $J^{\rm M}_{\boldsymbol{\kappa}}$ is constant on $(1,2)$ and between consecutive integers; the truncation at $2$ and the finite M\"obius inversion keep the support in the integers, so every cell defect vanishes and, by \CGeqref{eq:first-cell} and the proof of Theorem~\CGref{thm:joining-factor},
\begin{equation}\label{eq:free-masses-factor}
 \zeta_{\boldsymbol{\kappa}}(s)=\exp\bigl(\nu_6+\nu_7\,2^{-s}\bigr)\,\tilde\zeta_{\Z,\boldsymbol{\kappa}}(s)\qquad(\Re s>1)
\end{equation}
with real constants $\nu_6,\nu_7$.  The two factors on the right are the exponential $\exp(\nu_6+\nu_72^{-s})=e^{\nu_6}\sum_{r\ge0}\frac{\nu_7^r}{r!}2^{-rs}$, an absolutely convergent Dirichlet series on all of $\C$ supported on the powers of $2$, and $\tilde\zeta_{\Z,\boldsymbol{\kappa}}$, an absolutely convergent Dirichlet series in $\Re s>1$ by Corollary~\CGref{cor:abscissa-one}, which Theorem~\CGref{thm:transfer-O}(vi) extends to every list of $\mathscr O$; the product of two absolutely convergent Dirichlet series is one.  (Why the support stays in the integers: $J^{\rm M}_{\boldsymbol{\kappa}}$ constant between consecutive integers means $dJ^{\rm M}_{\boldsymbol{\kappa}}$ is a combination of $\delta_n$, $n\ge2$; the truncation $J^*$ at $2$ removes the cell $(1,2)$, on which $J^{\rm M}_{\boldsymbol{\kappa}}$ is a constant, and the finite M\"obius inversion \CGeqref{eq:pi-RVM} maps $\sum_n\mathsf w_n\delta_n$ to $\sum_m\frac{\mu(m)}m\sum_n\mathsf w_n\delta_{n^{1/m}}$, whose support is in the integers exactly when the cell masses vanish off the integers, which is the vanishing of every cell defect $\mathfrak h_n$; the constant on $(1,2)$ is what produces the factor $e^{\nu_6+\nu_72^{-s}}$ through \CGeqref{eq:first-cell}.)  Hence $\zeta_{\boldsymbol{\kappa}}$ satisfies (H1), and it satisfies (H0), (H2) and (H3) by Theorem~\CGref{thm:transfer-O}(i) (order one, functional equation \CGeqref{eq:Z1-FE}, pole at $1$ with residue $1$) for every list of $\mathscr O$; so Hamburger's theorem in the form of Subsection~\CGref{subsec:hamburger-table} gives $\zeta_{\boldsymbol{\kappa}}=C_H\zeta$, and $C_H=1$ by the residues at $s=1$.  The zeros of $\zeta_{\boldsymbol{\kappa}}$ are $\frac12\pm i\kappa_k$, so $\boldsymbol{\kappa}=(\gamma_k)$ and $\zeta_{Z1}=\zeta$.  Then $(\log x)\,dJ^{\rm M}=\sum_{p^m}\log p\,\delta_{p^m}$, which gives \CGeqref{eq:free-masses}.
\end{proof}

In Theorem~\CGref{thm:free-masses} the primes are not part of the data: the data are the function $\mathsf k$, which comes from the factor $G$ and the pole, the space $\mathscr O$, which is defined by the clock, and the condition that the support lie in the points $\log n$.  The masses of von Mangoldt are an output.  What the theorem proves is \emph{uniqueness}: a solution, if it exists, is the zeta list with the von Mangoldt masses.  \emph{Existence} of a solution is exactly (D) and is not proved; the theorem is a reformulation of the final step, not a step towards it.  Experiment~\CGref{exp:clock-to-zeros} solves a finite section of \CGeqref{eq:rung-system} numerically, starting from the clock; nothing proved here says that finite sections converge to a solution of the infinite distributional system.

\begin{theorem}[Nondegeneracy of the solution]\label{thm:nondegenerate}
Let $(\boldsymbol{\kappa},(\mathsf a_n))$ be a solution of \CGeqref{eq:rung-system} with $\boldsymbol{\kappa}\in\mathscr O$ (Section~\CGref{sec:lists}).  Let $(\dot\kappa_k)$ be a bounded real sequence and $(\dot{\mathsf a}_n)_{n\ge2}$ real numbers with $\dot{\mathsf a}_n=O_\eps(n^{-1/2+\eps})$ for every $\eps>0$, such that
\begin{equation}\label{eq:rung-linearized}
 \sum_{k\ge1}\dot\kappa_k\,u\sin(\kappa_ku)=\sum_{n\ge2}\dot{\mathsf a}_n\,\delta(u-\log n)\qquad\text{in }\mathcal D'(0,\infty),
\end{equation}
the derivative of \CGeqref{eq:rung-system} in the direction $\bigl((\dot\kappa_k),(\dot{\mathsf a}_n)\bigr)$.  Then $\dot{\mathsf a}_n=0$ for every $n$, and $\sum_{k:\,\kappa_k=t}\dot\kappa_k=0$ for every value $t$ of the list.  If the ordinates are simple, $\dot\kappa_k=0$ for every $k$.
\end{theorem}

\begin{proof}
\emph{Step 1.}  The series $\mathsf J:=\sum_k\dot\kappa_k\,u\sin(\kappa_ku)$ converges in $\mathcal S'(\R)$: for $\phi\in\mathcal S$, three integrations by parts give $\int\phi(u)u\sin(\kappa u)\,du=\kappa^{-3}\int(u\phi)'''(u)\cos(\kappa u)\,du$ (the boundary terms vanish since $u\phi$ and its derivatives are Schwartz), so $|\int\phi\,u\sin(\kappa u)|\le\|(u\phi)'''\|_{L^1}\kappa^{-3}$, and $\sum_k|\dot\kappa_k|\kappa_k^{-3}<\infty$ since $\kappa_k\asymp k/\log k$.  Each term $u\sin(\kappa_ku)$ is even, so $\mathsf J$ is even.  By \CGeqref{eq:rung-linearized}, $\mathsf J$ vanishes on $(0,\log2)$, hence by evenness on $(-\log2,0)$, so the restriction of $\mathsf J$ to $(-\log2,\log2)$ is supported at $\{0\}$.  We use the following standard fact.

\emph{Lemma (point support).}  A distribution $\mathsf T\in\mathcal D'(-a,a)$ with $\operatorname{supp}\mathsf T\subset\{0\}$ is a finite combination $\mathsf T=\sum_{j\le m}c_j\delta_0^{(j)}$, with $m$ the order of $\mathsf T$ \cite[Thm.~2.3.4]{Hormander1990}; consequently $\langle\mathsf T,\chi\rangle=\sum_{j\le m}(-1)^jc_j\chi^{(j)}(0)=0$ for every $\chi\in C_c^\infty(-a,a)$ vanishing to order $m$ at $0$, and the same formula defines $\langle\mathsf T,\chi\rangle$ for $\chi\in C^m_c(-a,a)$.  Applied to $\mathsf J|_{(-\log2,\log2)}$ and using evenness, $\mathsf J=\sum_{j\le m,\ j\text{ even}}c_j\delta_0^{(j)}$ on $(-\log2,\log2)$ for some finite $m$.  Fix an even $\mathsf x_0\in C_c^\infty(-\log2,\log2)$ equal to $1$ near $0$, and for $\Re z>0$ put
\begin{equation}\label{eq:nondeg-H1}
 \mathsf H_1(z):=\sum_{k\ge1}\dot\kappa_k\int_0^\infty\bigl(1-\mathsf x_0(u)\bigr)\,u\sin(\kappa_ku)\,e^{-zu}\,du=\tfrac12\bigl\langle\mathsf J,(1-\mathsf x_0)\,e^{-z|u|}\bigr\rangle .
\end{equation}
Approximating by test functions with compact support in $(0,\infty)$ and using \CGeqref{eq:rung-linearized}, $\mathsf H_1(z)=\sum_{n\ge2}\dot{\mathsf a}_nn^{-z}$ for $\Re z>\frac12$, where the series converges absolutely.

\emph{Step 2.}  For $\Re z>0$, $\int_0^\infty u\sin(\kappa u)e^{-zu}du=2\kappa z/(z^2+\kappa^2)^2$, so
\begin{equation}\label{eq:nondeg-split}
 \mathsf H_1(z)=\sum_{k\ge1}\dot\kappa_k\,\frac{2\kappa_kz}{(z^2+\kappa_k^2)^2}-\mathsf y(z),\qquad \mathsf y(z):=\sum_{k\ge1}\dot\kappa_k\int_0^{\log2}\mathsf x_0(u)\,u\sin(\kappa_ku)\,e^{-zu}\,du .
\end{equation}
\emph{$\mathsf y$ is a polynomial.}  Let $m$ be the order of $\mathsf J$ at $0$ from Step~1 and fix an integer $A\ge\max(m,3)$.  Write $e^{-zu}=\sum_{j\le A}\frac{(-z)^j}{j!}u^j+\mathsf R_A(z,u)$ with $\mathsf R_A(z,u)=\frac{(-z)^{A+1}}{A!}\int_0^1(1-\vartheta)^Au^{A+1}e^{-\vartheta zu}\,d\vartheta$, so that $\mathsf R_A(z,\cdot)$ is smooth on $[0,\log2]$ and vanishes to order $A+1$ at $u=0$.  Accordingly $\mathsf y=\mathsf y_{\rm pol}+\mathsf y_{\rm rem}$.  \emph{Polynomial part.}  $\mathsf y_{\rm pol}(z)=\sum_{j\le A}\frac{(-z)^j}{j!}\sum_k\dot\kappa_k\,\mathsf i_{j,k}$ with $\mathsf i_{j,k}:=\int_0^{\log2}\mathsf x_0(u)u^{j+1}\sin(\kappa_ku)\,du$.  For $g\in C^\infty_c[0,\infty)$ two integrations by parts give $\int_0^\infty g(u)\sin(\kappa u)\,du=g(0)/\kappa-\kappa^{-2}\int_0^\infty g''(u)\sin(\kappa u)\,du=g(0)/\kappa-g''(0)/\kappa^3+O_g(\kappa^{-5})$; with $g=\mathsf x_0u^{j+1}$ one has $g(0)=0$ for all $j\ge0$, so $\mathsf i_{j,k}=O_j(\kappa_k^{-3})$ and each inner series converges absolutely since $\sum_k|\dot\kappa_k|\kappa_k^{-3}<\infty$.  Thus $\mathsf y_{\rm pol}$ is a polynomial of degree at most $A$.  \emph{Remainder part.}  Put $\chi_z(u):=\mathsf x_0(u)\mathsf R_A(z,|u|)$, an even function with compact support in $(-\log2,\log2)$, of class $C^A$ (because $|u|^{A+1}$ is $C^A$) and vanishing to order $A+1$ at $0$.  Since $u\sin(\kappa u)$ is even, $\mathsf y_{\rm rem}(z)=\frac12\sum_k\dot\kappa_k\int_{\R}\chi_z(u)\,u\sin(\kappa_ku)\,du$.  Let $\chi_{z,\epsilon}:=\chi_z*\varrho_\epsilon$ with a standard mollifier; for small $\epsilon$ these are in $C_c^\infty(-\log2,\log2)$, $\|(u\chi_{z,\epsilon})'''\|_{L^1}$ is bounded uniformly in $\epsilon$ because $A\ge3$, and $\chi_{z,\epsilon}^{(j)}(0)\to\chi_z^{(j)}(0)=0$ for $j\le A$.  By Step~1 and the point-support lemma, $\sum_k\dot\kappa_k\int\chi_{z,\epsilon}\,u\sin(\kappa_ku)\,du=\langle\mathsf J,\chi_{z,\epsilon}\rangle=\sum_{j\le m}(-1)^jc_j\chi_{z,\epsilon}^{(j)}(0)\to0$ as $\epsilon\to0$; on the other side, each term converges to $\dot\kappa_k\int\chi_z\,u\sin(\kappa_ku)\,du$ and the series is dominated by $\sum_k|\dot\kappa_k|\kappa_k^{-3}\sup_\epsilon\|(u\chi_{z,\epsilon})'''\|_{L^1}<\infty$, so the limit may be taken termwise.  Hence $\mathsf y_{\rm rem}\equiv0$ and $\mathsf y=\mathsf y_{\rm pol}$ is a polynomial.

\emph{$\mathsf y\equiv0$.}  As $z=\sigma\to+\infty$: (a) $\mathsf H_1(\sigma)=\sum_{n\ge2}\dot{\mathsf a}_nn^{-\sigma}\to0$, since for $\sigma\ge1$ the terms are dominated by $|\dot{\mathsf a}_n|n^{-1}$, $\sum_n|\dot{\mathsf a}_n|n^{-1}<\infty$ because $\dot{\mathsf a}_n=O_\eps(n^{-1/2+\eps})$, and each term tends to $0$; (b) the first sum in \CGeqref{eq:nondeg-split} is $O(\sigma^{-1}\log\sigma)$: the terms with $\kappa_k\le\sigma$ are at most $2\kappa_k\sigma/\sigma^4$ and add up to at most $2\sigma^{-3}\sum_{\kappa_k\le\sigma}\kappa_k\le2\sigma^{-2}N_{\boldsymbol{\kappa}}(\sigma)=O(\sigma^{-1}\log\sigma)$, while the terms with $\kappa_k>\sigma$ are at most $2\sigma\kappa_k^{-3}$ and, by partial summation with $N_{\boldsymbol{\kappa}}(t)=O(t\log t)$, $\sum_{\kappa_k>\sigma}\kappa_k^{-3}=\int_\sigma^\infty t^{-3}dN_{\boldsymbol{\kappa}}(t)=O(\sigma^{-2}\log\sigma)$ (the boundary term $N_{\boldsymbol{\kappa}}(\sigma)\sigma^{-3}$ included), which gives $O(\sigma^{-1}\log\sigma)$ again; here $(\dot\kappa_k)$ is bounded.  Hence the polynomial $\mathsf y(\sigma)=\sum_k\dot\kappa_k2\kappa_k\sigma/(\sigma^2+\kappa_k^2)^2-\mathsf H_1(\sigma)\to0$ as $\sigma\to+\infty$, so $\mathsf y\equiv0$.

\emph{Step 3.}  Put $\mathsf H(z):=\sum_{n\ge2}\frac{\dot{\mathsf a}_n}{\log n}n^{-z}$ for $\Re z>\frac12$, so that $\mathsf H'=-\mathsf H_1$.  By Step~2 there is a real constant $\nu_8$ with
\begin{equation}\label{eq:nondeg-H}
 \mathsf H(z)=\nu_8+\sum_{k\ge1}\dot\kappa_k\Bigl(\frac{\kappa_k}{z^2+\kappa_k^2}-\frac1{\kappa_k}\Bigr),
\end{equation}
which continues $\mathsf H$ to a meromorphic function in $\C$, even in $z$, with at most simple poles at $z=\pm i\kappa_k$, of residue $\pm\frac1{2i}\sum_{k':\kappa_{k'}=\kappa_k}\dot\kappa_{k'}$.

\emph{Step 4: finite order.}  By Theorem~\CGref{thm:free-masses}, $\zeta_{\boldsymbol{\kappa}}=\zeta$, with zeros $\frac12\pm i\kappa_k$ of the multiplicity of the list.  Put $\mathsf Z(s):=\mathsf H(s-\frac12)\,\zeta(s)$.  At $s=\frac12\pm i\kappa$ the function $\mathsf H(s-\frac12)$ has at most a simple pole (the residues of the entries of the list with the same value $\kappa$ add up, \CGeqref{eq:nondeg-H}) and $\zeta$ has a zero of order equal to the multiplicity of $\kappa$ in the list, which is at least $1$; so $\mathsf Z$ is holomorphic there and $(s-1)\mathsf Z(s)$ is entire.  \emph{Admissible radii.}  Call $r\ge3$ admissible if $\operatorname{dist}(r,\{\kappa_k\})\ge1/\log r$.  Since $N_{\boldsymbol{\kappa}}(2R+1)-N_{\boldsymbol{\kappa}}(R-1)=(\frac R{2\pi}+o(R))\log R$ (bounded displacement from the clock), the union of the intervals of radius $1/\log R$ about the $\kappa_k\in[R-1,2R+1]$ has total length at most $(\frac1\pi+o(1))R<R$, so every interval $[R,2R]$ with $R\ge R_0$ contains an admissible radius; choose admissible $r_j\in[2^jR_0,2^{j+1}R_0]$, so that $r_j\to\infty$ and $r_{j+1}\le4r_j$.  \emph{Bound on admissible circles.}  Write \CGeqref{eq:nondeg-H} as $\mathsf H(z)=\nu_8-\sum_k\dot\kappa_k\frac{z^2}{\kappa_k(z^2+\kappa_k^2)}$ and let $|z|=r$ be admissible.  Then $|z\mp i\kappa_k|\ge|r-\kappa_k|$ and $|z-i\kappa_k|+|z+i\kappa_k|\ge2\kappa_k$.  Shell $\kappa_k\ge2r$: $|z^2+\kappa_k^2|\ge\kappa_k^2-r^2\ge\frac34\kappa_k^2$, terms $\le\frac43r^2\kappa_k^{-3}\sup|\dot\kappa|$, sum $O(r^2\cdot r^{-2}\log r)=O(\log r)$ by partial summation.  Shell $\kappa_k\le r/2$: $|z^2+\kappa_k^2|\ge\frac34r^2$, terms $\le\frac43\kappa_k^{-1}\sup|\dot\kappa|$, sum $O(\sum_{\kappa_k\le r}\kappa_k^{-1})=O(\log^2r)$.  Shell $r/2<\kappa_k<2r$: one of $|z\pm i\kappa_k|$ is at least $\kappa_k\ge r/2$ and the other is at least $\max(|r-\kappa_k|,1/\log r)$, so the term is at most $\frac{4r}{\kappa_k}\cdot\frac{\sup|\dot\kappa|}{\max(|r-\kappa_k|,1/\log r)}\le\frac{8\sup|\dot\kappa|}{\max(|r-\kappa_k|,1/\log r)}$; the $O(\log r)$ ordinates with $|\kappa_k-r|\le1$ contribute $O(\log^2r)$, and for $i\ge0$ the $O(2^i\log r)$ ordinates with $2^i<|\kappa_k-r|\le2^{i+1}$ contribute $O(\log r)$ each shell, $i\le\log_2r$, so $O(\log^2r)$ in all (multiplicities are counted by $N_{\boldsymbol{\kappa}}$ throughout).  Hence $|\mathsf H(z)|\le C_{\mathsf H}\log^2r$ on every admissible circle $|z|=r\ge R_0$.  \emph{From circles to the plane.}  $(s-1)\zeta(s)$ is entire of order one, so $|(s-1)\zeta(s)|\le\exp(|s|^{1+\eps})$ for $|s|$ large.  By the maximum principle for the entire function $(s-1)\mathsf Z(s)$ on the disc $|s-\frac12|\le r_j$,
\begin{equation}\label{eq:nondeg-order}
 \max_{|s-\frac12|\le r_j}|(s-1)\mathsf Z(s)|\le C_{\mathsf H}\log^2r_j\cdot\max_{|s-\frac12|=r_j}|(s-1)\zeta(s)|\le\exp\bigl(2r_j^{1+\eps}\bigr)
\end{equation}
for $j\ge j_0(\eps)$.  Given $R\ge r_{j_0}$ choose the least $j$ with $r_j\ge R$; then $r_j\le4R$ and \CGeqref{eq:nondeg-order} gives $\max_{|s-\frac12|\le R}|(s-1)\mathsf Z|\le\exp(2(4R)^{1+\eps})$.  So $(s-1)\mathsf Z$ is entire of order at most one: this is (H0).  $\mathsf Z$ is the product of $\sum_n\dot{\mathsf a}_n\sqrt n(\log n)^{-1}n^{-s}$ and $\zeta(s)$, both absolutely convergent in $\Re s>1$: this is (H1).  Since $\mathsf H$ is even, $\mathsf Z(1-s)=\mathsf H(s-\frac12)\zeta(1-s)$, so $\mathsf Z$ satisfies the functional equation of $\zeta$: this is (H2) and (H3) with the second function equal to $\mathsf Z$.  Hamburger's theorem gives $\mathsf Z=C_H\zeta$, that is, $\mathsf H\equiv C_H$.  Since $\mathsf H(\sigma)\to0$ as $\sigma\to+\infty$, $\mathsf H\equiv0$: every $\dot{\mathsf a}_n$ and every residue vanishes.
\end{proof}

\begin{remark}[Scope of Theorem~\CGref{thm:nondegenerate}]\label{rem:nondegenerate-scope}
The theorem is conditional on the existence of a solution; it shows that a solution is isolated to first order, the Jacobian of \CGeqref{eq:rung-system} being injective there on the directions considered, except for the splitting of multiple ordinates.  This is a necessary endpoint nondegeneracy condition for any continuation scheme; it is not, by itself, an implicit-function theorem, because no isomorphism between fixed Banach spaces is proved.  Steps~1--3 do not use that $\boldsymbol{\kappa}$ is a solution: at any list of $\mathscr O$ the kernel of the linearization consists of the Dirichlet series $\mathsf H$ that are even meromorphic functions with simple poles at $\pm i\kappa_k$ of the form \CGeqref{eq:nondeg-H}.  Such series exist for lists in arithmetic progression, for instance $(2^{z}+2^{-z}-2\cos\nu_9)^{-1}$, which is even, a power series in $2^{-z}$ for $\Re z>0$, with simple poles at $z=\pm i(\nu_9+2\pi m)/\log2$, $m\in\Z$; these lists do not belong to $\mathscr O$.  For perturbations $\dot{\mathsf a}_n$ of polynomial growth the same proof applies with Hamburger's theorem for a finite abscissa of absolute convergence \cite{BochnerChandrasekharan1956}.  The injectivity of the Jacobian at the clock and along a path, and with it the existence of a solution, are not proved.
\end{remark}

\subsection{The gradient of the vertical defect}
\label{subsec:gradient}

For Dirichlet series $A=\sum_na_nn^{-s}$, $B=\sum_nb_nn^{-s}$ and $\sigma>1$ put $\langle A,B\rangle_\sigma:=\sum_{n\ge1}a_n\bar b_n\,n^{-2\sigma}$, and for $\boldsymbol{\kappa}\in\mathscr O$ (Section~\CGref{sec:lists})
\begin{equation}\label{eq:K-kappa}
 \mathsf K_{\boldsymbol{\kappa}}:=\frac{\tilde\zeta_{\Z,\boldsymbol{\kappa}}}\zeta=\sum_{n\ge1}k_n(\boldsymbol{\kappa})\,n^{-s},\qquad \mathsf V_\sigma(\boldsymbol{\kappa})=\langle\mathsf K_{\boldsymbol{\kappa}},\mathsf K_{\boldsymbol{\kappa}}\rangle_\sigma-1,\qquad \mathsf r_m(\boldsymbol{\kappa}):=\Re\langle\mathsf K_{\boldsymbol{\kappa}}\mathsf L_m,\mathsf K_{\boldsymbol{\kappa}}\rangle_\sigma ,
\end{equation}
with $\mathsf L_m(s)=-\log(1-m^{-s})$ as in Proposition~\CGref{prop:one-ordinate}.  By Proposition~\CGref{prop:one-ordinate} the derivative $\partial\tilde c_m(\boldsymbol{\kappa})/\partial\kappa_k$ is a function of $m$ and of the value $\kappa_k$ alone; denote it by $\mathsf u_m(\kappa_k)$, so that $\mathsf u_m=2\,\Im V_m$ when $m$ is not a power of $2$ \CGeqref{eq:first-variation}.

\begin{proposition}[Critical points of the vertical defect]\label{prop:critical-points}
Let $\sigma>1$ and $\boldsymbol{\kappa}\in\mathscr O$ (Section~\CGref{sec:lists}).  Then $|\mathsf r_m(\boldsymbol{\kappa})|\le2m^{-\sigma}\bigl(1+\mathsf V_\sigma(\boldsymbol{\kappa})\bigr)$, $|\mathsf u_m(t)|\le Cm^{-1/2}$ off the powers of $2$ and $|\mathsf u_m(t)|\le C$ on them, and
\begin{equation}\label{eq:gradient}
 \frac{\partial\mathsf V_\sigma}{\partial\kappa_k}(\boldsymbol{\kappa})=2\,\mathsf Q_{\boldsymbol{\kappa}}(\kappa_k),\qquad \mathsf Q_{\boldsymbol{\kappa}}(t):=\sum_{m\ge2}\mathsf r_m(\boldsymbol{\kappa})\,\mathsf u_m(t),
\end{equation}
the series converging absolutely and uniformly in $t$.  If $\mathsf V_\sigma(\boldsymbol{\kappa})=0$, then $\mathsf r_m(\boldsymbol{\kappa})=0$ for every $m$ and $\mathsf Q_{\boldsymbol{\kappa}}\equiv0$.
\end{proposition}

\begin{proof}
Since $\log\mathsf K_{\boldsymbol{\kappa}}=\sum_m\bigl(\tilde c_m(\boldsymbol{\kappa})-\mathbf 1_{\rm prime}(m)\bigr)\mathsf L_m$, $\partial\mathsf K_{\boldsymbol{\kappa}}/\partial\tilde c_m=\mathsf K_{\boldsymbol{\kappa}}\mathsf L_m$ and $\partial\mathsf V_\sigma/\partial\tilde c_m=2\mathsf r_m$.  Multiplication by $m^{-rs}$ multiplies the norm $\langle\cdot,\cdot\rangle_\sigma^{1/2}$ by $m^{-r\sigma}$, so $\langle\mathsf K\mathsf L_m,\mathsf K\mathsf L_m\rangle_\sigma^{1/2}\le\langle\mathsf K,\mathsf K\rangle_\sigma^{1/2}\sum_{r\ge1}m^{-r\sigma}/r\le2m^{-\sigma}\langle\mathsf K,\mathsf K\rangle_\sigma^{1/2}$, and the bound for $\mathsf r_m$ follows by the Cauchy--Schwarz inequality.  In \CGeqref{eq:VW} the support of $\varphi_m(y^j)$ has length at most $\frac2jm^{1/j-1}$ and the integrand is at most $m^{-1/(2j)}$ in absolute value, so $|V_m(t)|\le\sum_j\frac2{j^2}m^{-1+1/(2j)}\le Cm^{-1/2}$; for $m=2^j$ the cut term is $-\frac{\mu(j)}j\int_2^3J_{\boldsymbol{\kappa}}(y)\frac d{dy}\varphi_m(y^j)\,dy$ (proof of Theorem~\CGref{thm:mean-square-cells}), and its coordinate derivative needs $\partial_{\kappa_k}J_{\boldsymbol{\kappa}}$ on $(2,3)$: in the series \CGeqref{eq:J-X} only the $k$-th term depends on $\kappa_k$, so the derivative exists and equals $\partial_{\kappa_k}J_{\boldsymbol{\kappa}}(y)=-2\,\Re\bigl(i\,y^{1/2+i\kappa_k}/(\tfrac12+i\kappa_k)\bigr)$, a bounded continuous function of $y\in(2,3)$ with $|\partial_{\kappa_k}J_{\boldsymbol{\kappa}}(y)|\le2\sqrt3/\kappa_k\le2\sqrt3/\tau_0^+$ uniformly on $\mathscr O$ (Lemma~\CGref{lem:L2-partial-sums}(b) is the same computation for the direction $\eta_k$); hence $|\mathsf u_m(t)|\le\frac2j\cdot2\sqrt3/t\cdot\mathrm{TV}(\varphi_m(y^j))\le C$ for $m=2^j$, with $\mathrm{TV}(\varphi_m(y^j))=2\|\varphi\|_\infty$.  This gives the two bounds for $\mathsf u_m$, and they are uniform in $t\ge\tau_0^+$.  The chain rule is justified first for finite Euler products $\prod_{m\le M}$ and then by passage to the limit $M\to\infty$ in the weighted space $\ell^1(n^{-\sigma'})$, $1<\sigma'<\sigma$, of cell masses, in which the coordinate derivatives $\partial_{\kappa_k}\tilde c_m=\mathsf u_m(\kappa_k)$ are bounded uniformly in $m$ and $M$ by the bounds just proved, so that $\sum_nn^{-2\sigma}\partial_{\kappa_k}|k_n|^2$ converges uniformly in the truncation.  It gives \CGeqref{eq:gradient}, with $\sum_m|\mathsf r_m||\mathsf u_m(t)|\le C\sum_mm^{-\sigma-1/2}+C\sum_j2^{-j\sigma}<\infty$ uniformly in $t$, which is the asserted absolute and uniform convergence on all of $t\ge\tau_0^+$, not only on compact sets.  If $\mathsf V_\sigma=0$, then $\mathsf K=1$ and $\mathsf r_m=\Re\langle\mathsf L_m,1\rangle_\sigma=0$, because $\mathsf L_m$ has no constant term.
\end{proof}

A list $\boldsymbol{\kappa}\in\mathscr O$ (Section~\CGref{sec:lists}) at which all coordinate derivatives of $\mathsf V_\sigma$ vanish is therefore a list all of whose ordinates are zeros of the single real-analytic function $\mathsf Q_{\boldsymbol{\kappa}}$.  The minimum of $\mathsf V_\sigma$ on $\mathscr O$ is attained (Corollary~\CGref{cor:existence-form}); if every minimizer satisfied the appropriate coordinate/KKT stationarity conditions and those conditions forced $\mathsf V_\sigma=0$, the minimum would be $0$ and (D) would follow.  This is not proved.  A count of zeros does not exclude critical points with $\mathsf V_\sigma>0$: $\mathsf Q_{\boldsymbol{\kappa}}(t)$ is a sum of terms of frequency $\log m$ with $m\gtrsim t/2\pi$, so it oscillates at the scale $1/\log t$, which is the scale of the gaps between the ordinates.


\section{Acceptance criteria and models in which the final step is a theorem}
\label{sec:models}

In the construction of this paper the continuous function $\zeta_\lambda$ is built from a list of ordinates and the discrete function $\tilde\zeta_{\Z,\lambda}$ from the cells, and the two are joined by $\tilde h_\lambda$.  This section states, without $\zeta$, the general criterion under which a list is accepted by both sides, and gives three models: a non-periodic list accepted by a discrete category with two generators, where the acceptance is proved by one unitary matrix; a finite Euler product over the primes, to which the constructions of Sections~\CGref{sec:clock}--\CGref{sec:lists} apply verbatim and for which the final step is a theorem; and the diagonal operator over all primes, which shows where the argument of the finite models stops.  The complex variable of the first two subsections is denoted by $\mathsf z$, and the half-plane $\Im\mathsf z>0$ plays the role of $\Re s>1$.

\subsection{The acceptance criterion: window and frequencies}
\label{subsec:acceptance}

Let $\mathfrak k>0$ and let $\boldsymbol{\kappa}=(\kappa_k)_{k\ge1}$ be a nondecreasing list of positive numbers with $|\kappa_k-k\pi/\mathfrak k|$ bounded.  The \emph{continuous side} is the canonical product
\begin{equation}\label{eq:model-product}
 \mathfrak C_{\boldsymbol{\kappa}}(\mathsf z):=\mathsf z\prod_{k\ge1}\Bigl(1-\frac{\mathsf z^2}{\kappa_k^2}\Bigr),
\end{equation}
an odd entire function of order one with real zeros; the oddness plays the role of the functional equation.  We assume that $\mathfrak C_{\boldsymbol{\kappa}}$ has exponential type $\mathfrak k$ and polynomial growth on the real axis, as $\sin(\mathfrak k\mathsf z)$ has for the list $\kappa_k=k\pi/\mathfrak k$.  The \emph{discrete side} is given by a closed, locally finite set $\mathscr F\subset(0,\infty)$ of frequencies ($\mathscr F\cap[0,R]$ is finite for every $R$; the set need not be uniformly discrete, and in the two-edge model below it is not) that is closed under addition; put $\mathscr F^{+}:=\mathscr F\cup\{0\}$.  The discrete category consists of the generalized Dirichlet series $\sum_{\mathfrak l\in\mathscr F^{+}}\mathfrak c_{\mathfrak l}\,e^{i\mathfrak l\mathsf z}$, absolutely convergent in some half-plane $\Im\mathsf z>\nu_{10}$ (the number $\nu_{10}$ is not prescribed and may be enlarged), and of their Euler forms $\exp\bigl(\sum_{\mathfrak l\in\mathscr F}\mathfrak b_{\mathfrak l}\,e^{i\mathfrak l\mathsf z}\bigr)$.  For $\Im\mathsf z>0$,
\begin{equation}\label{eq:model-logder}
 \frac{\mathfrak C_{\boldsymbol{\kappa}}'}{\mathfrak C_{\boldsymbol{\kappa}}}(\mathsf z)=\frac1{\mathsf z}+\sum_{k\ge1}\frac{2\mathsf z}{\mathsf z^2-\kappa_k^2}=-i\int_0^\infty e^{i\mathsf zu}\Bigl(1+2\sum_{k\ge1}\cos(\kappa_ku)\Bigr)du ,
\end{equation}
the bracket being a tempered distribution in $u$: it is the spectral density of the list.

\begin{theorem}[Acceptance criterion]\label{thm:acceptance}
For a list $\boldsymbol{\kappa}$ as above the following are equivalent:
\begin{enumerate}[label=\textup{(\roman*)}]
\item the spectral density is atomic on the frequencies: there are real $\mathfrak w_{\mathfrak l}$ with $\sum_{\mathfrak l\in\mathscr F}|\mathfrak w_{\mathfrak l}|e^{-\nu_{10}\mathfrak l}<\infty$ for some admissible $\nu_{10}$ and
\begin{equation}\label{eq:acceptance-density}
 1+2\sum_{k\ge1}\cos(\kappa_ku)=\sum_{\mathfrak l\in\mathscr F}\mathfrak w_{\mathfrak l}\,\delta(u-\mathfrak l)\qquad\text{in }\mathcal D'(0,\infty);
\end{equation}
\item the continuous function has an Euler form: for $\Im\mathsf z>\nu_{10}$, with a constant $\mathfrak c_0\ne0$,
\begin{equation}\label{eq:acceptance-euler}
 e^{i\mathfrak k\mathsf z}\,\mathfrak C_{\boldsymbol{\kappa}}(\mathsf z)=\mathfrak c_0\exp\Bigl(\sum_{\mathfrak l\in\mathscr F}\mathfrak b_{\mathfrak l}\,e^{i\mathfrak l\mathsf z}\Bigr),\qquad\mathfrak b_{\mathfrak l}=-\frac{\mathfrak w_{\mathfrak l}}{\mathfrak l};
\end{equation}
\item the continuous function is a finite sum in the discrete category:
\begin{equation}\label{eq:acceptance-window}
 e^{i\mathfrak k\mathsf z}\,\mathfrak C_{\boldsymbol{\kappa}}(\mathsf z)=\sum_{\mathfrak l\in\mathscr F^{+}\cap[0,2\mathfrak k]}\mathfrak c_{\mathfrak l}\,e^{i\mathfrak l\mathsf z},\qquad\mathfrak c_0\ne0 .
\end{equation}
\end{enumerate}
\end{theorem}

\begin{proof}
(i)$\Leftrightarrow$(ii): the two sides of \CGeqref{eq:acceptance-density} extend to distributions on $\R$ supported in $[0,\infty)$ (the left side as in \CGeqref{eq:model-logder}, the right side as the measure $\sum_{\mathfrak l}\mathfrak w_{\mathfrak l}\delta_{\mathfrak l}$), and (i) says that their difference vanishes on $(0,\infty)$; being supported at $\{0\}$ and of finite order, the difference is a finite combination of derivatives of $\delta_0$ \cite[Thm.~2.3.4]{Hormander1990}, whose transform $u\mapsto e^{i\mathsf zu}$ is a polynomial $\mathsf p(\mathsf z)$.  Hence, by \CGeqref{eq:model-logder}, (i) is the identity $\mathfrak C'_{\boldsymbol{\kappa}}/\mathfrak C_{\boldsymbol{\kappa}}=-i\sum_{\mathfrak l}\mathfrak w_{\mathfrak l}e^{i\mathfrak l\mathsf z}+\mathsf p(\mathsf z)$ in $\Im\mathsf z>\nu_{10}$.  Along $\mathsf z=iy$, $y\to+\infty$, the exponential sum tends to $0$ and $\mathfrak C'_{\boldsymbol{\kappa}}/\mathfrak C_{\boldsymbol{\kappa}}(iy)=-i\bigl(\frac1y+\sum_k\frac{2y}{y^2+\kappa_k^2}\bigr)\to-i\mathfrak k$ (type $\mathfrak k$); so the polynomial $\mathsf p$ has a finite limit along a ray, hence is constant, and the constant is $-i\mathfrak k$.  Integrating, $\log\mathfrak C_{\boldsymbol{\kappa}}=\text{const}-i\mathfrak k\mathsf z-\sum_{\mathfrak l}\frac{\mathfrak w_{\mathfrak l}}{\mathfrak l}e^{i\mathfrak l\mathsf z}$, and conversely.  (ii)$\Rightarrow$(iii): expanding the exponential, $\mathsf N(\mathsf z):=e^{i\mathfrak k\mathsf z}\mathfrak C_{\boldsymbol{\kappa}}(\mathsf z)=\sum_{\mathfrak l\in\mathscr F^{+}}\mathfrak c_{\mathfrak l}e^{i\mathfrak l\mathsf z}$, absolutely convergent in $\Im\mathsf z>\nu_{10}$.  On the other hand $\mathsf N$ satisfies $|\mathsf N(\mathsf z)|\le C(1+|\mathsf z|)^{N}e^{2\mathfrak k\max(0,-\Im\mathsf z)}$: for $\Im\mathsf z\ge0$ it is bounded by a polynomial, since $\mathfrak C_{\boldsymbol{\kappa}}$ has type $\mathfrak k$ and polynomial growth on the real axis and $|e^{i\mathfrak k\mathsf z}|=e^{-\mathfrak k\Im\mathsf z}$ absorbs the exponential growth in the upper half-plane, and for $\Im\mathsf z<0$ the factor $e^{i\mathfrak k\mathsf z}$ contributes $e^{\mathfrak k|\Im\mathsf z|}$ on top of the type $\mathfrak k$ of $\mathfrak C_{\boldsymbol{\kappa}}$. By the Paley--Wiener--Schwartz theorem \cite[Thm.~7.3.1]{Hormander1990}, used with the transform $\mathsf T\mapsto\langle\mathsf T,e^{i\mathsf zu}\rangle$ of \CGeqref{eq:model-logder} (for which a distribution supported in $[0,2\mathfrak k]$ has a transform bounded by $C(1+|\mathsf z|)^Ne^{2\mathfrak k\max(0,-\Im\mathsf z)}$, and conversely), an entire function with this bound is the transform of a distribution supported in $[0,2\mathfrak k]$, so the Fourier transform of the restriction of $\mathsf N$ to a horizontal line is supported in $[0,2\mathfrak k]$.  On the line $\Im\mathsf z=\nu_{10}+1$ this transform is $\sum_{\mathfrak l}\mathfrak c_{\mathfrak l}e^{-\mathfrak l(\nu_{10}+1)}\delta_{\mathfrak l}$, so $\mathfrak c_{\mathfrak l}=0$ for $\mathfrak l>2\mathfrak k$, and the sum is finite because $\mathscr F^{+}\cap[0,2\mathfrak k]$ is finite.  (iii)$\Rightarrow$(ii): write $\mathsf N=\mathfrak c_0(1+\mathsf w)$ with $\mathsf w(\mathsf z)=\sum_{0<\mathfrak l\le2\mathfrak k}(\mathfrak c_{\mathfrak l}/\mathfrak c_0)e^{i\mathfrak l\mathsf z}$; for $\Im\mathsf z\ge\nu'$ with $\sum_{0<\mathfrak l\le2\mathfrak k}|\mathfrak c_{\mathfrak l}/\mathfrak c_0|e^{-\mathfrak l\nu'}\le\frac12$ one has $|\mathsf w|\le\frac12$, so $\log(1+\mathsf w)=\sum_{r\ge1}(-1)^{r+1}\mathsf w^r/r$ converges absolutely, and expanding the powers gives an absolutely convergent series of exponentials whose frequencies are finite sums of elements of $\mathscr F$, hence elements of $\mathscr F$ (closure under addition); this proves (ii) in the half-plane $\Im\mathsf z>\nu'$, after increasing $\nu_{10}$ to $\nu'$ if necessary.
\end{proof}

In model identities of the form $2+2\sum_k\cos(\kappa_k u)$, the leading $2$ records a double zero at the origin; the normalized criterion \CGeqref{eq:acceptance-density} starts with $1$ because $\mathfrak C_{\boldsymbol{\kappa}}$ has a simple zero there.

The continuous side contributes a \emph{window}: by its type and its symmetry, the frequencies of $e^{i\mathfrak k\mathsf z}\mathfrak C_{\boldsymbol{\kappa}}$ lie in $[0,2\mathfrak k]$.  The discrete side contributes the \emph{admissible frequencies} $\mathscr F^{+}$.  A list is accepted by both sides if and only if its function lives in the intersection, and it is then determined by the finitely many coefficients $\mathfrak c_{\mathfrak l}$, $\mathfrak l\in\mathscr F^{+}\cap[0,\mathfrak k)$, the others being fixed by the oddness.  If $\mathscr F=2\mathfrak k\N$, the intersection is $\{0,2\mathfrak k\}$, the only function is a multiple of $\sin(\mathfrak k\mathsf z)$, and the only accepted list is the clock $\kappa_k=k\pi/\mathfrak k$.  Hamburger's theorem is the statement that for the window fixed by the factor $\Gamma(\frac s2)\pi^{-s/2}$ and the frequencies $\log n$ the space of accepted functions has dimension one.  The criterion says nothing about whether a given list, the zeta list in particular, belongs to the intersection; for $\zeta$ that membership is (D).

\subsection{A non-periodic accepted list: one unitary matrix}
\label{subsec:graph-model}

Let $\mathfrak p,\mathfrak q>0$ with $\mathfrak p/\mathfrak q\notin\mathbb Q$, $0<\mathfrak v<1$, $\mathfrak k=\mathfrak p+\mathfrak q$, and
\begin{equation}\label{eq:graph-matrices}
 \mathsf U:=\begin{pmatrix}\mathfrak v&\sqrt{1-\mathfrak v^2}\\ \sqrt{1-\mathfrak v^2}&-\mathfrak v\end{pmatrix},\qquad
 \mathsf X(\mathsf z):=\begin{pmatrix}e^{2i\mathfrak p\mathsf z}&0\\0&e^{2i\mathfrak q\mathsf z}\end{pmatrix},
\end{equation}
\begin{equation}\label{eq:graph-det}
 \mathfrak D(\mathsf z):=\det\bigl(I-\mathsf U\mathsf X(\mathsf z)\bigr)=1-\mathfrak v\,e^{2i\mathfrak p\mathsf z}+\mathfrak v\,e^{2i\mathfrak q\mathsf z}-e^{2i(\mathfrak p+\mathfrak q)\mathsf z}=-2i\,e^{i(\mathfrak p+\mathfrak q)\mathsf z}\,\mathfrak S(\mathsf z),
\end{equation}
\begin{equation}\label{eq:graph-sine}
 \mathfrak S(\mathsf z):=\sin\bigl((\mathfrak p+\mathfrak q)\mathsf z\bigr)+\mathfrak v\sin\bigl((\mathfrak p-\mathfrak q)\mathsf z\bigr).
\end{equation}

\begin{proposition}[The two-edge model]\label{prop:graph-model}
\begin{enumerate}[label=\textup{(\roman*)}]
\item All zeros of $\mathfrak S$ are real.  Each interval between consecutive points $(m+\frac12)\pi/\mathfrak k$, $m\in\Z$, contains a zero, and the number of zeros in $[0,T]$, with multiplicity, is $\mathfrak kT/\pi+O(1)$; hence the positive zeros form a list $\boldsymbol{\kappa}$ with $|\kappa_k-k\pi/\mathfrak k|$ bounded.  The set of zeros is not periodic.
\item $\mathfrak S(-\mathsf z)=-\mathfrak S(\mathsf z)$, and $\mathfrak D(\mathsf z)=-e^{2i\mathfrak k\mathsf z}\,\overline{\mathfrak D(\bar{\mathsf z})}$.
\item For $\Im\mathsf z>0$,
\begin{equation}\label{eq:graph-euler}
 \log\mathfrak D(\mathsf z)=-\sum_{r\ge1}\frac1r\operatorname{tr}\bigl(\mathsf U\mathsf X(\mathsf z)\bigr)^r=\sum_{\mathfrak l\in\mathscr F}\mathfrak b_{\mathfrak l}\,e^{i\mathfrak l\mathsf z},\qquad\mathscr F=\{2(m\mathfrak p+m'\mathfrak q):m,m'\ge0\}\setminus\{0\},
\end{equation}
absolutely convergent, and the spectral density of the zeros of $\mathfrak S$, counted with multiplicity, is atomic on $\mathscr F$ in the sense of \CGeqref{eq:acceptance-density}.
\item $\mathfrak D(\mathsf z)=-2i\,\mathfrak S'(0)\,e^{i\mathfrak k\mathsf z}\,\mathfrak C_{\boldsymbol{\kappa}}(\mathsf z)$ for the list $\boldsymbol{\kappa}$ of all positive zeros of $\mathfrak S$: the continuous and the discrete functions coincide up to the factor $e^{i\mathfrak k\mathsf z}$, and the list is accepted.
\end{enumerate}
\end{proposition}

\begin{proof}
(i) For $\Im\mathsf z>0$ the matrix $\mathsf U\mathsf X(\mathsf z)$ has norm at most $\max(e^{-2\mathfrak p\Im\mathsf z},e^{-2\mathfrak q\Im\mathsf z})<1$, because $\mathsf U$ is unitary; so $I-\mathsf U\mathsf X(\mathsf z)$ is invertible and $\mathfrak D(\mathsf z)\ne0$.  $\mathfrak S$ is odd, so it has no zeros in $\Im\mathsf z<0$ either.  At $\mathsf z_m=(m+\frac12)\pi/\mathfrak k$, $\sin(\mathfrak k\mathsf z_m)=(-1)^m$ and $|\mathfrak v\sin((\mathfrak p-\mathfrak q)\mathsf z_m)|<1$, so $\mathfrak S(\mathsf z_m)$ has the sign $(-1)^m$ and there is a zero in $(\mathsf z_{m-1},\mathsf z_m)$.  For the count, let $\mathsf M(t):=\mathsf U\mathsf X(t)$, unitary and real-analytic in $t\in\R$, with $\mathsf M'(t)=2i\,\mathsf M(t)\mathsf D$, $\mathsf D:=\operatorname{diag}(\mathfrak p,\mathfrak q)$ (diagonal matrices commute).  $\mathsf M(t)$ is never a scalar matrix, since $\mathsf U$ has nonzero off-diagonal entries; so its two eigenvalues $e^{i\theta_1(t)},e^{i\theta_2(t)}$, which can be labelled as continuous functions of $t$ (roots of a quadratic with continuous coefficients), coincide only on a discrete set (the discriminant $\operatorname{tr}^2-4\det=\mathfrak v^2(e^{2i\mathfrak pt}-e^{2i\mathfrak qt})^2+4e^{2i\mathfrak kt}$ is real-analytic and equals $4$ at $t=0$).  Where an eigenvalue $\mu=e^{i\theta_j}$ is simple it is analytic with an analytic unit eigenvector $v_j$, and differentiating $\mathsf Mv_j=\mu v_j$ and pairing with $v_j$ gives $\mu'=\langle\mathsf M'v_j,v_j\rangle$ (the term $\langle\mathsf Mv_j',v_j\rangle=\langle v_j',\mathsf M^*v_j\rangle=\mu\langle v_j',v_j\rangle$ cancels against $\mu\langle v_j',v_j\rangle$), i.e.
\begin{equation}\label{eq:graph-eigenphase}
 \theta_j'(t)=2\langle\mathsf Dv_j,v_j\rangle=2\bigl(\mathfrak p|v_{j,1}|^2+\mathfrak q|v_{j,2}|^2\bigr)\ \ge\ 2\min(\mathfrak p,\mathfrak q)>0 .
\end{equation}
Hence each continuous eigenphase is strictly increasing with $\theta_j(b)-\theta_j(a)\ge2\min(\mathfrak p,\mathfrak q)(b-a)$, and $\theta_1+\theta_2=2\mathfrak kt+\mathrm{const}$ because $\det\mathsf M(t)=-e^{2i\mathfrak kt}$.  For real $t$, $\mathfrak D(t)=(1-e^{i\theta_1(t)})(1-e^{i\theta_2(t)})$, and $1$ is never a double eigenvalue ($\mathsf M(t)=I$ is impossible), so every real zero of $\mathfrak D$ is a simple zero of exactly one factor, at which that eigenphase crosses $2\pi\Z$ with positive speed: the real zeros of $\mathfrak S$ are simple, and their number in $[0,T]$ is $\sum_j\#\{t\in[0,T]:\theta_j(t)\in2\pi\Z\}=\sum_j\bigl(\frac{\theta_j(T)-\theta_j(0)}{2\pi}+O(1)\bigr)=\frac{\mathfrak kT}\pi+O(1)$.  If the zero set had a period $y\ne0$, then $\mathfrak S(\mathsf z+y)$ and $\mathfrak S(\mathsf z)$ would be entire functions of order one with the same zero multiset, so by Hadamard's factorization $\mathfrak S(\mathsf z+y)=e^{a\mathsf z+b}\mathfrak S(\mathsf z)$ for constants $a,b$. As $\Im\mathsf z\to+\infty$, \CGeqref{eq:graph-sine} and $\mathfrak k>|\mathfrak p-\mathfrak q|$ give $\mathfrak S(\mathsf z)=-\frac1{2i}e^{-i\mathfrak k\mathsf z}(1+o(1))$ uniformly in $\Re\mathsf z$, so $e^{a\mathsf z+b}=e^{-i\mathfrak ky}(1+o(1))$ as $\Im\mathsf z\to+\infty$, which forces $a=0$. With $a=0$, comparing the coefficients of $e^{\pm i(\mathfrak p+\mathfrak q)\mathsf z}$ and of $e^{\pm i(\mathfrak p-\mathfrak q)\mathsf z}$ in $\mathfrak S(\mathsf z+y)=e^b\mathfrak S(\mathsf z)$ gives $e^{\pm i(\mathfrak p+\mathfrak q)y}=e^b=e^{\pm i(\mathfrak p-\mathfrak q)y}$, hence $(\mathfrak p+\mathfrak q)y\in\pi\Z$ and $(\mathfrak p-\mathfrak q)y\in\pi\Z$, so $2\mathfrak py\in\pi\Z$ and $2\mathfrak qy\in\pi\Z$, against $\mathfrak p/\mathfrak q\notin\mathbb Q$.  (ii) is immediate; the second identity follows from $\det\mathsf U=-1$.  (iii) The series is the logarithm of the determinant, convergent because the norm is less than $1$; each trace is a combination of exponentials with frequencies in $\mathscr F$.  The statement on the density is Theorem~\CGref{thm:acceptance}, (ii)$\Rightarrow$(i), applied to (iv).  (iv) $\mathfrak S$ is an odd entire function of exponential type $\mathfrak k$, bounded on the real axis, with real zeros; by Hadamard's factorization and the oddness, $\mathfrak S(\mathsf z)=\mathfrak S'(0)\,\mathfrak C_{\boldsymbol{\kappa}}(\mathsf z)$.
\end{proof}

In this model nothing has to be joined: the continuous function and the discrete function are the same determinant.  The Euler form comes from the expansion of $\det(I-\mathsf U\mathsf X)$ in closed paths, and the real zeros and the symmetry come from the unitarity of $\mathsf U$, through the contraction in $\Im\mathsf z>0$.  The acceptance criterion at this level is the existence of the common object $\mathsf U$.  In finite rank this is general: the zeros of $\mathsf P(e^{i\mathfrak l_1\mathsf z},\dots,e^{i\mathfrak l_r\mathsf z})$, for a polynomial $\mathsf P$ without zeros in the product of the open unit discs and in the product of their exteriors, are real and their spectral density is atomic \cite{KurasovSarnak2020}, and conversely the lists with unit masses whose spectral density is atomic with locally finite support are zero sets of exponential polynomials with real zeros \cite{OlevskiiUlanovskii2020}.

\subsection{A finite Euler product: the constructions of this paper with the final step proved}
\label{subsec:finite-euler-model}

Let
\begin{equation}\label{eq:finite-euler}
 \mathsf F(s):=\bigl(1-2^{-s}\bigr)\bigl(1-3^{-s}\bigr)=1-2^{-s}-3^{-s}+6^{-s}=\sum_{n\mid6}\mu(n)\,n^{-s}.
\end{equation}
It is a Dirichlet series over the integers with an Euler product over the primes $2$ and $3$, and $\mathsf F(-s)=6^{s}\,\mathsf F(s)$, with axis of symmetry $\Re s=0$.  Its zeros are $s=0$, double, and $s=\pm i\mathfrak y_k$, where
\begin{equation}\label{eq:finite-euler-zeros}
 \{\mathfrak y_k\}_{k\ge1}=\Bigl\{\frac{2\pi m}{\log2}\Bigr\}_{m\ge1}\cup\Bigl\{\frac{2\pi m}{\log3}\Bigr\}_{m\ge1}\quad\text{in increasing order},
\end{equation}
all simple, a non-periodic list, and, with $2\sinh(\frac{s\log p}2)=s\log p\prod_{m\ge1}\bigl(1+\frac{s^2\log^2p}{4\pi^2m^2}\bigr)$,
\begin{equation}\label{eq:finite-euler-hadamard}
 \mathsf F(s)=6^{-s/2}\,(\log2\log3)\,s^{2}\prod_{k\ge1}\Bigl(1+\frac{s^2}{\mathfrak y_k^2}\Bigr).
\end{equation}
The constructions of this paper apply as follows.  The count of the ordinates is $\lfloor\frac{T\log2}{2\pi}\rfloor+\lfloor\frac{T\log3}{2\pi}\rfloor=\frac{T\log6}{2\pi}-1+O(1)$ with a remainder of absolute value at most $1$; the midpoint clock is $\frac{\mathfrak x_k\log6}{2\pi}-1=k-\frac12$, that is, $\mathfrak x_k=(2k+1)\pi/\log6$, and $\mathfrak x_k-\mathfrak y_k$ is bounded.  For a list $\boldsymbol{\kappa}$ at bounded distance from the clock,
\begin{equation}\label{eq:finite-euler-list}
 \mathsf F_{\boldsymbol{\kappa}}(s):=6^{-s/2}\,(\log2\log3)\,s^{2}\prod_{k\ge1}\Bigl(1+\frac{s^2}{\kappa_k^2}\Bigr),\qquad\mathsf F_{\boldsymbol{\kappa}}(-s)=6^{s}\,\mathsf F_{\boldsymbol{\kappa}}(s),
\end{equation}
and the family is $\mathsf F_\lambda:=\mathsf F_{\boldsymbol{\kappa}(\lambda)}$ with $\kappa_k(\lambda)=(1-\lambda)\mathfrak y_k+\lambda\mathfrak x_k$, so that $\mathsf F_0=\mathsf F$ and, since $\prod_{k\ge0}\bigl(1+s^2/\mathfrak x_k^2\bigr)=\cosh(\frac{s\log6}2)$,
\begin{equation}\label{eq:finite-euler-clock}
 \mathsf F_1(s)=\frac{\log2\log3}2\cdot\frac{s^{2}\,(1+6^{-s})}{1+s^2\log^26/\pi^2}.
\end{equation}

\begin{theorem}[The final step in the model]\label{thm:finite-euler}
\begin{enumerate}[label=\textup{(\roman*)}]
\item \textup{(Poisson, prime by prime.)}  In $\mathcal D'(\R)$,
\begin{equation}\label{eq:finite-euler-poisson}
 2+2\sum_{k\ge1}\cos(\mathfrak y_ku)=\sum_{p\in\{2,3\}}\ \sum_{m\in\Z}\cos\frac{2\pi mu}{\log p}=\sum_{p\in\{2,3\}}\log p\sum_{m\in\Z}\delta(u-m\log p).
\end{equation}
The spectral density of the list $(\mathfrak y_k)$ is atomic on the points $m\log p$: the continuous function built from the list is the Dirichlet series with Euler product, $\mathsf F_0=\mathsf F$.
\item \textup{(Hamburger's theorem of the model.)}  Let $f$ be entire of finite order with $f(-s)=6^{s}f(s)$ and $f(0)=0$, and let $f(s)=\sum_{n\ge1}a_nn^{-s}$ converge absolutely in some half-plane $\Re s>\nu_{12}$.  Then $f=a_1\mathsf F$.
\item \textup{(The family.)}  $\mathsf F_\lambda$ is an absolutely convergent Dirichlet series in a half-plane if and only if $\lambda=0$.  In particular the clock function $\mathsf F_1$ is not a Dirichlet series.
\end{enumerate}
\end{theorem}

\begin{proof}
(i) is the Poisson summation formula $\sum_{m\in\Z}e^{2\pi imu/\log p}=\log p\sum_{m\in\Z}\delta(u-m\log p)$ for $p=2$ and $p=3$; the second statement is \CGeqref{eq:finite-euler-hadamard}.
(ii) Put $\sigma_1:=\nu_{12}+1$ and $\sigma_2:=-\sigma_1$.  On $\Re s\ge\sigma_1$, $f(s)=\sum_na_nn^{-s}$ converges absolutely and $|f|\le C:=\sum_n|a_n|n^{-\sigma_1}$; on $\Re s\le\sigma_2$, the functional equation gives $f(s)=6^{-s}f(-s)=\sum_na_n(n/6)^{s}$, again absolutely convergent, with $|f(s)|\le C\,6^{-\sigma}$.  By the Phragm\'en--Lindel\"of theorem $f$ is bounded in the closed strip $\mathbb S:=\{\sigma_2\le\Re s\le\sigma_1\}$: $f$ has finite order $\rho$, so $|f(s)|\le\exp(|s|^{\rho+1})\le\exp(e^{c|\Im s|})$ for every $c>0$ and $|s|$ large, which is the growth condition of the theorem in a strip of any finite width; let $\mathrm M:=\sup_{\mathbb S}|f|$.  \emph{Comparison of the two Dirichlet expansions.}  For real $\lambda$, $\sigma\in[\sigma_2,\sigma_1]$ and $T>0$ put $\mathsf A_{\lambda,T}(\sigma):=\frac1{2T}\int_{-T}^Tf(\sigma+it)e^{i\lambda t}\,dt$.  Since $f'(\sigma+it)=-i\,\partial_tf(\sigma+it)$, an integration by parts gives
\begin{equation}\label{eq:finite-euler-coefficients}
 \frac{d}{d\sigma}\mathsf A_{\lambda,T}(\sigma)=\frac{-i}{2T}\Bigl[f(\sigma+it)e^{i\lambda t}\Bigr]_{t=-T}^{t=T}-\lambda\,\mathsf A_{\lambda,T}(\sigma),
 \qquad\text{so}\qquad
 \Bigl|\frac{d}{d\sigma}\bigl(e^{\lambda\sigma}\mathsf A_{\lambda,T}(\sigma)\bigr)\Bigr|\le\frac{\mathrm M}{T}e^{\lambda\sigma},
\end{equation}
and integrating over $[\sigma_2,\sigma_1]$, $|e^{\lambda\sigma_1}\mathsf A_{\lambda,T}(\sigma_1)-e^{\lambda\sigma_2}\mathsf A_{\lambda,T}(\sigma_2)|\le\frac{\mathrm M}T\int_{\sigma_2}^{\sigma_1}e^{\lambda\sigma}d\sigma\to0$ as $T\to\infty$.  On the two boundary lines the limits $\mathsf A_\lambda(\sigma_i):=\lim_{T\to\infty}\mathsf A_{\lambda,T}(\sigma_i)$ exist and are computed termwise from the absolutely convergent expansions (dominated convergence; the mean of $e^{i\mu t}$ is $1$ for $\mu=0$ and $0$ otherwise): $\mathsf A_\lambda(\sigma_1)=a_nn^{-\sigma_1}$ if $\lambda=\log n$ and $0$ otherwise; $\mathsf A_\lambda(\sigma_2)=a_m(m/6)^{\sigma_2}$ if $\lambda=\log(6/m)$ and $0$ otherwise.  Hence
\begin{equation}\label{eq:finite-euler-spectrum}
 e^{\lambda\sigma_1}\mathsf A_\lambda(\sigma_1)=e^{\lambda\sigma_2}\mathsf A_\lambda(\sigma_2)\qquad(\lambda\in\R).
\end{equation}
Take $\lambda=\log n$ with $a_n\ne0$: the left side is $a_n\ne0$, so the right side is nonzero, which forces $\log n=\log(6/m)$ for some $m$ with $a_m\ne0$, i.e.\ $n\mid6$ and $m=6/n$, and then $e^{\lambda\sigma_2}\mathsf A_\lambda(\sigma_2)=a_m(m/6)^{\sigma_2}n^{\sigma_2}=a_m$, so $a_n=a_{6/n}$.  Thus $a_n=0$ unless $n\mid6$, $a_6=a_1$ and $a_3=a_2$, and $f(0)=2a_1+2a_2=0$ gives $a_2=-a_1$ and $f=a_1(1-2^{-s}-3^{-s}+6^{-s})=a_1\mathsf F$.  (This is the elementary case of Bohr's theorem that the Dirichlet exponents of a function almost periodic in a strip are the same on every line of the strip \cite{Bohr1926}; the earlier text displaced the coefficient formula to the lines $\Re s<0$ without justification, which is exactly what \CGeqref{eq:finite-euler-coefficients} supplies.)
(iii) $\mathsf F_\lambda$ is entire of order one, satisfies the functional equation and vanishes at $0$.  If it is a Dirichlet series, (ii) gives $\mathsf F_\lambda=a_1\mathsf F$, the lists of zeros coincide, and $\lambda(\mathfrak x_k-\mathfrak y_k)=0$ for every $k$, so $\lambda=0$.
\end{proof}

In the model every object of the programme exists and the analogue of (D) is proved.  What makes $\lambda=0$ close is a property of the list that the clock and the lists with $\lambda\ne0$ do not have,
\begin{equation}\label{eq:finite-euler-superposition}
 \{\mathfrak y_k\}=\frac{2\pi}{\log2}\,\N\ \cup\ \frac{2\pi}{\log3}\,\N ,
\end{equation}
a superposition of one clock for each prime; it passes the scale test of Remark~\CGref{rem:scale-test} and it is the hypothesis of the Poisson formula \CGeqref{eq:finite-euler-poisson}.  The constructions extend to $\prod_{p\le X}(1-p^{-s})=\sum_{n\mid X\#}\mu(n)n^{-s}$, with $X\#$ the product of the primes up to $X$, the functional equation with the factor $(-1)^{\pi(X)}(X\#)^{s}$, a zero of order $\pi(X)$ at the origin and the list $\bigcup_{p\le X}\frac{2\pi}{\log p}\N$: part (i) of Theorem~\CGref{thm:finite-euler} holds prime by prime, and the clock and the family are defined in the same way. The uniqueness statement (ii) does not extend verbatim: for $\pi(X)=3$ the coefficient conditions $a_{X\#/n}=-a_n$ leave four free coefficients, $f(0)=0$ is then automatic, and the two remaining conditions $f'(0)=f''(0)=0$ leave a two-dimensional space, so the analogue of (ii) is false as stated, and for larger $X$ the defect $2^{\pi(X)-1}-\pi(X)$ only grows; since (iii) was proved through (ii), it is proved here only for the two primes $2,3$.  As $X\to\infty$ these functions tend to $1/\zeta$ in $\Re s>1$ only; their axis of symmetry is $\Re s=0$, not $\Re s=\frac12$, and their lists do not converge to the ordinates of $\zeta$, which are not a superposition of one clock for each prime.

\subsection{The diagonal operator over all primes and the three abscissas}
\label{subsec:three-abscissas}

On $\mathscr L:=\ell^2(\{p\ \text{prime}\})$ let $\mathsf X_{\mathfrak P}(s):=\operatorname{diag}(p^{-s})$.  The closed paths of this model are the primes themselves, without scattering between them, and
\begin{equation}\label{eq:prime-det}
 \det\bigl(I-\mathsf X_{\mathfrak P}(s)\bigr)=\prod_p\bigl(1-p^{-s}\bigr)=\frac1{\zeta(s)}\qquad(\Re s>1).
\end{equation}
$\mathsf X_{\mathfrak P}(s)$ has norm $2^{-\sigma}<1$ for $\sigma>0$, it is of trace class if and only if $\sigma>1$, and Hilbert--Schmidt if and only if $\sigma>\frac12$.  With $\mathfrak P(s):=\sum_pp^{-s}$ for $\Re s>1$, the regularized determinant is the function
\begin{equation}\label{eq:prime-det2}
 D_2(s):=\det\nolimits_2\bigl(I-\mathsf X_{\mathfrak P}(s)\bigr)=\prod_p\bigl(1-p^{-s}\bigr)e^{p^{-s}}=\exp\Bigl(-\sum_{m\ge2}\frac1m\,\mathfrak P(ms)\Bigr)\qquad(\Re s>\tfrac12),
\end{equation}
holomorphic and without zeros in $\Re s>\frac12$ because $I-\mathsf X_{\mathfrak P}(s)$ is invertible there; the last expression converges for $\Re s>\frac12$ since $m\sigma>1$ for $m\ge2$. For $\Re s>1$, where $\mathfrak P$ and the Euler product converge, $D_2(s)=e^{\mathfrak P(s)}/\zeta(s)$; in the strip $\frac12<\Re s\le1$ the primary object is $D_2$, and $e^{\mathfrak P}/\zeta$ is defined there only through it, since $\mathfrak P$ itself is not given by its series.

\begin{proposition}[No common object of finite rank in the window of $\zeta$]\label{prop:no-finite-rank}
Let $X\ge2$.  There is no function $f\not\equiv0$ that satisfies (H0) and the functional equation \CGeqref{eq:Z1-FE} and is an absolutely convergent Dirichlet series in $\Re s>1$ supported on the integers all of whose prime factors are at most $X$.
\end{proposition}

\begin{proof}
By Hamburger's theorem $f=C_H\zeta$, and the coefficients of $\zeta$ do not vanish at the integers with a prime factor larger than $X$; so $C_H=0$.
\end{proof}

In the models of Subsections~\CGref{subsec:graph-model} and~\CGref{subsec:finite-euler-model} three boundaries coincide: the boundary of the domain where the matrix is a strict contraction, the boundary of the domain of convergence of the Euler product, and the axis of symmetry.  The real zeros come from this coincidence.  For \CGeqref{eq:prime-det} they are different:
\begin{equation}\label{eq:three-abscissas}
 \text{contraction: }\sigma>0;\qquad\text{convergence of the determinant: }\sigma>1;\qquad\text{axis of symmetry: }\sigma=\tfrac12 .
\end{equation}
The contraction, together with the trace-class determinant identity \CGeqref{eq:prime-det} valid for $\sigma>1$, proves that $\zeta$ has no zeros in $\sigma>1$; the contraction together with the Hilbert--Schmidt property proves that $D_2=\det_2(I-\mathsf X_{\mathfrak P})$ is holomorphic without zeros in $\sigma>\frac12$, and $D_2=e^{\mathfrak P}/\zeta$ there only in the sense that $D_2$ is the continuation of that quotient from $\sigma>1$. Neither statement reaches the axis, and in the strip $\frac12<\sigma\le1$ the regularized determinant moves the question to the continuation of $\mathfrak P$ (equivalently of $\zeta$).  In the language of this paper, $\log\tilde\zeta_{\Z,\lambda}=\sum_n\tilde c_n(\lambda)\mathsf L_n$ is the determinant \CGeqref{eq:prime-det} with real multiplicities $\tilde c_n(\lambda)$ at all integers, and the strip $\frac12<\sigma\le1$ is where the mean square of the cell masses (Theorem~\CGref{thm:mean-square-cells}) and the continuation of the joining factor (Theorem~\CGref{thm:atomized-continuation}) are used.  By Proposition~\CGref{prop:no-finite-rank} the common object of the window of $\zeta$ has no finite-rank approximants in the discrete category; a unitary object of infinite rank from which both $\zeta_{Z1}$ and $\tilde\zeta_{\Z,0}$ descend, as $\mathsf U$ does in Proposition~\CGref{prop:graph-model}, is not constructed in this paper.  A model without $\zeta$ with the same separation of the three boundaries is a graph with infinitely many edges of lengths $\mathfrak l_j$ such that $\sum_je^{-\mathfrak l_jy}$ diverges for $y\le y_0$, with $y_0>0$.

%======================================================================
\section{Route B for finite sections: compactness, prescribed masses, certificates, and a Kadec--Avdonin criterion with blocks of variable length}
\label{sec:routeB-finite}

This section incorporates, with the notation of this paper, the companion note \emph{Route B for the midpoint RvM system} (version 4; the file \texttt{supplementary/lemas\_B\_v4/} of the repository). Its subsections are self-contained and use the preceding sections only through the results quoted by label. The numerical part (Section~\CGref{subsec:rb-numerics}) is a computation and is not part of any proof.

\subsection{Introduction to the section}\label{subsec:rb-intro}

This paper constructs a synthetic zeta system on the midpoint clock $\tau_k$, defined by $\NR(\tau_k)=k-\frac12$, and reduces the remaining step to condition (D). Route B approaches (D) without constructing the zeta ordinates. One solves finite sections of the rung system \CGeqref{eq:rb-rung}, which imposes simultaneously a real critical spectrum and atomic arithmetic on $\{\log n\}$. If the residuals of these sections tend to zero in a fixed box, compactness yields an exact solution. The uniqueness theorem \CGref{thm:free-masses} then identifies that solution as the list of zeta ordinates with the von Mangoldt masses. In this scheme the ordinates $\gamma_k$ are the output of an inverse problem, not an input. Section~\CGref{subsec:rb-routeB} makes this precise and isolates the two statements that remain open: existence of sections with vanishing residual, and uniqueness in the box without the third condition of $\mathscr O$. Section~\CGref{subsec:rb-prescribed} shows that prescribing the von Mangoldt masses removes the second: the Laplace transform of the residual is identified (Lemma~\CGref{lem:rb-laplace}), and uniqueness follows from Hadamard's theorem and the functional equation, so that route B$'$ is equivalent to the Riemann hypothesis. Section~\CGref{subsec:rb-staged} gives the existence part a certificate form: two computable quantities per stage, a right inverse of the Jacobian and the residual of the new tests, imply an exact solution of the stage. Section~\CGref{subsec:rb-gonek} shows that Gonek's symmetrised Euler products give explicit lists in a tube around the clock. Section~\CGref{subsec:rb-numerics} reports the computations that measure these quantities.

Section~\CGref{subsec:rb-apriori} proves an a priori bound for quasi-solutions. Section~\CGref{subsec:rb-riesz} proves a Kadec--Avdonin criterion with blocks of variable length, which is the tool needed for a Fredholm/stability formulation of the linearised problem. Section~\CGref{subsec:rb-application} applies that criterion to lists near the clock and proves that, for the zeta list, the uniform version fails.

References of the form Theorem~\CGref{thm:free-masses} and the other results of the preceding sections are used as stated there. Except in Section~\CGref{subsec:rb-obstruction}, where the zeta list is the object of study, nothing below uses the ordinates of $\zeta$.

%======================================================================
\subsection{Setting}\label{subsec:rb-setting}

Throughout, $q(s):=s(s-1)$ and $G(s):=s(s-1)\pi^{-s/2}\Gamma(s/2)$. The smooth clock count is
\begin{equation}\label{eq:rb-NRVM}
 \NR(t):=\frac t{2\pi}\log\frac t{2\pi e}+\frac78,\qquad
 \NR'(t)=\frac1{2\pi}\log\frac t{2\pi}\qquad(t>0).
\end{equation}
It is strictly increasing on $[2\pi,\infty)$, with $\NR(2\pi)=-\frac18$, and the clock points are the solutions $\tau_k>2\pi$ of $\NR(\tau_k)=k-\frac12$, $k\ge1$ ($\tau_1=14.5213\ldots$). We write $n_{\RVM}(t):=\#\{k\ge1:\tau_k\le t\}$. Since $\NR$ is increasing,
\begin{equation}\label{eq:rb-Ntau-floor}
 n_{\RVM}(t)=\Bigl\lfloor\NR(t)+\tfrac12\Bigr\rfloor,\qquad
 \bigl|n_{\RVM}(t)-\NR(t)\bigr|\le\tfrac12\qquad(t\ge2\pi).
\end{equation}

\begin{definition}[Admissible lists]\label{def:rb-lists}
For $C>0$, a list $\boldsymbol\kappa=(\kappa_k)_{k\ge1}$ is \emph{$C$-admissible} if it is real and nondecreasing, $\kappa_1\ge2\pi$, and $|\kappa_k-\tau_k|\le C$ for all $k$. Put $N_{\boldsymbol\kappa}(t):=\#\{k:\kappa_k\le t\}$ (with multiplicity) and
\begin{equation}\label{eq:rb-Pk}
 P_{\boldsymbol\kappa}(s):=\prod_{k\ge1}\Bigl(1+\frac{q(s)}{\frac14+\kappa_k^2}\Bigr),\qquad
 \zeta_{\boldsymbol\kappa}:=\frac{P_{\boldsymbol\kappa}}{G}.
\end{equation}
The clock list $\boldsymbol\tau=(\tau_k)$ gives $\zeta_{\boldsymbol\tau}=\zeta_{\RVM}$.
\end{definition}

$P_{\boldsymbol\kappa}$ is entire and real on $\mathbb R$, and $P_{\boldsymbol\kappa}(1-s)=P_{\boldsymbol\kappa}(s)$. Its zeros are exactly $\frac12\pm i\kappa_k$, with multiplicity, and it is positive on $[-1,2]$, because $q([-1,2])=[-\frac14,2]$ and $\kappa_k^2\ge4\pi^2$. Writing $s=\frac12+w$, one has $\frac14+\kappa^2+q(s)=\kappa^2+w^2$.

The \emph{rung system} \CGref{eq:rung-system} for a list $\boldsymbol\kappa$ and real masses $\mathsf a=(\mathsf a_n)_{n\ge2}$ is
\begin{equation}\label{eq:rb-rung}
 \sum_{k\ge1}\cos(\kappa_ku)+\sum_{n\ge2}\mathsf a_n\,\delta(u-\log n)=\mathsf k(u)\qquad\text{in }\mathcal D'(0,\infty),
\end{equation}
with the fixed distribution $\mathsf k$ of \CGref{eq:transport-defs}. Its residual is the distribution $\mathfrak R(\boldsymbol\kappa,\mathsf a)$ defined, for $\varphi\in C_c^\infty(0,\infty)$, by
\begin{equation}\label{eq:rb-residual}
 \langle\mathfrak R(\boldsymbol\kappa,\mathsf a),\varphi\rangle:=\sum_{k\ge1}\mathsf c_\varphi(\kappa_k)+\sum_{n\ge2}\mathsf a_n\varphi(\log n)-\langle\mathsf k,\varphi\rangle,\qquad
 \mathsf c_\varphi(t):=\int_0^\infty\varphi(u)\cos(tu)\,du .
\end{equation}
Two integrations by parts give $|\mathsf c_\varphi(t)|\le\|\varphi''\|_{L^1}\,t^{-2}$. Since $\kappa_k\ge\tau_k-C$ and $\tau_k\gg k/\log(k+1)$, the series in \CGeqref{eq:rb-residual} converges absolutely. The mass sum is finite, because $\operatorname{supp}\varphi\subset[u_1,u^*]$ with $0<u_1<u^*$ leaves only $n\le e^{u^*}$.

%======================================================================
\subsection{Route B: compactness, selection and the canonical box}\label{subsec:rb-routeB}

For $C,\mathbb M>0$ let $\mathbb K(C,\mathbb M)$ be the set of pairs $(\boldsymbol\kappa,\mathsf a)$ with $\boldsymbol\kappa$ $C$-admissible and $|\mathsf a_n|\le \mathbb M$ for all $n$. A \emph{finite section} is an element of $\mathbb K(C,\mathbb M)$ with $\kappa_k=\tau_k$ for $k>k^*$ and $\mathsf a_n=0$ for $n>n^*$, for some $k^*,n^*$.

Fix finite sets $\mathscr T_1\subset\mathscr T_2\subset\cdots\subset C_c^\infty(0,\infty)$ whose union $\mathscr T_\infty$ is dense in $C_c^\infty(0,\infty)$. Such sets exist because $C_c^\infty(0,\infty)$ is separable for its usual LF topology: it is the countable strict inductive limit of the separable Fr\'echet spaces $C^\infty_{[1/j,j]}$, and a countable union of countable dense subsets of these is dense. For $(\boldsymbol\kappa,\mathsf a)\in\mathbb K(C,\mathbb M)$ put
\begin{equation}\label{eq:rb-rho}
 \mathsf s_m(\boldsymbol\kappa,\mathsf a):=\max_{\varphi\in\mathscr T_m}\bigl|\langle\mathfrak R(\boldsymbol\kappa,\mathsf a),\varphi\rangle\bigr| .
\end{equation}

\begin{definition}[Statement B]\label{def:rb-B}
\emph{Statement B$(C,\mathbb M)$} says that there are finite sections $(\boldsymbol\kappa^{(m)},\mathsf a^{(m)})\in\mathbb K(C,\mathbb M)$ with $\mathsf s_m(\boldsymbol\kappa^{(m)},\mathsf a^{(m)})\to0$.
\end{definition}

\begin{lemma}[Canonical box of the exact solution]\label{lem:rb-canonical-box}
Let $\mathbb M_*:=\max_{n\ge2}\Lambda(n)/(2\sqrt n)$. Then $\mathbb M_*=\log7/(2\sqrt7)=0.367742\ldots<1/e$. If condition (D) holds, the pair $(\boldsymbol\gamma,\mathsf a^\zeta)$, with $\mathsf a^\zeta_n:=\Lambda(n)/(2\sqrt n)$, solves \CGeqref{eq:rb-rung} and lies in $\mathbb K(C_\eta,\mathbb M_*)$, where $C_\eta$ is the constant of \CGref{lem:bounded-displacement}.
\end{lemma}

\begin{proof}
$x\mapsto\log x/(2\sqrt x)$ increases on $[1,e^2]$ and decreases on $[e^2,\infty)$, with $e^2=7.389\ldots$. Hence among primes its maximum is at $p=7$ or $p=11$; the values are $0.367742\ldots$ and $0.361500\ldots$. A prime power $p^j$ with $j\ge2$ gives $\log p/(2p^{j/2})<\log p/(2\sqrt p)\le \mathbb M_*$. The bound by $1/e$ is the continuous maximum, attained at $x=e^2$. Under (D), $\zeta_{Z1}=\zeta$, so all zeros of $\zeta$ are $\frac12\pm i\gamma_k$. The rung system with list $\boldsymbol\gamma$ and masses $\mathsf a^\zeta$ is then the distributional explicit formula \CGref{thm:distributional-EF}. Finally, $|\gamma_k-\tau_k|\le C_\eta$ by \CGref{lem:bounded-displacement}, and $\gamma_1=14.13\ldots\ge2\pi$.
\end{proof}

\begin{corollary}[(D) implies B]\label{cor:rb-D-implies-B}
If (D) holds, then B$(C,\mathbb M)$ holds for every $C\ge C_\eta$ and $\mathbb M\ge \mathbb M_*$.
\end{corollary}

\begin{proof}
Let $(\boldsymbol\kappa^{(k^*,n^*)},\mathsf a^{(k^*,n^*)})$ be the truncation of $(\boldsymbol\gamma,\mathsf a^\zeta)$. Its list is the nondecreasing rearrangement of the multiset $\{\gamma_1,\dots,\gamma_{k^*}\}\cup\{\tau_{k^*+1},\tau_{k^*+2},\dots\}$, and its masses are $\mathsf a^\zeta_n$ for $n\le n^*$ and $0$ beyond. By Lemma~\CGref{lem:rb-matching} (Section~\CGref{subsec:rb-staged}) the rearranged list satisfies $|\kappa_k-\tau_k|\le C_\eta$, coincides with $(\tau_k)$ beyond some index, and $\mathfrak R$ depends only on the multiset. So this is a finite section in $\mathbb K(C_\eta,\mathbb M_*)$. For fixed $\varphi$, $\langle\mathfrak R(\boldsymbol\kappa^{(k^*,n^*)},\mathsf a^{(k^*,n^*)}),\varphi\rangle-\langle\mathfrak R(\boldsymbol\gamma,\mathsf a^\zeta),\varphi\rangle=\sum_{k>k^*}(\mathsf c_\varphi(\tau_k)-\mathsf c_\varphi(\gamma_k))$ as soon as $e^{u^*}\le n^*$. This tends to $0$ as $k^*\to\infty$. Since $\mathscr T_m$ is finite and $\mathfrak R(\boldsymbol\gamma,\mathsf a^\zeta)=0$ by Lemma~\CGref{lem:rb-canonical-box}, $k^*_m,n^*_m$ can be chosen with $\mathsf s_m\le1/m$.
\end{proof}

\begin{proposition}[Compactness and passage to the limit]\label{prop:rb-compactness}
Let $(\boldsymbol\kappa^{(m)},\mathsf a^{(m)})\in\mathbb K(C,\mathbb M)$ with $\mathsf s_m(\boldsymbol\kappa^{(m)},\mathsf a^{(m)})\to0$. Then a subsequence converges coordinatewise to some $(\boldsymbol\kappa,\mathsf a)\in\mathbb K(C,\mathbb M)$, and every such limit satisfies $\mathfrak R(\boldsymbol\kappa,\mathsf a)=0$ in $\mathcal D'(0,\infty)$.
\end{proposition}

\begin{proof}
Each $\kappa^{(m)}_k$ lies in the compact interval $[\tau_k-C,\tau_k+C]$ and each $\mathsf a^{(m)}_n$ in $[-\mathbb M,\mathbb M]$. A diagonal extraction over the countably many coordinates gives a coordinatewise convergent subsequence. The limit is nondecreasing, satisfies $\kappa_1\ge2\pi$, $|\kappa_k-\tau_k|\le C$ and $|\mathsf a_n|\le \mathbb M$, and therefore lies in $\mathbb K(C,\mathbb M)$.

\emph{Continuity of each pairing.} Fix $\varphi$ with support in $[u_1,u^*]$. The mass part of \CGeqref{eq:rb-residual} involves only $n\le e^{u^*}$ and converges. For the spectral part, $|\mathsf c_\varphi(\kappa^{(m)}_k)|\le\|\varphi''\|_1(\tau_k-C)^{-2}$ when $\tau_k>2C$, and $|\mathsf c_\varphi|\le\|\varphi\|_1$ for the finitely many other $k$. The majorant is summable, so dominated convergence for series gives $\sum_k\mathsf c_\varphi(\kappa_k^{(m)})\to\sum_k\mathsf c_\varphi(\kappa_k)$. Hence $\langle\mathfrak R(\boldsymbol\kappa^{(m)},\mathsf a^{(m)}),\varphi\rangle\to\langle\mathfrak R(\boldsymbol\kappa,\mathsf a),\varphi\rangle$ along the subsequence.

\emph{Vanishing.} If $\varphi\in\mathscr T_\infty$, then $\varphi\in\mathscr T_m$ for all large $m$, so $|\langle\mathfrak R(\boldsymbol\kappa^{(m)},\mathsf a^{(m)}),\varphi\rangle|\le\mathsf s_m\to0$ and $\langle\mathfrak R(\boldsymbol\kappa,\mathsf a),\varphi\rangle=0$. The limit $\mathfrak R(\boldsymbol\kappa,\mathsf a)$ is a single distribution, hence continuous on $C_c^\infty(0,\infty)$. It vanishes on the dense set $\mathscr T_\infty$, so it vanishes identically. No uniform bound on the residuals is needed for this step.
\end{proof}

\begin{hypothesis}[Uniqueness in the box]\label{hyp:rb-U}
\emph{U$(C,\mathbb M)$}: the only $(\boldsymbol\kappa,\mathsf a)\in\mathbb K(C,\mathbb M)$ with $\mathfrak R(\boldsymbol\kappa,\mathsf a)=0$ is $(\boldsymbol\gamma,\mathsf a^\zeta)$.
\end{hypothesis}

\begin{remark}[What is proved and what is open in U]\label{rem:rb-U-status}
\CGref{thm:free-masses} proves uniqueness for lists in $\mathscr O$. Besides the tube and the ordering, $\mathscr O$ (\CGref{eq:list-space}) carries a third condition, $|\int(N_{\boldsymbol\kappa}-\NR)|\le C\log \mathrm X$, and a coordinatewise limit of sections does not inherit it. Hypothesis~\CGref{hyp:rb-U} is therefore the statement that \CGref{thm:free-masses} holds for exact solutions in $\mathbb K(C,\mathbb M)$ without the third condition. That condition is used to bound $\zeta_{\boldsymbol\kappa}$ on half-planes. For an exact solution, Lemma~\CGref{lem:rb-factor} below with $\widehat{\mathsf R}\equiv0$ gives that bound directly on $\Re s\ge2$. Whether this suffices for the rest of the proof of \CGref{thm:free-masses} has to be checked line by line against that proof. This is not done here, and Hypothesis~\CGref{hyp:rb-U} is not claimed.
\end{remark}

\begin{corollary}[Selection of the zeta list]\label{cor:rb-selection}
Assume B$(C,\mathbb M)$ and U$(C,\mathbb M)$, and let $(\boldsymbol\kappa^{(m)},\mathsf a^{(m)})$ be any sequence in $\mathbb K(C,\mathbb M)$ with $\mathsf s_m\to0$. Then the whole sequence converges coordinatewise to $(\boldsymbol\gamma,\mathsf a^\zeta)$. In particular $\mathfrak R(\boldsymbol\gamma,\mathsf a^\zeta)=0$, that is, (D) holds.
\end{corollary}

\begin{proof}
By Proposition~\CGref{prop:rb-compactness}, every subsequence has a further subsequence converging to a zero of $\mathfrak R$ in $\mathbb K(C,\mathbb M)$. By U$(C,\mathbb M)$ that zero is $(\boldsymbol\gamma,\mathsf a^\zeta)$, so the whole sequence converges to it. Finally we show that $\mathfrak R(\boldsymbol\gamma,\mathsf a^\zeta)=0$ implies (D). By Lemma~\CGref{lem:rb-laplace} (Section~\CGref{subsec:rb-prescribed}), $\widehat{\mathsf R}$ of \CGeqref{eq:rb-Rhat} is the Laplace transform of a distribution of finite order on $[0,\infty)$ whose restriction to $(0,\infty)$ is $\mathfrak R$. If $\mathfrak R=0$, that distribution is supported at $\{0\}$, so $\widehat{\mathsf R}$ is a polynomial $p(s)$. Since $\sum_n\mathsf a^\zeta_nn^{\frac12-s}=-\frac12\zeta'/\zeta(s)$, this reads $P_{\boldsymbol\gamma}'/P_{\boldsymbol\gamma}-(G\zeta)'/(G\zeta)=2p$, so $P_{\boldsymbol\gamma}=c\,G\zeta\,e^{Q}$ with $Q'=2p$. Here $G\zeta=2\xi$ and $P_{\boldsymbol\gamma}$ are entire of order $1$, so $e^Q=P_{\boldsymbol\gamma}/(2c\,\xi)$ is entire, zero-free and of order at most $1$, and $\deg Q\le1$. Both $P_{\boldsymbol\gamma}$ and $\xi$ are invariant under $s\mapsto1-s$, so $e^{Q(1-s)}=e^{Q(s)}$, which forces $Q$ to be constant. Hence $\zeta_{\boldsymbol\gamma}=c'\zeta$. Both functions have residue $1$ at $s=1$ ($P_{\boldsymbol\gamma}(1)=1$ and $G(s)\sim s-1$), so $c'=1$. This is (D), because $\zeta_{\boldsymbol\gamma}=\zeta_{Z1}$.
\end{proof}

\begin{remark}[Logical status of route B]\label{rem:rb-B-status}
Corollaries~\CGref{cor:rb-D-implies-B} and~\CGref{cor:rb-selection} give (D) $\Rightarrow$ B$(C_\eta,\mathbb M_*)$ and [B$(C,\mathbb M)$ and U$(C,\mathbb M)$] $\Rightarrow$ (D). Route B therefore consists of two independent pieces. The first is existence, B, which is not weaker than (D). The second is uniqueness without the third condition of $\mathscr O$, Hypothesis~\CGref{hyp:rb-U}, which does not involve $\zeta$. Compactness itself is not an issue in the box: Proposition~\CGref{prop:rb-compactness} uses only the imposed bounds. In this scheme the ordinates $\gamma_k$ never enter the construction; they appear only as the unique admissible zero of $\mathfrak R$. Lemma~\CGref{lem:rb-canonical-box} does \emph{not} show that near-solutions of sections lie in $\mathbb K(C_\eta,\mathbb M_*)$: it locates only the exact solution. Restricting the sections to that box is a choice, justified a posteriori by Corollary~\CGref{cor:rb-D-implies-B}.
\end{remark}

\subsubsection{Route B with prescribed von Mangoldt masses}\label{subsec:rb-prescribed}

The masses $\mathsf a^\zeta_n=\Lambda(n)/(2\sqrt n)$ are known and do not depend on (D) or on the ordinates of $\zeta$. Prescribing them leaves the ordinates as the only unknowns, and the uniqueness question of Hypothesis~\CGref{hyp:rb-U} disappears. The tool is the identification of the Laplace transform of the residual. In version 3 of the companion note this was a convention; here it is a lemma.

\emph{Laplace transforms of distributions.} Let $\mathscr T\in\mathcal D'(\mathbb R)$ have support in $[0,\infty)$ and suppose that $e^{-au}\mathscr T$ is a tempered distribution for some $a\in\mathbb R$. For $\Re s>a+\frac12$ put $z:=s-\frac12$ and
\begin{equation}\label{eq:rb-laplace-def}
 \mathcal L\mathscr T(s):=\bigl\langle\mathscr T,\ \mathsf j(u)\,e^{-zu}\bigr\rangle,
\end{equation}
where $\mathsf j\in C^\infty(\mathbb R)$ equals $1$ on $[-1,\infty)$ and $0$ on $(-\infty,-2]$; the value does not depend on $\mathsf j$, and $\mathcal L\mathscr T$ is holomorphic there \cite[Ch.~8]{Zemanian1965}. The shift by $\frac12$ in $z$ is the normalisation of \CGeqref{eq:rb-Rhat}. Since $\mathsf j'$ vanishes on the support of $\mathscr T$, $\mathcal L(\partial_u\mathscr T)(s)=z\,\mathcal L\mathscr T(s)$; also $\mathcal L\delta_0=1$, $\mathcal L\delta_0'=z$, and for a locally integrable $g$ on $[0,\infty)$ with $|g(u)|\le ce^{au}$, $\mathcal L(\mathfrak g\mathbf 1_+)(s)=\int_0^\infty e^{-zu}g(u)\,du$. We write $\mathbf 1_+:=\mathbf 1_{[0,\infty)}$ and $\mathrm{fp}\,\frac{\mathbf 1_+}u$ for the distribution $\varphi\mapsto\int_0^1\frac{\varphi(u)-\varphi(0)}u\,du+\int_1^\infty\frac{\varphi(u)}u\,du$, of order $1$, supported in $[0,\infty)$, with restriction $1/u$ to $(0,\infty)$.

\begin{lemma}[The Archimedean distribution]\label{lem:rb-archimedean}
Let $\mathfrak g(u):=\frac{e^{-u/2}}{e^{2u}-1}-\frac1{2u}$ for $u>0$; $\mathfrak g$ is bounded, $\mathfrak g(0^+)=-\frac34$, and $\mathfrak g(u)=-\frac1{2u}+O(e^{-5u/2})$ as $u\to\infty$, in particular bounded. Put $\mathrm c_{\mathscr K}:=-\frac14\log(2\pi e^{\gamma_E})$, $\gamma_E$ Euler's constant, and
\begin{equation}\label{eq:rb-K-dist}
 \mathscr K:=\tfrac12e^{u/2}\mathbf 1_+-\tfrac12\,\mathfrak g\,\mathbf 1_+-\tfrac14\,\mathrm{fp}\,\frac{\mathbf 1_+}u+\mathrm c_{\mathscr K}\,\delta_0 .
\end{equation}
Then $\mathscr K$ is a distribution of order $1$ supported in $[0,\infty)$, its restriction to $(0,\infty)$ is the function $\mathsf k$ of \CGref{eq:transport-defs}, and
\begin{equation}\label{eq:rb-LK}
 \mathcal L\mathscr K(s)=\tfrac12\,\frac{G'}G(s)\qquad(\Re s>1).
\end{equation}
\end{lemma}

\begin{proof}
On $(0,\infty)$, $\frac12e^{u/2}-\frac12\mathfrak g(u)-\frac1{4u}=\frac12\bigl(e^{u/2}-\frac{e^{-u/2}}{e^{2u}-1}\bigr)=\mathsf k(u)$. The expansion $\frac{e^{-u/2}}{e^{2u}-1}=\frac1{2u}-\frac34+O(u)$ gives the behaviour of $\mathfrak g$ at $0$. With $z=s-\frac12$ and $\Re s>1$: $\mathcal L(\frac12e^{u/2}\mathbf 1_+)=\frac1{2(z-\frac12)}=\frac1{2(s-1)}$. Gauss's integral $\frac{\Gamma'}{\Gamma}(x)=\int_0^\infty\bigl(\frac{e^{-t}}t-\frac{e^{-xt}}{1-e^{-t}}\bigr)dt$ ($\Re x>0$) for the logarithmic derivative of $\Gamma$ \cite[Ch.~XII]{WhittakerWatson1927}, with $t=2u$ and $x=\frac s2$, reads
\begin{equation}\label{eq:rb-aux-12}
 \tfrac12\frac{\Gamma'}{\Gamma}\bigl(\tfrac s2\bigr)=\int_0^\infty\Bigl(\frac{e^{-2u}}{2u}-\frac{e^{-su}}{1-e^{-2u}}\Bigr)du
 =\int_0^\infty\Bigl(\frac{e^{-2u}}{2u}-e^{-zu}\frac{e^{-u/2}}{e^{2u}-1}\Bigr)du-\frac1s,
\end{equation}
because $\frac{e^{-su}}{1-e^{-2u}}=e^{-zu}\bigl(\frac{e^{-u/2}}{e^{2u}-1}+e^{-u/2}\bigr)$ and $\int_0^\infty e^{-zu}e^{-u/2}du=\frac1s$. Writing $\frac{e^{-u/2}}{e^{2u}-1}=\mathfrak g(u)+\frac1{2u}$ and using Frullani's integral $\int_0^\infty\frac{e^{-2u}-e^{-zu}}u\,du=\log\frac z2$ (both sides vanish at $z=2$ and have derivative $\int_0^\infty e^{-zu}du=1/z$ in $z$, for $\Re z>0$), the integral in \CGeqref{eq:rb-aux-12} equals $-\mathcal L(\mathfrak g\mathbf 1_+)(s)+\frac12\log\frac z2$. Hence
\begin{equation}\label{eq:rb-aux-13}
 \tfrac12\frac{G'}G(s)=\frac1{2s}+\frac1{2(s-1)}-\frac14\log\pi+\frac14\frac{\Gamma'}{\Gamma}\bigl(\tfrac s2\bigr)
 =\mathcal L\bigl(\tfrac12e^{u/2}\mathbf 1_+\bigr)(s)-\tfrac12\mathcal L(\mathfrak g\mathbf 1_+)(s)+\tfrac14\log\frac z2-\tfrac14\log\pi .
\end{equation}
Finally $\mathcal L\bigl(\mathrm{fp}\,\frac{\mathbf 1_+}u\bigr)(s)=\int_0^1\frac{e^{-zu}-1}u\,du+\int_1^\infty\frac{e^{-zu}}u\,du=-\gamma_E-\log z$: the left side has derivative $-\int_0^\infty e^{-zu}du=-1/z$ in $z$, and its value at $z=1$ is $\int_0^1\frac{e^{-u}-1}u\,du+\int_1^\infty\frac{e^{-u}}u\,du=-\gamma_E$, the classical integral representation of Euler's constant \cite[Ch.~XII]{WhittakerWatson1927}. Substituting, $\frac14\log\frac z2=-\frac14\mathcal L(\mathrm{fp}\,\frac{\mathbf 1_+}u)-\frac14(\gamma_E+\log2)$, and \CGeqref{eq:rb-aux-13} becomes \CGeqref{eq:rb-LK} with the constant $\mathrm c_{\mathscr K}$. (The identity was also checked numerically to $25$ digits at four points $s$; file \texttt{laplace\_check.py}.)
\end{proof}

\begin{lemma}[Laplace transform of the residual]\label{lem:rb-laplace}
Let $\boldsymbol\kappa$ be $C$-admissible and put $\mathfrak H_{\boldsymbol\kappa}(u):=\sum_k\kappa_k^{-2}\cos(\kappa_ku)$, which converges absolutely and uniformly on $\mathbb R$ since $\mathfrak H_{\boldsymbol\kappa}(0)=\sum_k\kappa_k^{-2}<\infty$, and
\begin{equation}\label{eq:rb-S-A-dist}
 \mathscr S_{\boldsymbol\kappa}:=-\partial_u^2\bigl(\mathfrak H_{\boldsymbol\kappa}\mathbf 1_+\bigr)+\mathfrak H_{\boldsymbol\kappa}(0)\,\delta_0',\qquad
 \mathscr A_\zeta:=\sum_{n\ge2}\mathsf a^\zeta_n\,\delta_{\log n},\qquad
 \mathscr E_{\boldsymbol\kappa}:=\mathscr S_{\boldsymbol\kappa}+\mathscr A_\zeta-\mathscr K .
\end{equation}
Then $\mathscr E_{\boldsymbol\kappa}$ is a distribution of local order at most $2$ (its restriction to every bounded interval is of order $\le2$; the support is not compact, since $\mathscr A_\zeta$ is a locally finite measure on $[\log2,\infty)$), supported in $[0,\infty)$, Laplace-transformable on $\Re s>1$, whose restriction to $(0,\infty)$ is $\mathfrak R(\boldsymbol\kappa,\mathsf a^\zeta)$, and for $\Re s>1$
\begin{equation}\label{eq:rb-LR}
 \mathcal L\mathscr E_{\boldsymbol\kappa}(s)=\widehat{\mathsf R}(s)=\tfrac12\Bigl(\frac{P_{\boldsymbol\kappa}'}{P_{\boldsymbol\kappa}}-\frac{\xi'}{\xi}\Bigr)(s),\qquad\xi:=\tfrac12G\zeta ,
\end{equation}
where $\widehat{\mathsf R}$ is \CGeqref{eq:rb-Rhat} with the masses $\mathsf a^\zeta$ for all $n$. The same holds for a finite section $(\boldsymbol\kappa,\mathsf a)$ with $\mathscr A_\zeta$ replaced by $\sum_{n\le n^*}\mathsf a_n\delta_{\log n}$ and $-\frac12\xi'/\xi$ by $-\frac12G'/G+\sum_{n\le n^*}\mathsf a_nn^{\frac12-s}$.
\end{lemma}

\begin{proof}
\emph{Spectral part.} $\kappa_k\ge\tau_k-C\gg k/\log(k+1)$, so $\mathfrak H_{\boldsymbol\kappa}(0)<\infty$ and $\mathfrak H_{\boldsymbol\kappa}$ is continuous and bounded by $\mathfrak H_{\boldsymbol\kappa}(0)$. In $\mathcal D'(0,\infty)$ a uniformly convergent series may be differentiated termwise, so $-\mathfrak H_{\boldsymbol\kappa}''=\sum_k\cos(\kappa_ku)$ there; the term $\mathfrak H_{\boldsymbol\kappa}(0)\delta_0'$ is supported at $0$. Hence $\mathscr S_{\boldsymbol\kappa}|_{(0,\infty)}=\sum_k\cos(\kappa_ku)$, and $\mathscr S_{\boldsymbol\kappa}$ has order at most $2$. The Laplace transform of a distribution $\mathscr T$ supported in $[0,\infty)$ is $\mathcal L\mathscr T(s):=\langle\mathscr T,\chi e^{-zu}\rangle$, $z=s-\frac12$, with $\chi\in C^\infty(\mathbb R)$ equal to $1$ on a neighbourhood of the support and to $0$ on $(-\infty,-1]$; thus $\mathcal L(\delta_0')(s)=-\frac d{du}e^{-zu}\big|_{u=0}=z$ and $\mathcal L(\mathbf 1_+f)(s)=\int_0^\infty f(u)e^{-zu}du$. For $\Re z>0$, Fubini (the double integral $\sum_k\kappa_k^{-2}\int_0^\infty e^{-\Re z\,u}du$ is finite) gives $\mathcal L(\mathfrak H_{\boldsymbol\kappa}\mathbf 1_+)(s)=\sum_k\kappa_k^{-2}\frac z{z^2+\kappa_k^2}$, so
\begin{equation}\label{eq:rb-aux-14}
 \mathcal L\mathscr S_{\boldsymbol\kappa}(s)=-z^2\sum_k\frac z{\kappa_k^2(z^2+\kappa_k^2)}+z\sum_k\frac1{\kappa_k^2}
 =\sum_k\frac{z}{z^2+\kappa_k^2}=\tfrac12\frac{P_{\boldsymbol\kappa}'}{P_{\boldsymbol\kappa}}(s),
\end{equation}
because $\frac z{\kappa^2}-\frac{z^3}{\kappa^2(z^2+\kappa^2)}=\frac z{z^2+\kappa^2}$, and $P_{\boldsymbol\kappa}=\prod_k\frac{z^2+\kappa_k^2}{\frac14+\kappa_k^2}$ converges locally uniformly, so its logarithmic derivative is the termwise sum $\sum_k\frac{2z}{z^2+\kappa_k^2}$.

\emph{Arithmetic part.} On $[0,u^*]$ only $n\le e^{u^*}$ contribute, so $\mathscr A_\zeta$ is a locally finite measure, of order $0$, and $e^{-au}\mathscr A_\zeta$ is a finite measure for $a>\frac12$. For $\Re s>1$, $\mathcal L\mathscr A_\zeta(s)=\sum_n\frac{\Lambda(n)}{2\sqrt n}n^{-z}=\frac12\sum_n\Lambda(n)n^{-s}=-\frac12\frac{\zeta'}\zeta(s)$.

\emph{Assembly.} By Lemma~\CGref{lem:rb-archimedean}, $\mathcal L\mathscr E_{\boldsymbol\kappa}=\frac12P'_{\boldsymbol\kappa}/P_{\boldsymbol\kappa}-\frac12\zeta'/\zeta-\frac12G'/G$, and $(G\zeta)'/(G\zeta)=\xi'/\xi$ since $\xi=\frac12G\zeta$. The restriction to $(0,\infty)$ is $\sum_k\cos(\kappa_ku)+\sum_n\mathsf a^\zeta_n\delta(u-\log n)-\mathsf k(u)$, which is \CGeqref{eq:rb-residual}. The finite-section case is identical with a finite sum.
\end{proof}

\begin{remark}[The boundary term]\label{rem:rb-boundary-term}
Without the correction $\mathfrak H_{\boldsymbol\kappa}(0)\delta_0'$ in \CGeqref{eq:rb-S-A-dist}, the transform of $-\partial_u^2(\mathfrak H_{\boldsymbol\kappa}\mathbf 1_+)$ is $\frac12P'_{\boldsymbol\kappa}/P_{\boldsymbol\kappa}-(s-\frac12)\mathfrak H_{\boldsymbol\kappa}(0)$: the second derivative of the cut-off function produces $\mathfrak H_{\boldsymbol\kappa}(0)\delta_0'=\mathfrak H_{\boldsymbol\kappa}(0)\delta_0'$. Either normalisation gives the same restriction to $(0,\infty)$ and differs by a polynomial in $s$, which is harmless in the argument below. Both $\mathscr K$ and $\mathscr S_{\boldsymbol\kappa}$ carry singular terms at $u=0$; this is why the rung system is posed in $\mathcal D'(0,\infty)$ and not on $[0,\infty)$.
\end{remark}

\begin{proposition}[Uniqueness with prescribed masses]\label{prop:rb-prescribed-uniqueness}
Let $\boldsymbol\kappa$ be $C$-admissible for some $C$, and suppose $\mathfrak R(\boldsymbol\kappa,\mathsf a^\zeta)=0$ in $\mathcal D'(0,\infty)$. Then all nontrivial zeros of $\zeta$ lie on $\Re s=\frac12$, and $\boldsymbol\kappa=\boldsymbol\gamma$.
\end{proposition}

\begin{proof}
By Lemma~\CGref{lem:rb-laplace}, $\mathscr E_{\boldsymbol\kappa}$ is a distribution of finite order supported in $[0,\infty)$ whose restriction to $(0,\infty)$ vanishes, so it is supported at $\{0\}$. A distribution of finite order supported at a point is a finite linear combination of $\delta_0$ and its derivatives \cite[Thm.~2.3.4]{Hormander1990}, so $\widehat{\mathsf R}=\mathcal L\mathscr E_{\boldsymbol\kappa}$ is a polynomial $p$ in $z=s-\frac12$. By \CGeqref{eq:rb-LR}, $(P_{\boldsymbol\kappa}e^{-Q}/\xi)'=0$ on $\Re s>1$ with $Q'=2p$, $Q$ a polynomial. Hence $P_{\boldsymbol\kappa}=c\,\xi\,e^{Q}$ on $\Re s>1$ with $c\ne0$, and, both sides being entire, the identity holds on all of $\mathbb C$.  Only now is the quotient discussed: the identity shows that $P_{\boldsymbol\kappa}$ and $\xi$ have the same zeros with the same multiplicities, since $e^Q$ has none.  In the variable $w=s-\frac12$, $P_{\boldsymbol\kappa}$ is a constant multiple of $\prod_k(1+w^2/\kappa_k^2)$, and $N_{\boldsymbol\kappa}(t)\ll t\log t$, so $P_{\boldsymbol\kappa}$ has order $1$; $\xi$ has order $1$. Comparing the Hadamard factorizations $P_{\boldsymbol\kappa}=e^{A_1+B_1s}\Pi(s)$ and $\xi=e^{A_2+B_2s}\Pi(s)$, with the same canonical product $\Pi$ over the common zeros, gives $e^{Q}=c^{-1}e^{(A_1-A_2)+(B_1-B_2)s}$, so $\deg Q\le1$. Since $P_{\boldsymbol\kappa}(1-s)=P_{\boldsymbol\kappa}(s)$ and $\xi(1-s)=\xi(s)$, $e^{Q(1-s)}=e^{Q(s)}$, that is, $Q(1-s)-Q(s)\in2\pi i\mathbb Z$ for every $s$; the left side is continuous, hence constant, and for $Q(s)=a+bs$ it equals $b(1-2s)$, so $b=0$. Therefore $P_{\boldsymbol\kappa}=c'\xi$. The zeros of $P_{\boldsymbol\kappa}$ are $\frac12\pm i\kappa_k$, so every zero of $\xi$, that is, every nontrivial zero of $\zeta$, has real part $\frac12$, and the ordinates counted with multiplicity are the $\kappa_k$. Both lists are nondecreasing, so $\boldsymbol\kappa=\boldsymbol\gamma$.
\end{proof}

\begin{definition}[Statement B$'$]\label{def:rb-Bprime}
\emph{Statement B$'(C)$} says that there are $C$-admissible lists $\boldsymbol\kappa^{(m)}$, each equal to the clock beyond some finite index, with $\mathsf s_m(\boldsymbol\kappa^{(m)},\mathsf a^\zeta)\to0$.
\end{definition}

\begin{corollary}[Route B$'$]\label{cor:rb-Bprime}
\begin{enumerate}[label=\textup{(\alph*)}]
\item If B$'(C)$ holds for some $C$, then every sequence as in Definition~\CGref{def:rb-Bprime} converges coordinatewise to $\boldsymbol\gamma$, and all nontrivial zeros of $\zeta$ lie on $\Re s=\frac12$.
\item If all nontrivial zeros of $\zeta$ lie on $\Re s=\frac12$, then B$'(C)$ holds for every $C\ge C_\eta$.
\end{enumerate}
Hence the Riemann hypothesis is equivalent to B$'(C)$ for some $C<\infty$.
\end{corollary}

\begin{proof}
(a) $|\mathsf a^\zeta_n|\le \mathbb M_*$ by the elementary computation in Lemma~\CGref{lem:rb-canonical-box}, so $(\boldsymbol\kappa^{(m)},\mathsf a^\zeta)\in\mathbb K(C,\mathbb M_*)$. By Proposition~\CGref{prop:rb-compactness}, every subsequence has a further subsequence converging coordinatewise to some $(\boldsymbol\kappa,\mathsf a^\zeta)$ with $\mathfrak R(\boldsymbol\kappa,\mathsf a^\zeta)=0$ and $\boldsymbol\kappa$ $C$-admissible; the masses do not move. By Proposition~\CGref{prop:rb-prescribed-uniqueness}, $\boldsymbol\kappa=\boldsymbol\gamma$ and the zeros lie on the line. The product $\prod_k[\tau_k-C,\tau_k+C]$ is compact and metrizable, for instance by $\mathrm d(\boldsymbol\kappa,\boldsymbol\kappa')=\sum_k2^{-k}\min(1,|\kappa_k-\kappa'_k|)$, which induces the product topology; in a compact metric space a sequence all of whose convergent subsequences have the same limit converges to it. Hence the whole sequence converges. (b) Under this hypothesis the rung system with $(\boldsymbol\gamma,\mathsf a^\zeta)$ is the distributional explicit formula \CGref{thm:distributional-EF}, and $|\gamma_k-\tau_k|\le C_\eta$ by \CGref{lem:bounded-displacement}. The truncations in the proof of Corollary~\CGref{cor:rb-D-implies-B}, with the masses kept equal to $\mathsf a^\zeta$, give the sections.
\end{proof}

\begin{remark}[What prescribing the masses changes, and what it does not]\label{rem:rb-prescribed}
Hypothesis~\CGref{hyp:rb-U} is no longer needed: uniqueness holds for every list at bounded distance from the clock, not only in $\mathscr O$, and it follows from Hadamard's theorem and the functional equation instead of \CGref{thm:free-masses}. Route B$'$ has a single open piece, the existence statement B$'$, which by Corollary~\CGref{cor:rb-Bprime} is equivalent to the Riemann hypothesis. With prescribed masses the whole window $(0,u^*)$ constrains the ordinates; the part above $\log\Omega$ can no longer be absorbed by the masses, as it could in Heuristic~\CGref{heur:rb-informative}.

The equivalence itself is classical in substance. The prime side of the rung system is the prime side of the explicit formula of Riemann and von Mangoldt in the form of Guinand and Weil \cite{Guinand1948,Weil1952}, and the uniqueness is Hadamard's factorisation. Arias de Reyna \cite{AriasdeReyna2024} proves that the Riemann hypothesis is equivalent to the temperedness of the measure $-\sum_n\frac{\Lambda(n)}{\sqrt n}(\delta_{\log n}+\delta_{-\log n})+2\cosh(x/2)\,dx$, the same prime side up to the factor $2$ and the symmetrisation, and identifies its Fourier transform with $\sum_\gamma\delta_{\gamma/2\pi}-2\vartheta'(2\pi t)\,dt$ under that hypothesis; by partial summation his criterion is von Koch's $\sum_{n\le x}\Lambda(n)=x+O(x^{1/2}\log^2x)$. Statement B$'$ is the finite-section form of the same equivalence: it asks for a real spectrum in a tube around the clock that reproduces the prime side against a dense family of tests. It is not a weaker intermediate statement, and the finite sections are not more constrained than the global statement; what they offer is a formulation in which each stage is a finite-dimensional problem with computable certificates (Section~\CGref{subsec:rb-staged}) and explicit candidates (Section~\CGref{subsec:rb-gonek}). Dyson's suggestion that the zeros form a one-dimensional quasicrystal, and the classification of crystalline measures by Kurasov and Sarnak \cite{KurasovSarnak2020} and Olevskii and Ulanovskii \cite{OlevskiiUlanovskii2020}, concern the same object; the zeta list is not in those classes because of its density, as Arias de Reyna points out.
\end{remark}

\subsubsection{Staged construction: a certificate theorem and a graded extension}\label{subsec:rb-staged}

\begin{definition}[Staged sections]\label{def:rb-staged}
Fix $u^*_1<u^*_2<\cdots\to\infty$, finite test sets $\mathscr T_1\subset\mathscr T_2\subset\cdots$ supported in $(0,u^*_m)$ with dense union as in Section~\CGref{subsec:rb-routeB}, $\mathrm t_m:=|\mathscr T_m|$, and $k^*_m\to\infty$. Stage $m$ starts from a list $\boldsymbol\kappa^{(m,0)}$ that equals the clock beyond some finite index, moves the coordinates $k\le k^*_m$ by $x\in\mathbb R^{k^*_m}$, and has masses $\mathsf a^\zeta$. Its residual map is
\begin{equation}\label{eq:rb-Fm}
 \mathbb F_m:\mathbb R^{k^*_m}\to\mathbb R^{\mathrm t_m},\qquad \mathbb F_m(x)_\varphi:=\bigl\langle\mathfrak R(\boldsymbol\kappa^{(m,0)}+x,\mathsf a^\zeta),\varphi\bigr\rangle\quad(\varphi\in\mathscr T_m),
\end{equation}
where $(\boldsymbol\kappa+x)_k=\kappa_k+x_k$ for $k\le k^*_m$ and $\kappa_k$ otherwise, and its Jacobian at $0$ is
\begin{equation}\label{eq:rb-jacobian}
 (\mathbb J_m)_{\varphi,k}=\partial_{x_k}\mathbb F_m(0)_\varphi=-\int_0^\infty\varphi(u)\,u\sin(\kappa^{(m,0)}_ku)\,du\qquad(\varphi\in\mathscr T_m,\ k\le k^*_m).
\end{equation}
We use the norms $\|x\|_\infty$ on $\mathbb R^{k^*_m}$ and $\mathbb R^{\mathrm t_m}$ and the induced operator norms $\|\cdot\|_{\infty\to\infty}$ (maximal absolute row sum). Put
\begin{equation}\label{eq:rb-eta-H}
 \mathsf h_m:=\|\mathbb F_m(0)\|_\infty,\qquad \mathbb H_m:=k^*_m\max_{\varphi\in\mathscr T_m}\int_0^\infty|\varphi(u)|\,u^2\,du .
\end{equation}
\end{definition}

\begin{lemma}[Monotone matching]\label{lem:rb-matching}
Let $(\mathrm a_k)_{k\ge1}$ be real with $|\mathrm a_k-\tau_k|\le C$ for all $k$, and let $(\mathrm a_{(k)})$ be its nondecreasing rearrangement. Then $|\mathrm a_{(k)}-\tau_k|\le C$ for all $k$. If moreover $\mathrm a_k=\tau_k$ for $k>k^*$, then $\mathrm a_{(k)}=\tau_k$ for all $k>k^*+j$, where $j:=\#\{k>k^*:\tau_k\le\tau_{k^*}+C\}<\infty$.
\end{lemma}

\begin{proof}
At least $k$ of the $\mathrm a_i$ (those with $i\le k$) are $\le\tau_k+C$, so $\mathrm a_{(k)}\le\tau_k+C$; at least $k$ of them (those with $i\ge k$) are $\ge\tau_k-C$, hence at most $k-1$ are $<\tau_k-C$, so $\mathrm a_{(k)}\ge\tau_k-C$. For the second claim, the elements of the multiset not exceeding $\tau_{k^*}+C$ are the $k^*$ moved ones (each is at most $\tau_{k^*}+C$) and the $j$ clock points $\tau_k$, $k^*<k\le k^*+j$; the remaining elements are exactly $\tau_k$, $k>k^*+j$, in increasing order.
\end{proof}

\begin{lemma}[Lipschitz bound for the Jacobian]\label{lem:rb-lipschitz}
For all $x,x'\in\mathbb R^{k^*_m}$, $\|D\mathbb F_m(x)-D\mathbb F_m(x')\|_{\infty\to\infty}\le \mathbb H_m\|x-x'\|_\infty$, and
\begin{equation}\label{eq:rb-taylor}
 \|\mathbb F_m(x)-\mathbb F_m(0)-\mathbb J_mx\|_\infty\le\tfrac12\mathbb H_m\|x\|_\infty^2 .
\end{equation}
\end{lemma}

\begin{proof}
$\mathbb F_m(x)_\varphi=\sum_{k\le k^*_m}\mathsf c_\varphi(\kappa^{(m,0)}_k+x_k)+\text{const}$, so $D^2\mathbb F_m(x)_\varphi$ is diagonal with entries $\partial_{x_k}^2\mathbb F_m(x)_\varphi=-\int\varphi(u)u^2\cos((\kappa_k+x_k)u)\,du$, of modulus at most $\int|\varphi|u^2$. The row $\varphi$ of $D\mathbb F_m(x)-D\mathbb F_m(x')$ has absolute sum at most $\sum_{k\le k^*_m}\int|\varphi|u^2\,|x_k-x'_k|\le \mathbb H_m\|x-x'\|_\infty$. The Taylor bound follows by integrating $D\mathbb F_m$ along the segment.
\end{proof}

\begin{theorem}[Certificate for one stage]\label{thm:rb-certificate}
Suppose $\mathbb J_m$ has a right inverse $\mathbb Q_m:\mathbb R^{\mathrm t_m}\to\mathbb R^{k^*_m}$, $\mathbb J_m\mathbb Q_m=I_{\mathrm t_m}$, with $\|\mathbb Q_m\|_{\infty\to\infty}\le \mathbb L_m$, and
\begin{equation}\label{eq:rb-Am}
 \mathsf l_m:=2\mathbb L_m^2\mathbb H_m\mathsf h_m<1 .\tag{A$_m$}
\end{equation}
Then the map $\mathbb T_m(y):=y-\mathbb F_m(\mathbb Q_my)$ is a contraction of the ball $\{\|y\|_\infty\le2\mathsf h_m\}\subset\mathbb R^{\mathrm t_m}$ with constant $\mathsf l_m$, and $x_*:=\mathbb Q_my_*$, $y_*$ its unique fixed point in that ball, satisfies
\begin{equation}\label{eq:rb-cert-conclusion}
 \mathbb F_m(x_*)=0,\qquad\|x_*\|_\infty\le2\mathbb L_m\mathsf h_m .
\end{equation}
The uniqueness is that of $x_*$ among the zeros of $\mathbb F_m$ in $\mathbb Q_m\{\|y\|_\infty\le2\mathsf h_m\}\subset\operatorname{Ran}\mathbb Q_m$; nothing is asserted about zeros of $\mathbb F_m$ outside $\operatorname{Ran}\mathbb Q_m$.
\end{theorem}

\begin{proof}
$D\mathbb T_m(y)=I-D\mathbb F_m(\mathbb Q_my)\mathbb Q_m=(\mathbb J_m-D\mathbb F_m(\mathbb Q_my))\mathbb Q_m$ because $\mathbb J_m\mathbb Q_m=I$, so by Lemma~\CGref{lem:rb-lipschitz}, for $\|y\|_\infty\le2\mathsf h_m$, $\|D\mathbb T_m(y)\|_{\infty\to\infty}\le \mathbb H_m\|\mathbb Q_my\|_\infty \mathbb L_m\le2\mathbb H_m\mathbb L_m^2\mathsf h_m=\mathsf l_m<1$. Moreover $\mathbb T_m(y)=-\mathbb F_m(0)-\bigl(\mathbb F_m(\mathbb Q_my)-\mathbb F_m(0)-\mathbb J_m\mathbb Q_my\bigr)$, and by \CGeqref{eq:rb-taylor} the bracket has norm at most $\frac12\mathbb H_m\mathbb L_m^2\|y\|_\infty^2\le2\mathbb H_m\mathbb L_m^2\mathsf h_m^2$. Hence $\|\mathbb T_m(y)\|_\infty\le\mathsf h_m(1+\mathsf l_m)<2\mathsf h_m$: the ball is mapped into itself. Banach's fixed point theorem gives a unique fixed point $y_*$ in the ball, $y_*=y_*-\mathbb F_m(\mathbb Q_my_*)$, that is, $\mathbb F_m(x_*)=0$ with $\|x_*\|_\infty\le \mathbb L_m\|y_*\|_\infty\le2\mathbb L_m\mathsf h_m$. The uniqueness is that of $y_*$ in the ball of $\mathbb R^{\mathrm t_m}$, hence of $x_*$ in $\operatorname{Ran}\mathbb Q_m$; when $k^*_m>\mathrm t_m$ the Jacobian has a kernel and $\mathbb F_m$ may have other zeros in the ball of $\mathbb R^{k^*_m}$, about which the theorem says nothing. The iteration $y_{q+1}=\mathbb T_m(y_q)$, $y_0=0$, converges to $y_*$ geometrically with ratio $\mathsf l_m$; in the variable $x$ it is $x_{q+1}=x_q-\mathbb Q_m\mathbb F_m(x_q)$, a \emph{modified} Newton iteration, the right inverse being frozen at $\mathbb Q_m$ instead of being recomputed at each $x_q$.
\end{proof}

\begin{remark}[Evaluation errors and order of the ordinates]\label{rem:rb-certificate-errors}
Theorem~\CGref{thm:rb-certificate} is stated for exact values of $\mathsf h_m$, $\mathbb L_m$ and $\mathbb H_m$.  If the residual is only known up to an error, $\|\tilde{\mathbb F}_m(0)-\mathbb F_m(0)\|_\infty\le\epsilon_m$, the theorem applies with $\mathsf h_m$ replaced by $\|\tilde{\mathbb F}_m(0)\|_\infty+\epsilon_m$; if the computed matrix $\tilde{\mathbb Q}_m$ satisfies only $\|I-\mathbb J_m\tilde{\mathbb Q}_m\|_{\infty\to\infty}\le\delta_m<1$, then $\mathbb Q_m:=\tilde{\mathbb Q}_m(\mathbb J_m\tilde{\mathbb Q}_m)^{-1}$ is an exact right inverse with $\|\mathbb Q_m\|\le\|\tilde{\mathbb Q}_m\|/(1-\delta_m)$, and this is the $\mathbb L_m$ to use.  A rigorous application therefore needs enclosures of $\mathbb F_m(0)$ (quadrature and truncation of the ordinate sums), of $\mathbb J_m\tilde{\mathbb Q}_m$ (rounding) and of $\mathbb H_m$; none of these enclosures is produced in Section~\CGref{subsec:rb-numerics}, where the word \texttt{CERTIFIED} in the console output means only that the floating-point value of $\mathsf l_m$ is below $1$.  The theorem also does not guarantee that the corrected list $\boldsymbol\kappa^{(m,0)}+x_*$ is nondecreasing; this holds when $2\mathbb L_m\mathsf h_m$ is smaller than half the least gap of the moved ordinates, and otherwise the list is rearranged by Lemma~\CGref{lem:rb-matching}, which preserves the tube but changes the section.  Finally, certifying each stage separately gives no compatibility between stages: the solutions $x_*^{(m)}$ need not be restrictions of one another, and the uniform estimates that would be needed to pass to an infinite sequence (a stage-independent bound on $\mathbb L_m\mathbb H_m\mathsf h_m$ together with the compactness of Proposition~\CGref{prop:rb-compactness}) are not available.
\end{remark}

\begin{theorem}[Staged certificate implies route B$'$]\label{thm:rb-staged}
Let $C<\tau_1-2\pi$. Start at $\boldsymbol\kappa^{(1,0)}:=\boldsymbol\tau$ and, for $m\ge1$, let $\boldsymbol\kappa^{(m)}$ be the nondecreasing rearrangement of $\boldsymbol\kappa^{(m,0)}+x_*^{(m)}$ given by Theorem~\CGref{thm:rb-certificate}, and $\boldsymbol\kappa^{(m+1,0)}:=\boldsymbol\kappa^{(m)}$. Suppose \textup{(A$_m$)} holds for every $m$ and
\begin{equation}\label{eq:rb-B}
 \sum_{j\ge1}2\mathbb L_j\mathsf h_j=\sup_{m\ge1}\ \sum_{j\le m}2\mathbb L_j\mathsf h_j\le C .\tag{B}
\end{equation}
Then every $\boldsymbol\kappa^{(m)}$ is $C$-admissible and equals the clock beyond a finite index $\widetilde k^*_m$, $\mathsf s_m(\boldsymbol\kappa^{(m)},\mathsf a^\zeta)=0$ for every $m$, B$'(C)$ holds, and therefore (Corollary~\CGref{cor:rb-Bprime}) $\boldsymbol\kappa^{(m)}\to\boldsymbol\gamma$ coordinatewise and the Riemann hypothesis holds. In \textup{(A$_{m+1}$)} and \textup{(B)} the quantity $\mathsf h_{m+1}$ is $\max_{\varphi\in\mathscr T_{m+1}\setminus\mathscr T_m}|\langle\mathfrak R(\boldsymbol\kappa^{(m)},\mathsf a^\zeta),\varphi\rangle|$, the residual of the new tests only.
\end{theorem}

\begin{proof}
Induction on $m$. Suppose $\boldsymbol\kappa^{(m,0)}$ is $C_{m-1}$-admissible with $C_{m-1}:=\sum_{j<m}2\mathbb L_j\mathsf h_j$ ($C_0=0$) and equals the clock beyond a finite index. By Theorem~\CGref{thm:rb-certificate}, $\boldsymbol\kappa^{(m,0)}+x_*^{(m)}$ satisfies $|(\cdot)_k-\tau_k|\le C_{m-1}+2\mathbb L_m\mathsf h_m=C_m$ for all $k$, and only finitely many coordinates differ from the clock. Lemma~\CGref{lem:rb-matching} shows that its rearrangement $\boldsymbol\kappa^{(m)}$ satisfies $|\kappa^{(m)}_k-\tau_k|\le C_m\le C$ and equals the clock beyond some $\widetilde k^*_m$; it is nondecreasing, and $\kappa^{(m)}_1\ge\tau_1-C>2\pi$. So $\boldsymbol\kappa^{(m)}$ is $C$-admissible. The residual $\mathfrak R$ depends only on the multiset of ordinates, so $\langle\mathfrak R(\boldsymbol\kappa^{(m)},\mathsf a^\zeta),\varphi\rangle=\mathbb F_m(x_*^{(m)})_\varphi=0$ for all $\varphi\in\mathscr T_m$, that is, $\mathsf s_m=0$. Since $\mathscr T_m\subset\mathscr T_{m+1}$, the start of stage $m+1$ has zero residual on $\mathscr T_m$, which gives the stated form of $\mathsf h_{m+1}$. B$'(C)$ holds with $\mathsf s_m=0$, and Corollary~\CGref{cor:rb-Bprime}(a) gives the rest.
\end{proof}

\begin{remark}[What the certificate theorem does and does not do]\label{rem:rb-certificate-status}
Theorem~\CGref{thm:rb-staged} replaces the extension lemma of version 3 of the companion note (there Hypothesis~E) by two computable quantities per stage, $\mathsf l_m=2\mathbb L_m^2\mathbb H_m\mathsf h_m$ and $2\mathbb L_m\mathsf h_m$. It requires no contraction factor $q$, no separate control of the perturbation of old tests, and no prior ``entry into the basin'': old tests are kept in $\mathscr T_{m+1}$ and are solved again exactly. The two conditions are not independent of the height: $\mathbb H_m\ge k^*_m\int|\varphi|u^2$ grows with the number of unknowns, and $\mathbb L_m$ is the $\ell^\infty$ norm of a right inverse of the Jacobian, whose growth is measured in Section~\CGref{subsec:rb-numerics}. The conditions (A$_m$) for all $m$ and (B) together imply the Riemann hypothesis, so they are not a weaker intermediate step; what changes is that each stage is a statement about one finite matrix and one finite residual vector, verifiable in principle by interval arithmetic. The restriction $C<\tau_1-2\pi=8.24\ldots$ only ensures $\kappa_1\ge2\pi$ and can be removed by treating the first coordinate separately.
\end{remark}

\begin{proposition}[Graded extension]\label{prop:rb-graded}
Let $\mathsf s_m\ge0$ satisfy $\mathsf s_{m+1}\le \mathrm q_m\mathsf s_m+\mathsf h_m$ with $0\le \mathrm q_m\le1$ and $\sum_m\mathsf h_m<\infty$, and assume
\begin{equation}\label{eq:rb-graded-tail}
 \prod_{j=L}^{m}\mathrm q_j\longrightarrow0\quad(m\to\infty)\qquad\text{for every }L\ge1 .
\end{equation}
Then $\mathsf s_m\to0$. Condition \CGeqref{eq:rb-graded-tail} holds if $\mathrm q_m\le\mathrm q<1$ for all large $m$, and more generally if $\sum_m(1-\mathrm q_m)=\infty$. If the lists $\boldsymbol\kappa^{(m)}$ producing $\mathsf s_m=\mathsf s_m(\boldsymbol\kappa^{(m)},\mathsf a^\zeta)$ are $C$-admissible sections for a fixed $C$, then B$'(C)$ holds.
\end{proposition}

\begin{proof}
Iterating, $\mathsf s_m\le\mathsf s_1\prod_{j<m}\mathrm q_j+\sum_{\ell<m}\mathsf h_\ell\prod_{\ell<j<m}\mathrm q_j$. Given $\epsilon>0$ choose $L$ with $\sum_{\ell\ge L}\mathsf h_\ell<\epsilon$. For $m>L$ and $\ell<L$, every $\mathrm q_j$ is at most $1$, so $\prod_{\ell<j<m}\mathrm q_j\le\prod_{j=L}^{m-1}\mathrm q_j$, and the same bound holds for $\prod_{j<m}\mathrm q_j$; hence the first term and the terms with $\ell<L$ are together at most $\bigl(\mathsf s_1+\sum_{\ell<L}\mathsf h_\ell\bigr)\prod_{j=L}^{m-1}\mathrm q_j$, which tends to $0$ by \CGeqref{eq:rb-graded-tail}. The terms with $\ell\ge L$ are at most $\sum_{\ell\ge L}\mathsf h_\ell<\epsilon$. Hence $\limsup_m\mathsf s_m\le\epsilon$, and $\mathsf s_m\to0$. The weaker hypothesis $\prod_{m\ge1}\mathrm q_m=0$ would not suffice: with $\mathrm q_1=0$ and $\mathrm q_m=1$ for $m\ge2$ the product vanishes, while $\mathsf s_{m+1}\le\mathsf s_m+\mathsf h_m$ allows $\mathsf s_m\to\mathsf s_2>0$. For the two sufficient conditions, $\prod_{j=L}^m\mathrm q_j\le\mathrm q^{\,m-\max(L,m_0)}$ when $\mathrm q_j\le\mathrm q$ for $j\ge m_0$, and $\prod_{j=L}^m\mathrm q_j\le\exp\bigl(-\sum_{j=L}^m(1-\mathrm q_j)\bigr)$ in general, since $\mathrm q\le e^{-(1-\mathrm q)}$ on $[0,1]$. The second claim is Definition~\CGref{def:rb-Bprime}.
\end{proof}

\begin{remark}[Tube budget, and the two geometries]\label{rem:rb-graded-status}
In Proposition~\CGref{prop:rb-graded} the tube $|\kappa^{(m)}_k-\tau_k|\le C$, uniform in $m$, is a separate hypothesis: summability of the coordinate increments $\sum_m|\kappa^{(m+1)}_k-\kappa^{(m)}_k|$ gives convergence of each coordinate but not that the limit stays in a prescribed tube, so a construction must carry a displacement budget $\sum_{m\ge m(k)}|\kappa^{(m+1)}_k-\kappa^{(m)}_k|\le C-|\kappa^{(m(k))}_k-\tau_k|$. For B$'(C)$ itself no convergence of the whole sequence is needed: compactness in the product of the intervals $[\tau_k-C,\tau_k+C]$ (Proposition~\CGref{prop:rb-compactness}) extracts a convergent subsequence, and uniqueness identifies the limit. A second point concerns norms. Theorem~\CGref{thm:rb-certificate} is stated in $\ell^\infty$, where $\mathbb H_m$ has the simple form \CGeqref{eq:rb-eta-H} because the Hessian of each $\mathsf c_\varphi$ is diagonal. A bound $\|\mathbb J_m^\dagger\|_{2\to2}\le1/\sigma_{\min}(\mathbb J_m)$ is a Euclidean statement; passing to $\ell^\infty$ costs factors $\sqrt{\mathrm t_m}$ and $\sqrt{K_m}$ unless the right inverse is controlled directly in $\ell^\infty_{\mathrm t_m}\to\ell^\infty_{K_m}$. The numerical constants of Section~\CGref{subsec:rb-numerics} are computed directly in $\ell^\infty$.
\end{remark}

\begin{heuristic}[Frame bound from exponential sums of the clock]\label{heur:rb-frame}
In the opposite regime $k^*_m\gg \mathrm t_m$ one may take $\mathbb Q_m:=\mathbb J_m^{\!T}(\mathbb J_m\mathbb J_m^{\!T})^{-1}$. If $\|\frac1{K_m}\mathbb J_m\mathbb J_m^{\!T}-D\|_{\infty\to\infty}\le\epsilon<d_-$ for a diagonal $D\ge d_-I$, and $|(\mathbb J_m)_{\varphi,k}|\le\mathrm a_J$, then $\mathbb L_m\le \mathrm t_mA/(k^*_m(d_--\epsilon))=O(\mathrm t_m/k^*_m)$ and $\mathbb L_m^2\mathbb H_m=O(\mathrm t_m^2B/k^*_m)$ with $\mathrm B_0:=\max_\varphi\int|\varphi|u^2$, so (A$_m$) becomes $k^*_m\gg \mathrm t_m^2B\mathsf h_m$: a wide block compensates a moderate residual. In the point model $\varphi\approx\delta_{u_j}$ the entries of $\frac1{k^*}\mathbb J\mathbb J^{\!T}$ are $u_ju_\ell\cdot\frac1k^*\sum_{k\in\mathfrak B}\sin(\tau_ku_j)\sin(\tau_ku_\ell)$, so the frame property reduces to cancellation in the exponential sums $\frac1k^*\sum_{k\in\mathfrak B}e^{i\mathrm v\tau_k}$ at the frequencies $\mathrm v=u_j\pm u_\ell\ne0$. On a block of $k^*$ consecutive clock points at height $T$ the spacing is $\Delta_T\approx2\pi/\log T$ with curvature $O((k^*)^2/(T\log^3T))$, which suggests $|\frac1k^*\sum e^{i\mathrm v\tau_k}|\lesssim\frac{\log T}{K|v|}+\frac{|v|(k^*)^2}{T\log^3T}$ and a window $\frac{\log T}h\ll k^*\ll(T\log^3T/u^*)^{1/2}$ for tests on a grid of step $\Delta u$. None of this is proved here, and the point model omits a decisive feature of smooth tests: a test of width $\mathrm w$ in $u$ has $|\partial_\kappa \mathsf c_\varphi(\kappa)|\lesssim\mathrm w^{-1}e^{-(\mathrm w\kappa)^2/2}$ for Gaussian $\varphi$, so ordinates with $\mathrm w\kappa\gg1$ are invisible to it and the corresponding columns of $\mathbb J_m$ are exponentially small. The underdetermined right inverse then has $\mathbb L_m$ dominated by the invisible columns; in the computations of Section~\CGref{subsec:rb-numerics} this produced $\mathbb L_m$ of size $10^{11}$ within four stages, and the stable regime was the opposite one, $\mathrm t_m\ge k^*_m$ with the unknowns restricted to the visible ordinates (a least-squares regime; the certificate theorem applies to its square subsystems only). The quasi-orthogonality heuristic is therefore a statement about tests that are not in $C_c^\infty(0,\infty)$, and it is not pursued further here.
\end{heuristic}

\subsubsection{Explicit lists in the tube: Gonek's symmetrised Euler products}\label{subsec:rb-gonek}

Gonek \cite{Gonek2012} introduced, for $\mathrm X\ge2$, the truncated Euler product $P_{\mathrm X}(s):=\exp\bigl(\sum_{n\le \mathrm X^2}\frac{\Lambda_{\mathrm X}(n)}{n^s\log n}\bigr)$ with $\Lambda_{\mathrm X}(n):=\Lambda(n)$ for $n\le \mathrm X$ and $\Lambda_{\mathrm X}(n):=\Lambda(n)\log(\mathrm X^2/n)/\log \mathrm X$ for $\mathrm X<n\le \mathrm X^2$, and the symmetrised function $\zeta_{\mathrm X}(s):=P_{\mathrm X}(s)+\chi(s)P_{\mathrm X}(1-s)$, where $\chi$ is the factor of the functional equation $\zeta(s)=\chi(s)\zeta(1-s)$. On the critical line $\chi(\frac12+it)=e^{-2i\vartheta(t)}$ and $P_{\mathrm X}(\frac12-it)=\overline{P_{\mathrm X}(\frac12+it)}$, so with
\begin{equation}\label{eq:rb-alphaX}
 \mathrm p_{\mathrm X}(t):=\arg P_{\mathrm X}\bigl(\tfrac12+it\bigr)=-\sum_{n\le \mathrm X^2}\frac{\Lambda_{\mathrm X}(n)\sin(t\log n)}{\sqrt n\log n},
\end{equation}
and the two constants
\begin{equation}\label{eq:rb-aux-16}
 \mathrm m_{\mathrm X}:=\sum_{n\le \mathrm X^2}\frac{\Lambda_{\mathrm X}(n)}{\sqrt n\log n},\qquad \mathrm n_{\mathrm X}:=\sum_{n\le \mathrm X^2}\frac{\Lambda_{\mathrm X}(n)}{\sqrt n},
\end{equation}
one has $\zeta_{\mathrm X}(\frac12+it)=2|P_{\mathrm X}(\frac12+it)|\,e^{-i\vartheta(t)}\cos\bigl(\vartheta(t)+\mathrm p_{\mathrm X}(t)\bigr)$. Since $P_{\mathrm X}$ has no zeros, the zeros of $\zeta_{\mathrm X}$ on the line are exactly the solutions of
\begin{equation}\label{eq:rb-gonek-phase}
 \varTheta_{\mathrm X}(t):=\vartheta(t)+\mathrm p_{\mathrm X}(t)\in\tfrac\pi2+\pi\mathbb Z .
\end{equation}
Gonek proves \cite[Thm.~6.2 of the arXiv version]{Gonek2012} that every zero of $\zeta_{\mathrm X}$ in $0\le\Re s\le1$ with imaginary part at least an absolute constant $C_0$ lies on $\Re s=\frac12$, unconditionally, and that under the Riemann hypothesis the zeros of $\zeta_{\mathrm X}$ approach those of $\zeta$ as $\mathrm X\to\infty$. The phase \CGeqref{eq:rb-gonek-phase} is the clock phase plus a truncated prime sum: $\NR(t)=1+\vartheta(t)/\pi+\mathrm r_\vartheta(t)$ with $|\mathrm r_\vartheta(t)|\le \mathrm c_r/t$ for $t\ge1$, by Stirling's formula \cite[\S12.33]{WhittakerWatson1927}, so \CGeqref{eq:rb-gonek-phase} reads $\NR(t)+\mathrm p_{\mathrm X}(t)/\pi\in\frac12+\mathbb Z+\mathrm r_\vartheta(t)$, to be compared with $\NR(\tau_k)=k-\frac12$.

\begin{lemma}[Gonek lists lie in a tube around the clock]\label{lem:rb-gonek-tube}
Fix $\mathrm X\ge2$ and put $\mathfrak t_{\mathrm X}:=\max\bigl(2\pi e^{2\mathrm n_{\mathrm X}+2},\,t_1\bigr)$, where $t_1$ is such that $|\vartheta'(t)-\frac12\log\frac t{2\pi}|\le\frac12$ for $t\ge t_1$. Then $\varTheta_{\mathrm X}$ is strictly increasing on $[\mathfrak t_{\mathrm X},\infty)$ with $\varTheta_{\mathrm X}'\ge\frac12$, the solutions of \CGeqref{eq:rb-gonek-phase} in $[\mathfrak t_{\mathrm X},\infty)$ form a strictly increasing sequence $\mathrm t_1^{(\mathrm X)}<\mathrm t_2^{(\mathrm X)}<\cdots$, each a simple zero of $\zeta_{\mathrm X}$, and there is an integer $\mathrm j_{\mathrm X}$ such that the list
\begin{equation}\label{eq:rb-gonek-list}
 \kappa^{(\mathrm X)}_k:=\begin{cases}\tau_k,&k\le \mathrm j_{\mathrm X},\\ \mathrm t^{(\mathrm X)}_{k-\mathrm j_{\mathrm X}},&k>\mathrm j_{\mathrm X},\end{cases}
\end{equation}
is, after the nondecreasing rearrangement of its prefix $k\le\mathrm j_{\mathrm X}$ by Lemma~\CGref{lem:rb-matching} if the junction $k=\mathrm j_{\mathrm X}$ breaks monotonicity, nondecreasing and $\mathrm c_{\mathrm X}$-admissible for some $\mathrm c_{\mathrm X}<\infty$, with
\begin{equation}\label{eq:rb-gonek-tube}
 \bigl|\kappa^{(\mathrm X)}_k-\tau_k\bigr|\le\frac{2\mathrm m_{\mathrm X}+o(1)}{\log(\tau_k/2\pi)}\qquad(k\to\infty).
\end{equation}
\end{lemma}

\begin{proof}
$|\mathrm p_{\mathrm X}'(t)|=|\sum_{n\le \mathrm X^2}\Lambda_{\mathrm X}(n)n^{-1/2}\cos(t\log n)|\le \mathrm n_{\mathrm X}$ and $|\mathrm p_{\mathrm X}|\le \mathrm m_{\mathrm X}$. For $t\ge \mathfrak t_{\mathrm X}$, $\vartheta'(t)\ge\frac12\log\frac t{2\pi}-\frac12\ge \mathrm n_{\mathrm X}+\frac12$, so $\varTheta_{\mathrm X}'\ge\frac12$. Hence $\varTheta_{\mathrm X}$ is a strictly increasing bijection of $[\mathfrak t_{\mathrm X},\infty)$ onto $[\varTheta_{\mathrm X}(\mathfrak t_{\mathrm X}),\infty)$, and the solutions of \CGeqref{eq:rb-gonek-phase} are the preimages of the points $\frac\pi2+\pi j$, $j\ge j_0$, where $j_0$ is the least integer with $\frac\pi2+\pi j_0\ge\varTheta_{\mathrm X}(\mathfrak t_{\mathrm X})$; they are simple zeros because $\frac d{dt}\cos\varTheta_{\mathrm X}=-\varTheta_{\mathrm X}'\sin\varTheta_{\mathrm X}\ne0$ there. The $i$-th solution $\mathrm t^{(\mathrm X)}_i$ satisfies $\varTheta_{\mathrm X}(\mathrm t^{(\mathrm X)}_i)=\frac\pi2+\pi(j_0+i-1)=\pi(j_0+i-\frac12)$; we set $\mathrm j_{\mathrm X}:=j_0+1$, so that $k=\mathrm j_{\mathrm X}+i$ gives $\varTheta_{\mathrm X}(\kappa^{(\mathrm X)}_k)=\pi(k-\frac32)$ for $k>\mathrm j_{\mathrm X}$. Since $\NR=1+\vartheta/\pi+\mathrm r_\vartheta$, the constant $1$ raises the count by one unit and
\begin{equation}\label{eq:rb-aux-15}
 \NR\bigl(\kappa^{(\mathrm X)}_k\bigr)=k-\tfrac12-\frac{\mathrm p_{\mathrm X}(\kappa^{(\mathrm X)}_k)}\pi+\mathrm r_\vartheta\bigl(\kappa^{(\mathrm X)}_k\bigr)\qquad(k>\mathrm j_{\mathrm X}),
\end{equation}
so, since $\NR(\tau_k)=k-\frac12$,
\begin{equation}\label{eq:rb-aux-17}
 \bigl|\NR(\kappa^{(\mathrm X)}_k)-\NR(\tau_k)\bigr|\le\frac{\mathrm m_{\mathrm X}}\pi+\frac{\mathrm c_r}{\kappa^{(\mathrm X)}_k}
\end{equation}
(the choice $\mathrm j_{\mathrm X}=j_0+1$, and not $j_0$, is what makes the two counts agree; in the computations of Section~\CGref{subsec:rb-numerics} the offset was determined by matching). The right-hand side is bounded, and $\NR(2t)-\NR(t)\to\infty$, so $\kappa^{(\mathrm X)}_k/\tau_k\to1$. By the mean value theorem, $|\kappa^{(\mathrm X)}_k-\tau_k|\le(\mathrm m_{\mathrm X}/\pi+o(1))/\NR'(\mathrm s_k)$ with $\mathrm s_k$ between $\kappa^{(\mathrm X)}_k$ and $\tau_k$, and $\NR'(\mathrm s_k)=\frac1{2\pi}\log\frac{\tau_k}{2\pi}(1+o(1))$. This is \CGeqref{eq:rb-gonek-tube}; it gives $|\kappa^{(\mathrm X)}_k-\tau_k|\to0$, so the displacements are bounded, the list is eventually nondecreasing, and finitely many initial indices (where the list is the clock) are harmless. If monotonicity of the concatenated list fails at the junction $k=\mathrm j_{\mathrm X}$, one replaces the prefix $\tau_k$, $k\le \mathrm j_{\mathrm X}$, by the nondecreasing rearrangement of Lemma~\CGref{lem:rb-matching}, which keeps the tube.
\end{proof}

\begin{remark}[What the Gonek lists give and what remains]\label{rem:rb-gonek-status}
The lists $\boldsymbol\kappa^{(\mathrm X)}$ use the primes up to $\mathrm X^2$ and the clock, and nothing about the zeros of $\zeta$. Their prime side is built in: formally, $\sum_k\delta(t-\kappa^{(\mathrm X)}_k)=\frac{\varTheta_{\mathrm X}'(t)}\pi\bigl[1+2\sum_{j\ge1}(-1)^j\cos(2j\varTheta_{\mathrm X}(t))\bigr]$ where $\varTheta_{\mathrm X}'>0$, and the $j=0$ term carries $\mathrm p_{\mathrm X}'/\pi=-\frac1\pi\sum_{n\le \mathrm X^2}\Lambda_{\mathrm X}(n)n^{-1/2}\cos(t\log n)$, whose cosine transform is a sum of masses $\Lambda_{\mathrm X}(n)/(2\sqrt n)$ at $u=\log n$, equal to $\mathsf a^\zeta_n$ for $n\le \mathrm X$. Two statements would make them a solution of B$'$: \textup{(G1)} a tube uniform in $\mathrm X$, $\sup_{\mathrm X\ge\mathrm X_0}\sup_k|\kappa^{(\mathrm X)}_k-\tau_k|\le C$, which by \CGeqref{eq:rb-gonek-tube} concerns the size of $\mathrm p_{\mathrm X}$ at its own crossing points and not its supremum ($\mathrm m_{\mathrm X}\asymp\sqrt{\mathrm X}/\log^2\mathrm X$ by the prime number theorem, and $\mathrm n_{\mathrm X}\asymp\sqrt{\mathrm X}$; the supremum is useless for this); and \textup{(G2)} $\langle\mathfrak R(\boldsymbol\kappa^{(\mathrm X)},\mathsf a^\zeta),\varphi\rangle\to0$ as $\mathrm X\to\infty$ for each test $\varphi$, whose content is the vanishing of the harmonics $j\ge1$ against $\varphi$. Neither is proved here. The monotonicity threshold $\mathfrak t_{\mathrm X}$ of Lemma~\CGref{lem:rb-gonek-tube} grows like $e^{2\mathrm n_{\mathrm X}}$, so the global construction cannot rely on monotonicity of $\varTheta_{\mathrm X}$ at fixed height as $\mathrm X\to\infty$; Gonek's theorem on the location of the zeros does not need it. For fixed $\mathrm X$ the tube constant in \CGeqref{eq:rb-gonek-tube} is $2\mathrm m_{\mathrm X}$, and the proof only uses $\sup_t|\mathrm p_{\mathrm X}(t)|\le\mathrm m_{\mathrm X}$; by hand $\mathrm m_3=0.707+0.577+0.185+0.239+0.087+0.013=1.81$, while the observed suprema of $|\mathrm p_{\mathrm X}|$ on the computed ranges are $1.42$ ($\mathrm X=3$), $3.02$ ($\mathrm X=8$) and $5.98$ ($\mathrm X=20$), which are the constants that matter numerically. Numerically (Section~\CGref{subsec:rb-numerics}) the lists with $3\le \mathrm X\le20$ lie within $0.19$ to $0.035$ mean spacings of the zeta ordinates for $k\le400$, whereas the clock is within $0.66$ spacings.
\end{remark}

%======================================================================
\subsection{An a priori bound for quasi-solutions}\label{subsec:rb-apriori}

Throughout this section $(\boldsymbol\kappa,\mathsf a)$ is a finite section in $\mathbb K(C,\mathbb M)$ with parameters $k^*,n^*$.

\begin{definition}[Laplace residual]\label{def:rb-laplace-residual}
For $\Re s>1$ put
\begin{equation}\label{eq:rb-Rhat}
 \widehat{\mathsf R}(s):=\tfrac12\,\frac{P_{\boldsymbol\kappa}'}{P_{\boldsymbol\kappa}}(s)-\tfrac12\,\frac{G'}{G}(s)+\sum_{n=2}^{n^*}\mathsf a_n\,n^{\frac12-s}.
\end{equation}
For $T\ge0$ put
\begin{equation}\label{eq:rb-E}
 \mathbb E(T):=\int_{-T}^{T}\bigl|\widehat{\mathsf R}(2+iv)\bigr|\,dv,\qquad
 \upsilon:=\sup_{T\ge2}\frac{\mathbb E(T)}{\log T}\in[0,\infty].
\end{equation}
Since $\int_0^\infty e^{-(s-\frac12)u}\cos(\kappa u)\,du=\frac{s-\frac12}{(s-\frac12)^2+\kappa^2}=\frac12\frac d{ds}\log\bigl(1+\frac{q(s)}{\frac14+\kappa^2}\bigr)$ and $\int_0^\infty e^{-(s-\frac12)u}\delta(u-\log n)\,du=n^{\frac12-s}$, $\widehat{\mathsf R}$ is formally the Laplace transform of the residual distribution of \CGeqref{eq:rb-rung}, with the normalisation of $\mathsf k$ fixed in \CGref{eq:transport-defs}; this is made precise in Lemma~\CGref{lem:rb-laplace}. We take \CGeqref{eq:rb-Rhat} as the definition. Because the coefficients are real, $\widehat{\mathsf R}(\bar s)=\overline{\widehat{\mathsf R}(s)}$, so $\mathbb E(T)=2\int_0^T|\widehat{\mathsf R}(2+iv)|\,dv$. We assume $\upsilon<\infty$ throughout this section; if $\upsilon=\infty$ every bound below is void. By Lemma~\CGref{lem:rb-laplace}, $\widehat{\mathsf R}$ is the Laplace transform of the residual distribution of the section.
\end{definition}

\begin{lemma}[When the Laplace residual grows logarithmically]\label{lem:rb-upsilon-infinite}
For a finite section, $\upsilon<\infty$ if and only if $\mathsf a_n=0$ for every $n$. More precisely, with $\mathcal D(v):=\sum_{n=2}^{n^*}\mathsf a_nn^{-3/2-iv}$,
\begin{equation}\label{eq:rb-upsilon-split}
 \mathbb E(T)=\int_{-T}^{T}|\mathcal D(v)|\,dv+O_{C,k^*}(\log T),\quad
 \liminf_{T\to\infty}\frac1{2T}\int_{-T}^{T}|\mathcal D(v)|\,dv\ge\frac{\sum_{n\le n^*}|\mathsf a_n|^2n^{-3}}{\sum_{n\le n^*}|\mathsf a_n|n^{-3/2}}>0\ \ (\mathsf a\ne0).
\end{equation}
\end{lemma}

\begin{proof}
With $z=\frac32+iv$, the spectral part of $\widehat{\mathsf R}(2+iv)$ is $\frac12\zeta_{\RVM}'/\zeta_{\RVM}(2+iv)+\sum_{k\le k^*}\bigl[\frac z{z^2+\kappa_k^2}-\frac z{z^2+\tau_k^2}\bigr]$, because the list equals the clock beyond $k^*$. Each term of the finite sum is $z(\tau_k^2-\kappa_k^2)/((z^2+\kappa_k^2)(z^2+\tau_k^2))=O_{C,k}(|v|^{-3})$, integrable over $\mathbb R$. By Theorem~\CGref{thm:transition}(a) with $B=\frac32$, $\zeta_{\RVM}=\varpi_{\RVM}(1+w)(1+\mathcal E_0)$ on the band $1\le\sigma'\le2$, $t\ge t_0$, with $|w|\ll t^{-3/2}$, $w'=-\log(\tau_*/2\pi)\,w$ (Corollary~\CGref{cor:derivative-zeros}) and $|\mathcal E_0|\le C/t$; Cauchy's estimate on discs of radius $\frac14$ inside the band gives $|\mathcal E_0'|\le C'/t$, so $|\zeta_{\RVM}'/\zeta_{\RVM}(2+iv)|\ll(1+|v|)^{-1}$ and its integral over $[-T,T]$ is $O(\log T)$. This is the first statement of \CGeqref{eq:rb-upsilon-split}. For the second, $\mathcal D$ is a Dirichlet polynomial with distinct frequencies $\log n$, so $\frac1{2T}\int_{-T}^T|\mathcal D|^2dv\to\sum_n|\mathsf a_n|^2n^{-3}$ (Montgomery--Vaughan, or directly: the off-diagonal terms are $O(1/T)$), while $|\mathcal D|\le\sum_n|\mathsf a_n|n^{-3/2}$; hence $\int_{-T}^T|\mathcal D|\ge\int_{-T}^T|\mathcal D|^2/\sup|\mathcal D|$ gives the stated lower bound. If $\mathsf a\ne0$, then $\mathbb E(T)\asymp T$ and $\upsilon=\infty$; if $\mathsf a=0$, then $\mathbb E(T)=O(\log T)$ and $\upsilon<\infty$.
\end{proof}

\begin{remark}[Scope of the counting bound]\label{rem:rb-count-scope}
Lemma~\CGref{lem:rb-upsilon-infinite} restricts Theorem~\CGref{thm:rb-count} and Corollary~\CGref{cor:rb-displacement}, as stated for finite sections, to the sections without masses, for which $|N_{\boldsymbol\kappa}-n_{\RVM}|\le k^*+1$ is immediate. The reason is that on a vertical line the Laplace transform of the prime side is almost periodic and does not decay, so an $L^1$ residual of logarithmic growth on $\Re s=2$ demands that the spectral side reproduce the prime side at all heights, which no finite section does; the natural examples with $\upsilon<\infty$ are infinite lists with the full mass sequence, such as $(\boldsymbol\gamma,\mathsf a^\zeta)$, for which $\widehat{\mathsf R}\equiv0$ and $\upsilon=0$. The theorem is kept as the template of the count-to-displacement mechanism (Lemmas~\CGref{lem:rb-PL}--\CGref{lem:rb-argument}), which is reused in Appendix~\CGref{sec:appendix-endpoint-v18}; its extension to infinite lists needs a polynomial growth bound for $P_{\boldsymbol\kappa}/P_{\boldsymbol\tau}$ on $\Re s\ge2$, which is not given here (Remark~\CGref{rem:rb-tube}), and Heuristic~\CGref{heur:rb-window} explains why the Euclidean residual of a numerical section is a different quantity.
\end{remark}

Constants. Put $\mathrm c_\ell:=2\sum_{n\ge2}n^{-3/2}/\log n<\infty$, and, for $s\in\mathbb C$ and real $x\ge2\pi$, $\mathfrak f(x;s):=\log\bigl(1+\frac{q(s)}{\frac14+x^2}\bigr)$, with the principal branch, which is defined whenever $|s-\frac12|<x$. Let $\mathbb J:=\{2\}\cup\{\sigma+i:\ -1\le\sigma\le2\}$, a compact set on which $|s-\frac12|\le\frac52<2\pi$, and
\begin{equation}\label{eq:rb-beta}
 \mathfrak B(C):=\sum_{k\ge1}\ \sup_{x\in[\max(\tau_k-C,2\pi),\,\tau_k+C]}\ \sup_{s\in\mathbb J}\bigl|\mathfrak f(x;s)-\mathfrak f(\tau_k;s)\bigr| .
\end{equation}
This is finite. On $\mathbb J$, $|q(s)|\le5$ and $|\partial_xf(x;s)|=\bigl|\frac{2xq(s)}{(\frac14+x^2)(x^2+(s-\frac12)^2)}\bigr|\le\frac{20}{x^3}$ for $x\ge4$. The $k$-th term is therefore at most $160C\tau_k^{-3}$ when $\tau_k\ge2C$, and $\sum_k\tau_k^{-3}<\infty$. Further, $\mathrm c_\tau:=\max_{-1\le\sigma\le2}|\log|\zeta_{\RVM}(\sigma+i)||$ and $\mathrm c_{\RVM}:=\sup_{\Re s\ge2}|\log\zeta_{\RVM}(s)|$. Both are finite: the first because $\zeta_{\RVM}$ has no zeros or poles on that segment, the second by \CGref{thm:halfplane-RVM}.

\begin{lemma}[Phragm\'en--Lindel\"of, two forms]\label{lem:rb-PL}
\textup{(a)} Let $\mathbb S:=\{\mathrm l_1\le\Re s\le\mathrm l_2,\ \Im s\ge t_1\}$ have width $a:=\mathrm l_2-\mathrm l_1$. Let $u$ be upper semicontinuous on $\mathbb S$ and subharmonic in its interior, with $u\le \mathrm M_0$ on $\partial\mathbb S$ and $u(s)\le \mathrm A_0e^{c\Im s}$ for some $c<\pi/a$. Then $u\le \mathrm M_0$ on $\mathbb S$.\\
\textup{(b)} Let $\mathfrak g$ be holomorphic in $\Re s>1$, continuous on $\Re s\ge2$, with $|g|\le1$ on $\Re s=2$ and $|\mathfrak g(s)|\le \mathrm A_0(2+|s|)^{\mathrm L}$ on $\Re s\ge2$. Then $|g|\le1$ on $\Re s\ge2$.
\end{lemma}

\begin{proof}
(a) Choose $c<c'<\pi/a$ and $m:=\frac{\mathrm l_1+\mathrm l_2}2$, and put $\omega(s):=e^{c'\Im s}\cos(c'(\Re s-m))$. This is harmonic, being the real part of $e^{-ic'(s-m)}$. On $\mathbb S$, $\omega\ge\cos(c'a/2)\,e^{c'\Im s}$, and $\cos(c'a/2)>0$. For $\mathrm p>0$, $u-\mathrm M_0-\mathrm p\omega$ is subharmonic, $\le0$ on $\partial\mathbb S$, and tends to $-\infty$ as $\Im s\to\infty$. The maximum principle on $\mathbb S\cap\{\Im s\le t_2\}$, with $t_2$ large, gives $u\le \mathrm M_0+\mathrm p\omega$. Let $\mathrm p\to0$.\\
(b) For $\mathrm p>0$, $\mathfrak g_\eta(s):=\mathfrak g(s)\exp(-\mathrm p(s-1)^{1/2})$, with the principal branch, satisfies $|\mathfrak g_\eta|\le|g|$ on $\Re s\ge2$, because $\Re(s-1)^{1/2}\ge\cos(\pi/4)|s-1|^{1/2}>0$ there. Moreover $\mathfrak g_\eta(s)\to0$ as $|s|\to\infty$ in $\Re s\ge2$. The maximum principle on large half-discs gives $|\mathfrak g_\eta|\le1$. Let $\mathrm p\to0$.
\end{proof}

\begin{lemma}[The Dirichlet factor]\label{lem:rb-factor}
Put $\mathcal D_{\mathsf a}(s):=\sum_{n=2}^{n^*}\frac{2\mathsf a_n}{\log n}n^{\frac12-s}$ and $\mathcal Y:=\zeta_{\boldsymbol\kappa}e^{-\mathcal D_{\mathsf a}}$. Then $\mathcal Y$ is holomorphic and zero-free in $\Re s>1$, and $|\mathcal D_{\mathsf a}|\le \mathbb M\mathrm c_\ell$ on $\Re s\ge2$. For the branch of $\log \mathcal Y$ that is real on $(1,\infty)$,
\begin{equation}\label{eq:rb-logY-deriv}
 (\log \mathcal Y)'=2\,\widehat{\mathsf R},\qquad
 |\log \mathcal Y(2)|\le \mathbb B_0:=|\log\zeta_{\RVM}(2)|+\mathfrak B(C)+\mathbb M\mathrm c_\ell,
\end{equation}
and, for all real $t$,
\begin{equation}\label{eq:rb-line}
 |\log \mathcal Y(2+it)|\le \mathbb B_0+\mathbb E(|t|)\le \mathbb B_0+\upsilon\log(2+|t|).
\end{equation}
\end{lemma}

\begin{proof}
$P_{\boldsymbol\kappa}$ vanishes only on $\Re s=\frac12$, and $G$ is holomorphic and zero-free in $\Re s>1$. Hence $\zeta_{\boldsymbol\kappa}$ and $\mathcal Y$ are holomorphic and zero-free there. $|\mathcal D_{\mathsf a}(s)|\le\sum_n2Mn^{\frac12-\Re s}/\log n\le \mathbb M\mathrm c_\ell$ for $\Re s\ge2$. Differentiating, $(\log \mathcal Y)'=\frac{P'_{\boldsymbol\kappa}}{P_{\boldsymbol\kappa}}-\frac{G'}G+2\sum_n\mathsf a_nn^{\frac12-s}=2\widehat{\mathsf R}$. At $s=2$, $\log\zeta_{\boldsymbol\kappa}(2)=\log\zeta_{\RVM}(2)+\sum_{k\le k^*}(\mathfrak f(\kappa_k;2)-\mathfrak f(\tau_k;2))$, a finite sum of real numbers bounded by $\mathfrak B(C)$, and $|\mathcal D_{\mathsf a}(2)|\le \mathbb M\mathrm c_\ell$. Finally, $\log \mathcal Y(2+it)-\log \mathcal Y(2)=2i\int_0^t\widehat{\mathsf R}(2+iv)\,dv$ has modulus at most $\mathbb E(|t|)$. For $|t|\ge2$, $\mathbb E(|t|)\le\upsilon\log|t|$; for $|t|<2$, $\mathbb E(|t|)\le \mathbb E(2)\le\upsilon\log2$.
\end{proof}

\begin{lemma}[Bound on the half-plane]\label{lem:rb-halfplane}
With $\mathbb W:=\log|\zeta_{\boldsymbol\kappa}|=\log|P_{\boldsymbol\kappa}|-\log|G|$ and $\mathbb B_1:=\mathbb B_0+\upsilon\log(2\sqrt2)$, $\mathbb B_2:=\mathbb B_1+\mathbb M\mathrm c_\ell$, for $\Re s\ge2$,
\begin{equation}\label{eq:rb-halfplane}
 \bigl|\log|\mathcal Y(s)|\bigr|\le \mathbb B_1+\upsilon\log|s-1|,\qquad |\mathbb W(s)|\le \mathbb B_2+\upsilon\log|s-1| .
\end{equation}
\end{lemma}

\begin{proof}
On $\Re s=2$, $|s-1|=\sqrt{1+t^2}\ge(2+|t|)/(2\sqrt2)$, so \CGeqref{eq:rb-line} gives $|\mathcal Y(s)|^{\pm1}\le e^{B_1}|s-1|^{\upsilon}$. A priori growth: on $\Re s\ge2$, $|\log|\zeta_{\RVM}||\le \mathrm c_{\RVM}$. The quotient $\zeta_{\boldsymbol\kappa}/\zeta_{\RVM}$ is a finite product of $k^*$ factors $\frac{(\kappa_k^2+w^2)(\frac14+\tau_k^2)}{(\frac14+\kappa_k^2)(\tau_k^2+w^2)}$ with $w=s-\frac12$. Since $|\kappa^2+w^2|=|\kappa-\Im s+i\Re w|\,|\kappa+\Im s-i\Re w|\ge(\Re w)^2\ge\frac94$, each factor and its inverse are $O((2+|s|)^4)$. Hence $|\mathcal Y|^{\pm1}\le \mathrm A_0(2+|s|)^{4k^*}$. Apply Lemma~\CGref{lem:rb-PL}(b) to $\mathfrak g_+:=\mathcal Y\,(s-1)^{-\upsilon}e^{-\mathbb B_1}$ and to $\mathfrak g_-:=\mathcal Y^{-1}(s-1)^{-\upsilon}e^{-\mathbb B_1}$, with the principal branch of $(s-1)^{-\upsilon}$, holomorphic in $\Re s>1$ with $|(s-1)^{-\upsilon}|=|s-1|^{-\upsilon}$; both are holomorphic in $\Re s>1$ and of polynomial growth on $\Re s\ge2$. Finally $\mathbb W=\log|\mathcal Y|+\Re\mathcal D_{\mathsf a}$.
\end{proof}

\begin{lemma}[Bound on the half-strip]\label{lem:rb-strip}
Let $\mathbb S:=\{-1\le\Re s\le2,\ \Im s\ge1\}$ and $\mathbb B_3:=\max\bigl(\mathbb B_2+\frac34\log2,\ \mathrm c_\tau+\mathfrak B(C)\bigr)$. Then
\begin{equation}\label{eq:rb-strip}
 \mathbb W(s)\le \mathbb B_3+\bigl(\tfrac32+\upsilon\bigr)\log|s+2|\qquad(s\in\mathbb S).
\end{equation}
\end{lemma}

\begin{proof}
$\mathbb W$ is subharmonic in the interior of $\mathbb S$, since $\log|G|$ is harmonic there ($G$ has no zeros or poles with $\Im s\ge1$), and it is upper semicontinuous on $\mathbb S$. The majorant $\mathrm h_0(s):=(\frac32+\upsilon)\log|s+2|$ is harmonic and nonnegative on $\mathbb S$. We check the three sides.

\emph{Right side.} By \CGeqref{eq:rb-halfplane}, $\mathbb W(2+it)\le \mathbb B_2+\upsilon\log|1+it|\le \mathbb B_2+\mathrm h_0(2+it)$.

\emph{Left side.} $P_{\boldsymbol\kappa}$ and $G$ are real, so $\mathbb W(\bar s)=\mathbb W(s)$. Combined with $P_{\boldsymbol\kappa}(1-s)=P_{\boldsymbol\kappa}(s)$, this gives $\mathbb W(-1+it)=\mathbb W(2+it)+\log|G(2+it)/G(-1+it)|$. By $\Gamma(z+1)=z\Gamma(z)$ and the reflection formula, $|\Gamma(\frac12+iy)|^2=\pi/\cosh\pi y$ and $|\Gamma(1+iy)|^2=\pi y/\sinh\pi y$, so
\begin{equation}\label{eq:rb-G-ratio}
 \Bigl|\frac{G(2+it)}{G(-1+it)}\Bigr|=\frac{|1+it|}2\,\pi^{-3/2}\Bigl(\frac t2\coth\frac{\pi t}2\Bigr)^{1/2}\le(1+t)^{3/2}\qquad(t\ge1).
\end{equation}
For the inequality: $|1+it|\le1+t$, $\frac t2\coth\frac{\pi t}2\le\frac{1+t}2\coth\frac\pi2$, and $\frac12\pi^{-3/2}(\frac12\coth\frac\pi2)^{1/2}<1$. Since $(1+t)\le\sqrt2|1+it|$ and $|s+2|=|1+it|$ at $s=-1+it$, we get $\mathbb W(-1+it)\le \mathbb B_2+\frac34\log2+\mathrm h_0(-1+it)$.

\emph{Bottom side.} For $s=\sigma+i$, $\mathbb W(s)=\log|\zeta_{\RVM}(s)|+\sum_{k\le k^*}\Re(\mathfrak f(\kappa_k;s)-\mathfrak f(\tau_k;s))\le \mathrm c_\tau+\mathfrak B(C)\le \mathrm c_\tau+\mathfrak B(C)+\mathrm h_0(s)$.

\emph{Growth.} $\log|P_{\boldsymbol\kappa}(s)|\le\sum_k\log(1+|q(s)|/\kappa_k^2)$. Split at $\kappa_k\le2|s|$ and $\kappa_k>2|s|$: the first part has $N_{\boldsymbol\kappa}(2|s|)=O(|s|\log|s|)$ terms, each at most $\log(1+|q(s)|/4\pi^2)=O(\log|s|)$; in the second, $\log(1+|q(s)|/\kappa_k^2)\le|q(s)|\kappa_k^{-2}$, and $\sum_{\kappa_k>2|s|}\kappa_k^{-2}=O(\log|s|/|s|)$ by partial summation with $N_{\boldsymbol\kappa}(t)=O(t\log t)$. Hence $\log|P_{\boldsymbol\kappa}(s)|=O(|s|\log^2|s|)$, and $-\log|G(s)|\le\frac\pi4|\Im s|+O(\log|s|)$ on $\mathbb S$ by Stirling's formula \cite[\S12.33]{WhittakerWatson1927}. Hence $\mathbb W-\mathrm h_0\le \mathrm A_0e^{c\Im s}$ for any $c>0$. Lemma~\CGref{lem:rb-PL}(a), with width $3$ and $c<\pi/3$, applied to $\mathbb W-\mathrm h_0$ gives the claim.
\end{proof}

\begin{lemma}[Backlund--Jensen on a horizontal segment]\label{lem:rb-backlund}
There are absolute constants $T_0\ge3$ and $\mathrm c_5$ such that, for $T\ge T_0$ with $T\notin\{\kappa_k\}$, the variation of $\arg P_{\boldsymbol\kappa}$ along the segment from $2+iT$ to $\frac12+iT$ is at most $\pi(n_T+1)$, where
\begin{equation}\label{eq:rb-nu}
 n_T\le\frac{(\frac52+\upsilon)\log(T+8)+\mathbb E(T)+\mathbb B_0+\mathbb B_3+\mathbb M\mathrm c_\ell+2\mathrm c_5}{\log(4/3)} .
\end{equation}
\end{lemma}

\begin{proof}
Choose $\mathrm a_T\in\mathbb R$ with $e^{i\mathrm a_T}P_{\boldsymbol\kappa}(2+iT)>0$ and put $\mathfrak F_T(z):=\frac12\bigl(e^{i\mathrm a_T}P_{\boldsymbol\kappa}(z+iT)+e^{-i\mathrm a_T}P_{\boldsymbol\kappa}(z-iT)\bigr)$, which is entire. Since $P_{\boldsymbol\kappa}$ is real, $P_{\boldsymbol\kappa}(x-iT)=\overline{P_{\boldsymbol\kappa}(x+iT)}$ for real $x$, so $\mathfrak F_T(x)=\Re(e^{i\mathrm a_T}P_{\boldsymbol\kappa}(x+iT))$ and $\mathfrak F_T(2)=|P_{\boldsymbol\kappa}(2+iT)|$. If $\mathfrak F_T$ has $n_T$ zeros on $[\frac12,2]$, the segment splits into at most $n_T+1$ pieces on which $\Re(e^{i\mathrm a_T}P_{\boldsymbol\kappa})$ has constant sign. On each piece $\arg$ varies by less than $\pi$, which gives the first claim \cite[\S9.4]{Titchmarsh1986}.

The segment $[\frac12,2]$ lies in the closed disc $|z-2|\le\frac32$. For $\frac32<r<2$, Jensen's formula on $|z-2|\le R=2$ gives $n_T\log\frac Rr\le\max_{|z-2|=2}\log|\mathfrak F_T(z)|-\log|\mathfrak F_T(2)|$, because the zeros in the closed disc of radius $r$, which include the $n_T$ zeros on the segment, contribute at least $n_T\log(R/r)$; letting $r\downarrow\frac32$ gives the inequality with $\log\frac43$. Moreover $|\mathfrak F_T(z)|\le\max_{|\zeta-2|\le2}|P_{\boldsymbol\kappa}(\zeta+iT)|$, because $|P_{\boldsymbol\kappa}(z-iT)|=|P_{\boldsymbol\kappa}(\bar z+iT)|$. Write $\log|P_{\boldsymbol\kappa}|=\mathbb W+\log|G|$. On the disc $\{|s-(2+iT)|\le2\}$, the part with $\Re s\le2$ lies in $\mathbb S$ and $|s+2|\le T+6$, so Lemma~\CGref{lem:rb-strip} gives $\mathbb W\le \mathbb B_3+(\frac32+\upsilon)\log(T+6)$. The part with $\Re s\ge2$ satisfies $\mathbb W\le \mathbb B_2+\upsilon\log(T+5)$ by Lemma~\CGref{lem:rb-halfplane}, a smaller bound. At the centre, \CGeqref{eq:rb-line} gives $\mathbb W(2+iT)\ge-(\mathbb B_0+\mathbb E(T)+\mathbb M\mathrm c_\ell)$. For the oscillation of $\log|G|$ on the disc: by Stirling's formula for $\Gamma'/\Gamma$ in $|\arg z|\le\frac34\pi$ \cite[\S12.33]{WhittakerWatson1927}, $|G'/G(s)|=|\frac1s+\frac1{s-1}-\frac12\log\pi+\frac12\frac{\Gamma'}{\Gamma}(\frac s2)|\le\frac12\log(T+8)+\mathrm c_5$ on the disc for $T\ge T_0$. So the oscillation is at most $\log(T+8)+2\mathrm c_5$. Collecting terms gives \CGeqref{eq:rb-nu}.
\end{proof}

\begin{lemma}[Argument principle]\label{lem:rb-argument}
For $T>0$, $T\notin\{\kappa_k\}$, let $\mathbb G_T$ be the path $2\to2+iT\to\frac12+iT$. Then
\begin{equation}\label{eq:rb-arg-principle}
 \pi N_{\boldsymbol\kappa}(T)=\Delta_{\mathbb G_T}\arg P_{\boldsymbol\kappa}
 =\pi+\vartheta(T)+\Delta_{\mathbb G_T}\arg\zeta_{\boldsymbol\kappa},
\end{equation}
where $\vartheta(T)=\arg\bigl(\pi^{-iT/2}\Gamma(\frac14+\frac{iT}2)\bigr)$ is the Riemann--Siegel theta function, defined by continuous variation from $\vartheta(0)=0$.
\end{lemma}

\begin{proof}
Take the rectangle $\mathfrak R_T$ with vertices $2\pm iT$, $-1\pm iT$. $P_{\boldsymbol\kappa}$ has no zeros on $\partial\mathfrak R_T$: its zeros lie on $\Re s=\frac12$, and $T\notin\{\kappa_k\}$, so $\frac12\pm iT$ are not zeros. Its zeros inside $\mathfrak R_T$ are $\frac12\pm i\kappa_k$ with $\kappa_k<T$, counted with multiplicity (coincident $\kappa_j=\kappa_k$ give a multiple zero), $2N_{\boldsymbol\kappa}(T)$ in total. Reflection $s\mapsto1-\bar s$ maps $\mathbb G_T$ to the path $-1\to-1+iT\to\frac12+iT$. Since $P_{\boldsymbol\kappa}(1-\bar s)=\overline{P_{\boldsymbol\kappa}(s)}$, the variation of $\arg P_{\boldsymbol\kappa}$ along the reversed image path $\frac12+iT\to-1+iT\to-1$ equals $\Delta_{\mathbb G_T}\arg P_{\boldsymbol\kappa}$. The two lower quarters are handled in the same way with $s\mapsto\bar s$. The four quarter paths make up $\partial\mathfrak R_T$, so $4\Delta_{\mathbb G_T}\arg P_{\boldsymbol\kappa}=2\pi\cdot2N_{\boldsymbol\kappa}(T)$.

For the second equality, write $P_{\boldsymbol\kappa}=G\zeta_{\boldsymbol\kappa}$. Along $\mathbb G_T$, $s(s-1)$ moves from $2$ to $-(\frac14+T^2)$ inside the closed upper half-plane, so its argument varies by $\pi$. $\pi^{-s/2}\Gamma(s/2)$ is holomorphic and zero-free in $\Re s>0$, so its argument variation along $\mathbb G_T$ equals that along $[2,\frac12]\cup[\frac12,\frac12+iT]$, which is $0+\vartheta(T)$.
\end{proof}

\begin{theorem}[Counting bound for quasi-solutions]\label{thm:rb-count}
There are absolute constants $T_0$ and $\mathrm c_5$ such that, for every finite section in $\mathbb K(C,\mathbb M)$ and every $T\ge T_0$,
\begin{equation}\label{eq:rb-count}
 \bigl|N_{\boldsymbol\kappa}(T)-n_{\RVM}(T)\bigr|\le(\mathrm c_1+\mathrm c_2\upsilon)\log(T+8)+\mathrm c_3\mathbb B+\mathrm c_4,
\end{equation}
with
\begin{equation}\label{eq:rb-a-constants}
 \mathrm c_1=\frac5{2\log(4/3)}=8.690\ldots,\quad \mathrm c_2=\frac2{\log(4/3)}+\frac1\pi=7.270\ldots,\quad \mathrm c_3=\frac1{\log(4/3)}+\frac1\pi=3.794\ldots,
\end{equation}
$\mathrm c_4:=\frac32+\frac{\mathrm c_6}\pi+\frac1{T_0}+\frac{2\mathrm c_5}{\log(4/3)}$ with $\mathrm c_5$ the constant of Lemma~\CGref{lem:rb-backlund} and $\mathrm c_6:=\frac32$, and $\mathbb B:=\mathbb B_0+\mathbb B_3+\mathbb M\mathrm c_\ell$. The constants $\mathrm c_1,\dots,\mathrm c_6$ do not depend on $k^*$, $n^*$, $C$, $\mathbb M$ or $\upsilon$; $\mathbb B$ depends on $C$ (through $\mathfrak B(C)$), on $\mathbb M$, and linearly on $\upsilon$ (through $\mathbb B_2$), so the total coefficient of $\upsilon$ in \CGeqref{eq:rb-count} is $\mathrm c_2\log(T+8)+\mathrm c_3\log(2\sqrt2)$. In Corollary~\CGref{cor:rb-displacement} the parameters $C$, $\mathbb M$, $\upsilon$ are fixed before $k\to\infty$.
\end{theorem}

\begin{proof}
By continuity from the right it suffices to take $T\notin\{\kappa_k\}$. By Lemma~\CGref{lem:rb-argument}, $N_{\boldsymbol\kappa}(T)-1-\vartheta(T)/\pi=\pi^{-1}\Delta_{\mathbb G_T}\arg\zeta_{\boldsymbol\kappa}$.

\emph{Vertical piece.} $\arg\zeta_{\boldsymbol\kappa}=\Im\log \mathcal Y+\Im\mathcal D_{\mathsf a}$, with $\Im\log \mathcal Y(2)=\Im\mathcal D_{\mathsf a}(2)=0$. Its variation is at most $\mathbb E(T)+\mathbb M\mathrm c_\ell$ by Lemma~\CGref{lem:rb-factor}.

\emph{Horizontal piece.} The variation is $\Delta\arg P_{\boldsymbol\kappa}-\Delta\arg G$. The first is at most $\pi(n_T+1)$ by Lemma~\CGref{lem:rb-backlund}. The second is $-\Im\int_{1/2}^2(G'/G)(\sigma+iT)\,d\sigma$. Here $|\Im(G'/G)|\le|\frac1s|+|\frac1{s-1}|+\frac12|\arg\frac s2|+O(1/T)\le1$ for $T\ge T_0$, by Stirling's formula for $\Gamma'/\Gamma$ \cite[\S12.33]{WhittakerWatson1927}. So the second variation is at most $\mathrm c_6:=\frac32$.

\emph{Main term.} By the Stirling expansion of $\vartheta$ \cite[\S12.33]{WhittakerWatson1927}, $1+\vartheta(T)/\pi=\NR(T)+O(1/T)$, and this error is at most $1/T$ for $T\ge T_0$. With \CGeqref{eq:rb-Ntau-floor}:
\begin{equation}\label{eq:rb-aux-1}
 |N_{\boldsymbol\kappa}(T)-n_{\RVM}(T)|\le n_T+1+\frac{\mathbb E(T)+\mathbb M\mathrm c_\ell+\mathrm c_6}{\pi}+\frac1T+\frac12 .
\end{equation}
Insert \CGeqref{eq:rb-nu} and $\mathbb E(T)\le\upsilon\log(T+8)$: the coefficient of $\log(T+8)$ is $\frac{5/2+2\upsilon}{\log(4/3)}+\frac\upsilon\pi=\mathrm c_1+\mathrm c_2\upsilon$, the terms in $\mathbb B_0+\mathbb B_3+\mathbb M\mathrm c_\ell$ are at most $(\frac1{\log(4/3)}+\frac1\pi)\mathbb B=\mathrm c_3\mathbb B$, and the remaining constant is $\frac{2\mathrm c_5}{\log(4/3)}+1+\frac{\mathrm c_6}\pi+\frac1{T_0}+\frac12=\mathrm c_4$.
\end{proof}

\begin{corollary}[Displacement bound]\label{cor:rb-displacement}
Assume $T_0\ge2\pi e$. For every $k$ with $\tau_k-C\ge T_0$,
\begin{equation}\label{eq:rb-disp}
 |\kappa_k-\tau_k|\le\frac{2\pi\bigl[(\mathrm c_1+\mathrm c_2\upsilon)\log(\tau_k+C+8)+\mathrm c_3\mathbb B+\mathrm c_4\bigr]}{\log\bigl((\tau_k-C)/2\pi\bigr)} .
\end{equation}
In particular $\limsup_{k\to\infty}|\kappa_k-\tau_k|\le2\pi(\mathrm c_1+\mathrm c_2\upsilon)$ uniformly over sections with given $C,\mathbb M,\upsilon$. For every $\mathrm p>0$ there is $k_0(C,\mathbb M,\upsilon,\mathrm p)$ such that $|\kappa_k-\tau_k|\le2\pi(\mathrm c_1+\mathrm c_2\upsilon)+\mathrm p$ for $k\ge k_0$.
\end{corollary}

\begin{proof}
Let $\mathbb D(T)$ be the right-hand side of \CGeqref{eq:rb-count}; it is nondecreasing. Suppose $\kappa_k>\tau_k+\mathrm q$ with $0<\mathrm q\le C$, and let $T:=\tau_k+\mathrm q-\mathrm h'$ with small $\mathrm h'>0$. The list is nondecreasing, so $N_{\boldsymbol\kappa}(T)\le k-1$. By \CGeqref{eq:rb-Ntau-floor} and because $\NR'$ is increasing, $n_{\RVM}(T)\ge\NR(T)-\frac12\ge k-1+(\mathrm q-\mathrm h')\NR'(\tau_k)$. Hence $(\mathrm q-\mathrm h')\frac1{2\pi}\log\frac{\tau_k}{2\pi}\le\mathbb D(\tau_k+C)$. Now suppose $\kappa_k<\tau_k-\mathrm q$, and let $T:=\tau_k-\mathrm q+\mathrm h'$. Then $N_{\boldsymbol\kappa}(T)\ge k$ and $n_{\RVM}(T)\le\NR(T)+\frac12\le k-(\mathrm q-\mathrm h')\NR'(\tau_k-C)$. Let $\mathrm h'\to0$ in both cases.
\end{proof}

\begin{remark}[What Corollary~\ref{cor:rb-displacement} does not give]\label{rem:rb-tube}
The proof uses the imposed tube $|\kappa_k-\tau_k|\le C$ twice: through $\mathfrak B(C)$ in $\mathbb B$, and to replace $\mathbb D(\tau_k+\mathrm q)$ by $\mathbb D(\tau_k+C)$. The conclusion is that, inside any finite tube, the coordinates of sufficiently high index lie in the smaller tube of radius $2\pi(\mathrm c_1+\mathrm c_2\upsilon)+\mathrm p$. So the constraint is inactive there once $C>2\pi(\mathrm c_1+\mathrm c_2\upsilon)+\mathrm p$. For the indices $k<k_0(C,\mathbb M,\upsilon,\mathrm p)$ nothing is proved, and the corollary does not show that unconstrained near-minimisers lie in a universal tube. The proof uses finite sections: Lemma~\CGref{lem:rb-halfplane} needs the polynomial growth $\mathrm A_0(2+|s|)^{4k^*}$. Extending it to infinite lists requires a separate growth bound for $P_{\boldsymbol\kappa}/P_{\boldsymbol\tau}$ on $\Re s\ge2$, which is not given here. The constant $2\pi \mathrm c_1=54.60\ldots$ is far from optimal.
\end{remark}

\begin{heuristic}[What a finite section controls]\label{heur:rb-window}
The quantity in Theorem~\CGref{thm:rb-count} is $\mathbb E(T)$, the $L^1$ norm on $\Re s=2$ of the Laplace transform of the residual. It is not the Euclidean residual that a numerical section minimises over finitely many tests. Let the tests be $\mathsf j_0(u)\varphi(u)$ with a fixed window $\mathsf j_0\in C_c^\infty(0,\infty)$, $\mathsf j_0\equiv1$ on $[u_0,u^*]$, and $\varphi$ band-limited to $|\omega|\le\mathrm w^\sharp$. Then the kernel $\mathsf j_0(u)e^{-\frac32u}e^{-ivu}$ is, approximately, in the span of the tests only for $|v|\lesssim\mathrm w^\sharp$. The residual on $(0,u_0)$ is then not controlled, and for a list at distance $C$ from the clock it is expected to be of size $C\log(1/u)$ as $u\to0$. Its contribution to $\mathbb E(T)$ is then $O(CTu_0\log(1/u_0))$. Since $u_0\gtrsim1/\mathrm w^\sharp$ in a band-limited section, Theorem~\CGref{thm:rb-count} would apply with a given $\upsilon_0$ only at heights $T\lesssim\upsilon_0\mathrm w^\sharp/(C\log\mathrm w^\sharp)$. None of these size estimates is proved here. They indicate which norm a section must make small. Lemma~\CGref{lem:rb-factor} shows that this norm is the right one.
\end{heuristic}

\begin{heuristic}[Informative part of a section]\label{heur:rb-informative}
Where the atoms are denser than the resolution of the tests, $e^{-u}\ll\pi/\mathrm w^\sharp$, the masses alone can approximate a smooth target by a Riemann sum $\mathsf a_n\approx \mathfrak g(\log n)/n$ without moving any ordinate. That part of the window then carries no information on the ordinates. Counting unknowns against equations in the remaining window, $u^*\le\log\mathrm w^\sharp+O(1)$, gives an overdetermined system. This proves nothing about compatibility. It only explains why B is an existence statement and not a dimension count.
\end{heuristic}

%======================================================================
\subsection{A Kadec--Avdonin criterion with blocks of variable length}\label{subsec:rb-riesz}

\subsubsection{Discrepancy and generating function}\label{subsec:rb-generating}

\begin{definition}[Discrepancy, Poisson and conjugate operators]\label{def:rb-discrepancy}
Let $\mathscr X=(\mathrm x_k)_{k\in\mathbb Z}$ be a real nondecreasing sequence with
\begin{equation}\label{eq:rb-H}
 |\mathrm x_k-(k-\tfrac12)|\le \mathrm c_{\mathscr X}\log(2+|k|)\qquad(k\in\mathbb Z).
\end{equation}
Its \emph{discrepancy} with respect to the half-integers is
\begin{equation}\label{eq:rb-psi}
 \varPsi_{\mathscr X}(\mathrm x):=\sum_{k\in\mathbb Z}\bigl(\mathbf 1_{[\mathrm x_k,\infty)}(\mathrm x)-\mathbf 1_{[k-\frac12,\infty)}(\mathrm x)\bigr).
\end{equation}
The $k$-th term is $\pm1$ on the interval between $\mathrm x_k$ and $k-\frac12$ and $0$ elsewhere. By \CGeqref{eq:rb-H}, at most $O(\log(2+|\mathrm x|))$ terms are nonzero at $\mathrm x$, so $|\varPsi_{\mathscr X}(\mathrm x)|\le \mathrm c'_{\mathscr X}\log(2+|\mathrm x|)$ and $\varPsi_{\mathscr X}\in L^1(d\mathrm x/(1+\mathrm x^2))$. For $\mathrm g\in L^1(d\mathrm x/(1+\mathrm x^2))$ and $\mathrm y>0$ put
\begin{equation}\label{eq:rb-PQ}
 \mathbb P_{\mathrm y}\mathrm g(\mathrm x):=\frac1\pi\int\frac{\mathrm y\,\mathrm g(\mathrm s)\,d\mathrm s}{(\mathrm x-\mathrm s)^2+\mathrm y^2},\qquad
 \mathbb Y_{\mathrm y}\mathrm g(\mathrm x):=\frac1\pi\int \mathrm g(\mathrm s)\Bigl[\frac{\mathrm x-\mathrm s}{(\mathrm x-\mathrm s)^2+\mathrm y^2}+\frac{\mathrm s}{1+\mathrm s^2}\Bigr]d\mathrm s .
\end{equation}
The last term in $\mathbb Y_{\mathrm y}$ only makes the integral converge. Changing it changes $\mathbb Y_{\mathrm y}\mathrm g$ by a constant, which is irrelevant for all statements below.
\end{definition}

\begin{lemma}[Generating function]\label{lem:rb-generating}
Let $\mathscr X=(\mathrm x_k)_{k\in\mathbb Z}$ satisfy \CGeqref{eq:rb-H}, understood as a sequence indexed by $\mathbb Z$ in which a point of multiplicity $m$ occurs $m$ times. Let $\mathcal F_0(\mathfrak z):=\cos\pi \mathfrak z$, and fix $\mathfrak z_0\in\mathbb R\setminus(\mathscr X\cup(\mathbb Z+\frac12))$. The series
\begin{equation}\label{eq:rb-aux-2}
 \mathfrak G_{\mathscr X}(\mathfrak z):=\sum_{k\in\mathbb Z}\Bigl(\frac1{\mathfrak z-\mathrm x_k}-\frac1{\mathfrak z-k+\frac12}\Bigr)
\end{equation}
converges absolutely and locally uniformly off $\mathscr X\cup(\mathbb Z+\frac12)$, and
\begin{equation}\label{eq:rb-aux-3}
 \mathcal F_{\mathscr X}(\mathfrak z):=\mathcal F_0(\mathfrak z)\exp\int_{\mathfrak z_0}^{\mathfrak z}\mathfrak G_{\mathscr X}(\mathrm z)\,d\mathrm z
\end{equation}
is a real entire function whose zero set is $\mathscr X$, with multiplicities. Moreover:
\begin{enumerate}[label=\textup{(\roman*)}]
\item for $\mathrm y>0$ there is a constant $\mathrm c_{\mathscr X}''$ such that
\begin{equation}\label{eq:rb-F-formula}
 \log|\mathcal F_{\mathscr X}(\mathrm x+iy)|=\log|\mathcal F_0(\mathrm x+iy)|+\pi\,\mathbb Y_{\mathrm y}\varPsi_{\mathscr X}(\mathrm x)+\mathrm c_{\mathscr X}'' ;
\end{equation}
\item $\mathcal F_{\mathscr X}$ is of exponential type $\pi$, and its indicator diagram is the segment $[-i\pi,i\pi]$;
\item if $\mathscr X$ is symmetric, $\mathrm x_{1-k}=-\mathrm x_k$, and $\mathfrak z_0=0\notin \mathscr X$, then $\mathcal F_{\mathscr X}$ is even.
\end{enumerate}
\end{lemma}

\begin{proof}
Each term of $\mathfrak G_{\mathscr X}$ equals $\frac{x_k-k+\frac12}{(\mathfrak z-\mathrm x_k)(\mathfrak z-k+\frac12)}=O(\log(2+|k|)/k^2)$ locally uniformly, so the series converges. $\mathfrak G_{\mathscr X}$ is meromorphic with simple poles, with residue $+m$ at a point of $\mathscr X$ of multiplicity $m$ and $-1$ at each half-integer; coincident poles combine their residues. The residues are integers, so $\exp\int_{\mathfrak z_0}^{\mathfrak z}\mathfrak G_{\mathscr X}$ is single-valued. It has a pole of order one at each half-integer not in $\mathscr X$, cancelled by the zero of $\mathcal F_0$. Hence $\mathcal F_{\mathscr X}$ is entire with zero set $\mathscr X$. $\mathfrak G_{\mathscr X}(\bar{\mathfrak z})=\overline{\mathfrak G_{\mathscr X}(\mathfrak z)}$ and $\mathfrak z_0$ is real, so $\mathcal F_{\mathscr X}$ is real on $\mathbb R$. Under the symmetry in (iii), $\mathfrak G_{\mathscr X}$ is odd, so $\int_0^{\mathfrak z}\mathfrak G_{\mathscr X}$ is even, and $\mathcal F_{\mathscr X}$ is even.

(i) For $\Im \mathfrak z\ne0$ put $\mathbb A_{\mathscr X}(\mathfrak z):=\int\varPsi_{\mathscr X}(\mathrm s)\bigl[\frac1{\mathfrak z-\mathrm s}+\frac{\mathrm s}{1+\mathrm s^2}\bigr]d\mathrm s$. For $f(\mathrm s):=\frac1{\mathfrak z-\mathrm s}$, the identity $f(\mathrm x_k)-f(k-\frac12)=-\int f'(\mathrm s)\bigl(\mathbf 1_{[\mathrm x_k,\infty)}-\mathbf 1_{[k-\frac12,\infty)}\bigr)(\mathrm s)\,d\mathrm s$ summed over $k$ (Fubini applies because $\sum_k\int|f'||\cdots|<\infty$ by \CGeqref{eq:rb-H}) gives $\mathfrak G_{\mathscr X}(\mathfrak z)=-\int\varPsi_{\mathscr X}(\mathrm s)(\mathfrak z-\mathrm s)^{-2}d\mathrm s=\mathbb A_{\mathscr X}'(\mathfrak z)$. So $\mathcal F_{\mathscr X}=\mathrm c_+\mathcal F_0e^{\mathbb A_{\mathscr X}}$ on $\Im \mathfrak z>0$, with $\mathrm c_+\ne0$. Since $\Re\frac1{\mathfrak z-\mathrm s}=\frac{\mathrm x-\mathrm s}{(\mathrm x-\mathrm s)^2+\mathrm y^2}$, $\Re\mathbb A_{\mathscr X}(\mathrm x+iy)=\pi\mathbb Y_{\mathrm y}\varPsi_{\mathscr X}(\mathrm x)$ and $\Im\mathbb A_{\mathscr X}(\mathrm x+iy)=-\pi \mathbb P_{\mathrm y}\varPsi_{\mathscr X}(\mathrm x)$.

(ii) Split $\mathbb Y_{\mathrm y}\varPsi_{\mathscr X}(\mathrm x)$ into $|\mathrm s-\mathrm x|\le1$, $1<|\mathrm s-\mathrm x|$ with $|\mathrm s|\le2|\mathfrak z|+2$, and $|\mathrm s|>2|\mathfrak z|+2$. Using $|\varPsi_{\mathscr X}(\mathrm s)|\le \mathrm c'_{\mathscr X}\log(2+|\mathrm s|)$, $\int_{|u|\le1}\frac{|u|\,du}{u^2+\mathrm y^2}=\log(1+\mathrm y^{-2})$, and, for $|\mathrm s|>2|\mathfrak z|+2$, $|\frac{\mathrm x-\mathrm s}{(\mathrm x-\mathrm s)^2+\mathrm y^2}+\frac{\mathrm s}{1+\mathrm s^2}|\le8(1+|\mathfrak z|)/\mathrm s^2$, one gets
\begin{equation}\label{eq:rb-Q-growth}
 |\mathbb Y_{\mathrm y}\varPsi_{\mathscr X}(\mathrm x)|\le \mathrm c_{\mathscr X}'''\bigl(1+\log(1+\mathrm y^{-2})\bigr)\log^2(4+|\mathfrak z|)\qquad(\mathfrak z=\mathrm x+iy,\ \mathrm y>0).
\end{equation}
For $|\mathrm y|\ge1$, $\log|\mathcal F_{\mathscr X}(\mathfrak z)|\le\pi|\mathrm y|+O(\log^2(4+|\mathfrak z|))$ by \CGeqref{eq:rb-F-formula}, \CGeqref{eq:rb-Q-growth} and $|\mathcal F_{\mathscr X}(\bar{\mathfrak z})|=|\mathcal F_{\mathscr X}(\mathfrak z)|$. For $|\mathrm y_1|<1$, apply the sub-mean-value property of $\log|\mathcal F_{\mathscr X}|$ on the circle $|\mathfrak z-\mathfrak z_1|=\frac12$. The factor $\log(1+\mathrm y^{-2})$ is integrable along the circle uniformly in $\mathrm y_1$, since it has at most two logarithmic singularities. This gives $\log|\mathcal F_{\mathscr X}(\mathfrak z_1)|\le\frac{3\pi}2+O(\log^2(4+|\mathfrak z_1|))$ for $|\mathrm y_1|<1$. Hence the type is at most $\pi$ and the indicator satisfies $\mathrm h(\mathrm u)\le\pi|\sin\mathrm u|$. Since \CGeqref{eq:rb-F-formula} is an equality and \CGeqref{eq:rb-Q-growth} is two-sided, $\log|\mathcal F_{\mathscr X}(iy)|=\log|\cos(i\pi \mathrm y)|+\pi\mathbb Y_{\mathrm y}\varPsi_{\mathscr X}(0)+\mathrm c_{\mathscr X}''=\pi|\mathrm y|+O(\log^2(4+|\mathrm y|))$ for $|\mathrm y|\ge1$, so $\mathrm h(\pm\frac\pi2)=\pi$. The indicator diagram lies in $[-i\pi,i\pi]$, whose support function is $\pi|\sin\mathrm u|$, and it contains $\pm i\pi$, so it equals that segment.
\end{proof}

\subsubsection{The criterion}\label{subsec:rb-criterion}

A weight $\mathrm w>0$ on $\mathbb R$ is in $(A_2)$ if
\begin{equation}\label{eq:rb-aux-4}
 [\mathrm w]_{A_2}:=\sup_I\Bigl(\frac1{|I|}\int_Iw\Bigr)\Bigl(\frac1{|I|}\int_Iw^{-1}\Bigr)<\infty ,
\end{equation}
the supremum being over bounded intervals.

\begin{theorem}[Block $\tfrac14$-criterion]\label{thm:rb-criterion}
Let $\mathscr X$ satisfy \CGeqref{eq:rb-H} and suppose $\varPsi_{\mathscr X}=\mathrm A+\mathrm B$ with measurable $\mathrm A,\mathrm B$ such that:
\begin{enumerate}[label=\textup{(\roman*)}]
\item $\|\mathrm A\|_\infty\le \mathrm d_\star<\frac14$;
\item $\sup_x|\mathbb Y_1\mathrm B(\mathrm x)-\mathrm c_{\mathrm B}|\le \mathbb U$ for some constant $\mathrm c_{\mathrm B}$.
\end{enumerate}
Then $\mathrm w(\mathrm x):=|\mathcal F_{\mathscr X}(\mathrm x+i)|^2$ is in $(A_2)$ and
\begin{equation}\label{eq:rb-A2const}
 [\mathrm w]_{A_2}\le\frac{\pi^2\coth^2(\pi)\,e^{4\pi \mathbb U}}{\cos^2(2\pi\mathrm d_\star)} .
\end{equation}
\end{theorem}

\begin{proof}
$\mathrm B=\varPsi_{\mathscr X}-\mathrm A\in L^1(d\mathrm x/(1+\mathrm x^2))$, so $\mathbb Y_1\varPsi_{\mathscr X}=\mathbb Y_1\mathrm A+\mathbb Y_1\mathrm B$. By \CGeqref{eq:rb-F-formula}, $\mathrm w=e^{2\mathrm c''_{\mathscr X}}|\mathcal F_0(\mathrm x+i)|^2e^{2\pi\mathbb Y_1\mathrm A}e^{2\pi\mathbb Y_1\mathrm B}$.

Put $\mathbb A_{\mathrm A}(\mathfrak z):=\int \mathrm A(\mathrm s)\bigl[\frac1{\mathfrak z-\mathrm s}+\frac{\mathrm s}{1+\mathrm s^2}\bigr]d\mathrm s$. It is holomorphic on $\Im \mathfrak z>0$, and $\mathbb A_{\mathrm A}(\mathfrak o+i)$ is holomorphic on $\Im\mathfrak o>-1$. Let $\mathcal W(\mathfrak o):=\exp(2\mathbb A_{\mathrm A}(\mathfrak o+i))$. For $\mathfrak o=\mathrm x_1+i\mathrm y_1$ with $\mathrm y_1\ge0$,
\begin{equation}\label{eq:rb-aux-5}
 |\mathcal W(\mathfrak o)|=e^{2\pi\mathbb Y_{1+\mathrm y_1}\mathrm A(\mathrm x_1)},\qquad\arg\mathcal W(\mathfrak o)=-2\pi \mathbb P_{1+\mathrm y_1}\mathrm A(\mathrm x_1)\in[-2\pi\mathrm d_\star,2\pi\mathrm d_\star].
\end{equation}
So $\mathcal W$ maps the closed upper half-plane into the sector $|\arg|\le \mathrm b_\star:=2\pi\mathrm d_\star<\frac\pi2$, and $|\mathcal W(\mathrm x)|=e^{2\pi\mathbb Y_1\mathrm A(\mathrm x)}$ on $\mathbb R$.

$\Re\mathcal W$ is a positive harmonic function on $\Im\mathfrak o>0$, continuous up to $\mathbb R$. By the Herglotz representation it equals $\mathrm c\,\mathrm y_1+\int \mathbb P_{\mathfrak o}(t)\,d\mathrm n(t)$ with $\mathrm c\ge0$, $\mathrm n\ge0$, and $\mathrm n=\Re\mathcal W(t)\,dt$ by continuity. Here $\mathbb P_{\mathfrak o}(t):=\frac1\pi\frac{\mathrm y_1}{(t-\mathrm x_1)^2+\mathrm y_1^2}$. Hence $\Re\mathcal W(\mathfrak o)\ge\int \mathbb P_{\mathfrak o}(t)\Re\mathcal W(t)\,dt$. For an interval $I$ with centre $\mathrm x_I$, put $\mathfrak o_I:=\mathrm x_I+i|I|/2$. Then $\mathbb P_{\mathfrak o_I}(t)\ge(\pi|I|)^{-1}$ for $t\in I$, and
\begin{equation}\label{eq:rb-aux-6}
 \frac1{|I|}\int_I|\mathcal W|\le\frac1{\cos\mathrm b_\star}\,\frac1{|I|}\int_I\Re\mathcal W\le\frac\pi{\cos\mathrm b_\star}\Re\mathcal W(\mathfrak o_I)\le\frac\pi{\cos\mathrm b_\star}|\mathcal W(\mathfrak o_I)| .
\end{equation}
$1/\mathcal W$ maps into the same sector, so the same holds for it. Multiplying, $[|\mathcal W|]_{A_2}\le\pi^2/\cos^2\mathrm b_\star$.

Finally $\mathrm w=\mathrm g\,|\mathcal W|$ with $\mathrm g:=e^{2\mathrm c''_{\mathscr X}}|\mathcal F_0(\mathrm x+i)|^2e^{2\pi\mathbb Y_1\mathrm B}$. Since $|\cos\pi(\mathrm x+i)|^2=\cos^2(\pi \mathrm x)\cosh^2\pi+\sin^2(\pi \mathrm x)\sinh^2\pi\in[\sinh^2\pi,\cosh^2\pi]$ and $e^{2\pi\mathbb Y_1\mathrm B}\in[e^{2\pi(\mathrm c_{\mathrm B}-\mathbb U)},e^{2\pi(\mathrm c_{\mathrm B}+\mathbb U)}]$, we have $\sup \mathrm g/\inf \mathrm g\le\coth^2(\pi)e^{4\pi \mathbb U}$. For any weight $\mathrm w_1$, $[\mathrm g\mathrm w_1]_{A_2}\le(\sup \mathrm g/\inf \mathrm g)[\mathrm w_1]_{A_2}$.
\end{proof}

\begin{corollary}[Riesz bases]\label{cor:rb-riesz}
Under the hypotheses of Theorem~\CGref{thm:rb-criterion}:
\begin{enumerate}[label=\textup{(\alph*)}]
\item if $\mathscr X$ is uniformly discrete (in particular without repeated points), $\{e^{i\mathrm x_k\mathrm v}\}_{k\in\mathbb Z}$ is a Riesz basis of $L^2(-\pi,\pi)$;
\item if $\mathscr X$ is a finite union of uniformly discrete sets, the generalised divided differences of $\{e^{i\mathrm x_k\mathrm v}\}$ in the sense of \cite{AvdoninIvanov2002} form a Riesz basis of $L^2(-\pi,\pi)$. Consequently, for every partition of $\mathbb Z$ into consecutive blocks $\mathfrak B_m$ that are unions of the groups of \cite{AvdoninIvanov2002}, the subspaces spanned by $\{\mathrm v^je^{i\mathrm x\mathrm v}\}$, with $\mathrm x=\mathrm x_k$ for $k\in\mathfrak B_m$ and $j$ below the multiplicity of $\mathrm x$, form a Riesz basis of subspaces.
\end{enumerate}
\end{corollary}

\begin{proof}
By Lemma~\CGref{lem:rb-generating}, $\mathcal F_{\mathscr X}$ is entire of exponential type with indicator diagram of width $2\pi$ and zero set $\mathscr X$. By Theorem~\CGref{thm:rb-criterion}, $|\mathcal F_{\mathscr X}(\mathrm x+i)|^2\in(A_2)$, which is equivalent to the Helson--Szeg\H{o} condition \cite[Ch.~VI]{Garnett2007}. Then (a) is Pavlov's criterion \cite[Thm.~1]{AvdoninIvanov2002}, and the first statement of (b) is \cite[Thm.~3]{AvdoninIvanov2002}. Both are stated for $L^2(0,2\pi)$, and translation by $\pi$ multiplies each exponential by a unimodular constant. Grouping the vectors of a Riesz basis into disjoint groups gives a Riesz basis of subspaces.
\end{proof}

\subsubsection{Kadec, Avdonin, and variable block length}\label{subsec:rb-variable}

Write $\mathrm x_k=k-\frac12+\mathrm d_k$, $I_k:=[k-1,k)$, $\mathsf i_k:=\mathbf 1_{[\mathrm x_k,\infty)}-\mathbf 1_{[k-\frac12,\infty)}$, so that $\varPsi_{\mathscr X}=\sum_k\mathsf i_k$ and $\int\mathsf i_k=-\mathrm d_k$.

\begin{proposition}[Bounded conjugate for block-supported oscillation]\label{prop:rb-variable-blocks}
Let $\{\mathfrak J_m\}_{m\in\mathbb Z}$ be a partition of $\mathbb R$ into consecutive intervals of lengths $\mathfrak n_m\ge1$. Let $\mathrm B=\sum_m\mathrm B_m$ with $\operatorname{supp}\mathrm B_m\subset \mathfrak J_m$, $\int \mathrm B_m=0$ and $|\mathrm B_m|\le \mathrm D'_m$, $\sup_m\mathrm D'_m<\infty$. Put
\begin{equation}\label{eq:rb-L0}
 \mathbb U_{\rm loc}:=\sup_m\ \sup_{\mathrm x\in \mathfrak J_{m-1}\cup \mathfrak J_m\cup \mathfrak J_{m+1}}\Bigl|\int_{\mathfrak J_m}\frac{(\mathrm x-\mathrm s)\,\mathrm B_m(\mathrm s)}{\pi((\mathrm x-\mathrm s)^2+1)}\,d\mathrm s\Bigr|,\qquad
 \mathbb O:=\sup_{m_0}\sum_{|m-m_0|\ge2}\frac{\mathfrak n_m^2}{\operatorname{dist}(\mathfrak J_{m_0},\mathfrak J_m)^2}.
\end{equation}
If $\mathbb U_{\rm loc}<\infty$ and $\mathbb O<\infty$, then $\sup_x|\mathbb Y_1\mathrm B(\mathrm x)-\mathrm c_{\mathrm B}|\le3\mathbb U_{\rm loc}+\frac{\mathbb O}{2\pi}\sup_m\mathrm D'_m$ for a constant $\mathrm c_{\mathrm B}$.
\end{proposition}

\begin{proof}
The term $\frac1\pi\int\frac{\mathrm s}{1+\mathrm s^2}\mathrm B_m(\mathrm s)\,d\mathrm s$ equals $\frac1\pi\int(\frac{\mathrm s}{1+\mathrm s^2}-\frac{\mathrm m_m}{1+\mathrm m_m^2})\mathrm B_m(\mathrm s)\,d\mathrm s$, where $\mathrm m_m$ is the centre of $\mathfrak J_m$. Since $|\frac d{d\mathrm s}\frac{\mathrm s}{1+\mathrm s^2}|\le\frac1{1+\mathrm s^2}$, its modulus is at most $\mathrm D'_m\mathfrak n_m^2/(2\pi(1+\operatorname{dist}(0,\mathfrak J_m)^2))$. If $m_0$ is the index with $0\in \mathfrak J_{m_0}$, then $\operatorname{dist}(0,\mathfrak J_m)\ge\operatorname{dist}(\mathfrak J_{m_0},\mathfrak J_m)$, so these terms are summable by $\mathbb O<\infty$. Their sum is the constant $\mathrm c_{\mathrm B}$. Let $\mathrm x\in \mathfrak J_{m_0}$. The three blocks with $|m-m_0|\le1$ contribute at most $3\mathbb U_{\rm loc}$. For $|m-m_0|\ge2$, $\int \mathrm B_m=0$ gives $\bigl|\int_{\mathfrak J_m}\mathrm k(\mathrm x-\mathrm s)\mathrm B_m(\mathrm s)\,d\mathrm s\bigr|\le \mathrm D'_m\mathfrak n_m\cdot\frac{\mathfrak n_m}2\sup_{\mathrm s\in \mathfrak J_m}|\mathrm k'(\mathrm x-\mathrm s)|$ with $\mathrm k(\mathrm v):=\frac{\mathrm v}{\pi(\mathrm v^2+1)}$. Moreover $|\mathrm k'(\mathrm v)|\le\frac1{\pi\mathrm v^2}$ and $|\mathrm x-\mathrm s|\ge\operatorname{dist}(\mathfrak J_{m_0},\mathfrak J_m)$. Sum over $m$.
\end{proof}

\begin{corollary}[Kadec and Avdonin, with the constant $\tfrac14$]\label{cor:rb-kadec-avdonin}
\textup{(a)} \cite{Kadec1964} If $|\mathrm d_k|\le \mathrm d_\star<\frac14$ for all $k$, then $\{e^{i\mathrm x_k\mathrm v}\}$ is a Riesz basis of $L^2(-\pi,\pi)$.\\
\textup{(b)} \cite{Avdonin1974} If $\mathscr X$ is uniformly discrete, $\sup_k|\mathrm d_k|\le \mathrm D$, and for a fixed $\mathfrak n$ the means of $\mathrm d_k$ over consecutive blocks of $\mathfrak n$ indices are at most $\mathrm d_\star<\frac14$ in modulus, then $\{e^{i\mathrm x_k\mathrm v}\}$ is a Riesz basis of $L^2(-\pi,\pi)$.
\end{corollary}

\begin{proof}
Both statements follow from Theorem~\CGref{thm:rb-variable-N} below: in (a) take $\mathfrak n_m\equiv1$, and in (b) take $\mathfrak n_m\equiv \mathfrak n$. In both cases the blocks at index distance $j\ge2$ from $\mathfrak J_{m_0}$ are at distance $(j-1)\mathfrak n$, so $\mathbb O\le2\sum_{j\ge2}\mathfrak n^2/((j-1)\mathfrak n)^2=\pi^2/3$ and $\mathbb V\le2\mathrm D\log(2+3\mathfrak n)$; for $\mathfrak n=1$, $\mathrm D_m=0$ and $\mathbb V=0$. For (a), $\mathrm d_\star<\frac14$ gives gaps $\ge1-2\mathrm d_\star>\frac12$, so $\mathscr X$ is uniformly discrete.
\end{proof}

\begin{theorem}[Avdonin's theorem with variable block length]\label{thm:rb-variable-N}
Let $\mathscr X$ be uniformly discrete with $\sup_k|\mathrm d_k|\le \mathrm D$. Let $\{\mathfrak B_m\}$ be a partition of $\mathbb Z$ into consecutive blocks of $\mathfrak n_m$ indices, and $\mathfrak J_m:=\bigcup_{k\in\mathfrak B_m}I_k$. Let $\bar{\mathrm d}_m$ be the mean of $\mathrm d_k$ over $\mathfrak B_m$, $\mathrm D_m:=\max_{k\in\mathfrak B_m}|\mathrm d_k-\bar{\mathrm d}_m|$, and $\mathfrak n_m^*:=\mathfrak n_{m-1}+\mathfrak n_m+\mathfrak n_{m+1}$. Assume
\begin{enumerate}[label=\textup{(\alph*)}]
\item $\sup_m|\bar{\mathrm d}_m|\le \mathrm d_\star<\frac14$;
\item $\mathbb O<\infty$, with $\mathbb O$ as in \CGeqref{eq:rb-L0};
\item $\mathbb V:=\sup_m\mathrm D_m\log(2+\mathfrak n_m^*)<\infty$.
\end{enumerate}
Then $\{e^{i\mathrm x_k\mathrm v}\}$ is a Riesz basis of $L^2(-\pi,\pi)$, and \CGeqref{eq:rb-A2const} holds with
\begin{equation}\label{eq:rb-U-variable}
 \mathbb U=16(1+\mathrm D)^3+\frac{6\mathbb V}\pi+\frac{\mathbb V\mathbb O}{2\pi\log2} .
\end{equation}
\end{theorem}

\begin{proof}
Write $\varPsi_{\mathscr X}=\mathrm A+\mathrm B_1+\mathrm B_2$ with $\mathrm A:=-\sum_m\bar{\mathrm d}_m\mathbf 1_{\mathfrak J_m}$, $\mathrm B_1:=\sum_k\mathrm b_k$ where $\mathrm b_k:=\mathsf i_k+\mathrm d_k\mathbf 1_{I_k}$, and $\mathrm B_2:=-\sum_m\sum_{k\in\mathfrak B_m}(\mathrm d_k-\bar{\mathrm d}_m)\mathbf 1_{I_k}$. Then $\|\mathrm A\|_\infty\le \mathrm d_\star$ by (a).

\emph{$\mathrm B_1$.} $\mathrm b_k$ has mean zero, $|\mathrm b_k|\le1+\mathrm D$, and support in an interval of length at most $2\mathrm D+1$ containing $\mathrm m_k:=k-\frac12$. With $\mathrm k$ as above, $\sup|\mathrm k|=\frac1{2\pi}$ and $|\mathrm k'(\mathrm v)|\le\frac1{\pi(1+\mathrm v^2)}$, so $|\int \mathrm k(\mathrm x-\mathrm s)\mathrm b_k(\mathrm s)\,d\mathrm s|=|\int(\mathrm k(\mathrm x-\mathrm s)-\mathrm k(\mathrm x-\mathrm m_k))\mathrm b_k(\mathrm s)\,d\mathrm s|\le(1+\mathrm D)(2\mathrm D+1)\min\bigl(\frac1\pi,(2\mathrm D+1)\sup_{\mathrm s\in\operatorname{supp}\mathrm b_k}|\mathrm k'(\mathrm x-\mathrm s)|\bigr)$. At most $8\mathrm D+5\le8(1+\mathrm D)$ indices have $|\mathrm x-\mathrm m_k|<2(2\mathrm D+1)$, and they contribute at most $16(1+\mathrm D)^3/\pi$. For the others $|\mathrm x-\mathrm s|\ge|\mathrm x-\mathrm m_k|/2$ on the support, so the term is at most $4(1+\mathrm D)^3\cdot\frac4\pi(4+(\mathrm x-\mathrm m_k)^2)^{-1}$. Since $\sum_k(4+(\mathrm x-\mathrm m_k)^2)^{-1}\le\frac12+\frac\pi2$, they contribute at most $\frac{16}\pi(\frac12+\frac\pi2)(1+\mathrm D)^3$. The total is below $16(1+\mathrm D)^3$. The modifying term of $\mathbb Y_1$ contributes a convergent sum of constants, as in Proposition~\CGref{prop:rb-variable-blocks}.

\emph{$\mathrm B_2$.} On $\mathfrak J_m$ it has mean zero and modulus at most $\mathrm D_m$. For $\mathrm x$ in $\mathfrak J_{m-1}\cup \mathfrak J_m\cup \mathfrak J_{m+1}$, $|\int_{\mathfrak J_m}\mathrm k(\mathrm x-\mathrm s)\mathrm B_2(\mathrm s)\,d\mathrm s|\le\frac{\mathrm D_m}\pi\int_{|u|\le \mathfrak n_m^*}\frac{|u|\,du}{1+u^2}\le\frac2\pi \mathrm D_m\log(2+\mathfrak n_m^*)\le\frac{2\mathbb V}\pi$. Moreover $\sup_m\mathrm D_m\le \mathbb V/\log2$. Proposition~\CGref{prop:rb-variable-blocks}, applied with the blocks $\mathfrak J_m$, $\mathrm D'_m=\mathrm D_m$ and $\mathbb U_{\rm loc}\le\frac{2\mathbb V}\pi$, gives $|\mathbb Y_1\mathrm B_2-\mathrm c|\le\frac{6\mathbb V}\pi+\frac{\mathbb V\mathbb O}{2\pi\log2}$.

Apply Theorem~\CGref{thm:rb-criterion} and Corollary~\CGref{cor:rb-riesz}(a).
\end{proof}

\begin{remark}[The price of variable length]\label{rem:rb-price}
For fixed $\mathfrak n$, conditions (b) and (c) hold automatically, and Theorem~\CGref{thm:rb-variable-N} is Avdonin's theorem. If $\mathfrak n_m\to\infty$, the averaging condition (a) is still only $|\bar{\mathrm d}_m|\le \mathrm d_\star<\frac14$. However, (c) then requires the oscillation inside the block to decay like $1/\log \mathfrak n_m$. In Theorem~\CGref{thm:rb-criterion} the oscillation enters only through the bounded-conjugate condition (ii), which allows cancellation inside the blocks and is weaker than (c). Without some condition of this kind the conclusion fails for growing blocks: a block of length $\mathfrak n$ with oscillation of size $\mathrm D_m$ can make $|\mathbb Y_1\mathrm B|$ of size $\mathrm D_m\log \mathfrak n$. Lyubarskii and Seip derive Kadec's theorem from their Muckenhoupt characterisation (\cite[Thm.~2]{LyubarskiiSeip1997} for $p=2$), and indicate that the same technique gives results of Avdonin type. I have not searched the literature for Theorem~\CGref{thm:rb-variable-N} or for the bounded-conjugate formulation of Theorem~\CGref{thm:rb-criterion}, so their novelty is not claimed. The proof of Theorem~\CGref{thm:rb-criterion} is the sufficiency half of the Helson--Szeg\H{o} argument, applied on the line $\Im \mathfrak z=1$.
\end{remark}

%======================================================================
\subsection{Application to lists near the RvM clock}\label{subsec:rb-application}

\subsubsection{Unfolding}\label{subsec:rb-unfolding}

\begin{definition}[Unfolded list]\label{def:rb-unfolding}
Let $\mathfrak t_0:=9.6769\ldots$ be the zero of $\NR$ in $(2\pi,\infty)$. For a $C$-admissible list $\boldsymbol\kappa$ with $\kappa_1>\mathfrak t_0$ (true for $\boldsymbol\tau$ and $\boldsymbol\gamma$), put $\mathrm x_k:=\NR(\kappa_k)>0$ for $k\ge1$ and $\mathrm x_{1-k}:=-\mathrm x_k$. The resulting symmetric sequence $\mathscr X_{\boldsymbol\kappa}$ has discrepancy $\varPsi_{\boldsymbol\kappa}:=\varPsi_{\mathscr X_{\boldsymbol\kappa}}$, and $\mathcal F_{\boldsymbol\kappa}:=F_{\mathscr X_{\boldsymbol\kappa}}$ with $\mathfrak z_0=0$. For the clock, $\mathscr X_{\boldsymbol\tau}=\mathbb Z+\frac12$ and $\mathcal F_{\boldsymbol\tau}=\mathcal F_0$.
\end{definition}

\begin{lemma}[Unfolded discrepancy]\label{lem:rb-psi-bound}
For a $C$-admissible list with $\kappa_1>\mathfrak t_0$, $|\mathrm x_k-(k-\frac12)|\le\frac C{2\pi}\log\frac{\tau_k+C}{2\pi}$ for $k\ge1$, so \CGeqref{eq:rb-H} holds. Moreover, for $\mathrm x=\NR(t)$ with $\mathrm x>0$,
\begin{equation}\label{eq:rb-psi-t}
 \varPsi_{\boldsymbol\kappa}(\mathrm x)=N_{\boldsymbol\kappa}(t)-n_{\RVM}(t),\qquad
 |N_{\boldsymbol\kappa}(t)-n_{\RVM}(t)|\le\frac C\pi\log\frac{t+C}{2\pi}+2,
\end{equation}
and $\varPsi_{\boldsymbol\kappa}(-\mathrm x)=-\varPsi_{\boldsymbol\kappa}(\mathrm x)$ for almost every $\mathrm x$.
\end{lemma}

\begin{proof}
By the mean value theorem and the monotonicity of $\NR'$, $|\NR(\kappa_k)-\NR(\tau_k)|\le C\NR'(\tau_k+C)$. For $\mathrm x>0$ the terms of \CGeqref{eq:rb-psi} with $k\le0$ vanish, because both $\mathrm x_k$ and $k-\frac12$ are negative. For $k\ge1$, $\mathrm x_k\le \mathrm x$ if and only if $\kappa_k\le t$, by monotonicity of $\NR$. If $N_{\boldsymbol\kappa}(t)=j$, then $\tau_j\le t+C$ and $\tau_{j+1}>t-C$, so $n_{\RVM}(t-C)\le j\le n_{\RVM}(t+C)$. By \CGeqref{eq:rb-Ntau-floor}, $n_{\RVM}(t+C)-n_{\RVM}(\max(t-C,2\pi))\le\NR(t+C)-\NR(\max(t-C,2\pi))+1$. This is at most $\frac{2C}{2\pi}\log\frac{t+C}{2\pi}+2$, the extra $1$ covering $t-C<2\pi$, where $\NR\ge-\frac18$ and $n_{\RVM}=0$. Oddness follows from the symmetry of $\mathscr X_{\boldsymbol\kappa}$ and of $\mathbb Z+\frac12$.
\end{proof}

\begin{proposition}[The condition of $\mathscr O$ gives the averaging hypothesis with blocks $\asymp\log^2T$]\label{prop:rb-O-blocks}
Let $\boldsymbol\kappa$ be $C$-admissible with $\kappa_1>\mathfrak t_0$, and assume the third condition of $\mathscr O$ (\CGref{eq:list-space}) in the form
\begin{equation}\label{eq:rb-O3}
 \Bigl|\int_{2\pi}^{R}\bigl(N_{\boldsymbol\kappa}(t)-\NR(t)\bigr)dt\Bigr|\le C_2\log R\qquad(R\ge3),
\end{equation}
with a constant $C_2>0$; this is no restriction, since \CGeqref{eq:rb-O3} with one constant holds with every larger one, and it is needed because the block length $\mathrm H(T):=16C_2\log T$ must be positive. Put $\mathrm H(T):=16C_2\log T$. There is $T_1=T_1(C,C_2)$ such that the following holds for the points $T_{m+1}:=T_m+\mathrm H(T_m)$ ($m\ge1$) and the blocks $\mathfrak J_m:=[\NR(T_m),\NR(T_{m+1}))$ with their mirror images $-\mathfrak J_m$. The step function $\mathrm A_{\boldsymbol\kappa}$, equal to the mean of $\varPsi_{\boldsymbol\kappa}$ on each $\pm \mathfrak J_m$ and to $0$ on $(-\NR(T_1),\NR(T_1))$, satisfies
\begin{equation}\label{eq:rb-block-mean}
 \|\mathrm A_{\boldsymbol\kappa}\|_\infty\le\tfrac3{16},\qquad |\mathfrak J_m|=\frac{8C_2}\pi\log^2T_m\,(1+o(1))\quad(m\to\infty).
\end{equation}
The remainder $\mathrm B_{\boldsymbol\kappa}:=\varPsi_{\boldsymbol\kappa}-\mathrm A_{\boldsymbol\kappa}$ has mean zero on every block. For $\boldsymbol\kappa=\boldsymbol\gamma$, \CGeqref{eq:rb-O3} holds unconditionally for some $C_2$.
\end{proposition}

\begin{proof}
Write $\mathfrak J=[\NR(T),\NR(T+\mathrm H))$. Then
\begin{equation}\label{eq:rb-aux-7}
\begin{split}
 \int_{\mathfrak J}\varPsi_{\boldsymbol\kappa}\,d\mathrm x&=\int_T^{T+\mathrm H}(N_{\boldsymbol\kappa}-n_{\RVM})(t)\,\NR'(t)\,dt\\
 &=\NR'(T)\!\int_T^{T+\mathrm H}\!(N_{\boldsymbol\kappa}-n_{\RVM})+\int_T^{T+\mathrm H}\!(N_{\boldsymbol\kappa}-n_{\RVM})(\NR'(t)-\NR'(T))\,dt .
\end{split}
\end{equation}
The last integral is at most $\mathrm H\cdot\frac{\mathrm H}{2\pi T}\bigl(\frac C\pi\log\frac{T+\mathrm H+C}{2\pi}+1\bigr)$, by \CGeqref{eq:rb-psi-t} and $\NR''(t)=\frac1{2\pi t}$.

Next, $\int_T^{T+\mathrm H}(N_{\boldsymbol\kappa}-n_{\RVM})=\int_T^{T+\mathrm H}(N_{\boldsymbol\kappa}-\NR)+\int_T^{T+\mathrm H}(\NR-n_{\RVM})$. The first piece is at most $2C_2\log(T+\mathrm H)$ by \CGeqref{eq:rb-O3}. For the second, substitute $\mathrm x=\NR(t)$. By \CGeqref{eq:rb-Ntau-floor}, $\NR-n_{\RVM}=\mathrm{saw}(\mathrm x):=\mathrm x-\lfloor \mathrm x+\frac12\rfloor$, which is $1$-periodic with mean zero and $|\mathrm{saw}|\le\frac12$, and $dt=\mathfrak g_1(\mathrm x)\,d\mathrm x$ with $\mathfrak g_1:=1/\NR'(t(\mathrm x))$ positive and decreasing. On a full period $[j-\frac12,j+\frac12)$, $|\int\mathrm{saw}\,\mathfrak g_1|=|\int\mathrm{saw}\,(\mathfrak g_1-\mathfrak g_1(j))|\le\frac14(\mathfrak g_1(j-\frac12)-\mathfrak g_1(j+\frac12))$, which telescopes. The two partial periods contribute at most $\mathfrak g_1(\NR(T))$. Hence $|\int_T^{T+\mathrm H}(\NR-n_{\RVM})|\le\frac5{4\NR'(T)}$.

Since $|\mathfrak J|\ge\NR'(T)\mathrm H$,
\begin{equation}\label{eq:rb-aux-8}
 \Bigl|\frac1{|\mathfrak J|}\int_{\mathfrak J}\varPsi_{\boldsymbol\kappa}\Bigr|\le\frac{2C_2\log(T+\mathrm H)}{\mathrm H}+\frac5{4\NR'(T)\mathrm H}+\frac{\mathrm H}{2\pi T\NR'(T)}\Bigl(\frac C\pi\log\frac{T+\mathrm H+C}{2\pi}+2\Bigr).
\end{equation}
With $\mathrm H=16C_2\log T$, the first term tends to $\frac18$ and the others tend to $0$, so the bound is at most $\frac3{16}$ for $T\ge T_1(C,C_2)$. The length is $|\mathfrak J_m|=\NR'(\mathrm s)\mathrm H(T_m)$ with $\mathrm s\in[T_m,T_{m+1}]$. The mirror blocks follow by oddness. For $\boldsymbol\gamma$: $N_{\boldsymbol\gamma}(t)=1+\vartheta(t)/\pi+S(t)$ \cite[\S9.3]{Titchmarsh1986}, $\int_0^RS=O(\log R)$ \cite[Thm.~9.9]{Titchmarsh1986}, and $1+\vartheta/\pi-\NR=O(1/t)$. Finally, $\boldsymbol\gamma$ is $C_\eta$-admissible, unconditionally, by \CGref{lem:bounded-displacement}. In particular \CGeqref{eq:rb-H} holds for $\mathscr X_{\boldsymbol\gamma}$, by Lemma~\CGref{lem:rb-psi-bound}.
\end{proof}

\begin{remark}[Correction of an earlier formulation]\label{rem:rb-not-equivalent}
Proposition~\CGref{prop:rb-O-blocks} is an implication. The averaging condition on blocks of length $\asymp\log^2T$ does not imply \CGeqref{eq:rb-O3}: summing $|\int_{\mathfrak J_m}\varPsi|\le\frac3{16}|\mathfrak J_m|$ over the blocks up to height $R$ gives only $O(R\log R)$, not $O(\log R)$. The two conditions do not coincide.
\end{remark}

\begin{corollary}[Reduction to one condition]\label{cor:rb-one-condition}
For every $\boldsymbol\kappa$ as in Proposition~\CGref{prop:rb-O-blocks}: if $\mathbb Y_1\mathrm B_{\boldsymbol\kappa}$ is bounded up to a constant, then $|\mathcal F_{\boldsymbol\kappa}(\mathrm x+i)|^2\in(A_2)$, with constant \CGeqref{eq:rb-A2const} for $\mathrm d_\star=\frac3{16}$. Corollary~\CGref{cor:rb-riesz}(a) then holds for $\mathscr X_{\boldsymbol\kappa}$ if it is uniformly discrete, and Corollary~\CGref{cor:rb-riesz}(b) if it is a finite union of uniformly discrete sets.
\end{corollary}

\begin{proof}
Lemma~\CGref{lem:rb-psi-bound}, Proposition~\CGref{prop:rb-O-blocks} and Theorem~\CGref{thm:rb-criterion}.
\end{proof}

\subsubsection{The weight is a quotient of synthetic zetas}\label{subsec:rb-weight}

\begin{proposition}[Weight identity]\label{prop:rb-weight}
Let $\boldsymbol\kappa$ be $C$-admissible and satisfy \CGeqref{eq:rb-O3}. For $0<\hbar\le1$ and $t\ge2\pi$,
\begin{equation}\label{eq:rb-weight-identity}
 \log\Bigl|\frac{\zeta_{\boldsymbol\kappa}}{\zeta_{\RVM}}\Bigl(\tfrac12+\hbar+it\Bigr)\Bigr|=\int_{2\pi}^{\infty}\frac{(t-\mathrm s)\,(N_{\boldsymbol\kappa}-n_{\RVM})(\mathrm s)}{(t-\mathrm s)^2+\hbar^2}\,d\mathrm s+\mathrm c_{\boldsymbol\kappa}+\mathrm m_{\boldsymbol\kappa,\hbar}(t),
\end{equation}
where the integral converges as an improper integral and $\mathrm c_{\boldsymbol\kappa}$ is a constant. The remainder satisfies $|\mathrm m_{\boldsymbol\kappa,\hbar}(t)|\le \mathrm c(C,C_2)\log(2+t)/(1+t)$.
\end{proposition}

\begin{proof}
For $s=\frac12+\hbar+it$, $|\frac14+\kappa^2+q(s)|=|\kappa^2+(\hbar+it)^2|=|\kappa-t+i\hbar|\,|\kappa+t-i\hbar|$. Hence, with $\mathrm l_\hbar(u):=\frac12\log(u^2+\hbar^2)$,
\begin{equation}\label{eq:rb-aux-9}
 \log\Bigl|\frac{P_{\boldsymbol\kappa}}{P_{\boldsymbol\tau}}(s)\Bigr|=\sum_k\bigl[\mathrm l_\hbar(t-\kappa_k)-\mathrm l_\hbar(t-\tau_k)\bigr]+\sum_k\bigl[\mathrm l_\hbar(t+\kappa_k)-\mathrm l_\hbar(t+\tau_k)\bigr]-\sum_k\log\frac{\frac14+\kappa_k^2}{\frac14+\tau_k^2}.
\end{equation}
Put $\mathfrak g_1:=N_{\boldsymbol\kappa}-n_{\RVM}$ and $\mathfrak G_1(\mathrm s):=\int_{2\pi}^{\mathrm s}\mathfrak g_1$. By \CGeqref{eq:rb-O3}, by the sawtooth estimate in the proof of Proposition~\CGref{prop:rb-O-blocks} on $[\tau_1,\infty)$, and by $|\NR-n_{\RVM}|\le\frac12$ on $[2\pi,\tau_1]$, $|\mathfrak G_1(\mathrm s)|\le\mathrm c\log(2+\mathrm s)$. For $\mathrm f$ smooth with $\mathrm f'(\mathrm s)=O(1/\mathrm s)$ and $\mathrm f''(\mathrm s)=O(1/\mathrm s^2)$, the sums $\sum_{k\le k^*}[\mathrm f(\kappa_k)-\mathrm f(\tau_k)]=-\int \mathrm f'\mathfrak g_{1,k^*}$ converge as $k^*\to\infty$ to $-\int^\infty \mathrm f'\mathfrak g_1$, an improper integral. Integration by parts against $\mathfrak G_1$ shows both the convergence and that the tail contributions vanish. With $\mathrm f=\mathrm l_\hbar(t-\cdot)$ this gives the integral in \CGeqref{eq:rb-weight-identity}. With $\mathrm f=\log(\frac14+(\cdot)^2)$ it gives the constant $\mathrm c_{\boldsymbol\kappa}$. With $\mathrm f=\mathrm l_\hbar(t+\cdot)$ it gives $\mathrm m_{\boldsymbol\kappa,\hbar}(t)=\int_{2\pi}^\infty \mathrm k_1'(\mathrm s)\mathfrak G_1(\mathrm s)\,d\mathrm s$, where $\mathrm k_1(\mathrm s):=\frac{t+\mathrm s}{(t+\mathrm s)^2+\hbar^2}$ and $|\mathrm k_1'(\mathrm s)|\le(t+\mathrm s)^{-2}$. This yields the bound.
\end{proof}

\begin{remark}[Local scale]\label{rem:rb-local-scale}
With $\hbar=1$ and $\hbar_T:=\hbar/\NR'(T)=2\pi/\log(T/2\pi)$ (the notation of the proof of Theorem~\CGref{thm:rb-obstruction}), the substitution $\mathrm x=\NR(t)$ maps the kernel $\frac{(t-\mathrm s)\,d\mathrm s}{(t-\mathrm s)^2+\hbar_T^2}$ to $\frac{(\mathrm x-\mathrm x')\,d\mathrm x'}{(\mathrm x-\mathrm x')^2+1}$ to first order near $t=T$. Comparing \CGeqref{eq:rb-weight-identity} with \CGeqref{eq:rb-F-formula} at $\mathrm y=1$, the unfolded weight $|\mathcal F_{\boldsymbol\kappa}(\mathrm x+i)|^2$ near $\mathrm x=\NR(T)$ is therefore, to first order, $|\zeta_{\boldsymbol\kappa}/\zeta_{\RVM}|^2$ on the line $\Re s=\frac12+\hbar_T$, up to the bounded factor $|\mathcal F_0(\mathrm x+i)|^2$. The identity \CGeqref{eq:rb-weight-identity} is exact. The identification with the unfolded weight is local, and its quantitative form for the Poisson part is Step~3 of the proof of Theorem~\CGref{thm:rb-obstruction}.
\end{remark}

\subsubsection{Obstruction for the zeta list}\label{subsec:rb-obstruction}

\begin{theorem}[No $(A_2)$ for the unfolded zeta ordinates]\label{thm:rb-obstruction}
Let $\mathscr X_{\boldsymbol\gamma}$ be the unfolded list of the ordinates $\gamma_k>0$ of the nontrivial zeros of $\zeta$, counted with multiplicity, and let $\mathcal F$ be any entire function of exponential type whose zero set is $\mathscr X_{\boldsymbol\gamma}$, with multiplicities. Then $|\mathcal F(\mathrm x+i\hbar)|^2\notin(A_2)$ for every real $\hbar\ne0$. The proof does not assume the Riemann hypothesis.
\end{theorem}

\begin{proof}
Write $\varPsi:=\varPsi_{\boldsymbol\gamma}$ and $\mathcal F_{\boldsymbol\gamma}$ for the function of Definition~\CGref{def:rb-unfolding}. It is even and real by Lemma~\CGref{lem:rb-generating}(iii), because $\mathrm x_1=\NR(\gamma_1)=0.45\ldots>0$, so $0\notin \mathscr X_{\boldsymbol\gamma}$.

\emph{Step 1: reduction to $\mathcal F_{\boldsymbol\gamma}$ and $\hbar>0$.} $\mathcal F$ and $\mathcal F_{\boldsymbol\gamma}$ are entire of order at most $1$ (exponential type) with the same zeros, with multiplicity, so by Hadamard's factorisation theorem both are $e^{\text{linear}}$ times the same canonical product, and $\mathcal F/\mathcal F_{\boldsymbol\gamma}=e^{\mathrm a_1+\mathrm b_1\mathfrak z}$. Hence $|\mathcal F(\mathrm x+i\hbar)|^2=\mathrm c\,e^{2\mathrm b_1'\mathrm x}\mathrm w_0(\mathrm x)$ with $\mathrm b_1'=\Re\mathrm b_1$, a constant $\mathrm c>0$, and $\mathrm w_0(\mathrm x):=|\mathcal F_{\boldsymbol\gamma}(\mathrm x+i\hbar)|^2$, which is even in $\mathrm x$. Moreover $|\mathcal F_{\boldsymbol\gamma}(\mathrm x-i\hbar)|=|\mathcal F_{\boldsymbol\gamma}(\mathrm x+i\hbar)|$. Suppose $e^{2\mathrm b_1'\mathrm x}\mathrm w_0\in(A_2)$. Reflection preserves $(A_2)$, so $e^{-2\mathrm b_1'\mathrm x}\mathrm w_0\in(A_2)$ as well. By Step~2 below both logarithms are in BMO, so $4\mathrm b_1'\mathrm x\in\mathrm{BMO}$, which forces $\mathrm b_1'=0$. So we may assume $\mathcal F=\mathcal F_{\boldsymbol\gamma}$ and $\hbar>0$, and that $\mathrm w_0\in(A_2)$.

\emph{Step 2: BMO.} Let $\mathrm f:=\log \mathrm w_0$ and $\mathrm f_I$ its mean on $I$. By Jensen's inequality, $\frac1{|I|}\int_Ie^{\pm(f-\mathrm f_I)}\ge1$. The product of the two averages is at most $[\mathrm w_0]_{A_2}$, so each is at most $[\mathrm w_0]_{A_2}$. Since $u^2\le2e^{|u|}\le2(e^u+e^{-u})$,
\begin{equation}\label{eq:rb-aux-10}
 \frac1{|I|}\int_I|f-\mathrm f_I|^2\le4[\mathrm w_0]_{A_2}\qquad\text{for every interval }I .
\end{equation}
By \CGeqref{eq:rb-F-formula} and the boundedness of $\log|\mathcal F_0(\mathrm x+i\hbar)|$, $\mathrm g_0:=\pi\mathbb Y_\hbar\varPsi\in\mathrm{BMO}$. Now take $\mathbb A_\hbar(\mathfrak z):=\mathbb A_{\mathscr X}(\mathfrak z+i\hbar)$, in the notation of the proof of Lemma~\CGref{lem:rb-generating}, which is holomorphic on $\Im \mathfrak z>-\hbar$. Its real part on $\mathbb R$ is $\mathrm g_0$, and its imaginary part on $\mathbb R$ is $-\pi \mathbb P_\hbar\varPsi$. By \CGeqref{eq:rb-Q-growth}, $\Re\mathbb A_\hbar(\mathfrak z)=O(\log^2(4+|\mathfrak z|))$ on $\Im \mathfrak z\ge0$, and so is the Poisson integral $\mathbb P[\mathrm g_0]$. Hence $\Re\mathbb A_\hbar-\mathbb P[\mathrm g_0]$ is harmonic on $\Im \mathfrak z>0$, continuous with zero boundary values, and $O(\log^2|\mathfrak z|)$. After odd reflection it is harmonic on $\mathbb C$ and $o(|\mathfrak z|)$, hence constant, hence zero. Therefore $\mathbb A_\hbar=\mathbb P[\mathrm g_0]+i\mathbb Y[\mathrm g_0]+ic$ on $\Im \mathfrak z>0$, where $\mathbb P[\mathrm g_0]+i\mathbb Y[\mathrm g_0]$ is the holomorphic function with real part $\mathbb P[\mathrm g_0]$ and $\mathbb Y[\mathrm g_0](\mathrm x+iy)=\mathbb Y_{\mathrm y}\mathrm g_0(\mathrm x)$. Letting $\mathrm y\to0$, $-\pi \mathbb P_\hbar\varPsi=\tilde{\mathrm g}_0+c$ almost everywhere on $\mathbb R$, where $\tilde{\mathrm g}_0$ is the conjugate function of $\mathrm g_0$, the a.e.\ boundary value of $\mathbb Y_{\mathrm y}u$ \cite[Ch.~III]{Garnett2007}. Conjugation is bounded on BMO \cite[Ch.~VI]{Garnett2007}, so $\mathbb P_\hbar\varPsi\in\mathrm{BMO}$. By the John--Nirenberg inequality \cite[Ch.~VI]{Garnett2007}, there is $\mathrm c_{\rm JN}<\infty$ with
\begin{equation}\label{eq:rb-BMO-L2}
 \frac1{|I|}\int_I\bigl|\mathbb P_\hbar\varPsi-(\mathbb P_\hbar\varPsi)_I\bigr|^2\le\mathrm c_{\rm JN}\qquad\text{for every interval }I .
\end{equation}

\emph{Step 3: back to the ordinates.} By \cite[\S9.3]{Titchmarsh1986}, $N_{\boldsymbol\gamma}(t)=1+\vartheta(t)/\pi+S(t)$, with $S(t)=O(\log t)$ \cite[Thm.~9.4]{Titchmarsh1986}. Lemma~\CGref{lem:rb-psi-bound} and \CGeqref{eq:rb-Ntau-floor} then give, for $\mathrm x=\NR(t)>0$, $\varPsi(\mathrm x)=S(t)+\mathrm r(\mathrm x)$ with $|\mathrm r|\le \mathrm c_e$. Here $\mathrm c_e$ bounds $|1+\vartheta/\pi-\NR|+\frac12$ on $t\ge2\pi$. Let $\mathfrak t_0:=\NR^{-1}(0)$ and $T\to\infty$, with $t\in[T,2T]$ and $\hbar_t:=\hbar/\NR'(t)$, $\hbar_T:=\hbar/\NR'(T)$.
\begin{enumerate}[label=(\alph*)]
\item The part of $\mathbb P_\hbar\varPsi(\NR(t))$ coming from $\mathrm x'<0$ is $O(\log T/T)$, since $|\varPsi(\mathrm x')|\le \mathrm c\log(2+|\mathrm x'|)$. The part coming from $\mathrm r$ is at most $\mathrm c_e$. Substituting $\mathrm x'=\NR(\mathrm s)$, the rest is $\int_{\mathfrak t_0}^\infty \mathbb P^{\NR}_t(\mathrm s)S(\mathrm s)\,d\mathrm s$ with $\mathbb P^{\NR}_t(\mathrm s):=\mathbb P_\hbar(\NR(t)-\NR(\mathrm s))\NR'(\mathrm s)$, where $\mathbb P_\hbar(\mathrm v):=\frac \hbar{\pi(\mathrm v^2+\hbar^2)}$.
\item Put $\mathbb P^{\circ}_t(\mathrm s):=\mathbb P_{\hbar_t}(t-\mathrm s)=\mathbb P_\hbar(\NR'(t)(t-\mathrm s))\NR'(t)$. For $|\mathrm s-t|\le T^{1/3}$, Taylor's formula with $\NR''=\frac1{2\pi \mathrm s}$ gives $|\NR(t)-\NR(\mathrm s)-\NR'(t)(t-\mathrm s)|\le\mathrm d_T:=T^{-1/3}/\pi$. Hence $|\mathbb P^{\NR}_t(\mathrm s)-\mathbb P^{\circ}_t(\mathrm s)|\le(5\mathrm d_T/\hbar)\mathbb P^{\circ}_t(\mathrm s)$ for $T$ large. With $|S|\le \mathrm c\log(3T)$ on that range, the core difference is $O(T^{-1/3}\log T)$. For $|\mathrm s-t|>T^{1/3}$ we use $\mathbb P_\hbar(\mathrm v)\le\frac \hbar{\pi \mathrm v^2}$, the convexity of $\NR$ and $|S(\mathrm s)|\le \mathrm c\log(2+\mathrm s)$. If $\mathrm s>t+T^{1/3}$, then $\NR(\mathrm s)-\NR(t)\ge\NR'(t)(\mathrm s-t)\ge \mathrm s-t$ for $T$ large, so $\int_{t+T^{1/3}}^\infty \mathbb P^{\NR}_t(\mathrm s)|S(\mathrm s)|\,d\mathrm s\le\frac{\mathrm c\hbar}\pi\int_{T^{1/3}}^\infty\frac{\NR'(t+\mathrm v)\log(2+t+\mathrm v)}{\mathrm v^2}\,d\mathrm v=O(T^{-1/3}\log^2T)$, splitting the $\mathrm v$-integral at $\mathrm v=T$. If $\mathfrak t_0\le \mathrm s<t-T^{1/3}$, then $\NR(t)-\NR(\mathrm s)\ge\NR'(\mathrm s)(t-\mathrm s)$ and $\NR'(\mathrm s)\ge\NR'(\mathfrak t_0)>0$, so $\mathbb P^{\NR}_t(\mathrm s)\le\frac{h}{\pi\NR'(\mathrm s)(t-\mathrm s)^2}\le\frac{h}{\pi\NR'(\mathfrak t_0)(t-\mathrm s)^2}$ and $\int_{\mathfrak t_0}^{t-T^{1/3}}\mathbb P^{\NR}_t|S|\le\frac{\mathrm c\hbar\log(3T)}{\pi\NR'(\mathfrak t_0)}\int_{T^{1/3}}^\infty \mathrm v^{-2}\,d\mathrm v=O(T^{-1/3}\log T)$. The same bounds hold for $\mathbb P^{\circ}_t$ with $\hbar_t$ in place of $\hbar$ and $\NR'(t)$ in place of $\NR'(\mathrm s)$, the region $\mathrm s\gg T$ being handled by the split at $\mathrm v=T$. So that part tends to $0$ as well.
\item $\|\mathbb P_a-\mathbb P_b\|_{L^1}\le\frac2\pi|\log(a/b)|$, because $\int|\partial_c\mathbb P_c|=\frac2{\pi c}$. Since $\hbar_t/\hbar_T\in[1-\frac{2\log2}{\log(T/2\pi)},1]$, $|\int(\mathbb P_{\hbar_t}-\mathbb P_{\hbar_T})(t-\mathrm s)S(\mathrm s)\,d\mathrm s|=O(\log T/\log T)=O(1)$, again with the tails controlled by $|S(\mathrm s)|\le \mathrm c\log(2+\mathrm s)$.
\item $\|\mathbb P_{\hbar_T}\ast S-S\|_{L^2[T,2T]}=O(\sqrt T)$, with $S:=0$ on $(-\infty,\mathfrak t_0)$; this is Lemma~\CGref{lem:rb-smoothing} below.
\end{enumerate}
Together, (a)--(d) give $\|\mathbb P_\hbar\varPsi\circ\NR-S\|_{L^2[T,2T]}=O(\sqrt T)$.

\emph{Step 4: contradiction.} Let $\mathfrak I_T:=[\NR(T),\NR(2T)]$ and $\bar{\mathrm p}_T:=(\mathbb P_\hbar\varPsi)_{\mathfrak I_T}$. By \CGeqref{eq:rb-BMO-L2}, $\int_T^{2T}|\mathbb P_\hbar\varPsi(\NR(t))-\bar{\mathrm p}_T|^2\NR'(t)\,dt\le\mathrm c_{\rm JN}|\mathfrak I_T|\le\mathrm c_{\rm JN}T\NR'(2T)$, so $\int_T^{2T}|\mathbb P_\hbar\varPsi(\NR(t))-\bar{\mathrm p}_T|^2\,dt\le2\mathrm c_{\rm JN}T$. With Step~3, $\int_T^{2T}|S-\bar{\mathrm p}_T|^2=O(T)$. But
\begin{equation}\label{eq:rb-aux-11}
 \int_T^{2T}|S-\mathrm c|^2\ge\int_T^{2T}S^2-\frac1T\Bigl(\int_T^{2T}S\Bigr)^2=\frac1{2\pi^2}T\log\log T\,(1+o(1))
\end{equation}
for every constant $\mathrm c$, by Selberg's theorem $\int_0^TS^2=\frac{T}{2\pi^2}\log\log T+O(T(\log\log T)^{1/2})$ \cite{Selberg1946}, with the normalisation $S(t)=\pi^{-1}\arg\zeta(\frac12+it)$ used here and in \cite[p.~225]{Fujii1974}, and $\int_T^{2T}S=O(\log T)$ \cite[Thm.~9.9]{Titchmarsh1986}. This is a contradiction.
\end{proof}

\begin{lemma}[Smoothing $S$ at scale $1/\log T$]\label{lem:rb-smoothing}
Let $\hbar>0$ be fixed, $\hbar_T:=2\pi \hbar/\log(T/2\pi)$, and $S:=0$ on $(-\infty,\mathfrak t_0)$. Then $\|\mathbb P_{\hbar_T}\ast S-S\|_{L^2[T,2T]}=O_\hbar(\sqrt T)$, unconditionally.
\end{lemma}

\begin{proof}
Fujii's Main Theorem \cite[Main Theorem]{Fujii1974}, proved unconditionally by modifying Selberg's method, states for $q=1$, $k=1$, $T^{1/2+a}\le H\le T$ and $0<u\le H-(H/\sqrt T)^{1/8}$ that $\int_T^{T+H}(S(t+u)-S(t))^2dt=c_1H\cdot2\log(2+u\log T)\,(1+O((\log(2+u\log T))^{-1/2}))$ with an absolute $c_1$. The relative error is $O(1)$ when $u\log T$ is bounded, so the statement carries no asymptotic information there, and we use it only as the upper bound it contains: for $0<u\le1$, $T\ge T_0$ and $H=T$,
\begin{equation}\label{eq:rb-fujii}
 \int_T^{2T}\bigl(S(t+u)-S(t)\bigr)^2dt\le \mathrm c_FT\log(2+u\log T),
\end{equation}
with an absolute constant $\mathrm c_F$.
The range $0<u\le H-(H/\sqrt T)^{1/8}=T-T^{1/16}$ of \cite{Fujii1974} contains $0<u\le1$, and the same bound holds for $-1\le u<0$: the substitution $t\mapsto t-u$ moves the integral to $[T-|u|,2T-|u|]\subset[T/2,2T]$, covered by \CGeqref{eq:rb-fujii} at the base points $T/2$ (with $H=T/2$) and $T$ (with $H=T$). By Minkowski's inequality, $\|\mathbb P_{\hbar_T}\ast S-S\|_{L^2[T,2T]}\le\int \mathbb P_{\hbar_T}(u)\|S(\cdot-u)-S\|_{L^2[T,2T]}\,du$. For $0<|u|\le1$ the norm is at most $(2\mathrm c_FT\log(2+|u|\log T))^{1/2}$. For $1<|u|\le T/2$, Selberg's theorem $\int_0^TS^2=\frac1{2\pi^2}T\log\log T+O(T(\log\log T)^{1/2})$ \cite{Selberg1946} bounds it by $\mathrm c(T\log\log T)^{1/2}$. For $|u|>T/2$, $|S(t)|\le \mathrm c\log(2+t)$ \cite[Thm.~9.4]{Titchmarsh1986} bounds it by $\mathrm c\sqrt T\log(3T+|u|)$. Substituting $u=\hbar_Tv$ in the first range, where $\hbar_T\log T\to2\pi \hbar$, and using $\int_{|u|>1}\mathbb P_{\hbar_T}=O(\hbar_T)$ and $\int_{|u|>T/2}\mathbb P_{\hbar_T}(u)\log(3T+|u|)\,du=O(\hbar_T\log T/T)$, all three contributions are $O(\sqrt T)$.
\end{proof}

\begin{remark}[A second proof]\label{rem:rb-fujii}
Lemma~\CGref{lem:rb-smoothing} also follows from Selberg's mean-square approximation of $S(t)$ by $S_{\mathrm x_0}(t):=-\frac1\pi\sum_{p\le \mathrm x}p^{-1/2}\sin(t\log p)$, $\mathrm x_0=T^\delta$ \cite{Selberg1946}, and the Montgomery--Vaughan mean value theorem \cite[Cor.~3]{MontgomeryVaughan1974}. Indeed $\mathbb P_{\hbar_T}\ast S_{\mathrm x_0}-S_{\mathrm x_0}=-\frac1\pi\sum_{p\le \mathrm x}p^{-1/2}(e^{-\hbar_T\log p}-1)\sin(t\log p)$ has mean square $\ll T\,\hbar_T^2\log^2\mathrm x\ll T$ on $[T,2T]$, and the remainder $S-S_{\mathrm x_0}$ is handled by Minkowski's inequality. Fujii's announcement \cite{Fujii1975} states the moment lemma of \cite{Fujii1974} in the form used in \cite[Lemma~1]{Fujii1975b}.
\end{remark}

\begin{corollary}[Consequences for the zeta list]\label{cor:rb-obstruction}
\textup{(a)} $\{e^{ix_k\xi}\}_{x_k\in \mathscr X_{\boldsymbol\gamma}}$ is not a Riesz basis of $L^2(-\pi,\pi)$.\\
\textup{(b)} If $\mathscr X_{\boldsymbol\gamma}$ is a finite union of uniformly discrete sets, its generalised divided differences are not a Riesz basis of $L^2(-\pi,\pi)$.\\
\textup{(c)} For every decomposition $\varPsi_{\boldsymbol\gamma}=A+B$ with $\|A\|_\infty<\frac14$, $\mathbb Y_1B$ is unbounded. In particular this holds for the blocks of Proposition~\CGref{prop:rb-O-blocks}, where the averaging hypothesis is satisfied.
\end{corollary}

\begin{proof}
(a) and (b): the criteria \cite[Thm.~1, Thm.~3]{AvdoninIvanov2002} are necessary and sufficient, and require a generating function satisfying Helson--Szeg\H{o}, equivalently $(A_2)$, on some line $\Im \mathfrak z=\hbar$. The line $\hbar=0$ is excluded, since $|\mathcal F(\mathrm x)|^{-2}$ is not locally integrable at a real zero. All lines with $\hbar\ne0$ are excluded by Theorem~\CGref{thm:rb-obstruction}. (c): otherwise Theorem~\CGref{thm:rb-criterion} would give $|\mathcal F_{\boldsymbol\gamma}(\mathrm x+i)|^2\in(A_2)$.
\end{proof}

\begin{heuristic}[Graduated version]\label{heur:rb-graduated}
The proof of Theorem~\CGref{thm:rb-obstruction} uses intervals $\mathfrak I_T$ of length $\asymp T\log T$. On a window of $L$ unfolded points at height $T$, the mean square oscillation of $\mathbb P_\hbar\varPsi$ is governed by the variance of increments of $S$ over $t$-intervals of length $L/\NR'(T)$. By Fujii's asymptotics this is expected to be of size $\frac1{\pi^2}\log L$. This suggests that local $(A_2)$ constants on windows of $L$ points grow with $L$ but stay bounded for $L$ bounded. That is the graduated stability that a Fredholm formulation would need. None of this is proved here. The computations of Section~\CGref{subsec:rb-numerics} measure the $\ell^\infty$ norm of the inverse of square finite sections directly and find growth like $T^{1.8}$ up to $T=487$.
\end{heuristic}

\begin{remark}[Bearing on route B]\label{rem:rb-routeB-stability}
Theorem~\CGref{thm:rb-obstruction} rules out a uniform Riesz-basis (Fredholm) structure for the unfolded zeta list on the whole line, and hence a linear stability argument with constants independent of the height. It says nothing about statement B or B$'$. The compactness step of route B (Proposition~\CGref{prop:rb-compactness}) uses no stability; the open part of route B is existence (Definition~\CGref{def:rb-B}) together with uniqueness in the box (Hypothesis~\CGref{hyp:rb-U}), and the open part of route B$'$ is existence alone (Corollary~\CGref{cor:rb-Bprime}).
\end{remark}

%======================================================================
\subsection{Numerical certificates for finite sections}\label{subsec:rb-numerics}

This section reports computations. They are not part of any proof; their role is to measure the quantities of Theorems~\CGref{thm:rb-certificate} and~\CGref{thm:rb-staged} on actual sections and to locate where the staged construction fails. The programs and their outputs are listed in Appendix~\CGref{app:code-routeB} and in \texttt{code/routeB/} of the repository. Floating point arithmetic was used throughout; no interval arithmetic was run.

\subsubsection{Set-up}\label{subsec:rb-num-setup}

The tests are Gaussians $\varphi_{c,\mathrm w}(u):=e^{-(u-c)^2/2\mathrm w^2}$ with centre $c\ge9\mathrm w$, so that their restriction to $u>0$ is, to working precision, supported in $(0,\infty)$; their cosine transform is $\mathsf c_{\varphi_{c,\mathrm w}}(t)=\mathrm w\sqrt{2\pi}\,e^{-(\mathrm wt)^2/2}\cos(ct)$. A Gaussian is not compactly supported, so these tests are not elements of $C_c^\infty(0,\infty)$ and the computed quantities are not literally those of Theorem~\CGref{thm:rb-certificate}, whose tests are in that space. The pairings against a compactly supported bump $\varphi_{c,\mathrm w}\chi$, with $\chi\in C_c^\infty(0,\infty)$ equal to $1$ on $[c/2,3c/2]$, differ from the computed ones by the pairings against the tail $\varphi_{c,\mathrm w}(1-\chi)$, which were not computed; the alternative, an extension of Theorem~\CGref{thm:rb-certificate} to Schwartz tests, is not carried out here. For this reason, and because of the floating-point arithmetic, the certificates reported below are \emph{candidate certificates}: floating-point evidence for the inequalities of the theorem evaluated on tests outside its hypothesis, not verified statements. Stage $m$ uses the width $\mathrm w_m:=0.08\cdot0.85^{m-1}$, the height $T_m:=4/\mathrm w_m$, and the new tests $\varphi_{c,\mathrm w_m}$ with centres $c\in[9\mathrm w_m,\,\log(\max(T_m,7)/2\pi)+2]$ in steps of $\mathrm w_m$; the test set $\mathscr T_m$ is the union of the levels $\le m$. The unknowns of stage $m$ are the ordinates with $\tau_k\le2.4/\mathrm w_m$ (``interior''); the factor $e^{-(\mathrm w_m\kappa)^2/2}$ makes the others invisible to the tests of level $m$. The masses are $\mathsf a^\zeta_n$ for all $n\le e^{c+9\mathrm w}$, and $\mathsf k$ is \CGref{eq:transport-defs}. The pairings $\langle\mathsf k,\varphi\rangle$ were computed by adaptive quadrature. The zeta ordinates $\gamma_k$, $k\le400$ ($\gamma_{400}=679.74$), computed with \texttt{mpmath}, enter only in the diagnostic columns; no construction uses them.

Two checks fix the normalisation: $\langle\mathfrak R(\boldsymbol\gamma,\mathsf a^\zeta),\varphi\rangle\approx10^{-16}$ for all tests tried, and $\langle\mathfrak R(\boldsymbol\tau,\mathsf a^\zeta),\varphi\rangle$ between $10^{-3}$ and $10^{-1}$.

\subsubsection{Initialisation: the clock is outside the basin, the Gonek lists are inside}\label{subsec:rb-num-init}

\emph{Receding horizon from the clock.} Starting at $\boldsymbol\kappa=\boldsymbol\tau$, each stage solved the overdetermined least-squares problem for the interior ordinates with all tests of levels $\le m$, with bounds $|\kappa_k-\tau_k|\le3$ and the non-interior ordinates frozen at their current values. This is the overdetermined geometry $\mathrm t_m\ge k^*_m$: the Jacobian has a left inverse (the pseudo-inverse) and Gauss--Newton is a least-squares iteration. Theorem~\CGref{thm:rb-certificate} requires a right inverse, $\mathbb J_m\mathbb Q_m=I_{\mathrm t_m}$, which forces $\mathrm t_m\le k^*_m$; the two geometries meet only in the square case $\mathrm t_m=k^*_m$, which is the one used for the certificates of Section~\CGref{subsec:rb-num-certified}. The receding-horizon runs are therefore diagnostics of the basin, not evaluations of the theorem. The iteration stalls at the first Gram-type failure of the clock: the fit returns $\kappa_{34}=109.75$ where $\gamma_{34}=111.03$ ($\tau_{34}=110.04$), and $\kappa_{41}=125.50$ where $\gamma_{41}=124.26$; from $m=9$ on the residual on the lagged levels $\le m-4$ grows from $2\cdot10^{-5}$ to $4\cdot10^{-2}$ and the front of ordinates within $10^{-2}$ of $\boldsymbol\gamma$ stays at $k\approx31$.

\emph{Receding horizon from a Gonek list.} The same iteration started at $\boldsymbol\kappa^{(\mathrm X)}$ with $\mathrm X=8$ ($983$ roots up to $t=1400$, that is, every solution of \CGeqref{eq:rb-gonek-phase} from $t\approx14$ on; this is not the list of Lemma~\CGref{lem:rb-gonek-tube}, whose prefix is the clock below $\mathfrak t_8\approx2\pi e^{2\mathrm n_8+2}$ with $\mathrm n_8\approx7$, so that the lemma's tube does not cover the list used here, whose monotonicity was checked directly; the roots were computed from the Riemann--Siegel expansion of $\vartheta$ and the sum \CGeqref{eq:rb-alphaX}) converged through $m=16$, $T_{16}=572$, $165$ interior ordinates and $4985$ tests: all interior ordinates within $10^{-2}$ of $\gamma_k$, $90$ of them within $10^{-4}$, the first twenty within $5\cdot10^{-6}$, residual on the lagged levels between $10^{-6}$ and $9\cdot10^{-6}$ at every stage, and $\max_k|\kappa_k-\tau_k|=0.99$ throughout. The residual over all tests of level $\le m$ stays at $2\cdot10^{-3}$, which is the contribution of the frozen edge.

\emph{Why.} The basin of the exact solution, measured by random perturbations of $\boldsymbol\gamma$ in the interior block, has radius $0.1$ to $0.25$ mean spacings. The clock lies at up to $0.66$ spacings from $\boldsymbol\gamma$ for $k\le400$, with $50$ indices beyond $0.25$; the lists $\boldsymbol\kappa^{(\mathrm X)}$ lie within $0.184$ ($\mathrm X=3$), $0.081$ ($\mathrm X=5$), $0.053$ ($\mathrm X=8$), $0.047$ ($\mathrm X=12$) and $0.035$ ($\mathrm X=20$) spacings for $k\le160$, and within $0.103$ spacings for $k\le400$ at $\mathrm X=8$ (Figure~\CGref{fig:rb-gonek}). The $\ell^2$ conditioning of the interior block at $\boldsymbol\gamma$ is moderate and roughly uniform up to $T=299$: condition numbers $17$ to $75$, smallest singular value $0.08$ to $0.15$. The coupling of the interior to the frozen edge, $\|\mathbb J_I^+\mathbb J_E\|_{\infty\to\infty}$, is between $0.25$ and $0.47$, and the shift it induces on the low ordinates decays from $2\cdot10^{-2}$ at $m=2$ to $10^{-5}$ for $m\ge10$.

\begin{figure}[tbp]
\centering
\includegraphics[width=0.9\textwidth]{fig18_gonek_vs_clock.png}
\caption{Distance to the zeta ordinates, in units of the local mean spacing $2\pi/\log(\gamma_k/2\pi)$, of the midpoint clock and of the Gonek lists $\boldsymbol\kappa^{(\mathrm X)}$ for $\mathrm X=3,8,20$ ($k\le160$ for $\mathrm X=3,20$; $k\le400$ for $\mathrm X=8$). The dashed line is the observed radius of the Newton basin. The ordinates $\gamma_k$ are used here only to measure the distance.}
\label{fig:rb-gonek}
\end{figure}

\subsubsection{Certificates at the Gonek list: the direct bounds fail}\label{subsec:rb-num-kantorovich}

At $\boldsymbol\kappa^{(8)}$, with the interior unknowns of stage $m$ and three choices of the right inverse, the quantities of Theorem~\CGref{thm:rb-certificate} were computed for $m=1,\dots,14$ ($T$ from $50$ to $414$).

\begingroup\small
\begin{longtable}{p{0.40\textwidth}p{0.22\textwidth}p{0.30\textwidth}}
\toprule
Right inverse and norm & Certificate index & Range over $m=1,\dots,14$\\
\midrule
square subsystem (tests chosen by pivoted QR), $\ell^\infty$ & $2\mathbb L_m^2\mathbb H_m\mathsf h_m$ & $56\to1.2\cdot10^{7}$\\
pseudo-inverse of the full system, $\ell^2$, $\beta_2=1/\sigma_{\min}$ & $\beta_2\mathrm K_2\delta_2$ (Kantorovich) & $0.47\to233$\\
pseudo-inverse of the full system, $\ell^\infty$, diagonal Hessian & $2\mathrm c_{\rm diag}\mathsf h$ & $1.3\to1.6\cdot10^{3}$\\
\bottomrule
\end{longtable}
\endgroup

None of the certificates is satisfied, not even at $T=50$. Yet one Gauss--Newton step from $\boldsymbol\kappa^{(8)}$ reduces $\max_k|\kappa_k-\gamma_k|$ by a factor between $4$ and $15$ at every height ($0.088\to0.016$ at $T=572$), and $\sigma_{\min}$ in $\ell^2$ stays between $1/17$ and $1/5$. The worst-case bounds lose a factor $10^3$ to $10^4$ through cancellation: $\mathbb H_m$ counts every unknown with the full weight $\int|\varphi|u^2$, while the actual second-order term of $\mathbb F_m$ along the Newton step is a sum of $k^*_m$ oscillating terms.

\subsubsection{Floating-point candidate certificates for exact square sections}\label{subsec:rb-num-certified}

The iterate of Section~\CGref{subsec:rb-num-init} was polished, stage by stage, by Newton's method on the square system ($n$ tests selected by pivoted QR on $\mathbb J^{\!T}$, $n$ interior unknowns) to residual $10^{-13}$. At the polished point $x_0$ the certificate of Theorem~\CGref{thm:rb-certificate} was evaluated in the sharper form that uses the separable structure of $F_m$: with $Q$ the inverse of the square Jacobian, $B$ the matrix of second derivatives $\partial^2_{\kappa_k}\mathsf c_\varphi$ at $x_0$, and $A_3(r)$ an entrywise bound for the third derivatives on $|x-x_0|\le\mathrm r$, the map $\mathbb T(x)=x-\mathbb Q\mathbb F(x)$ contracts the ball of radius $\mathrm r$ in $\ell^\infty$ if
\begin{equation}\label{eq:rb-num-banach}
 \mathsf l(\mathrm r):=\|\mathbb Q\mathbb J'\|_{\infty\to\infty}\,\mathrm r+\tfrac12\,\bigl\||\mathbb Q|\mathbb A^{(3)}(\mathrm r)\bigr\|_{\infty\to\infty}\,\mathrm r^2<1,\qquad \mathsf h\le(1-\mathsf l(\mathrm r))\,\mathrm r,
\end{equation}
in which case the ball contains an exact solution of the square system; here $\mathrm t_m=k^*_m$, so the right inverse of Theorem~\CGref{thm:rb-certificate} is the inverse. The margin of a stage is $\max_{\mathrm r}(1-\mathsf l(\mathrm r))\mathrm r$, the largest residual that \CGeqref{eq:rb-num-banach} could certify at that point.

\begingroup\small
\begin{longtable}{rrrrrrr}
\toprule
$m$ & $T_m$ & $n$ & $\mathbb L_m=\|Q\|_{\infty\to\infty}$ & margin & $\mathsf s_m$ on all tests $\le m$ & $\max_k|\kappa_k-\gamma_k|$\\
\midrule
1 & 50 & 4 & $41$ & $7.0\cdot10^{-3}$ & $7.7\cdot10^{-8}$ & $7\cdot10^{-7}$\\
4 & 81 & 9 & $50$ & $3.1\cdot10^{-3}$ & $2.6\cdot10^{-7}$ & $1\cdot10^{-6}$\\
8 & 156 & 26 & $150$ & $7.5\cdot10^{-4}$ & $9.6\cdot10^{-7}$ & $1\cdot10^{-5}$\\
12 & 299 & 68 & $460$ & $2.2\cdot10^{-4}$ & $5.5\cdot10^{-6}$ & $4\cdot10^{-5}$\\
15 & 487 & 133 & $1600$ & $6.7\cdot10^{-5}$ & $3.3\cdot10^{-3}$ & $4\cdot10^{-2}$\\
16 & 572 & 165 & -- & square Jacobian singular & -- & --\\
\bottomrule
\end{longtable}
\endgroup

For stages $1$ to $15$ the floating-point evaluation of \CGeqref{eq:rb-num-banach} holds with $\mathrm r$ equal to the polishing residual ($10^{-15}$ to $5\cdot10^{-13}$). If the inequalities held in exact arithmetic for compactly supported tests, Theorem~\CGref{thm:rb-certificate} would give an exact solution of the square section within $\mathrm r$ of $x_0$, inside the tube with $C=1$; what has been established is the floating-point inequality, without rounding bounds and with Gaussian tests, so these stages carry candidate certificates and not verified ones. The margins exceed the residuals by $10^{7}$ or more, which is far above the expected rounding error, so the same inequalities are expected, but not verified, to survive interval arithmetic. The candidate solutions are not unique: at $m=15$ the polished point is $4\cdot10^{-2}$ from $\boldsymbol\gamma$. At $m=16$ the square Jacobian is singular to working precision and the stage has no candidate certificate. The scaling with height (Figure~\CGref{fig:rb-scaling}) is
\begin{equation}\label{eq:rb-num-scaling}
 \mathbb L_m\sim T^{1.8},\qquad\|\mathbb Q\mathbb J'\|_{\infty\to\infty}\sim T^{2.1},\qquad\text{margin}\sim T^{-2.1}\qquad(T\le487).
\end{equation}

\subsubsection{The entry step}\label{subsec:rb-num-entry}

The open point of the staged construction, in the form that Theorem~\CGref{thm:rb-staged} makes precise, is whether the residual $\mathsf h_{m+1}$ of the new tests at the start of stage $m+1$ is below the margin of that stage. This was tested as follows: start from the polished point of stage $m-1$ (Section~\CGref{subsec:rb-num-certified}), put the new band of unknowns (those with $\tau_k\in(2.4/\mathrm w_{m-1},2.4/\mathrm w_m]$) and everything beyond at their Gonek values, and evaluate \CGeqref{eq:rb-num-banach} (i) for the band alone, with the new tests of level $m$, and (ii) for the full square system of stage $m$, without polishing.

The band-only test passes (in the same floating-point sense) for $T\le96$ and fails from $T=113$ on; its margin falls from $9\cdot10^{-2}$ to $3\cdot10^{-3}$ while the Gonek residual of the band stays between $10^{-2}$ and $3\cdot10^{-1}$ (the Gonek error in the band is $0.02$ to $0.09$). The full-system test passes only at $T=59$. The margin decreases as a power of $T$ by \CGeqref{eq:rb-num-scaling}, whereas the error of the Gonek lists with fixed $\mathrm X$ does not decrease with $T$, and with $\mathrm X=T^{\mathrm x}$ it would decrease at most like $1/\log T$. Within the range computed, no worst-case certificate in $\ell^\infty$ closes the entry step, and the trend of the constants says that none will at larger $T$. This is a numerical statement, not a theorem.

\begin{figure}[tbp]
\centering
\includegraphics[width=0.9\textwidth]{fig19_cert_scaling.png}
\caption{Constants of the Banach certificate \CGeqref{eq:rb-num-banach} for the square sections of Section~\CGref{subsec:rb-num-certified} (solid lines, with fitted slopes in $\log T$), and the band-only entry test of Section~\CGref{subsec:rb-num-entry} (dashed): the reciprocal of the band margin and the reciprocal of the residual of the new band at its Gonek values. The entry test passes, in floating point, where the triangles lie below the squares.}
\label{fig:rb-scaling}
\end{figure}

\subsubsection{Reading of the computations}\label{subsec:rb-num-reading}

\begin{enumerate}[label=(\roman*)]
\item The constructions that use only the clock and the primes (the Gonek lists, the square sections, the candidate certificates) behave as predicted by the tube lemma and the certificate theorem, and the candidate certificates hold in floating point up to $T=487$. Heuristically, such sections would exist equally if the Riemann hypothesis were false: a zero $\frac12+\mathrm b_0+i\gamma_0$ off the line contributes $2\cosh(\mathrm b_0u)\cos(\gamma_0u)$ to the prime side, which real ordinates near $\gamma_0$ can imitate against tests supported in $(0,u^*)$ as long as $\mathrm b_0u^*$ is small. On this heuristic the information about the hypothesis sits at heights $T\gtrsim(C\log\gamma_0)^{1/\mathrm b_0}$, which no finite computation reaches; this is a reading of the computations, not a theorem.
\item The quantities $\mathbb L_m$ and $\|QB\|$ grow polynomially in $T$. This is the quantitative face of the obstruction theorem of Section~\CGref{subsec:rb-obstruction}: no height-uniform linear stability, hence no height-uniform Newton basin. A certificate for the infinite construction needs a basin that shrinks no faster than the initial error of a $\gamma$-free list decreases, and the computed trend is $T^{-2}$ against $1/\log T$.
\item The observed contraction of Gauss--Newton (factor $4$ to $15$ per step at every height) is a mean-square phenomenon that the worst-case constants do not see. A proof of the entry step would have to bound $\|\mathbb Q\mathbb J'\,\mathrm{diag}(\mathrm e)\,\mathrm e\|_\infty$ for the actual error vector $\mathrm e$ of the Gonek list, not for an arbitrary vector of the same norm; that error vector is $\gamma$-free, and its residual $\mathbb F_m(\boldsymbol\kappa^{(\mathrm X)})$ is controlled through the explicit formula by the moments of $S-S_{\mathrm X}$, which are unconditional (Selberg, Fujii). Whether this closes is not known; any proof along this line would contain a proof of the Riemann hypothesis, by Corollary~\CGref{cor:rb-Bprime}.
\end{enumerate}

%======================================================================
\subsection{Status of the statements}\label{subsec:rb-status}

\begingroup\small
\begin{longtable}{p{0.42\textwidth}p{0.5\textwidth}}
\toprule
Statement & Status\\
\midrule
Lemma~\CGref{lem:rb-canonical-box}, Corollary~\CGref{cor:rb-D-implies-B} & proved; conditional on (D) as stated\\
Proposition~\CGref{prop:rb-compactness}, Lemma~\CGref{lem:rb-matching} & proved, unconditional\\
Hypothesis~\CGref{hyp:rb-U} & open; it extends \CGref{thm:free-masses}; not needed with prescribed masses\\
Corollary~\CGref{cor:rb-selection} & proved, assuming B and U\\
Lemmas~\CGref{lem:rb-archimedean}, \CGref{lem:rb-laplace} & proved, unconditional (the former Hypothesis~(L))\\
Proposition~\CGref{prop:rb-prescribed-uniqueness}, Corollary~\CGref{cor:rb-Bprime} & proved, unconditional: RH $\Leftrightarrow$ B$'(C)$; classical in substance\\
Lemma~\CGref{lem:rb-lipschitz}, Theorems~\CGref{thm:rb-certificate}, \CGref{thm:rb-staged}, Proposition~\CGref{prop:rb-graded} & proved; their hypotheses (A$_m$) for all $m$ and (B) are open and imply RH\\
Heuristic~\CGref{heur:rb-frame} & not proved; its point model is not a test of $C_c^\infty(0,\infty)$\\
Lemma~\CGref{lem:rb-gonek-tube} & proved, unconditional; uses \cite{Gonek2012} only for the interpretation of the list\\
(G1), (G2) of Remark~\CGref{rem:rb-gonek-status} & open; together they imply B$'$ and RH\\
Lemmas~\CGref{lem:rb-PL}--\CGref{lem:rb-argument}, Theorem~\CGref{thm:rb-count}, Corollary~\CGref{cor:rb-displacement} & proved, unconditional; uses \CGref{thm:halfplane-RVM}; by Lemma~\CGref{lem:rb-upsilon-infinite} the hypothesis $\upsilon<\infty$ holds for a finite section only when all masses vanish, so the bound is a template, not a bound for sections with primes (Remark~\CGref{rem:rb-count-scope})\\
Heuristics~\CGref{heur:rb-window}, \CGref{heur:rb-informative}, \CGref{heur:rb-graduated} & not proved\\
Lemma~\CGref{lem:rb-generating}, Theorem~\CGref{thm:rb-criterion}, Corollary~\CGref{cor:rb-riesz} & proved; Corollary~\CGref{cor:rb-riesz} uses \cite{AvdoninIvanov2002}\\
Corollary~\CGref{cor:rb-kadec-avdonin}, Theorem~\CGref{thm:rb-variable-N}, Proposition~\CGref{prop:rb-variable-blocks} & proved; novelty not checked\\
Lemma~\CGref{lem:rb-psi-bound}, Proposition~\CGref{prop:rb-O-blocks}, Proposition~\CGref{prop:rb-weight}, Corollary~\CGref{cor:rb-one-condition} & proved, unconditional\\
Lemma~\CGref{lem:rb-smoothing}, Theorem~\CGref{thm:rb-obstruction}, Corollary~\CGref{cor:rb-obstruction} & proved, unconditional; uses the upper bound contained in \cite[Main Theorem]{Fujii1974}, \cite{Selberg1946,Titchmarsh1986,Garnett2007}\\
Section~\CGref{subsec:rb-numerics} & computations in floating point with Gaussian tests; candidate certificates, not verified; not part of any proof\\
\bottomrule
\end{longtable}
\endgroup

%======================================================================

\section{Numerical illustration with \texorpdfstring{$2\,001\,052$}{2001052} zeros}
\label{sec:numerics}

This section shows what is being constructed and checks that the theorems are consistent with computation.  It uses the ordinates $\gamma_1,\dots,\gamma_N$ of the first $N=2\,001\,052$ zeros of $\zeta$ (file \texttt{zeross.bin}, $\gamma_N=1\,132\,490.6587$) and the midpoint ordinates $\tau_1,\dots,\tau_N$ computed from \CGeqref{eq:lambert} (file \texttt{zLamb.bin}, $\tau_N=1\,132\,490.7226$).  The programs are listed in Appendix~\CGref{app:code} and their complete console output in Appendix~\CGref{app:output}; all figures have linear axes.  Each item is an \emph{experiment}: it reports a computation and does not replace a proof.

Products over the ordinates are evaluated as follows, in double precision throughout: the first $10^5$ factors one by one, the factors $10^5<k\le N$ through power sums of $(\tfrac14+o_k^2)^{-1}$, and the factors $k>N$ by the midpoint rule for the clock, which gives the tail $\frac1{2\pi}\bigl(q\frac{L_U+1}U-\frac{q^2}2\frac{3L_U+1}{9U^3}+\frac{q^3}3\frac{5L_U+1}{25U^5}\bigr)$ with $\NR(U)=N$, $L_U=\log(U/2\pi)$.  For the products over the $\gamma_k$ ($P_{Z1}$, $\zeta_{Z1}$) this means that the unknown tail of the zeta ordinates beyond the file is \emph{replaced by the clock tail}; every agreement reported for $\zeta_{Z1}$ against $\zeta$ includes that substitution and is an agreement of the construction with itself beyond height $\gamma_N$.  ``All $N$ ordinates'' or ``all zeros'' below always means all zeros in the supplied file, $\gamma_k\le\gamma_N$.  No quantity in this section is certified: the words ``exact'' and ``certified'' are not used for floating-point values, and the console output's \texttt{CERTIFIED} (Appendix~\CGref{app:output}) is a floating-point test, as explained in Remark~\CGref{rem:rb-certificate-errors}.

\subsection{The clock and the zeros}
\label{subsec:num-clock}

\begin{experiment}[Clock, midpoints and displacements]\label{exp:clock}
\emph{Object:} $\NR$, $N_\zeta$, $\tau_k$, $\eta_k=\tau_k-\gamma_k$.  \emph{Range:} $k\le N$.  \emph{Expected:} \CGeqref{eq:midpoint-condition} to rounding, and $|\eta_k|\le C_\eta$ (Lemma~\CGref{lem:bounded-displacement}).  \emph{Observed:} $\max_k|\NR(\tau_k)-(k-\tfrac12)|=6.985\cdot10^{-10}$; $\max_k|\eta_k|=0.9937$ (at $k=34$), mean $|\eta_k|=0.1512$, root mean square $0.1881$, mean $\eta_k=1.7\cdot10^{-7}$.  By blocks of $k$, $\max|\eta_k|$ is $0.840,0.994,0.915,0.834,0.786,0.763,0.717$ on $[1,10]$, $[11,10^2]$, $[10^2,10^3]$, \dots, $[10^6,N]$.
\end{experiment}

Figure~\CGref{fig:clock-zoom} shows how the clock places the midpoint ordinates.  The staircase $N_\zeta$ jumps by one at each $\gamma_k$; the smooth clock $\NR$ crosses the half-integer level $k-\tfrac12$ at $\tau_k$.  Since $N_\zeta-\NR=S(T)+O(1/T)$ (Figure~\CGref{fig:remainder}), the horizontal distance between the open square $(\gamma_k,k-\tfrac12)$ and the red point $(\tau_k,k-\tfrac12)$ is the remainder $S$ converted into ordinate units by the slope $\NR'=\frac1{2\pi}\log\frac T{2\pi}$.  The right panel repeats the picture at height $1.13\cdot10^6$, where the mean spacing is $0.519$: every $\tau_k$ still lies next to its $\gamma_k$.

\begin{figure}[htbp]
\centering
\includegraphics[width=0.86\textwidth]{fig00_clock_zoom.png}
\caption{The midpoint clock $\NR$ (red), the zero staircase $N_\zeta$ (black), the midpoint ordinates $\tau_k$ (red points) and the points $(\gamma_k,k-\frac12)$ (squares).  Left: the first nine zeros.  Right: the last twelve zeros of the file.}
\label{fig:clock-zoom}
\end{figure}

\begin{figure}[htbp]
\centering
\includegraphics[width=0.86\textwidth]{fig00b_remainder.png}
\caption{The remainder $N_\zeta(T)-\NR(T)=S(T)+O(1/T)$ of \CGeqref{eq:RvM-full} for $T\le200$.}
\label{fig:remainder}
\end{figure}

\subsection{The \texorpdfstring{$\lambda$}{lambda}-family and the endpoint}
\label{subsec:num-family}

Figure~\CGref{fig:family} shows the core of the method.  The lower left panel draws the moving ordinates $\gamma_k(\lambda)$ in clock units, $\NR(\gamma_k(\lambda))-(k-\frac12)$: at $\lambda=1$ all lines meet at $0$ (the $\tau_k$), at $\lambda=0$ they sit at the ordinates of $\zeta$; for $\lambda>0$ the ordinates $\gamma_k(\lambda)$ remain strictly increasing in $k$ (Theorem~\CGref{thm:path}).  The lower right panel draws the real functions $Z_\lambda(t)=e^{i\vartheta(t)}\zeta_\lambda(\frac12+it)$ for $\lambda=0,0.25,0.5,0.75,1$ together with Hardy's function $Z(t)=e^{i\vartheta(t)}\zeta(\frac12+it)$: the zeros of every $Z_\lambda$ are real, they move continuously with $\lambda$, and the curve for $\lambda=0$ lies on top of $Z$.

\begin{figure}[htbp]
\centering
\includegraphics[width=0.80\textwidth]{fig01_clock_family.png}
\caption{Top: $\eta_k$ for every 50th $k\le N$ and its distribution.  Bottom: moving ordinates in clock units, $k\le30$, and the functions $Z_\lambda(t)$ on the critical line.}
\label{fig:family}
\end{figure}

\begin{experiment}[The endpoint $\zeta_{Z1}$ against $\zeta$]\label{exp:Z1}
\emph{Object:} $\zeta_{Z1}=P_{Z1}/G$ computed from the $N$ ordinates $\gamma_k$, and $\zeta$ computed by mpmath.  \emph{Range:} $\sigma\in\{2,1.25,0.75\}$, $1\le t\le2000$ ($499$ points per line), and the critical line $0.5\le t\le60$.  \emph{Expected:} if $\zeta_{Z1}=\zeta$, the ratio is $1$ up to the truncation of the product.  \emph{Observed:} $\max|\zeta_{Z1}/\zeta-1|=8.7\cdot10^{-10}$, $1.1\cdot10^{-9}$, $2.6\cdot10^{-9}$ on the three lines, and $\max|Z_{Z1}-Z|/\max(1,|Z|)=8.4\cdot10^{-10}$ on the critical line (Figure~\CGref{fig:Z1}); the tail beyond $\gamma_N$ is the clock tail, so the $10^{-9}$ figure measures the file truncation and the tail substitution, not (D).
\end{experiment}

\begin{figure}[htbp]
\centering
\includegraphics[width=0.86\textwidth]{fig02_zetaZ1_vs_zeta.png}
\caption{Left: $|\zeta_{Z1}/\zeta-1|$ on three vertical lines.  Right: $Z_{Z1}$ and Hardy's $Z$ on the critical line.}
\label{fig:Z1}
\end{figure}

Experiment~\CGref{exp:Z1} shows numerical agreement at the displayed truncation precision.  This does not test (D): Proposition~\CGref{prop:varpi-verified} already forces $|\log\varpi_{Z1}|<4\cdot10^{-13}$ from verified zeros, while the observed $10^{-9}$ scale is dominated by truncating the ordered clock tail.

\begin{experiment}[The value at infinity: the product identity]\label{exp:product-identity}
\emph{Object:} the ordered partial products $P_K=\prod_{k\le K}(\frac14+\gamma_k^2)/(\frac14+\tau_k^2)$ of \CGeqref{eq:product-identity}.  \emph{Range:} $K\le N$.  \emph{Expected:} by Corollary~\CGref{cor:product-identity}, $\zeta_{Z1}=\zeta$ is the statement $P_K\to\varpi_{\RVM}$.  \emph{Observed:} $P_K=0.96256$, $0.96248$, $0.96229$, $0.961951$, $0.961930$, $0.9619418189$, $0.9619415390$ for $K=10,10^2,\dots,10^6,N$, against $\varpi_{\RVM}=0.9619415807$ from Experiment~\CGref{exp:transition}; for $K\ge10^6$, $P_K$ stays in $[0.9619389,0.9619440]$ with mean $0.9619416633$ (Figure~\CGref{fig:product-identity}).  Proposition~\CGref{prop:varpi-verified} already bounds the endpoint discrepancy by $4\cdot10^{-13}$ under the verified-zero input, so the $10^{-9}$--$10^{-6}$ discrepancies in these finite products are truncation/tail errors and are not numerical evidence for (D).  The flow at infinity \CGeqref{eq:Upsilon} is $\Upsilon(\lambda)=0.0436950$, $0.0412445$, $0.0387987$, $0.0363572$, $0.0339198$ for $\lambda=0,\frac14,\frac12,\frac34,1$, decreasing as in \CGeqref{eq:varpi-derivative} with slope close to $-\sum_k2\eta_k^2/\gamma_k^2=-0.0098211$, and $\int_0^1\Upsilon\,d\lambda=\sum_{k\le N}\log\frac{1/4+\tau_k^2}{1/4+\gamma_k^2}=0.0388016004$ against $\log L=0.0388015571$.
\end{experiment}

\begin{figure}[htbp]
\centering
\includegraphics[width=0.86\textwidth]{fig11_product_identity.png}
\caption{The partial products $P_K$ of \CGeqref{eq:product-identity} against $\varpi_{\RVM}$ (dashed).  Left: $K\le10^4$.  Right: $(P_K-\varpi_{\RVM})\cdot10^6$ for $10^5\le K\le N$.  The value of the endpoint at $+\infty$ is $\varpi_{Z1}=\varpi_{\RVM}/\lim_KP_K$ (Theorem~\CGref{thm:Z1-clock}).}
\label{fig:product-identity}
\end{figure}

\subsection{The transition theorem and \texorpdfstring{$\zeta_{\RVM}$}{zeta\_RVM} against \texorpdfstring{$\zeta$}{zeta}}
\label{subsec:num-transition}

\begin{experiment}[Transition theorem]\label{exp:transition}
\emph{Object:} $\zeta_{\RVM}$, $w(s)$, $\varpi_{\RVM}$, $L$.  \emph{Range:} $\sigma'\in\{0.25,0.5,1.5\}$, $10\le t\le10^4$.  \emph{Expected:} Theorem~\CGref{thm:transition} and Corollary~\CGref{cor:vertical-limit}: $t\,|\zeta_{\RVM}/(\varpi_{\RVM}(1+w))-1|$ bounded, $|w|=(t/2\pi)^{-\sigma'}(1+O(1/t))$.  \emph{Observed:} Richardson extrapolation on $\Re s=2$ (which presupposes an expansion in $1/T$ and $1/T^2$ finer than the $O(1/T)$ proved) gives $\varpi_{\RVM}=0.9619415807$ and $L=1.0395641690$, values without error bound; $\max_{t>100}t\,|\zeta_{\RVM}/(\varpi_{\RVM}(1+w))-1|=0.0617,0.0618,0.0620$ for the three values of $\sigma'$, and $0.0790$ for $\varpi_{\RVM}(1+\chi)$ at $\sigma'=0.25$; $|w|\,t^{\sigma'}=1.5832$, $2.5066$, $15.66$--$15.75$, i.e.\ $(2\pi)^{\sigma'}$.
\end{experiment}

\begin{figure}[htbp]
\centering
\includegraphics[width=0.80\textwidth]{fig03_transition.png}
\caption{Top: the transition error multiplied by $t$, and the rescaled vertical limit $t^{\min(\sigma',1)}|1/\zeta_{\RVM}-L|$.  Bottom: $|\zeta|$ and $|\zeta_{\RVM}|$ on $\sigma=2$ and $\sigma=0.75$.}
\label{fig:transition}
\end{figure}

The bottom panels of Figure~\CGref{fig:transition} compare $|\zeta|$ and $|\zeta_{\RVM}|$ on two vertical lines: $\zeta$ carries the Euler product and oscillates with the primes, whereas $\zeta_{\RVM}$ approaches the constant $\varpi_{\RVM}$ on $\sigma=2$ and $\varpi_{\RVM}|1+\chi|$ on $\sigma=0.75$.  This is the analytic content of Theorem~\CGref{thm:no-Dirichlet-RVM}.

\subsection{Exponential sums}
\label{subsec:num-expsums}

\begin{experiment}[Exponential sums over the two spectra]\label{exp:expsums}
\emph{Object:} $S(N;\log x)=\sum_{k\le N}x^{io_k}$ for $o_k=\tau_k$ and $o_k=\gamma_k$.  \emph{Range:} $N\le2\,001\,052$, $x\in\{2,3,4,6,10\}$.  \emph{Expected:} $|S_{\RVM}|\ll1+\log N$ (Theorem~\CGref{thm:exp-sum}); for $\gamma_k$, Landau's term $-\frac T{2\pi}\frac{\Lambda(x)}{\sqrt x}$.  \emph{Observed:} $\max_N|S_{\RVM}(N;\log x)|=3.57,3.98,3.95,4.53,9.70$; at $N=2\,001\,052$, $|\sum x^{i\gamma_k}|=88\,301$, $114\,268$, $62\,435$, $0.26$, $2.42$, against Landau's $88\,342$, $114\,324$, $62\,467$, $0$, $0$.
\end{experiment}

\begin{figure}[htbp]
\centering
\includegraphics[width=0.86\textwidth]{fig04_exponential_sums.png}
\caption{The midpoint spectrum carries no prime frequencies (left); the spectrum of $\zeta$ resonates at prime powers (right).}
\label{fig:expsums}
\end{figure}

\subsection{Prime counts: \texorpdfstring{$\pi$}{pi} and \texorpdfstring{$\pi_{\RVM}$}{pi\_RVM}}
\label{subsec:num-primes}

\begin{experiment}[Prime counts and RVM primes]\label{exp:primes}
\emph{Object:} $J_{\RVM}$, $\pi_{\RVM}$, $c_n^{\RVM}$, $p_j^{\RVM}$ against $J$, $\pi$, the primes.  \emph{Range:} $2\le x\le100$ with all $N$ ordinates of each list ($3000$ points); $2\le x\le1000$ with the first $K=17\,061$ ordinates and a smooth cut ($40\,000$ points).  \emph{Expected:} the explicit formula with the $\gamma_k$ reproduces $\pi(x)$; $\pi_{\RVM}$ is continuous and $\pi_{\RVM}(x)-\li(x)=O(\sqrt x(\log\log x)^2/\log x)$ (Theorem~\CGref{thm:Lambert-estimate}).  \emph{Observed:} away from the primes, $|\pi_{\gamma}-\pi|\le0.40$ on $[2,100]$; $\pi_{\RVM}(100)=26.08$, $\pi(100)=25$, $\operatorname{li}(100)=30.13$; $\pi_{\RVM}(1000)=169.98$, $\pi(1000)=168$, $\operatorname{li}(1000)=177.61$.  $\pi_{\RVM}$ is not monotone: it carries a ripple of unit scale, with minimum grid increment $-0.25$ on $[2,100]$.  The first crossings of the levels $j=1,\dots,8$ are $p_j^{\RVM}=2.21,3.16,5.19,7.16,11.14,14.02,17.26,19.18$ against $p_j=2,3,5,7,11,13,17,19$.
\end{experiment}

\begin{figure}[htbp]
\centering
\includegraphics[width=0.86\textwidth]{fig05_prime_counts.png}
\caption{$\pi(x)$ (black staircase), the explicit formula with the zeros of $\zeta$ (blue), $\pi_{\RVM}(x)$ (red) and $\operatorname{li}(x)-\operatorname{li}(2)$ (green).  The distance between the red and green curves is the quantity bounded in Theorem~\CGref{thm:Lambert-estimate}; the ratios are in Figure~\CGref{fig:pi-ratios}.}
\label{fig:prime-counts}
\end{figure}

\begin{figure}[htbp]
\centering
\includegraphics[width=0.86\textwidth]{fig06_rvm_primes.png}
\caption{Left: cell masses $c_n^{\RVM}$, the prime indicator and $1/\log n$.  Right: real RVM prime levels against the primes up to $1000$.}
\label{fig:rvm-primes}
\end{figure}

Figure~\CGref{fig:prime-counts} shows $\pi$ and $\pi_{\RVM}$ in the usual format.  The RVM count follows the prime staircase much more closely than $\operatorname{li}$: replacing each $\gamma_k$ by $\tau_k$ moves every zero by less than one unit (Experiment~\CGref{exp:clock}), and the low ordinates, which shape $\pi(x)$ for small $x$, are moved least in relative terms.  The steps are replaced by a continuous ripple, so the cell masses $c_n^{\RVM}$ of Figure~\CGref{fig:rvm-primes} are real numbers, occasionally negative, and are not a prime indicator.

\begin{experiment}[The three ratios and the square-root scale]\label{exp:pi-ratios}
\emph{Object:} $\li/\pi_{\RVM}$, $\pi/\li$, $\pi/\pi_{\RVM}$, and $(\pi_{\RVM}-\li)/B$, $(\pi-\li)/B$ with $B(x)=\sqrt x(\log\log x)^2/\log x$.  \emph{Range:} the grid of Experiment~\CGref{exp:primes} on $[8,1000]$ ($K=17\,061$ ordinates).  \emph{Expected:} the three ratios tend to $1$ (Corollary~\CGref{cor:three-limits}) and $(\pi_{\RVM}-\li)/B$ stays bounded (Theorem~\CGref{thm:Lambert-estimate}).  \emph{Observed:} at $x=100,500,1000$, $\li/\pi_{\RVM}=1.155,1.086,1.045$, $\pi/\li=0.830,0.933,0.946$, $\pi/\pi_{\RVM}=0.959,1.013,0.988$; $\max|\pi_{\RVM}-\li|/B=2.72$ on $[8,1000]$, $1.35$ on $[50,1000]$ and $0.72$ on $[500,1000]$, while $\max|\pi-\li|/B=2.45$ on $[8,1000]$ (Figure~\CGref{fig:pi-ratios}).
\end{experiment}

\begin{figure}[htbp]
\centering
\includegraphics[width=0.86\textwidth]{fig10_prime_ratios.png}
\caption{Left: the three ratios of Corollary~\CGref{cor:three-limits}.  Right: $\pi_{\RVM}-\li$ and $\pi-\li$ in units of the scale $B(x)$ of Theorem~\CGref{thm:Lambert-estimate}.  The ratio $\pi/\pi_{\RVM}$ (black) is the closest of the three to $1$; the ratios with $\li$ approach $1$ at the rate $\tfrac12\li(\sqrt x)/\li(x)$ of the bias $J-\pi$.}
\label{fig:pi-ratios}
\end{figure}

On $[2,1000]$ the two normalized errors in Figure~\CGref{fig:pi-ratios} have the same size and sign: both counts lie below $\li$ by about $\tfrac12\li(\sqrt x)$, the $m=2$ term of the M\"obius inversion \CGeqref{eq:pi-minus-li}.  This is why $\pi/\pi_{\RVM}$ is much closer to $1$ than either ratio with $\li$.

\subsection{The Mertens functions and the ternary division}
\label{subsec:num-mertens}

The Mertens function is the sum of the ternary sequence $\mu(n)\in\{-1,0,+1\}$.  Its values are specified by two pieces of data:
\begin{enumerate}[label=\textup{(\alph*)}]
\item \emph{the zero set} $\{n:\mu(n)=0\}$, the non-squarefree integers, of density $1-6/\pi^2=0.392073$;
\item \emph{the signs} on the squarefree integers, $\mu(n)=(-1)^{\omega(n)}$, each sign of density $3/\pi^2=0.303964$.
\end{enumerate}
The ternary walk $\mu_h$ of \CGeqref{eq:intro-muh}--\CGeqref{eq:intro-quantizer} keeps (a) exactly, $\mu_h(n)=0\Leftrightarrow\mu(n)=0$, and replaces (b): at a squarefree $n$ it steps up when $M_h(n-1)$ lies below the target $\mathcal M(n)=M_{\RVM}^{\crit}(n)$ and down otherwise.  By construction the walk follows its target within one unit, so its closeness to $M_{\RVM}^{\crit}$ is a property of the quantizer and carries no information about $M$; any coincidence of signs between $\mu_h$ and $\mu$ is reported as an observation without a statistical analysis and is not claimed to have theoretical content.  Thus
\begin{equation}\label{eq:ternary-update}
 M_h(n)=M_h(n-1)+\mu(n)^2\,\varsigma_n,\qquad |M_h(n)-\mathcal M(n)|\ \text{small},
\end{equation}
so that $M_h$ is an integer walk with the zero set of $\mu$ that follows $M_{\RVM}^{\crit}$, in the same way as $M$ is the integer walk with the zero set of $\mu$ whose signs are $(-1)^{\omega(n)}$.

\begin{table}[htbp]
\centering
\small
\begin{tabular}{rrrrrcrrr}
\toprule
$n$ & $\mu(n)$ & $\mu(n)^2$ & $\mathcal M(n)$ & $M_h(n-1)$ & $\varsigma_n$ & $\mu_h(n)$ & $M_h(n)$ & $M(n)$\\
\midrule
1 & 1 & 1 & 0.2925 & 0 & $+1$ & 1 & 1 & 1\\
2 & $-1$ & 1 & 0.4563 & 1 & $-1$ & $-1$ & 0 & 0\\
3 & $-1$ & 1 & $-0.6523$ & 0 & $-1$ & $-1$ & $-1$ & $-1$\\
4 & 0 & 0 & $-0.8720$ & $-1$ & -- & 0 & $-1$ & $-1$\\
5 & $-1$ & 1 & $-1.7255$ & $-1$ & $-1$ & $-1$ & $-2$ & $-2$\\
6 & 1 & 1 & $-1.3063$ & $-2$ & $+1$ & 1 & $-1$ & $-1$\\
7 & $-1$ & 1 & $-1.5931$ & $-1$ & $-1$ & $-1$ & $-2$ & $-2$\\
8 & 0 & 0 & $-2.4029$ & $-2$ & -- & 0 & $-2$ & $-2$\\
9 & 0 & 0 & $-1.3292$ & $-2$ & -- & 0 & $-2$ & $-2$\\
10 & 1 & 1 & $-1.3280$ & $-2$ & $+1$ & 1 & $-1$ & $-1$\\
11 & $-1$ & 1 & $-2.3272$ & $-1$ & $-1$ & $-1$ & $-2$ & $-2$\\
12 & 0 & 0 & $-2.2067$ & $-2$ & -- & 0 & $-2$ & $-2$\\
13 & $-1$ & 1 & $-1.8468$ & $-2$ & $+1$ & 1 & $-1$ & $-3$\\
14 & 1 & 1 & $-1.9324$ & $-1$ & $-1$ & $-1$ & $-2$ & $-2$\\
15 & 1 & 1 & $-1.5435$ & $-2$ & $+1$ & 1 & $-1$ & $-1$\\
\bottomrule
\end{tabular}
\caption{The ternary division for $n\le15$ (the first $40$ values are in Appendix~\CGref{app:out-s05-mertens}).  The zero columns of $\mu$ and $\mu_h$ coincide; the signs of $\mu_h$ are chosen by comparing $M_h(n-1)$ with $\mathcal M(n)$.  For $n\le12$ the two walks agree; at $n=13$ the target $-1.85$ lies above $M_h(12)=-2$, so $\mu_h(13)=+1$ while $\mu(13)=-1$.}
\label{tab:ternary}
\end{table}

\begin{experiment}[Mertens functions and the ternary walk]\label{exp:mertens}
\emph{Object:} $M$, $M_{\RVM}^{\crit}$, $M_h$; the explicit formulas.  \emph{Range:} $n\le10^6$; explicit formulas on $[2,2000]$ with all $N$ ordinates of each list.  \emph{Expected:} the explicit formula with the zeros of $\zeta$ reproduces $M$; $M_{\RVM}^{\crit}(x)=O(x^{1/2+\eps})$ (Theorem~\CGref{thm:unsmoothed-RVM}); $M_h$ tracks $M_{\RVM}^{\crit}$ with the densities of $\mu$.  \emph{Observed:}
\begin{itemize}
\item explicit formula for $M$ with the $\gamma_k$: root mean square error $0.0005$ on $[2,2000]$ off the integers;
\item $\zeta_{\RVM}'(\rho_k^{\RVM})/\log(\tau_k/2\pi)=0.96194$ at $k=1000$, against $\varpi_{\RVM}=0.96194$ (Corollary~\CGref{cor:derivative-zeros}); the values $|\zeta'(\rho_k)|/\log(\gamma_k/2\pi)$ are spread over $(0,3)$ (Figure~\CGref{fig:explicit});
\item spectral sums $2\sum|a|^2$: $0.029034$ for $\zeta$, $0.027203$ for the midpoint system;
\item the integer target agrees with the full explicit formula for $M_{\RVM}^{\crit}$ to $0.019$ at $100$ integers $\le10^5$;
\item $\max|M_h-\mathcal M|=2.62$ on $[1,10^6]$ and $1.65$ on $[10^4,10^6]$;
\item densities of $(-1,0,+1)$ for $\mu$: $0.303857$, $0.392074$, $0.304069$;
\item densities of $(-1,0,+1)$ for $\mu_h$: $0.304003$, $0.392074$, $0.303923$;
\item $\mu_h(n)=\mu(n)$ at $50.03\%$ of the squarefree $n\le10^6$, and $\sum_{n\le10^6}\mu(n)\mu_h(n)/\#\{n\le10^6:\mu(n)\ne0\}=0.0007$: the empirical normalized correlation of the sign sequences is $0.0007$ on this range (no asymptotic statement is made), although both walks have the same zero set and comparable size;
\item $\max_{10\le n\le10^6}|W(n)|/\sqrt n=0.832,0.754,0.688$ for $W=M,M_{\RVM}^{\crit},M_h$;
\item energies $\frac1N\sum_{n\le N}W(n)^2$ at $N=10^6$: $14\,851$, $12\,538$, $12\,538$.
\end{itemize}
\end{experiment}

\begin{figure}[htbp]
\centering
\includegraphics[width=0.86\textwidth]{fig07_mertens.png}
\caption{Left: $M$ (black), $M_{\RVM}^{\crit}$ (red) and $M_h$ (green) at every integer $n\le2000$.  Right: range of $M$ and mean of $M_{\RVM}^{\crit}$ in blocks of $500$ integers up to $10^6$, with $\pm\sqrt n$.}
\label{fig:mertens}
\end{figure}

\begin{figure}[htbp]
\centering
\includegraphics[width=0.86\textwidth]{fig08_ternary_walk.png}
\caption{Left: $\mu(n)$ and $\mu_h(n)$ for $n\le60$: same zeros, different signs.  Right: $M$, the target $M_{\RVM}^{\crit}$ and $M_h$ for $n\le300$.}
\label{fig:ternary}
\end{figure}

\begin{figure}[htbp]
\centering
\includegraphics[width=0.86\textwidth]{fig09_explicit_weights.png}
\caption{Left: $M(x)$, its explicit formula from the zeros of $\zeta$ and $M_{\RVM}^{\crit}(x)$ from the $\tau_k$ on $[2,300]$.  Right: normalized derivatives at the zeros: erratic for $\zeta$, constant $\varpi_{\RVM}$ for $\zeta_{\RVM}$.}
\label{fig:explicit}
\end{figure}

Figures~\CGref{fig:mertens}--\CGref{fig:explicit} summarize the Mertens side.  $M_{\RVM}^{\crit}$ has the size and the visual texture of $M$, it stays inside $\pm\sqrt n$ in agreement with the square-root-scale behaviour of Theorem~\CGref{thm:unsmoothed-RVM}, and its residues are deterministic; $M_h$ is its ternary integer version with the zero set of $\mu$.  The passage from these synthetic objects to $M$ itself is not numerical: it is the identity $\zeta_{Z1}=\zeta$ of Step~3, which Experiment~\CGref{exp:Z1} is consistent with.


\begin{remark}[{The opposite phase of $M$ and $M_{\RVM}^{\crit}$ between $10^3$ and $2\cdot10^3$}]\label{rem:anticorrelation}
In the left panel of Figure~\CGref{fig:mertens} the curves of $M$ and $M_{\RVM}^{\crit}$ are in opposite phase between $n\approx10^3$ and $n\approx2\cdot10^3$.  Heuristically, the visible anticorrelation is largely explained by the first mode, $2\sqrt x\,\Re\bigl(x^{i\gamma_1}/(\rho_1\zeta'(\rho_1))\bigr)$ and $2\sqrt x\,\Re\bigl(x^{i\tau_1}/(\rho_1^{\RVM}\zeta_{\RVM}'(\rho_1^{\RVM}))\bigr)$, two oscillations in $\log x$ with the frequencies $\gamma_1$ and $\tau_1$.  Their phase difference is $\eta_1\log x$ with $\eta_1=\tau_1-\gamma_1=14.5213-14.1347=0.3866$, up to the constant difference of the arguments of the two residues, which is not computed here; $\eta_1\log x$ equals $\pi$ at $x=e^{\pi/\eta_1}\approx3.4\cdot10^3$ and lies between $2.67$ and $2.94$ on $[10^3,2\cdot10^3]$.  The beat has period $2\pi/\eta_1\approx16.3$ in $\log x$.  This is the rotation \CGeqref{eq:Xi-rotation} seen on the Mertens side: the passage from the clock to $\zeta$ turns the mode $k$ by the angle $\eta_k\log x$.  The higher modes turn at different rates, and at large scale neither a correlation nor an anticorrelation remains, in agreement with the value $0.0007$ of Experiment~\CGref{exp:mertens}.  The effect depends on $\eta_1$ alone and changes sign with $x$; it does not distinguish $\lambda=0$.
\end{remark}

\subsection{The rung equations: transport, the defect of the clock, and the solution from the clock}
\label{subsec:num-rungs}

The four experiments of this subsection evaluate the objects of Subsections~\CGref{subsec:transport} and~\CGref{subsec:free-masses}.  In the first two, $\mathsf f$ is the function $\exp\bigl(-1/(1-y^2)\bigr)$, $y=(u-0.36)/0.26$, on $[0.10,0.62]\subset(0,\log2)$, normalized to $\int\mathsf f=1$; it tests the first rung.  In the programs $\mathsf f$, $\mathsf k_{\mathsf f}$ and $\mathsf c_{\mathsf f}$ are called \texttt{psi}, \texttt{W\_psi} and \texttt{c\_psi}.

\begin{experiment}[Transport on the first rung]\label{exp:transport}
\emph{Object:} $\mathsf R_{\mathsf f}(1)$, $\mathsf S_{\mathsf f}(\lambda)$ and the moments $\mathsf m_j(\mathsf f)$ of \CGeqref{eq:rung-functional}--\CGeqref{eq:transport-moments}.  \emph{Range:} the first $20\,000$ ordinates of each list; $\max_{k\ge1000}|\mathsf c_{\mathsf f}(\tau_k)|=2.5\cdot10^{-10}$ and $\max_{k\ge3000}|\mathsf c_{\mathsf f}(\tau_k)|=2.1\cdot10^{-13}$.  \emph{Expected:} if $\zeta_{Z1}=\zeta$, then $\mathsf S_{\mathsf f}(0)=\mathsf R_{\mathsf f}(1)$ (Proposition~\CGref{prop:rung-functional}); the series \CGeqref{eq:transport-moments} converges to $\mathsf S_{\mathsf f}(0)$ with terms bounded by $C(1+j)^2(C_\eta\log2)^j/j!$ (Theorem~\CGref{thm:transport-band}).  \emph{Observed:}
\begin{itemize}
\item $\mathsf k_{\mathsf f}=0.133643530916$, $\sum_k\mathsf c_{\mathsf f}(\tau_k)=0.176106813076$, $\sum_k\mathsf c_{\mathsf f}(\gamma_k)=0.133643530929$;
\item defect of the clock $\mathsf R_{\mathsf f}(1)=-0.042463282161$; transport $\mathsf S_{\mathsf f}(0)=-0.042463282147$; $\mathsf R_{\mathsf f}(0)=\mathsf R_{\mathsf f}(1)-\mathsf S_{\mathsf f}(0)=-1.4\cdot10^{-11}$;
\item moment series (Table~\CGref{tab:transport}): the term $j=1$ is $95.1\%$ of $\mathsf R_{\mathsf f}(1)$, the partial sum with $j\le2$ is $100.8\%$, and the partial sums with $j\le3$, $j\le5$, $j\le8$ differ from $\mathsf R_{\mathsf f}(1)$ by $3.7\cdot10^{-6}$, $6.5\cdot10^{-9}$, $1.0\cdot10^{-11}$;
\item along the family (Table~\CGref{tab:transport}) $\mathsf R_{\mathsf f}(\lambda)$ is monotone and almost linear on $[-1,2]$, with slope about $-0.044$ and a single zero, at $\lambda=0$;
\item the partial sums $\sum_{k\le K_0}\bigl(\mathsf c_{\mathsf f}(\gamma_k)-\mathsf c_{\mathsf f}(\tau_k)\bigr)$ are $-0.02086$, $-0.04381$, $-0.04649$, $-0.04221$, $-0.04249$, $-0.0424634$ for $K_0=1,5,10,50,100,400$: the first ordinate gives half of the transport, and the first rung is decided by about fifty ordinates.
\end{itemize}
\end{experiment}

\begin{table}[htbp]
\centering
\small
\begin{tabular}{rrrr@{\qquad\qquad}rr}
\toprule
$j$ & $\mathsf m_j(\mathsf f)$ & $(-1)^j\mathsf m_j(\mathsf f)/j!$ & partial sum $-\ \mathsf R_{\mathsf f}(1)$ & $\lambda$ & $\mathsf R_{\mathsf f}(\lambda)$\\
\midrule
1 & $+4.0395\cdot10^{-2}$ & $-4.0395\cdot10^{-2}$ & $+2.07\cdot10^{-3}$ & $-1.00$ & $+0.045303554710$\\
2 & $-4.8098\cdot10^{-3}$ & $-2.4049\cdot10^{-3}$ & $-3.37\cdot10^{-4}$ & $-0.50$ & $+0.022413444752$\\
3 & $-2.0432\cdot10^{-3}$ & $+3.4053\cdot10^{-4}$ & $+3.69\cdot10^{-6}$ & $-0.25$ & $+0.011133149816$\\
4 & $-5.5184\cdot10^{-5}$ & $-2.2993\cdot10^{-6}$ & $+1.39\cdot10^{-6}$ & $0.00$ & $-0.000000000014$\\
5 & $+1.6734\cdot10^{-4}$ & $-1.3945\cdot10^{-6}$ & $-6.47\cdot10^{-9}$ & $+0.25$ & $-0.010956709511$\\
6 & $+5.5721\cdot10^{-6}$ & $+7.7390\cdot10^{-9}$ & $+1.27\cdot10^{-9}$ & $+0.50$ & $-0.021706821780$\\
7 & $+7.2608\cdot10^{-6}$ & $-1.4406\cdot10^{-9}$ & $-1.69\cdot10^{-10}$ & $+1.00$ & $-0.042463282161$\\
8 & $+6.4039\cdot10^{-6}$ & $+1.5883\cdot10^{-10}$ & $-9.77\cdot10^{-12}$ & $+2.00$ & $-0.080112190380$\\
\bottomrule
\end{tabular}
\caption{Experiment~\CGref{exp:transport}.  Left: the moments $\mathsf m_j(\mathsf f)=\sum_k\eta_k^{\,j}\mathsf c^{(j)}_{\mathsf f}(\tau_k)$ and the partial sums of the series \CGeqref{eq:transport-moments} for $\mathsf S_{\mathsf f}(0)$, compared with the defect $\mathsf R_{\mathsf f}(1)=-0.042463282161$ of the clock.  Right: the rung functional along the family.}
\label{tab:transport}
\end{table}

\begin{experiment}[The defect of the clock by Poisson summation]\label{exp:clock-defect}
\emph{Object:} the two parts of $\mathsf R_{\mathsf f}(1)$ in Proposition~\CGref{prop:clock-defect-poisson}.  \emph{Expected:} $\mathsf R_{\mathsf f}(1)=(\mathsf k_{\mathsf f}-\mathsf I_0)-2\sum_{m\ge1}(-1)^m\mathsf I_m$.  \emph{Observed:} $\mathfrak t_0=9.676906787$; $\mathsf I_0=0.063533498702$ from \CGeqref{eq:clock-defect-I0}; the mismatch of the smooth densities is $\mathsf k_{\mathsf f}-\mathsf I_0=+0.070110032214$, and the oscillating part, obtained as the difference $\mathsf R_{\mathsf f}(1)-(\mathsf k_{\mathsf f}-\mathsf I_0)$ and not by summing the $\mathsf I_m$, is $-0.112573314375$.  The stationary heights $2\pi e^u$, $u$ in the support of $\mathsf f$, fill $[6.944,11.680]$; only those with $u>0.4319$ lie above $\mathfrak t_0$.
\end{experiment}

\begin{experiment}[From the clock to the ordinates and to the masses of the primes]\label{exp:clock-to-zeros}
\emph{Object:} a finite section of the system \CGeqref{eq:rung-system}: the unknowns are the ordinates $\kappa_k$ below a height and the masses $\mathsf a_n$, $2\le n\le11$; the equations are the pairings of \CGeqref{eq:rung-system} with $\omega_{\rm test}(u)\cos(\omega u)$ and $\omega_{\rm test}(u)\sin(\omega u)$, where $\omega_{\rm test}$ is a fixed smooth function supported in $(0,\log13.5)$ and $\omega$ runs over a grid of step $0.1$ in $[0,\Omega]$.  The ordinates above $\Omega+12$ are kept on the clock; only the first $KT=6000$ clock ordinates are summed numerically.  Although the support window $(0,\log13.5)$ contains $\log12$ and $\log13$, the implementation drops mass columns whose bump weight is below $0.05$, so those two masses are fixed to zero in this run.  These are explicit truncation sources in the residual.  \emph{Method:} Levenberg--Marquardt iteration started at $\kappa_k=\tau_k$; the masses enter linearly and are solved for at each step by ridge regression (the normal matrix $P_n^{\mathsf T}P_n$ is regularized by $10^{-10}I$ in the code, so they are not eliminated exactly), with continuation in height $\Omega=18,22,\dots,50$.  The proximity to $\gamma_k$ is never used inside the solver, only as a diagnostic afterwards.  The file \texttt{zeross.bin} is read only after the iteration, for the comparison.  \emph{Expected:} by Theorem~\CGref{thm:free-masses}, if $\zeta_{Z1}=\zeta$ the system has the unique solution $\kappa_k=\gamma_k$, $\mathsf a_n=\log p/(2\sqrt n)$ at prime powers and $0$ elsewhere.  \emph{Observed:} the iteration converges at every stage (residual norms between $4\cdot10^{-4}$ and $5\cdot10^{-3}$, of the size of the contribution of the ordinates kept on the clock); the solution is in Table~\CGref{tab:clock-to-zeros}: the first ten ordinates agree with $\gamma_1,\dots,\gamma_{10}$ within $4.4\cdot10^{-4}$, the eleventh within $2.2\cdot10^{-3}$, and the masses agree with \CGeqref{eq:free-masses} within $8\cdot10^{-4}$, including the values $0$ at $n=6,10$ and the prime-power values at $n=4,8,9$.  The ordinates $k=12,13,14$, which lie above the band $\omega\le50$, are not determined (errors $0.10$, $0.77$, $1.73$).  With masses up to $n\le34$ and the band $\omega\le60$ the section is not resolved and the iteration does not converge: resolving the ordinates up to a height $\Omega$ requires a strict margin $\log N>(1+\varkappa)\log(\Omega/2\pi)$ for some fixed $\varkappa>0$ in a stable finite solve and resolving the masses up to $N$ requires $\Omega$ of the order of $\pi N$, in agreement with the heights $2\pi x^{1/m}$ of Proposition~\CGref{prop:clock-defect-poisson}.  The experiment illustrates the existence form of the final step; it is a computation and not a proof of existence.
\end{experiment}

\begin{table}[htbp]
\centering
\small
\begin{tabular}{rrrr@{\qquad}rrr}
\toprule
$k$ & $\tau_k$ (start) & solution & $\gamma_k$ & $n$ & $\mathsf a_n$ (solution) & \CGeqref{eq:free-masses}\\
\midrule
1 & 14.521346953 & 14.134660056 & 14.134725142 & 2 & $+0.245076213$ & 0.245064536\\
2 & 20.655740356 & 21.022091631 & 21.022039639 & 3 & $+0.317145240$ & 0.317142050\\
3 & 25.492675432 & 25.010827129 & 25.010857580 & 4 & $+0.173268937$ & 0.173286795\\
4 & 29.739411632 & 30.424917668 & 30.424876126 & 5 & $+0.359884977$ & 0.359881258\\
5 & 33.624531889 & 32.935072953 & 32.935061588 & 6 & $-0.000004422$ & 0\\
6 & 37.257370087 & 37.586146548 & 37.586178159 & 7 & $+0.367785301$ & 0.367742452\\
7 & 40.700551604 & 40.918919600 & 40.918719012 & 8 & $+0.122615622$ & 0.122532268\\
8 & 43.994013953 & 43.327322117 & 43.327073281 & 9 & $+0.183143344$ & 0.183102048\\
9 & 47.165128427 & 48.004716753 & 48.005150881 & 10 & $-0.000286111$ & 0\\
10 & 50.233653391 & 49.773657726 & 49.773832478 & 11 & $+0.362276961$ & 0.361496314\\
11 & 53.214422874 & 52.972535790 & 52.970321478 & & &\\
\bottomrule
\end{tabular}
\caption{Experiment~\CGref{exp:clock-to-zeros}.  Solution of a finite section of the rung system \CGeqref{eq:rung-system}, started at the clock.  Left: the ordinates.  Right: the masses at the points $\log n$, which are not assumed, against $\log p/(2\sqrt n)$ at prime powers $n=p^m$.}
\label{tab:clock-to-zeros}
\end{table}


\begin{experiment}[Multiscale Jacobian of the clock]\label{exp:multiscale-jacobian}
\emph{Object:} the minimum midpoint-clock gap, its bound \CGeqref{eq:clock-gap-lower}, the Ingham critical window $2\pi/\Delta_T$, and the adjacent-pair bound \CGeqref{eq:fixed-window-pair}.  \emph{Input:} none of the genuine zeta ordinates; the program \texttt{s12\_multiscale\_jacobian.py} generates the $\tau_k$ directly from the Lambert formula up to $K=2\,001\,052$.  \emph{Observed:}
\begin{center}
\begin{tabular}{rrrr}
\toprule
$K$ & $T=\tau_K$ & $\Delta_T$ & $2\pi/\log(T/2\pi)$\\
\midrule
$10^2$ & $235.987267$ & $1.73462445$ & $1.73286190$\\
$10^3$ & $1419.517765$ & $1.15930485$ & $1.15921748$\\
$10^4$ & $9877.629616$ & $0.85368118$ & $0.85367617$\\
$10^5$ & $74920.891033$ & $0.66939912$ & $0.66939880$\\
$5\cdot10^5$ & $319387.322632$ & $0.57982852$ & $0.57982848$\\
$2\,001\,052$ & $1132490.722588$ & $0.51918344$ & $0.51918343$\\
\bottomrule
\end{tabular}
\end{center}
At the last height, $2\pi/\Delta_T=12.102053$ while $\log(T/2\pi)=12.10205288$.  Thus the rigorous scale $U\asymp\log T$ is numerically almost exact for the midpoint clock.  On every fixed window the adjacent-pair upper bound decreases with the clock gap, illustrating the obstruction in Theorem~\CGref{thm:clock-jacobian}(iii).  These calculations use no $\gamma_k$ and no $\eta_k$.
\end{experiment}

\begin{figure}[htbp]
\centering
\includegraphics[width=0.82\textwidth]{fig14_multiscale_gap.png}
\caption{Minimum gap of the midpoint clock up to height $T$ and the explicit lower bound $2\pi/\log(T/2\pi)$.  Their ratio is already $1+2\cdot10^{-8}$ at $K=2\,001\,052$.}
\label{fig:multiscale-gap}
\end{figure}

\begin{figure}[htbp]
\centering
\includegraphics[width=0.82\textwidth]{fig15_fixed_window_degeneracy.png}
\caption{Upper bound in \CGeqref{eq:fixed-window-pair} for adjacent clock frequencies on four fixed $u$-windows.  The decay is slow, of order $1/\log\tau_k$, but it is enough to exclude a height-uniform lower frame bound on any fixed window.}
\label{fig:fixed-window-degeneracy}
\end{figure}

\clearpage
\section{Conclusions and open problems}
\label{sec:conclusions}

\subsection*{What is new, what is proved, what is assumed, what is open}

\emph{New in this paper.}  The objects and results that, to the author's knowledge, do not appear elsewhere are: the midpoint clock $\tau_k$ as an exact Lambert inversion of the smooth Riemann--von Mangoldt law and the synthetic zeta function $\zeta_{\RVM}$ built on it (Sections~\CGref{sec:clock}--\CGref{sec:zeta-rvm}); the unconditional square-root estimate $\pi_{\RVM}(x)-\li(x)=O(\sqrt x(\log\log x)^2/\log x)$ for the synthetic prime count, with its continuous prime levels and cell masses (Sections~\CGref{sec:explicit-formula}--\CGref{sec:integer-weights}); the transition theorem $\zeta_{\RVM}=\varpi_{\RVM}(1+w)(1+O(1/|t|))$ with its constant $\varpi_{\RVM}$ and the synthetic Mertens bound $M^{\crit}_{\RVM}(x)=L+O_\eps(x^{1/2+\eps})$ (Section~\CGref{sec:transition}); the $\lambda$-family from the clock to the zeta ordinates, with the endpoint $\zeta_{Z1}$, the value $\varpi_{Z1}$ at infinity as an ordered product over the displacements, and the equivalences (D)$\Leftrightarrow\zeta_{Z1}=\zeta\Leftrightarrow\varpi_{Z1}=1\Leftrightarrow$RH (Sections~\CGref{sec:lambda-family}--\CGref{sec:euler-products}); the atomization of every member of the family by smooth cells, the mean-square bound $X^{1+\eps}$ for the cell masses uniformly in $\lambda$, and Hamburger's theorem for the pair (continuous, atomized) in its symmetric, theta and vertical mean-square forms (Sections~\CGref{sec:atomized}--\CGref{sec:criteria}); the compact space $\mathscr O$ of lists with the rigidity, existence-form and free-mass theorems, and the nondegeneracy theorem for the rung system (Section~\CGref{sec:lists}); the acceptance criterion and the two finite models in which the final step is a theorem (Section~\CGref{sec:models}); in Route B for finite sections, the identification of the Laplace transform of the residual as $\frac12(P'_{\boldsymbol\kappa}/P_{\boldsymbol\kappa}-\xi'/\xi)$ with the boundary term $\mathfrak H(0)\delta_0'$, the certificate theorem for staged sections, Gonek's lists in a tube around the clock, and the Kadec--Avdonin criterion with blocks of variable length (Section~\CGref{sec:routeB-finite}); and the a priori displacement proposition with its drift constant $\mathrm b$ (Appendix~\CGref{subsec:rung-apriori-v18}).  The equivalence of B$'$ with the Riemann hypothesis is classical in substance (Remark~\CGref{rem:rb-prescribed}), and no novelty is claimed for the variable-block variant beyond its formulation.

\emph{Proved without hypotheses.}  Every numbered statement (1)--(40) below is a theorem of this paper, proved with no hypothesis on $\zeta$, including the secondary lemmas and corollaries that carry them (the bounded displacement lemma, the weighted clock sum, the Lipschitz control of $J_{\RVM}$, the mean-square convergence of the ordinate series, the transfer theorem to $\mathscr O$, the compactness and continuity on $\mathscr O$, the point-support and finite-order steps of the nondegeneracy theorem, the Laplace-transform lemma of Route B, the monotone matching lemma); Section~\CGref{sec:routeB-finite} has its own status table (Section~\CGref{subsec:rb-status}).

\emph{Conditional on the one hypothesis.}  Hypothesis~\CGref{hyp:varpi-less-one}, $\varpi_{\RVM}<1$, is numerically supported and not proved (the paragraph after it says what a proof would need).  Exactly five statements use it and say so: Proposition~\CGref{prop:atomized-clock}(ii), the sign $\mathsf W(1)>0$ in Proposition~\CGref{prop:one-number}, Proposition~\CGref{prop:open-path-varpi}, Corollary~\CGref{cor:ladder-pair} at the clock end, and the remark after Corollary~\CGref{cor:positivity}.  Hypothesis~\CGref{hyp:cell-remainder} is proved (Corollary~\CGref{cor:clock-cells}) and is a hypothesis only in name.  The jump-growth condition in Theorem~\CGref{thm:cell-form}(ii) is part of a statement equivalent to (D), not an assumption of the paper.

\emph{Equivalent to (D) and not proved.}  The identity $\zeta_{Z1}=\zeta$, (D), RH, and each of the statements listed in Subsection~\CGref{subsec:status} and after item~(40) below: $\varpi_{Z1}=1$; the product identity \CGeqref{eq:product-identity}; the cell identity \CGeqref{eq:cell-identity}; $\mathsf W(0)\ge0$; the vanishing of the cell defects at $\lambda=0$; the symmetry of $e^{\tilde h_0}$; $T_0\equiv0$; $\mathsf V_\sigma(0)=0$; the existence of a coincidence list in $\mathscr O$; the solvability of the rung system in $\mathscr O$; statement B$'$.  These are reductions, each proved as an equivalence; none of them is resolved.

\emph{Open and not claimed.}  The existence of a solution of the rung system; the injectivity of the clock-only linearization and any continuation from the clock; the convergence of finite sections to a distributional limit; a rigorous enclosure of $\varpi_{\RVM}$; and any computer-assisted proof (the certificates of Section~\CGref{subsec:rb-numerics} are floating-point candidates).

\emph{Dependence table.}  The following table records what each group of results depends on.

\begingroup\small
\begin{longtable}{p{0.46\textwidth}p{0.46\textwidth}}
\toprule
Results & Depend on\\
\midrule
\endfirsthead
\toprule
Results & Depend on\\
\midrule
\endhead
\bottomrule
\endfoot
Clock, $\zeta_{\RVM}$, exponential sums, $J_{\RVM}$, $\pi_{\RVM}$, unit cell, transition theorem, synthetic Mertens bound (Sections~\CGref{sec:clock}--\CGref{sec:transition}) & the clock only; no property of $\zeta$ beyond the Riemann--von Mangoldt formula (for Lemma~\CGref{lem:bounded-displacement})\\
$\lambda$-family, $\zeta_{Z1}$, $\varpi_{Z1}$, Mellin counts, explicit formula, cell form (Sections~\CGref{sec:lambda-family}--\CGref{sec:euler-products}) & the ordinates $\gamma_k$ as a list in the tube, i.e.\ membership in $\mathscr O$; no hypothesis\\
Atomized family, mean square of cell masses, joining factor, pair criteria (Sections~\CGref{sec:atomized}--\CGref{sec:criteria}) & membership in $\mathscr O$; no hypothesis\\
Five statements listed above & Hypothesis~\CGref{hyp:varpi-less-one}\\
Step~3 and the transfer of the Mertens bound to $M$ (Theorems~\CGref{thm:step3}, \CGref{thm:transfer}) & (D), as an explicit hypothesis of the statement\\
Rigidity, existence form, free masses, nondegeneracy (Section~\CGref{sec:lists}) & membership in $\mathscr O$; nondegeneracy assumes the existence of a solution\\
Route B: compactness, prescribed-mass uniqueness, certificate, Gonek lists, $(A_2)$ criterion (Section~\CGref{sec:routeB-finite}) & $C$-admissibility (tube) only; B$'$ is equivalent to RH and not proved\\
A priori displacement (Appendix~\CGref{subsec:rung-apriori-v18}) & the existence of an exact solution with the local-density and mean-square hypotheses stated there\\
$|\log\varpi_{Z1}|<4\cdot10^{-13}$ (Proposition~\CGref{prop:varpi-verified}) & the verified zeros up to $T_{\rm ver}$ (input data)\\
Section~\CGref{sec:numerics}, Subsection~\CGref{subsec:rb-numerics} & floating-point computations with the supplied files; not part of any proof\\
\end{longtable}
\endgroup

\subsection*{The proved statements}

The following statements are proved (the results of Section~\CGref{sec:routeB-finite} are summarised in its status table, Section~\CGref{subsec:rb-status}):
\begin{enumerate}[label=\textup{(\arabic*)}]
\item the midpoint ordinates $\tau_k$ are given by the Lambert formula \CGeqref{eq:lambert} and stay within a bounded distance of the ordinates $\gamma_k$ of $\zeta$ (Lemma~\CGref{lem:bounded-displacement});
\item $P_{\RVM}^{\crit}$ is entire and symmetric, and $\zeta_{\RVM}$ has exactly the simple zeros $\tfrac12\pm i\tau_k$, the trivial zeros and a residue-one pole at $1$ (Theorem~\CGref{thm:canonical});
\item $S_{\RVM}(N;u)\ll1+\log N$ uniformly on compact positive frequency intervals, so the point centroid and $M_H^{\RVM}(x)$ tend to $0$ (Section~\CGref{sec:exp-sums});
\item $J_{\RVM}$ exists compact-uniformly and is related to $\pi_{\RVM}$ by exact finite M\"obius inversion (Section~\CGref{sec:explicit-formula});
\item $J_{\RVM}(x)-\li(x)$ and $\pi_{\RVM}(x)-\li(x)$ are $O(\sqrt x(\log\log x)^2/\log x)$, so $\li/\pi_{\RVM}\to1$, and with the prime number theorem $\pi/\li\to1$ and $\pi/\pi_{\RVM}\to1$ (Theorem~\CGref{thm:Lambert-estimate}, Corollary~\CGref{cor:three-limits});
\item $J_{\RVM}$ is Lipschitz on $[2,\infty)$, $c_n^{\RVM}=O(1)$, and $\zeta_{\Z,\RVM}$ is an absolutely convergent Dirichlet series in $\Re s>1$ (Section~\CGref{sec:integer-weights});
\item $\zeta_{\RVM}=\varpi_{\RVM}(1+w)(1+O(1/|t|))$ on every band to the right of $\tfrac12$ and $\zeta_{\RVM}=\varpi_{\RVM}(1+\chi)(1+O(1/|t|))$ on closed substrips of $\tfrac12<\Re s<1$ (Theorem~\CGref{thm:transition});
\item $1/\zeta_{\RVM}\to L$ on every vertical line to the right of $\tfrac12$, $\zeta_{\RVM}^{\pm1}$ have no absolutely convergent Dirichlet expansion, and $M_{\RVM}^{\crit}(x)=L+O_\eps(x^{1/2+\eps})$ (Section~\CGref{sec:transition});
\item $2\xi$ equals its Hadamard product over $\rho(1-\rho)$ exactly (Theorem~\CGref{thm:Hadamard}), and $|\zeta|\asymp|\mathcal Q_{\RVM}|$ on $\frac12+\eps\le\sigma\le c'$, $|t|\ge t_0$ (Corollary~\CGref{cor:Q-control});
\item the $\lambda$-family from $\zeta_{\RVM}$ to $\zeta_{Z1}$ converges jointly, has simple zeros for $\lambda>0$, an explicit generator and total multiplier, and its members differ only through $\gamma_k-\tau_k$ (Section~\CGref{sec:lambda-family});
\item $\zeta_{Z1}$ has the functional equation, the pole, the residue and the order of $\zeta$, and $\zeta_{Z1}=\zeta\Leftrightarrow$(D)$\Leftrightarrow$RH (Theorem~\CGref{thm:Hamburger-Z1});
\item $\zeta_{\RVM}$ is bounded above and below on every half-plane $\Re s\ge\sigma_0>1$ and tends to $\varpi_{\RVM}$ as $\Re s\to+\infty$ (Theorem~\CGref{thm:halfplane-RVM}); $1/\zeta_{\Z,\RVM}$ is an absolutely convergent Dirichlet series in $\Re s>1$, so (D) is equivalent to an absolutely convergent Dirichlet expansion of $\mathcal K$ in $\Re s>1$ (Theorem~\CGref{thm:D-multiplier});
\item no member of the $\lambda$-family, $\zeta_{Z1}$ included, can have a zero off the critical line other than the trivial zeros (Theorem~\CGref{thm:path});
\item the endpoint has the value $\varpi_{Z1}=\varpi_{\RVM}\lim_{K\to\infty}\prod_{k\le K}(\frac14+\tau_k^2)/(\frac14+\gamma_k^2)\ge1$ at $+\infty$, and $\zeta_{Z1}=\zeta$ if and only if $\varpi_{Z1}=1$ (Theorems~\CGref{thm:Z1-clock}, \CGref{thm:varpi-Z1}, Corollary~\CGref{cor:product-identity});
\item the prime-power counts of $\zeta_{\RVM}$ and $\zeta_{Z1}$ are the Mellin transforms of their logarithms and satisfy a distributional explicit formula; this proves the Mellin--Stieltjes identification and the weighted factorization $\zeta_{\Z,\RVM}=\zeta_{\RVM}e^{h_{\RVM}}$, and, with the jump-growth hypothesis of Theorem~\CGref{thm:cell-form}(ii), (D) holds if and only if the prime-power count of $\zeta_{Z1}$ is constant between consecutive integers (Theorems~\CGref{thm:mellin-counts}, \CGref{thm:distributional-EF}, \CGref{thm:MS}, \CGref{thm:cell-form});
\item along the family $\frac d{d\lambda}\log\varpi_\lambda=-\Upsilon(\lambda)$, and (D) is the identity $\int_0^1\Upsilon\,d\lambda=\log L$; with the verified zeros, $0\le\log\varpi_{Z1}<4\cdot10^{-13}$ (Propositions~\CGref{prop:family-flow}, \CGref{prop:varpi-verified});
\item the atomized family $\tilde\zeta_{\Z,\lambda}=\zeta_\lambda e^{\tilde h_\lambda}$ consists of absolutely convergent Dirichlet series in $\Re s>1$ for every $\lambda$, the endpoint included, because the cell masses have mean square $X^{1+\eps}$ on dyadic blocks uniformly in $\lambda$; the joining factor is explicit cell by cell and holomorphic in $\Re s>\frac12$ (in $\Re s>0$ for $\lambda=1$), and (D) is equivalent to the vanishing of every cell defect at $\lambda=0$ and of $J_{Z1}$ on the first cell $(1,2)$ (Proposition~\CGref{prop:first-cell}, Theorems~\CGref{thm:joining-factor}, \CGref{thm:mean-square-cells}, \CGref{thm:atomized-continuation}, \CGref{thm:closing});
\item at the clock endpoint the joining factor is not constant, $\tilde h_1(\sigma)\to\log L$; the atomized reciprocals have M\"obius coefficients $\tilde\mu_\lambda$, and the endpoint Mertens function has a Mellin form $B_{Z1}$ with $B_{Z1}=M$ if and only if (D) (Propositions~\CGref{prop:atomized-clock}, \CGref{prop:Mellin-Mertens});
\item Step~3: $\zeta_{Z1}=\zeta$ implies $M(x)=O_\eps(x^{1/2+\eps})$ (Theorem~\CGref{thm:step3}); equivalently, the synthetic bound $M_{\RVM}^{\crit}(x)=L+O_\eps(x^{1/2+\eps})$ transfers to $M$ through $1/\zeta=F_{\RVM}\Pi_{\rm tr}^{-1}$ (Theorem~\CGref{thm:transfer});
\item for $0<\lambda\le1$ no $\zeta_\lambda$ is an absolutely convergent Dirichlet series in $\Re s>1$ and no joining factor $\tilde h_\lambda$ is constant, and $\lambda\mapsto\log\zeta_\lambda(\sigma)$ is strictly convex for every real $\sigma>1$ (Propositions~\CGref{prop:nonconstant}, \CGref{prop:convexity});
\item if the joining factor $e^{\tilde h_\lambda}$ is meromorphic of finite order with finitely many poles and symmetric under $s\mapsto1-s$, Hamburger's theorem gives $\tilde\zeta_{\Z,\lambda}=\zeta$, and for $\lambda=0$ also $\zeta_{Z1}=\zeta$ (Theorem~\CGref{thm:pair-criterion});
\item if $\zeta_\lambda$ is the Mellin--Stieltjes transform of a positive measure, then $\varpi_\lambda\le1$; (D) holds if and only if $\zeta_{Z1}$ has such a representation, if and only if $J^{\rm M}_{Z1}$ is non-decreasing (Theorem~\CGref{thm:positivity}, Corollary~\CGref{cor:positivity});
\item $J_{\RVM}$ is continuously differentiable on $(2,\infty)$ with $J_{\RVM}'(y)=\frac1{\log y}+2\cos(2\pi y)-2^{-1/2}y^{-1/4}\cos(4\pi\sqrt y+\frac\pi4)+O(y^{-1/3})$; consequently $\tilde c_n(1)=\frac1{\log n}+O(n^{-1/4})$ off the powers of $2$, and Hypothesis~\CGref{hyp:cell-remainder} holds off the integers $2^j$ with $\mu(j)\ne0$ (Theorem~\CGref{thm:clock-density}, Corollary~\CGref{cor:clock-cells});
\item the cell masses are sums over the ordinates with the height kernel $\hat{\varphi}(\gamma/2\pi n)$, which vanishes at the heights $2\pi qn$; their flow is the pairing of $V_n$ with the count difference $N_\zeta-n_{\RVM}$; and the Gaussian defect of the atomization is $\frac12(e^{\tilde h_\lambda(1)}-1)y^{-1/2}+O(y^{-1/4-\delta})$ (Propositions~\CGref{prop:height-kernels}, \CGref{prop:flow-count}, \CGref{prop:gaussian-defect});
\item $\zeta_{Z1}(\sigma)\ge\zeta(\sigma)$ for every real $\sigma>1$, with equality at one point if and only if $\zeta_{Z1}=\zeta$; in particular $\prod_k\bigl(1+2/(\frac14+\gamma_k^2)\bigr)\ge\pi/3$ (Theorem~\CGref{thm:real-axis});
\item every $\zeta_\lambda$ is the Mellin transform of a function $\mathsf A_\lambda$ that satisfies the summation formula of $\zeta$, and $\varpi_\lambda-1=-\mathsf W(\lambda)/\hat{\mathsf g}(1)$ for one explicit number $\mathsf W(\lambda)$; (D) is equivalent to $\mathsf W(0)\ge0$ and, under Hypothesis~\CGref{hyp:varpi-less-one}, to $\varpi_\lambda<1$ for all $\lambda\in(0,1]$ (Lemma~\CGref{lem:summation-all}, Propositions~\CGref{prop:one-number}, \CGref{prop:open-path-varpi});
\item Hamburger's theorem for the pair holds without continuation of the joining factor: $\tilde\zeta_{\Z,\lambda}=\zeta$ if and only if the atomized coefficients satisfy the theta relation, if and only if the vertical mean square of $\tilde\zeta_{\Z,\lambda}/\zeta$ on a line $\Re s=\sigma>1$ is $1$; both defects are continuous in $\lambda$ (Theorem~\CGref{thm:double-hamburger});
\item at the clock end the joining factor has no symmetric meromorphic continuation, and $\tilde\zeta_{\Z,\lambda}=\zeta$ for at most finitely many $\lambda\in[0,1]$, not for $\lambda=1$ (Propositions~\CGref{prop:clock-asymmetry}, \CGref{prop:exceptional-set});
\item the cell masses are entire functions of $\lambda$ of exponential type $O(\log n)$, and the cell defects have the height kernel $\hat{\varphi}'(\gamma/2\pi n)$ (Propositions~\CGref{prop:entire-lambda}, \CGref{prop:defect-kernel});
\item in the compact space $\mathscr O$ of ordinate lists at bounded distance from the clock, at most one list has coinciding continuous and atomized zeta functions; such a list exists if and only if $\zeta_{Z1}=\zeta$, and it is then the list of the ordinates of $\zeta$ (Theorem~\CGref{thm:rigidity}, Proposition~\CGref{prop:lists-compact}, Corollary~\CGref{cor:existence-form});
\item (D) holds if and only if $\tilde c_n(0)=\mathbf 1_{\rm prime}(n)$ for every $n\ge2$, one real number for each cell; every $\lambda$ with $\tilde\zeta_{\Z,\lambda}=\zeta$ is a witness of (D), so the exceptional set is empty or contains $0$; cell by cell the remaining identity is \CGeqref{eq:per-cell} (Corollary~\CGref{cor:cell-mass-form}, Propositions~\CGref{prop:witness}, \CGref{prop:per-cell});
\item for a pair joined by a symmetric factor Hamburger's theorem gives a constant factor when the continuous function has order one, and not in general; at $\lambda=0$ the symmetry of the joining factor is equivalent to (D) (Proposition~\CGref{prop:dichotomy}, Remarks~\CGref{rem:pair-no-slack}, \CGref{rem:dichotomy-counterexample});
\item the ladder: for every $N\ge2$, $\tilde h_\lambda(\sigma)=O(\sigma N^{-\sigma})$ as $\sigma\to+\infty$ if and only if the Mellin count of $\zeta_\lambda$ has no mass outside the integers below $N$; (D) holds if and only if this holds at $\lambda=0$ for every $N$, and at the clock end the first rung fails (Theorems~\CGref{thm:ladder}, \CGref{thm:halfline-hamburger}, Corollary~\CGref{cor:ladder-pair});
\item the bridge theorem: if the imaginary part of the joining factor of a pair extends continuously to $\Re s\ge\frac12$, vanishes on the critical line, and the factor satisfies the temperate bound \CGeqref{eq:bridge-B2}, then the factor is constant, the discrete function is $\zeta$ and the continuous one is a multiple of $\zeta$; for the family the two conditions hold for no $\lambda\in(0,1]$, and at $\lambda=0$ they are equivalent to (D) (Theorem~\CGref{thm:bridge}, Corollary~\CGref{cor:bridge-family});
\item on the critical line $\Im\log\zeta_\lambda=-\vartheta(t)+\pi(N_\lambda(t)-1)$, and the cell masses of two lists differ by a sum over the ordinates of one function of one ordinate (Propositions~\CGref{prop:continuous-phase}, \CGref{prop:one-ordinate});
\item on each rung the passage from the clock to the ordinates of $\zeta$ is the transport $\mathsf S_{\mathsf f}(\lambda)$, an entire function of $\lambda$ of exponential type at most $C_\eta\log N$ whose value at $\lambda=0$ is the sum of the series of the moments $\sum_k\eta_k^{\,j}\mathsf c_{\mathsf f}^{(j)}(\tau_k)$; (D) holds if and only if the transport equals the defect of the clock on every rung, and this defect is given exactly by Poisson summation, with stationary heights $2\pi x^{1/m}$ (Proposition~\CGref{prop:rung-functional}, Theorem~\CGref{thm:transport-band}, Proposition~\CGref{prop:clock-defect-poisson});
\item (D) holds if and only if the rung system \CGeqref{eq:rung-system}, with free masses at the points $\log n$, has a solution in $\mathscr O$; the solution is then the list of $\zeta$ with the masses of von Mangoldt, and its linearization is injective up to the splitting of multiple ordinates (Theorems~\CGref{thm:free-masses}, \CGref{thm:nondegenerate});
\item at the midpoint clock, the finite ordinate Jacobian up to height $T$ has a uniform lower frame bound on windows of length $(1+\varkappa_{\rm win})\log(T/2\pi)$, no uniform lower bound is possible on a fixed window, the infinite ordinate block has dense range on every finite window, and smooth bounded-displacement transport tails are $O_{A,\mathsf f}(T^{1-A}\log T)$ (Theorem~\CGref{thm:clock-jacobian});
\item the gradient of the vertical defect on $\mathscr O$ is $\partial\mathsf V_\sigma/\partial\kappa_k=2\mathsf Q_{\boldsymbol{\kappa}}(\kappa_k)$ with one real-analytic function $\mathsf Q_{\boldsymbol{\kappa}}$ (Proposition~\CGref{prop:critical-points});
\item a list is accepted by a continuous side and a discrete side if and only if its function lies in the intersection of the window of the first with the frequencies of the second; a non-periodic list is accepted through one unitary matrix; for the finite Euler product $(1-2^{-s})(1-3^{-s})$ the constructions of the paper apply and the endpoint $\lambda=0$ is the only member of the family that is a Dirichlet series; for $\prod_{p\le X}(1-p^{-s})$ with more primes the Poisson part extends and the uniqueness part does not (Theorem~\CGref{thm:acceptance}, Proposition~\CGref{prop:graph-model}, Theorem~\CGref{thm:finite-euler} and the remark after it);
\item there is no function with the functional equation of $\zeta$ and an absolutely convergent Dirichlet series supported on the integers with bounded prime factors, and for the diagonal operator over the primes the contraction, the convergence of the determinant and the symmetry hold in $\sigma>0$, in $\sigma>1$ and on $\sigma=\frac12$ respectively (Proposition~\CGref{prop:no-finite-rank}, Subsection~\CGref{subsec:three-abscissas}).
\end{enumerate}
\emph{Not proved.}  The identity $\zeta_{Z1}=\zeta$ is not proved; neither is the Riemann hypothesis, to which it is equivalent.  The statements proved to be equivalent to it are listed in Subsection~\CGref{subsec:status} and above; in the form of Corollary~\CGref{cor:existence-form} it is the existence of one list in $\mathscr O$ whose continuous and atomized zeta functions coincide, and in the form of Corollary~\CGref{cor:bridge-family} it is the vanishing on the critical line of the imaginary part of the joining factor $\tilde h_0$, with the bound \CGeqref{eq:bridge-B2}.  The theorems of Subsections~\CGref{subsec:cell-mass-form}--\CGref{subsec:transport} and of Sections~\CGref{sec:lists} and~\CGref{sec:models} are formulations of this step and tools for it; none of them evaluates the pairing \CGeqref{eq:per-cell} of the count difference $N_\zeta-n_{\RVM}$, on which every one of them depends at $\lambda=0$.

The computations of Section~\CGref{sec:numerics} with $2\,001\,052$ zeros illustrate only the specific numerical statements explicitly tested there; they do not provide evidence for the unproved equivalent forms of (D).

\begin{remark}[Open problems]\label{rem:open-problems}
\begin{enumerate}[label=\textup{(\roman*)}]
\item Prove the identity $\zeta_{Z1}=\zeta$, equivalently condition (D), equivalently an absolutely convergent Dirichlet expansion of $\mathcal K=\zeta_{Z1}/\zeta_{\Z,\RVM}$ in $\Re s>1$, equivalently the product identity \CGeqref{eq:product-identity}, equivalently the cell identity \CGeqref{eq:cell-identity} with the jump bound of Theorem~\CGref{thm:cell-form}, equivalently step A7 of the programme in Subsection~\CGref{subsec:programme}, equivalently the positivity of the measure representing $\zeta_{Z1}$ (Corollary~\CGref{cor:positivity}), equivalently the existence of a coincidence list in $\mathscr O$ (Corollary~\CGref{cor:existence-form}, step A19).  This is the final step of the method.
\item Hypothesis~\CGref{hyp:cell-remainder} is proved for the integers $n$ that are not of the form $2^j$ with $\mu(j)\ne0$ (Corollary~\CGref{cor:clock-cells}); at $n=2^j$ the cell mass is $(1+\mu(j)J_{\RVM}(2)\log2)/\log n+O(n^{-1/4})$.  Determine the sign of $1\pm J_{\RVM}(2)\log2$.
\item Decide, for $0<\lambda<1$ outside the finite set $\mathsf E$, which hypothesis of Theorem~\CGref{thm:pair-criterion} fails, and which of \textup{(B1)}, \textup{(B2)} fails in Theorem~\CGref{thm:bridge}; decide whether the vanishing of all admissible coordinate derivatives (with the corresponding one-sided/KKT conditions at active constraints) forces $\mathsf V_\sigma=0$ (Proposition~\CGref{prop:critical-points}); and decide whether the atomized coefficients $\tilde a_N(0)$ control the sign of $\mathsf W(0)$ in Proposition~\CGref{prop:one-number}.
\item Determine the monotone branches of $\pi_{\RVM}$, on which the real prime levels $p_j^{\RVM}$ are defined, and explain the observed proximity $p_j^{\RVM}\approx p_j$ for small $j$ (Experiment~\CGref{exp:primes}).
\item Derive the bound \textup{(B2)} for $\tilde h_\lambda$ from the bounds of the cell defects, and turn Heuristic~\CGref{heur:phase-rotation} into a lemma with remainder; decide \textup{(B1)} at $\lambda=0$.
\item Evaluate from the construction the right side of \CGeqref{eq:per-cell}, or the moments $\mathsf m_j(\mathsf f)$ of \CGeqref{eq:transport-moments}; on the first rung the moments $j=1,2$ carry the transport up to $4\cdot10^{-6}$ (Experiment~\CGref{exp:transport}).
\item Theorem~\CGref{thm:clock-jacobian} gives the finite-height lower frame bound on the natural scale $U\asymp\log T$, the failure of fixed-window uniform invertibility, dense range on every fixed window, and rapidly summable smooth tails.  The exact-atomic interpolation tested later has no established interior path.  A remaining numerical existence programme is instead to control bounded endpoint finite sections on an exhausting nested family of tests and pass to a distributional limit; proving the required residual-to-zero and compactness statement would already be an existence theorem for the final rung system.
\item Under (D), decide whether $\mathsf E=\{0\}$, that is, whether the functions $\mathsf D_n$ of \CGeqref{eq:witness-zeros} have a common zero in $(0,1]$.
\item Construct an object of infinite rank from which $\zeta_{Z1}$ and $\tilde\zeta_{\Z,0}$ both descend, as the matrix $\mathsf U$ of Proposition~\CGref{prop:graph-model} does for the two-edge model (Subsection~\CGref{subsec:three-abscissas}).
\end{enumerate}
\end{remark}

\begin{remark}[Statements derived in part or conjectural, not proved here]\label{rem:not-proved-list}
The following statements arose in the work on the final step; none of them is used in this paper.
\begin{enumerate}[label=\textup{(\alph*)}]
\item \emph{Rates in $\lambda$.}  Off the powers of $2$ the formulas of Proposition~\CGref{prop:entire-lambda} suggest
\begin{equation}\label{eq:rates-lambda}
 \tilde h_\lambda(s)-\tilde h_0(s)=O_s(\lambda),\qquad\bigl|\mathfrak h_n(\lambda;s)-\mathfrak h_n(0;s)\bigr|\ll_{s,\eps}\lambda\,n^{-\sigma-\frac12+\eps};
\end{equation}
the cells $n=2^j$ and the uniformity in $n$ of the first-cell term are not written.
\item \emph{Flow form of the density condition.}  $\tilde h_0(1)=0$ if and only if $\int_0^1\partial_\lambda\tilde h_\lambda(1)\,d\lambda=\tilde h_1(1)$; the differentiability at $s=1$ uniformly in $\lambda$ is not written.
\item \emph{Pointwise bound.}  $|\tilde c_n(\lambda)|\le C_\eps n^\eps$ uniformly for $\lambda$ in a neighbourhood of $0$ is open; by Corollary~\CGref{cor:cell-mass-form} it is not needed for the final step.
\item \emph{Leading term of the clock defect.}  The form $\mathfrak h_n(1;s)\approx\sqrt2\,\pi\,\mathsf M_2\,s\,n^{-s-7/4}\sin(4\pi\sqrt n+\frac\pi4)$ is not proved; the bound \CGeqref{eq:clock-defect} is.
\item \emph{A heuristic model of the atomized clock.}  By \CGeqref{eq:clock-cell-smooth}, $\tilde\zeta_{\Z,1}(s)$ is close to $\prod_{n\ge2}(1-n^{-s})^{-1/\log n}$, whose logarithm has the leading part $\sum_n n^{-s}/\log n=\int_s^\infty(\zeta(y)-1)\,dy$; no error term is claimed.
\item The comparison between the two columns of Subsection~\CGref{subsec:hamburger-table} and the use of the classification of \cite{LevOlevskii2015} remain as discussed in the open-problem paragraph following Proposition~\CGref{prop:critical-points} and are not proved.
\end{enumerate}
\end{remark}

\appendix
\section{Python programs}
\label{app:code}

The programs below produce every number, table and figure of Section~\CGref{sec:numerics}.  They read the two data files \texttt{data/zeross.bin} (ordinates $\gamma_k$ of the zeros of $\zeta$) and \texttt{data/zLamb.bin} (midpoint ordinates $\tau_k$), both little-endian float64 with $2\,001\,052$ entries, and write their console output, CSV, MAT and binary files to \texttt{output/} and the figures to \texttt{figures/}.  They require Python~3 with NumPy, SciPy, mpmath and Matplotlib, and are run from the directory \texttt{code/} in the order \texttt{s01}--\texttt{s10}; the programs \texttt{s08}--\texttt{s10} use NumPy only and read the same two data files.

\subsection{\texttt{rvm\_common.py}}
\label{app:code-rvm-common}

\begin{lstlisting}[language=Python]
"""Common routines for the midpoint Riemann-von Mangoldt (RVM) zeta system.

Data files (little-endian float64, one ordinate per entry):
    data/zeross.bin  ordinates gamma_k of the nontrivial zeros of zeta (2,001,052 values)
    data/zLamb.bin   midpoint RVM ordinates tau_k, N_RVM(tau_k) = k - 1/2

Notation follows the paper: N_RVM is the clock, tau_k the midpoint ordinates,
gamma_k the ordinates of zeta, eta_k = tau_k - gamma_k, q(s) = s(s-1),
G(s) = s(s-1) Gamma(s/2) pi^(-s/2), zeta_lambda = P_lambda / G with ordinates
gamma_k(lambda) = (1-lambda) gamma_k + lambda tau_k.
"""
import os
import sys
import time
import numpy as np
from scipy import special

ROOT = os.path.dirname(os.path.dirname(os.path.abspath(__file__)))
DATA = os.path.join(ROOT, "data")
OUT = os.path.join(ROOT, "output")
FIG = os.path.join(ROOT, "figures")
os.makedirs(OUT, exist_ok=True)
os.makedirs(FIG, exist_ok=True)

TWO_PI = 2.0 * np.pi


class Tee:
    """Write console output both to the screen and to output/<name>_console.txt."""

    def __init__(self, name):
        self.f = open(os.path.join(OUT, name + "_console.txt"), "w")
        self.t0 = time.time()

    def __call__(self, fmt, *args):
        s = fmt % args if args else fmt
        sys.stdout.write(s + "\n")
        sys.stdout.flush()
        self.f.write(s + "\n")
        self.f.flush()

    def close(self):
        self("total run time %.1f s" % (time.time() - self.t0))
        self.f.close()


def read_bin(name):
    return np.fromfile(os.path.join(DATA, name), dtype="<f8")


def load_ordinates():
    gam = read_bin("zeross.bin")
    tau = read_bin("zLamb.bin")
    n = min(gam.size, tau.size)
    return gam[:n], tau[:n]


def clock(T):
    """N_RVM(T) = T/(2 pi) log(T/(2 pi e)) + 7/8 (real or complex T, principal log)."""
    return T / TWO_PI * (np.log(T / TWO_PI) - 1.0) + 0.875


def tail_coefficients(U):
    """Midpoint-rule tail of sum_{k>N} log(1 + q/tau_k^2), N = N_RVM(U):
    returns (c1, c2, c3) with tail(q) = c1 q - c2 q^2 + c3 q^3."""
    LU = np.log(U / TWO_PI)
    c1 = (LU + 1.0) / U / TWO_PI
    c2 = (3.0 * LU + 1.0) / (9.0 * U ** 3) / 2.0 / TWO_PI
    c3 = (5.0 * LU + 1.0) / (25.0 * U ** 5) / 3.0 / TWO_PI
    return c1, c2, c3


def clock_inverse(y):
    """Solve N_RVM(U) = y by Newton iteration."""
    U = TWO_PI * max(y, 2.0) / np.log(max(y, 3.0))
    for _ in range(80):
        U -= (clock(U) - y) / (np.log(U / TWO_PI) / TWO_PI)
    return U


class Product:
    """log P(s) = sum_k log(1 + q(s)/(1/4 + o_k^2)) for an ordinate list o_k.

    Exact summation for k <= K1, power sums for K1 < k <= N, and the clock tail
    for k > N.  Accurate for |q(s)| << o_{K1}^2."""

    def __init__(self, ords, K1=100000, nterms=6):
        self.o = np.asarray(ords, dtype=float)
        self.N = self.o.size
        self.K1 = min(K1, self.N)
        A = 0.25 + self.o[self.K1:] ** 2
        self.pw = [np.sum(A ** (-j)) for j in range(1, nterms + 1)]
        self.A1 = 0.25 + self.o[: self.K1] ** 2
        self.tail = tail_coefficients(clock_inverse(float(self.N)))

    def logP(self, s):
        s = np.atleast_1d(np.asarray(s, dtype=complex))
        q = s * (s - 1.0)
        out = np.empty(s.shape, dtype=complex)
        for i, qi in enumerate(q):
            v = np.sum(np.log1p(qi / self.A1))
            for j, p in enumerate(self.pw, start=1):
                v += (-1) ** (j + 1) * qi ** j / j * p
            c1, c2, c3 = self.tail
            v += c1 * qi - c2 * qi ** 2 + c3 * qi ** 3
            out[i] = v
        return out


def logG(s):
    s = np.asarray(s, dtype=complex)
    return np.log(s) + np.log(s - 1.0) + special.loggamma(s / 2.0) - s / 2.0 * np.log(np.pi)


def zeta_from_product(prod, s):
    s = np.atleast_1d(np.asarray(s, dtype=complex))
    return np.exp(prod.logP(s) - logG(s))


def chi(s):
    s = np.asarray(s, dtype=complex)
    return np.exp((s - 0.5) * np.log(np.pi) + special.loggamma((1 - s) / 2) - special.loggamma(s / 2))


def w_transition(s):
    """w(s) = exp(-2 pi i N_RVM(tau_*)), tau_* = -i (s - 1/2)."""
    s = np.asarray(s, dtype=complex)
    tstar = -1j * (s - 0.5)
    return np.exp(-2j * np.pi * clock(tstar))


def theta_rs(t):
    t = np.asarray(t, dtype=float)
    return t / 2 * np.log(t / TWO_PI) - t / 2 - np.pi / 8 + 1 / (48 * t) + 7 / (5760 * t ** 3)


def gamma_lambda(gam, tau, lam):
    return (1.0 - lam) * gam + lam * tau


def ei_upper(w):
    """Ei on Im w != 0 along the horizontal line (paper, eq. Ei-def): Ei(w) = -E1(-w)."""
    return -special.exp1(-np.asarray(w, dtype=complex))


def mobius_sieve(X):
    mu = np.ones(X + 1, dtype=np.int8)
    mu[0] = 0
    is_comp = np.zeros(X + 1, dtype=bool)
    for p in range(2, X + 1):
        if not is_comp[p]:
            is_comp[2 * p::p] = True
            mu[p::p] *= -1
            if p * p <= X:
                mu[p * p::p * p] = 0
    return mu


def primes_upto(X):
    s = np.ones(X + 1, dtype=bool)
    s[:2] = False
    for p in range(2, int(X ** 0.5) + 1):
        if s[p]:
            s[p * p::p] = False
    return np.nonzero(s)[0]


def write_csv(name, header, cols):
    A = np.column_stack(cols)
    np.savetxt(os.path.join(OUT, name), A, delimiter=",", header=",".join(header), comments="", fmt="%.12g")


def savefig(fig, name):
    fig.savefig(os.path.join(FIG, name), dpi=150, bbox_inches="tight")
\end{lstlisting}

\subsection{\texttt{s01\_clock\_and\_family.py}}
\label{app:code-s01-clock-and-family}

\begin{lstlisting}[language=Python]
"""Experiment 1: the midpoint clock, the displacements eta_k = tau_k - gamma_k,
the lambda-family zeta_lambda, and the numerical comparison zeta_Z1 versus zeta."""
import numpy as np
import mpmath as mp
import scipy.io as sio
import matplotlib
matplotlib.use("Agg")
import matplotlib.pyplot as plt
from rvm_common import (Tee, load_ordinates, clock, Product, zeta_from_product, gamma_lambda,
                        write_csv, savefig, OUT)
import os

say = Tee("s01_clock_and_family")
gam, tau = load_ordinates()
N = gam.size
k = np.arange(1, N + 1)
say("zeros loaded: N = %d, gamma_1 = %.9f, gamma_N = %.6f, tau_1 = %.9f, tau_N = %.6f",
    N, gam[0], gam[-1], tau[0], tau[-1])

# ---- clock check and displacements
clock_err = np.max(np.abs(clock(tau) - (k - 0.5)))
eta = tau - gam
say("max |N_RVM(tau_k) - (k - 1/2)| = %.3e", clock_err)
say("eta_k = tau_k - gamma_k:  max |eta| = %.6f (k = %d),  mean |eta| = %.6f,  rms = %.6f,  mean eta = %.3e",
    np.max(np.abs(eta)), int(np.argmax(np.abs(eta))) + 1, np.mean(np.abs(eta)),
    np.sqrt(np.mean(eta ** 2)), np.mean(eta))
blocks = [(1, 10), (11, 100), (101, 1000), (1001, 10 ** 4), (10 ** 4 + 1, 10 ** 5), (10 ** 5 + 1, 10 ** 6), (10 ** 6 + 1, N)]
rows = []
say("%12s %12s %12s %12s", "k from", "k to", "max|eta|", "rms eta")
for a, b in blocks:
    e = eta[a - 1:b]
    rows.append([a, b, np.max(np.abs(e)), np.sqrt(np.mean(e ** 2))])
    say("%12d %12d %12.6f %12.6f", a, b, rows[-1][2], rows[-1][3])
rows = np.array(rows)
write_csv("s01_eta_blocks.csv", ["k_from", "k_to", "max_abs_eta", "rms_eta"], [rows[:, i] for i in range(4)])
say("first ordinates:")
say("%4s %16s %16s %12s", "k", "gamma_k", "tau_k", "eta_k")
for i in range(12):
    say("%4d %16.9f %16.9f %12.6f", i + 1, gam[i], tau[i], eta[i])

# ---- zeta_Z1 versus zeta on vertical lines
PZ = Product(gam)
PR = Product(tau)
tline = np.concatenate([np.linspace(1.0, 100.0, 199), np.linspace(101.0, 2000.0, 300)])
mp.mp.dps = 20
cmp_rows = []
for sigma in (2.0, 1.25, 0.75):
    s = sigma + 1j * tline
    zz1 = zeta_from_product(PZ, s)
    ze = np.array([complex(mp.zeta(complex(x))) for x in s])
    rel = np.abs(zz1 / ze - 1.0)
    cmp_rows.append([sigma, np.max(rel), np.median(rel)])
    say("sigma = %.2f, t in [1,2000] (%d points): max |zeta_Z1/zeta - 1| = %.3e, median = %.3e",
        sigma, tline.size, np.max(rel), np.median(rel))
    write_csv("s01_zetaZ1_vs_zeta_sigma%.2f.csv" % sigma, ["t", "abs_zeta", "abs_zetaZ1", "rel_err"],
              [tline, np.abs(ze), np.abs(zz1), rel])
cmp_rows = np.array(cmp_rows)

# ---- Hardy-type functions Z_lambda(t) = exp(i theta(t)) zeta_lambda(1/2+it) on the critical line
def theta_exact(t):
    return np.array([float(mp.siegeltheta(x)) for x in t])

tc = np.linspace(0.5, 60.0, 1200)
th = theta_exact(tc)
lams = [0.0, 0.25, 0.5, 0.75, 1.0]
Zl = {}
for lam in lams:
    P = PZ if lam == 0.0 else (PR if lam == 1.0 else Product(gamma_lambda(gam, tau, lam)))
    Zl[lam] = np.real(np.exp(1j * th) * zeta_from_product(P, 0.5 + 1j * tc))
Zhardy = np.array([float(mp.siegelz(x)) for x in tc])
dZ = np.max(np.abs(Zl[0.0] - Zhardy) / np.maximum(1.0, np.abs(Zhardy)))
say("critical line, t in [0.5,60]: max |Z_Z1(t) - Z(t)| / max(1,|Z(t)|) = %.3e", dZ)
write_csv("s01_Z_lambda_critical_line.csv", ["t", "Z_hardy"] + ["Z_lambda_%.2f" % l for l in lams],
          [tc, Zhardy] + [Zl[l] for l in lams])

# ---- figures
dec = np.arange(0, N, 50)
fig, ax = plt.subplots(2, 2, figsize=(12, 8))
ax[0, 0].plot(k[dec], eta[dec], ".", ms=1.5, color="tab:blue")
ax[0, 0].axhline(0, color="k", lw=0.5)
ax[0, 0].set_xlabel("k")
ax[0, 0].set_ylabel(r"$\eta_k=\tau_k-\gamma_k$")
ax[0, 0].set_title(r"Bounded displacement, $N=%d$ zeros" % N)
ax[0, 1].hist(eta, bins=200, density=True, color="tab:blue")
ax[0, 1].set_xlabel(r"$\eta_k$")
ax[0, 1].set_title(r"Distribution of $\eta_k$ (all $k\leq N$)")
lam = np.linspace(0, 1, 101)
for kk in range(30):
    gl = (1 - lam) * gam[kk] + lam * tau[kk]
    ax[1, 0].plot(lam, clock(gl) - (kk + 0.5), "-", lw=1)
ax[1, 0].axhline(0, color="k", lw=0.5)
ax[1, 0].set_xlabel(r"$\lambda$")
ax[1, 0].set_ylabel(r"$\mathcal{N}_{\rm RVM}(\gamma_k(\lambda))-(k-\frac{1}{2})$")
ax[1, 0].set_title(r"Moving ordinates in clock units, $k\leq30$ ($\lambda=0$: $\zeta_{Z1}$, $\lambda=1$: $\zeta_{\rm RVM}$)", fontsize=9)
ax[1, 1].plot(tc, Zhardy, "k-", lw=2.2, label=r"Hardy $Z(t)$ of $\zeta$")
cols = ["tab:red", "tab:orange", "tab:green", "tab:cyan", "tab:blue"]
for lam, c in zip(lams, cols):
    ax[1, 1].plot(tc, Zl[lam], "-", color=c, lw=1.0, label=r"$Z_\lambda$, $\lambda=%.2f$" % lam)
ax[1, 1].axhline(0, color="k", lw=0.5)
ax[1, 1].set_xlabel("t")
ax[1, 1].set_title(r"$e^{i\vartheta(t)}\zeta_\lambda(\frac{1}{2}+it)$ along the family")
ax[1, 1].legend(fontsize=7, ncol=2)
fig.tight_layout()
savefig(fig, "fig01_clock_family.png")

fig, ax = plt.subplots(1, 2, figsize=(12, 4))
for sigma, c in zip((2.0, 1.25, 0.75), ("tab:blue", "tab:green", "tab:red")):
    d = np.loadtxt(os.path.join(OUT, "s01_zetaZ1_vs_zeta_sigma%.2f.csv" % sigma), delimiter=",", skiprows=1)
    ax[0].plot(d[:, 0], d[:, 3], ".", ms=2, color=c, label=r"$\sigma=%.2f$" % sigma)
ax[0].set_xlabel("t")
ax[0].set_ylabel(r"$|\zeta_{Z1}/\zeta-1|$")
ax[0].set_title(r"$\zeta_{Z1}$ from %d ordinates versus $\zeta$" % N)
ax[0].legend()
ax[1].plot(tc, Zhardy, "k-", lw=2.5, label=r"$Z(t)$ (mpmath)")
ax[1].plot(tc, Zl[0.0], "r--", lw=1.0, label=r"$Z_{Z1}(t)$ (product)")
ax[1].set_xlabel("t")
ax[1].legend()
ax[1].set_title(r"Critical line: $\lambda=0$ endpoint and Hardy $Z$")
fig.tight_layout()
savefig(fig, "fig02_zetaZ1_vs_zeta.png")

sio.savemat(os.path.join(OUT, "s01_results.mat"), {
    "N": N, "clock_err": clock_err, "eta_blocks": rows, "zetaZ1_cmp": cmp_rows,
    "t_critical": tc, "Z_hardy": Zhardy, "Z_lambda": np.array([Zl[l] for l in lams]), "lambdas": np.array(lams)})
say.close()

# ---- zoom: how the smooth clock places the midpoint ordinates next to the zeros
say2 = Tee("s01b_clock_zoom")
fig, ax = plt.subplots(1, 2, figsize=(12, 4))
T = np.linspace(9.0, 52.0, 4000)
ax[0].step(T, np.searchsorted(gam, T, side="right"), where="post", color="k", lw=1.5, label=r"$N_\zeta(T)$")
ax[0].plot(T, clock(T), "r-", lw=1.5, label=r"$\mathcal{N}_{\rm RVM}(T)$")
for i in range(9):
    ax[0].plot([gam[i], gam[i]], [0, i + 1], "k:", lw=0.8)
ax[0].plot(tau[:9], np.arange(9) + 0.5, "ro", ms=5, label=r"$\tau_k$: $\mathcal{N}_{\rm RVM}(\tau_k)=k-\frac{1}{2}$")
ax[0].plot(gam[:9], np.arange(9) + 0.5, "ks", ms=4, mfc="none", label=r"$(\gamma_k,\,k-\frac{1}{2})$")
ax[0].set_xlabel("T")
ax[0].set_title("Clock and staircase at the first zeros")
ax[0].legend(loc="upper left", fontsize=8)
kz = np.arange(N - 12, N)
T3 = np.linspace(gam[kz[0]] - 0.3, gam[kz[-1]] + 0.3, 6000)
ax[1].step(T3 - 1132480, np.searchsorted(gam, T3, side="right") - 2001000, where="post", color="k", lw=1.5, label=r"$N_\zeta(T)$")
ax[1].plot(T3 - 1132480, clock(T3) - 2001000, "r-", lw=1.5, label=r"$\mathcal{N}_{\rm RVM}(T)$")
ax[1].plot(tau[kz] - 1132480, kz + 0.5 - 2001000, "ro", ms=5, label=r"$\tau_k$")
ax[1].plot(gam[kz] - 1132480, kz + 0.5 - 2001000, "ks", ms=4, mfc="none", label=r"$(\gamma_k,\,k-\frac{1}{2})$")
ax[1].set_xlabel(r"$T-1132480$")
ax[1].set_ylabel(r"count $-\,2001000$")
ax[1].set_title(r"Zoom at the last 12 zeros of the file (mean spacing %.3f)" % (2 * np.pi / np.log(gam[-1] / (2 * np.pi))))
ax[1].legend(loc="upper left", fontsize=8)
fig.tight_layout()
savefig(fig, "fig00_clock_zoom.png")
fig, ax = plt.subplots(1, 1, figsize=(12, 3))
T2 = np.linspace(9.0, 200.0, 20000)
ax.plot(T2, np.searchsorted(gam, T2, side="right") - clock(T2), "k-", lw=0.8)
ax.axhline(0, color="r", lw=0.8)
ax.set_xlabel("T")
ax.set_ylabel(r"$S(T)+O(1/T)$")
ax.set_title(r"$N_\zeta(T)-\mathcal{N}_{\rm RVM}(T)$: the remainder of the Riemann-von Mangoldt formula (jump $+1$ at each $\gamma_k$)")
fig.tight_layout()
savefig(fig, "fig00b_remainder.png")
say2("zoom table near the top of the file:")
say2("%9s %18s %18s %10s %10s", "k", "gamma_k", "tau_k", "eta_k", "N*-N_RVM")
for i in kz:
    say2("%9d %18.6f %18.6f %10.6f %10.6f", i + 1, gam[i], tau[i], tau[i] - gam[i],
         np.searchsorted(gam, gam[i], side="right") - 0.5 - clock(gam[i]))
say2.close()
\end{lstlisting}

\subsection{\texttt{s02\_transition.py}}
\label{app:code-s02-transition}

\begin{lstlisting}[language=Python]
"""Experiment 2: the transition theorem zeta_RVM = varpi (1 + w)(1 + O(1/t)),
the constant varpi_RVM, the vertical limit of 1/zeta_RVM, and zeta_RVM versus zeta."""
import os
import numpy as np
import mpmath as mp
import scipy.io as sio
import matplotlib
matplotlib.use("Agg")
import matplotlib.pyplot as plt
from rvm_common import (Tee, load_ordinates, Product, zeta_from_product, chi, w_transition,
                        write_csv, savefig, OUT)

say = Tee("s02_transition")
gam, tau = load_ordinates()
PR = Product(tau)

# ---- varpi_RVM by Richardson extrapolation on Re s = 2
Ts = 1.0e4 * np.array([0.25, 0.5, 1.0])
v = np.array([zeta_from_product(PR, 2 + 1j * T)[0] / (1 + w_transition(2 + 1j * T)) for T in Ts])
A = np.column_stack([np.ones(3), 1 / Ts, 1 / Ts ** 2])
c = np.linalg.solve(A.astype(complex), v)
varpi = c[0].real
L = 1.0 / varpi
say("varpi_RVM (Richardson, sigma = 2, T = 2500, 5000, 10000) = %.10f %+.2e i", c[0].real, c[0].imag)
say("L = 1/varpi_RVM = %.10f", L)

# ---- relative error of the transition formula
tt = np.exp(np.linspace(np.log(10.0), np.log(1.0e4), 160))
rows = [tt]
hdr = ["t"]
fits = []
for sp in (0.25, 0.5, 1.5):
    s = 0.5 + sp + 1j * tt
    zr = zeta_from_product(PR, s)
    err_w = np.abs(zr / (varpi * (1 + w_transition(s))) - 1)
    rows.append(err_w)
    hdr.append("err_w_sigmap%.2f" % sp)
    sel = tt > 100
    slope = np.polyfit(np.log(tt[sel]), np.log(err_w[sel]), 1)[0]
    fits.append([sp, slope, np.max(err_w[sel] * tt[sel])])
    say("sigma' = %.2f: |zeta_RVM/(varpi(1+w)) - 1|: log-log slope on t>100 = %.3f, max t*err = %.4f",
        sp, slope, np.max(err_w[sel] * tt[sel]))
    if sp < 0.5:
        err_c = np.abs(zr / (varpi * (1 + chi(s))) - 1)
        rows.append(err_c)
        hdr.append("err_chi_sigmap%.2f" % sp)
        say("sigma' = %.2f: |zeta_RVM/(varpi(1+chi)) - 1|: max t*err (t>100) = %.4f",
            sp, np.max(err_c[sel] * tt[sel]))
    wmod = np.abs(w_transition(s))
    rows.append(wmod)
    hdr.append("abs_w_sigmap%.2f" % sp)
    Fv = 1 / zr
    rows.append(np.abs(Fv - L))
    hdr.append("abs_F_minus_L_sigmap%.2f" % sp)
    say("          |w| * t^sigma' in [%.4f, %.4f];  |1/zeta_RVM - L| at t = 1e4: %.3e",
        np.min(wmod * tt ** sp), np.max(wmod * tt ** sp), np.abs(Fv[-1] - L))
write_csv("s02_transition_errors.csv", hdr, rows)

# ---- zeta_RVM versus zeta on two vertical lines and on the critical line
mp.mp.dps = 20
tl = np.linspace(1.0, 120.0, 900)
cmp = {}
for sigma in (2.0, 0.75):
    s = sigma + 1j * tl
    zr = zeta_from_product(PR, s)
    ze = np.array([complex(mp.zeta(complex(x))) for x in s])
    cmp[sigma] = (zr, ze)
    write_csv("s02_zetaRVM_vs_zeta_sigma%.2f.csv" % sigma, ["t", "abs_zeta", "abs_zetaRVM", "abs_varpi_1_plus_w"],
              [tl, np.abs(ze), np.abs(zr), np.abs(varpi * (1 + w_transition(s)))])

fig, ax = plt.subplots(2, 2, figsize=(12, 8))
cols = {0.25: "tab:red", 0.5: "tab:green", 1.5: "tab:blue"}
j = 1
for sp in (0.25, 0.5, 1.5):
    ax[0, 0].plot(tt, tt * rows[j], "-", color=cols[sp], label=r"$\sigma'=%.2f$" % sp)
    j += 4 if sp < 0.5 else 3
ax[0, 0].set_xlabel("t")
ax[0, 0].set_ylabel(r"$t\,|\zeta_{\rm RVM}/(\varpi_{\rm RVM}(1+w))-1|$")
ax[0, 0].set_title(r"Transition theorem: error times $t$ stays bounded, $\varpi_{\rm RVM}=%.10f$" % varpi, fontsize=9)
ax[0, 0].legend()
for sp in (0.25, 0.5, 1.5):
    s = 0.5 + sp + 1j * tt
    ax[0, 1].plot(tt, np.abs(1 / zeta_from_product(PR, s) - L) * tt ** min(sp, 1.0), "-", color=cols[sp], label=r"$\sigma'=%.2f$" % sp)
ax[0, 1].set_xlabel("t")
ax[0, 1].set_ylabel(r"$t^{\min(\sigma',1)}\,|1/\zeta_{\rm RVM}(\sigma+it)-L|$")
ax[0, 1].set_title(r"Vertical limit $L=1/\varpi_{\rm RVM}$, rescaled error", fontsize=9)
ax[0, 1].legend()
for a, sigma in zip((ax[1, 0], ax[1, 1]), (2.0, 0.75)):
    zr, ze = cmp[sigma]
    a.plot(tl, np.abs(ze), "k-", lw=1.2, label=r"$|\zeta(\sigma+it)|$")
    a.plot(tl, np.abs(zr), "r-", lw=1.2, label=r"$|\zeta_{\rm RVM}(\sigma+it)|$")
    a.set_xlabel("t")
    a.set_title(r"$\sigma=%.2f$" % sigma)
    a.legend()
fig.tight_layout()
savefig(fig, "fig03_transition.png")

sio.savemat(os.path.join(OUT, "s02_results.mat"), {"varpi": varpi, "L": L, "richardson": c,
                                                   "slopes": np.array(fits), "t": tt})
say.close()
\end{lstlisting}

\subsection{\texttt{s03\_exponential\_sums.py}}
\label{app:code-s03-exponential-sums}

\begin{lstlisting}[language=Python]
"""Experiment 3: exponential sums S(N;u) = sum_{k<=N} exp(i u o_k) over the midpoint
ordinates tau_k and over the ordinates gamma_k of zeta, and the height transforms M_H(x)."""
import os
import numpy as np
import scipy.io as sio
import matplotlib
matplotlib.use("Agg")
import matplotlib.pyplot as plt
from rvm_common import Tee, load_ordinates, write_csv, savefig, OUT

say = Tee("s03_exponential_sums")
gam, tau = load_ordinates()
N = gam.size
step = 2000
Nk = np.arange(step, N + 1, step)
xs = [2.0, 3.0, 4.0, 6.0, 10.0]
Lam = {2.0: np.log(2), 3.0: np.log(3), 4.0: np.log(2), 6.0: 0.0, 10.0: 0.0}
res = {}
say("%6s %14s %14s %14s %16s %14s" % ("x", "max|S_RVM|", "max|S_RVM|/log", "|S_zeta(N)|",
                                     "Landau -T/2pi*L/sqrt(x)", "|M_H^RVM(x)|"))
for x in xs:
    u = np.log(x)
    SR = np.cumsum(np.exp(1j * u * tau))[Nk - 1]
    SZ = np.cumsum(np.exp(1j * u * gam))[Nk - 1]
    res[x] = (SR, SZ)
    landau = -gam[-1] / (2 * np.pi) * Lam[x] / np.sqrt(x)
    say("%6.1f %14.4f %14.4f %14.2f %16.2f %14.3e" % (x, np.max(np.abs(SR)), np.max(np.abs(SR) / np.log(Nk)),
                                                     np.abs(SZ[-1]), landau, np.abs(SR[-1]) / N))
    write_csv("s03_expsum_x%g.csv" % x, ["N", "re_S_RVM", "im_S_RVM", "re_S_zeta", "im_S_zeta"],
              [Nk, SR.real, SR.imag, SZ.real, SZ.imag])

fig, ax = plt.subplots(1, 2, figsize=(12, 4))
cols = {2.0: "tab:blue", 3.0: "tab:green", 4.0: "tab:purple", 6.0: "tab:orange", 10.0: "tab:red"}
for x in xs:
    ax[0].plot(Nk, np.abs(res[x][0]), "-", color=cols[x], lw=0.8, label=r"$x=%g$" % x)
ax[0].plot(Nk, np.log(Nk), "k--", lw=1, label=r"$\log N$")
ax[0].set_xlabel("N")
ax[0].set_ylabel(r"$|S_{\rm RVM}(N;\log x)|$")
ax[0].set_title(r"Midpoint ordinates: $|\sum_{k\leq N}x^{i\tau_k}|$ stays $O(\log N)$")
ax[0].legend(fontsize=8, ncol=2)
for x in xs:
    ax[1].plot(Nk, res[x][1].real, "-", color=cols[x], lw=0.8, label=r"$x=%g$" % x)
T = np.interp(Nk, np.arange(1, N + 1), gam)
for x in (2.0, 3.0, 4.0):
    ax[1].plot(Nk, -T / (2 * np.pi) * Lam[x] / np.sqrt(x), ":", color=cols[x], lw=1.5)
ax[1].set_xlabel("N")
ax[1].set_ylabel(r"$\mathrm{Re}\sum_{k\leq N}x^{i\gamma_k}$")
ax[1].set_title(r"Zeros of $\zeta$: Landau term $-\frac{T}{2\pi}\frac{\Lambda(x)}{\sqrt{x}}$ (dotted)")
ax[1].legend(fontsize=8, ncol=2)
fig.tight_layout()
savefig(fig, "fig04_exponential_sums.png")
sio.savemat(os.path.join(OUT, "s03_results.mat"), {"N": Nk, "x": np.array(xs),
                                                   "S_RVM": np.array([res[x][0] for x in xs]),
                                                   "S_zeta": np.array([res[x][1] for x in xs])})
say.close()
\end{lstlisting}

\subsection{\texttt{s04\_primes.py}}
\label{app:code-s04-primes}

\begin{lstlisting}[language=Python]
"""Experiment 4: the prime-power count J_RVM, the synthetic prime count pi_RVM, the cell
masses c_n^RVM and the real RVM prime levels p_j^RVM, compared with J, pi(x) and the primes.

J_o(x) = li(x) - 2 Re sum_k Ei((1/2 + i o_k) log x) - log 2 + int_x^inf dt/(t(t^2-1) log t)
for an ordinate list o_k (o = tau: J_RVM; o = gamma: Riemann's explicit formula for J),
pi_o(x) = sum_{m <= log2 x} mu(m)/m J_o(x^(1/m)).
The zero sums are evaluated on a uniform grid in u = log x by the recurrence
exp(i o_k (u + du)) = exp(i o_k u) exp(i o_k du); Ei is exact (scipy exp1) for k <= 300 and
uses Ei(w) = e^w (1/w + 1/w^2 + 2/w^3 + 6/w^4) for k > 300, where |w| > 400."""
import os
import numpy as np
import scipy.io as sio
from scipy import special, integrate
import matplotlib
matplotlib.use("Agg")
import matplotlib.pyplot as plt
from rvm_common import (Tee, load_ordinates, ei_upper, mobius_sieve, primes_upto, write_csv,
                        savefig, clock_inverse, OUT)

say = Tee("s04_primes")
gam, tau = load_ordinates()
KE = 300


def zero_sum_ei(ords, u, K=None, taper=None):
    """sum_k Ei((1/2 + i o_k) u) on a uniform grid u, k <= K, optional weights."""
    o = ords if K is None else ords[:K]
    wts = np.ones(o.size) if taper is None else taper
    rho = 0.5 + 1j * o
    out = np.zeros(u.size, dtype=complex)
    # exact part
    for i in range(0, KE):
        out += wts[i] * ei_upper(rho[i] * u)
    # asymptotic part by recurrence
    r = rho[KE:]
    b = [wts[KE:] / r ** m for m in (1, 2, 3, 4)]
    du = u[1] - u[0]
    stepf = np.exp(1j * o[KE:] * du)
    for j in range(u.size):
        if j % 200 == 0:
            E = np.exp(1j * o[KE:] * u[j])
        A = [E @ bm for bm in b]
        uj = u[j]
        out[j] += np.exp(uj / 2) * (A[0] / uj + A[1] / uj ** 2 + 2 * A[2] / uj ** 3 + 6 * A[3] / uj ** 4)
        E = E * stepf
    return out


def tail_integral(x):
    return integrate.quad(lambda t: 1.0 / (t * (t * t - 1.0) * np.log(t)), x, np.inf, limit=200)[0]


def J_from(ords, u, K=None, taper=None):
    x = np.exp(u)
    S = zero_sum_ei(ords, u, K, taper)
    ti = np.array([tail_integral(v) for v in x])
    return special.expi(u) - 2 * S.real - np.log(2.0) + ti


def pi_from_J(u, J):
    mu = mobius_sieve(64)
    out = np.zeros(u.size)
    for i, ui in enumerate(u):
        m = 1
        while ui / m >= np.log(2.0) - 1e-12:
            if mu[m] != 0:
                out[i] += mu[m] / m * np.interp(ui / m, u, J)
            m += 1
    return out


def exact_J_pi(x):
    P = primes_upto(int(x.max()) + 1)
    pi = np.searchsorted(P, x, side="right").astype(float)
    J = np.zeros(x.size)
    for m in range(1, 40):
        J += np.searchsorted(P, x ** (1.0 / m) + 1e-12, side="right") / m
    return J, pi


def run(xmax, npts, K, label):
    u = np.linspace(np.log(2.0), np.log(xmax), npts)
    if K is None:
        tg = tr = None
        Kg = Kr = None
    else:
        Kg = Kr = K
        v = np.arange(K) / K
        tg = tr = np.where(v < 0.6, 1.0, np.cos(np.pi / 2 * (v - 0.6) / 0.4) ** 2)
    JR = J_from(tau, u, Kr, tr)
    JG = J_from(gam, u, Kg, tg)
    PR = pi_from_J(u, JR)
    PG = pi_from_J(u, JG)
    x = np.exp(u)
    Je, Pe = exact_J_pi(x)
    write_csv("s04_prime_counts_%s.csv" % label, ["x", "J_exact", "J_from_zeta_zeros", "J_RVM",
                                                  "pi_exact", "pi_from_zeta_zeros", "pi_RVM", "li"],
              [x, Je, JG, JR, Pe, PG, PR, special.expi(u)])
    return x, Je, JG, JR, Pe, PG, PR


# x <= 100 with all zeros; x <= 1000 with a smooth cut at K zeros
x1, Je1, JG1, JR1, Pe1, PG1, PR1 = run(100.0, 3000, None, "x100_allzeros")
say("x <= 100, all %d zeros of each list, %d grid points:" % (gam.size, x1.size))
say("  max |pi_from_zeta_zeros - pi| away from primes (distance > 0.2): %.4f",
    np.max(np.abs(PG1 - Pe1)[np.min(np.abs(x1[:, None] - primes_upto(101)[None, :]), axis=1) > 0.2]))
say("  pi_RVM(100) = %.4f, pi(100) = %d, li(100) = %.4f, J_RVM(100) = %.4f, J(100) = %.4f",
    PR1[-1], Pe1[-1], special.expi(np.log(100.0)), JR1[-1], Je1[-1])
dPR = np.diff(PR1)
say("  pi_RVM monotone on [2,100]: %s (min increment on the grid %.3e)", bool(np.all(dPR > 0)), dPR.min())

tauK = 2 * np.pi * 1000.0 * 2.5
K = int(np.searchsorted(tau, tauK))
x2, Je2, JG2, JR2, Pe2, PG2, PR2 = run(1000.0, 40000, K, "x1000_K%d" % K)
say("x <= 1000, smooth cut at K = %d zeros (tau_K = %.1f), %d grid points:" % (K, tau[K - 1], x2.size))
say("  pi_RVM(1000) = %.4f, pi(1000) = %d, li(1000) = %.4f", PR2[-1], Pe2[-1], special.expi(np.log(1000.0)))

# cell masses and real prime levels from the fine x<=1000 run
n = np.arange(2, 1001)
pin = np.interp(n, x2, PR2)
c = np.diff(np.concatenate([[0.0], pin]))
P = primes_upto(1000)
isp = np.isin(n, P).astype(float)
say("cell masses c_n = pi_RVM(n) - pi_RVM(n-1):")
say("%5s %10s %10s %8s" % ("n", "c_n", "1/log n", "prime"))
for i in range(0, 40):
    say("%5d %10.5f %10.5f %8d" % (n[i], c[i], 1 / np.log(n[i]), isp[i]))
say("sum_{n<=1000} c_n = %.4f,  pi(1000) = %d", c.sum(), len(P))
levels = np.arange(1, int(np.floor(PR2.max())) + 1)
pj = np.array([x2[np.argmax(PR2 >= j)] for j in levels])
say("real RVM prime levels p_j^RVM (first crossing of level j) versus p_j:")
say("%5s %12s %8s %10s" % ("j", "p_j^RVM", "p_j", "diff"))
for j in range(min(30, levels.size)):
    say("%5d %12.4f %8d %10.4f" % (levels[j], pj[j], P[j], pj[j] - P[j]))
write_csv("s04_cell_masses.csv", ["n", "c_n_RVM", "inv_log_n", "is_prime"], [n, c, 1 / np.log(n), isp])
m = min(levels.size, P.size)
write_csv("s04_rvm_prime_levels.csv", ["j", "p_j_RVM", "p_j"], [levels[:m], pj[:m], P[:m]])

# ---- figures (linear axes)
fig, ax = plt.subplots(1, 2, figsize=(12, 4))
ax[0].step(x1, Pe1, where="post", color="k", lw=1.6, label=r"$\pi(x)$")
ax[0].plot(x1, PG1, "b-", lw=0.6, label=r"$\pi(x)$ from %d zeros of $\zeta$" % gam.size)
ax[0].plot(x1, PR1, "r-", lw=1.4, label=r"$\pi_{\rm RVM}(x)$ from %d $\tau_k$" % tau.size)
ax[0].plot(x1, special.expi(np.log(x1)) - special.expi(np.log(2.0)), "g--", lw=1, label=r"$\mathrm{li}(x)-\mathrm{li}(2)$")
ax[0].set_xlabel("x")
ax[0].set_title(r"Prime counts on $[2,100]$")
ax[0].legend(fontsize=7)
ax[1].step(x2, Pe2, where="post", color="k", lw=1.2, label=r"$\pi(x)$")
ax[1].plot(x2, PR2, "r-", lw=1.0, label=r"$\pi_{\rm RVM}(x)$")
ax[1].set_xlabel("x")
ax[1].set_title(r"Prime counts on $[2,1000]$ ($K=%d$ ordinates, smooth cut)" % K)
ax[1].legend(fontsize=8)
fig.tight_layout()
savefig(fig, "fig05_prime_counts.png")

fig, ax = plt.subplots(1, 2, figsize=(12, 4))
nn = n[:98]
ax[0].bar(nn, isp[:98], width=0.8, color="0.75", label="prime indicator")
ax[0].plot(nn, c[:98], "r.-", lw=0.8, ms=4, label=r"$c_n^{\rm RVM}$")
ax[0].plot(nn, 1 / np.log(nn), "g--", lw=1, label=r"$1/\log n$")
ax[0].set_xlabel("n")
ax[0].set_title(r"Cell masses $c_n^{\rm RVM}=\pi_{\rm RVM}(n)-\pi_{\rm RVM}(n-1)$")
ax[0].legend(fontsize=8)
ax[1].plot(P[:m], pj[:m], "r.", ms=4, label=r"$(p_j,\,p_j^{\rm RVM})$")
ax[1].plot([0, 1000], [0, 1000], "k-", lw=0.8, label="diagonal")
ax[1].set_xlabel(r"$p_j$")
ax[1].set_ylabel(r"$p_j^{\rm RVM}$")
ax[1].set_title(r"Real RVM prime levels $\pi_{\rm RVM}(p_j^{\rm RVM})=j$ versus primes")
ax[1].legend(fontsize=8)
fig.tight_layout()
savefig(fig, "fig06_rvm_primes.png")

sio.savemat(os.path.join(OUT, "s04_results.mat"), {"x100": x1, "pi100": Pe1, "piZeta100": PG1, "piRVM100": PR1,
                                                   "x1000": x2, "pi1000": Pe2, "piRVM1000": PR2,
                                                   "c_n": c, "p_RVM": pj, "K": K})
say.close()
\end{lstlisting}

\subsection{\texttt{s05\_mertens.py}}
\label{app:code-s05-mertens}

\begin{lstlisting}[language=Python]
"""Experiment 5: the Mertens function M, the synthetic Mertens function M_RVM^crit and the
ternary walk M_h with values in {-1,0,+1} on the squarefree support.

Explicit formulas (simple zeros, residues a = 1/(rho zeta'(rho))):
  M_zeta(x) = -2 + 2 Re sum_k a_k x^(rho_k) + trivial terms        (zeros of zeta)
  M_RVM(x)  = -2 + 2 Re sum_k b_k x^(1/2 + i tau_k) + trivial terms (b_k = 1/(rho_k zeta_RVM'(rho_k)))
Ternary walk: mu_h(n) = mu(n)^2 * sign, sign = +1 if M_h(n-1) < target(n) else -1."""
import os
import math
import numpy as np
import mpmath as mp
import scipy.io as sio
from scipy.interpolate import PchipInterpolator
import matplotlib
matplotlib.use("Agg")
import matplotlib.pyplot as plt
from rvm_common import (Tee, load_ordinates, Product, logG, clock_inverse, tail_coefficients,
                        theta_rs, mobius_sieve, write_csv, savefig, OUT)

say = Tee("s05_mertens")
gam, tau = load_ordinates()
Nz = gam.size
XMAX = 10 ** 6
K0 = 1000

# ---------------------------------------------------------------- RVM residues
c1, c2, c3 = tail_coefficients(clock_inverse(float(Nz)))
t2 = tau ** 2
zpR = np.zeros(K0, dtype=complex)
for k in range(K0):
    tk = tau[k]
    fac = (t2 - tk ** 2) / (t2 + 0.25)
    fac[k] = 1.0
    q = -0.25 - tk ** 2
    lmod = np.sum(np.log(np.abs(fac))) + c1 * q - c2 * q ** 2 + c3 * q ** 3
    rho = 0.5 + 1j * tk
    lP = np.log(2j * tk / (0.25 + tk ** 2)) + lmod
    zpR[k] = (-1) ** k * np.exp(lP - logG(rho))
ratio = zpR / np.log(tau[:K0] / (2 * np.pi))
PR = Product(tau)
Ts = 1.0e4 * np.array([0.25, 0.5, 1.0])
from rvm_common import zeta_from_product, w_transition
v = np.array([zeta_from_product(PR, 2 + 1j * T)[0] / (1 + w_transition(2 + 1j * T)) for T in Ts])
varpi = np.linalg.solve(np.column_stack([np.ones(3), 1 / Ts, 1 / Ts ** 2]).astype(complex), v)[0].real
sel = np.arange(K0 // 2, K0)
cfit = np.mean((ratio[sel] / varpi - 1) * tau[sel])
say("zeta_RVM'(rho_k) exact for k <= %d; zeta_RVM'(rho_k)/log(tau_k/2pi) at k = %d: %.10f %+.2e i",
    K0, K0, ratio[-1].real, ratio[-1].imag)
say("varpi_RVM = %.10f; fit zeta_RVM'(rho_k) = varpi log(tau_k/2pi)(1 + c/tau_k), c = %.4f %+.4f i",
    varpi, cfit.real, cfit.imag)
zR = varpi * np.log(tau / (2 * np.pi)) * (1 + cfit / tau)
zR[:K0] = zpR
bR = 1.0 / ((0.5 + 1j * tau) * zR)

# ---------------------------------------------------------------- residues of zeta
def rs_Z(t):
    a = np.sqrt(t / (2 * np.pi))
    Nn = np.floor(a).astype(int)
    p = a - Nn
    nmax = Nn.max()
    n = np.arange(1, nmax + 1)
    ph = theta_rs(t)[:, None] - t[:, None] * np.log(n)[None, :]
    mask = n[None, :] <= Nn[:, None]
    Z = 2 * np.sum(np.cos(ph) * mask / np.sqrt(n)[None, :], axis=1)
    cc = np.cos(2 * np.pi * p)
    C0 = np.cos(2 * np.pi * (p ** 2 - p - 1.0 / 16)) / cc
    return Z + (-1.0) ** (Nn - 1) * (t / (2 * np.pi)) ** (-0.25) * C0


mp.mp.dps = 20
KM = 1000
zG = np.zeros(Nz, dtype=complex)
for k in range(KM + 50):
    zG[k] = complex(mp.zeta(mp.mpc(0.5, gam[k]), 1, 1))
h = 1e-5
for i0 in range(KM, Nz, 4000):
    i1 = min(Nz, i0 + 4000)
    t = gam[i0:i1]
    Zp = (rs_Z(t + h) - rs_Z(t - h)) / (2 * h)
    val = -1j * np.exp(-1j * theta_rs(t)) * Zp
    if i0 == KM:
        chk = np.max(np.abs(val[:50] / zG[KM:KM + 50] - 1))
        say("Riemann-Siegel zeta'(rho) versus mpmath at k = %d..%d: max relative difference %.2e",
            KM + 1, KM + 50, chk)
        zG[KM + 50:i1] = val[50:]
    else:
        zG[i0:i1] = val
aG = 1.0 / ((0.5 + 1j * gam) * zG)
S2G = 2 * np.sum(np.abs(aG) ** 2)
S2R = 2 * np.sum(np.abs(bR) ** 2)
say("median |zeta'(rho)| = %.4f;  spectral sums 2 sum|a|^2: zeta %.6f, RVM %.6f", np.median(np.abs(zG)), S2G, S2R)


def trivial_zeta(x, n=6):
    v = np.zeros_like(x)
    for k in range(1, n + 1):
        v += (-1) ** (k - 1) * (2 * np.pi / x) ** (2 * k) / (math.factorial(2 * k) * k * float(mp.zeta(2 * k + 1)))
    return v


def trivial_rvm(x, n=6):
    v = np.zeros_like(x)
    for k in range(1, n + 1):
        s = -2.0 * k
        lP = np.real(PR.logP(s)[0])
        v -= x ** (-2 * k) * (2 * k + 1) * np.pi ** k * 2 * (-1) ** k / (math.factorial(k) * np.exp(lP))
    return v


def explicit_uniform(ords, wts, u):
    """2 Re sqrt(x) sum_k w_k exp(i o_k u) on a uniform u grid (recurrence)."""
    out = np.zeros(u.size)
    st = np.exp(1j * ords * (u[1] - u[0]))
    for j in range(u.size):
        if j % 200 == 0:
            E = np.exp(1j * ords * u[j])
        out[j] = 2 * np.exp(u[j] / 2) * np.real(E @ wts)
        E = E * st
    return out


# ---------------------------------------------------------------- exact M by sieve
mu = mobius_sieve(XMAX)
M = np.cumsum(mu.astype(np.int64))
sqf = mu != 0
say("M(10^6) = %d;  squarefree density Q(10^6)/10^6 = %.6f (6/pi^2 = %.6f)", M[XMAX], sqf[1:].mean(), 6 / np.pi ** 2)

# ---------------------------------------------------------------- explicit formulas on [2, 2000]
u = np.linspace(np.log(2.0), np.log(2000.0), 2500)
xg = np.exp(u)
Mz = -2 + explicit_uniform(gam, aG, u) + trivial_zeta(xg)
Mc = -2 + explicit_uniform(tau, bR, u) + trivial_rvm(xg)
Mfl = M[np.floor(xg).astype(int)]
off = np.abs(xg - np.round(xg)) > 0.05
say("explicit formula with %d zeros of zeta on [2,2000]: rms(M_zeta - M) = %.4f off integers (|x-n|>0.05)",
    Nz, np.sqrt(np.mean((Mz - Mfl)[off] ** 2)))
write_csv("s05_explicit_formulas.csv", ["x", "M_floor", "M_zeta_explicit", "M_RVM_explicit"], [xg, Mfl, Mz, Mc])

# ---------------------------------------------------------------- integer target for M_RVM^crit
def target_grid(X, div=24):
    x = [1.0]
    while x[-1] < X:
        x.append(min(X, x[-1] + max(1.0, np.floor(np.sqrt(x[-1]) / div))))
    return np.array(x)


def tapered(x, ords, a, C1=2.0, C2=4.0):
    S = np.zeros(x.size, dtype=complex)
    lx = np.log(x)
    i0 = 0
    while i0 < x.size:
        i1 = min(x.size, i0 + 200)
        K = max(1, int(np.searchsorted(ords, C2 * 2 * np.pi * np.sqrt(x[i1 - 1]))))
        uu = ords[None, :K] / (2 * np.pi * np.sqrt(x[i0:i1, None]))
        w = np.where(uu <= C1, 1.0, np.where(uu >= C2, 0.0, np.cos(np.pi / 2 * (uu - C1) / (C2 - C1)) ** 2))
        E = np.exp(1j * lx[i0:i1, None] * ords[None, :K])
        S[i0:i1] = np.sqrt(x[i0:i1]) * np.sum(E * w * a[None, :K], axis=1)
        i0 = i1
    return S


xT = target_grid(float(XMAX))
TT = -2 + 2 * np.real(tapered(xT, tau, bR)) + trivial_rvm(xT)
nn = np.arange(1, XMAX + 1)
T = np.concatenate([[0.0], PchipInterpolator(xT, TT)(nn)])
nchk = np.arange(1000, 100001, 997)
full = -2 + 2 * np.real(np.array([np.sqrt(n) * np.sum(np.exp(1j * np.log(n) * tau) * bR) for n in nchk])) + trivial_rvm(nchk.astype(float))
say("target on %d grid points; at %d integers <= 1e5: max |M_RVM(n) (all zeros) - target(n)| = %.4f",
    xT.size, nchk.size, np.max(np.abs(full - T[nchk])))

# ---------------------------------------------------------------- ternary walk
mu_h = np.zeros(XMAX + 1, dtype=np.int8)
Mh = np.zeros(XMAX + 1)
sgn = np.zeros(XMAX + 1, dtype=np.int8)
cur = 0.0
for n in range(1, XMAX + 1):
    if sqf[n]:
        sgn[n] = 1 if cur < T[n] else -1
        cur += sgn[n]
        mu_h[n] = sgn[n]
    Mh[n] = cur
err = Mh[1:] - T[1:]
say("ternary walk: max |M_h - target| on [1,1e6] = %.4f, on [1e4,1e6] = %.4f", np.max(np.abs(err)), np.max(np.abs(err[9999:])))
dens = lambda a: [np.mean(a[1:] == -1), np.mean(a[1:] == 0), np.mean(a[1:] == 1)]
say("densities of (-1, 0, +1):  mu   %.6f %.6f %.6f", *dens(mu))
say("                           mu_h %.6f %.6f %.6f   (3/pi^2 = %.6f, 1 - 6/pi^2 = %.6f)", *dens(mu_h), 3 / np.pi ** 2, 1 - 6 / np.pi ** 2)
say("sign agreement mu_h = mu on squarefrees: %.6f;  sum mu mu_h / Q = %.6f",
    np.mean(mu_h[sqf] == mu[sqf]), np.sum(mu_h.astype(int) * mu.astype(int)) / sqf[1:].sum())
say("max |M(n)|/sqrt(n) on [10,1e6] = %.4f;  max |M_RVM(n)|/sqrt(n) = %.4f;  max |M_h(n)|/sqrt(n) = %.4f",
    np.max(np.abs(M[10:]) / np.sqrt(nn[9:])), np.max(np.abs(T[10:]) / np.sqrt(nn[9:])), np.max(np.abs(Mh[10:]) / np.sqrt(nn[9:])))
say("\nNUMERICAL EXAMPLE OF THE TERNARY DIVISION (first 40 integers)")
say("%4s %4s %6s %10s %8s %6s %6s %6s %6s" % ("n", "mu", "mu^2", "target", "M_h(n-1)", "sign", "mu_h", "M_h", "M"))
rows = []
for n in range(1, 41):
    say("%4d %4d %6d %10.4f %8d %6s %6d %6d %6d" % (n, mu[n], abs(mu[n]), T[n], Mh[n - 1],
                                                     ("+1" if sgn[n] > 0 else ("-1" if sgn[n] < 0 else "-")),
                                                     mu_h[n], Mh[n], M[n]))
    rows.append([n, mu[n], abs(mu[n]), T[n], Mh[n - 1], sgn[n], mu_h[n], Mh[n], M[n]])
rows = np.array(rows, dtype=float)
write_csv("s05_ternary_example.csv", ["n", "mu", "mu2", "target", "Mh_prev", "sign", "mu_h", "Mh", "M"],
          [rows[:, i] for i in range(rows.shape[1])])

Ns = 10 ** np.arange(3, 7)
say("\nENERGY E_N = (1/N) sum_{n<=N} W(n)^2")
say("%10s %14s %14s %14s" % ("N", "E_N(M)", "E_N(M_RVM)", "E_N(M_h)"))
EN = []
for Nn in Ns:
    EN.append([Nn, np.mean(M[1:Nn + 1].astype(float) ** 2), np.mean(T[1:Nn + 1] ** 2), np.mean(Mh[1:Nn + 1] ** 2)])
    say("%10d %14.4f %14.4f %14.4f" % tuple(EN[-1]))
EN = np.array(EN)
write_csv("s05_energy.csv", ["N", "E_M", "E_MRVM", "E_Mh"], [EN[:, i] for i in range(4)])
nall = np.arange(1, XMAX + 1)
write_csv("s05_walks.csv", ["n", "M", "M_RVM_target", "M_h"], [nall, M[nall], T[nall], Mh[nall]])
mu_h[1:].astype(np.int8).tofile(os.path.join(OUT, "mu_h_int8_1e6.bin"))

# ---------------------------------------------------------------- figures (linear axes)
fig, ax = plt.subplots(1, 2, figsize=(12, 4))
w = np.arange(1, 2001)
ax[0].step(w, M[w], where="post", color="k", lw=1.0, label=r"$M(n)$")
ax[0].plot(w, T[w], "r-", lw=1.2, label=r"$M_{\rm RVM}^{\rm crit}(n)$")
ax[0].step(w, Mh[w], where="post", color="tab:green", lw=0.8, label=r"$M_h(n)$")
ax[0].set_xlabel("n")
ax[0].set_title(r"Mertens functions on $[1,2000]$, every integer")
ax[0].legend(fontsize=8)
nb = 2000
bw = XMAX // nb
xb = (np.arange(nb) + 0.5) * bw
Mb = M[1:].reshape(nb, bw)
Tb = T[1:].reshape(nb, bw)
ax[1].fill_between(xb, Mb.min(1), Mb.max(1), color="0.6", label=r"$M(n)$ (range per block)")
ax[1].plot(xb, Tb.mean(1), "r-", lw=1.0, label=r"$M_{\rm RVM}^{\rm crit}(n)$")
ax[1].plot(xb, np.sqrt(xb), "k--", lw=0.8, label=r"$\pm\sqrt{n}$")
ax[1].plot(xb, -np.sqrt(xb), "k--", lw=0.8)
ax[1].set_xlabel("n")
ax[1].set_title(r"$[1,10^6]$, blocks of %d integers" % bw)
ax[1].legend(fontsize=8)
fig.tight_layout()
savefig(fig, "fig07_mertens.png")

fig, ax = plt.subplots(1, 2, figsize=(12, 4))
w = np.arange(1, 61)
ax[0].vlines(w - 0.15, 0, mu[w], color="k", lw=2, label=r"$\mu(n)$")
ax[0].vlines(w + 0.15, 0, mu_h[w], color="tab:green", lw=2, label=r"$\mu_h(n)$")
ax[0].axhline(0, color="k", lw=0.5)
ax[0].set_xlabel("n")
ax[0].set_yticks([-1, 0, 1])
ax[0].set_title(r"Ternary values: same zero set $\{\mu(n)=0\}$, signs chosen by tracking")
ax[0].legend(fontsize=8)
w = np.arange(1, 301)
ax[1].step(w, M[w], where="post", color="k", lw=1.0, label=r"$M(n)$")
ax[1].plot(w, T[w], "r.-", lw=0.8, ms=3, label=r"target $M_{\rm RVM}^{\rm crit}(n)$")
ax[1].step(w, Mh[w], where="post", color="tab:green", lw=1.2, label=r"$M_h(n)$")
ax[1].set_xlabel("n")
ax[1].set_title(r"The walk $M_h$ follows $M_{\rm RVM}^{\rm crit}$ within $\pm1$ step")
ax[1].legend(fontsize=8)
fig.tight_layout()
savefig(fig, "fig08_ternary_walk.png")

fig, ax = plt.subplots(1, 2, figsize=(12, 4))
sel = xg <= 300
ax[0].step(np.arange(1, 301), M[1:301], where="post", color="k", lw=1.0, label=r"$M(x)$")
ax[0].plot(xg[sel], Mz[sel], "b-", lw=0.7, label=r"explicit formula, %d zeros of $\zeta$" % Nz)
ax[0].plot(xg[sel], Mc[sel], "r-", lw=1.0, label=r"$M_{\rm RVM}^{\rm crit}(x)$, %d $\tau_k$" % Nz)
ax[0].set_xlabel("x")
ax[0].set_title("Explicit formulas on [2,300]")
ax[0].legend(fontsize=7)
kk = np.arange(5000)
ax[1].plot(gam[kk], np.abs(zG[kk]) / np.log(gam[kk] / (2 * np.pi)), "b.", ms=1.5, label=r"$|\zeta'(\rho_k)|/\log(\gamma_k/2\pi)$")
ax[1].plot(tau[kk], np.abs(zR[kk]) / np.log(tau[kk] / (2 * np.pi)), "r.", ms=1.5, label=r"$|\zeta_{\rm RVM}'(\rho^{\rm RVM}_k)|/\log(\tau_k/2\pi)$")
ax[1].axhline(varpi, color="k", lw=0.8, ls="--", label=r"$\varpi_{\rm RVM}$")
ax[1].set_xlabel("ordinate")
ax[1].set_ylim(0, 4)
ax[1].set_title(r"Derivatives at the zeros, $k\leq5000$")
ax[1].legend(fontsize=7)
fig.tight_layout()
savefig(fig, "fig09_explicit_weights.png")

sio.savemat(os.path.join(OUT, "s05_results.mat"), {"varpi": varpi, "cfit": cfit, "S2G": S2G, "S2R": S2R,
                                                   "energy": EN, "ternary_example": rows,
                                                   "x_explicit": xg, "M_zeta_explicit": Mz, "M_RVM_explicit": Mc})
say.close()
\end{lstlisting}

\subsection{\texttt{s06\_prime\_ratios.py}}
\label{app:code-s06-prime-ratios}

\begin{lstlisting}[language=Python]
"""Experiment 6: the three ratios li/pi_RVM, pi/li, pi/pi_RVM and the square-root estimate
pi_RVM(x) - li(x) = O(sqrt(x) (log log x)^2 / log x), from the prime counts written by s04_primes.py.

Input: output/s04_prime_counts_x1000_K17061.csv (columns x, pi_exact, pi_RVM, li, ...).
B(x) = sqrt(x) (log log x)^2 / log x is the scale of the square-root estimate."""
import os
import numpy as np
import matplotlib
matplotlib.use("Agg")
import matplotlib.pyplot as plt
from rvm_common import Tee, savefig, write_csv, OUT

say = Tee("s06_prime_ratios")
D = np.genfromtxt(os.path.join(OUT, "s04_prime_counts_x1000_K17061.csv"), delimiter=",", names=True)
x, pe, pr, li = D["x"], D["pi_exact"], D["pi_RVM"], D["li"]
say("grid: %d points on [%.0f, %.0f]; pi_RVM from the first 17061 ordinates tau_k (smooth cut)"
    % (x.size, x[0], x[-1]))

sel = x >= 10
r1, r2, r3 = li[sel] / pr[sel], pe[sel] / li[sel], pe[sel] / pr[sel]
xs = x[sel]
say("")
say("ratios at selected x")
say("%8s %12s %12s %12s %12s %12s %12s" % ("x", "pi", "pi_RVM", "li", "li/pi_RVM", "pi/li", "pi/pi_RVM"))
for X in [10, 20, 50, 100, 200, 500, 1000]:
    i = int(np.argmin(np.abs(x - X)))
    say("%8.1f %12.4f %12.4f %12.4f %12.6f %12.6f %12.6f"
        % (x[i], pe[i], pr[i], li[i], li[i] / pr[i], pe[i] / li[i], pe[i] / pr[i]))

s8 = x >= 8
B = np.sqrt(x[s8]) * np.log(np.log(x[s8])) ** 2 / np.log(x[s8])
e_rvm = (pr[s8] - li[s8]) / B
e_pi = (pe[s8] - li[s8]) / B
say("")
say("square-root scale B(x) = sqrt(x) (log log x)^2 / log x on [8, 1000]")
say("max |pi_RVM - li| / B = %.4f   (at x = %.2f)" % (np.max(np.abs(e_rvm)), x[s8][np.argmax(np.abs(e_rvm))]))
say("max |pi - li| / B     = %.4f   (at x = %.2f)" % (np.max(np.abs(e_pi)), x[s8][np.argmax(np.abs(e_pi))]))
for X in [100, 1000]:
    m = x[s8] >= X / 2
    say("max |pi_RVM - li| / B on [%d, 1000] = %.4f" % (X // 2, np.max(np.abs(e_rvm[m]))))

# CSV files: full grid
write_csv("s06_prime_ratios.csv", ["x", "li_over_pi_RVM", "pi_over_li", "pi_over_pi_RVM"],
          [xs, r1, r2, r3])
write_csv("s06_square_root_scale.csv", ["x", "pi_RVM_minus_li_over_B", "pi_minus_li_over_B"],
          [x[s8], e_rvm, e_pi])

# ---- figure (linear axes)
fig, ax = plt.subplots(1, 2, figsize=(12, 4))
ax[0].plot(xs, r3, "k-", lw=0.8, label=r"$\pi(x)/\pi_{\rm RVM}(x)$")
ax[0].plot(xs, r2, "b-", lw=0.8, label=r"$\pi(x)/\mathrm{li}(x)$")
ax[0].plot(xs, r1, "r-", lw=1.4, label=r"$\mathrm{li}(x)/\pi_{\rm RVM}(x)$")
ax[0].axhline(1.0, color="0.5", lw=0.8, ls="--")
ax[0].set_xlabel("x")
ax[0].set_title(r"The three ratios on $[10,1000]$")
ax[0].legend(fontsize=8)
ax[1].plot(x[s8], e_pi, "b-", lw=0.6, label=r"$(\pi-\mathrm{li})/B$")
ax[1].plot(x[s8], e_rvm, "r-", lw=1.4, label=r"$(\pi_{\rm RVM}-\mathrm{li})/B$")
ax[1].axhline(0.0, color="0.5", lw=0.8, ls="--")
ax[1].set_xlabel("x")
ax[1].set_title(r"Square-root scale $B(x)=\sqrt{x}\,(\log\log x)^2/\log x$")
ax[1].legend(fontsize=8)
fig.tight_layout()
savefig(fig, "fig10_prime_ratios.png")
say.close()
\end{lstlisting}

\subsection{\texttt{s07\_product\_identity.py}}
\label{app:code-s07-product-identity}

\begin{lstlisting}[language=Python]
"""Experiment 7: the value at infinity of the endpoint.

varpi_Z1 = lim_{sigma -> +inf} zeta_Z1(sigma) = varpi_RVM * prod_k (1/4 + tau_k^2) / (1/4 + gamma_k^2),
so zeta_Z1 = zeta holds exactly when the ordered partial products
P_K = prod_{k <= K} (1/4 + gamma_k^2) / (1/4 + tau_k^2)
converge to varpi_RVM.  The partial products are computed from data/zeross.bin (gamma_k) and
data/zLamb.bin (tau_k) and compared with varpi_RVM from s02_transition.py."""
import numpy as np
import matplotlib
matplotlib.use("Agg")
import matplotlib.pyplot as plt
from rvm_common import Tee, load_ordinates, savefig, write_csv

say = Tee("s07_product_identity")
gam, tau = load_ordinates()
VARPI = 0.9619415807          # Richardson value of varpi_RVM (s02_transition.py)
N = gam.size
logP = np.cumsum(np.log1p(0.25 / gam ** 2) + 2 * np.log(gam) - np.log1p(0.25 / tau ** 2) - 2 * np.log(tau))
P = np.exp(logP)
say("N = %d ordinates; varpi_RVM = %.10f" % (N, VARPI))
say("%10s %16s %14s" % ("K", "P_K", "P_K - varpi"))
for K in [10, 100, 1000, 10 ** 4, 10 ** 5, 10 ** 6, N]:
    say("%10d %16.10f %14.3e" % (K, P[K - 1], P[K - 1] - VARPI))
for K0 in [10 ** 4, 10 ** 5, 10 ** 6]:
    seg = P[K0 - 1:]
    say("K >= %7d: min P_K = %.10f, max P_K = %.10f, mean = %.10f" % (K0, seg.min(), seg.max(), seg.mean()))
# the flow along the family: Upsilon(lam) = sum_k 2 gamma_k(lam) eta_k / (1/4 + gamma_k(lam)^2),
# eta_k = tau_k - gamma_k, gamma_k(lam) = (1 - lam) gamma_k + lam tau_k; d/dlam log varpi_lam = -Upsilon(lam)
eta = tau - gam
say("")
say("flow along the lambda-family")
for lam in [0.0, 0.25, 0.5, 0.75, 1.0]:
    gl = (1 - lam) * gam + lam * tau
    say("Upsilon(%.2f) = %.10f" % (lam, np.sum(2 * gl * eta / (0.25 + gl ** 2))))
say("sum_k 2 eta_k^2 / gamma_k^2 = %.10f" % np.sum(2 * eta ** 2 / gam ** 2))
total = np.sum(np.log((0.25 + tau ** 2) / (0.25 + gam ** 2)))
say("integral_0^1 Upsilon = sum_k log((1/4+tau_k^2)/(1/4+gamma_k^2)) = %.10f" % total)
say("log L = log(1/varpi_RVM)                               = %.10f" % np.log(1 / VARPI))
step = 50
idx = np.arange(step - 1, N, step)
write_csv("s07_partial_products.csv", ["K", "P_K"], [idx + 1, P[idx]])

fig, ax = plt.subplots(1, 2, figsize=(12, 4))
k1 = np.arange(1, 10 ** 4 + 1)
ax[0].plot(k1, P[:10 ** 4], "r-", lw=0.6, label=r"$P_K$")
ax[0].axhline(VARPI, color="k", lw=0.8, ls="--", label=r"$\varpi_{\rm RVM}$")
ax[0].set_xlabel("K")
ax[0].set_title(r"Partial products $P_K$, $K\leq 10^4$")
ax[0].legend(fontsize=8)
sel = idx >= 10 ** 5 - 1
ax[1].plot(idx[sel] + 1, (P[idx[sel]] - VARPI) * 1e6, "r-", lw=0.6)
ax[1].axhline(0.0, color="k", lw=0.8, ls="--")
ax[1].set_xlabel("K")
ax[1].set_ylabel(r"$(P_K-\varpi_{\rm RVM})\cdot 10^{6}$")
ax[1].set_title(r"$10^5\leq K\leq %d$" % N)
fig.tight_layout()
savefig(fig, "fig11_product_identity.png")
say.close()
\end{lstlisting}


\subsection{\texttt{s08\_transport\_moments.py}}
\label{app:code-s08-transport-moments}

\begin{lstlisting}[language=Python]
#!/usr/bin/env python3
"""s08_transport_moments.py
Transport tau -> gamma on the first rung (support of psi inside (0, log 2)).

Input (not recomputed): data/zeross.bin (gamma_k), data/zLamb.bin (tau_k), little-endian float64.

For a fixed normalized smooth bump psi supported in [A,B] subset (0, log 2):
  c_psi^{(j)}(t) = int psi(u) u^j cos(t u + j pi/2) du
  W_psi    = (1/2) int psi(u) ( e^{u/2} - e^{-u/2}/(e^{2u}-1) ) du       (lambda-independent part)
  R_psi(1) = W_psi - sum_k c_psi(tau_k)                                   (clock defect; clock side only)
  S_psi(lambda) = sum_k [ c_psi(gamma_k + lambda*eta_k) - c_psi(tau_k) ]  (transport), eta_k = tau_k - gamma_k
  m_j(psi) = sum_k eta_k^j c_psi^{(j)}(tau_k)                             (clock-sampled moments)
and the partial sums of  S_psi(0) = sum_{j>=1} (-1)^j m_j / j!.
Outputs: output/s08_transport_console.txt, s08_moments.csv, s08_transport_lambda.csv, s08_first_moment_terms.csv
"""
import os, math
import numpy as np

K = 20000
A, B = 0.10, 0.62            # support of psi, inside (0, log 2 = 0.6931)
NU = 4001                    # quadrature nodes
JMAX = 12
HERE = os.path.dirname(os.path.abspath(__file__))
OUT = os.path.join(HERE, '..', 'output')
DATA = os.path.join(HERE, '..', 'data')
os.makedirs(OUT, exist_ok=True)

u = np.linspace(A, B, NU)
v = (u - 0.5 * (A + B)) / (0.5 * (B - A))
psi = np.zeros_like(u)
ins = np.abs(v) < 1
psi[ins] = np.exp(-1.0 / (1.0 - v[ins] ** 2))
du = u[1] - u[0]
psi /= psi.sum() * du                      # int psi = 1

def c(t, j=0, chunk=2000):
    t = np.atleast_1d(np.asarray(t, dtype=float))
    w = psi * u ** j * du
    out = np.empty_like(t)
    for a in range(0, len_{\RVM}(t), chunk):
        out[a:a + chunk] = np.cos(np.outer(t[a:a + chunk], u) + j * np.pi / 2) @ w
    return out

def main():
    lines = []
    def P(s=''):
        print(s); lines.append(s)
    g_all = np.fromfile(os.path.join(DATA, 'zeross.bin'), dtype='<f8')
    t_all = np.fromfile(os.path.join(DATA, 'zLamb.bin'), dtype='<f8')
    g = g_all[:K]; tau = t_all[:K]; eta = tau - g
    W = 0.5 * (psi * (np.exp(u / 2) - np.exp(-u / 2) / (np.exp(2 * u) - 1))).sum() * du
    c_tau = c(tau); c_gam = c(g)
    R1 = W - c_tau.sum()
    S0 = (c_gam - c_tau).sum()
    P('s08: transport tau -> gamma on the first rung')
    P(f'entries in zeross.bin / zLamb.bin: {len(g_all)} / {len(t_all)};  used K = {K}')
    P(f'gamma_1 = {g[0]:.9f}, tau_1 = {tau[0]:.9f};  gamma_K = {g[-1]:.6f}, tau_K = {tau[-1]:.6f}')
    P(f'psi: normalized bump on [{A}, {B}] (log 2 = {math.log(2):.6f}), {NU} nodes')
    P(f'max|eta_k| = {np.abs(eta).max():.6f} at k = {np.abs(eta).argmax()+1};  rms eta = {np.sqrt((eta**2).mean()):.6f}')
    for k0 in (100, 300, 1000, 3000, 10000, K - 1):
        P(f'  max_(k>={k0:5d}) |c_psi(tau_k)| = {np.abs(c_tau[k0:]).max():.3e}')
    P()
    P(f'W_psi                          = {W:+.12f}')
    P(f'sum_k c_psi(tau_k)             = {c_tau.sum():+.12f}')
    P(f'sum_k c_psi(gamma_k)           = {c_gam.sum():+.12f}')
    P(f'R_psi(1) = W - sum c(tau)      = {R1:+.12f}   (clock defect)')
    P(f'S_psi(0) = sum[c(gam)-c(tau)]  = {S0:+.12f}   (transport)')
    P(f'R_psi(0) = R_psi(1) - S_psi(0) = {R1 - S0:+.3e}')
    P()
    P(' j        m_j(psi)        (-1)^j m_j/j!      partial sum      partial - R_psi(1)')
    rows = []; part = 0.0
    for j in range(1, JMAX + 1):
        mj = (eta ** j * c(tau, j)).sum()
        term = (-1) ** j * mj / math.factorial(j)
        part += term
        rows.append((j, mj, term, part, part - R1))
        P(f'{j:2d}  {mj:+.10e}  {term:+.10e}  {part:+.12f}  {part - R1:+.3e}')
    np.savetxt(os.path.join(OUT, 's08_moments.csv'), np.array(rows), delimiter=',',
               header='j,m_j,term,partial_sum,partial_minus_R1', comments='', fmt=['%d'] + ['%.15e'] * 4)
    # which k carry the first two moments
    t1 = -eta * c(tau, 1); t2 = 0.5 * eta ** 2 * c(tau, 2)
    kk = np.arange(1, K + 1)
    np.savetxt(os.path.join(OUT, 's08_first_moment_terms.csv'), np.c_[kk, tau, eta, c_tau, c_gam - c_tau, t1, t2][:400], delimiter=',',
               header='k,tau_k,eta_k,c_psi(tau_k),c_psi(gamma_k)-c_psi(tau_k),-eta_k*c1(tau_k),0.5*eta_k^2*c2(tau_k)', comments='',
               fmt=['%d'] + ['%.12e'] * 6)
    P()
    P(' cumulative over k (first K0 ordinates):  K0   sum[c(gam)-c(tau)]   W - sum c(tau)   -m_1 part   +m_2/2 part')
    for K0 in (1, 2, 3, 5, 10, 20, 50, 100, 200, 400):
        P(f'   {K0:4d}   {(c_gam-c_tau)[:K0].sum():+.10f}   {W - c_tau[:K0].sum():+.10f}   {t1[:K0].sum():+.10f}   {t2[:K0].sum():+.10f}')
    P()
    P(' lambda     S_psi(lambda)     R_psi(lambda) = R_psi(1) - S_psi(lambda)')
    lam = np.linspace(-1.0, 2.0, 61); rows = []
    for L in lam:
        S = (c(g + L * eta) - c_tau).sum()
        rows.append((L, S, R1 - S))
    for L, S, R in rows[::5]:
        P(f'{L:+.2f}   {S:+.12f}   {R:+.12f}')
    np.savetxt(os.path.join(OUT, 's08_transport_lambda.csv'), np.array(rows), delimiter=',',
               header='lambda,S_psi,R_psi', comments='', fmt='%.15e')
    with open(os.path.join(OUT, 's08_transport_console.txt'), 'w') as f:
        f.write('\n'.join(lines) + '\n')

if __name__ == '__main__':
    main()
\end{lstlisting}

\subsection{\texttt{s09\_clock\_defect\_split.py}}
\label{app:code-s09-clock-defect-split}

\begin{lstlisting}[language=Python]
#!/usr/bin/env python3
"""s09_clock_defect_split.py
Splits the clock defect R_psi(1) of s08 (same psi) by Poisson summation on the clock:
  sum_{k>=1} c_psi(tau_k) = I_0 + 2 sum_{m>=1} (-1)^m I_m,
  I_m = int_{t_0}^infty c_psi(t) cos(2 pi m N_RVM(t)) N_RVM'(t) dt,   N_RVM(t_0) = 0,
  I_0 = int psi(u) K_0(u) du,  K_0(u) = -(1/(2 pi u)) [ sin(u t_0) log(t_0/(2 pi)) + pi/2 - Si(u t_0) ].
Reports W_psi, I_0, the smooth mismatch W_psi - I_0 and the oscillating part  -2 sum_{m>=1}(-1)^m I_m
(obtained as R_psi(1) - (W_psi - I_0), with sum_k c_psi(tau_k) read from data/zLamb.bin).
Output: output/s09_clock_defect_split_console.txt
"""
import os, math
import numpy as np
HERE = os.path.dirname(os.path.abspath(__file__))
OUT = os.path.join(HERE, '..', 'output'); DATA = os.path.join(HERE, '..', 'data')
A, B, NU, K = 0.10, 0.62, 4001, 20000
u = np.linspace(A, B, NU); v = (u - 0.5*(A+B)) / (0.5*(B-A))
psi = np.zeros_like(u); ins = np.abs(v) < 1; psi[ins] = np.exp(-1.0/(1.0 - v[ins]**2))
du = u[1]-u[0]; psi /= psi.sum()*du

def NR(t): return t/(2*np.pi)*np.log(t/(2*np.pi*np.e)) + 7.0/8.0
t0 = 9.7
for _ in range(60): t0 -= NR(t0) / (np.log(t0/(2*np.pi))/(2*np.pi))

def Si(x, n=20001):
    s = np.linspace(0.0, 1.0, n)[None, :] * np.atleast_1d(x)[:, None]
    f = np.ones_like(s); nz = s != 0; f[nz] = np.sin(s[nz]) / s[nz]
    h = np.atleast_1d(x) / (n - 1)
    return (f.sum(axis=1) - 0.5*(f[:, 0] + f[:, -1])) * h

K0 = -(1.0/(2*np.pi*u)) * (np.sin(u*t0)*np.log(t0/(2*np.pi)) + np.pi/2 - Si(u*t0))
I0 = (psi*K0).sum()*du
W = 0.5*(psi*(np.exp(u/2) - np.exp(-u/2)/(np.exp(2*u)-1))).sum()*du
tau = np.fromfile(os.path.join(DATA, 'zLamb.bin'), dtype='<f8')[:K]
w = psi*du; ctau = sum((np.cos(np.outer(tau[a:a+2000], u)) @ w).sum() for a in range(0, K, 2000))
R1 = W - ctau
L = [f's09: Poisson split of the clock defect (psi of s08 on [{A},{B}])',
     f't_0 (N_RVM(t_0)=0)                 = {t0:.12f}   (2 pi = {2*np.pi:.6f}, tau_1 = {tau[0]:.6f})',
     f'W_psi                              = {W:+.12f}',
     f'I_0  (m = 0, smooth clock density) = {I0:+.12f}',
     f'sum_k c_psi(tau_k)                 = {ctau:+.12f}',
     f'R_psi(1)                           = {R1:+.12f}',
     f'smooth mismatch  W_psi - I_0       = {W - I0:+.12f}',
     f'oscillating part -2 sum_(m>=1)(-1)^m I_m = {R1 - (W - I0):+.12f}',
     f'stationary heights 2 pi e^u, u in supp psi: [{2*np.pi*np.exp(A):.4f}, {2*np.pi*np.exp(B):.4f}];  part above t_0: u > {np.log(t0/(2*np.pi)):.6f}']
print('\n'.join(L)); open(os.path.join(OUT, 's09_clock_defect_split_console.txt'), 'w').write('\n'.join(L)+'\n')
\end{lstlisting}

\subsection{\texttt{s10\_clock\_to\_zeros.py}}
\label{app:code-s10-clock-to-zeros}

\begin{lstlisting}[language=Python]
#!/usr/bin/env python3
"""s10_clock_to_zeros.py
Solves, starting from the clock (eta = 0), the system of rung equations

   sum_k c_j(gamma_k) + sum_n a_n psi_j(log n) = W_j ,   j = 1..J,

for the unknown ordinates gamma_k (below a height) and the unknown masses a_n at the log-integers.
  psi_j(u) = w(u) cos(omega_j u), w(u) sin(omega_j u);  w = smooth bump on (0, U), U = log(N + 1/2)
  c_j(t)   = int psi_j(u) cos(t u) du
  W_j      = int psi_j(u) (1/2)( e^{u/2} - e^{-u/2}/(e^{2u}-1) ) du      (lambda-independent part)
Ordinates above the free range are kept on the clock (tau_k from data/zLamb.bin).
The masses enter linearly and are eliminated exactly at every step (ridge RIDGE).
Continuation in height: the band omega <= Omega is enlarged step by step (OMEGAS); at each stage the
ordinates with tau_k <= Omega + MARGIN are free.  Start: gamma_k = tau_k.  Method: Levenberg-Marquardt.
data/zeross.bin is used only AFTER the iteration, to compare.
Outputs: output/s10_clock_to_zeros_console.txt, s10_ordinates.csv, s10_masses.csv
"""
import os, math
import numpy as np
HERE = os.path.dirname(os.path.abspath(__file__))
OUT = os.path.join(HERE, '..', 'output'); DATA = os.path.join(HERE, '..', 'data')
N = 13; U = math.log(N + 0.5)
KT = 6000; DOM = 0.1; NU = 8001; ITER = 400
OMEGAS = [18, 22, 26, 30, 34, 38, 42, 46, 50]; MARGIN = 12.0
WMIN = 5e-2; STEP = 0.2; RIDGE = 1e-10

u = np.linspace(0.0, U, NU)[1:-1]; du = U / (NU - 1)
v = (u - U/2) / (U/2); w = np.exp(-1.0/(1.0 - v**2)); w /= w.sum()*du
Wfun = 0.5*(np.exp(u/2) - np.exp(-u/2)/(np.exp(2*u) - 1))
ns = np.array([n for n in range(2, N + 1) if np.interp(math.log(n), u, w)/w.max() >= WMIN])
tau = np.fromfile(os.path.join(DATA, 'zLamb.bin'), dtype='<f8')[:KT]
COS_TAU = np.zeros_like(u)                       # sum over all clock ordinates of cos(tau_k u)
for a in range(0, KT, 500): COS_TAU += np.cos(np.outer(u, tau[a:a+500])).sum(axis=1)

lines = []
def P(s=''):
    print(s); lines.append(s)
P('s10: from the clock to the ordinates through the rung equations (continuation in height)')
P(f'N = {N}, U = log({N}+1/2) = {U:.6f}; free masses n = {ns[0]}..{ns[-1]}; start gamma_k = tau_k')
g_all = tau.copy()
for OM in OMEGAS:
    om = np.arange(0.0, OM + 1e-9, DOM)
    Psi = np.vstack([w*np.cos(o*u) for o in om] + [w*np.sin(o*u) for o in om[1:]]) * du
    Wj = Psi @ Wfun
    Pn = np.vstack([np.interp(np.log(ns), u, Psi[j]/du) for j in range(Psi.shape[0])])
    Q = np.linalg.solve(Pn.T @ Pn + RIDGE*np.eye(len(ns)), Pn.T)
    K = int((tau <= OM + MARGIN).sum())
    g = g_all[:K].copy()
    fixed = Psi @ (COS_TAU - np.cos(np.outer(u, tau[:K])).sum(axis=1)) - Wj
    def r0(g): return Psi @ np.cos(np.outer(u, g)).sum(axis=1) + fixed
    def res(g):
        x = r0(g); a = -Q @ x; return x + Pn @ a, a
    r, a = res(g); n0 = np.linalg.norm(r); damp = 1e-3
    for it in range(1, ITER + 1):
        Jg = Psi @ (-u[:, None]*np.sin(np.outer(u, g))); Jg = Jg - Pn @ (Q @ Jg)
        A = Jg.T @ Jg; b = -Jg.T @ r
        while True:
            d = np.linalg.solve(A + damp*np.diag(np.diag(A) + 1e-14), b)
            d *= min(1.0, STEP/max(np.abs(d).max(), 1e-300))
            r2, a2 = res(g + d)
            if np.linalg.norm(r2) < np.linalg.norm(r) or damp > 1e12: break
            damp *= 4
        g = g + d; r, a = r2, a2; damp = max(damp/3, 1e-14)
        if np.abs(d).max() < 1e-11: break
    g_all[:K] = g
    P(f'band omega <= {OM:2d}: {K:2d} free ordinates, {Psi.shape[0]} tests, |residual| {n0:.3e} -> {np.linalg.norm(r):.3e} in {it} iterations')
P()
gam = np.fromfile(os.path.join(DATA, 'zeross.bin'), dtype='<f8')[:K]
P('  k      tau_k (start)    solved ordinate    gamma_k (zeross.bin)   solved - gamma_k    tau_k - gamma_k')
for k in range(K):
    P(f'{k+1:3d}   {tau[k]:14.9f}   {g_all[k]:14.9f}   {gam[k]:14.9f}   {g_all[k]-gam[k]:+.3e}   {tau[k]-gam[k]:+.6f}')
P()
def Lam(n):
    for p in range(2, n + 1):
        if n % p == 0:
            m = n
            while m % p == 0: m //= p
            return math.log(p) if m == 1 else 0.0
    return 0.0
P('  n     solved mass a_n     Lambda(n)/(2 sqrt n)     difference     w(log n)/max w')
for i, n in enumerate(ns):
    ex = Lam(int(n))/(2*math.sqrt(n))
    P(f'{n:3d}   {a[i]:+.9f}   {ex:+.9f}   {a[i]-ex:+.3e}   {np.interp(math.log(n), u, w)/w.max():.3e}')
np.savetxt(os.path.join(OUT, 's10_ordinates.csv'), np.c_[np.arange(1, K+1), tau[:K], g_all[:K], gam], delimiter=',',
           header='k,tau_k,solved,gamma_k', comments='', fmt=['%d'] + ['%.12f']*3)
np.savetxt(os.path.join(OUT, 's10_masses.csv'), np.c_[ns, a, [Lam(int(n))/(2*math.sqrt(n)) for n in ns]], delimiter=',',
           header='n,solved_mass,Lambda(n)/(2sqrt(n))', comments='', fmt=['%d'] + ['%.12f']*2)
open(os.path.join(OUT, 's10_clock_to_zeros_console.txt'), 'w').write('\n'.join(lines) + '\n')
\end{lstlisting}


\subsection{\texttt{s11\_rvm\_z1\_difference\_analysis.py}}
\label{app:code-s11-rvm-z1}
This supplemental diagnostic reproduces the RVM--$Z1$ difference, shell and convergence calculations.  It uses the genuine ordinates only for post-construction diagnostics and is not used in any proof.
\lstinputlisting[language=Python]{code/s11_rvm_z1_difference_analysis.py}

\subsection{\texttt{s12\_multiscale\_jacobian.py}}
\label{app:code-s12-multiscale}
This program produces Experiment~\CGref{exp:multiscale-jacobian} and Figures~\CGref{fig:multiscale-gap} and its companion fixed-window diagnostic.  It generates the clock from the Lambert formula and does not read the genuine ordinates.
\lstinputlisting[language=Python]{code/s12_multiscale_jacobian.py}

\subsection{\texttt{s13\_nested\_bounded\_endpoint.py}}
\label{app:code-s13-bounded}
This program produces the bounded and nested endpoint diagnostics of Appendix~\CGref{sec:appendix-endpoint-v18}.  The solver uses the midpoint clock; genuine ordinates and von Mangoldt masses, when displayed, are read only after the solve as diagnostics.
\lstinputlisting[language=Python]{code/s13_nested_bounded_endpoint.py}

\section{Console output of the programs}
\label{app:output}

The complete console output of each program, as written to \texttt{output/*\_console.txt}.

{\sloppy The post-v16 numerical programs that produce the multiscale and bounded-endpoint figures (files \texttt{fig14}--\texttt{fig17}) and the bounded/nested endpoint tables are also distributed with their machine-readable outputs in the repository.  In particular, \texttt{s12\_\allowbreak multiscale\_\allowbreak jacobian.py} writes \texttt{output/\allowbreak s12\_\allowbreak multiscale\_\allowbreak jacobian\_\allowbreak console.txt}, \texttt{output/\allowbreak s12\_\allowbreak multiscale\_\allowbreak summary.csv} and \texttt{output/\allowbreak s12\_\allowbreak fixed\_\allowbreak window\_\allowbreak pairs.csv}; \texttt{s13\_\allowbreak nested\_\allowbreak bounded\_\allowbreak endpoint.py} writes \texttt{output/\allowbreak s13\_\allowbreak bounded\_\allowbreak mass\_\allowbreak summary.csv}, \texttt{output/\allowbreak s13\_\allowbreak bounded\_\allowbreak mass\_\allowbreak ordinates.csv}, \texttt{output/\allowbreak s13\_\allowbreak bounded\_\allowbreak mass\_\allowbreak masses.csv}, \texttt{output/\allowbreak s13\_\allowbreak cross\_\allowbreak validation.csv} and \texttt{output/\allowbreak s13\_\allowbreak nested\_\allowbreak metric\_\allowbreak summary.csv}.  Here ``normalized residual'' means the Euclidean residual divided by the square root of the number of test equations.  The bump and permanent-block conventions are specified in the corresponding code and in Appendix~\CGref{sec:appendix-endpoint-v18}.\par}

\subsection{\texttt{s01\_clock\_and\_family}}
\label{app:out-s01-clock-and-family}

\begin{lstlisting}[style=out]
zeros loaded: N = 2001052, gamma_1 = 14.134725142, gamma_N = 1132490.658714, tau_1 = 14.521346953, tau_N = 1132490.722588
max |N_RVM(tau_k) - (k - 1/2)| = 6.985e-10
eta_k = tau_k - gamma_k:  max |eta| = 0.993711 (k = 34),  mean |eta| = 0.151209,  rms = 0.188066,  mean eta = 1.731e-07
      k from         k to     max|eta|      rms eta
           1           10     0.840022     0.545768
          11          100     0.993711     0.453166
         101         1000     0.915450     0.345177
        1001        10000     0.833881     0.274177
       10001       100000     0.785835     0.226352
      100001      1000000     0.762816     0.192411
     1000001      2001052     0.717359     0.178900
first ordinates:
   k          gamma_k            tau_k        eta_k
   1     14.134725142     14.521346953     0.386622
   2     21.022039639     20.655740356    -0.366299
   3     25.010857580     25.492675432     0.481818
   4     30.424876126     29.739411632    -0.685464
   5     32.935061588     33.624531889     0.689470
   6     37.586178159     37.257370087    -0.328808
   7     40.918719012     40.700551604    -0.218167
   8     43.327073281     43.994013953     0.666941
   9     48.005150881     47.165128427    -0.840022
  10     49.773832478     50.233653391     0.459821
  11     52.970321478     53.214422874     0.244101
  12     56.446247697     56.118921948    -0.327326
sigma = 2.00, t in [1,2000] (499 points): max |zeta_Z1/zeta - 1| = 8.706e-10, median = 2.685e-10
sigma = 1.25, t in [1,2000] (499 points): max |zeta_Z1/zeta - 1| = 1.128e-09, median = 3.693e-10
sigma = 0.75, t in [1,2000] (499 points): max |zeta_Z1/zeta - 1| = 2.629e-09, median = 5.179e-10
critical line, t in [0.5,60]: max |Z_Z1(t) - Z(t)| / max(1,|Z(t)|) = 8.370e-10
total run time 30.6 s
\end{lstlisting}

\subsection{\texttt{s01b\_clock\_zoom}}
\label{app:out-s01b-clock-zoom}

\begin{lstlisting}[style=out]
zoom table near the top of the file:
        k            gamma_k              tau_k      eta_k   N*-N_RVM
  2001041     1132485.303923     1132485.011569  -0.292355  -0.563104
  2001042     1132485.662638     1132485.530752  -0.131886  -0.254025
  2001043     1132486.244063     1132486.049936  -0.194127  -0.373908
  2001044     1132486.392256     1132486.569120   0.176863   0.340656
  2001045     1132486.874881     1132487.088303   0.213422   0.411072
  2001046     1132487.571868     1132487.607487   0.035619   0.068605
  2001047     1132487.950694     1132488.126670   0.175977   0.338949
  2001048     1132488.905328     1132488.645854  -0.259474  -0.499774
  2001049     1132489.162780     1132489.165037   0.002257   0.004348
  2001050     1132489.585215     1132489.684221   0.099006   0.190696
  2001051     1132490.165305     1132490.203404   0.038099   0.073383
  2001052     1132490.658714     1132490.722588   0.063873   0.123026
total run time 0.6 s
\end{lstlisting}

\subsection{\texttt{s02\_transition}}
\label{app:out-s02-transition}

\begin{lstlisting}[style=out]
varpi_RVM (Richardson, sigma = 2, T = 2500, 5000, 10000) = 0.9619415807 -1.28e-10 i
L = 1/varpi_RVM = 1.0395641690
sigma' = 0.25: |zeta_RVM/(varpi(1+w)) - 1|: log-log slope on t>100 = -1.181, max t*err = 0.0617
sigma' = 0.25: |zeta_RVM/(varpi(1+chi)) - 1|: max t*err (t>100) = 0.0790
          |w| * t^sigma' in [1.5832, 1.5832];  |1/zeta_RVM - L| at t = 1e4: 1.425e-01
sigma' = 0.50: |zeta_RVM/(varpi(1+w)) - 1|: log-log slope on t>100 = -1.181, max t*err = 0.0618
          |w| * t^sigma' in [2.5061, 2.5066];  |1/zeta_RVM - L| at t = 1e4: 2.543e-02
sigma' = 1.50: |zeta_RVM/(varpi(1+w)) - 1|: log-log slope on t>100 = -1.182, max t*err = 0.0620
          |w| * t^sigma' in [15.6619, 15.7496];  |1/zeta_RVM - L| at t = 1e4: 1.598e-05
total run time 11.5 s
\end{lstlisting}

\subsection{\texttt{s03\_exponential\_sums}}
\label{app:out-s03-exponential-sums}

\begin{lstlisting}[style=out]
     x     max|S_RVM| max|S_RVM|/log    |S_zeta(N)| Landau -T/2pi*L/sqrt(x)   |M_H^RVM(x)|
   2.0         3.5745         0.2695       88300.77        -88341.59      1.722e-06
   3.0         3.9750         0.3759      114267.92       -114324.31      1.175e-06
   4.0         3.9523         0.3732       62435.29        -62466.94      8.004e-07
   6.0         4.5264         0.4547           0.26            -0.00      2.262e-06
  10.0         9.6989         1.1495           2.42            -0.00      4.031e-06
total run time 1.7 s
\end{lstlisting}

\subsection{\texttt{s04\_primes}}
\label{app:out-s04-primes}

\begin{lstlisting}[style=out]
x <= 100, all 2001052 zeros of each list, 3000 grid points:
  max |pi_from_zeta_zeros - pi| away from primes (distance > 0.2): 0.4035
  pi_RVM(100) = 26.0782, pi(100) = 25, li(100) = 30.1261, J_RVM(100) = 29.5889, J(100) = 28.5333
  pi_RVM monotone on [2,100]: False (min increment on the grid -2.528e-01)
x <= 1000, smooth cut at K = 17061 zeros (tau_K = 15707.6), 40000 grid points:
  pi_RVM(1000) = 169.9750, pi(1000) = 168, li(1000) = 177.6097
cell masses c_n = pi_RVM(n) - pi_RVM(n-1):
    n        c_n    1/log n    prime
    2    0.52944    1.44270        1
    3    1.04188    0.91024        1
    4    0.53593    0.72135        0
    5    0.46739    0.62133        1
    6    0.36089    0.55811        0
    7    0.63254    0.51390        1
    8    0.62393    0.48090        0
    9   -0.37669    0.45512        0
   10    0.17151    0.43429        0
   11    0.69779    0.41703        1
   12    0.32703    0.40243        0
   13    0.36227    0.38987        1
   14    0.57662    0.37892        0
   15    0.19671    0.36927        0
   16   -0.05284    0.36067        0
   17    0.39958    0.35296        1
   18    0.76865    0.34598        0
   19    0.43500    0.33962        1
   20   -0.03342    0.33381        0
   21    0.03009    0.32846        0
   22    0.37466    0.32352        0
   23    0.40277    0.31893        1
   24    0.11768    0.31466        0
   25    0.01006    0.31067        0
   26    0.27350    0.30693        0
   27    0.55508    0.30341        0
   28    0.47627    0.30010        0
   29    0.12577    0.29697        1
   30   -0.08367    0.29401        0
   31    0.07329    0.29121        1
   32    0.38040    0.28854        0
   33    0.43927    0.28600        0
   34    0.34619    0.28358        0
   35    0.06419    0.28127        0
   36   -0.02475    0.27906        0
   37    0.15448    0.27694        1
   38    0.40992    0.27491        0
   39    0.49628    0.27296        0
   40    0.35226    0.27109        0
   41    0.14030    0.26928        1
sum_{n<=1000} c_n = 169.9750,  pi(1000) = 168
real RVM prime levels p_j^RVM (first crossing of level j) versus p_j:
    j      p_j^RVM      p_j       diff
    1       2.2139        2     0.2139
    2       3.1550        3     0.1550
    3       5.1887        5     0.1887
    4       7.1593        7     0.1593
    5      11.1406       11     0.1406
    6      14.0208       13     1.0208
    7      17.2552       17     0.2552
    8      19.1781       19     0.1781
    9      26.0536       23     3.0536
   10      28.0447       29    -0.9553
   11      33.0707       31     2.0707
   12      38.0931       37     1.0931
   13      40.2282       41    -0.7718
   14      44.1861       43     1.1861
   15      46.1221       47    -0.8779
   16      57.1779       53     4.1779
   17      59.1201       59     0.1201
   18      61.0809       61     0.0809
   19      67.0384       67     0.0384
   20      74.0931       71     3.0931
   21      76.1235       73     3.1235
   22      83.0822       79     4.0822
   23      86.0513       83     3.0513
   24      92.0825       89     3.0825
   25      94.0781       97    -2.9219
   26      96.1319      101    -4.8681
   27     103.2219      103     0.2219
   28     112.1163      107     5.1163
   29     115.0277      109     6.0277
   30     118.1431      113     5.1431
total run time 107.7 s
\end{lstlisting}

\subsection{\texttt{s05\_mertens}}
\label{app:out-s05-mertens}

\begin{lstlisting}[style=out]
zeta_RVM'(rho_k) exact for k <= 1000; zeta_RVM'(rho_k)/log(tau_k/2pi) at k = 1000: 0.9619388429 -1.41e-05 i
varpi_RVM = 0.9619415807; fit zeta_RVM'(rho_k) = varpi log(tau_k/2pi)(1 + c/tau_k), c = -0.0052 -0.0208 i
Riemann-Siegel zeta'(rho) versus mpmath at k = 1001..1050: max relative difference 1.19e-05
median |zeta'(rho)| = 7.8899;  spectral sums 2 sum|a|^2: zeta 0.029034, RVM 0.027203
M(10^6) = 212;  squarefree density Q(10^6)/10^6 = 0.607926 (6/pi^2 = 0.607927)
explicit formula with 2001052 zeros of zeta on [2,2000]: rms(M_zeta - M) = 0.0005 off integers (|x-n|>0.05)
target on 49896 grid points; at 100 integers <= 1e5: max |M_RVM(n) (all zeros) - target(n)| = 0.0191
ternary walk: max |M_h - target| on [1,1e6] = 2.6155, on [1e4,1e6] = 1.6539
densities of (-1, 0, +1):  mu   0.303857 0.392074 0.304069
                           mu_h 0.304003 0.392074 0.303923   (3/pi^2 = 0.303964, 1 - 6/pi^2 = 0.392073)
sign agreement mu_h = mu on squarefrees: 0.500344;  sum mu mu_h / Q = 0.000688
max |M(n)|/sqrt(n) on [10,1e6] = 0.8321;  max |M_RVM(n)|/sqrt(n) = 0.7541;  max |M_h(n)|/sqrt(n) = 0.6882

NUMERICAL EXAMPLE OF THE TERNARY DIVISION (first 40 integers)
   n   mu   mu^2     target M_h(n-1)   sign   mu_h    M_h      M
   1    1      1     0.2925        0     +1      1      1      1
   2   -1      1     0.4563        1     -1     -1      0      0
   3   -1      1    -0.6523        0     -1     -1     -1     -1
   4    0      0    -0.8720       -1      -      0     -1     -1
   5   -1      1    -1.7255       -1     -1     -1     -2     -2
   6    1      1    -1.3063       -2     +1      1     -1     -1
   7   -1      1    -1.5931       -1     -1     -1     -2     -2
   8    0      0    -2.4029       -2      -      0     -2     -2
   9    0      0    -1.3292       -2      -      0     -2     -2
  10    1      1    -1.3280       -2     +1      1     -1     -1
  11   -1      1    -2.3272       -1     -1     -1     -2     -2
  12    0      0    -2.2067       -2      -      0     -2     -2
  13   -1      1    -1.8468       -2     +1      1     -1     -3
  14    1      1    -1.9324       -1     -1     -1     -2     -2
  15    1      1    -1.5435       -2     +1      1     -1     -1
  16    0      0    -1.0414       -1      -      0     -1     -1
  17   -1      1    -1.6970       -1     -1     -1     -2     -2
  18    0      0    -3.0095       -2      -      0     -2     -2
  19   -1      1    -3.2869       -2     -1     -1     -3     -3
  20    0      0    -2.2604       -3      -      0     -3     -3
  21    1      1    -1.2515       -3     +1      1     -2     -2
  22    1      1    -1.1120       -2     +1      1     -1     -1
  23   -1      1    -1.3607       -1     -1     -1     -2     -2
  24    0      0    -1.3669       -2      -      0     -2     -2
  25    0      0    -1.3810       -2      -      0     -2     -2
  26    1      1    -1.9805       -2     +1      1     -1     -1
  27    0      0    -3.0141       -1      -      0     -1     -1
  28    0      0    -3.6155       -1      -      0     -1     -1
  29   -1      1    -3.2066       -1     -1     -1     -2     -2
  30   -1      1    -2.1599       -2     -1     -1     -3     -3
  31   -1      1    -1.3649       -3     +1      1     -2     -4
  32    0      0    -1.2973       -2      -      0     -2     -4
  33    1      1    -1.6435       -2     +1      1     -1     -3
  34    1      1    -1.7783       -1     -1     -1     -2     -2
  35    1      1    -1.4711       -2     +1      1     -1     -1
  36    0      0    -1.0723       -1      -      0     -1     -1
  37   -1      1    -1.0774       -1     -1     -1     -2     -2
  38    1      1    -1.5897       -2     +1      1     -1     -1
  39    1      1    -2.2227       -1     -1     -1     -2      0
  40    0      0    -2.4881       -2      -      0     -2      0

ENERGY E_N = (1/N) sum_{n<=N} W(n)^2
         N         E_N(M)     E_N(M_RVM)       E_N(M_h)
      1000        14.6400        18.6280        18.5960
     10000       163.6747       195.6965       196.0683
    100000      1622.3189      1308.3041      1308.8186
   1000000     14851.0857     12537.9008     12538.1174
total run time 148.2 s
\end{lstlisting}

\subsection{\texttt{s06\_prime\_ratios}}
\label{app:out-s06-prime-ratios}

\begin{lstlisting}[style=out]
grid: 40000 points on [2, 1000]; pi_RVM from the first 17061 ordinates tau_k (smooth cut)

ratios at selected x
       x           pi       pi_RVM           li    li/pi_RVM        pi/li    pi/pi_RVM
    10.0       4.0000       3.9876       6.1657     1.546218     0.648747     1.003105
    20.0       8.0000       7.6636       9.9052     1.292496     0.807657     1.043893
    50.0      15.0000      14.7686      18.4695     1.250592     0.812151     1.015670
   100.0      25.0000      26.0813      30.1265     1.155098     0.829835     0.958541
   200.0      46.0000      43.3656      50.1910     1.157393     0.916499     1.060749
   500.0      95.0000      93.7489     101.7975     1.085854     0.933225     1.013346
  1000.0     168.0000     169.9750     177.6097     1.044916     0.945895     0.988381

square-root scale B(x) = sqrt(x) (log log x)^2 / log x on [8, 1000]
max |pi_RVM - li| / B = 2.7213   (at x = 9.73)
max |pi - li| / B     = 2.4490   (at x = 11.00)
max |pi_RVM - li| / B on [50, 1000] = 1.3455
max |pi_RVM - li| / B on [500, 1000] = 0.7216
total run time 1.2 s
\end{lstlisting}

\subsection{\texttt{s07\_product\_identity}}
\label{app:out-s07-product-identity}

\begin{lstlisting}[style=out]
N = 2001052 ordinates; varpi_RVM = 0.9619415807
         K              P_K    P_K - varpi
        10     0.9625640008      6.224e-04
       100     0.9624787401      5.372e-04
      1000     0.9622891707      3.476e-04
     10000     0.9619513787      9.798e-06
    100000     0.9619301865     -1.139e-05
   1000000     0.9619418189      2.382e-07
   2001052     0.9619415390     -4.166e-08
K >=   10000: min P_K = 0.9618023968, max P_K = 0.9620641927, mean = 0.9619418623
K >=  100000: min P_K = 0.9619210208, max P_K = 0.9619605076, mean = 0.9619417591
K >= 1000000: min P_K = 0.9619389488, max P_K = 0.9619439825, mean = 0.9619416633

flow along the lambda-family
Upsilon(0.00) = 0.0436949578
Upsilon(0.25) = 0.0412445375
Upsilon(0.50) = 0.0387987180
Upsilon(0.75) = 0.0363572218
Upsilon(1.00) = 0.0339197739
sum_k 2 eta_k^2 / gamma_k^2 = 0.0098210688
integral_0^1 Upsilon = sum_k log((1/4+tau_k^2)/(1/4+gamma_k^2)) = 0.0388016004
log L = log(1/varpi_RVM)                               = 0.0388015571
total run time 1.5 s
\end{lstlisting}

\subsection{\texttt{s08\_transport\_moments}}
\label{app:out-s08-transport-moments}

\begin{lstlisting}[style=out]
s08: transport tau -> gamma on the first rung
entries in zeross.bin / zLamb.bin: 2001052 / 2001052;  used K = 20000
gamma_1 = 14.134725142, tau_1 = 14.521346953;  gamma_K = 18046.464296, tau_K = 18046.544665
psi: normalized bump on [0.1, 0.62] (log 2 = 0.693147), 4001 nodes
max|eta_k| = 0.993711 at k = 34;  rms eta = 0.266395
  max_(k>=  100) |c_psi(tau_k)| = 6.880e-05
  max_(k>=  300) |c_psi(tau_k)| = 7.263e-07
  max_(k>= 1000) |c_psi(tau_k)| = 2.527e-10
  max_(k>= 3000) |c_psi(tau_k)| = 2.074e-13
  max_(k>=10000) |c_psi(tau_k)| = 2.074e-13
  max_(k>=19999) |c_psi(tau_k)| = 1.616e-15

W_psi                          = +0.133643530916
sum_k c_psi(tau_k)             = +0.176106813076
sum_k c_psi(gamma_k)           = +0.133643530929
R_psi(1) = W - sum c(tau)      = -0.042463282161   (clock defect)
S_psi(0) = sum[c(gam)-c(tau)]  = -0.042463282147   (transport)
R_psi(0) = R_psi(1) - S_psi(0) = -1.376e-11

 j        m_j(psi)        (-1)^j m_j/j!      partial sum      partial - R_psi(1)
 1  +4.0395231784e-02  -4.0395231784e-02  -0.040395231784  +2.068e-03
 2  -4.8097910277e-03  -2.4048955139e-03  -0.042800127298  -3.368e-04
 3  -2.0431952768e-03  +3.4053254613e-04  -0.042459594751  +3.687e-06
 4  -5.5184329118e-05  -2.2993470466e-06  -0.042461894098  +1.388e-06
 5  +1.6734348340e-04  -1.3945290283e-06  -0.042463288628  -6.467e-09
 6  +5.5720624242e-06  +7.7389755891e-09  -0.042463280889  +1.272e-09
 7  +7.2607702259e-06  -1.4406290131e-09  -0.042463282329  -1.686e-10
 8  +6.4039357525e-06  +1.5882777164e-10  -0.042463282170  -9.773e-12
 9  -8.6516463634e-06  +2.3841618065e-11  -0.042463282147  +1.407e-11
10  -2.6791658297e-06  -7.3830628022e-13  -0.042463282147  +1.333e-11
11  +2.8798049874e-06  -7.2145186673e-14  -0.042463282147  +1.326e-11
12  +8.0540801176e-07  +1.6814307338e-15  -0.042463282147  +1.326e-11

 cumulative over k (first K0 ordinates):  K0   sum[c(gam)-c(tau)]   W - sum c(tau)   -m_1 part   +m_2/2 part
      1   -0.0208609181   +0.0155048033   -0.0168182977   -0.0039851021
      2   -0.0179065615   +0.0336283174   -0.0153170893   -0.0024894395
      3   -0.0230428015   -0.0613123463   -0.0185197872   -0.0044492860
      5   -0.0438112092   -0.0723587090   -0.0422769516   -0.0020012704
     10   -0.0464863677   -0.0403092255   -0.0440974925   -0.0029001855
     20   -0.0406060545   -0.0470689184   -0.0387914922   -0.0021299305
     50   -0.0422141654   -0.0432549118   -0.0401314609   -0.0024151466
    100   -0.0424897738   -0.0426540294   -0.0404216192   -0.0024058875
    200   -0.0424647471   -0.0424770847   -0.0403966024   -0.0024050814
    400   -0.0424633643   -0.0424636133   -0.0403953106   -0.0024049002

 lambda     S_psi(lambda)     R_psi(lambda) = R_psi(1) - S_psi(lambda)
-1.00   -0.087766836871   +0.045303554710
-0.75   -0.076275893121   +0.033812610960
-0.50   -0.064876726912   +0.022413444752
-0.25   -0.053596431977   +0.011133149816
+0.00   -0.042463282147   -0.000000000014
+0.25   -0.031506572650   -0.010956709511
+0.50   -0.020756460380   -0.021706821780
+0.75   -0.010243803437   -0.032219478724
+1.00   +0.000000000000   -0.042463282161
+1.25   +0.009943173538   -0.052406455698
+1.50   +0.019553725448   -0.062017007609
+1.75   +0.028799612928   -0.071262895088
+2.00   +0.037648908219   -0.080112190380
\end{lstlisting}

\subsection{\texttt{s09\_clock\_defect\_split}}
\label{app:out-s09-clock-defect-split}

\begin{lstlisting}[style=out]
s09: Poisson split of the clock defect (psi of s08 on [0.1,0.62])
t_0 (N_RVM(t_0)=0)                 = 9.676906787166   (2 pi = 6.283185, tau_1 = 14.521347)
W_psi                              = +0.133643530916
I_0  (m = 0, smooth clock density) = +0.063533498702
sum_k c_psi(tau_k)                 = +0.176106813076
R_psi(1)                           = -0.042463282161
smooth mismatch  W_psi - I_0       = +0.070110032214
oscillating part -2 sum_(m>=1)(-1)^m I_m = -0.112573314375
stationary heights 2 pi e^u, u in supp psi: [6.9440, 11.6800];  part above t_0: u > 0.431865
\end{lstlisting}

\subsection{\texttt{s10\_clock\_to\_zeros}}
\label{app:out-s10-clock-to-zeros}

\begin{lstlisting}[style=out]
s10: from the clock to the ordinates through the rung equations (continuation in height)
N = 13, U = log(13+1/2) = 2.602690; free masses n = 2..11; start gamma_k = tau_k
band omega <= 18:  4 free ordinates, 361 tests, |residual| 8.136e-01 -> 3.892e-04 in 37 iterations
band omega <= 22:  5 free ordinates, 441 tests, |residual| 1.824e-02 -> 1.520e-03 in 20 iterations
band omega <= 26:  6 free ordinates, 521 tests, |residual| 1.985e-02 -> 8.002e-04 in 25 iterations
band omega <= 30:  7 free ordinates, 601 tests, |residual| 2.363e-02 -> 2.506e-03 in 21 iterations
band omega <= 34:  8 free ordinates, 681 tests, |residual| 4.635e-02 -> 2.827e-03 in 47 iterations
band omega <= 38:  9 free ordinates, 761 tests, |residual| 3.868e-02 -> 1.683e-03 in 11 iterations
band omega <= 42: 11 free ordinates, 841 tests, |residual| 1.881e-01 -> 1.595e-03 in 11 iterations
band omega <= 46: 12 free ordinates, 921 tests, |residual| 1.878e-02 -> 4.591e-03 in 12 iterations
band omega <= 50: 14 free ordinates, 1001 tests, |residual| 5.189e-02 -> 2.885e-03 in 22 iterations

  k      tau_k (start)    solved ordinate    gamma_k (zeross.bin)   solved - gamma_k    tau_k - gamma_k
  1     14.521346953     14.134660056     14.134725142   -6.509e-05   +0.386622
  2     20.655740356     21.022091631     21.022039639   +5.199e-05   -0.366299
  3     25.492675432     25.010827129     25.010857580   -3.045e-05   +0.481818
  4     29.739411632     30.424917668     30.424876126   +4.154e-05   -0.685464
  5     33.624531889     32.935072953     32.935061588   +1.137e-05   +0.689470
  6     37.257370087     37.586146548     37.586178159   -3.161e-05   -0.328808
  7     40.700551604     40.918919600     40.918719012   +2.006e-04   -0.218167
  8     43.994013953     43.327322117     43.327073281   +2.488e-04   +0.666941
  9     47.165128427     48.004716753     48.005150881   -4.341e-04   -0.840022
 10     50.233653391     49.773657726     49.773832478   -1.748e-04   +0.459821
 11     53.214422874     52.972535790     52.970321478   +2.214e-03   +0.244101
 12     56.118921948     56.548052093     56.446247697   +1.018e-01   -0.327326
 13     58.956266008     58.576523603     59.347044003   -7.705e-01   -0.390778
 14     61.733838843     62.561880181     60.831778525   +1.730e+00   +0.902060

  n     solved mass a_n     Lambda(n)/(2 sqrt n)     difference     w(log n)/max w
  2   +0.245076213   +0.245064536   +1.168e-05   7.562e-01
  3   +0.317145240   +0.317142050   +3.190e-06   9.754e-01
  4   +0.173268937   +0.173286795   -1.786e-05   9.957e-01
  5   +0.359884977   +0.359881258   +3.719e-06   9.423e-01
  6   -0.000004422   +0.000000000   -4.422e-06   8.474e-01
  7   +0.367785301   +0.367742452   +4.285e-05   7.225e-01
  8   +0.122615622   +0.122532268   +8.335e-05   5.733e-01
  9   +0.183143344   +0.183102048   +4.130e-05   4.062e-01
 10   -0.000286111   +0.000000000   -2.861e-04   2.344e-01
 11   +0.362276961   +0.361496314   +7.806e-04   8.642e-02
\end{lstlisting}

\begin{thebibliography}{99}

\bibitem{Ahlfors1979}
L.~V. Ahlfors, \emph{Complex Analysis}, 3rd ed., McGraw--Hill, New York, 1979.

\bibitem{AriasdeReyna2024}
J.~Arias de Reyna, \emph{Explicit formula and quasicrystal definition}, arXiv:2402.10604 (2024, v2 2025), \href{https://doi.org/10.48550/arXiv.2402.10604}{doi:10.48550/arXiv.2402.10604}.

\bibitem{Avdonin1974}
S.~A. Avdonin, \emph{On the question of Riesz bases of exponential functions in $L^2$}, Vestnik Leningrad. Univ. Mat. Mekh. Astronom. (1974), no.~13, 5--12. (No DOI.)

\bibitem{AvdoninIvanov2002}
S.~A. Avdonin and S.~A. Ivanov, \emph{Exponential Riesz bases of subspaces and divided differences}, St. Petersburg Math. J. \textbf{13} (2002), no.~3, 339--351; arXiv:math/0103160, \href{https://doi.org/10.48550/arXiv.math/0103160}{doi:10.48550/arXiv.math/0103160}.

\bibitem{BatemanDiamond1969}
P.~T. Bateman and H.~G. Diamond, Asymptotic distribution of Beurling's generalized prime numbers, in: \emph{Studies in Number Theory}, MAA Studies in Mathematics~6, Mathematical Association of America, 1969.

\bibitem{BenedettoPowellYilmaz2006}
J.~J. Benedetto, A.~M. Powell and \"O.~Y{\i}lmaz, Sigma-delta ($\Sigma\Delta$) quantization and finite frames, \emph{IEEE Trans. Inform. Theory} \textbf{52} (2006), no.~5, 1990--2005.
\href{https://doi.org/10.1109/TIT.2006.872849}{doi:10.1109/TIT.2006.872849}.

\bibitem{Fujii1974}
A.~Fujii, \emph{On the zeros of Dirichlet $L$-functions. I}, Trans. Amer. Math. Soc. \textbf{196} (1974), 225--235, \href{https://doi.org/10.1090/S0002-9947-1974-0349603-2}{doi:10.1090/S0002-9947-1974-0349603-2}.

\bibitem{Fujii1975}
A.~Fujii, \emph{On the distribution of the zeros of the Riemann zeta function in short intervals}, Bull. Amer. Math. Soc. \textbf{81} (1975), 139--142, \href{https://doi.org/10.1090/S0002-9904-1975-13674-3}{doi:10.1090/S0002-9904-1975-13674-3}.

\bibitem{Fujii1975b}
A.~Fujii, \emph{On the difference between $r$ consecutive ordinates of the zeros of the Riemann zeta function}, Proc. Japan Acad. \textbf{51} (1975), no.~10, 741--743, \href{https://doi.org/10.3792/pja/1195518466}{doi:10.3792/pja/1195518466}.

\bibitem{Garnett2007}
J.~B. Garnett, \emph{Bounded Analytic Functions}, revised first ed., Graduate Texts in Mathematics 236, Springer, New York, 2007, \href{https://doi.org/10.1007/0-387-49763-3}{doi:10.1007/0-387-49763-3}.

\bibitem{Gonek2012}
S.~M. Gonek, \emph{Finite Euler products and the Riemann hypothesis}, Trans. Amer. Math. Soc. \textbf{364} (2012), 2157--2191, \href{https://doi.org/10.1090/S0002-9947-2011-05546-7}{doi:10.1090/S0002-9947-2011-05546-7}; arXiv:0704.3448.

\bibitem{Guinand1948}
A.~P. Guinand, \emph{A summation formula in the theory of prime numbers}, Proc. London Math. Soc. (2) \textbf{50} (1948), 107--119, \href{https://doi.org/10.1112/plms/s2-50.2.107}{doi:10.1112/plms/s2-50.2.107}.

\bibitem{Hormander1990}
L.~H\"ormander, \emph{The Analysis of Linear Partial Differential Operators I}, 2nd ed., Grundlehren der mathematischen Wissenschaften 256, Springer, Berlin, 1990, \href{https://doi.org/10.1007/978-3-642-61497-2}{doi:10.1007/978-3-642-61497-2}.

\bibitem{Ingham1936}
A.~E. Ingham, Some trigonometrical inequalities with applications to the theory of series, \emph{Math. Z.} \textbf{41} (1936), 367--379.

\bibitem{Beurling1937}
A.~Beurling, Analyse de la loi asymptotique de la distribution des nombres premiers g\'en\'eralis\'es.~I, \emph{Acta Math.} \textbf{68} (1937), 255--291.
\href{https://doi.org/10.1007/BF02546666}{doi:10.1007/BF02546666}.

\bibitem{Bohr1926}
H.~Bohr, Zur Theorie der fastperiodischen Funktionen. III. Dirichletentwicklung analytischer Funktionen, \emph{Acta Math.} \textbf{47} (1926), 237--281.
\href{https://doi.org/10.1007/BF02543846}{doi:10.1007/BF02543846}.

\bibitem{BochnerChandrasekharan1956}
S.~Bochner and K.~Chandrasekharan, On Riemann's functional equation, \emph{Ann. of Math. (2)} \textbf{63} (1956), 336--360.
\href{https://doi.org/10.2307/1969614}{doi:10.2307/1969614}.

\bibitem{BrouckeDebruyneVindas2020}
F.~Broucke, G.~Debruyne and J.~Vindas, Beurling integers with RH and large oscillation, \emph{Adv. Math.} \textbf{370} (2020), 107240.
\href{https://doi.org/10.1016/j.aim.2020.107240}{doi:10.1016/j.aim.2020.107240}.

\bibitem{BrouckeVindas2024}
F.~Broucke and J.~Vindas, A new generalized prime random approximation procedure and some of its applications, \emph{Math. Z.} \textbf{307} (2024), Paper No.~62.
\href{https://doi.org/10.1007/s00209-024-03526-4}{doi:10.1007/s00209-024-03526-4}.

\bibitem{CohnElkies2003}
H.~Cohn and N.~Elkies, New upper bounds on sphere packings~I, \emph{Ann. of Math. (2)} \textbf{157} (2003), no.~2, 689--714.
\href{https://doi.org/10.4007/annals.2003.157.689}{doi:10.4007/annals.2003.157.689}.

\bibitem{Corless1996}
R.~M. Corless, G.~H. Gonnet, D.~E.~G. Hare, D.~J. Jeffrey and D.~E. Knuth, On the Lambert $W$ function, \emph{Adv. Comput. Math.} \textbf{5} (1996), 329--359.
\href{https://doi.org/10.1007/BF02124750}{doi:10.1007/BF02124750}.

\bibitem{DaubechiesDeVore2003}
I.~Daubechies and R.~DeVore, Approximating a bandlimited function using very coarsely quantized data: a family of stable sigma-delta modulators of arbitrary order, \emph{Ann. of Math.} \textbf{158} (2003), 679--710.
\href{https://doi.org/10.4007/annals.2003.158.679}{doi:10.4007/annals.2003.158.679}.

\bibitem{Davenport1980}
H.~Davenport, \emph{Multiplicative Number Theory}, 2nd ed., revised by H.~L. Montgomery, Graduate Texts in Mathematics~74, Springer, New York, 1980.
\href{https://doi.org/10.1007/978-1-4757-5927-3}{doi:10.1007/978-1-4757-5927-3}.

\bibitem{delaValleePoussin1896}
C.-J. de la Vall\'ee Poussin, Recherches analytiques sur la th\'eorie des nombres premiers, \emph{Ann. Soc. Sci. Bruxelles} \textbf{20} (1896), 183--256.

\bibitem{Diamond1970}
H.~G. Diamond, Asymptotic distribution of Beurling's generalized integers, \emph{Illinois J. Math.} \textbf{14} (1970).
\href{https://doi.org/10.1215/ijm/1256053295}{doi:10.1215/ijm/1256053295}.

\bibitem{DiamondMontgomeryVorhauer2006}
H.~G. Diamond, H.~L. Montgomery and U.~M.~A. Vorhauer, Beurling primes with large oscillation, \emph{Math. Ann.} \textbf{334} (2006), 1--36.
\href{https://doi.org/10.1007/s00208-005-0638-2}{doi:10.1007/s00208-005-0638-2}.

\bibitem{DiamondZhang2016}
H.~G. Diamond and W.-B. Zhang, \emph{Beurling Generalized Numbers}, Mathematical Surveys and Monographs~213, American Mathematical Society, Providence, RI, 2016.
\href{https://doi.org/10.1090/surv/213}{doi:10.1090/surv/213}.

\bibitem{DLMF2024}
NIST Digital Library of Mathematical Functions, F.~W.~J. Olver, A.~B. Olde Daalhuis, D.~W. Lozier et al. (eds.), \S4.13 (Lambert $W$-function) and \S5.11 (Stirling series).
\url{https://dlmf.nist.gov/}.

\bibitem{Edwards1974}
H.~M. Edwards, \emph{Riemann's Zeta Function}, Academic Press, New York, 1974.

\bibitem{FrancaLeClair2014}
G.~Fran\c{c}a and A.~LeClair, Statistical and other properties of Riemann zeros based on an explicit equation for the $n$-th zero on the critical line, arXiv:1307.8395.
\href{https://doi.org/10.48550/arXiv.1307.8395}{doi:10.48550/arXiv.1307.8395}.

\bibitem{GrosswaldSchnitzer1978}
E.~Grosswald and F.~J. Schnitzer, A class of modified $\zeta$ and $L$-functions, \emph{Pacific J. Math.} \textbf{74} (1978), no.~2, 357--364.
\href{https://doi.org/10.2140/pjm.1978.74.357}{doi:10.2140/pjm.1978.74.357}.

\bibitem{GunturkThao2003}
C.~S. G\"unt\"urk and N.~T. Thao, Ergodic dynamics in sigma-delta quantization: tiling invariant sets and spectral analysis of error, arXiv:math/0308223.
\href{https://doi.org/10.48550/arXiv.math/0308223}{doi:10.48550/arXiv.math/0308223}.

\bibitem{Hadamard1896}
J.~Hadamard, Sur la distribution des z\'eros de la fonction $\zeta(s)$ et ses cons\'equences arithm\'etiques, \emph{Bull. Soc. Math. France} \textbf{24} (1896), 199--220.
\href{https://doi.org/10.24033/bsmf.545}{doi:10.24033/bsmf.545}.

\bibitem{Hamburger1921}
H.~Hamburger, \"Uber die Riemannsche Funktionalgleichung der $\zeta$-Funktion, \emph{Math. Z.} \textbf{10} (1921), 240--254.
\href{https://doi.org/10.1007/BF01211612}{doi:10.1007/BF01211612}.

\bibitem{KaczorowskiPerelli1999}
J.~Kaczorowski and A.~Perelli, On the structure of the Selberg class, I: $0\le d\le1$, \emph{Acta Math.} \textbf{182} (1999), 207--241.
\href{https://doi.org/10.1007/BF02392574}{doi:10.1007/BF02392574}.

\bibitem{Kadec1964}
M.~I. Kadec, \emph{The exact value of the Paley--Wiener constant}, Dokl. Akad. Nauk SSSR \textbf{155} (1964), 1253--1254; English transl. Soviet Math. Dokl. \textbf{5} (1964), 559--561. (No DOI.)

\bibitem{Knopp1994}
M.~I. Knopp, On Dirichlet series satisfying Riemann's functional equation, \emph{Invent. Math.} \textbf{117} (1994), 361--372.
\href{https://doi.org/10.1007/BF01232248}{doi:10.1007/BF01232248}.

\bibitem{KurasovSarnak2020}
P.~Kurasov and P.~Sarnak, Stable polynomials and crystalline measures, \emph{J. Math. Phys.} \textbf{61} (2020), no.~8, 083501.
\href{https://doi.org/10.1063/5.0012286}{doi:10.1063/5.0012286}.

\bibitem{LapidusvanFrankenhuysen2000}
M.~L. Lapidus and M.~van Frankenhuysen, \emph{Fractal Geometry and Number Theory: Complex Dimensions of Fractal Strings and Zeros of Zeta Functions}, Birkh\"auser, Boston, 2000.
\href{https://doi.org/10.1007/978-1-4612-5314-3}{doi:10.1007/978-1-4612-5314-3}.

\bibitem{LevOlevskii2015}
N.~Lev and A.~Olevskii, Quasicrystals and Poisson's summation formula, \emph{Invent. Math.} \textbf{200} (2015), no.~2, 585--606.
\href{https://doi.org/10.1007/s00222-014-0542-z}{doi:10.1007/s00222-014-0542-z}.

\bibitem{LyubarskiiSeip1997}
Yu.~I. Lyubarskii and K.~Seip, \emph{Complete interpolating sequences for Paley--Wiener spaces and Muckenhoupt's $(A_p)$ condition}, Rev. Mat. Iberoam. \textbf{13} (1997), no.~2, 361--376, \href{https://doi.org/10.4171/RMI/224}{doi:10.4171/RMI/224}.

\bibitem{MontgomeryVaughan1974}
H.~L. Montgomery and R.~C. Vaughan, \emph{Hilbert's inequality}, J. London Math. Soc. (2) \textbf{8} (1974), 73--82, \href{https://doi.org/10.1112/jlms/s2-8.1.73}{doi:10.1112/jlms/s2-8.1.73}.

\bibitem{OlevskiiUlanovskii2020}
A.~Olevskii and A.~Ulanovskii, Fourier quasicrystals with unit masses, \emph{C. R. Math. Acad. Sci. Paris} \textbf{358} (2020), no.~11--12, 1207--1211.
\href{https://doi.org/10.5802/crmath.142}{doi:10.5802/crmath.142}.

\bibitem{OdlyzkoTeRiele1985}
A.~M. Odlyzko and H.~J.~J. te~Riele, Disproof of the Mertens conjecture, \emph{J. reine angew. Math.} \textbf{357} (1985), 138--160.
\href{https://doi.org/10.1515/crll.1985.357.138}{doi:10.1515/crll.1985.357.138}.

\bibitem{Nakamura2020}
T.~Nakamura, Dirichlet series with periodic coefficients, Riemann's functional equation and real zeros of Dirichlet $L$-functions, preprint, arXiv:2008.02570 (2020).
\href{https://doi.org/10.48550/arXiv.2008.02570}{doi:10.48550/arXiv.2008.02570}.

\bibitem{Perelli2016}
A.~Perelli, Converse theorems: from the Riemann zeta function to the Selberg class, arXiv:1605.02354.
\href{https://doi.org/10.48550/arXiv.1605.02354}{doi:10.48550/arXiv.1605.02354}.

\bibitem{PlattTrudgian2021}
D.~J. Platt and T.~S. Trudgian, The Riemann hypothesis is true up to $3\cdot10^{12}$, \emph{Bull. London Math. Soc.} \textbf{53} (2021); arXiv:2004.09765.
\href{https://doi.org/10.48550/arXiv.2004.09765}{doi:10.48550/arXiv.2004.09765}.

\bibitem{Riemann1859}
B.~Riemann, \"Uber die Anzahl der Primzahlen unter einer gegebenen Gr\"osse, \emph{Monatsberichte der K\"oniglich Preussischen Akademie der Wissenschaften zu Berlin} (1859), 671--680.

\bibitem{Rudin1987}
W.~Rudin, \emph{Real and Complex Analysis}, 3rd ed., McGraw--Hill, New York, 1987. (No DOI.)

\bibitem{Selberg1946}
A.~Selberg, \emph{Contributions to the theory of the Riemann zeta-function}, Arch. Math. Naturvid. \textbf{48} (1946), no.~5, 89--155. (No DOI.)

\bibitem{Siegel1922}
C.~L. Siegel, Bemerkung zu einem Satz von Hamburger \"uber die Funktionalgleichung der Riemannschen Zetafunktion, \emph{Math. Ann.} \textbf{86} (1922), 276--279.
\href{https://doi.org/10.1007/BF01457989}{doi:10.1007/BF01457989}.

\bibitem{Titchmarsh1939}
E.~C. Titchmarsh, \emph{The Theory of Functions}, 2nd ed., Oxford University Press, 1939.

\bibitem{Titchmarsh1986}
E.~C. Titchmarsh, \emph{The Theory of the Riemann Zeta-Function}, 2nd ed., revised by D.~R. Heath-Brown, Clarendon Press, Oxford, 1986.

\bibitem{vonMangoldt1895}
H.~von Mangoldt, Zu Riemanns Abhandlung ``\"Uber die Anzahl der Primzahlen unter einer gegebenen Gr\"osse'', \emph{J. reine angew. Math.} \textbf{114} (1895), 255--305.
\href{https://doi.org/10.1515/crll.1895.114.255}{doi:10.1515/crll.1895.114.255}.

\bibitem{Weil1952}
A.~Weil, \emph{Sur les ``formules explicites'' de la th\'eorie des nombres premiers}, Comm. S\'em. Math. Univ. Lund [Medd. Lunds Univ. Mat. Sem.] (1952), Tome Suppl\'ementaire, 252--265. (No DOI.)

\bibitem{WhittakerWatson1927}
E.~T. Whittaker and G.~N. Watson, \emph{A Course of Modern Analysis}, 4th ed., Cambridge University Press, 1927; reissued in the Cambridge Mathematical Library, 1996, \href{https://doi.org/10.1017/CBO9780511608759}{doi:10.1017/CBO9780511608759}.

\bibitem{Zemanian1965}
A.~H. Zemanian, \emph{Distribution Theory and Transform Analysis}, McGraw-Hill, New York, 1965; reprinted by Dover, New York, 1987. (No DOI.)

\end{thebibliography}


\clearpage
\makeatletter
\expandafter\gdef\csname CGTip@prop:additive-level-v18\endcsname{v18 appendix: the additive family is exactly a level-set reformulation through the quotient zeta/zeta_RVM.}
\expandafter\gdef\csname CGTip@eq:additive-family-v18\endcsname{v18 appendix: G_epsilon=(1-epsilon) zeta_RVM+epsilon zeta.}
\expandafter\gdef\csname CGTip@eq:additive-level-v18\endcsname{v18 appendix: G_epsilon=zeta_RVM((1-epsilon)+epsilon Q_RVM).}
\expandafter\gdef\csname CGTip@eq:Q-negative-level-v18\endcsname{v18 appendix: zeros of the additive family satisfy Q_RVM=-(1-epsilon)/epsilon.}
\expandafter\gdef\csname CGTip@eq:rvm-first-rs-v18\endcsname{v18 appendix: transition theorem, zeta_RVM=varpi_RVM(1+chi)(1+O(1/|t|)).}
\expandafter\gdef\csname CGTip@eq:first-rs-hardy-v18\endcsname{v18 appendix: on the critical line, exp(i theta)(1+chi)=2 cos theta.}
\expandafter\gdef\csname CGTip@eq:finite-additive-v18\endcsname{v18 appendix: finite Dirichlet-polynomial completion with the exact functional equation.}
\expandafter\gdef\csname CGTip@eq:finite-additive-hardy-v18\endcsname{v18 appendix: Hardy form of the finite additive completion.}
\expandafter\gdef\csname CGTip@eq:dominance-v18\endcsname{v18 appendix: sufficient n=1 dominance condition at Gram points.}
\expandafter\gdef\csname CGTip@eq:clock-injectivity-v18\endcsname{v18 appendix: autonomous midpoint-clock injectivity problem in meromorphic/Dirichlet form.}
\expandafter\gdef\csname CGTip@eq:bounded-section-v18\endcsname{v18 appendix: numerical regularization bounds for the finite-section experiment.}
\makeatother
\section{Python programs of Section~\CGref{sec:routeB-finite}}
\label{app:code-routeB}

The programs of \texttt{code/routeB/} produce the numbers of Section~\CGref{subsec:rb-numerics} and Figures~\CGref{fig:rb-gonek} and~\CGref{fig:rb-scaling}. They use only the clock, the primes and the Gonek lists; \texttt{zeros.py} computes the zeta ordinates used in the diagnostic columns with \texttt{mpmath}. The exploratory scripts (\texttt{staged.py}, \texttt{lsq*.py}, \texttt{basin*.py}, \texttt{coupling.py}, \texttt{horizon.py}, \texttt{certG.py}, \texttt{cert2--6.py}) are distributed in the same directory with their logs and are described in \texttt{RESULTS.txt}.

\subsection{\texttt{core.py}}
\label{app:code-routeB-core}

\begin{lstlisting}[language=Python]
import numpy as np
from scipy.special import lambertw
from scipy.integrate import quad
SQ2PI=np.sqrt(2*np.pi)
def clock(K):
    k=np.arange(1,K+1); a=(k-11/8)/np.e
    return np.real(2*np.pi*(k-11/8)/lambertw(a))
def vonmangoldt(Nmax):
    L=np.zeros(Nmax+1)
    for p in range(2,Nmax+1):
        if all(p%q for q in range(2,int(p**0.5)+1)):
            q=p
            while q<=Nmax: L[q]=np.log(p); q*=p
    return L
LAM=vonmangoldt(200000)
def kfun(u): return np.cosh(u/2)-np.exp(-u/2)/(2*(1-np.exp(-2*u)))
def gauss(c,s,u): return np.exp(-(u-c)**2/(2*s*s))
def Fphi(c,s,kap):  # int_R gauss cos(kap u) du
    return s*SQ2PI*np.exp(-0.5*(s*kap)**2)*np.cos(c*kap)
def dFphi(c,s,kap):
    return s*SQ2PI*np.exp(-0.5*(s*kap)**2)*(-(s*s)*kap*np.cos(c*kap)-c*np.sin(c*kap))
_kcache={}
def kpair(c,s):
    key=(round(c,12),round(s,12))
    if key not in _kcache:
        a=max(c-9*s,1e-9); b=c+9*s
        _kcache[key]=quad(lambda u:kfun(u)*gauss(c,s,u),a,b,limit=400,epsabs=1e-13,epsrel=1e-12)[0]
    return _kcache[key]
def masspair(c,s):
    nmax=int(np.exp(c+9*s))+1
    n=np.arange(2,nmax+1); a=LAM[n]/(2*np.sqrt(n))
    return np.sum(a*gauss(c,s,np.log(n)))
def residual(tests,kap):
    """tests: list of (c,s); kap: full list (must extend far enough)"""
    out=[]
    for c,s in tests:
        out.append(np.sum(Fphi(c,s,kap))+masspair(c,s)-kpair(c,s))
    return np.array(out)
def jac(tests,kap,K):
    return np.array([dFphi(c,s,kap[:K]) for c,s in tests])
\end{lstlisting}

\subsection{\texttt{zeros.py}}
\label{app:code-routeB-zeros}

\begin{lstlisting}[language=Python]
import mpmath as mp, numpy as np, sys
mp.mp.dps=20
N=int(sys.argv[1])
g=[float(mp.zetazero(n).imag) for n in range(1,N+1)]
np.save('gamma.npy',np.array(g)); print(len(g), g[-1])
\end{lstlisting}

\subsection{\texttt{gonek\_fast.py}}
\label{app:code-routeB-gonek-fast}

\begin{lstlisting}[language=Python]
import numpy as np
from scipy.optimize import brentq
from core import LAM
def theta(t): t=np.asarray(t,float); return t/2*np.log(t/(2*np.pi))-t/2-np.pi/8+1/(48*t)+7/(5760*t**3)
def make_alpha(X):
    n=np.arange(2,int(X*X)+1); L=LAM[n]; w=np.where(n<=X,1.0,2-np.log(n)/np.log(X)); c=L*w/(np.sqrt(n)*np.log(n)); ln=np.log(n)
    return lambda t: -(np.sin(np.multiply.outer(np.atleast_1d(t),ln))@c)
def gonek_list(X,tmax,t0=10.0,h=0.005):
    al=make_alpha(X); ph=lambda t: theta(t)+al(t)
    ts=np.arange(t0,tmax,h); v=(ph(ts)-np.pi/2)/np.pi; j=np.floor(v); idx=np.where(np.diff(j)!=0)[0]; roots=[]
    for i in idx:
        lev=max(j[i],j[i+1]); f=lambda t:(ph(np.array([t]))[0]-np.pi/2)/np.pi-lev
        a,b=ts[i],ts[i+1]
        if f(a)*f(b)<0: roots.append(brentq(f,a,b,xtol=1e-13))
    return np.array(roots)
if __name__=='__main__':
    import sys
    X=float(sys.argv[1]); r=gonek_list(X,float(sys.argv[2])); np.save(f'gonekF_X{int(X)}.npy',r); print(len(r),r[:3],r[-1])
\end{lstlisting}

\subsection{\texttt{horizonG.py}}
\label{app:code-routeB-horizonG}

\begin{lstlisting}[language=Python]
import numpy as np, json, sys
from scipy.optimize import least_squares
from core import *
g=np.load('gamma.npy'); tau=clock(60000)
q,s0,sp,M,INT,LAG=0.85,0.08,1.0,int(sys.argv[1]),float(sys.argv[2]),int(sys.argv[3])
r=np.load('gonekF_X8.npy'); kap=tau.copy(); kap[:len(r)]=r; tests=[]; level=[]; log=[]
for m in range(1,M+1):
    s=s0*q**(m-1); T=4.0/s; U=np.log(max(T,7)/(2*np.pi))+2.0
    new=[(float(c),s) for c in np.arange(9*s,U+1e-9,sp*s)]; tests+=new; level+=[m]*len(new)
    cs=np.array([t[0] for t in tests]); ss=np.array([t[1] for t in tests])
    ki=int(np.sum(tau<=INT/s))
    const=np.array([masspair(c,s_)-kpair(c,s_) for c,s_ in tests])
    tail=np.array([np.sum(Fphi(c,s_,kap[ki:])) for c,s_ in tests])  # edges & beyond: current values (clock)
    F=lambda y: np.sum(Fphi(cs[:,None],ss[:,None],y[None,:]),axis=1)+tail+const
    Jf=lambda y: dFphi(cs[:,None],ss[:,None],y[None,:])
    y0=kap[:ki].copy()
    sol=least_squares(F,y0,jac=Jf,method='trf',bounds=(tau[:ki]-3,tau[:ki]+3),xtol=1e-15,ftol=1e-15,gtol=1e-15,max_nfev=3000)
    kap=kap.copy(); kap[:ki]=np.sort(sol.x)
    Fall=F(kap[:ki]); lv=np.array(level)
    lag=np.max(np.abs(Fall[lv<=m-LAG])) if np.any(lv<=m-LAG) else float('nan')
    nk=min(ki,len(g)); dev=np.abs(kap[:nk]-g[:nk])
    front=lambda tol: (int(np.where(dev>tol)[0][0]) if np.any(dev>tol) else nk)
    rec=dict(m=m,s=s,T=T,interior=ki,tests=len(tests),res_all=float(np.max(np.abs(Fall))),res_lag=float(lag),
             dev5=float(dev[:5].max()),dev10=float(dev[:10].max()),dev20=float(dev[:min(20,nk)].max()),
             f6=front(1e-6),f4=front(1e-4),f2=front(1e-2),dev_tau=float(np.max(np.abs(kap[:ki]-tau[:ki]))))
    log.append(rec)
    print(f"m={m:2d} T={T:5.0f} interior={ki:3d} tests={len(tests):4d} res_all={rec['res_all']:.1e} res_lag{LAG}={rec['res_lag']:.1e} |k-g| k<=5:{rec['dev5']:.1e} k<=10:{rec['dev10']:.1e} k<=20:{rec['dev20']:.1e} front(1e-6/1e-4/1e-2)={rec['f6']}/{rec['f4']}/{rec['f2']} |k-tau|max={rec['dev_tau']:.2f}",flush=True)
json.dump(log,open(f'horizonG_M{M}_I{INT}_L{LAG}.json','w'),indent=1); np.save(f'horizonG_kap_M{M}_I{INT}.npy',kap)
\end{lstlisting}

\subsection{\texttt{cert7.py}}
\label{app:code-routeB-cert7}

\begin{lstlisting}[language=Python]
# Computer-assisted (gamma-free) existence certificate for finite square sections, at the
# receding-horizon iterate kappa^H (horizonG, Gonek X=8 initialization, M=16).
# Stage m: tests = levels <= m (spacing 1.0 s), unknowns = interior ordinates (tau_k <= 2.4/s_m),
# other ordinates frozen at kappa^H. Square subset of tests chosen by pivoted QR on J^T.
# F_j separable in kappa_k  =>  J(x)-J(x0) has columns f_j''(x0_k) delta_k + R_jk, |R_jk|<=A3_jk delta_k^2/2.
# x0 = kappa^H polished by 8 Newton steps on the square system (gamma-free).
# Banach for Psi(x)=x-Q F_sel(x) on {||x-x0||_inf<=r}:  k(r)=||QB||r + 0.5*|| |Q|A3 || r^2 < 1,  eta<=(1-k(r)) r.
import numpy as np, json
from scipy.linalg import qr
from core import *
g=np.load('gamma.npy'); tau=clock(60000); kH=np.load('horizonG_kap_M16_I2.4.npy')
q,s0,sp,INT=0.85,0.08,1.0,2.4
def d2(c,s,k):
    a=s*s; e=s*SQ2PI*np.exp(-a*k*k/2)
    return e*(((a*k)**2-a-c*c)*np.cos(c*k)+2*a*k*c*np.sin(c*k))
def A3b(c,s,k,r):
    a=s*s; kl=np.maximum(np.abs(k)-r,0); ku=np.abs(k)+r; w=a*ku+c
    return s*SQ2PI*np.exp(-a*kl*kl/2)*(w**3+3*a*w)
tests=[]; out=[]
for m in range(1,17):
    s=s0*q**(m-1); T=4.0/s; U=np.log(max(T,7)/(2*np.pi))+2.0
    tests+=[(float(c),s) for c in np.arange(9*s,U+1e-9,sp*s)]
    n=int(np.sum(tau<=INT/s))
    cs=np.array([t[0] for t in tests]); ss=np.array([t[1] for t in tests])
    x0=kH[:n]
    J=dFphi(cs[:,None],ss[:,None],x0[None,:])
    _,_,piv=qr(J.T,pivoting=True); sel=np.sort(piv[:n])
    kk=kH.copy(); tsel=[tests[i] for i in sel]
    try:
        for it in range(8):                      # gamma-free Newton polishing on the square system
            Jc=dFphi(cs[sel][:,None],ss[sel][:,None],kk[None,:n]); kk[:n]=kk[:n]-np.linalg.solve(Jc,residual(tsel,kk))
    except np.linalg.LinAlgError:
        print(f"m={m:2d} Newton polishing failed (singular square Jacobian)"); break
    x0=kk[:n]; J=dFphi(cs[:,None],ss[:,None],x0[None,:])
    Q=np.linalg.inv(J[sel]); F=residual(tsel,kk)
    eta=np.max(np.abs(Q@F))
    QB=np.max(np.sum(np.abs(Q@d2(cs[sel][:,None],ss[sel][:,None],x0[None,:])),axis=1))
    best=None
    for r in np.geomspace(max(eta,1e-15)*1.0001,1.0,400):
        C3=np.max(np.sum(np.abs(Q)@A3b(cs[sel][:,None],ss[sel][:,None],x0[None,:],r),axis=1)) if best is None or r<10*best[0] else C3
        kr=QB*r+0.5*C3*r*r
        if kr<1 and eta<=(1-kr)*r: best=(r,kr); break
    margin=max((1-(QB*r+0.5*np.max(np.sum(np.abs(Q)@A3b(cs[sel][:,None],ss[sel][:,None],x0[None,:],r),axis=1))*r*r))*r for r in np.geomspace(1e-12,1e-1,60))
    rho_all=float(np.max(np.abs(residual(tests,kk))))
    Linf=np.max(np.sum(np.abs(Q),axis=1))
    nk=min(n,len(g)); devg=float(np.max(np.abs(x0[:nk]-g[:nk]))); devt=float(np.max(np.abs(x0-tau[:n])))
    rec=dict(m=m,T=T,n=n,Linf=float(Linf),eta=float(eta),QB=float(QB),certified=best is not None,margin=float(margin),rho_all_levels=rho_all,
             r=(float(best[0]) if best else None),k_r=(float(best[1]) if best else None),dev_gamma=devg,dev_tau=devt)
    out.append(rec)
    print(f"m={m:2d} T={T:5.0f} n={n:3d} Linf={Linf:.1e} eta={eta:.1e} ||QB||={QB:.1e} "
          + (f"CERTIFIED r={best[0]:.1e} k(r)={best[1]:.2f}" if best else "not certified")
          + f" margin={margin:.1e} rho(all tests<=m)={rho_all:.1e} | max|x0-tau|={devt:.2f} max|x0-gamma|={devg:.1e}",flush=True)
json.dump(out,open('cert7.json','w'),indent=1)
\end{lstlisting}

\subsection{\texttt{entry.py}}
\label{app:code-routeB-entry}

\begin{lstlisting}[language=Python]
# Block-entry test (item (i) of the review): from the certified stage-m solution x^(m)
# (square section polished as in cert7), put the NEW band (tau_k in (2.4/s_m, 2.4/s_{m+1}]) and
# everything beyond at Gonek X=8 values, and ask whether a Banach certificate holds for
# (a) the band-only square system (new band unknowns, tests from level m+1),
# (b) the full square system of stage m+1, started there (no polishing).
import numpy as np, json
from scipy.linalg import qr
from core import *
exec(open('cert5.py').read().split("tests=[]; out=[]")[0].split("import numpy as np, json")[1])  # d2, A3b, data
rG=np.load('gonekF_X8.npy'); kG=tau.copy(); kG[:len(rG)]=rG
def certify(Q,F,QBn,A3fun):
    eta=np.max(np.abs(Q@F))
    for r in np.geomspace(max(eta,1e-15)*1.0001,1.0,300):
        kr=QBn*r+0.5*np.max(np.sum(np.abs(Q)@A3fun(r),axis=1))*r*r
        if kr<1 and eta<=(1-kr)*r: return eta,r
    return eta,None
def margin(QBn,Q,A3fun): return max((1-(QBn*r+0.5*np.max(np.sum(np.abs(Q)@A3fun(r),axis=1))*r*r))*r for r in np.geomspace(1e-9,1e-1,50))
tests=[]; lev=[]; out=[]; prev=None
for m in range(1,16):
    s=s0*q**(m-1); T=4.0/s; U=np.log(max(T,7)/(2*np.pi))+2.0
    new=[(float(c),s) for c in np.arange(9*s,U+1e-9,sp*s)]; tests+=new; lev+=[m]*len(new)
    n=int(np.sum(tau<=INT/s)); cs=np.array([t[0] for t in tests]); ss=np.array([t[1] for t in tests]); lv=np.array(lev)
    if prev is not None:
        x0=kG.copy(); x0[:prev[1]]=prev[0][:prev[1]]          # old interior: certified stage-(m-1) solution; rest Gonek
        band=np.arange(prev[1],n); idx=np.where(lv==m)[0]
        Jb=dFphi(cs[idx][:,None],ss[idx][:,None],x0[None,band]); _,_,pv=qr(Jb.T,pivoting=True); sb=idx[np.sort(pv[:len(band)])]
        Qb=np.linalg.inv(dFphi(cs[sb][:,None],ss[sb][:,None],x0[None,band])); Fb=residual([tests[i] for i in sb],x0)
        QBb=np.max(np.sum(np.abs(Qb@d2(cs[sb][:,None],ss[sb][:,None],x0[None,band])),axis=1))
        A3f=lambda r: A3b(cs[sb][:,None],ss[sb][:,None],x0[None,band],r)
        eb,rb=certify(Qb,Fb,QBb,A3f); mb=margin(QBb,Qb,A3f)
        J=dFphi(cs[:,None],ss[:,None],x0[None,:n]); _,_,pv=qr(J.T,pivoting=True); sf=np.sort(pv[:n])
        Qf=np.linalg.inv(J[sf]); Ff=residual([tests[i] for i in sf],x0)
        QBf=np.max(np.sum(np.abs(Qf@d2(cs[sf][:,None],ss[sf][:,None],x0[None,:n])),axis=1))
        A3g=lambda r: A3b(cs[sf][:,None],ss[sf][:,None],x0[None,:n],r)
        ef,rf=certify(Qf,Ff,QBf,A3g); mf=margin(QBf,Qf,A3g)
        bandG=float(np.max(np.abs(x0[band]-g[band]))) if band[-1]<len(g) else float('nan')
        rec=dict(m=m,T=T,band=len(band),eta_band=float(eb),margin_band=float(mb),cert_band=rb is not None,
                 eta_full=float(ef),margin_full=float(mf),cert_full=rf is not None,gonek_band_err=bandG)
        out.append(rec)
        print(f"m={m:2d} T={T:5.0f} band={len(band):2d} | band-only: eta={eb:.1e} margin={mb:.1e} {'CERT' if rb else 'no'} "
              f"| full: eta={ef:.1e} margin={mf:.1e} {'CERT' if rf else 'no'} | Gonek band err={bandG:.1e}",flush=True)
    kk=kG.copy() if prev is None else x0.copy(); tsel=None
    J=dFphi(cs[:,None],ss[:,None],kk[None,:n]); _,_,pv=qr(J.T,pivoting=True); sel=np.sort(pv[:n]); tsel=[tests[i] for i in sel]
    try:
        for it in range(10):
            Jc=dFphi(cs[sel][:,None],ss[sel][:,None],kk[None,:n]); kk[:n]-=np.linalg.solve(Jc,residual(tsel,kk))
    except np.linalg.LinAlgError: print("singular at",m); break
    prev=(kk.copy(),n)
json.dump(out,open('entry.json','w'),indent=1)
\end{lstlisting}

\subsection{\texttt{laplace\_check.py}}
\label{app:code-routeB-laplace-check}

\begin{lstlisting}[language=Python]
# Numerical check of the identity  (1/2) G'/G(s) = L[K](s)  with
# K = (1/2)e^{u/2} 1_+ - (1/2) g 1_+ - (1/4) fp(1_+/u) + c0 delta_0,
# g(u)=e^{-u/2}/(e^{2u}-1)-1/(2u),  c0=-(1/4)log(2 pi e^{gamma_E}),  z=s-1/2,
# L[fp(1_+/u)](z) = -log z - gamma_E.   Also checks k(u)=cosh(u/2)-e^{-u/2}/(2(1-e^{-2u})) = K|_(0,inf).
from mpmath import mp, mpf, exp, expm1, log, pi, euler, quad, digamma, cosh, inf, mpc
mp.dps=50
def g(u):
    if u<mpf('1e-12'): return -mpf(3)/4       # g(0+)=-3/4; avoids cancellation
    return exp(-u/2)/expm1(2*u)-1/(2*u)
def k_v18(u): return (exp(u/2)-exp(-u/2)/(exp(2*u)-1))/2
def k_mine(u): return cosh(u/2)-exp(-u/2)/(2*(1-exp(-2*u)))
def K_restr(u): return exp(u/2)/2-g(u)/2-1/(4*u)
for u in [mpf('0.1'),mpf(1),mpf('2.5')]:
    print('u=%s  k_v18=%s  k_mine=%s  K|=%s'%(u,k_v18(u),k_mine(u),K_restr(u)))
def halfGG(s): return (1/s+1/(s-1)-log(pi)/2+digamma(s/2)/2)/2
def LK(s):
    z=s-mpf(1)/2
    c0=-(euler+log(2)+log(pi))/4
    I1=1/(2*(z-mpf(1)/2))                      # L[(1/2)e^{u/2}]
    I2=-quad(lambda u: exp(-z*u)*g(u),[0,1,inf])/2
    I3=-(-log(z)-euler)/4
    return I1+I2+I3+c0
for s in [mpf(2),mpf('1.3'),mpc(2,5),mpc('1.5',-3)]:
    print('s=%s  (1/2)G\'/G=%s  L[K]=%s'%(s,halfGG(s),LK(s)))
\end{lstlisting}

\subsection{\texttt{figs.py}}
\label{app:code-routeB-figs}

\begin{lstlisting}[language=Python]
# Figures for lemas_B_v4 / v19: (1) certificate scaling, (2) distance to gamma of clock and Gonek lists.
import numpy as np, json, matplotlib
matplotlib.use('Agg'); import matplotlib.pyplot as plt
from core import clock
g=np.load('gamma.npy'); K=400; tau=clock(2000)
def align(r,K):
    best=None
    for off in range(0,6):
        a=r[off:off+K]
        if len(a)<K: continue
        d=np.abs(a-g[:K]).mean()
        if best is None or d<best[0]: best=(d,off,a)
    return best[2]
sp=2*np.pi/np.log(g[:K]/(2*np.pi))
fig,ax=plt.subplots(figsize=(7.2,4.2))
ax.plot(np.arange(1,K+1),np.abs(tau[:K]-g[:K])/sp,'.',ms=3,color='0.3',label='midpoint clock $\\tau_k$')
rF=np.load('gonekF_X8.npy'); a=align(rF,K)
ax.plot(np.arange(1,K+1),np.abs(a-g[:K])/sp,'.',ms=3,color='C3',label='Gonek list, $X=8$')
K2=160; sp2=sp[:K2]
for X,c in [(3,'C0'),(20,'C2')]:
    r=np.load(f'gonek_X{X}.npy'); a=align(r,K2)
    ax.plot(np.arange(1,K2+1),np.abs(a-g[:K2])/sp2,'.',ms=3,color=c,label=f'Gonek list, $X={X}$')
ax.axhline(0.25,ls='--',color='k',lw=0.8); ax.text(5,0.26,'$0.25$ spacings',fontsize=8)
ax.set_xlabel('$k$'); ax.set_ylabel('$|\\kappa_k-\\gamma_k|/\\mathrm{spacing}$'); ax.set_ylim(0,0.75)
ax.legend(fontsize=8,loc='upper right'); ax.set_title('Distance to the zeta ordinates, in local mean spacings $2\\pi/\\log(\\gamma_k/2\\pi)$',fontsize=9)
fig.tight_layout(); fig.savefig('fig_gonek_vs_clock.png',dpi=160)
d=json.load(open('cert7.json'))
T=np.array([x['T'] for x in d]); L=np.array([x['Linf'] for x in d]); QB=np.array([x['QB'] for x in d]); mg=np.array([x['margin'] for x in d])
fig,ax=plt.subplots(figsize=(7.2,4.2))
for y,lab,c in [(L,'$L_m=\\|Q_m\\|_{\\infty\\to\\infty}$','C0'),(QB,'$\\|Q_mB_m\\|_{\\infty\\to\\infty}$','C1'),(1/mg,'1 / (Banach margin)','C3')]:
    p=np.polyfit(np.log(T[3:]),np.log(y[3:]),1); ax.loglog(T,y,'o-',ms=4,color=c,label=lab+f'  (slope {p[0]:.2f})')
e=json.load(open('entry.json')); Te=np.array([x['T'] for x in e]); eb=np.array([x['eta_band'] for x in e]); mb=np.array([x['margin_band'] for x in e])
ax.loglog(Te,1/mb,'s--',ms=4,color='C2',label='1 / (band-only margin)')
ax.loglog(Te,1/eb,'^--',ms=4,color='C4',label='1 / (Gonek entry residual $\\eta$, band)')
ax.set_xlabel('height $T_m$'); ax.legend(fontsize=7.5); ax.set_title('Certificate constants of the square finite sections (floating point)',fontsize=9)
fig.tight_layout(); fig.savefig('fig_cert_scaling.png',dpi=160)
print('ok')
\end{lstlisting}

\section{Console output of the programs of Section~\CGref{sec:routeB-finite}}
\label{app:out-routeB}

\subsection{\texttt{RESULTS.txt}}
\label{app:out-routeB-results}

\begin{lstlisting}[style=out]
RVM rung system - staged certificate numerics (2026-10-10)
Tests: Gaussians phi(u)=exp(-(u-c)^2/(2 s^2)), c>=9s; stage m: s_m=0.08*0.85^(m-1), T_m=4/s_m,
centres spaced 1.0*s_m in [9 s_m, log(max(T_m,7)/2pi)+2]. Interior unknowns: tau_k <= 2.4/s_m.
Masses a_n = Lambda(n)/(2 sqrt n); k(u)=cosh(u/2)-e^{-u/2}/(2(1-e^{-2u})).
Gonek list X=8: roots of (theta(t)+alpha_X(t)-pi/2)/pi in Z (gonek_fast.py).
gamma (400 zeros, mpmath) is used ONLY for diagnostics (columns 'gamma'), never in the constructions.

FILES
 core.py          residual, Jacobian, clock, k(u), mass pairing
 validate.py      R(gamma)~1e-16 normalisation check            -> out_validate.txt
 gonek*.py        Gonek lists, distance to gamma                 -> out_gonek.txt, gonek*.npy
 horizon.py       receding horizon from the clock (fails)        -> horizon_M16_I2.4_L4.json
 horizonG.py      receding horizon from Gonek X=8 (works)        -> horizonG_M16_I2.4_L4.json, horizonG_kap_M16_I2.4.npy
 certG.py         l-inf Kantorovich at Gonek list (fails)         -> out_certG.txt
 cert2/3/4.py     l2 / Gauss-Newton / diagonal l-inf constants at Gonek list -> out_cert2/3/4.txt
 cert5/6/7.py     Banach certificates of square finite sections (cert7 = final) -> out_cert5/6/7.txt, cert*.json
 entry.py         block-entry test (item (i))                    -> out_entry.txt, entry.json
 staged.py, lsq*.py, floor.py, basin*.py, reg.py, coupling.py: earlier exploratory runs (logs *.json)

horizon_M16_I2.4_L4.json
 m= 1 T=   50 n=  4 res_all=1.4e-02 res_lag4=nan max|k-g|(k<=20)=8.1e-02 front(1e-2)=3 max|k-tau|=0.77
 m= 2 T=   59 n=  5 res_all=1.4e-02 res_lag4=nan max|k-g|(k<=20)=4.6e-03 front(1e-2)=5 max|k-tau|=0.69
 m= 3 T=   69 n=  7 res_all=2.1e-02 res_lag4=nan max|k-g|(k<=20)=1.9e-02 front(1e-2)=5 max|k-tau|=0.69
 m= 4 T=   81 n=  9 res_all=1.9e-02 res_lag4=nan max|k-g|(k<=20)=1.2e-02 front(1e-2)=7 max|k-tau|=0.83
 m= 5 T=   96 n= 12 res_all=1.9e-02 res_lag4=2.9e-05 max|k-g|(k<=20)=3.2e-02 front(1e-2)=10 max|k-tau|=0.84
 m= 6 T=  113 n= 16 res_all=1.3e-02 res_lag4=1.7e-04 max|k-g|(k<=20)=5.0e-02 front(1e-2)=11 max|k-tau|=0.88
 m= 7 T=  133 n= 20 res_all=1.7e-02 res_lag4=6.9e-05 max|k-g|(k<=20)=1.1e-01 front(1e-2)=17 max|k-tau|=0.90
 m= 8 T=  156 n= 26 res_all=1.3e-02 res_lag4=8.6e-05 max|k-g|(k<=20)=3.5e-03 front(1e-2)=23 max|k-tau|=0.90
 m= 9 T=  183 n= 34 res_all=1.2e-02 res_lag4=1.8e-05 max|k-g|(k<=20)=3.6e-04 front(1e-2)=31 max|k-tau|=0.90
 m=10 T=  216 n= 43 res_all=2.6e-02 res_lag4=1.5e-04 max|k-g|(k<=20)=9.1e-04 front(1e-2)=30 max|k-tau|=0.90
 m=11 T=  254 n= 54 res_all=3.3e-02 res_lag4=7.2e-04 max|k-g|(k<=20)=1.4e-03 front(1e-2)=31 max|k-tau|=0.90
 m=12 T=  299 n= 68 res_all=4.2e-02 res_lag4=3.1e-03 max|k-g|(k<=20)=2.0e-03 front(1e-2)=30 max|k-tau|=0.90
 m=13 T=  352 n= 85 res_all=5.6e-02 res_lag4=9.0e-03 max|k-g|(k<=20)=1.2e-03 front(1e-2)=33 max|k-tau|=0.90
 m=14 T=  414 n=107 res_all=6.3e-02 res_lag4=1.7e-02 max|k-g|(k<=20)=1.3e-03 front(1e-2)=31 max|k-tau|=0.90
 m=15 T=  487 n=133 res_all=8.0e-02 res_lag4=3.0e-02 max|k-g|(k<=20)=1.4e-03 front(1e-2)=32 max|k-tau|=0.90
 m=16 T=  572 n=165 res_all=8.6e-02 res_lag4=3.7e-02 max|k-g|(k<=20)=1.3e-03 front(1e-2)=32 max|k-tau|=0.90
horizonG_M16_I2.4_L4.json
 m= 1 T=   50 n=  4 res_all=1.1e-04 res_lag4=nan max|k-g|(k<=20)=4.4e-04 front(1e-2)=4 max|k-tau|=0.69
 m= 2 T=   59 n=  5 res_all=1.7e-04 res_lag4=nan max|k-g|(k<=20)=3.1e-04 front(1e-2)=5 max|k-tau|=0.69
 m= 3 T=   69 n=  7 res_all=2.7e-04 res_lag4=nan max|k-g|(k<=20)=1.6e-03 front(1e-2)=7 max|k-tau|=0.69
 m= 4 T=   81 n=  9 res_all=3.8e-04 res_lag4=nan max|k-g|(k<=20)=3.1e-03 front(1e-2)=9 max|k-tau|=0.84
 m= 5 T=   96 n= 12 res_all=7.4e-04 res_lag4=3.0e-06 max|k-g|(k<=20)=2.1e-03 front(1e-2)=12 max|k-tau|=0.84
 m= 6 T=  113 n= 16 res_all=5.7e-04 res_lag4=1.2e-06 max|k-g|(k<=20)=5.9e-04 front(1e-2)=16 max|k-tau|=0.90
 m= 7 T=  133 n= 20 res_all=1.1e-03 res_lag4=2.8e-06 max|k-g|(k<=20)=1.8e-03 front(1e-2)=20 max|k-tau|=0.90
 m= 8 T=  156 n= 26 res_all=1.8e-03 res_lag4=2.4e-06 max|k-g|(k<=20)=2.5e-04 front(1e-2)=26 max|k-tau|=0.90
 m= 9 T=  183 n= 34 res_all=1.1e-03 res_lag4=4.2e-06 max|k-g|(k<=20)=1.4e-04 front(1e-2)=33 max|k-tau|=0.96
 m=10 T=  216 n= 43 res_all=2.0e-03 res_lag4=3.9e-06 max|k-g|(k<=20)=3.0e-05 front(1e-2)=43 max|k-tau|=0.99
 m=11 T=  254 n= 54 res_all=1.5e-03 res_lag4=4.5e-06 max|k-g|(k<=20)=2.3e-05 front(1e-2)=54 max|k-tau|=0.99
 m=12 T=  299 n= 68 res_all=2.3e-03 res_lag4=6.4e-06 max|k-g|(k<=20)=1.7e-05 front(1e-2)=68 max|k-tau|=0.99
 m=13 T=  352 n= 85 res_all=1.9e-03 res_lag4=3.4e-06 max|k-g|(k<=20)=6.4e-06 front(1e-2)=85 max|k-tau|=0.99
 m=14 T=  414 n=107 res_all=2.5e-03 res_lag4=2.1e-06 max|k-g|(k<=20)=2.8e-06 front(1e-2)=107 max|k-tau|=0.99
 m=15 T=  487 n=133 res_all=2.1e-03 res_lag4=3.8e-06 max|k-g|(k<=20)=4.1e-06 front(1e-2)=133 max|k-tau|=0.99
 m=16 T=  572 n=165 res_all=1.9e-03 res_lag4=8.7e-06 max|k-g|(k<=20)=4.7e-06 front(1e-2)=165 max|k-tau|=0.99

Scaling fits (cert7, m>=4): Linf ~ T^1.81, ||QB|| ~ T^2.10, Banach margin ~ T^-2.08
\end{lstlisting}

\subsection{\texttt{out\_cert7.txt}}
\label{app:out-routeB-cert7}

\begin{lstlisting}[style=out]
m= 1 T=   50 n=  4 Linf=4.1e+01 eta=1.7e-15 ||QB||=3.1e+01 CERTIFIED r=1.7e-15 k(r)=0.00 margin=7.0e-03 rho(all tests<=m)=7.7e-08 | max|x0-tau|=0.69 max|x0-gamma|=7.4e-07
m= 2 T=   59 n=  5 Linf=2.8e+01 eta=1.3e-15 ||QB||=2.2e+01 CERTIFIED r=1.3e-15 k(r)=0.00 margin=1.0e-02 rho(all tests<=m)=1.1e-07 | max|x0-tau|=0.69 max|x0-gamma|=2.2e-07
m= 3 T=   69 n=  7 Linf=3.8e+01 eta=1.8e-15 ||QB||=2.9e+01 CERTIFIED r=1.8e-15 k(r)=0.00 margin=7.1e-03 rho(all tests<=m)=1.6e-07 | max|x0-tau|=0.69 max|x0-gamma|=1.0e-06
m= 4 T=   81 n=  9 Linf=5.0e+01 eta=4.3e-15 ||QB||=7.3e+01 CERTIFIED r=4.3e-15 k(r)=0.00 margin=3.1e-03 rho(all tests<=m)=2.6e-07 | max|x0-tau|=0.84 max|x0-gamma|=1.1e-06
m= 5 T=   96 n= 12 Linf=6.7e+01 eta=2.0e-14 ||QB||=1.4e+02 CERTIFIED r=2.0e-14 k(r)=0.00 margin=1.7e-03 rho(all tests<=m)=9.4e-08 | max|x0-tau|=0.84 max|x0-gamma|=7.5e-07
m= 6 T=  113 n= 16 Linf=1.1e+02 eta=1.0e-14 ||QB||=1.4e+02 CERTIFIED r=1.0e-14 k(r)=0.00 margin=1.7e-03 rho(all tests<=m)=3.5e-07 | max|x0-tau|=0.90 max|x0-gamma|=2.6e-06
m= 7 T=  133 n= 20 Linf=1.8e+02 eta=3.4e-14 ||QB||=2.0e+02 CERTIFIED r=3.4e-14 k(r)=0.00 margin=1.2e-03 rho(all tests<=m)=1.0e-06 | max|x0-tau|=0.90 max|x0-gamma|=2.0e-05
m= 8 T=  156 n= 26 Linf=1.5e+02 eta=2.2e-14 ||QB||=3.2e+02 CERTIFIED r=2.2e-14 k(r)=0.00 margin=7.5e-04 rho(all tests<=m)=9.6e-07 | max|x0-tau|=0.90 max|x0-gamma|=1.3e-05
m= 9 T=  183 n= 34 Linf=1.6e+02 eta=2.5e-14 ||QB||=3.7e+02 CERTIFIED r=2.5e-14 k(r)=0.00 margin=6.6e-04 rho(all tests<=m)=1.1e-06 | max|x0-tau|=0.99 max|x0-gamma|=8.7e-06
m=10 T=  216 n= 43 Linf=3.4e+02 eta=4.0e-14 ||QB||=5.5e+02 CERTIFIED r=4.0e-14 k(r)=0.00 margin=4.5e-04 rho(all tests<=m)=2.3e-06 | max|x0-tau|=0.99 max|x0-gamma|=1.8e-05
m=11 T=  254 n= 54 Linf=4.9e+02 eta=6.8e-14 ||QB||=9.2e+02 CERTIFIED r=6.8e-14 k(r)=0.00 margin=2.6e-04 rho(all tests<=m)=3.2e-06 | max|x0-tau|=0.99 max|x0-gamma|=4.5e-05
m=12 T=  299 n= 68 Linf=4.6e+02 eta=1.8e-13 ||QB||=1.1e+03 CERTIFIED r=1.8e-13 k(r)=0.00 margin=2.2e-04 rho(all tests<=m)=5.5e-06 | max|x0-tau|=0.99 max|x0-gamma|=4.5e-05
m=13 T=  352 n= 85 Linf=6.8e+02 eta=2.2e-13 ||QB||=1.5e+03 CERTIFIED r=2.2e-13 k(r)=0.00 margin=1.6e-04 rho(all tests<=m)=3.6e-05 | max|x0-tau|=0.99 max|x0-gamma|=5.0e-04
m=14 T=  414 n=107 Linf=1.0e+03 eta=2.2e-13 ||QB||=2.4e+03 CERTIFIED r=2.2e-13 k(r)=0.00 margin=1.0e-04 rho(all tests<=m)=3.7e-04 | max|x0-tau|=0.99 max|x0-gamma|=5.4e-03
m=15 T=  487 n=133 Linf=1.6e+03 eta=5.3e-13 ||QB||=3.6e+03 CERTIFIED r=5.3e-13 k(r)=0.00 margin=6.7e-05 rho(all tests<=m)=3.3e-03 | max|x0-tau|=0.99 max|x0-gamma|=4.3e-02
m=16 Newton polishing failed (singular square Jacobian)
\end{lstlisting}

\subsection{\texttt{out\_entry.txt}}
\label{app:out-routeB-entry}

\begin{lstlisting}[style=out]
m= 2 T=   59 band= 1 | band-only: eta=1.7e-03 margin=8.8e-02 CERT | full: eta=5.5e-03 margin=1.0e-02 CERT | Gonek band err=5.2e-03
m= 3 T=   69 band= 2 | band-only: eta=1.5e-02 margin=3.6e-02 CERT | full: eta=1.1e-02 margin=7.1e-03 no | Gonek band err=5.7e-03
m= 4 T=   81 band= 2 | band-only: eta=1.1e-02 margin=4.4e-02 CERT | full: eta=1.5e-02 margin=3.3e-03 no | Gonek band err=1.7e-02
m= 5 T=   96 band= 3 | band-only: eta=1.1e-02 margin=3.1e-02 CERT | full: eta=1.2e-02 margin=1.8e-03 no | Gonek band err=1.0e-02
m= 6 T=  113 band= 4 | band-only: eta=3.1e-02 margin=2.7e-02 no | full: eta=1.8e-02 margin=1.3e-03 no | Gonek band err=2.4e-02
m= 7 T=  133 band= 4 | band-only: eta=3.2e-02 margin=2.0e-02 no | full: eta=3.3e-02 margin=1.2e-03 no | Gonek band err=4.4e-02
m= 8 T=  156 band= 6 | band-only: eta=7.0e-02 margin=1.9e-02 no | full: eta=4.7e-02 margin=8.8e-04 no | Gonek band err=4.0e-02
m= 9 T=  183 band= 8 | band-only: eta=1.5e-01 margin=1.7e-02 no | full: eta=9.3e-02 margin=5.5e-04 no | Gonek band err=8.5e-02
m=10 T=  216 band= 9 | band-only: eta=2.6e-01 margin=8.0e-03 no | full: eta=1.1e-01 margin=3.8e-04 no | Gonek band err=5.1e-02
m=11 T=  254 band=11 | band-only: eta=1.6e-01 margin=3.2e-03 no | full: eta=2.9e-01 margin=2.9e-04 no | Gonek band err=6.7e-02
singular at 11
\end{lstlisting}

\subsection{\texttt{out\_certG.txt}}
\label{app:out-routeB-certG}

\begin{lstlisting}[style=out]
m= 1 T=   50 interior=  4 L=4.16e+01 H=1.23e+01 eta_sel=1.31e-03 eta_all=1.31e-03 A=2L^2H eta=5.61e+01  2L*eta=1.09e-01
m= 2 T=   59 interior=  5 L=2.83e+01 H=1.42e+01 eta_sel=1.99e-03 eta_all=2.12e-03 A=2L^2H eta=4.52e+01  2L*eta=1.13e-01
m= 3 T=   69 interior=  7 L=4.01e+01 H=1.96e+01 eta_sel=2.48e-03 eta_all=2.53e-03 A=2L^2H eta=1.56e+02  2L*eta=1.99e-01
m= 4 T=   81 interior=  9 L=5.63e+01 H=2.26e+01 eta_sel=2.81e-03 eta_all=3.19e-03 A=2L^2H eta=4.03e+02  2L*eta=3.16e-01
m= 5 T=   96 interior= 12 L=6.77e+01 H=2.80e+01 eta_sel=4.47e-03 eta_all=4.47e-03 A=2L^2H eta=1.15e+03  2L*eta=6.05e-01
m= 6 T=  113 interior= 16 L=9.78e+01 H=3.32e+01 eta_sel=4.61e-03 eta_all=5.00e-03 A=2L^2H eta=2.93e+03  2L*eta=9.02e-01
m= 7 T=  133 interior= 20 L=1.41e+02 H=3.97e+01 eta_sel=4.56e-03 eta_all=6.93e-03 A=2L^2H eta=7.17e+03  2L*eta=1.28e+00
m= 8 T=  156 interior= 26 L=1.84e+02 H=5.16e+01 eta_sel=8.89e-03 eta_all=8.89e-03 A=2L^2H eta=3.09e+04  2L*eta=3.26e+00
m= 9 T=  183 interior= 34 L=2.60e+02 H=6.45e+01 eta_sel=9.13e-03 eta_all=1.01e-02 A=2L^2H eta=7.97e+04  2L*eta=4.75e+00
m=10 T=  216 interior= 43 L=3.59e+02 H=8.53e+01 eta_sel=1.45e-02 eta_all=1.45e-02 A=2L^2H eta=3.17e+05  2L*eta=1.04e+01
m=11 T=  254 interior= 54 L=4.69e+02 H=1.02e+02 eta_sel=1.51e-02 eta_all=1.83e-02 A=2L^2H eta=6.80e+05  2L*eta=1.41e+01
m=12 T=  299 interior= 68 L=7.53e+02 H=1.29e+02 eta_sel=1.66e-02 eta_all=2.11e-02 A=2L^2H eta=2.43e+06  2L*eta=2.50e+01
m=13 T=  352 interior= 85 L=6.45e+02 H=1.61e+02 eta_sel=2.36e-02 eta_all=2.61e-02 A=2L^2H eta=3.16e+06  2L*eta=3.04e+01
m=14 T=  414 interior=107 L=1.02e+03 H=2.03e+02 eta_sel=2.83e-02 eta_all=2.83e-02 A=2L^2H eta=1.20e+07  2L*eta=5.79e+01
\end{lstlisting}

\subsection{\texttt{out\_cert3.txt}}
\label{app:out-routeB-cert3}

\begin{lstlisting}[style=out]
m= 1 T=   50 n=  4 tests=  42 beta=  8.38 K=6.36e+00 delta=8.85e-03 h=4.72e-01 |kG-g|=8.75e-03 after1GN=4.77e-04
m= 2 T=   59 n=  5 tests=  96 beta=  5.19 K=9.92e+00 delta=1.00e-02 h=5.16e-01 |kG-g|=1.02e-02 after1GN=3.30e-04
m= 3 T=   69 n=  7 tests= 164 beta=  6.83 K=1.31e+01 delta=1.27e-02 h=1.13e+00 |kG-g|=1.18e-02 after1GN=1.65e-03
m= 4 T=   81 n=  9 tests= 248 beta=  6.99 K=1.59e+01 delta=2.03e-02 h=2.25e+00 |kG-g|=2.26e-02 after1GN=2.98e-03
m= 5 T=   96 n= 12 tests= 353 beta=  7.16 K=1.87e+01 delta=2.47e-02 h=3.31e+00 |kG-g|=2.50e-02 after1GN=2.10e-03
m= 6 T=  113 n= 16 tests= 482 beta=  8.00 K=2.12e+01 delta=4.45e-02 h=7.55e+00 |kG-g|=4.51e-02 after1GN=7.97e-04
m= 7 T=  133 n= 20 tests= 641 beta=  8.18 K=2.36e+01 delta=7.09e-02 h=1.37e+01 |kG-g|=7.14e-02 after1GN=9.39e-04
m= 8 T=  156 n= 26 tests= 836 beta=  8.06 K=2.59e+01 delta=9.47e-02 h=1.98e+01 |kG-g|=9.58e-02 after1GN=5.21e-03
m= 9 T=  183 n= 34 tests=1074 beta=  9.83 K=2.80e+01 delta=1.60e-01 h=4.42e+01 |kG-g|=1.53e-01 after1GN=4.14e-02
m=10 T=  216 n= 43 tests=1364 beta= 10.26 K=3.00e+01 delta=1.72e-01 h=5.28e+01 |kG-g|=1.80e-01 after1GN=1.60e-02
m=11 T=  254 n= 54 tests=1717 beta= 17.15 K=3.18e+01 delta=2.33e-01 h=1.27e+02 |kG-g|=2.43e-01 after1GN=2.26e-02
m=12 T=  299 n= 68 tests=2146 beta= 15.39 K=3.36e+01 delta=2.66e-01 h=1.37e+02 |kG-g|=2.79e-01 after1GN=2.68e-02
m=13 T=  352 n= 85 tests=2667 beta= 10.99 K=3.52e+01 delta=3.38e-01 h=1.31e+02 |kG-g|=3.55e-01 after1GN=3.56e-02
m=14 T=  414 n=107 tests=3298 beta= 15.67 K=3.67e+01 delta=4.05e-01 h=2.33e+02 |kG-g|=4.27e-01 after1GN=5.19e-02
\end{lstlisting}

\subsection{\texttt{out\_gonek.txt}}
\label{app:out-routeB-gonek}

\begin{lstlisting}[style=out]
clock : mean|.-g|/sp=0.187 max=0.506  #>0.25sp=50
Gonek X= 3 (offset -3): mean|.-g|/sp=0.048 max=0.184 #>0.25sp=0  k=34: 110.923 (g=111.030)
Gonek X= 5 (offset -3): mean|.-g|/sp=0.029 max=0.081 #>0.25sp=0  k=34: 111.034 (g=111.030)
Gonek X= 8 (offset -3): mean|.-g|/sp=0.019 max=0.053 #>0.25sp=0  k=34: 111.071 (g=111.030)
Gonek X=12 (offset -3): mean|.-g|/sp=0.014 max=0.047 #>0.25sp=0  k=34: 111.092 (g=111.030)
Gonek X=20 (offset -3): mean|.-g|/sp=0.008 max=0.035 #>0.25sp=0  k=34: 111.075 (g=111.030)
\end{lstlisting}

\subsection{\texttt{out\_laplace\_check.txt}}
\label{app:out-routeB-laplace-check}

\begin{lstlisting}[style=out]
u=0.1  k_v18=-1.6225522892294519278916331593609088459137814857449  k_mine=-1.6225522892294519278916331593609088459137814857449  K|=-1.6225522892294519278916331593609088459137814857449
u=1.0  k_v18=0.77689426079325358170733757683887017228268849806797  k_mine=0.77689426079325358170733757683887017228268849806797  K|=0.77689426079325358170733757683887017228268849806797
u=2.5  k_v18=1.7441997038954705252460103768783663870122500505333  k_mine=1.7441997038954705252460103768783663870122500505333  K|=1.7441997038954705252460103768783663870122500505333
s=2.0  (1/2)G'/G=0.31951361231226674131251513964113471432763646278619  L[K]=0.31951361231226674131251526910712335613061628485609
s=1.3  (1/2)G'/G=1.4225122146003985739483096016703817041300334693697  L[K]=1.4225122146003985739483097311363703459957308649348
s=(2.0 + 5.0j)  (1/2)G'/G=(-0.0000044839256965157547625412009849844988519201439850855 + 0.16033845735403173017774194624597807348599151073349j)  L[K]=(-0.000004483925696515754762411734996342695870171439857178 + 0.1603384573540317301777419462459780730380104234455j)
s=(1.5 - 3.0j)  (1/2)G'/G=(-0.092240558736457654753515338039303026644761224786081 - 0.054136617330789788303900306539987383185752403110009j)  L[K]=(-0.092240558736457654753515208573314384796983337328152 - 0.054136617330789788303900306539987382916962760406523j)
\end{lstlisting}

\section{Endpoint-only closure, interpolation tests, and autonomous residual questions}\label{sec:appendix-endpoint-v18}
\label{sec:v18-addendum}

This addendum records the outcome of the subsequent interpolation and finite-section tests.  It does not alter any theorem above.  Its purpose is to separate the statements that are exact but equivalent to the final step from the one clock-only subproblem that remains autonomous.

\subsection{Three interpolations and their status}\label{subsec:three-interpolations-v18}

There are three natural ways to interpolate between the midpoint clock and the Riemann zeta function.

\emph{The ordinate interpolation.}  The family $\zeta_\lambda$ of Sections~\CGref{sec:lambda-family}--\CGref{sec:criteria} moves the ordinates while keeping every nontrivial zero on $\Re s=\frac12$.  It is an exact analytic family, but for every $0<\lambda\le1$ its continuous prime measure is not in the ordinary Dirichlet category.  Thus the property that is needed at the endpoint is not preserved along the path; this is exactly the role of condition (D).

\emph{The exact-atomic interpolation.}  One may instead ask for real ordinates $\kappa_k$ and masses $\mathsf a_n$ supported at $u=\log n$ throughout an interpolation.  Finite sections of this problem were tested numerically after v16.  The endpoints can be resolved, whereas the linear interpolation of the target distribution develops a persistent interior residual in the tested sections.  This is evidence against that particular interpolation, not a proof that no infinite exact-atomic path exists.  In any case the existence of such a path is not supplied by the results of this paper.

\emph{The additive interpolation of functions.}  There is an exact elementary interpolation
\begin{equation}\label{eq:additive-family-v18}
 \mathcal G_\varepsilon(s):=(1-\varepsilon)\zeta_{\RVM}(s)+\varepsilon\zeta(s),\qquad 0\le\varepsilon\le1.
\end{equation}
Both summands satisfy the same functional equation and reality symmetry, so the same is true of $\mathcal G_\varepsilon$.  This family exists for every $\varepsilon$, but it does not create a new control quantity: it is a level-set representation of the already defined quotient $\mathcal Q_{\RVM}=\zeta/\zeta_{\RVM}$.

\begin{proposition}[The additive family is a level-set reformulation]\label{prop:additive-level-v18}
Away from the zeros of $\zeta_{\RVM}$,
\begin{equation}\label{eq:additive-level-v18}
 \mathcal G_\varepsilon(s)=\zeta_{\RVM}(s)\bigl((1-\varepsilon)+\varepsilon\mathcal Q_{\RVM}(s)\bigr).
\end{equation}
Hence for $0<\varepsilon<1$ every zero of $\mathcal G_\varepsilon$ which is not a zero of $\zeta_{\RVM}$ lies on the level set
\begin{equation}\label{eq:Q-negative-level-v18}
 \mathcal Q_{\RVM}(s)=-\frac{1-\varepsilon}{\varepsilon}.
\end{equation}
On the critical line $e^{i\vartheta(t)}\mathcal G_\varepsilon(\frac12+it)$ is real.  At $\varepsilon=1$ the statement that all its nontrivial zeros lie on the critical line is exactly RH.
\end{proposition}

\begin{proof}
Equation~\CGeqref{eq:additive-level-v18} is the definition of $\mathcal Q_{\RVM}$.  The functional equation and the reality statement follow by linearity from those of the two summands.  The last assertion is $\mathcal G_1=\zeta$.
\end{proof}

Thus the additive family is useful for visualizing the zero motion, but without an additional monotonicity or flow law it is circular as a route to RH: the terminal level is precisely the original zero problem.

\subsection{The clock as the first transition term}\label{subsec:first-transition-v18}

The transition theorem gives the following $1+\chi$ form on every closed subband of $\tfrac12<\Re s<1$:
\begin{equation}\label{eq:rvm-first-rs-v18}
 \zeta_{\RVM}(s)=\varpi_{\RVM}(1+\chi(s))\bigl(1+O(|t|^{-1})\bigr).
\end{equation}
Independently, on $s=\frac12+it$, where $\chi(s)=e^{-2i\vartheta(t)}$, the exact trigonometric identity
\begin{equation}\label{eq:first-rs-hardy-v18}
 e^{i\vartheta(t)}(1+\chi(\tfrac12+it))=2\cos\vartheta(t).
\end{equation}
Consequently the midpoint clock is asymptotically the zero set selected by the $n=1$ term of the Hardy--Littlewood/Riemann--Siegel approximate functional equation.  This is a structural interpretation of the transition theorem, not an exact identification: the $O(1/t)$ term is visible already in the first ordinates.

For a finite Dirichlet polynomial
\begin{equation}\label{eq:finite-additive-v18}
 \mathsf P_{\varepsilon,M}(s):=1+\varepsilon\sum_{2\le n\le M}n^{-s},\qquad
 \mathsf R_{\varepsilon,M}(s):=\mathsf P_{\varepsilon,M}(s)+\chi(s)\mathsf P_{\varepsilon,M}(1-s),
\end{equation}
the functional equation is exact.  On the critical line,
\begin{equation}\label{eq:finite-additive-hardy-v18}
 e^{i\vartheta(t)}\mathsf R_{\varepsilon,M}(\tfrac12+it)
 =2\Re\!\left(e^{i\vartheta(t)}\mathsf P_{\varepsilon,M}(\tfrac12+it)\right).
\end{equation}
At Gram points the $n=1$ term fixes the sign whenever
\begin{equation}\label{eq:dominance-v18}
 \varepsilon\sum_{2\le n\le M}n^{-1/2}<1,
\end{equation}
so in every interval where $\vartheta$ is monotone there is a critical-line zero between consecutive Gram points.  This dominance statement is only sufficient; it does not count all zeros of $\mathsf R_{\varepsilon,M}$ and is not used elsewhere in the paper.

\subsection{What remains genuinely live}\label{subsec:live-v18}

After these tests the status is narrower than the programme suggested by the interpolation pictures.

\begin{enumerate}[label=\textup{(L\arabic*)}]
\item \emph{The endpoint problem remains the endpoint problem.}  Proving that the cell defect vanishes at $\lambda=0$, equivalently proving (D), remains the only step that closes the main route.  The additive family and the exact-atomic interpolation do not simplify this statement.

\item \emph{The clock-only injectivity problem is autonomous.}  The one subproblem that does not contain the genuine ordinates is the injectivity of the linearized rung system at the midpoint clock.  In the meromorphic form of Remark~\CGref{rem:nondegenerate-scope}, the target statement is: if a bounded real sequence $(\dot\kappa_k)$ and coefficients $\dot{\mathsf a}_n=O_\eps(n^{-1/2+\eps})$ give a Dirichlet series $\mathsf H$ satisfying
\begin{equation}\label{eq:clock-injectivity-v18}
 \mathsf H(z)=\nu_8+\sum_{k\ge1}\dot\kappa_k\left(\frac{\tau_k}{z^2+\tau_k^2}-\frac1{\tau_k}\right)
     =\sum_{n\ge2}\frac{\dot{\mathsf a}_n}{\log n}\,n^{-z}
\end{equation}
through a common right half-plane, under the growth hypotheses of the linearization, then $\mathsf H\equiv0$.  Theorem~\CGref{thm:nondegenerate} proves the analogous conclusion at an actual solution, while Remark~\CGref{rem:nondegenerate-scope} explains why it is not automatic for a general list.  Equation~\CGeqref{eq:clock-injectivity-v18} is therefore a closed clock-only rigidity problem, naturally related to Fourier crystalline uniqueness, but by itself it does not prove (D).

\item \emph{Bounded finite sections are a numerical existence programme, not yet an analytic route.}  Direct finite sections of the endpoint rung system can be solved while constraining both the ordinate displacement and the free masses.  If one could prove, for an exhausting nested family of tests, a uniform boundedness and residual-to-zero statement and then pass to a distributional limit, Theorem~\CGref{thm:free-masses} would identify the limit.  Such a theorem would already be an existence theorem for the final system and is therefore not a shortcut around (D).
\end{enumerate}

In particular, the individual displacements $\gamma_k-\tau_k$ are not autonomous data.  They are used below only for post-solution numerical diagnostics.  None of the live statement \CGeqref{eq:clock-injectivity-v18}, nor the definition of the bounded finite-section problem, requires them as input.


\subsection{A priori bounded displacement for exact rung solutions}\label{subsec:rung-apriori-v18}

The following statement is independent of the existence of an interior exact-atomic interpolation.  It says only that, if an exact real solution exists, has the stated local density and masses of the stated mean square, then its ordinates follow the midpoint clock up to a linear drift governed by one constant $\mathrm b$, and stay within a bounded distance of it exactly when $\mathrm b=0$.

\begin{proposition}[A priori displacement bound for an exact rung solution]\label{prop:rung-apriori-v18}
Fix $\mathsf t\in[0,1]$.  Let $\boldsymbol{\kappa}=(\kappa_k)_{k\ge1}$ satisfy $0<\kappa_1\le\kappa_2\le\cdots$ and $\kappa_k\to\infty$, and let $(\mathsf a_n)_{n\ge2}$ be real masses satisfying, in $\mathcal D'(0,\infty)$,
\begin{equation}\label{eq:rung-apriori-system-v18}
 \sum_{k\ge1}\cos(\kappa_k u)
 -(1-\mathsf t)\sum_{k\ge1}\cos(\tau_k u)-\mathsf t\,\mathsf k(u)
 =-\sum_{n\ge2}\mathsf a_n\,\delta(u-\log n).
\end{equation}
Assume that for every $\epsilon_0>0$ the local count and the masses satisfy
\begin{equation}\label{eq:rung-apriori-density-v18}
 N_{\boldsymbol{\kappa}}(T+1)-N_{\boldsymbol{\kappa}}(T)\ll_{\epsilon_0}(1+T)^{\epsilon_0}\qquad(T\ge0),
\end{equation}
\begin{equation}\label{eq:rung-apriori-mass-v18}
 \sum_{X<n\le2X}|\mathsf a_n|^2\ll_{\epsilon_0}X^{1+\epsilon_0}\qquad(X\ge1).
\end{equation}
Then the limit
\begin{equation}\label{eq:rung-apriori-b-v18}
 \mathrm b:=\lim_{\sigma\to+\infty}\Bigl[\sum_{k\ge1}\frac{\sigma-\frac12}{(\sigma-\frac12)^2+\kappa_k^2}-(1-\mathsf t)\sum_{k\ge1}\frac{\sigma-\frac12}{(\sigma-\frac12)^2+\tau_k^2}-\mathsf t\,\frac{G'}{2G}(\sigma)\Bigr]
\end{equation}
exists and is real, and there is a constant $C^*$, depending only on $\mathsf t$-independent data, namely the implied constants in \CGeqref{eq:rung-apriori-density-v18} and \CGeqref{eq:rung-apriori-mass-v18} for one fixed $\epsilon_0\in(0,\frac14)$, the number $\sum_k\kappa_k^{-2}$ and the number $|\mathrm b|$, such that
\begin{equation}\label{eq:rung-apriori-count-v18}
 \Bigl|N_{\boldsymbol{\kappa}}(T)-n_{\RVM}(T)-\frac{2\mathrm b}{\pi}T\Bigr|\le C^*\log(2+T)\qquad(T\ge0)
\end{equation}
and
\begin{equation}\label{eq:rung-apriori-bound-v18}
 \Bigl|\kappa_k-\tau_k+\frac{4\mathrm b\,\kappa_k}{\log(\kappa_k/2\pi)}\Bigr|\le C^*\qquad(k\ge1).
\end{equation}
In particular the displacement $\kappa_k-\tau_k$ is bounded if and only if $\mathrm b=0$, and then $|\kappa_k-\tau_k|\le C^*$ for all $k$, uniformly in $\mathsf t$ and over every class of solutions in which the four quantities above are bounded. The constant $\mathrm b$ is the mass at $u=0$ of the distribution on $[0,\infty)$ that extends the left side of \CGeqref{eq:rung-apriori-system-v18} with the extensions of the proof, and the identity on the open half-line does not determine it: a list whose count exceeds the clock by $2\mathrm b/\pi$ per unit height satisfies \CGeqref{eq:rung-apriori-b-v18} with that very $\mathrm b$.
\end{proposition}

\begin{proof}
\emph{Order of $P_{\boldsymbol{\kappa}}$.} The local-density hypothesis with any fixed $\epsilon_0<1$ gives $N_{\boldsymbol{\kappa}}(T)=O(T^{1+\epsilon_0})$; hence $\sum_k\kappa_k^{-2}<\infty$, the symmetric canonical product $P_{\boldsymbol{\kappa}}(s)=\prod_k(1+q(s)/(\frac14+\kappa_k^2))$ converges locally uniformly, is entire and real on $\mathbb R$, is positive for real $s>1$ (every factor exceeds $1$ there, since $q(s)>0$), and has order at most $1$.  The last point is proved directly, not through the convergence exponent: for $|s|\ge1$, $|q(s)|\le2|s|^2=:R^2$ and $\log|P_{\boldsymbol{\kappa}}(s)|\le\sum_k\log\bigl(1+R^2/\kappa_k^2\bigr)=\int_0^\infty\log(1+R^2/t^2)\,dN_{\boldsymbol{\kappa}}(t)=\int_0^\infty N_{\boldsymbol{\kappa}}(t)\frac{2R^2}{t(t^2+R^2)}\,dt$ (the boundary term $N_{\boldsymbol{\kappa}}(t)\log(1+R^2/t^2)\ll t^{1+\epsilon_0}R^2t^{-2}$ vanishes at $t\to\infty$ and at $t\to0$); with $N_{\boldsymbol{\kappa}}(t)\ll t^{1+\epsilon_0}$ and $t=Rv$ this is $\ll R^{1+\epsilon_0}\int_0^\infty\frac{2v^{\epsilon_0}}{1+v^2}dv\ll|s|^{1+\epsilon_0}$.  Since the local-density hypothesis holds for every $\epsilon_0>0$, $\log|P_{\boldsymbol{\kappa}}(s)|=O_{\epsilon}(|s|^{1+\epsilon})$ for every $\epsilon>0$, which is order at most $1$.

\emph{From $\mathcal D'(0,\infty)$ to the Laplace transform.} The identity \CGeqref{eq:rung-apriori-system-v18} holds on the open half-line and says nothing about the point $u=0$; the exponential $e^{-(s-1/2)u}$ is not a test function of $\mathcal D(0,\infty)$, so it cannot be paired with \CGeqref{eq:rung-apriori-system-v18} directly. We first extend each term to a distribution on $\mathbb R$ supported in $[0,\infty)$, as in Lemma~\CGref{lem:rb-laplace}. For a list with $\sum_k\kappa_k^{-2}<\infty$ put $\mathfrak H_{\boldsymbol{\kappa}}(u):=\sum_k\kappa_k^{-2}\cos(\kappa_ku)$, continuous and bounded, and $\mathscr S_{\boldsymbol{\kappa}}:=-\partial_u^2(\mathfrak H_{\boldsymbol{\kappa}}\mathbf 1_+)+\mathfrak H_{\boldsymbol{\kappa}}(0)\delta_0'$: a distribution of order at most $2$, supported in $[0,\infty)$, whose restriction to $(0,\infty)$ is $\sum_k\cos(\kappa_ku)$ (termwise differentiation of a uniformly convergent series in $\mathcal D'(0,\infty)$) and whose Laplace transform is, by the computation \CGeqref{eq:rb-aux-14}, which uses only $\sum_k\kappa_k^{-2}<\infty$,
\begin{equation}\label{eq:rung-laplace-S-v18}
 \mathcal L\mathscr S_{\boldsymbol{\kappa}}(s)=\sum_{k\ge1}\frac{s-\frac12}{(s-\frac12)^2+\kappa_k^2}=\tfrac12\,\frac{P_{\boldsymbol{\kappa}}'}{P_{\boldsymbol{\kappa}}}(s)\qquad(\Re s>\tfrac12).
\end{equation}
The same extension $\mathscr S_{\boldsymbol\tau}$ serves the clock, with $P^{\crit}_{\RVM}$ in place of $P_{\boldsymbol{\kappa}}$. The term $\mathsf k$ is extended by the Archimedean distribution $\mathscr K$ of Lemma~\CGref{lem:rb-archimedean}, with $\mathcal L\mathscr K=\frac12G'/G$. The mass term $\mathscr A:=\sum_{n\ge2}\mathsf a_n\delta_{\log n}$ is a locally finite measure on $\mathbb R$ supported in $[\log2,\infty)$, with no mass at $0$; by \CGeqref{eq:rung-apriori-mass-v18} and Cauchy--Schwarz on dyadic blocks, $\sum_n|\mathsf a_n|n^{1/2-\sigma}<\infty$ for every $\sigma>\frac32$, so $\mathcal L\mathscr A(s)=\sum_{n\ge2}\mathsf a_nn^{\frac12-s}$ converges absolutely on $\Re s>\frac32$. The distribution
\begin{equation}\label{eq:rung-E-v18}
 \mathscr E:=\mathscr S_{\boldsymbol{\kappa}}-(1-\mathsf t)\mathscr S_{\boldsymbol\tau}-\mathsf t\,\mathscr K+\mathscr A
\end{equation}
has finite order and support in $[0,\infty)$, and by \CGeqref{eq:rung-apriori-system-v18} its restriction to $(0,\infty)$ vanishes. Hence $\mathscr E$ is supported at $\{0\}$, so it is a finite combination $\sum_{j\le J}c_j\delta_0^{(j)}$ with real $c_j$ \cite[Thm.~2.3.4]{Hormander1990}, whose Laplace transform is the real polynomial $p(z)=\sum_{j\le J}c_jz^j$, $z=s-\frac12$. Taking Laplace transforms of \CGeqref{eq:rung-E-v18} on $\Re s>\frac32$,
\begin{equation}\label{eq:rung-laplace-identity-v18}
 \tfrac12\frac{P_{\boldsymbol{\kappa}}'}{P_{\boldsymbol{\kappa}}}(s)-(1-\mathsf t)\,\tfrac12\frac{(P^{\crit}_{\RVM})'}{P^{\crit}_{\RVM}}(s)-\mathsf t\,\tfrac12\frac{G'}{G}(s)+\sum_{n\ge2}\mathsf a_nn^{\frac12-s}=p\bigl(s-\tfrac12\bigr)\qquad(\Re s>\tfrac32).
\end{equation}
The polynomial $p$ is the price of the passage from the open half-line to the closed one (Remark~\CGref{rem:rb-boundary-term}).

\emph{The polynomial is a constant.} On the real axis, as $\sigma\to+\infty$, each term on the left of \CGeqref{eq:rung-laplace-identity-v18} is $O(\sigma^{\epsilon_0})$ with the fixed $\epsilon_0<\frac14$: integrating by parts against the count (the boundary term $N_{\boldsymbol{\kappa}}(t)\,z/(z^2+t^2)$ vanishes at $t=0$ and, since $N_{\boldsymbol{\kappa}}(t)\ll t^{1+\epsilon_0}$, at $t\to\infty$), $\sum_k\frac{z}{z^2+\kappa_k^2}=\int_0^\infty N_{\boldsymbol{\kappa}}(t)\frac{2zt}{(z^2+t^2)^2}\,dt\ll\int_0^\infty t^{1+\epsilon_0}\frac{zt\,dt}{(z^2+t^2)^2}=z^{\epsilon_0}\int_0^\infty\frac{v^{2+\epsilon_0}\,dv}{(1+v^2)^2}\ll z^{\epsilon_0}$, with the implied constant that of \CGeqref{eq:rung-apriori-density-v18}; the clock term is $O(\log\sigma)$ by the same computation with $n_{\RVM}$; $\frac12G'/G(\sigma)=O(\log\sigma)$ by Stirling's formula; and the mass sum tends to $0$ by absolute convergence. A polynomial that is $O(\sigma^{\epsilon_0})$ on the real axis with $\epsilon_0<1$ has degree zero, so $p\equiv\mathrm b$ with $\mathrm b\in\mathbb R$, and \CGeqref{eq:rung-laplace-identity-v18} read at real $\sigma\to+\infty$ is \CGeqref{eq:rung-apriori-b-v18}: the limit exists and equals the constant $\mathrm b$.

\emph{Integration.} All terms of \CGeqref{eq:rung-laplace-identity-v18} are real on the real axis $\sigma>\frac32$, where $P_{\boldsymbol{\kappa}}$, $P^{\crit}_{\RVM}$ and $G$ are positive. Integrating from $2$ to $s$ along a path in $\Re s>\frac32$ and exponentiating,
\begin{equation}\label{eq:rung-three-factor-v18}
 P_{\boldsymbol{\kappa}}(s)
 =C_\mathsf t\bigl(P_{\RVM}^{\crit}(s)\bigr)^{1-\mathsf t}\,G(s)^{\mathsf t}\,Y_\mathsf t(s)\,e^{2\mathrm b s}
 =C_\mathsf t\,G(s)\,\zeta_{\RVM}(s)^{1-\mathsf t}\,Y_\mathsf t(s)\,e^{2\mathrm b s}\qquad(\Re s>\tfrac32),
\end{equation}
where
\begin{equation}\label{eq:rung-Y-v18}
 Y_\mathsf t(s)=\exp\!\left(\sum_{n\ge2}\frac{2\mathsf a_n}{\log n}\,n^{1/2-s}\right),
\end{equation}
$C_\mathsf t>0$ is the normalization constant fixed by the values at $s=2$.  The fractional powers are defined globally, not along paths: the half-plane $\Re s>\frac32$ is simply connected, $P^{\crit}_{\RVM}$ and $G$ have no zeros there (the zeros of $P^{\crit}_{\RVM}$ lie on $\Re s=\frac12$, $G$ has none) and both are positive on $(\frac32,\infty)$; so each has a unique holomorphic logarithm on $\Re s>\frac32$ that is real on the real axis, and $(P^{\crit}_{\RVM})^{1-\mathsf t}:=\exp((1-\mathsf t)\log P^{\crit}_{\RVM})$, $G^{\mathsf t}:=\exp(\mathsf t\log G)$ with these logarithms.  The integration from $2$ to $s$ is then path-independent, and the zero count below is performed on the entire function $P_{\boldsymbol{\kappa}}$ itself, so no branch of a fractional factor enters it. By \CGeqref{eq:rung-apriori-mass-v18} and Cauchy--Schwarz on dyadic blocks, the Dirichlet series in \CGeqref{eq:rung-Y-v18} converges absolutely and uniformly on every half-plane $\Re s\ge\sigma_0>\frac32$: for $X<n\le2X$,
\[
 \sum_{X<n\le2X}|\mathsf a_n|n^{1/2-\sigma_0}
 \ll_{\epsilon_0}X^{3/2-\sigma_0+\epsilon_0/2},
\]
so with $\sigma_0:=2$ and $\epsilon_0<\frac14$ the dyadic sum converges and
\begin{equation}\label{eq:rung-B-v18}
 \sup_{\Re s\ge2}|\log Y_\mathsf t(s)|\le B_2<\infty,
\end{equation}
with $B_2$ depending only on the implied constant in \CGeqref{eq:rung-apriori-mass-v18}. The constant $C_\mathsf t$ is fixed by \CGeqref{eq:rung-three-factor-v18} at $s=2$, where $P_{\boldsymbol{\kappa}}(2)=\prod_k(1+2/(\frac14+\kappa_k^2))\in[1,e^{2\sum_k\kappa_k^{-2}}]$, $G(2)=2/\pi$ and $\zeta_{\RVM}(2)$ is a fixed positive number; so $|\log C_\mathsf t|$ is bounded in terms of $\sum_k\kappa_k^{-2}$, $B_2$ and $|\mathrm b|$. No displacement bound is used here.

\emph{Growth on the right boundary.} The quotient $P_{\boldsymbol{\kappa}}/G$ has a simple pole at $s=1$, so the pole-cleared function
\[
 \mathcal Z_{\boldsymbol{\kappa}}(s):=(s-1)\frac{P_{\boldsymbol{\kappa}}(s)}{G(s)}=\frac{P_{\boldsymbol{\kappa}}(s)}{2\Gamma(\frac s2+1)\pi^{-s/2}}
\]
is entire of order one.  On the line $\Re s=2$, \CGeqref{eq:rung-three-factor-v18}, Theorem~\CGref{thm:halfplane-RVM} ($|\log\zeta_{\RVM}|\le\mathrm c_{\RVM}$ on $\Re s\ge2$) and \CGeqref{eq:rung-B-v18} give
\begin{equation}\label{eq:rung-line-v18}
 \bigl|\log|P_{\boldsymbol{\kappa}}(s)|-\log|G(s)|\bigr|\le \mathbb K_{\rm E}:=|\log C_\mathsf t|+\mathrm c_{\RVM}+B_2+8|\mathrm b|\qquad(2\le\Re s\le4),
\end{equation}
uniformly in $\mathsf t$ and $t$; the factor $e^{2\mathrm bs}$ contributes $2\mathrm b\Re s$ to $\log|P_{\boldsymbol{\kappa}}|$, which is $4\mathrm b$ on the line $\Re s=2$ and at most $8|\mathrm b|$ in modulus on $0\le\Re s\le4$, the strip in which the Jensen disc below lies; this is why $\mathbb K_{\rm E}$ carries $8|\mathrm b|$ rather than $4|\mathrm b|$.  On the line $\Re s=2$ itself, by the symmetry $P_{\boldsymbol{\kappa}}(1-s)=P_{\boldsymbol{\kappa}}(s)$, $G(1-s)=G(s)$ and the reality of both functions on $\mathbb R$, the same bound holds on $\Re s=-1$. The function $u(s):=\log|\mathcal Z_{\boldsymbol{\kappa}}(s)/(s+2)|$ is subharmonic on the strip $-1\le\Re s\le2$, at most $\mathbb K_{\rm E}+\log2$ on its two boundary lines (there $|s-1|\le|s+2|$ on $\Re s=2$ and $|s-1|\le2|s+2|$ on $\Re s=-1$), and, $\mathcal Z_{\boldsymbol{\kappa}}$ being of order one, $u(s)\le|s|^2\le e^{|\Im s|/2}$ for $|s|$ large: this is the growth condition of Lemma~\CGref{lem:rb-PL}(a) with $c=\frac12<\pi/3=\pi/a$ for the strip of width $a=3$ (the lemma requires order-one growth only in this weak exponential form, so no estimate of the type is needed).  The division by $s+2$ introduces no singularity, since $s+2\ne0$ on $-1\le\Re s\le2$, and it is what makes the boundary values at most $\mathbb K_{\rm E}+\log2$.  Lemma~\CGref{lem:rb-PL}(a) applied to $\Im s\ge1$ and to $\Im s\le-1$ (with the compact part $|\Im s|\le1$ bounded directly) gives $u\le\mathbb K_{\rm E}+\mathbb K_{\rm E}'$ on the strip, with an absolute $\mathbb K_{\rm E}'$. Hence, on the strip,
\begin{equation}\label{eq:rung-strip-v18}
 \log|P_{\boldsymbol{\kappa}}(s)|\le\mathbb K_{\rm E}+\mathbb K_{\rm E}'+\log|s+2|+\log|G(s)|-\log|s-1| .
\end{equation}

\emph{Horizontal variation of the argument.} Fix $T\ge T_0$ with $T\notin\{\kappa_k\}$, where $T_0\ge3$ is absolute. Let $n_T$ be the number of zeros of $P_{\boldsymbol{\kappa}}$ in the closed disc $|s-(2+iT)|\le\frac32$, which contains the segment from $2+iT$ to $\frac12+iT$. Jensen's formula on the disc of radius $2$ about $c_T:=2+iT$ (contained in the strip $0\le\Re s\le4$, where \CGeqref{eq:rung-strip-v18} applies for $\Re s\le2$ and \CGeqref{eq:rung-line-v18}, valid on $2\le\Re s\le4$ by \CGeqref{eq:rung-three-factor-v18} and \CGeqref{eq:rung-B-v18}, applies for $\Re s\ge2$) reads, with $a_1,a_2,\dots$ the zeros of $P_{\boldsymbol{\kappa}}$ in $|s-c_T|<2$,
\[
 \sum_{|a_j-c_T|<2}\log\frac2{|a_j-c_T|}=\frac1{2\pi}\int_0^{2\pi}\log|P_{\boldsymbol{\kappa}}(c_T+2e^{i\phi})|\,d\phi-\log|P_{\boldsymbol{\kappa}}(c_T)| ;
\]
the $n_T$ zeros with $|a_j-c_T|\le\frac32$ contribute at least $\log\frac43$ each, the others contribute positively, and the mean over the circle is at most the maximum, a weaker but valid bound; hence
\[
 n_T\log\tfrac43\le\max_{|s-(2+iT)|=2}\log|P_{\boldsymbol{\kappa}}(s)|-\log|P_{\boldsymbol{\kappa}}(2+iT)| .
\]
By Stirling's formula the oscillation of $\log|G|$ on the disc is at most $\log(T+8)+2\mathrm c_5$ (Lemma~\CGref{lem:rb-backlund}), $\log|s+2|\le\log(T+6)$ and $-\log|s-1|\le\log2$ there, and the lower bound at the centre is $\log|P_{\boldsymbol{\kappa}}(2+iT)|\ge\log|G(2+iT)|-\mathbb K_{\rm E}$ by \CGeqref{eq:rung-line-v18}. Hence
\begin{equation}\label{eq:rung-nT-v18}
 n_T\le\frac{2\mathbb K_{\rm E}+\mathbb K_{\rm E}'+2\log(T+8)+2\mathrm c_5+\log2}{\log(4/3)},
\end{equation}
and, exactly as in Lemma~\CGref{lem:rb-backlund}, the variation of $\arg P_{\boldsymbol{\kappa}}$ along the segment from $2+iT$ to $\frac12+iT$ is at most $\pi(n_T+1)$, while that of $\arg G$ is at most $\mathrm c_6=\frac32$ for $T\ge T_0$.

\emph{Vertical variation of the argument.} On the segment from $2$ to $2+iT$, by \CGeqref{eq:rung-three-factor-v18},
\[
 \arg\zeta_{\boldsymbol{\kappa}}(2+it)=\arg\frac{P_{\boldsymbol{\kappa}}}{G}(2+it)=(1-\mathsf t)\Im\log\zeta_{\RVM}(2+it)+\Im\log Y_\mathsf t(2+it)+2\mathrm b\,t ,
\]
with the branches chosen continuous along the segment and real at $t=0$ (the branch of $\log\zeta_{\RVM}$ is the one of Theorem~\CGref{thm:halfplane-RVM}, continuous on $\Re s\ge2$, and $\log Y_{\mathsf t}$ is the series \CGeqref{eq:rung-Y-v18} itself).  What the argument principle uses is the \emph{net} variation of the argument along the segment, that is, the difference of the values of the continuous branch at its two ends, not its total variation; the first two terms are bounded by $\mathrm c_{\RVM}$ and $B_2$ in modulus at both ends, so the net variation along the segment is $2\mathrm bT+O(\mathrm c_{\RVM}+B_2)$. This is where the factor $e^{2\mathrm bs}$ enters the count: its argument is not bounded, it grows linearly.

\emph{Zero count.} Lemma~\CGref{lem:rb-argument} applies verbatim to $P_{\boldsymbol{\kappa}}$ (it uses only that $P_{\boldsymbol{\kappa}}$ is real on $\mathbb R$, symmetric under $s\mapsto1-s$, with zeros exactly at $\frac12\pm i\kappa_k$ and none on the contour): $\pi N_{\boldsymbol{\kappa}}(T)=\pi+\vartheta(T)+\Delta_{\mathbb G_T}\arg\zeta_{\boldsymbol{\kappa}}$ along the path $2\to2+iT\to\frac12+iT$. Combining the two variations, and using $1+\vartheta(T)/\pi=\NR(T)+O(1/T)$ (Stirling) and $n_{\RVM}(T)=\lfloor\NR(T)+\frac12\rfloor$,
\begin{equation}\label{eq:rung-count-v18}
 \Bigl|N_{\boldsymbol{\kappa}}(T)-n_{\RVM}(T)-\frac{2\mathrm b}\pi T\Bigr|
 \le n_T+1+\frac{2\mathrm c_{\RVM}+2B_2+\mathrm c_6}{\pi}+\frac1{T}+\frac12\qquad(T\ge T_0,\ T\notin\{\kappa_k\}),
\end{equation}
which with \CGeqref{eq:rung-nT-v18} is \CGeqref{eq:rung-apriori-count-v18} for $T\ge T_0$, with $C^*$ depending only on $\mathbb K_{\rm E}$, that is, on $|\log C_\mathsf t|$, $B_2$ and $|\mathrm b|$, hence on the four quantities of the statement; it extends to $T\notin\{\kappa_k\}$, $T\ge T_0$ by right continuity and to $0\le T<T_0$ by $N_{\boldsymbol{\kappa}}(T_0)\le C^*\log(2+T_0)+n_{\RVM}(T_0)+2|\mathrm b|T_0/\pi$, enlarging $C^*$. The argument uses $P_{\boldsymbol{\kappa}}$, which is entire, so no branch of a fractional factor enters the zero count.

\emph{From the count to the displacement.} At the $k$th ordinate, with multiplicity included, \CGeqref{eq:rung-apriori-count-v18} gives
\begin{equation}\label{eq:rung-quantile-v18}
 \NR(\kappa_k)=k-\frac{2\mathrm b}\pi\kappa_k+O\bigl(C^*\log(2+\kappa_k)\bigr),
\end{equation}
the jump of $N_{\boldsymbol{\kappa}}$ at a multiple ordinate being itself $O(C^*\log(2+\kappa_k))$ by \CGeqref{eq:rung-apriori-count-v18}. Since $\NR(\tau_k)=k-\frac12$ and $\NR(t)\asymp t\log t$ while the right side of \CGeqref{eq:rung-quantile-v18} is $k(1+o(1))$, first $\kappa_k\asymp\tau_k$ for $k\ge k_0$, with $k_0$ depending only on $C^*$ and $|\mathrm b|$; the mean-value theorem on the increasing branch then gives
\[
 (\kappa_k-\tau_k)\,\frac1{2\pi}\log\frac{\mathrm s_k}{2\pi}=-\frac{2\mathrm b}\pi\kappa_k+O\bigl(C^*\log(2+\kappa_k)\bigr),\qquad \mathrm s_k\ \text{between}\ \kappa_k\ \text{and}\ \tau_k ,
\]
and the comparability makes $\log(\mathrm s_k/2\pi)=\log(\kappa_k/2\pi)+O(1)$. This is \CGeqref{eq:rung-apriori-bound-v18} for $k\ge k_0$, after enlarging $C^*$. For $k<k_0$, $0<\kappa_k\le\kappa_{k_0}$ and $\tau_k\le\tau_{k_0}$, so the left side of \CGeqref{eq:rung-apriori-bound-v18} is bounded by $\kappa_{k_0}+\tau_{k_0}+4|\mathrm b|\kappa_{k_0}$, a bound in terms of the same constants; this is how the finitely many initial indices are absorbed into $C^*$ without any further dependence on the solution. All constants depend only on the implied constants of \CGeqref{eq:rung-apriori-density-v18} and \CGeqref{eq:rung-apriori-mass-v18} for the fixed $\epsilon_0$, on $\sum_k\kappa_k^{-2}$ and on $|\mathrm b|$, and not on $\mathsf t$.
\end{proof}

\begin{corollary}[Exact solutions without mass at the origin]\label{cor:rung-apriori-b0}
Under the hypotheses of Proposition~\CGref{prop:rung-apriori-v18}, if the extension of the left side of \CGeqref{eq:rung-apriori-system-v18} constructed in the proof has no mass at $u=0$, that is, $\mathrm b=0$, then $|N_{\boldsymbol{\kappa}}(T)-n_{\RVM}(T)|\le C^*\log(2+T)$ and $|\kappa_k-\tau_k|\le C^*$ for all $k$, with $C^*$ depending only on the implied constants of \CGeqref{eq:rung-apriori-density-v18}--\CGeqref{eq:rung-apriori-mass-v18} and on $\sum_k\kappa_k^{-2}$.
\end{corollary}

The mechanism behind the drift is elementary: a mass $\mathrm b\delta_0$ in the extended identity adds the constant $\mathrm b$ to the Laplace residual, hence the factor $e^{2\mathrm bs}$ to $P_{\boldsymbol{\kappa}}$, whose argument grows linearly along vertical lines; the argument principle converts this into an excess of $2\mathrm bT/\pi$ zeros up to height $T$, and dividing by the clock density $\NR'(\kappa)=\frac1{2\pi}\log\frac\kappa{2\pi}$ turns the excess into a displacement of size $4\mathrm b\kappa/\log(\kappa/2\pi)$.  The mean-square hypothesis \CGeqref{eq:rung-apriori-mass-v18} on the masses is independent of the local-density hypothesis and is not implied by it for a general list; both are assumptions on the solution, and the proposition is an a priori statement about hypothetical solutions, not an existence theorem.

\begin{remark}\label{rem:rung-apriori-scope-v18}
Proposition~\CGref{prop:rung-apriori-v18} is an a priori theorem, not an existence theorem.  The exact-atomic interpolation has no established interior solution for $0<\theta<1$; the proposition only says that any exact solution satisfying its hypotheses lies in one uniform bounded-displacement class. The mean-square bound \CGeqref{eq:rung-apriori-mass-v18} on the masses is a hypothesis: for the lists $(\gamma_k(\lambda))$ it is Theorem~\CGref{thm:mean-square-cells}, whose proof uses the tube through Lemma~\CGref{lem:mean-value-ordinates} and the term \CGeqref{eq:cell-mass-power-two}; for a general list satisfying only \CGeqref{eq:rung-apriori-density-v18} it is not derived here, and the proposition assumes it. The constant $\mathrm b$ of \CGeqref{eq:rung-apriori-b-v18} is the mass at $u=0$ that the identity on the open half-line leaves undetermined; its argument along $\Re s=2$ grows linearly and produces the drift $-4\mathrm b\kappa_k/\log(\kappa_k/2\pi)$ in \CGeqref{eq:rung-apriori-bound-v18}, so the proposition bounds the displacement only for solutions with $\mathrm b=0$, that is, for solutions of the identity extended to $[0,\infty)$ without a mass at the origin. The zero count follows the mechanism of Lemmas~\CGref{lem:rb-backlund} and~\CGref{lem:rb-argument}: Backlund--Jensen on a horizontal segment, the argument principle on the symmetric rectangle, and Stirling for $\vartheta$; a disc of fixed radius about $2+iT$ controls only the horizontal piece of the argument variation, and the count up to $T$ comes from the vertical piece and the contour.
\end{remark}

\subsection{Additional finite-section diagnostics}\label{subsec:finite-diagnostics-v18}

The post-v16 calculations solve bounded finite sections directly from the midpoint clock.  Here the normalized residual is the Euclidean residual divided by the square root of the number of test equations; a test bump is a fixed $C_c^\infty$ weight in the $u$-variable; and a permanent block is a set of test equations that, once introduced, is retained in every later nested solve.  The repository scripts and their console/CSV outputs are included with version~18.  A representative run imposes
\begin{equation}\label{eq:bounded-section-v18}
 |\kappa_k-\tau_k|\le1.2,\qquad |a_n|\le0.5,
\end{equation}
with the unresolved ordinate tail kept on the midpoint clock.  The bounds in \CGeqref{eq:bounded-section-v18} are numerical regularizations, not hypotheses of any theorem.  The genuine ordinates and the von Mangoldt masses are read only after the solve for diagnostics.

With free masses and the bounds \CGeqref{eq:bounded-section-v18}, the sections $N=20,34,50$ have normalized residuals of order $10^{-4}$, and the low ordinates and low masses remain close to their diagnostic zeta/von-Mangoldt values.  Cross-validation with test bumps not used in the solve gives residuals of the same order.  These endpoint calculations are distinct from the linearized clock problem: in the latter, it is the ordinate-direction coefficients $h_k$ that become large while a residual of roughly $15$--$18\%$ persists.  Imposing bounds there does not remove an artifact; it exposes the lack of an exact solution to that tested linear problem.

A second test uses a nested family of permanent test blocks: once a test block is introduced it remains in all later solves.  For the bounded sections $N=50,70,100$, the RMS residuals on the old blocks stay small and in the $N=100$ section are between $2.5\cdot10^{-8}$ and $6.1\cdot10^{-8}$ on the blocks $N=13,20,34,50,70$, while the newest $N=100$ block has RMS $2.2\cdot10^{-6}$.  This is computational evidence for local stabilization only; no monotone convergence theorem is claimed.

\begin{figure}[htbp]
\centering
\includegraphics[width=0.78\textwidth]{fig16_nested_bounded_endpoint.png}
\caption{Nested bounded endpoint sections.  Each curve shows the residual of one solved section on all permanent test blocks already introduced.  The ordinate shifts and free masses are bounded during the solve.  This is a finite computation and is not an existence proof for the infinite rung system.}
\label{fig:nested-bounded-v18}
\end{figure}

\begin{figure}[htbp]
\centering
\includegraphics[width=0.78\textwidth]{fig17_bounded_mass_stability.png}
\caption{Recovered free masses in the interior $n\le15$ of bounded endpoint sections.  The von Mangoldt values are shown only as a post-solution diagnostic and are not supplied to the solver; section-boundary masses are omitted from the plot.}
\label{fig:bounded-masses-v18}
\end{figure}

\medskip\noindent\textbf{Status.} The additive interpolation is exact but is a level-set reformulation of $\mathcal Q_{\RVM}$; the exact-atomic interpolation has no established interior path; the bounded finite-section computations do not prove an infinite solution.  The midpoint-clock injectivity problem \CGeqref{eq:clock-injectivity-v18} is a closed autonomous rigidity question, and Theorem~\CGref{thm:clock-jacobian} gives the accompanying multiscale linear information.  Proposition~\CGref{prop:rung-apriori-v18} gives an a priori displacement bound for any exact rung solution satisfying its hypotheses.  None of these statements supplies existence; the main endpoint identity (D) remains open.

\end{document}
